\documentclass[10pt,letterpaper,twoside]{book}
\usepackage[T1]{fontenc}
\usepackage[utf8]{inputenc}
\usepackage[margin=1.7cm,bottom=2.2cm,headsep=10pt]{geometry}
\usepackage{mathpazo}
\usepackage{amsmath,amssymb}
\usepackage{booktabs}
\usepackage{longtable}
\usepackage{array}
\usepackage{xcolor}
\usepackage[most]{tcolorbox}
\usepackage{graphicx}
\usepackage{fancyhdr}
\usepackage{enumitem}
\usepackage[hidelinks,pdfusetitle]{hyperref}
\usepackage{titlesec}

% --- angle macros used throughout the generated tables ---
\newcommand{\dg}{\ensuremath{{}^{\circ}}}
\newcommand{\pr}{\ensuremath{{}^{\prime}}}

\setlength{\tabcolsep}{4pt}
\renewcommand{\arraystretch}{1.05}
\emergencystretch=4em
\hyphenpenalty=1000

\pagestyle{fancy}
\fancyhf{}
\fancyhead[LE,RO]{\small\thepage}
\fancyhead[RE]{\small\itshape The Sextant Navigator's Almanac 2026}
\fancyhead[LO]{\small\itshape\leftmark}
\renewcommand{\headrulewidth}{0.3pt}

\titleformat{\chapter}[display]{\normalfont\huge\bfseries}{\chaptertitlename\ \thechapter}{16pt}{\Huge}

% inline data macros produced by generate_tables.wls
\newcommand{\exOneInt}{5.0}
\newcommand{\exOneDir}{Toward}
\newcommand{\exOneZn}{219\dg 19.2\pr}
\newcommand{\exOneHo}{69\dg 32.4\pr}
\newcommand{\exOneHc}{69\dg 27.4\pr}
\newcommand{\exOneLHA}{14\dg 00.0\pr}
\newcommand{\exTwoLat}{N40\dg 06.0\pr}
\newcommand{\exThreeHc}{17\dg 54.6\pr}
\newcommand{\exThreeZn}{122\dg 26.3\pr}
\newcommand{\genDate}{2026-07-02}


\title{The Sextant Navigator's Almanac \& Sight-Reduction Handbook (2026)}
\author{Generated from the open-source \textsc{SextantNavigation} ephemeris engine}
\date{Edition of \genDate}

\begin{document}
\frontmatter

% ============================ TITLE PAGE ============================
\begin{titlepage}
\centering
\vspace*{2.2cm}
{\Huge\bfseries The Sextant Navigator's\\[4pt] Almanac 2026\par}
\vspace{6pt}
{\Large\itshape \& Sight-Reduction Handbook\par}
\vspace{1.4cm}
{\large A self-contained nautical almanac and paper sight-reduction manual,\\
computed entirely from a validated, open-source celestial-navigation engine.\par}
\vfill
{\large Greenwich Hour Angle and Declination of the Sun,\\
GHA of Aries, and the 58 navigational stars, for the year 2026 (UT).\par}
\vspace{1.2cm}
{\normalsize Generated from \textsc{SextantNavigation}\\
\url{https://github.com/mthiel74/SextantNavigation}\par}
\vspace{0.6cm}
{\small Edition of \genDate\par}
\vspace{0.4cm}
{\footnotesize \textcopyright{} 2026 Marco Thiel. Licensed under Creative Commons
Attribution 4.0 International (CC BY 4.0).\par}
\vfill
{\footnotesize\bfseries FOR EDUCATION AND PRACTICE ONLY --- NOT FOR SOLE-MEANS NAVIGATION\par}
\end{titlepage}

% ============================ DISCLAIMER ============================
\thispagestyle{empty}
\vspace*{1.5cm}
\begin{tcolorbox}[colback=red!5!white,colframe=red!60!black,boxrule=1.2pt,
  arc=3pt,title=\bfseries IMPORTANT --- READ BEFORE USE,fonttitle=\large]
\large
This almanac and handbook are an \textbf{educational deliverable}. Every number
in it was computed offline from the open-source ephemeris engine at
\url{https://github.com/mthiel74/SextantNavigation}, not copied from an official
publication.

\medskip
\textbf{Accuracy.} The underlying engine agrees with JPL for the Sun's geographic
position to about \textbf{0.14$\pr$}, reproduces Meeus' lunar worked example to
0.002$\pr$ (geocentric), and precesses the star catalogue with IAU--1976
precession plus proper motion. Realistic end-to-end accuracy for a reduced sight
is on the order of \textbf{one arcminute (about one nautical mile)}. That is good
enough to \emph{learn and practise} celestial navigation, and to check your paper
arithmetic --- but it is \textbf{not} certified for life-safety use.

\medskip
\textbf{Do not use this document as your sole means of navigation at sea.} For
real voyages, use an official \emph{Nautical Almanac} (HM Nautical Almanac Office
/ US Naval Observatory) and cross-check every figure. The authors accept no
liability for any use of this material.
\end{tcolorbox}

\bigskip
\noindent{\small\textbf{Sources \& method.} The positions in this book are
\emph{computed with the Wolfram Language}, not reproduced from any published
almanac. The Sun's geographic
position is from the low-precision solar formulae of the \emph{Astronomical
Almanac}; the Moon from Meeus's abridged ELP theory (\emph{Astronomical
Algorithms}, ch.~47); Greenwich sidereal time from the IAU GMST formula; and the
star places are propagated from \emph{Hipparcos} catalogue (ESA) J2000 mean
positions using IAU--1976 precession and catalogue proper motions to epoch~2026.5.
Full source code and tests: \url{https://github.com/mthiel74/SextantNavigation}.}

\medskip
\noindent{\small\textbf{Acknowledgements.} All computations and tables were
produced with the \emph{Wolfram Language} (Wolfram Research). This almanac
gratefully builds on published astronomy: the star catalogue derives from the
ESA \emph{Hipparcos} mission; the lunar theory follows Jean Meeus,
\emph{Astronomical Algorithms}; and the solar and sidereal-time formulae follow
the conventions of the IAU and the \emph{Astronomical Almanac} (HM Nautical
Almanac Office / US Naval Observatory). It is an independent, open-source
companion to the article cited above, and is not affiliated with, or endorsed
by, any of these organisations.}

\vfill
\begin{center}\small
Source engine, tests, and this document: \url{https://github.com/mthiel74/SextantNavigation}\\
Companion article: \url{https://community.wolfram.com/groups/-/m/t/3743363}
\end{center}
\clearpage

% ============================ HOW TO USE ============================
\chapter*{How to use this almanac}
\addcontentsline{toc}{chapter}{How to use this almanac}
\markboth{How to use this almanac}{How to use this almanac}

All times in the data tables are \textbf{Universal Time (UT)}, the time kept by
your chronometer set to Greenwich. All angles are given in degrees and decimal
minutes to $0.1\pr$, the same precision as a commercial nautical almanac.

\paragraph{To reduce a sight you need four numbers for the instant of your
observation:} the body's \emph{Greenwich Hour Angle} (GHA) and \emph{Declination}
(Dec), your \emph{assumed position}, and your \emph{observed altitude} $H_o$.
This book supplies the first two:

\begin{itemize}[nosep]
\item \textbf{Sun.} Turn to Part~II, find the month, then the row for your day and
      the whole hour of UT. Read GHA and Dec directly. For the odd minutes and
      seconds, add the increment from Part~IV (the Sun's GHA grows a steady
      $15\dg$ per hour). Declination changes so slowly you may interpolate it by eye.
\item \textbf{A star.} Turn to Part~III. Read \textbf{GHA Aries} for your day and
      hour, then add the star's \textbf{SHA} (Sidereal Hour Angle) from the star
      table: $\text{GHA}_\star = \text{GHA Aries} + \text{SHA}$. The star's
      declination is taken straight from the table. Add the Part~IV increment for
      minutes and seconds (Aries runs at $15.041\dg$ per hour).
\end{itemize}

\paragraph{Then work the triangle.} Part~I explains, step by step and with fully
worked 2026 examples, how to correct your sextant reading and turn these numbers
into a line of position on your chart.

\tableofcontents

\mainmatter

% ============================================================
% PART I
% ============================================================
\part{Finding Yourself with a Sextant}
\chapter{The idea: circles of equal altitude}

When you measure the altitude of a celestial body --- the angle between it and
the horizon --- you learn one thing about where you are: your distance from the
body's \emph{geographic position} (GP), the point on the Earth's surface directly
beneath it. If the Sun's measured altitude is $H_o$, then your \emph{zenith
distance} is
\[
  z = 90\dg - H_o ,
\]
and $z$, expressed in arcminutes, is exactly your great-circle distance from the
GP in nautical miles (one minute of arc on the Earth $=$ one nautical mile). You
must therefore lie somewhere on a circle of radius $z$ centred on the GP: a
\textbf{circle of equal altitude}.

One sight gives a circle. A \emph{second} sight of another body (or the same body
later) gives a second circle, and two circles cross at two points. A \emph{third}
sight resolves which of the two intersections is really you. In practice we never
draw these enormous circles; near your estimated position each circle is
indistinguishable from its tangent straight line, the \textbf{line of position}
(LOP), and we plot those instead. That is the method of this chapter.

\section{What you need}
\begin{itemize}[nosep]
\item A \textbf{sextant} to measure altitude $H_s$ above the sea horizon.
\item An accurate \textbf{clock on UT} (a chronometer). An error of 4 seconds of
      time moves your fix about one nautical mile in longitude.
\item This \textbf{almanac}, for GHA and Dec of the body.
\item A \textbf{dead-reckoning (DR) position}: your best estimate of where you are,
      used to pick an ``assumed position'' close by.
\item Plotting sheet, dividers, parallel rule, and a pocket calculator with
      trigonometric functions (or the concise sight-reduction tables).
\end{itemize}

\chapter{Correcting the sight}

The raw sextant altitude $H_s$ is not yet the true altitude of the body's centre.
Apply, in order:

\begin{enumerate}[nosep]
\item \textbf{Index error} $I$. Your own sextant's zero error, found by sighting
      the horizon. Subtract if ``on the arc'', add if ``off the arc''.
\item \textbf{Dip} of the horizon. Because your eye is above the sea, the visible
      horizon is depressed below the true horizontal by
      \[ \text{Dip} = 1.76\sqrt{h}\ \text{arcminutes}, \qquad h=\text{height of eye in metres.} \]
      Always subtract. (See the table in Part~IV.)
\item \textbf{Refraction} $R$. The atmosphere bends the light down, raising the
      apparent altitude; subtract $R$. This almanac uses Bennett's formula,
      $R = \cot\!\big(H_a + 7.31/(H_a+4.4)\big)$ arcminutes, evaluated at the
      apparent altitude $H_a$ (altitude after index error and dip). Refraction is
      about $1\pr$ at $45\dg$ but $10\pr$ near the horizon.
\item \textbf{Semidiameter} (Sun and Moon). You bring the \emph{limb} (edge), not
      the centre, to the horizon. Add the semidiameter (SD) for a lower-limb sight,
      subtract it for an upper-limb sight. The Sun's SD is about $16\pr$; the exact
      monthly value is tabulated in Part~IV.
\item \textbf{Parallax} (Moon, and slightly the Sun). A near body is seen against
      the stars from a slightly different direction than it would be from the
      Earth's centre. The correction, always \emph{added}, is
      \[ P = \arcsin\!\big(\sin(\mathrm{HP})\cos H\big) \approx \mathrm{HP}\cos H, \]
      where HP is the horizontal parallax. For the Moon HP is about $57\pr$, so the
      parallax reaches nearly a full degree; for the Sun it is a negligible
      $0.1\pr$. Stars are effectively infinitely far away: no parallax, no
      semidiameter.
\end{enumerate}

The result is the \textbf{observed altitude}
\[
  H_o = H_s - I - \text{Dip} - R \;(\pm\,\text{SD})\;(+\,P).
\]

\chapter{Sight reduction: the Marcq St-Hilaire intercept}

We cannot solve for position directly, so we \emph{guess} an assumed position (AP)
close to our DR and ask: \emph{if I were at the AP, what altitude and bearing
would the body have?} Comparing that computed altitude $H_c$ with the observed
$H_o$ tells us how far, and in which direction, to shift from the AP onto our true
LOP.

\section{The navigational (PZX) triangle}
On the celestial sphere, three points form a spherical triangle: the elevated
pole $P$, your zenith $Z$, and the body $X$. Its sides are the co-latitude
$90\dg-\varphi$, the co-declination $90\dg-\delta$, and the co-altitude
$90\dg-H_c$; the angle at the pole is the \textbf{local hour angle}
$\text{LHA}=\text{GHA}-\text{(west longitude)}$ (equivalently $\text{GHA}+$ east
longitude). Applying the spherical law of cosines to the side $90\dg-H_c$:
\begin{equation}
  \boxed{\;\sin H_c = \sin\varphi\,\sin\delta + \cos\varphi\,\cos\delta\,\cos(\text{LHA})\;}
\end{equation}
where $\varphi$ is the assumed latitude and $\delta$ the declination. The body's
\textbf{azimuth} $Z_n$ (true bearing, measured clockwise from north) comes from
\begin{equation}
  \cos Z = \frac{\sin\delta - \sin H_c\,\sin\varphi}{\cos H_c\,\cos\varphi},
  \qquad
  Z_n = \begin{cases} Z & \text{if } \sin(\text{LHA})<0 \ (\text{LHA }180\dg\text{--}360\dg),\\
                      360\dg - Z & \text{if } \sin(\text{LHA})>0 \ (\text{LHA }0\dg\text{--}180\dg).\end{cases}
\end{equation}

\section{The intercept, and how to plot it}
The heart of the method is a single subtraction. Because one minute of altitude
equals one nautical mile of distance from the GP,
\begin{equation}
  \boxed{\; a = 60\,(H_o - H_c)\ \text{nautical miles}\;}
\end{equation}
(with $H_o,H_c$ in degrees). If $H_o>H_c$ you are \emph{closer} to the body than
the AP is, so step the intercept \textbf{Toward} the body along the azimuth $Z_n$;
if $H_o<H_c$, step it \textbf{Away} (along $Z_n+180\dg$). A useful mnemonic is
\emph{``Ho Mo To''} --- $H_o$ More means Toward.

\textbf{To plot:} mark the AP on your sheet. Draw the azimuth line through it in
the direction $Z_n$. Measure $a$ nautical miles along that line, Toward or Away.
Through that point draw the LOP \emph{perpendicular} to the azimuth --- that
perpendicular is your line of position. Repeat for a second body; where the two
LOPs cross is your fix. A third LOP turns the small ``cocked hat'' triangle into a
confident position.

\section{Choosing the assumed position}
To keep the arithmetic clean, navigators choose the AP cleverly: take the
\emph{assumed latitude} as the whole degree of latitude nearest your DR, and the
\emph{assumed longitude} with whatever minutes make the LHA come out to a whole
number of degrees. Both examples below use this trick.

\chapter{Worked example 1: a Sun sight}
\label{ch:ex1}

At sea near $40\dg$N, $30\dg$W on 2026 June~29, you take a lower-limb Sun sight.
Your chronometer reads exactly \textbf{15:00:00 UT}. Height of eye is $3.0$ m and
your sextant's index error is $2.0\pr$ on the arc. Work as follows (all almanac
figures are read from Parts~II and~IV of this book).

\begin{center}\begin{tabular}{l r}
\toprule
\multicolumn{2}{l}{\textbf{Sight and corrections}}\\
\midrule
Body & Sun (lower limb)\\
UT & 2026 June 29, 15:00:00\\
DR position & 40\dg N, 30\dg W (approx.)\\
Sextant altitude $H_s$ & 69\dg 22.1\pr\\
Index error (on the arc) & $-$2.0\pr\\
Dip (height of eye 3.0 m) & $-$3.0\pr\\
Apparent altitude $H_a$ & 69\dg 17.1\pr\\
Refraction & $-$0.4\pr\\
Semidiameter (lower limb) & $+$15.7\pr\\
\textbf{Observed altitude} $H_o$ & \textbf{69\dg 32.4\pr}\\
\midrule
\multicolumn{2}{l}{\textbf{From the almanac (15:00 UT, 29 June)}}\\
\midrule
GHA Sun & 44\dg 06.9\pr\\
Declination & N23\dg 12.0\pr\\
Assumed longitude (W) & 30\dg 06.9\pr\\
Local hour angle LHA & 14\dg 00.0\pr\\
Assumed latitude & N40\dg 00.0\pr\\
\midrule
\multicolumn{2}{l}{\textbf{Sight reduction}}\\
\midrule
Computed altitude $H_c$ & 69\dg 27.4\pr\\
Azimuth $Z_n$ & 219\dg 19.2\pr\\
Intercept $a=60(H_o-H_c)$ & 5.0 nm Toward\\
\bottomrule
\end{tabular}
\end{center}

\noindent\textbf{Reading the result.} The observed altitude exceeds the computed
altitude by $H_o-H_c=\exOneInt\pr$, so the intercept is \textbf{\exOneInt\ nautical
miles \exOneDir} the Sun. Plot the assumed position ($40\dg$N, $30\dg06.9\pr$W),
draw the azimuth line toward $Z_n=\exOneZn$, step \exOneInt~nm along it \exOneDir\
the Sun, and draw your line of position at right angles to that azimuth. A single
Sun sight gives one LOP; a second sight a few hours later (when the Sun's azimuth
has swung round) crosses it to give a running fix.

\emph{Check it yourself:} $\sin H_c = \sin40\dg\sin(23\dg12.0\pr) +
\cos40\dg\cos(23\dg12.0\pr)\cos(14\dg)$ gives $H_c=\exOneHc$, exactly the value in
the workform. The whole-degree LHA of $\exOneLHA$ is what the assumed-longitude
choice was designed to produce.

\chapter{Worked example 2: latitude by the noon Sun}

The simplest sight of all needs no azimuth and almost no arithmetic. When the Sun
crosses your meridian (\emph{local apparent noon}, LAN) it bears due north or due
south and reaches its greatest altitude. At that instant your latitude follows
from the meridian altitude alone.

At $30\dg$W the Sun transits your meridian when its GHA equals $30\dg$; from
Part~II that occurs on 2026 June~29 at the UT shown below. Measure the Sun's
maximum altitude, correct it to $H_o$ as before, and take the zenith distance
$z=90\dg-H_o$. Because the Sun is south of you (its declination $23\dg$N is less
than your latitude), your latitude is
\[
  \varphi = \delta + z \qquad (\text{body to the south}).
\]

\begin{center}\begin{tabular}{l r}
\toprule
Observer (near) & 40\dg N, 30\dg W\\
Date & 2026 June 29\\
Local apparent noon (UT) & 14:03.5\\
\midrule
Observed mer. altitude $H_o$ & 73\dg 06.1\pr\\
Zenith distance $z=90\dg-H_o$ & 16\dg 53.9\pr\\
Sun's declination (from almanac) & N23\dg 12.1\pr\\
Sun bears & South\\
\textbf{Latitude} $=$ Dec $+ z$ & \textbf{N40\dg 06.0\pr}\\
\bottomrule
\end{tabular}
\end{center}

\noindent The result, $\varphi=\exTwoLat$, agrees with the true position to a
tenth of a minute. (Had the Sun been \emph{north} of you, the rule would be
$\varphi=\delta-z$; if you straddle the sub-solar latitude,
$\varphi=z-\delta$ with the Sun to the south of an overhead-passing Sun. Always
sketch the meridian to get the sign right.)

\chapter{Worked example 3: a star sight}

Stars are reduced exactly like the Sun, but their GHA is built in two pieces. On
2026 January~15 at \textbf{22:00:00 UT}, from an assumed position $30\dg$N, you
shoot the star \textbf{Sirius}. Read GHA Aries from Part~III and Sirius's SHA and
Dec from the star table, then add:

\begin{center}\begin{tabular}{l r}
\toprule
Star & Sirius\\
UT & 2026 January 15, 22:00:00\\
Assumed position & 30\dg N, 40\dg 47.1\pr W\\
\midrule
GHA Aries (Part III) & 85\dg 21.8\pr\\
SHA Sirius (star table) & 258\dg 25.3\pr\\
GHA Sirius $=$ GHA$\gamma+$SHA & 343\dg 47.1\pr\\
Declination Sirius & S16\dg 45.3\pr\\
LHA $=$ GHA $-$ Long(W) & 303\dg 00.0\pr\\
\midrule
Computed altitude $H_c$ & 17\dg 54.6\pr\\
Azimuth $Z_n$ & 122\dg 26.3\pr\\
\bottomrule
\end{tabular}
\end{center}

\noindent Because a star has no semidiameter and no parallax, its sight correction
is only index error, dip, and refraction --- simpler than the Sun. The computed
altitude $H_c=\exThreeHc$ and azimuth $Z_n=\exThreeZn$ are then compared with your
observed $H_o$ exactly as in Example~1 to get the intercept.

\chapter{A note on longitude and time}

Latitude falls out of a noon sight almost for free, but \textbf{longitude is a
question of time}. The Earth turns $15\dg$ of longitude per hour, so every second
of chronometer error displaces your computed longitude by a quarter of a minute of
arc ($\approx 0.25\pr\cos\varphi$ nautical miles). This is why the marine
chronometer was one of history's great navigational prizes: to find longitude you
must know GMT precisely and compare it with your \emph{local} time as read from the
Sun or stars. In the intercept method above, the chronometer time enters through
the GHA you look up; an error in time is indistinguishable from an error in
longitude. Keep your clock honest, log its rate, and never reduce a sight without
the exact UT of the observation.

% ============================================================
% PART II  --- THE SUN
% ============================================================
\part{The Sun, 2026}
\chapter{Daily data: meridian passage, equation of time, semidiameter}

For each day of 2026 the table below gives the \textbf{meridian passage} of the
Sun at Greenwich (the UT at which the Sun's GHA is $0\dg$, i.e.\ it crosses the
Greenwich meridian), the \textbf{equation of time} (EoT $= 12^{\text{h}} -$
Mer.Pass.\ $=$ apparent $-$ mean solar time; a positive value means the sundial is
ahead of the clock), and the Sun's \textbf{semidiameter}. Note the semidiameter is
largest in early January (Earth at perihelion) and smallest in early July.

{\setlength{\tabcolsep}{3pt}{\footnotesize
\begin{longtable}{r r r r | r r r r}
\caption{Sun --- daily Meridian Passage at Greenwich, Equation of Time, and Semidiameter, 2026}\label{sundaily}\\
\toprule
\multicolumn{4}{c|}{} & \multicolumn{4}{c}{} \\
\textbf{Date} & \textbf{Mer.Pass} & \textbf{EoT} & \textbf{SD} & \textbf{Date} & \textbf{Mer.Pass} & \textbf{EoT} & \textbf{SD} \\
 & (UT) & (m:s) & ($\pr$) & & (UT) & (m:s) & ($\pr$) \\
\midrule
\endfirsthead
\multicolumn{8}{c}{\textit{Daily Sun data, 2026 (continued)}}\\
\toprule
\textbf{Date} & \textbf{Mer.Pass} & \textbf{EoT} & \textbf{SD} & \textbf{Date} & \textbf{Mer.Pass} & \textbf{EoT} & \textbf{SD} \\
\midrule
\endhead
\midrule\multicolumn{8}{r}{\textit{continued}}\\
\endfoot
\bottomrule
\endlastfoot
\multicolumn{8}{l}{\textbf{January 2026}}\\
Jan 1 & 12:03.6 & $-$3:34 & 16.3 & Jan 17 & 12:10.1 & $-$10:05 & 16.3 \\
Jan 2 & 12:04.0 & $-$4:02 & 16.3 & Jan 18 & 12:10.4 & $-$10:24 & 16.3 \\
Jan 3 & 12:04.5 & $-$4:30 & 16.3 & Jan 19 & 12:10.7 & $-$10:43 & 16.3 \\
Jan 4 & 12:05.0 & $-$4:57 & 16.3 & Jan 20 & 12:11.0 & $-$11:01 & 16.3 \\
Jan 5 & 12:05.4 & $-$5:24 & 16.3 & Jan 21 & 12:11.3 & $-$11:18 & 16.3 \\
Jan 6 & 12:05.8 & $-$5:51 & 16.3 & Jan 22 & 12:11.6 & $-$11:34 & 16.3 \\
Jan 7 & 12:06.3 & $-$6:17 & 16.3 & Jan 23 & 12:11.8 & $-$11:50 & 16.3 \\
Jan 8 & 12:06.7 & $-$6:42 & 16.3 & Jan 24 & 12:12.1 & $-$12:04 & 16.3 \\
Jan 9 & 12:07.1 & $-$7:07 & 16.3 & Jan 25 & 12:12.3 & $-$12:18 & 16.3 \\
Jan 10 & 12:07.5 & $-$7:31 & 16.3 & Jan 26 & 12:12.5 & $-$12:31 & 16.3 \\
Jan 11 & 12:07.9 & $-$7:55 & 16.3 & Jan 27 & 12:12.7 & $-$12:44 & 16.2 \\
Jan 12 & 12:08.3 & $-$8:19 & 16.3 & Jan 28 & 12:12.9 & $-$12:55 & 16.2 \\
Jan 13 & 12:08.7 & $-$8:41 & 16.3 & Jan 29 & 12:13.1 & $-$13:06 & 16.2 \\
Jan 14 & 12:09.1 & $-$9:03 & 16.3 & Jan 30 & 12:13.3 & $-$13:16 & 16.2 \\
Jan 15 & 12:09.4 & $-$9:24 & 16.3 & Jan 31 & 12:13.4 & $-$13:25 & 16.2 \\
Jan 16 & 12:09.8 & $-$9:45 & 16.3 &  & & &  \\
\multicolumn{8}{l}{\textbf{February 2026}}\\
Feb 1 & 12:13.6 & $-$13:33 & 16.2 & Feb 15 & 12:14.1 & $-$14:05 & 16.2 \\
Feb 2 & 12:13.7 & $-$13:41 & 16.2 & Feb 16 & 12:14.0 & $-$14:02 & 16.2 \\
Feb 3 & 12:13.8 & $-$13:47 & 16.2 & Feb 17 & 12:14.0 & $-$13:58 & 16.2 \\
Feb 4 & 12:13.9 & $-$13:53 & 16.2 & Feb 18 & 12:13.9 & $-$13:53 & 16.2 \\
Feb 5 & 12:14.0 & $-$13:58 & 16.2 & Feb 19 & 12:13.8 & $-$13:48 & 16.2 \\
Feb 6 & 12:14.0 & $-$14:03 & 16.2 & Feb 20 & 12:13.7 & $-$13:41 & 16.2 \\
Feb 7 & 12:14.1 & $-$14:06 & 16.2 & Feb 21 & 12:13.6 & $-$13:35 & 16.2 \\
Feb 8 & 12:14.1 & $-$14:09 & 16.2 & Feb 22 & 12:13.5 & $-$13:27 & 16.2 \\
Feb 9 & 12:14.2 & $-$14:10 & 16.2 & Feb 23 & 12:13.3 & $-$13:19 & 16.2 \\
Feb 10 & 12:14.2 & $-$14:12 & 16.2 & Feb 24 & 12:13.2 & $-$13:10 & 16.2 \\
Feb 11 & 12:14.2 & $-$14:12 & 16.2 & Feb 25 & 12:13.0 & $-$13:01 & 16.2 \\
Feb 12 & 12:14.2 & $-$14:11 & 16.2 & Feb 26 & 12:12.9 & $-$12:51 & 16.2 \\
Feb 13 & 12:14.2 & $-$14:10 & 16.2 & Feb 27 & 12:12.7 & $-$12:41 & 16.2 \\
Feb 14 & 12:14.1 & $-$14:08 & 16.2 & Feb 28 & 12:12.5 & $-$12:30 & 16.2 \\
\multicolumn{8}{l}{\textbf{March 2026}}\\
Mar 1 & 12:12.3 & $-$12:19 & 16.1 & Mar 17 & 12:08.3 & $-$8:20 & 16.1 \\
Mar 2 & 12:12.1 & $-$12:07 & 16.1 & Mar 18 & 12:08.0 & $-$8:03 & 16.1 \\
Mar 3 & 12:11.9 & $-$11:54 & 16.1 & Mar 19 & 12:07.8 & $-$7:45 & 16.1 \\
Mar 4 & 12:11.7 & $-$11:41 & 16.1 & Mar 20 & 12:07.5 & $-$7:28 & 16.1 \\
Mar 5 & 12:11.5 & $-$11:28 & 16.1 & Mar 21 & 12:07.2 & $-$7:10 & 16.1 \\
Mar 6 & 12:11.2 & $-$11:14 & 16.1 & Mar 22 & 12:06.9 & $-$6:52 & 16.1 \\
Mar 7 & 12:11.0 & $-$11:00 & 16.1 & Mar 23 & 12:06.6 & $-$6:34 & 16.1 \\
Mar 8 & 12:10.8 & $-$10:45 & 16.1 & Mar 24 & 12:06.3 & $-$6:16 & 16.0 \\
Mar 9 & 12:10.5 & $-$10:30 & 16.1 & Mar 25 & 12:06.0 & $-$5:58 & 16.0 \\
Mar 10 & 12:10.3 & $-$10:15 & 16.1 & Mar 26 & 12:05.7 & $-$5:40 & 16.0 \\
Mar 11 & 12:10.0 & $-$10:00 & 16.1 & Mar 27 & 12:05.4 & $-$5:22 & 16.0 \\
Mar 12 & 12:09.7 & $-$9:44 & 16.1 & Mar 28 & 12:05.1 & $-$5:04 & 16.0 \\
Mar 13 & 12:09.5 & $-$9:27 & 16.1 & Mar 29 & 12:04.8 & $-$4:46 & 16.0 \\
Mar 14 & 12:09.2 & $-$9:11 & 16.1 & Mar 30 & 12:04.5 & $-$4:28 & 16.0 \\
Mar 15 & 12:08.9 & $-$8:54 & 16.1 & Mar 31 & 12:04.2 & $-$4:10 & 16.0 \\
Mar 16 & 12:08.6 & $-$8:37 & 16.1 &  & & &  \\
\multicolumn{8}{l}{\textbf{April 2026}}\\
Apr 1 & 12:03.9 & $-$3:52 & 16.0 & Apr 16 & 11:59.8 & +0:12 & 15.9 \\
Apr 2 & 12:03.6 & $-$3:34 & 16.0 & Apr 17 & 11:59.6 & +0:26 & 15.9 \\
Apr 3 & 12:03.3 & $-$3:17 & 16.0 & Apr 18 & 11:59.3 & +0:39 & 15.9 \\
Apr 4 & 12:03.0 & $-$2:59 & 16.0 & Apr 19 & 11:59.1 & +0:52 & 15.9 \\
Apr 5 & 12:02.7 & $-$2:42 & 16.0 & Apr 20 & 11:58.9 & +1:05 & 15.9 \\
Apr 6 & 12:02.4 & $-$2:25 & 16.0 & Apr 21 & 11:58.7 & +1:17 & 15.9 \\
Apr 7 & 12:02.1 & $-$2:08 & 16.0 & Apr 22 & 11:58.5 & +1:29 & 15.9 \\
Apr 8 & 12:01.9 & $-$1:51 & 16.0 & Apr 23 & 11:58.3 & +1:40 & 15.9 \\
Apr 9 & 12:01.6 & $-$1:35 & 16.0 & Apr 24 & 11:58.1 & +1:51 & 15.9 \\
Apr 10 & 12:01.3 & $-$1:19 & 16.0 & Apr 25 & 11:58.0 & +2:01 & 15.9 \\
Apr 11 & 12:01.0 & $-$1:03 & 16.0 & Apr 26 & 11:57.8 & +2:11 & 15.9 \\
Apr 12 & 12:00.8 & $-$0:47 & 16.0 & Apr 27 & 11:57.7 & +2:21 & 15.9 \\
Apr 13 & 12:00.5 & $-$0:32 & 16.0 & Apr 28 & 11:57.5 & +2:30 & 15.9 \\
Apr 14 & 12:00.3 & $-$0:17 & 16.0 & Apr 29 & 11:57.4 & +2:38 & 15.9 \\
Apr 15 & 12:00.0 & $-$0:03 & 15.9 & Apr 30 & 11:57.2 & +2:46 & 15.9 \\
\multicolumn{8}{l}{\textbf{May 2026}}\\
May 1 & 11:57.1 & +2:53 & 15.9 & May 17 & 11:56.4 & +3:35 & 15.8 \\
May 2 & 11:57.0 & +3:00 & 15.9 & May 18 & 11:56.5 & +3:33 & 15.8 \\
May 3 & 11:56.9 & +3:06 & 15.9 & May 19 & 11:56.5 & +3:30 & 15.8 \\
May 4 & 11:56.8 & +3:12 & 15.9 & May 20 & 11:56.6 & +3:27 & 15.8 \\
May 5 & 11:56.7 & +3:17 & 15.9 & May 21 & 11:56.6 & +3:23 & 15.8 \\
May 6 & 11:56.6 & +3:22 & 15.9 & May 22 & 11:56.7 & +3:19 & 15.8 \\
May 7 & 11:56.6 & +3:26 & 15.9 & May 23 & 11:56.8 & +3:14 & 15.8 \\
May 8 & 11:56.5 & +3:29 & 15.9 & May 24 & 11:56.9 & +3:08 & 15.8 \\
May 9 & 11:56.5 & +3:32 & 15.8 & May 25 & 11:57.0 & +3:02 & 15.8 \\
May 10 & 11:56.4 & +3:35 & 15.8 & May 26 & 11:57.1 & +2:56 & 15.8 \\
May 11 & 11:56.4 & +3:36 & 15.8 & May 27 & 11:57.2 & +2:49 & 15.8 \\
May 12 & 11:56.4 & +3:38 & 15.8 & May 28 & 11:57.3 & +2:42 & 15.8 \\
May 13 & 11:56.4 & +3:38 & 15.8 & May 29 & 11:57.4 & +2:34 & 15.8 \\
May 14 & 11:56.4 & +3:38 & 15.8 & May 30 & 11:57.6 & +2:26 & 15.8 \\
May 15 & 11:56.4 & +3:38 & 15.8 & May 31 & 11:57.7 & +2:17 & 15.8 \\
May 16 & 11:56.4 & +3:37 & 15.8 &  & & &  \\
\multicolumn{8}{l}{\textbf{June 2026}}\\
Jun 1 & 11:57.9 & +2:08 & 15.8 & Jun 16 & 12:00.7 & $-$0:44 & 15.8 \\
Jun 2 & 11:58.0 & +1:59 & 15.8 & Jun 17 & 12:01.0 & $-$0:57 & 15.7 \\
Jun 3 & 11:58.2 & +1:49 & 15.8 & Jun 18 & 12:01.2 & $-$1:10 & 15.7 \\
Jun 4 & 11:58.4 & +1:39 & 15.8 & Jun 19 & 12:01.4 & $-$1:23 & 15.7 \\
Jun 5 & 11:58.5 & +1:28 & 15.8 & Jun 20 & 12:01.6 & $-$1:37 & 15.7 \\
Jun 6 & 11:58.7 & +1:17 & 15.8 & Jun 21 & 12:01.8 & $-$1:50 & 15.7 \\
Jun 7 & 11:58.9 & +1:06 & 15.8 & Jun 22 & 12:02.0 & $-$2:03 & 15.7 \\
Jun 8 & 11:59.1 & +0:55 & 15.8 & Jun 23 & 12:02.3 & $-$2:16 & 15.7 \\
Jun 9 & 11:59.3 & +0:43 & 15.8 & Jun 24 & 12:02.5 & $-$2:28 & 15.7 \\
Jun 10 & 11:59.5 & +0:31 & 15.8 & Jun 25 & 12:02.7 & $-$2:41 & 15.7 \\
Jun 11 & 11:59.7 & +0:19 & 15.8 & Jun 26 & 12:02.9 & $-$2:54 & 15.7 \\
Jun 12 & 11:59.9 & +0:07 & 15.8 & Jun 27 & 12:03.1 & $-$3:06 & 15.7 \\
Jun 13 & 12:00.1 & $-$0:06 & 15.8 & Jun 28 & 12:03.3 & $-$3:19 & 15.7 \\
Jun 14 & 12:00.3 & $-$0:19 & 15.8 & Jun 29 & 12:03.5 & $-$3:31 & 15.7 \\
Jun 15 & 12:00.5 & $-$0:31 & 15.8 & Jun 30 & 12:03.7 & $-$3:43 & 15.7 \\
\multicolumn{8}{l}{\textbf{July 2026}}\\
Jul 1 & 12:03.9 & $-$3:54 & 15.7 & Jul 17 & 12:06.2 & $-$6:12 & 15.7 \\
Jul 2 & 12:04.1 & $-$4:06 & 15.7 & Jul 18 & 12:06.3 & $-$6:17 & 15.7 \\
Jul 3 & 12:04.3 & $-$4:17 & 15.7 & Jul 19 & 12:06.3 & $-$6:21 & 15.7 \\
Jul 4 & 12:04.5 & $-$4:28 & 15.7 & Jul 20 & 12:06.4 & $-$6:24 & 15.7 \\
Jul 5 & 12:04.6 & $-$4:38 & 15.7 & Jul 21 & 12:06.5 & $-$6:27 & 15.7 \\
Jul 6 & 12:04.8 & $-$4:48 & 15.7 & Jul 22 & 12:06.5 & $-$6:30 & 15.7 \\
Jul 7 & 12:05.0 & $-$4:58 & 15.7 & Jul 23 & 12:06.5 & $-$6:32 & 15.7 \\
Jul 8 & 12:05.1 & $-$5:07 & 15.7 & Jul 24 & 12:06.5 & $-$6:33 & 15.8 \\
Jul 9 & 12:05.3 & $-$5:16 & 15.7 & Jul 25 & 12:06.6 & $-$6:34 & 15.8 \\
Jul 10 & 12:05.4 & $-$5:25 & 15.7 & Jul 26 & 12:06.6 & $-$6:34 & 15.8 \\
Jul 11 & 12:05.6 & $-$5:33 & 15.7 & Jul 27 & 12:06.6 & $-$6:33 & 15.8 \\
Jul 12 & 12:05.7 & $-$5:41 & 15.7 & Jul 28 & 12:06.5 & $-$6:32 & 15.8 \\
Jul 13 & 12:05.8 & $-$5:48 & 15.7 & Jul 29 & 12:06.5 & $-$6:30 & 15.8 \\
Jul 14 & 12:05.9 & $-$5:55 & 15.7 & Jul 30 & 12:06.5 & $-$6:28 & 15.8 \\
Jul 15 & 12:06.0 & $-$6:01 & 15.7 & Jul 31 & 12:06.4 & $-$6:25 & 15.8 \\
Jul 16 & 12:06.1 & $-$6:07 & 15.7 &  & & &  \\
\multicolumn{8}{l}{\textbf{August 2026}}\\
Aug 1 & 12:06.4 & $-$6:22 & 15.8 & Aug 17 & 12:04.1 & $-$4:05 & 15.8 \\
Aug 2 & 12:06.3 & $-$6:17 & 15.8 & Aug 18 & 12:03.9 & $-$3:51 & 15.8 \\
Aug 3 & 12:06.2 & $-$6:13 & 15.8 & Aug 19 & 12:03.6 & $-$3:38 & 15.8 \\
Aug 4 & 12:06.1 & $-$6:07 & 15.8 & Aug 20 & 12:03.4 & $-$3:24 & 15.8 \\
Aug 5 & 12:06.0 & $-$6:01 & 15.8 & Aug 21 & 12:03.1 & $-$3:09 & 15.8 \\
Aug 6 & 12:05.9 & $-$5:55 & 15.8 & Aug 22 & 12:02.9 & $-$2:54 & 15.8 \\
Aug 7 & 12:05.8 & $-$5:48 & 15.8 & Aug 23 & 12:02.6 & $-$2:38 & 15.8 \\
Aug 8 & 12:05.7 & $-$5:40 & 15.8 & Aug 24 & 12:02.4 & $-$2:22 & 15.8 \\
Aug 9 & 12:05.5 & $-$5:32 & 15.8 & Aug 25 & 12:02.1 & $-$2:06 & 15.8 \\
Aug 10 & 12:05.4 & $-$5:23 & 15.8 & Aug 26 & 12:01.8 & $-$1:49 & 15.8 \\
Aug 11 & 12:05.2 & $-$5:13 & 15.8 & Aug 27 & 12:01.5 & $-$1:32 & 15.8 \\
Aug 12 & 12:05.1 & $-$5:03 & 15.8 & Aug 28 & 12:01.2 & $-$1:14 & 15.8 \\
Aug 13 & 12:04.9 & $-$4:53 & 15.8 & Aug 29 & 12:00.9 & $-$0:56 & 15.8 \\
Aug 14 & 12:04.7 & $-$4:41 & 15.8 & Aug 30 & 12:00.6 & $-$0:38 & 15.8 \\
Aug 15 & 12:04.5 & $-$4:30 & 15.8 & Aug 31 & 12:00.3 & $-$0:19 & 15.8 \\
Aug 16 & 12:04.3 & $-$4:17 & 15.8 &  & & &  \\
\multicolumn{8}{l}{\textbf{September 2026}}\\
Sep 1 & 12:00.0 & $-$0:00 & 15.9 & Sep 16 & 11:54.8 & +5:09 & 15.9 \\
Sep 2 & 11:59.7 & +0:19 & 15.9 & Sep 17 & 11:54.5 & +5:31 & 15.9 \\
Sep 3 & 11:59.4 & +0:39 & 15.9 & Sep 18 & 11:54.1 & +5:52 & 15.9 \\
Sep 4 & 11:59.0 & +0:59 & 15.9 & Sep 19 & 11:53.8 & +6:14 & 15.9 \\
Sep 5 & 11:58.7 & +1:19 & 15.9 & Sep 20 & 11:53.4 & +6:35 & 15.9 \\
Sep 6 & 11:58.4 & +1:39 & 15.9 & Sep 21 & 11:53.1 & +6:56 & 15.9 \\
Sep 7 & 11:58.0 & +1:59 & 15.9 & Sep 22 & 11:52.7 & +7:18 & 15.9 \\
Sep 8 & 11:57.7 & +2:20 & 15.9 & Sep 23 & 11:52.4 & +7:39 & 15.9 \\
Sep 9 & 11:57.3 & +2:41 & 15.9 & Sep 24 & 11:52.0 & +8:00 & 15.9 \\
Sep 10 & 11:57.0 & +3:02 & 15.9 & Sep 25 & 11:51.7 & +8:20 & 16.0 \\
Sep 11 & 11:56.6 & +3:23 & 15.9 & Sep 26 & 11:51.3 & +8:41 & 16.0 \\
Sep 12 & 11:56.3 & +3:44 & 15.9 & Sep 27 & 11:51.0 & +9:02 & 16.0 \\
Sep 13 & 11:55.9 & +4:05 & 15.9 & Sep 28 & 11:50.6 & +9:22 & 16.0 \\
Sep 14 & 11:55.6 & +4:27 & 15.9 & Sep 29 & 11:50.3 & +9:42 & 16.0 \\
Sep 15 & 11:55.2 & +4:48 & 15.9 & Sep 30 & 11:50.0 & +10:02 & 16.0 \\
\multicolumn{8}{l}{\textbf{October 2026}}\\
Oct 1 & 11:49.6 & +10:21 & 16.0 & Oct 17 & 11:45.3 & +14:39 & 16.1 \\
Oct 2 & 11:49.3 & +10:40 & 16.0 & Oct 18 & 11:45.1 & +14:51 & 16.1 \\
Oct 3 & 11:49.0 & +10:59 & 16.0 & Oct 19 & 11:45.0 & +15:03 & 16.1 \\
Oct 4 & 11:48.7 & +11:18 & 16.0 & Oct 20 & 11:44.8 & +15:13 & 16.1 \\
Oct 5 & 11:48.4 & +11:36 & 16.0 & Oct 21 & 11:44.6 & +15:23 & 16.1 \\
Oct 6 & 11:48.1 & +11:54 & 16.0 & Oct 22 & 11:44.5 & +15:32 & 16.1 \\
Oct 7 & 11:47.8 & +12:11 & 16.0 & Oct 23 & 11:44.3 & +15:41 & 16.1 \\
Oct 8 & 11:47.5 & +12:28 & 16.0 & Oct 24 & 11:44.2 & +15:49 & 16.1 \\
Oct 9 & 11:47.3 & +12:45 & 16.0 & Oct 25 & 11:44.1 & +15:56 & 16.1 \\
Oct 10 & 11:47.0 & +13:01 & 16.0 & Oct 26 & 11:44.0 & +16:03 & 16.1 \\
Oct 11 & 11:46.7 & +13:16 & 16.0 & Oct 27 & 11:43.9 & +16:08 & 16.1 \\
Oct 12 & 11:46.5 & +13:32 & 16.0 & Oct 28 & 11:43.8 & +16:13 & 16.1 \\
Oct 13 & 11:46.2 & +13:46 & 16.0 & Oct 29 & 11:43.7 & +16:17 & 16.1 \\
Oct 14 & 11:46.0 & +14:00 & 16.0 & Oct 30 & 11:43.7 & +16:21 & 16.1 \\
Oct 15 & 11:45.8 & +14:14 & 16.0 & Oct 31 & 11:43.6 & +16:24 & 16.1 \\
Oct 16 & 11:45.5 & +14:27 & 16.0 &  & & &  \\
\multicolumn{8}{l}{\textbf{November 2026}}\\
Nov 1 & 11:43.6 & +16:26 & 16.1 & Nov 16 & 11:44.8 & +15:15 & 16.2 \\
Nov 2 & 11:43.6 & +16:27 & 16.1 & Nov 17 & 11:44.9 & +15:03 & 16.2 \\
Nov 3 & 11:43.6 & +16:27 & 16.1 & Nov 18 & 11:45.2 & +14:51 & 16.2 \\
Nov 4 & 11:43.6 & +16:26 & 16.1 & Nov 19 & 11:45.4 & +14:38 & 16.2 \\
Nov 5 & 11:43.6 & +16:25 & 16.1 & Nov 20 & 11:45.6 & +14:24 & 16.2 \\
Nov 6 & 11:43.6 & +16:23 & 16.1 & Nov 21 & 11:45.9 & +14:09 & 16.2 \\
Nov 7 & 11:43.7 & +16:20 & 16.1 & Nov 22 & 11:46.1 & +13:53 & 16.2 \\
Nov 8 & 11:43.7 & +16:16 & 16.1 & Nov 23 & 11:46.4 & +13:37 & 16.2 \\
Nov 9 & 11:43.8 & +16:11 & 16.2 & Nov 24 & 11:46.7 & +13:20 & 16.2 \\
Nov 10 & 11:43.9 & +16:06 & 16.2 & Nov 25 & 11:47.0 & +13:02 & 16.2 \\
Nov 11 & 11:44.0 & +15:59 & 16.2 & Nov 26 & 11:47.3 & +12:43 & 16.2 \\
Nov 12 & 11:44.1 & +15:52 & 16.2 & Nov 27 & 11:47.6 & +12:24 & 16.2 \\
Nov 13 & 11:44.3 & +15:44 & 16.2 & Nov 28 & 11:47.9 & +12:04 & 16.2 \\
Nov 14 & 11:44.4 & +15:35 & 16.2 & Nov 29 & 11:48.3 & +11:43 & 16.2 \\
Nov 15 & 11:44.6 & +15:25 & 16.2 & Nov 30 & 11:48.6 & +11:22 & 16.2 \\
\multicolumn{8}{l}{\textbf{December 2026}}\\
Dec 1 & 11:49.0 & +10:59 & 16.2 & Dec 17 & 11:56.1 & +3:55 & 16.3 \\
Dec 2 & 11:49.4 & +10:37 & 16.2 & Dec 18 & 11:56.6 & +3:25 & 16.3 \\
Dec 3 & 11:49.8 & +10:13 & 16.2 & Dec 19 & 11:57.1 & +2:56 & 16.3 \\
Dec 4 & 11:50.2 & +9:49 & 16.2 & Dec 20 & 11:57.6 & +2:26 & 16.3 \\
Dec 5 & 11:50.6 & +9:25 & 16.2 & Dec 21 & 11:58.1 & +1:56 & 16.3 \\
Dec 6 & 11:51.0 & +8:59 & 16.2 & Dec 22 & 11:58.6 & +1:26 & 16.3 \\
Dec 7 & 11:51.4 & +8:34 & 16.2 & Dec 23 & 11:59.1 & +0:57 & 16.3 \\
Dec 8 & 11:51.9 & +8:08 & 16.2 & Dec 24 & 11:59.6 & +0:27 & 16.3 \\
Dec 9 & 11:52.3 & +7:41 & 16.2 & Dec 25 & 12:00.0 & $-$0:03 & 16.3 \\
Dec 10 & 11:52.8 & +7:14 & 16.2 & Dec 26 & 12:00.5 & $-$0:33 & 16.3 \\
Dec 11 & 11:53.2 & +6:46 & 16.2 & Dec 27 & 12:01.0 & $-$1:02 & 16.3 \\
Dec 12 & 11:53.7 & +6:19 & 16.3 & Dec 28 & 12:01.5 & $-$1:32 & 16.3 \\
Dec 13 & 11:54.2 & +5:50 & 16.3 & Dec 29 & 12:02.0 & $-$2:01 & 16.3 \\
Dec 14 & 11:54.6 & +5:22 & 16.3 & Dec 30 & 12:02.5 & $-$2:30 & 16.3 \\
Dec 15 & 11:55.1 & +4:53 & 16.3 & Dec 31 & 12:03.0 & $-$2:59 & 16.3 \\
Dec 16 & 11:55.6 & +4:24 & 16.3 &  & & &  \\
\end{longtable}}
}

\chapter{Hourly GHA and Declination of the Sun}

Read the Greenwich Hour Angle and Declination of the Sun's centre for the day and
whole hour of UT; interpolate the minutes and seconds with the Sun increment table
in Part~IV. Entries run left to right, then top to bottom, in three columns of
day~/~hour~/~GHA~/~Dec.

{\footnotesize
\begin{longtable}{rrrr|rrrr|rrrr}
\caption{Sun --- GHA and Declination, January 2026 (UT)}\label{sun01}\\
\toprule
\textbf{d} & \textbf{h} & \textbf{GHA} & \textbf{Dec} & \textbf{d} & \textbf{h} & \textbf{GHA} & \textbf{Dec} & \textbf{d} & \textbf{h} & \textbf{GHA} & \textbf{Dec} \\
\midrule
\endfirsthead
\multicolumn{12}{c}{\textit{Sun --- GHA and Declination, January 2026 (UT) (continued)}}\\
\toprule
\textbf{d} & \textbf{h} & \textbf{GHA} & \textbf{Dec} & \textbf{d} & \textbf{h} & \textbf{GHA} & \textbf{Dec} & \textbf{d} & \textbf{h} & \textbf{GHA} & \textbf{Dec} \\
\midrule
\endhead
\midrule\multicolumn{12}{r}{\textit{continued on next page}}\\
\endfoot
\bottomrule
\endlastfoot
1 & 0 & 179\dg 10.0\pr & S23\dg 00.9\pr & 1 & 1 & 194\dg 09.7\pr & S23\dg 00.7\pr & 1 & 2 & 209\dg 09.4\pr & S23\dg 00.5\pr \\
1 & 3 & 224\dg 09.1\pr & S23\dg 00.3\pr & 1 & 4 & 239\dg 08.8\pr & S23\dg 00.1\pr & 1 & 5 & 254\dg 08.6\pr & S22\dg 59.9\pr \\
1 & 6 & 269\dg 08.3\pr & S22\dg 59.7\pr & 1 & 7 & 284\dg 08.0\pr & S22\dg 59.4\pr & 1 & 8 & 299\dg 07.7\pr & S22\dg 59.2\pr \\
1 & 9 & 314\dg 07.4\pr & S22\dg 59.0\pr & 1 & 10 & 329\dg 07.1\pr & S22\dg 58.8\pr & 1 & 11 & 344\dg 06.8\pr & S22\dg 58.6\pr \\
1 & 12 & 359\dg 06.5\pr & S22\dg 58.4\pr & 1 & 13 & 14\dg 06.2\pr & S22\dg 58.2\pr & 1 & 14 & 29\dg 05.9\pr & S22\dg 58.0\pr \\
1 & 15 & 44\dg 05.6\pr & S22\dg 57.8\pr & 1 & 16 & 59\dg 05.3\pr & S22\dg 57.5\pr & 1 & 17 & 74\dg 05.0\pr & S22\dg 57.3\pr \\
1 & 18 & 89\dg 04.7\pr & S22\dg 57.1\pr & 1 & 19 & 104\dg 04.4\pr & S22\dg 56.9\pr & 1 & 20 & 119\dg 04.1\pr & S22\dg 56.7\pr \\
1 & 21 & 134\dg 03.9\pr & S22\dg 56.5\pr & 1 & 22 & 149\dg 03.6\pr & S22\dg 56.2\pr & 1 & 23 & 164\dg 03.3\pr & S22\dg 56.0\pr \\
2 & 0 & 179\dg 03.0\pr & S22\dg 55.8\pr & 2 & 1 & 194\dg 02.7\pr & S22\dg 55.6\pr & 2 & 2 & 209\dg 02.4\pr & S22\dg 55.4\pr \\
2 & 3 & 224\dg 02.1\pr & S22\dg 55.1\pr & 2 & 4 & 239\dg 01.8\pr & S22\dg 54.9\pr & 2 & 5 & 254\dg 01.5\pr & S22\dg 54.7\pr \\
2 & 6 & 269\dg 01.2\pr & S22\dg 54.5\pr & 2 & 7 & 284\dg 00.9\pr & S22\dg 54.2\pr & 2 & 8 & 299\dg 00.6\pr & S22\dg 54.0\pr \\
2 & 9 & 314\dg 00.4\pr & S22\dg 53.8\pr & 2 & 10 & 329\dg 00.1\pr & S22\dg 53.6\pr & 2 & 11 & 343\dg 59.8\pr & S22\dg 53.3\pr \\
2 & 12 & 358\dg 59.5\pr & S22\dg 53.1\pr & 2 & 13 & 13\dg 59.2\pr & S22\dg 52.9\pr & 2 & 14 & 28\dg 58.9\pr & S22\dg 52.6\pr \\
2 & 15 & 43\dg 58.6\pr & S22\dg 52.4\pr & 2 & 16 & 58\dg 58.3\pr & S22\dg 52.2\pr & 2 & 17 & 73\dg 58.0\pr & S22\dg 51.9\pr \\
2 & 18 & 88\dg 57.7\pr & S22\dg 51.7\pr & 2 & 19 & 103\dg 57.5\pr & S22\dg 51.5\pr & 2 & 20 & 118\dg 57.2\pr & S22\dg 51.2\pr \\
2 & 21 & 133\dg 56.9\pr & S22\dg 51.0\pr & 2 & 22 & 148\dg 56.6\pr & S22\dg 50.7\pr & 2 & 23 & 163\dg 56.3\pr & S22\dg 50.5\pr \\
3 & 0 & 178\dg 56.0\pr & S22\dg 50.3\pr & 3 & 1 & 193\dg 55.7\pr & S22\dg 50.0\pr & 3 & 2 & 208\dg 55.4\pr & S22\dg 49.8\pr \\
3 & 3 & 223\dg 55.1\pr & S22\dg 49.5\pr & 3 & 4 & 238\dg 54.9\pr & S22\dg 49.3\pr & 3 & 5 & 253\dg 54.6\pr & S22\dg 49.1\pr \\
3 & 6 & 268\dg 54.3\pr & S22\dg 48.8\pr & 3 & 7 & 283\dg 54.0\pr & S22\dg 48.6\pr & 3 & 8 & 298\dg 53.7\pr & S22\dg 48.3\pr \\
3 & 9 & 313\dg 53.4\pr & S22\dg 48.1\pr & 3 & 10 & 328\dg 53.1\pr & S22\dg 47.8\pr & 3 & 11 & 343\dg 52.8\pr & S22\dg 47.6\pr \\
3 & 12 & 358\dg 52.6\pr & S22\dg 47.3\pr & 3 & 13 & 13\dg 52.3\pr & S22\dg 47.1\pr & 3 & 14 & 28\dg 52.0\pr & S22\dg 46.8\pr \\
3 & 15 & 43\dg 51.7\pr & S22\dg 46.6\pr & 3 & 16 & 58\dg 51.4\pr & S22\dg 46.3\pr & 3 & 17 & 73\dg 51.1\pr & S22\dg 46.1\pr \\
3 & 18 & 88\dg 50.8\pr & S22\dg 45.8\pr & 3 & 19 & 103\dg 50.6\pr & S22\dg 45.6\pr & 3 & 20 & 118\dg 50.3\pr & S22\dg 45.3\pr \\
3 & 21 & 133\dg 50.0\pr & S22\dg 45.1\pr & 3 & 22 & 148\dg 49.7\pr & S22\dg 44.8\pr & 3 & 23 & 163\dg 49.4\pr & S22\dg 44.5\pr \\
4 & 0 & 178\dg 49.1\pr & S22\dg 44.3\pr & 4 & 1 & 193\dg 48.8\pr & S22\dg 44.0\pr & 4 & 2 & 208\dg 48.6\pr & S22\dg 43.8\pr \\
4 & 3 & 223\dg 48.3\pr & S22\dg 43.5\pr & 4 & 4 & 238\dg 48.0\pr & S22\dg 43.2\pr & 4 & 5 & 253\dg 47.7\pr & S22\dg 43.0\pr \\
4 & 6 & 268\dg 47.4\pr & S22\dg 42.7\pr & 4 & 7 & 283\dg 47.1\pr & S22\dg 42.5\pr & 4 & 8 & 298\dg 46.9\pr & S22\dg 42.2\pr \\
4 & 9 & 313\dg 46.6\pr & S22\dg 41.9\pr & 4 & 10 & 328\dg 46.3\pr & S22\dg 41.7\pr & 4 & 11 & 343\dg 46.0\pr & S22\dg 41.4\pr \\
4 & 12 & 358\dg 45.7\pr & S22\dg 41.1\pr & 4 & 13 & 13\dg 45.4\pr & S22\dg 40.9\pr & 4 & 14 & 28\dg 45.2\pr & S22\dg 40.6\pr \\
4 & 15 & 43\dg 44.9\pr & S22\dg 40.3\pr & 4 & 16 & 58\dg 44.6\pr & S22\dg 40.0\pr & 4 & 17 & 73\dg 44.3\pr & S22\dg 39.8\pr \\
4 & 18 & 88\dg 44.0\pr & S22\dg 39.5\pr & 4 & 19 & 103\dg 43.8\pr & S22\dg 39.2\pr & 4 & 20 & 118\dg 43.5\pr & S22\dg 38.9\pr \\
4 & 21 & 133\dg 43.2\pr & S22\dg 38.7\pr & 4 & 22 & 148\dg 42.9\pr & S22\dg 38.4\pr & 4 & 23 & 163\dg 42.6\pr & S22\dg 38.1\pr \\
5 & 0 & 178\dg 42.4\pr & S22\dg 37.8\pr & 5 & 1 & 193\dg 42.1\pr & S22\dg 37.6\pr & 5 & 2 & 208\dg 41.8\pr & S22\dg 37.3\pr \\
5 & 3 & 223\dg 41.5\pr & S22\dg 37.0\pr & 5 & 4 & 238\dg 41.2\pr & S22\dg 36.7\pr & 5 & 5 & 253\dg 41.0\pr & S22\dg 36.4\pr \\
5 & 6 & 268\dg 40.7\pr & S22\dg 36.2\pr & 5 & 7 & 283\dg 40.4\pr & S22\dg 35.9\pr & 5 & 8 & 298\dg 40.1\pr & S22\dg 35.6\pr \\
5 & 9 & 313\dg 39.8\pr & S22\dg 35.3\pr & 5 & 10 & 328\dg 39.6\pr & S22\dg 35.0\pr & 5 & 11 & 343\dg 39.3\pr & S22\dg 34.7\pr \\
5 & 12 & 358\dg 39.0\pr & S22\dg 34.5\pr & 5 & 13 & 13\dg 38.7\pr & S22\dg 34.2\pr & 5 & 14 & 28\dg 38.4\pr & S22\dg 33.9\pr \\
5 & 15 & 43\dg 38.2\pr & S22\dg 33.6\pr & 5 & 16 & 58\dg 37.9\pr & S22\dg 33.3\pr & 5 & 17 & 73\dg 37.6\pr & S22\dg 33.0\pr \\
5 & 18 & 88\dg 37.3\pr & S22\dg 32.7\pr & 5 & 19 & 103\dg 37.1\pr & S22\dg 32.4\pr & 5 & 20 & 118\dg 36.8\pr & S22\dg 32.1\pr \\
5 & 21 & 133\dg 36.5\pr & S22\dg 31.8\pr & 5 & 22 & 148\dg 36.2\pr & S22\dg 31.5\pr & 5 & 23 & 163\dg 36.0\pr & S22\dg 31.3\pr \\
6 & 0 & 178\dg 35.7\pr & S22\dg 31.0\pr & 6 & 1 & 193\dg 35.4\pr & S22\dg 30.7\pr & 6 & 2 & 208\dg 35.1\pr & S22\dg 30.4\pr \\
6 & 3 & 223\dg 34.9\pr & S22\dg 30.1\pr & 6 & 4 & 238\dg 34.6\pr & S22\dg 29.8\pr & 6 & 5 & 253\dg 34.3\pr & S22\dg 29.5\pr \\
6 & 6 & 268\dg 34.0\pr & S22\dg 29.2\pr & 6 & 7 & 283\dg 33.8\pr & S22\dg 28.9\pr & 6 & 8 & 298\dg 33.5\pr & S22\dg 28.6\pr \\
6 & 9 & 313\dg 33.2\pr & S22\dg 28.3\pr & 6 & 10 & 328\dg 32.9\pr & S22\dg 28.0\pr & 6 & 11 & 343\dg 32.7\pr & S22\dg 27.7\pr \\
6 & 12 & 358\dg 32.4\pr & S22\dg 27.3\pr & 6 & 13 & 13\dg 32.1\pr & S22\dg 27.0\pr & 6 & 14 & 28\dg 31.8\pr & S22\dg 26.7\pr \\
6 & 15 & 43\dg 31.6\pr & S22\dg 26.4\pr & 6 & 16 & 58\dg 31.3\pr & S22\dg 26.1\pr & 6 & 17 & 73\dg 31.0\pr & S22\dg 25.8\pr \\
6 & 18 & 88\dg 30.8\pr & S22\dg 25.5\pr & 6 & 19 & 103\dg 30.5\pr & S22\dg 25.2\pr & 6 & 20 & 118\dg 30.2\pr & S22\dg 24.9\pr \\
6 & 21 & 133\dg 29.9\pr & S22\dg 24.6\pr & 6 & 22 & 148\dg 29.7\pr & S22\dg 24.3\pr & 6 & 23 & 163\dg 29.4\pr & S22\dg 23.9\pr \\
7 & 0 & 178\dg 29.1\pr & S22\dg 23.6\pr & 7 & 1 & 193\dg 28.9\pr & S22\dg 23.3\pr & 7 & 2 & 208\dg 28.6\pr & S22\dg 23.0\pr \\
7 & 3 & 223\dg 28.3\pr & S22\dg 22.7\pr & 7 & 4 & 238\dg 28.0\pr & S22\dg 22.4\pr & 7 & 5 & 253\dg 27.8\pr & S22\dg 22.0\pr \\
7 & 6 & 268\dg 27.5\pr & S22\dg 21.7\pr & 7 & 7 & 283\dg 27.2\pr & S22\dg 21.4\pr & 7 & 8 & 298\dg 27.0\pr & S22\dg 21.1\pr \\
7 & 9 & 313\dg 26.7\pr & S22\dg 20.8\pr & 7 & 10 & 328\dg 26.4\pr & S22\dg 20.4\pr & 7 & 11 & 343\dg 26.2\pr & S22\dg 20.1\pr \\
7 & 12 & 358\dg 25.9\pr & S22\dg 19.8\pr & 7 & 13 & 13\dg 25.6\pr & S22\dg 19.5\pr & 7 & 14 & 28\dg 25.4\pr & S22\dg 19.1\pr \\
7 & 15 & 43\dg 25.1\pr & S22\dg 18.8\pr & 7 & 16 & 58\dg 24.8\pr & S22\dg 18.5\pr & 7 & 17 & 73\dg 24.6\pr & S22\dg 18.2\pr \\
7 & 18 & 88\dg 24.3\pr & S22\dg 17.8\pr & 7 & 19 & 103\dg 24.0\pr & S22\dg 17.5\pr & 7 & 20 & 118\dg 23.8\pr & S22\dg 17.2\pr \\
7 & 21 & 133\dg 23.5\pr & S22\dg 16.9\pr & 7 & 22 & 148\dg 23.2\pr & S22\dg 16.5\pr & 7 & 23 & 163\dg 23.0\pr & S22\dg 16.2\pr \\
8 & 0 & 178\dg 22.7\pr & S22\dg 15.9\pr & 8 & 1 & 193\dg 22.4\pr & S22\dg 15.5\pr & 8 & 2 & 208\dg 22.2\pr & S22\dg 15.2\pr \\
8 & 3 & 223\dg 21.9\pr & S22\dg 14.9\pr & 8 & 4 & 238\dg 21.6\pr & S22\dg 14.5\pr & 8 & 5 & 253\dg 21.4\pr & S22\dg 14.2\pr \\
8 & 6 & 268\dg 21.1\pr & S22\dg 13.8\pr & 8 & 7 & 283\dg 20.8\pr & S22\dg 13.5\pr & 8 & 8 & 298\dg 20.6\pr & S22\dg 13.2\pr \\
8 & 9 & 313\dg 20.3\pr & S22\dg 12.8\pr & 8 & 10 & 328\dg 20.0\pr & S22\dg 12.5\pr & 8 & 11 & 343\dg 19.8\pr & S22\dg 12.1\pr \\
8 & 12 & 358\dg 19.5\pr & S22\dg 11.8\pr & 8 & 13 & 13\dg 19.3\pr & S22\dg 11.5\pr & 8 & 14 & 28\dg 19.0\pr & S22\dg 11.1\pr \\
8 & 15 & 43\dg 18.7\pr & S22\dg 10.8\pr & 8 & 16 & 58\dg 18.5\pr & S22\dg 10.4\pr & 8 & 17 & 73\dg 18.2\pr & S22\dg 10.1\pr \\
8 & 18 & 88\dg 18.0\pr & S22\dg 09.7\pr & 8 & 19 & 103\dg 17.7\pr & S22\dg 09.4\pr & 8 & 20 & 118\dg 17.4\pr & S22\dg 09.0\pr \\
8 & 21 & 133\dg 17.2\pr & S22\dg 08.7\pr & 8 & 22 & 148\dg 16.9\pr & S22\dg 08.3\pr & 8 & 23 & 163\dg 16.6\pr & S22\dg 08.0\pr \\
9 & 0 & 178\dg 16.4\pr & S22\dg 07.6\pr & 9 & 1 & 193\dg 16.1\pr & S22\dg 07.3\pr & 9 & 2 & 208\dg 15.9\pr & S22\dg 06.9\pr \\
9 & 3 & 223\dg 15.6\pr & S22\dg 06.6\pr & 9 & 4 & 238\dg 15.3\pr & S22\dg 06.2\pr & 9 & 5 & 253\dg 15.1\pr & S22\dg 05.9\pr \\
9 & 6 & 268\dg 14.8\pr & S22\dg 05.5\pr & 9 & 7 & 283\dg 14.6\pr & S22\dg 05.2\pr & 9 & 8 & 298\dg 14.3\pr & S22\dg 04.8\pr \\
9 & 9 & 313\dg 14.1\pr & S22\dg 04.5\pr & 9 & 10 & 328\dg 13.8\pr & S22\dg 04.1\pr & 9 & 11 & 343\dg 13.5\pr & S22\dg 03.7\pr \\
9 & 12 & 358\dg 13.3\pr & S22\dg 03.4\pr & 9 & 13 & 13\dg 13.0\pr & S22\dg 03.0\pr & 9 & 14 & 28\dg 12.8\pr & S22\dg 02.7\pr \\
9 & 15 & 43\dg 12.5\pr & S22\dg 02.3\pr & 9 & 16 & 58\dg 12.3\pr & S22\dg 01.9\pr & 9 & 17 & 73\dg 12.0\pr & S22\dg 01.6\pr \\
9 & 18 & 88\dg 11.7\pr & S22\dg 01.2\pr & 9 & 19 & 103\dg 11.5\pr & S22\dg 00.8\pr & 9 & 20 & 118\dg 11.2\pr & S22\dg 00.5\pr \\
9 & 21 & 133\dg 11.0\pr & S22\dg 00.1\pr & 9 & 22 & 148\dg 10.7\pr & S21\dg 59.7\pr & 9 & 23 & 163\dg 10.5\pr & S21\dg 59.4\pr \\
10 & 0 & 178\dg 10.2\pr & S21\dg 59.0\pr & 10 & 1 & 193\dg 10.0\pr & S21\dg 58.6\pr & 10 & 2 & 208\dg 09.7\pr & S21\dg 58.3\pr \\
10 & 3 & 223\dg 09.5\pr & S21\dg 57.9\pr & 10 & 4 & 238\dg 09.2\pr & S21\dg 57.5\pr & 10 & 5 & 253\dg 08.9\pr & S21\dg 57.1\pr \\
10 & 6 & 268\dg 08.7\pr & S21\dg 56.8\pr & 10 & 7 & 283\dg 08.4\pr & S21\dg 56.4\pr & 10 & 8 & 298\dg 08.2\pr & S21\dg 56.0\pr \\
10 & 9 & 313\dg 07.9\pr & S21\dg 55.7\pr & 10 & 10 & 328\dg 07.7\pr & S21\dg 55.3\pr & 10 & 11 & 343\dg 07.4\pr & S21\dg 54.9\pr \\
10 & 12 & 358\dg 07.2\pr & S21\dg 54.5\pr & 10 & 13 & 13\dg 06.9\pr & S21\dg 54.1\pr & 10 & 14 & 28\dg 06.7\pr & S21\dg 53.8\pr \\
10 & 15 & 43\dg 06.4\pr & S21\dg 53.4\pr & 10 & 16 & 58\dg 06.2\pr & S21\dg 53.0\pr & 10 & 17 & 73\dg 05.9\pr & S21\dg 52.6\pr \\
10 & 18 & 88\dg 05.7\pr & S21\dg 52.2\pr & 10 & 19 & 103\dg 05.4\pr & S21\dg 51.9\pr & 10 & 20 & 118\dg 05.2\pr & S21\dg 51.5\pr \\
10 & 21 & 133\dg 04.9\pr & S21\dg 51.1\pr & 10 & 22 & 148\dg 04.7\pr & S21\dg 50.7\pr & 10 & 23 & 163\dg 04.4\pr & S21\dg 50.3\pr \\
11 & 0 & 178\dg 04.2\pr & S21\dg 49.9\pr & 11 & 1 & 193\dg 03.9\pr & S21\dg 49.5\pr & 11 & 2 & 208\dg 03.7\pr & S21\dg 49.2\pr \\
11 & 3 & 223\dg 03.4\pr & S21\dg 48.8\pr & 11 & 4 & 238\dg 03.2\pr & S21\dg 48.4\pr & 11 & 5 & 253\dg 02.9\pr & S21\dg 48.0\pr \\
11 & 6 & 268\dg 02.7\pr & S21\dg 47.6\pr & 11 & 7 & 283\dg 02.4\pr & S21\dg 47.2\pr & 11 & 8 & 298\dg 02.2\pr & S21\dg 46.8\pr \\
11 & 9 & 313\dg 02.0\pr & S21\dg 46.4\pr & 11 & 10 & 328\dg 01.7\pr & S21\dg 46.0\pr & 11 & 11 & 343\dg 01.5\pr & S21\dg 45.6\pr \\
11 & 12 & 358\dg 01.2\pr & S21\dg 45.2\pr & 11 & 13 & 13\dg 01.0\pr & S21\dg 44.8\pr & 11 & 14 & 28\dg 00.7\pr & S21\dg 44.4\pr \\
11 & 15 & 43\dg 00.5\pr & S21\dg 44.0\pr & 11 & 16 & 58\dg 00.2\pr & S21\dg 43.6\pr & 11 & 17 & 73\dg 00.0\pr & S21\dg 43.2\pr \\
11 & 18 & 87\dg 59.8\pr & S21\dg 42.9\pr & 11 & 19 & 102\dg 59.5\pr & S21\dg 42.4\pr & 11 & 20 & 117\dg 59.3\pr & S21\dg 42.0\pr \\
11 & 21 & 132\dg 59.0\pr & S21\dg 41.6\pr & 11 & 22 & 147\dg 58.8\pr & S21\dg 41.2\pr & 11 & 23 & 162\dg 58.5\pr & S21\dg 40.8\pr \\
12 & 0 & 177\dg 58.3\pr & S21\dg 40.4\pr & 12 & 1 & 192\dg 58.1\pr & S21\dg 40.0\pr & 12 & 2 & 207\dg 57.8\pr & S21\dg 39.6\pr \\
12 & 3 & 222\dg 57.6\pr & S21\dg 39.2\pr & 12 & 4 & 237\dg 57.3\pr & S21\dg 38.8\pr & 12 & 5 & 252\dg 57.1\pr & S21\dg 38.4\pr \\
12 & 6 & 267\dg 56.8\pr & S21\dg 38.0\pr & 12 & 7 & 282\dg 56.6\pr & S21\dg 37.6\pr & 12 & 8 & 297\dg 56.4\pr & S21\dg 37.2\pr \\
12 & 9 & 312\dg 56.1\pr & S21\dg 36.8\pr & 12 & 10 & 327\dg 55.9\pr & S21\dg 36.4\pr & 12 & 11 & 342\dg 55.6\pr & S21\dg 35.9\pr \\
12 & 12 & 357\dg 55.4\pr & S21\dg 35.5\pr & 12 & 13 & 12\dg 55.2\pr & S21\dg 35.1\pr & 12 & 14 & 27\dg 54.9\pr & S21\dg 34.7\pr \\
12 & 15 & 42\dg 54.7\pr & S21\dg 34.3\pr & 12 & 16 & 57\dg 54.5\pr & S21\dg 33.9\pr & 12 & 17 & 72\dg 54.2\pr & S21\dg 33.5\pr \\
12 & 18 & 87\dg 54.0\pr & S21\dg 33.0\pr & 12 & 19 & 102\dg 53.7\pr & S21\dg 32.6\pr & 12 & 20 & 117\dg 53.5\pr & S21\dg 32.2\pr \\
12 & 21 & 132\dg 53.3\pr & S21\dg 31.8\pr & 12 & 22 & 147\dg 53.0\pr & S21\dg 31.4\pr & 12 & 23 & 162\dg 52.8\pr & S21\dg 30.9\pr \\
13 & 0 & 177\dg 52.6\pr & S21\dg 30.5\pr & 13 & 1 & 192\dg 52.3\pr & S21\dg 30.1\pr & 13 & 2 & 207\dg 52.1\pr & S21\dg 29.7\pr \\
13 & 3 & 222\dg 51.8\pr & S21\dg 29.3\pr & 13 & 4 & 237\dg 51.6\pr & S21\dg 28.8\pr & 13 & 5 & 252\dg 51.4\pr & S21\dg 28.4\pr \\
13 & 6 & 267\dg 51.1\pr & S21\dg 28.0\pr & 13 & 7 & 282\dg 50.9\pr & S21\dg 27.5\pr & 13 & 8 & 297\dg 50.7\pr & S21\dg 27.1\pr \\
13 & 9 & 312\dg 50.4\pr & S21\dg 26.7\pr & 13 & 10 & 327\dg 50.2\pr & S21\dg 26.3\pr & 13 & 11 & 342\dg 50.0\pr & S21\dg 25.8\pr \\
13 & 12 & 357\dg 49.7\pr & S21\dg 25.4\pr & 13 & 13 & 12\dg 49.5\pr & S21\dg 25.0\pr & 13 & 14 & 27\dg 49.3\pr & S21\dg 24.5\pr \\
13 & 15 & 42\dg 49.1\pr & S21\dg 24.1\pr & 13 & 16 & 57\dg 48.8\pr & S21\dg 23.7\pr & 13 & 17 & 72\dg 48.6\pr & S21\dg 23.2\pr \\
13 & 18 & 87\dg 48.4\pr & S21\dg 22.8\pr & 13 & 19 & 102\dg 48.1\pr & S21\dg 22.4\pr & 13 & 20 & 117\dg 47.9\pr & S21\dg 21.9\pr \\
13 & 21 & 132\dg 47.7\pr & S21\dg 21.5\pr & 13 & 22 & 147\dg 47.4\pr & S21\dg 21.1\pr & 13 & 23 & 162\dg 47.2\pr & S21\dg 20.6\pr \\
14 & 0 & 177\dg 47.0\pr & S21\dg 20.2\pr & 14 & 1 & 192\dg 46.7\pr & S21\dg 19.7\pr & 14 & 2 & 207\dg 46.5\pr & S21\dg 19.3\pr \\
14 & 3 & 222\dg 46.3\pr & S21\dg 18.9\pr & 14 & 4 & 237\dg 46.1\pr & S21\dg 18.4\pr & 14 & 5 & 252\dg 45.8\pr & S21\dg 18.0\pr \\
14 & 6 & 267\dg 45.6\pr & S21\dg 17.5\pr & 14 & 7 & 282\dg 45.4\pr & S21\dg 17.1\pr & 14 & 8 & 297\dg 45.2\pr & S21\dg 16.7\pr \\
14 & 9 & 312\dg 44.9\pr & S21\dg 16.2\pr & 14 & 10 & 327\dg 44.7\pr & S21\dg 15.8\pr & 14 & 11 & 342\dg 44.5\pr & S21\dg 15.3\pr \\
14 & 12 & 357\dg 44.2\pr & S21\dg 14.9\pr & 14 & 13 & 12\dg 44.0\pr & S21\dg 14.4\pr & 14 & 14 & 27\dg 43.8\pr & S21\dg 14.0\pr \\
14 & 15 & 42\dg 43.6\pr & S21\dg 13.5\pr & 14 & 16 & 57\dg 43.3\pr & S21\dg 13.1\pr & 14 & 17 & 72\dg 43.1\pr & S21\dg 12.6\pr \\
14 & 18 & 87\dg 42.9\pr & S21\dg 12.2\pr & 14 & 19 & 102\dg 42.7\pr & S21\dg 11.7\pr & 14 & 20 & 117\dg 42.5\pr & S21\dg 11.3\pr \\
14 & 21 & 132\dg 42.2\pr & S21\dg 10.8\pr & 14 & 22 & 147\dg 42.0\pr & S21\dg 10.4\pr & 14 & 23 & 162\dg 41.8\pr & S21\dg 09.9\pr \\
15 & 0 & 177\dg 41.6\pr & S21\dg 09.4\pr & 15 & 1 & 192\dg 41.3\pr & S21\dg 09.0\pr & 15 & 2 & 207\dg 41.1\pr & S21\dg 08.5\pr \\
15 & 3 & 222\dg 40.9\pr & S21\dg 08.1\pr & 15 & 4 & 237\dg 40.7\pr & S21\dg 07.6\pr & 15 & 5 & 252\dg 40.5\pr & S21\dg 07.2\pr \\
15 & 6 & 267\dg 40.2\pr & S21\dg 06.7\pr & 15 & 7 & 282\dg 40.0\pr & S21\dg 06.2\pr & 15 & 8 & 297\dg 39.8\pr & S21\dg 05.8\pr \\
15 & 9 & 312\dg 39.6\pr & S21\dg 05.3\pr & 15 & 10 & 327\dg 39.3\pr & S21\dg 04.8\pr & 15 & 11 & 342\dg 39.1\pr & S21\dg 04.4\pr \\
15 & 12 & 357\dg 38.9\pr & S21\dg 03.9\pr & 15 & 13 & 12\dg 38.7\pr & S21\dg 03.5\pr & 15 & 14 & 27\dg 38.5\pr & S21\dg 03.0\pr \\
15 & 15 & 42\dg 38.3\pr & S21\dg 02.5\pr & 15 & 16 & 57\dg 38.0\pr & S21\dg 02.1\pr & 15 & 17 & 72\dg 37.8\pr & S21\dg 01.6\pr \\
15 & 18 & 87\dg 37.6\pr & S21\dg 01.1\pr & 15 & 19 & 102\dg 37.4\pr & S21\dg 00.7\pr & 15 & 20 & 117\dg 37.2\pr & S21\dg 00.2\pr \\
15 & 21 & 132\dg 37.0\pr & S20\dg 59.7\pr & 15 & 22 & 147\dg 36.7\pr & S20\dg 59.2\pr & 15 & 23 & 162\dg 36.5\pr & S20\dg 58.8\pr \\
16 & 0 & 177\dg 36.3\pr & S20\dg 58.3\pr & 16 & 1 & 192\dg 36.1\pr & S20\dg 57.8\pr & 16 & 2 & 207\dg 35.9\pr & S20\dg 57.3\pr \\
16 & 3 & 222\dg 35.7\pr & S20\dg 56.9\pr & 16 & 4 & 237\dg 35.4\pr & S20\dg 56.4\pr & 16 & 5 & 252\dg 35.2\pr & S20\dg 55.9\pr \\
16 & 6 & 267\dg 35.0\pr & S20\dg 55.4\pr & 16 & 7 & 282\dg 34.8\pr & S20\dg 55.0\pr & 16 & 8 & 297\dg 34.6\pr & S20\dg 54.5\pr \\
16 & 9 & 312\dg 34.4\pr & S20\dg 54.0\pr & 16 & 10 & 327\dg 34.2\pr & S20\dg 53.5\pr & 16 & 11 & 342\dg 34.0\pr & S20\dg 53.1\pr \\
16 & 12 & 357\dg 33.7\pr & S20\dg 52.6\pr & 16 & 13 & 12\dg 33.5\pr & S20\dg 52.1\pr & 16 & 14 & 27\dg 33.3\pr & S20\dg 51.6\pr \\
16 & 15 & 42\dg 33.1\pr & S20\dg 51.1\pr & 16 & 16 & 57\dg 32.9\pr & S20\dg 50.6\pr & 16 & 17 & 72\dg 32.7\pr & S20\dg 50.2\pr \\
16 & 18 & 87\dg 32.5\pr & S20\dg 49.7\pr & 16 & 19 & 102\dg 32.3\pr & S20\dg 49.2\pr & 16 & 20 & 117\dg 32.1\pr & S20\dg 48.7\pr \\
16 & 21 & 132\dg 31.9\pr & S20\dg 48.2\pr & 16 & 22 & 147\dg 31.6\pr & S20\dg 47.7\pr & 16 & 23 & 162\dg 31.4\pr & S20\dg 47.2\pr \\
17 & 0 & 177\dg 31.2\pr & S20\dg 46.7\pr & 17 & 1 & 192\dg 31.0\pr & S20\dg 46.3\pr & 17 & 2 & 207\dg 30.8\pr & S20\dg 45.8\pr \\
17 & 3 & 222\dg 30.6\pr & S20\dg 45.3\pr & 17 & 4 & 237\dg 30.4\pr & S20\dg 44.8\pr & 17 & 5 & 252\dg 30.2\pr & S20\dg 44.3\pr \\
17 & 6 & 267\dg 30.0\pr & S20\dg 43.8\pr & 17 & 7 & 282\dg 29.8\pr & S20\dg 43.3\pr & 17 & 8 & 297\dg 29.6\pr & S20\dg 42.8\pr \\
17 & 9 & 312\dg 29.4\pr & S20\dg 42.3\pr & 17 & 10 & 327\dg 29.2\pr & S20\dg 41.8\pr & 17 & 11 & 342\dg 29.0\pr & S20\dg 41.3\pr \\
17 & 12 & 357\dg 28.8\pr & S20\dg 40.8\pr & 17 & 13 & 12\dg 28.5\pr & S20\dg 40.3\pr & 17 & 14 & 27\dg 28.3\pr & S20\dg 39.8\pr \\
17 & 15 & 42\dg 28.1\pr & S20\dg 39.3\pr & 17 & 16 & 57\dg 27.9\pr & S20\dg 38.8\pr & 17 & 17 & 72\dg 27.7\pr & S20\dg 38.3\pr \\
17 & 18 & 87\dg 27.5\pr & S20\dg 37.8\pr & 17 & 19 & 102\dg 27.3\pr & S20\dg 37.3\pr & 17 & 20 & 117\dg 27.1\pr & S20\dg 36.8\pr \\
17 & 21 & 132\dg 26.9\pr & S20\dg 36.3\pr & 17 & 22 & 147\dg 26.7\pr & S20\dg 35.8\pr & 17 & 23 & 162\dg 26.5\pr & S20\dg 35.3\pr \\
18 & 0 & 177\dg 26.3\pr & S20\dg 34.8\pr & 18 & 1 & 192\dg 26.1\pr & S20\dg 34.3\pr & 18 & 2 & 207\dg 25.9\pr & S20\dg 33.8\pr \\
18 & 3 & 222\dg 25.7\pr & S20\dg 33.3\pr & 18 & 4 & 237\dg 25.5\pr & S20\dg 32.8\pr & 18 & 5 & 252\dg 25.3\pr & S20\dg 32.3\pr \\
18 & 6 & 267\dg 25.1\pr & S20\dg 31.8\pr & 18 & 7 & 282\dg 24.9\pr & S20\dg 31.2\pr & 18 & 8 & 297\dg 24.7\pr & S20\dg 30.7\pr \\
18 & 9 & 312\dg 24.5\pr & S20\dg 30.2\pr & 18 & 10 & 327\dg 24.3\pr & S20\dg 29.7\pr & 18 & 11 & 342\dg 24.1\pr & S20\dg 29.2\pr \\
18 & 12 & 357\dg 23.9\pr & S20\dg 28.7\pr & 18 & 13 & 12\dg 23.7\pr & S20\dg 28.2\pr & 18 & 14 & 27\dg 23.5\pr & S20\dg 27.7\pr \\
18 & 15 & 42\dg 23.3\pr & S20\dg 27.1\pr & 18 & 16 & 57\dg 23.1\pr & S20\dg 26.6\pr & 18 & 17 & 72\dg 23.0\pr & S20\dg 26.1\pr \\
18 & 18 & 87\dg 22.8\pr & S20\dg 25.6\pr & 18 & 19 & 102\dg 22.6\pr & S20\dg 25.1\pr & 18 & 20 & 117\dg 22.4\pr & S20\dg 24.5\pr \\
18 & 21 & 132\dg 22.2\pr & S20\dg 24.0\pr & 18 & 22 & 147\dg 22.0\pr & S20\dg 23.5\pr & 18 & 23 & 162\dg 21.8\pr & S20\dg 23.0\pr \\
19 & 0 & 177\dg 21.6\pr & S20\dg 22.5\pr & 19 & 1 & 192\dg 21.4\pr & S20\dg 21.9\pr & 19 & 2 & 207\dg 21.2\pr & S20\dg 21.4\pr \\
19 & 3 & 222\dg 21.0\pr & S20\dg 20.9\pr & 19 & 4 & 237\dg 20.8\pr & S20\dg 20.4\pr & 19 & 5 & 252\dg 20.6\pr & S20\dg 19.8\pr \\
19 & 6 & 267\dg 20.4\pr & S20\dg 19.3\pr & 19 & 7 & 282\dg 20.2\pr & S20\dg 18.8\pr & 19 & 8 & 297\dg 20.1\pr & S20\dg 18.3\pr \\
19 & 9 & 312\dg 19.9\pr & S20\dg 17.7\pr & 19 & 10 & 327\dg 19.7\pr & S20\dg 17.2\pr & 19 & 11 & 342\dg 19.5\pr & S20\dg 16.7\pr \\
19 & 12 & 357\dg 19.3\pr & S20\dg 16.2\pr & 19 & 13 & 12\dg 19.1\pr & S20\dg 15.6\pr & 19 & 14 & 27\dg 18.9\pr & S20\dg 15.1\pr \\
19 & 15 & 42\dg 18.7\pr & S20\dg 14.6\pr & 19 & 16 & 57\dg 18.5\pr & S20\dg 14.0\pr & 19 & 17 & 72\dg 18.4\pr & S20\dg 13.5\pr \\
19 & 18 & 87\dg 18.2\pr & S20\dg 13.0\pr & 19 & 19 & 102\dg 18.0\pr & S20\dg 12.4\pr & 19 & 20 & 117\dg 17.8\pr & S20\dg 11.9\pr \\
19 & 21 & 132\dg 17.6\pr & S20\dg 11.4\pr & 19 & 22 & 147\dg 17.4\pr & S20\dg 10.8\pr & 19 & 23 & 162\dg 17.2\pr & S20\dg 10.3\pr \\
20 & 0 & 177\dg 17.0\pr & S20\dg 09.7\pr & 20 & 1 & 192\dg 16.9\pr & S20\dg 09.2\pr & 20 & 2 & 207\dg 16.7\pr & S20\dg 08.7\pr \\
20 & 3 & 222\dg 16.5\pr & S20\dg 08.1\pr & 20 & 4 & 237\dg 16.3\pr & S20\dg 07.6\pr & 20 & 5 & 252\dg 16.1\pr & S20\dg 07.0\pr \\
20 & 6 & 267\dg 15.9\pr & S20\dg 06.5\pr & 20 & 7 & 282\dg 15.8\pr & S20\dg 06.0\pr & 20 & 8 & 297\dg 15.6\pr & S20\dg 05.4\pr \\
20 & 9 & 312\dg 15.4\pr & S20\dg 04.9\pr & 20 & 10 & 327\dg 15.2\pr & S20\dg 04.3\pr & 20 & 11 & 342\dg 15.0\pr & S20\dg 03.8\pr \\
20 & 12 & 357\dg 14.8\pr & S20\dg 03.2\pr & 20 & 13 & 12\dg 14.7\pr & S20\dg 02.7\pr & 20 & 14 & 27\dg 14.5\pr & S20\dg 02.2\pr \\
20 & 15 & 42\dg 14.3\pr & S20\dg 01.6\pr & 20 & 16 & 57\dg 14.1\pr & S20\dg 01.1\pr & 20 & 17 & 72\dg 13.9\pr & S20\dg 00.5\pr \\
20 & 18 & 87\dg 13.8\pr & S20\dg 00.0\pr & 20 & 19 & 102\dg 13.6\pr & S19\dg 59.4\pr & 20 & 20 & 117\dg 13.4\pr & S19\dg 58.9\pr \\
20 & 21 & 132\dg 13.2\pr & S19\dg 58.3\pr & 20 & 22 & 147\dg 13.0\pr & S19\dg 57.8\pr & 20 & 23 & 162\dg 12.9\pr & S19\dg 57.2\pr \\
21 & 0 & 177\dg 12.7\pr & S19\dg 56.6\pr & 21 & 1 & 192\dg 12.5\pr & S19\dg 56.1\pr & 21 & 2 & 207\dg 12.3\pr & S19\dg 55.5\pr \\
21 & 3 & 222\dg 12.2\pr & S19\dg 55.0\pr & 21 & 4 & 237\dg 12.0\pr & S19\dg 54.4\pr & 21 & 5 & 252\dg 11.8\pr & S19\dg 53.9\pr \\
21 & 6 & 267\dg 11.6\pr & S19\dg 53.3\pr & 21 & 7 & 282\dg 11.5\pr & S19\dg 52.8\pr & 21 & 8 & 297\dg 11.3\pr & S19\dg 52.2\pr \\
21 & 9 & 312\dg 11.1\pr & S19\dg 51.6\pr & 21 & 10 & 327\dg 10.9\pr & S19\dg 51.1\pr & 21 & 11 & 342\dg 10.8\pr & S19\dg 50.5\pr \\
21 & 12 & 357\dg 10.6\pr & S19\dg 50.0\pr & 21 & 13 & 12\dg 10.4\pr & S19\dg 49.4\pr & 21 & 14 & 27\dg 10.2\pr & S19\dg 48.8\pr \\
21 & 15 & 42\dg 10.1\pr & S19\dg 48.3\pr & 21 & 16 & 57\dg 09.9\pr & S19\dg 47.7\pr & 21 & 17 & 72\dg 09.7\pr & S19\dg 47.1\pr \\
21 & 18 & 87\dg 09.5\pr & S19\dg 46.6\pr & 21 & 19 & 102\dg 09.4\pr & S19\dg 46.0\pr & 21 & 20 & 117\dg 09.2\pr & S19\dg 45.5\pr \\
21 & 21 & 132\dg 09.0\pr & S19\dg 44.9\pr & 21 & 22 & 147\dg 08.9\pr & S19\dg 44.3\pr & 21 & 23 & 162\dg 08.7\pr & S19\dg 43.7\pr \\
22 & 0 & 177\dg 08.5\pr & S19\dg 43.2\pr & 22 & 1 & 192\dg 08.4\pr & S19\dg 42.6\pr & 22 & 2 & 207\dg 08.2\pr & S19\dg 42.0\pr \\
22 & 3 & 222\dg 08.0\pr & S19\dg 41.5\pr & 22 & 4 & 237\dg 07.8\pr & S19\dg 40.9\pr & 22 & 5 & 252\dg 07.7\pr & S19\dg 40.3\pr \\
22 & 6 & 267\dg 07.5\pr & S19\dg 39.8\pr & 22 & 7 & 282\dg 07.3\pr & S19\dg 39.2\pr & 22 & 8 & 297\dg 07.2\pr & S19\dg 38.6\pr \\
22 & 9 & 312\dg 07.0\pr & S19\dg 38.0\pr & 22 & 10 & 327\dg 06.8\pr & S19\dg 37.5\pr & 22 & 11 & 342\dg 06.7\pr & S19\dg 36.9\pr \\
22 & 12 & 357\dg 06.5\pr & S19\dg 36.3\pr & 22 & 13 & 12\dg 06.3\pr & S19\dg 35.7\pr & 22 & 14 & 27\dg 06.2\pr & S19\dg 35.2\pr \\
22 & 15 & 42\dg 06.0\pr & S19\dg 34.6\pr & 22 & 16 & 57\dg 05.8\pr & S19\dg 34.0\pr & 22 & 17 & 72\dg 05.7\pr & S19\dg 33.4\pr \\
22 & 18 & 87\dg 05.5\pr & S19\dg 32.8\pr & 22 & 19 & 102\dg 05.4\pr & S19\dg 32.3\pr & 22 & 20 & 117\dg 05.2\pr & S19\dg 31.7\pr \\
22 & 21 & 132\dg 05.0\pr & S19\dg 31.1\pr & 22 & 22 & 147\dg 04.9\pr & S19\dg 30.5\pr & 22 & 23 & 162\dg 04.7\pr & S19\dg 29.9\pr \\
23 & 0 & 177\dg 04.5\pr & S19\dg 29.3\pr & 23 & 1 & 192\dg 04.4\pr & S19\dg 28.8\pr & 23 & 2 & 207\dg 04.2\pr & S19\dg 28.2\pr \\
23 & 3 & 222\dg 04.1\pr & S19\dg 27.6\pr & 23 & 4 & 237\dg 03.9\pr & S19\dg 27.0\pr & 23 & 5 & 252\dg 03.7\pr & S19\dg 26.4\pr \\
23 & 6 & 267\dg 03.6\pr & S19\dg 25.8\pr & 23 & 7 & 282\dg 03.4\pr & S19\dg 25.2\pr & 23 & 8 & 297\dg 03.3\pr & S19\dg 24.7\pr \\
23 & 9 & 312\dg 03.1\pr & S19\dg 24.1\pr & 23 & 10 & 327\dg 02.9\pr & S19\dg 23.5\pr & 23 & 11 & 342\dg 02.8\pr & S19\dg 22.9\pr \\
23 & 12 & 357\dg 02.6\pr & S19\dg 22.3\pr & 23 & 13 & 12\dg 02.5\pr & S19\dg 21.7\pr & 23 & 14 & 27\dg 02.3\pr & S19\dg 21.1\pr \\
23 & 15 & 42\dg 02.2\pr & S19\dg 20.5\pr & 23 & 16 & 57\dg 02.0\pr & S19\dg 19.9\pr & 23 & 17 & 72\dg 01.8\pr & S19\dg 19.3\pr \\
23 & 18 & 87\dg 01.7\pr & S19\dg 18.7\pr & 23 & 19 & 102\dg 01.5\pr & S19\dg 18.1\pr & 23 & 20 & 117\dg 01.4\pr & S19\dg 17.5\pr \\
23 & 21 & 132\dg 01.2\pr & S19\dg 16.9\pr & 23 & 22 & 147\dg 01.1\pr & S19\dg 16.3\pr & 23 & 23 & 162\dg 00.9\pr & S19\dg 15.7\pr \\
24 & 0 & 177\dg 00.8\pr & S19\dg 15.1\pr & 24 & 1 & 192\dg 00.6\pr & S19\dg 14.5\pr & 24 & 2 & 207\dg 00.5\pr & S19\dg 13.9\pr \\
24 & 3 & 222\dg 00.3\pr & S19\dg 13.3\pr & 24 & 4 & 237\dg 00.1\pr & S19\dg 12.7\pr & 24 & 5 & 252\dg 00.0\pr & S19\dg 12.1\pr \\
24 & 6 & 266\dg 59.8\pr & S19\dg 11.5\pr & 24 & 7 & 281\dg 59.7\pr & S19\dg 10.9\pr & 24 & 8 & 296\dg 59.5\pr & S19\dg 10.3\pr \\
24 & 9 & 311\dg 59.4\pr & S19\dg 09.7\pr & 24 & 10 & 326\dg 59.2\pr & S19\dg 09.1\pr & 24 & 11 & 341\dg 59.1\pr & S19\dg 08.5\pr \\
24 & 12 & 356\dg 58.9\pr & S19\dg 07.9\pr & 24 & 13 & 11\dg 58.8\pr & S19\dg 07.3\pr & 24 & 14 & 26\dg 58.6\pr & S19\dg 06.7\pr \\
24 & 15 & 41\dg 58.5\pr & S19\dg 06.1\pr & 24 & 16 & 56\dg 58.3\pr & S19\dg 05.5\pr & 24 & 17 & 71\dg 58.2\pr & S19\dg 04.9\pr \\
24 & 18 & 86\dg 58.1\pr & S19\dg 04.3\pr & 24 & 19 & 101\dg 57.9\pr & S19\dg 03.7\pr & 24 & 20 & 116\dg 57.8\pr & S19\dg 03.0\pr \\
24 & 21 & 131\dg 57.6\pr & S19\dg 02.4\pr & 24 & 22 & 146\dg 57.5\pr & S19\dg 01.8\pr & 24 & 23 & 161\dg 57.3\pr & S19\dg 01.2\pr \\
25 & 0 & 176\dg 57.2\pr & S19\dg 00.6\pr & 25 & 1 & 191\dg 57.0\pr & S19\dg 00.0\pr & 25 & 2 & 206\dg 56.9\pr & S18\dg 59.4\pr \\
25 & 3 & 221\dg 56.7\pr & S18\dg 58.8\pr & 25 & 4 & 236\dg 56.6\pr & S18\dg 58.1\pr & 25 & 5 & 251\dg 56.4\pr & S18\dg 57.5\pr \\
25 & 6 & 266\dg 56.3\pr & S18\dg 56.9\pr & 25 & 7 & 281\dg 56.2\pr & S18\dg 56.3\pr & 25 & 8 & 296\dg 56.0\pr & S18\dg 55.7\pr \\
25 & 9 & 311\dg 55.9\pr & S18\dg 55.0\pr & 25 & 10 & 326\dg 55.7\pr & S18\dg 54.4\pr & 25 & 11 & 341\dg 55.6\pr & S18\dg 53.8\pr \\
25 & 12 & 356\dg 55.5\pr & S18\dg 53.2\pr & 25 & 13 & 11\dg 55.3\pr & S18\dg 52.6\pr & 25 & 14 & 26\dg 55.2\pr & S18\dg 51.9\pr \\
25 & 15 & 41\dg 55.0\pr & S18\dg 51.3\pr & 25 & 16 & 56\dg 54.9\pr & S18\dg 50.7\pr & 25 & 17 & 71\dg 54.8\pr & S18\dg 50.1\pr \\
25 & 18 & 86\dg 54.6\pr & S18\dg 49.5\pr & 25 & 19 & 101\dg 54.5\pr & S18\dg 48.8\pr & 25 & 20 & 116\dg 54.3\pr & S18\dg 48.2\pr \\
25 & 21 & 131\dg 54.2\pr & S18\dg 47.6\pr & 25 & 22 & 146\dg 54.1\pr & S18\dg 46.9\pr & 25 & 23 & 161\dg 53.9\pr & S18\dg 46.3\pr \\
26 & 0 & 176\dg 53.8\pr & S18\dg 45.7\pr & 26 & 1 & 191\dg 53.6\pr & S18\dg 45.1\pr & 26 & 2 & 206\dg 53.5\pr & S18\dg 44.4\pr \\
26 & 3 & 221\dg 53.4\pr & S18\dg 43.8\pr & 26 & 4 & 236\dg 53.2\pr & S18\dg 43.2\pr & 26 & 5 & 251\dg 53.1\pr & S18\dg 42.5\pr \\
26 & 6 & 266\dg 53.0\pr & S18\dg 41.9\pr & 26 & 7 & 281\dg 52.8\pr & S18\dg 41.3\pr & 26 & 8 & 296\dg 52.7\pr & S18\dg 40.7\pr \\
26 & 9 & 311\dg 52.6\pr & S18\dg 40.0\pr & 26 & 10 & 326\dg 52.4\pr & S18\dg 39.4\pr & 26 & 11 & 341\dg 52.3\pr & S18\dg 38.8\pr \\
26 & 12 & 356\dg 52.2\pr & S18\dg 38.1\pr & 26 & 13 & 11\dg 52.0\pr & S18\dg 37.5\pr & 26 & 14 & 26\dg 51.9\pr & S18\dg 36.8\pr \\
26 & 15 & 41\dg 51.8\pr & S18\dg 36.2\pr & 26 & 16 & 56\dg 51.6\pr & S18\dg 35.6\pr & 26 & 17 & 71\dg 51.5\pr & S18\dg 34.9\pr \\
26 & 18 & 86\dg 51.4\pr & S18\dg 34.3\pr & 26 & 19 & 101\dg 51.2\pr & S18\dg 33.7\pr & 26 & 20 & 116\dg 51.1\pr & S18\dg 33.0\pr \\
26 & 21 & 131\dg 51.0\pr & S18\dg 32.4\pr & 26 & 22 & 146\dg 50.9\pr & S18\dg 31.7\pr & 26 & 23 & 161\dg 50.7\pr & S18\dg 31.1\pr \\
27 & 0 & 176\dg 50.6\pr & S18\dg 30.5\pr & 27 & 1 & 191\dg 50.5\pr & S18\dg 29.8\pr & 27 & 2 & 206\dg 50.3\pr & S18\dg 29.2\pr \\
27 & 3 & 221\dg 50.2\pr & S18\dg 28.5\pr & 27 & 4 & 236\dg 50.1\pr & S18\dg 27.9\pr & 27 & 5 & 251\dg 50.0\pr & S18\dg 27.2\pr \\
27 & 6 & 266\dg 49.8\pr & S18\dg 26.6\pr & 27 & 7 & 281\dg 49.7\pr & S18\dg 25.9\pr & 27 & 8 & 296\dg 49.6\pr & S18\dg 25.3\pr \\
27 & 9 & 311\dg 49.5\pr & S18\dg 24.7\pr & 27 & 10 & 326\dg 49.3\pr & S18\dg 24.0\pr & 27 & 11 & 341\dg 49.2\pr & S18\dg 23.4\pr \\
27 & 12 & 356\dg 49.1\pr & S18\dg 22.7\pr & 27 & 13 & 11\dg 49.0\pr & S18\dg 22.1\pr & 27 & 14 & 26\dg 48.8\pr & S18\dg 21.4\pr \\
27 & 15 & 41\dg 48.7\pr & S18\dg 20.8\pr & 27 & 16 & 56\dg 48.6\pr & S18\dg 20.1\pr & 27 & 17 & 71\dg 48.5\pr & S18\dg 19.5\pr \\
27 & 18 & 86\dg 48.3\pr & S18\dg 18.8\pr & 27 & 19 & 101\dg 48.2\pr & S18\dg 18.1\pr & 27 & 20 & 116\dg 48.1\pr & S18\dg 17.5\pr \\
27 & 21 & 131\dg 48.0\pr & S18\dg 16.8\pr & 27 & 22 & 146\dg 47.8\pr & S18\dg 16.2\pr & 27 & 23 & 161\dg 47.7\pr & S18\dg 15.5\pr \\
28 & 0 & 176\dg 47.6\pr & S18\dg 14.9\pr & 28 & 1 & 191\dg 47.5\pr & S18\dg 14.2\pr & 28 & 2 & 206\dg 47.4\pr & S18\dg 13.6\pr \\
28 & 3 & 221\dg 47.2\pr & S18\dg 12.9\pr & 28 & 4 & 236\dg 47.1\pr & S18\dg 12.2\pr & 28 & 5 & 251\dg 47.0\pr & S18\dg 11.6\pr \\
28 & 6 & 266\dg 46.9\pr & S18\dg 10.9\pr & 28 & 7 & 281\dg 46.8\pr & S18\dg 10.3\pr & 28 & 8 & 296\dg 46.7\pr & S18\dg 09.6\pr \\
28 & 9 & 311\dg 46.5\pr & S18\dg 08.9\pr & 28 & 10 & 326\dg 46.4\pr & S18\dg 08.3\pr & 28 & 11 & 341\dg 46.3\pr & S18\dg 07.6\pr \\
28 & 12 & 356\dg 46.2\pr & S18\dg 07.0\pr & 28 & 13 & 11\dg 46.1\pr & S18\dg 06.3\pr & 28 & 14 & 26\dg 46.0\pr & S18\dg 05.6\pr \\
28 & 15 & 41\dg 45.8\pr & S18\dg 05.0\pr & 28 & 16 & 56\dg 45.7\pr & S18\dg 04.3\pr & 28 & 17 & 71\dg 45.6\pr & S18\dg 03.6\pr \\
28 & 18 & 86\dg 45.5\pr & S18\dg 03.0\pr & 28 & 19 & 101\dg 45.4\pr & S18\dg 02.3\pr & 28 & 20 & 116\dg 45.3\pr & S18\dg 01.6\pr \\
28 & 21 & 131\dg 45.2\pr & S18\dg 01.0\pr & 28 & 22 & 146\dg 45.0\pr & S18\dg 00.3\pr & 28 & 23 & 161\dg 44.9\pr & S17\dg 59.6\pr \\
29 & 0 & 176\dg 44.8\pr & S17\dg 59.0\pr & 29 & 1 & 191\dg 44.7\pr & S17\dg 58.3\pr & 29 & 2 & 206\dg 44.6\pr & S17\dg 57.6\pr \\
29 & 3 & 221\dg 44.5\pr & S17\dg 57.0\pr & 29 & 4 & 236\dg 44.4\pr & S17\dg 56.3\pr & 29 & 5 & 251\dg 44.3\pr & S17\dg 55.6\pr \\
29 & 6 & 266\dg 44.2\pr & S17\dg 54.9\pr & 29 & 7 & 281\dg 44.0\pr & S17\dg 54.3\pr & 29 & 8 & 296\dg 43.9\pr & S17\dg 53.6\pr \\
29 & 9 & 311\dg 43.8\pr & S17\dg 52.9\pr & 29 & 10 & 326\dg 43.7\pr & S17\dg 52.2\pr & 29 & 11 & 341\dg 43.6\pr & S17\dg 51.6\pr \\
29 & 12 & 356\dg 43.5\pr & S17\dg 50.9\pr & 29 & 13 & 11\dg 43.4\pr & S17\dg 50.2\pr & 29 & 14 & 26\dg 43.3\pr & S17\dg 49.5\pr \\
29 & 15 & 41\dg 43.2\pr & S17\dg 48.9\pr & 29 & 16 & 56\dg 43.1\pr & S17\dg 48.2\pr & 29 & 17 & 71\dg 43.0\pr & S17\dg 47.5\pr \\
29 & 18 & 86\dg 42.9\pr & S17\dg 46.8\pr & 29 & 19 & 101\dg 42.8\pr & S17\dg 46.1\pr & 29 & 20 & 116\dg 42.7\pr & S17\dg 45.5\pr \\
29 & 21 & 131\dg 42.6\pr & S17\dg 44.8\pr & 29 & 22 & 146\dg 42.4\pr & S17\dg 44.1\pr & 29 & 23 & 161\dg 42.3\pr & S17\dg 43.4\pr \\
30 & 0 & 176\dg 42.2\pr & S17\dg 42.7\pr & 30 & 1 & 191\dg 42.1\pr & S17\dg 42.1\pr & 30 & 2 & 206\dg 42.0\pr & S17\dg 41.4\pr \\
30 & 3 & 221\dg 41.9\pr & S17\dg 40.7\pr & 30 & 4 & 236\dg 41.8\pr & S17\dg 40.0\pr & 30 & 5 & 251\dg 41.7\pr & S17\dg 39.3\pr \\
30 & 6 & 266\dg 41.6\pr & S17\dg 38.6\pr & 30 & 7 & 281\dg 41.5\pr & S17\dg 37.9\pr & 30 & 8 & 296\dg 41.4\pr & S17\dg 37.3\pr \\
30 & 9 & 311\dg 41.3\pr & S17\dg 36.6\pr & 30 & 10 & 326\dg 41.2\pr & S17\dg 35.9\pr & 30 & 11 & 341\dg 41.1\pr & S17\dg 35.2\pr \\
30 & 12 & 356\dg 41.0\pr & S17\dg 34.5\pr & 30 & 13 & 11\dg 40.9\pr & S17\dg 33.8\pr & 30 & 14 & 26\dg 40.8\pr & S17\dg 33.1\pr \\
30 & 15 & 41\dg 40.7\pr & S17\dg 32.4\pr & 30 & 16 & 56\dg 40.6\pr & S17\dg 31.7\pr & 30 & 17 & 71\dg 40.5\pr & S17\dg 31.1\pr \\
30 & 18 & 86\dg 40.4\pr & S17\dg 30.4\pr & 30 & 19 & 101\dg 40.3\pr & S17\dg 29.7\pr & 30 & 20 & 116\dg 40.2\pr & S17\dg 29.0\pr \\
30 & 21 & 131\dg 40.1\pr & S17\dg 28.3\pr & 30 & 22 & 146\dg 40.1\pr & S17\dg 27.6\pr & 30 & 23 & 161\dg 40.0\pr & S17\dg 26.9\pr \\
31 & 0 & 176\dg 39.9\pr & S17\dg 26.2\pr & 31 & 1 & 191\dg 39.8\pr & S17\dg 25.5\pr & 31 & 2 & 206\dg 39.7\pr & S17\dg 24.8\pr \\
31 & 3 & 221\dg 39.6\pr & S17\dg 24.1\pr & 31 & 4 & 236\dg 39.5\pr & S17\dg 23.4\pr & 31 & 5 & 251\dg 39.4\pr & S17\dg 22.7\pr \\
31 & 6 & 266\dg 39.3\pr & S17\dg 22.0\pr & 31 & 7 & 281\dg 39.2\pr & S17\dg 21.3\pr & 31 & 8 & 296\dg 39.1\pr & S17\dg 20.6\pr \\
31 & 9 & 311\dg 39.0\pr & S17\dg 19.9\pr & 31 & 10 & 326\dg 38.9\pr & S17\dg 19.2\pr & 31 & 11 & 341\dg 38.8\pr & S17\dg 18.5\pr \\
31 & 12 & 356\dg 38.7\pr & S17\dg 17.8\pr & 31 & 13 & 11\dg 38.7\pr & S17\dg 17.1\pr & 31 & 14 & 26\dg 38.6\pr & S17\dg 16.4\pr \\
31 & 15 & 41\dg 38.5\pr & S17\dg 15.7\pr & 31 & 16 & 56\dg 38.4\pr & S17\dg 15.0\pr & 31 & 17 & 71\dg 38.3\pr & S17\dg 14.3\pr \\
31 & 18 & 86\dg 38.2\pr & S17\dg 13.6\pr & 31 & 19 & 101\dg 38.1\pr & S17\dg 12.9\pr & 31 & 20 & 116\dg 38.0\pr & S17\dg 12.2\pr \\
31 & 21 & 131\dg 37.9\pr & S17\dg 11.5\pr & 31 & 22 & 146\dg 37.9\pr & S17\dg 10.8\pr & 31 & 23 & 161\dg 37.8\pr & S17\dg 10.1\pr \\
\end{longtable}}

{\footnotesize
\begin{longtable}{rrrr|rrrr|rrrr}
\caption{Sun --- GHA and Declination, February 2026 (UT)}\label{sun02}\\
\toprule
\textbf{d} & \textbf{h} & \textbf{GHA} & \textbf{Dec} & \textbf{d} & \textbf{h} & \textbf{GHA} & \textbf{Dec} & \textbf{d} & \textbf{h} & \textbf{GHA} & \textbf{Dec} \\
\midrule
\endfirsthead
\multicolumn{12}{c}{\textit{Sun --- GHA and Declination, February 2026 (UT) (continued)}}\\
\toprule
\textbf{d} & \textbf{h} & \textbf{GHA} & \textbf{Dec} & \textbf{d} & \textbf{h} & \textbf{GHA} & \textbf{Dec} & \textbf{d} & \textbf{h} & \textbf{GHA} & \textbf{Dec} \\
\midrule
\endhead
\midrule\multicolumn{12}{r}{\textit{continued on next page}}\\
\endfoot
\bottomrule
\endlastfoot
1 & 0 & 176\dg 37.7\pr & S17\dg 09.3\pr & 1 & 1 & 191\dg 37.6\pr & S17\dg 08.6\pr & 1 & 2 & 206\dg 37.5\pr & S17\dg 07.9\pr \\
1 & 3 & 221\dg 37.4\pr & S17\dg 07.2\pr & 1 & 4 & 236\dg 37.3\pr & S17\dg 06.5\pr & 1 & 5 & 251\dg 37.3\pr & S17\dg 05.8\pr \\
1 & 6 & 266\dg 37.2\pr & S17\dg 05.1\pr & 1 & 7 & 281\dg 37.1\pr & S17\dg 04.4\pr & 1 & 8 & 296\dg 37.0\pr & S17\dg 03.7\pr \\
1 & 9 & 311\dg 36.9\pr & S17\dg 02.9\pr & 1 & 10 & 326\dg 36.8\pr & S17\dg 02.2\pr & 1 & 11 & 341\dg 36.8\pr & S17\dg 01.5\pr \\
1 & 12 & 356\dg 36.7\pr & S17\dg 00.8\pr & 1 & 13 & 11\dg 36.6\pr & S17\dg 00.1\pr & 1 & 14 & 26\dg 36.5\pr & S16\dg 59.4\pr \\
1 & 15 & 41\dg 36.4\pr & S16\dg 58.7\pr & 1 & 16 & 56\dg 36.3\pr & S16\dg 57.9\pr & 1 & 17 & 71\dg 36.3\pr & S16\dg 57.2\pr \\
1 & 18 & 86\dg 36.2\pr & S16\dg 56.5\pr & 1 & 19 & 101\dg 36.1\pr & S16\dg 55.8\pr & 1 & 20 & 116\dg 36.0\pr & S16\dg 55.1\pr \\
1 & 21 & 131\dg 36.0\pr & S16\dg 54.3\pr & 1 & 22 & 146\dg 35.9\pr & S16\dg 53.6\pr & 1 & 23 & 161\dg 35.8\pr & S16\dg 52.9\pr \\
2 & 0 & 176\dg 35.7\pr & S16\dg 52.2\pr & 2 & 1 & 191\dg 35.6\pr & S16\dg 51.5\pr & 2 & 2 & 206\dg 35.6\pr & S16\dg 50.7\pr \\
2 & 3 & 221\dg 35.5\pr & S16\dg 50.0\pr & 2 & 4 & 236\dg 35.4\pr & S16\dg 49.3\pr & 2 & 5 & 251\dg 35.3\pr & S16\dg 48.6\pr \\
2 & 6 & 266\dg 35.3\pr & S16\dg 47.9\pr & 2 & 7 & 281\dg 35.2\pr & S16\dg 47.1\pr & 2 & 8 & 296\dg 35.1\pr & S16\dg 46.4\pr \\
2 & 9 & 311\dg 35.0\pr & S16\dg 45.7\pr & 2 & 10 & 326\dg 35.0\pr & S16\dg 45.0\pr & 2 & 11 & 341\dg 34.9\pr & S16\dg 44.2\pr \\
2 & 12 & 356\dg 34.8\pr & S16\dg 43.5\pr & 2 & 13 & 11\dg 34.7\pr & S16\dg 42.8\pr & 2 & 14 & 26\dg 34.7\pr & S16\dg 42.0\pr \\
2 & 15 & 41\dg 34.6\pr & S16\dg 41.3\pr & 2 & 16 & 56\dg 34.5\pr & S16\dg 40.6\pr & 2 & 17 & 71\dg 34.4\pr & S16\dg 39.9\pr \\
2 & 18 & 86\dg 34.4\pr & S16\dg 39.1\pr & 2 & 19 & 101\dg 34.3\pr & S16\dg 38.4\pr & 2 & 20 & 116\dg 34.2\pr & S16\dg 37.7\pr \\
2 & 21 & 131\dg 34.2\pr & S16\dg 36.9\pr & 2 & 22 & 146\dg 34.1\pr & S16\dg 36.2\pr & 2 & 23 & 161\dg 34.0\pr & S16\dg 35.5\pr \\
3 & 0 & 176\dg 33.9\pr & S16\dg 34.7\pr & 3 & 1 & 191\dg 33.9\pr & S16\dg 34.0\pr & 3 & 2 & 206\dg 33.8\pr & S16\dg 33.3\pr \\
3 & 3 & 221\dg 33.7\pr & S16\dg 32.5\pr & 3 & 4 & 236\dg 33.7\pr & S16\dg 31.8\pr & 3 & 5 & 251\dg 33.6\pr & S16\dg 31.1\pr \\
3 & 6 & 266\dg 33.5\pr & S16\dg 30.3\pr & 3 & 7 & 281\dg 33.5\pr & S16\dg 29.6\pr & 3 & 8 & 296\dg 33.4\pr & S16\dg 28.9\pr \\
3 & 9 & 311\dg 33.3\pr & S16\dg 28.1\pr & 3 & 10 & 326\dg 33.3\pr & S16\dg 27.4\pr & 3 & 11 & 341\dg 33.2\pr & S16\dg 26.6\pr \\
3 & 12 & 356\dg 33.1\pr & S16\dg 25.9\pr & 3 & 13 & 11\dg 33.1\pr & S16\dg 25.2\pr & 3 & 14 & 26\dg 33.0\pr & S16\dg 24.4\pr \\
3 & 15 & 41\dg 32.9\pr & S16\dg 23.7\pr & 3 & 16 & 56\dg 32.9\pr & S16\dg 22.9\pr & 3 & 17 & 71\dg 32.8\pr & S16\dg 22.2\pr \\
3 & 18 & 86\dg 32.8\pr & S16\dg 21.5\pr & 3 & 19 & 101\dg 32.7\pr & S16\dg 20.7\pr & 3 & 20 & 116\dg 32.6\pr & S16\dg 20.0\pr \\
3 & 21 & 131\dg 32.6\pr & S16\dg 19.2\pr & 3 & 22 & 146\dg 32.5\pr & S16\dg 18.5\pr & 3 & 23 & 161\dg 32.4\pr & S16\dg 17.7\pr \\
4 & 0 & 176\dg 32.4\pr & S16\dg 17.0\pr & 4 & 1 & 191\dg 32.3\pr & S16\dg 16.3\pr & 4 & 2 & 206\dg 32.3\pr & S16\dg 15.5\pr \\
4 & 3 & 221\dg 32.2\pr & S16\dg 14.8\pr & 4 & 4 & 236\dg 32.1\pr & S16\dg 14.0\pr & 4 & 5 & 251\dg 32.1\pr & S16\dg 13.3\pr \\
4 & 6 & 266\dg 32.0\pr & S16\dg 12.5\pr & 4 & 7 & 281\dg 32.0\pr & S16\dg 11.8\pr & 4 & 8 & 296\dg 31.9\pr & S16\dg 11.0\pr \\
4 & 9 & 311\dg 31.9\pr & S16\dg 10.3\pr & 4 & 10 & 326\dg 31.8\pr & S16\dg 09.5\pr & 4 & 11 & 341\dg 31.7\pr & S16\dg 08.8\pr \\
4 & 12 & 356\dg 31.7\pr & S16\dg 08.0\pr & 4 & 13 & 11\dg 31.6\pr & S16\dg 07.3\pr & 4 & 14 & 26\dg 31.6\pr & S16\dg 06.5\pr \\
4 & 15 & 41\dg 31.5\pr & S16\dg 05.8\pr & 4 & 16 & 56\dg 31.5\pr & S16\dg 05.0\pr & 4 & 17 & 71\dg 31.4\pr & S16\dg 04.3\pr \\
4 & 18 & 86\dg 31.3\pr & S16\dg 03.5\pr & 4 & 19 & 101\dg 31.3\pr & S16\dg 02.8\pr & 4 & 20 & 116\dg 31.2\pr & S16\dg 02.0\pr \\
4 & 21 & 131\dg 31.2\pr & S16\dg 01.2\pr & 4 & 22 & 146\dg 31.1\pr & S16\dg 00.5\pr & 4 & 23 & 161\dg 31.1\pr & S15\dg 59.7\pr \\
5 & 0 & 176\dg 31.0\pr & S15\dg 59.0\pr & 5 & 1 & 191\dg 31.0\pr & S15\dg 58.2\pr & 5 & 2 & 206\dg 30.9\pr & S15\dg 57.5\pr \\
5 & 3 & 221\dg 30.9\pr & S15\dg 56.7\pr & 5 & 4 & 236\dg 30.8\pr & S15\dg 55.9\pr & 5 & 5 & 251\dg 30.8\pr & S15\dg 55.2\pr \\
5 & 6 & 266\dg 30.7\pr & S15\dg 54.4\pr & 5 & 7 & 281\dg 30.7\pr & S15\dg 53.7\pr & 5 & 8 & 296\dg 30.6\pr & S15\dg 52.9\pr \\
5 & 9 & 311\dg 30.6\pr & S15\dg 52.1\pr & 5 & 10 & 326\dg 30.5\pr & S15\dg 51.4\pr & 5 & 11 & 341\dg 30.5\pr & S15\dg 50.6\pr \\
5 & 12 & 356\dg 30.4\pr & S15\dg 49.9\pr & 5 & 13 & 11\dg 30.4\pr & S15\dg 49.1\pr & 5 & 14 & 26\dg 30.3\pr & S15\dg 48.3\pr \\
5 & 15 & 41\dg 30.3\pr & S15\dg 47.6\pr & 5 & 16 & 56\dg 30.2\pr & S15\dg 46.8\pr & 5 & 17 & 71\dg 30.2\pr & S15\dg 46.0\pr \\
5 & 18 & 86\dg 30.1\pr & S15\dg 45.3\pr & 5 & 19 & 101\dg 30.1\pr & S15\dg 44.5\pr & 5 & 20 & 116\dg 30.0\pr & S15\dg 43.7\pr \\
5 & 21 & 131\dg 30.0\pr & S15\dg 43.0\pr & 5 & 22 & 146\dg 30.0\pr & S15\dg 42.2\pr & 5 & 23 & 161\dg 29.9\pr & S15\dg 41.4\pr \\
6 & 0 & 176\dg 29.9\pr & S15\dg 40.7\pr & 6 & 1 & 191\dg 29.8\pr & S15\dg 39.9\pr & 6 & 2 & 206\dg 29.8\pr & S15\dg 39.1\pr \\
6 & 3 & 221\dg 29.7\pr & S15\dg 38.4\pr & 6 & 4 & 236\dg 29.7\pr & S15\dg 37.6\pr & 6 & 5 & 251\dg 29.6\pr & S15\dg 36.8\pr \\
6 & 6 & 266\dg 29.6\pr & S15\dg 36.1\pr & 6 & 7 & 281\dg 29.6\pr & S15\dg 35.3\pr & 6 & 8 & 296\dg 29.5\pr & S15\dg 34.5\pr \\
6 & 9 & 311\dg 29.5\pr & S15\dg 33.7\pr & 6 & 10 & 326\dg 29.4\pr & S15\dg 33.0\pr & 6 & 11 & 341\dg 29.4\pr & S15\dg 32.2\pr \\
6 & 12 & 356\dg 29.4\pr & S15\dg 31.4\pr & 6 & 13 & 11\dg 29.3\pr & S15\dg 30.7\pr & 6 & 14 & 26\dg 29.3\pr & S15\dg 29.9\pr \\
6 & 15 & 41\dg 29.2\pr & S15\dg 29.1\pr & 6 & 16 & 56\dg 29.2\pr & S15\dg 28.3\pr & 6 & 17 & 71\dg 29.2\pr & S15\dg 27.6\pr \\
6 & 18 & 86\dg 29.1\pr & S15\dg 26.8\pr & 6 & 19 & 101\dg 29.1\pr & S15\dg 26.0\pr & 6 & 20 & 116\dg 29.1\pr & S15\dg 25.2\pr \\
6 & 21 & 131\dg 29.0\pr & S15\dg 24.4\pr & 6 & 22 & 146\dg 29.0\pr & S15\dg 23.7\pr & 6 & 23 & 161\dg 28.9\pr & S15\dg 22.9\pr \\
7 & 0 & 176\dg 28.9\pr & S15\dg 22.1\pr & 7 & 1 & 191\dg 28.9\pr & S15\dg 21.3\pr & 7 & 2 & 206\dg 28.8\pr & S15\dg 20.6\pr \\
7 & 3 & 221\dg 28.8\pr & S15\dg 19.8\pr & 7 & 4 & 236\dg 28.8\pr & S15\dg 19.0\pr & 7 & 5 & 251\dg 28.7\pr & S15\dg 18.2\pr \\
7 & 6 & 266\dg 28.7\pr & S15\dg 17.4\pr & 7 & 7 & 281\dg 28.7\pr & S15\dg 16.6\pr & 7 & 8 & 296\dg 28.6\pr & S15\dg 15.9\pr \\
7 & 9 & 311\dg 28.6\pr & S15\dg 15.1\pr & 7 & 10 & 326\dg 28.6\pr & S15\dg 14.3\pr & 7 & 11 & 341\dg 28.5\pr & S15\dg 13.5\pr \\
7 & 12 & 356\dg 28.5\pr & S15\dg 12.7\pr & 7 & 13 & 11\dg 28.5\pr & S15\dg 11.9\pr & 7 & 14 & 26\dg 28.4\pr & S15\dg 11.2\pr \\
7 & 15 & 41\dg 28.4\pr & S15\dg 10.4\pr & 7 & 16 & 56\dg 28.4\pr & S15\dg 09.6\pr & 7 & 17 & 71\dg 28.3\pr & S15\dg 08.8\pr \\
7 & 18 & 86\dg 28.3\pr & S15\dg 08.0\pr & 7 & 19 & 101\dg 28.3\pr & S15\dg 07.2\pr & 7 & 20 & 116\dg 28.3\pr & S15\dg 06.4\pr \\
7 & 21 & 131\dg 28.2\pr & S15\dg 05.7\pr & 7 & 22 & 146\dg 28.2\pr & S15\dg 04.9\pr & 7 & 23 & 161\dg 28.2\pr & S15\dg 04.1\pr \\
8 & 0 & 176\dg 28.1\pr & S15\dg 03.3\pr & 8 & 1 & 191\dg 28.1\pr & S15\dg 02.5\pr & 8 & 2 & 206\dg 28.1\pr & S15\dg 01.7\pr \\
8 & 3 & 221\dg 28.1\pr & S15\dg 00.9\pr & 8 & 4 & 236\dg 28.0\pr & S15\dg 00.1\pr & 8 & 5 & 251\dg 28.0\pr & S14\dg 59.3\pr \\
8 & 6 & 266\dg 28.0\pr & S14\dg 58.5\pr & 8 & 7 & 281\dg 28.0\pr & S14\dg 57.7\pr & 8 & 8 & 296\dg 27.9\pr & S14\dg 57.0\pr \\
8 & 9 & 311\dg 27.9\pr & S14\dg 56.2\pr & 8 & 10 & 326\dg 27.9\pr & S14\dg 55.4\pr & 8 & 11 & 341\dg 27.9\pr & S14\dg 54.6\pr \\
8 & 12 & 356\dg 27.8\pr & S14\dg 53.8\pr & 8 & 13 & 11\dg 27.8\pr & S14\dg 53.0\pr & 8 & 14 & 26\dg 27.8\pr & S14\dg 52.2\pr \\
8 & 15 & 41\dg 27.8\pr & S14\dg 51.4\pr & 8 & 16 & 56\dg 27.8\pr & S14\dg 50.6\pr & 8 & 17 & 71\dg 27.7\pr & S14\dg 49.8\pr \\
8 & 18 & 86\dg 27.7\pr & S14\dg 49.0\pr & 8 & 19 & 101\dg 27.7\pr & S14\dg 48.2\pr & 8 & 20 & 116\dg 27.7\pr & S14\dg 47.4\pr \\
8 & 21 & 131\dg 27.6\pr & S14\dg 46.6\pr & 8 & 22 & 146\dg 27.6\pr & S14\dg 45.8\pr & 8 & 23 & 161\dg 27.6\pr & S14\dg 45.0\pr \\
9 & 0 & 176\dg 27.6\pr & S14\dg 44.2\pr & 9 & 1 & 191\dg 27.6\pr & S14\dg 43.4\pr & 9 & 2 & 206\dg 27.6\pr & S14\dg 42.6\pr \\
9 & 3 & 221\dg 27.5\pr & S14\dg 41.8\pr & 9 & 4 & 236\dg 27.5\pr & S14\dg 41.0\pr & 9 & 5 & 251\dg 27.5\pr & S14\dg 40.2\pr \\
9 & 6 & 266\dg 27.5\pr & S14\dg 39.4\pr & 9 & 7 & 281\dg 27.5\pr & S14\dg 38.6\pr & 9 & 8 & 296\dg 27.4\pr & S14\dg 37.8\pr \\
9 & 9 & 311\dg 27.4\pr & S14\dg 37.0\pr & 9 & 10 & 326\dg 27.4\pr & S14\dg 36.2\pr & 9 & 11 & 341\dg 27.4\pr & S14\dg 35.4\pr \\
9 & 12 & 356\dg 27.4\pr & S14\dg 34.6\pr & 9 & 13 & 11\dg 27.4\pr & S14\dg 33.8\pr & 9 & 14 & 26\dg 27.4\pr & S14\dg 33.0\pr \\
9 & 15 & 41\dg 27.3\pr & S14\dg 32.2\pr & 9 & 16 & 56\dg 27.3\pr & S14\dg 31.3\pr & 9 & 17 & 71\dg 27.3\pr & S14\dg 30.5\pr \\
9 & 18 & 86\dg 27.3\pr & S14\dg 29.7\pr & 9 & 19 & 101\dg 27.3\pr & S14\dg 28.9\pr & 9 & 20 & 116\dg 27.3\pr & S14\dg 28.1\pr \\
9 & 21 & 131\dg 27.3\pr & S14\dg 27.3\pr & 9 & 22 & 146\dg 27.2\pr & S14\dg 26.5\pr & 9 & 23 & 161\dg 27.2\pr & S14\dg 25.7\pr \\
10 & 0 & 176\dg 27.2\pr & S14\dg 24.9\pr & 10 & 1 & 191\dg 27.2\pr & S14\dg 24.1\pr & 10 & 2 & 206\dg 27.2\pr & S14\dg 23.3\pr \\
10 & 3 & 221\dg 27.2\pr & S14\dg 22.4\pr & 10 & 4 & 236\dg 27.2\pr & S14\dg 21.6\pr & 10 & 5 & 251\dg 27.2\pr & S14\dg 20.8\pr \\
10 & 6 & 266\dg 27.2\pr & S14\dg 20.0\pr & 10 & 7 & 281\dg 27.2\pr & S14\dg 19.2\pr & 10 & 8 & 296\dg 27.1\pr & S14\dg 18.4\pr \\
10 & 9 & 311\dg 27.1\pr & S14\dg 17.6\pr & 10 & 10 & 326\dg 27.1\pr & S14\dg 16.7\pr & 10 & 11 & 341\dg 27.1\pr & S14\dg 15.9\pr \\
10 & 12 & 356\dg 27.1\pr & S14\dg 15.1\pr & 10 & 13 & 11\dg 27.1\pr & S14\dg 14.3\pr & 10 & 14 & 26\dg 27.1\pr & S14\dg 13.5\pr \\
10 & 15 & 41\dg 27.1\pr & S14\dg 12.7\pr & 10 & 16 & 56\dg 27.1\pr & S14\dg 11.9\pr & 10 & 17 & 71\dg 27.1\pr & S14\dg 11.0\pr \\
10 & 18 & 86\dg 27.1\pr & S14\dg 10.2\pr & 10 & 19 & 101\dg 27.1\pr & S14\dg 09.4\pr & 10 & 20 & 116\dg 27.1\pr & S14\dg 08.6\pr \\
10 & 21 & 131\dg 27.1\pr & S14\dg 07.8\pr & 10 & 22 & 146\dg 27.1\pr & S14\dg 06.9\pr & 10 & 23 & 161\dg 27.1\pr & S14\dg 06.1\pr \\
11 & 0 & 176\dg 27.1\pr & S14\dg 05.3\pr & 11 & 1 & 191\dg 27.0\pr & S14\dg 04.5\pr & 11 & 2 & 206\dg 27.0\pr & S14\dg 03.7\pr \\
11 & 3 & 221\dg 27.0\pr & S14\dg 02.8\pr & 11 & 4 & 236\dg 27.0\pr & S14\dg 02.0\pr & 11 & 5 & 251\dg 27.0\pr & S14\dg 01.2\pr \\
11 & 6 & 266\dg 27.0\pr & S14\dg 00.4\pr & 11 & 7 & 281\dg 27.0\pr & S13\dg 59.6\pr & 11 & 8 & 296\dg 27.0\pr & S13\dg 58.7\pr \\
11 & 9 & 311\dg 27.0\pr & S13\dg 57.9\pr & 11 & 10 & 326\dg 27.0\pr & S13\dg 57.1\pr & 11 & 11 & 341\dg 27.0\pr & S13\dg 56.3\pr \\
11 & 12 & 356\dg 27.0\pr & S13\dg 55.4\pr & 11 & 13 & 11\dg 27.0\pr & S13\dg 54.6\pr & 11 & 14 & 26\dg 27.0\pr & S13\dg 53.8\pr \\
11 & 15 & 41\dg 27.0\pr & S13\dg 53.0\pr & 11 & 16 & 56\dg 27.0\pr & S13\dg 52.1\pr & 11 & 17 & 71\dg 27.0\pr & S13\dg 51.3\pr \\
11 & 18 & 86\dg 27.0\pr & S13\dg 50.5\pr & 11 & 19 & 101\dg 27.1\pr & S13\dg 49.6\pr & 11 & 20 & 116\dg 27.1\pr & S13\dg 48.8\pr \\
11 & 21 & 131\dg 27.1\pr & S13\dg 48.0\pr & 11 & 22 & 146\dg 27.1\pr & S13\dg 47.2\pr & 11 & 23 & 161\dg 27.1\pr & S13\dg 46.3\pr \\
12 & 0 & 176\dg 27.1\pr & S13\dg 45.5\pr & 12 & 1 & 191\dg 27.1\pr & S13\dg 44.7\pr & 12 & 2 & 206\dg 27.1\pr & S13\dg 43.8\pr \\
12 & 3 & 221\dg 27.1\pr & S13\dg 43.0\pr & 12 & 4 & 236\dg 27.1\pr & S13\dg 42.2\pr & 12 & 5 & 251\dg 27.1\pr & S13\dg 41.3\pr \\
12 & 6 & 266\dg 27.1\pr & S13\dg 40.5\pr & 12 & 7 & 281\dg 27.1\pr & S13\dg 39.7\pr & 12 & 8 & 296\dg 27.1\pr & S13\dg 38.8\pr \\
12 & 9 & 311\dg 27.1\pr & S13\dg 38.0\pr & 12 & 10 & 326\dg 27.1\pr & S13\dg 37.2\pr & 12 & 11 & 341\dg 27.1\pr & S13\dg 36.3\pr \\
12 & 12 & 356\dg 27.2\pr & S13\dg 35.5\pr & 12 & 13 & 11\dg 27.2\pr & S13\dg 34.7\pr & 12 & 14 & 26\dg 27.2\pr & S13\dg 33.8\pr \\
12 & 15 & 41\dg 27.2\pr & S13\dg 33.0\pr & 12 & 16 & 56\dg 27.2\pr & S13\dg 32.2\pr & 12 & 17 & 71\dg 27.2\pr & S13\dg 31.3\pr \\
12 & 18 & 86\dg 27.2\pr & S13\dg 30.5\pr & 12 & 19 & 101\dg 27.2\pr & S13\dg 29.7\pr & 12 & 20 & 116\dg 27.2\pr & S13\dg 28.8\pr \\
12 & 21 & 131\dg 27.2\pr & S13\dg 28.0\pr & 12 & 22 & 146\dg 27.3\pr & S13\dg 27.1\pr & 12 & 23 & 161\dg 27.3\pr & S13\dg 26.3\pr \\
13 & 0 & 176\dg 27.3\pr & S13\dg 25.5\pr & 13 & 1 & 191\dg 27.3\pr & S13\dg 24.6\pr & 13 & 2 & 206\dg 27.3\pr & S13\dg 23.8\pr \\
13 & 3 & 221\dg 27.3\pr & S13\dg 23.0\pr & 13 & 4 & 236\dg 27.3\pr & S13\dg 22.1\pr & 13 & 5 & 251\dg 27.4\pr & S13\dg 21.3\pr \\
13 & 6 & 266\dg 27.4\pr & S13\dg 20.4\pr & 13 & 7 & 281\dg 27.4\pr & S13\dg 19.6\pr & 13 & 8 & 296\dg 27.4\pr & S13\dg 18.7\pr \\
13 & 9 & 311\dg 27.4\pr & S13\dg 17.9\pr & 13 & 10 & 326\dg 27.4\pr & S13\dg 17.1\pr & 13 & 11 & 341\dg 27.4\pr & S13\dg 16.2\pr \\
13 & 12 & 356\dg 27.5\pr & S13\dg 15.4\pr & 13 & 13 & 11\dg 27.5\pr & S13\dg 14.5\pr & 13 & 14 & 26\dg 27.5\pr & S13\dg 13.7\pr \\
13 & 15 & 41\dg 27.5\pr & S13\dg 12.8\pr & 13 & 16 & 56\dg 27.5\pr & S13\dg 12.0\pr & 13 & 17 & 71\dg 27.5\pr & S13\dg 11.1\pr \\
13 & 18 & 86\dg 27.6\pr & S13\dg 10.3\pr & 13 & 19 & 101\dg 27.6\pr & S13\dg 09.5\pr & 13 & 20 & 116\dg 27.6\pr & S13\dg 08.6\pr \\
13 & 21 & 131\dg 27.6\pr & S13\dg 07.8\pr & 13 & 22 & 146\dg 27.6\pr & S13\dg 06.9\pr & 13 & 23 & 161\dg 27.7\pr & S13\dg 06.1\pr \\
14 & 0 & 176\dg 27.7\pr & S13\dg 05.2\pr & 14 & 1 & 191\dg 27.7\pr & S13\dg 04.4\pr & 14 & 2 & 206\dg 27.7\pr & S13\dg 03.5\pr \\
14 & 3 & 221\dg 27.7\pr & S13\dg 02.7\pr & 14 & 4 & 236\dg 27.8\pr & S13\dg 01.8\pr & 14 & 5 & 251\dg 27.8\pr & S13\dg 01.0\pr \\
14 & 6 & 266\dg 27.8\pr & S13\dg 00.1\pr & 14 & 7 & 281\dg 27.8\pr & S12\dg 59.3\pr & 14 & 8 & 296\dg 27.9\pr & S12\dg 58.4\pr \\
14 & 9 & 311\dg 27.9\pr & S12\dg 57.6\pr & 14 & 10 & 326\dg 27.9\pr & S12\dg 56.7\pr & 14 & 11 & 341\dg 27.9\pr & S12\dg 55.9\pr \\
14 & 12 & 356\dg 27.9\pr & S12\dg 55.0\pr & 14 & 13 & 11\dg 28.0\pr & S12\dg 54.2\pr & 14 & 14 & 26\dg 28.0\pr & S12\dg 53.3\pr \\
14 & 15 & 41\dg 28.0\pr & S12\dg 52.5\pr & 14 & 16 & 56\dg 28.0\pr & S12\dg 51.6\pr & 14 & 17 & 71\dg 28.1\pr & S12\dg 50.7\pr \\
14 & 18 & 86\dg 28.1\pr & S12\dg 49.9\pr & 14 & 19 & 101\dg 28.1\pr & S12\dg 49.0\pr & 14 & 20 & 116\dg 28.2\pr & S12\dg 48.2\pr \\
14 & 21 & 131\dg 28.2\pr & S12\dg 47.3\pr & 14 & 22 & 146\dg 28.2\pr & S12\dg 46.5\pr & 14 & 23 & 161\dg 28.2\pr & S12\dg 45.6\pr \\
15 & 0 & 176\dg 28.3\pr & S12\dg 44.8\pr & 15 & 1 & 191\dg 28.3\pr & S12\dg 43.9\pr & 15 & 2 & 206\dg 28.3\pr & S12\dg 43.0\pr \\
15 & 3 & 221\dg 28.3\pr & S12\dg 42.2\pr & 15 & 4 & 236\dg 28.4\pr & S12\dg 41.3\pr & 15 & 5 & 251\dg 28.4\pr & S12\dg 40.5\pr \\
15 & 6 & 266\dg 28.4\pr & S12\dg 39.6\pr & 15 & 7 & 281\dg 28.5\pr & S12\dg 38.8\pr & 15 & 8 & 296\dg 28.5\pr & S12\dg 37.9\pr \\
15 & 9 & 311\dg 28.5\pr & S12\dg 37.0\pr & 15 & 10 & 326\dg 28.6\pr & S12\dg 36.2\pr & 15 & 11 & 341\dg 28.6\pr & S12\dg 35.3\pr \\
15 & 12 & 356\dg 28.6\pr & S12\dg 34.4\pr & 15 & 13 & 11\dg 28.7\pr & S12\dg 33.6\pr & 15 & 14 & 26\dg 28.7\pr & S12\dg 32.7\pr \\
15 & 15 & 41\dg 28.7\pr & S12\dg 31.9\pr & 15 & 16 & 56\dg 28.7\pr & S12\dg 31.0\pr & 15 & 17 & 71\dg 28.8\pr & S12\dg 30.1\pr \\
15 & 18 & 86\dg 28.8\pr & S12\dg 29.3\pr & 15 & 19 & 101\dg 28.9\pr & S12\dg 28.4\pr & 15 & 20 & 116\dg 28.9\pr & S12\dg 27.5\pr \\
15 & 21 & 131\dg 28.9\pr & S12\dg 26.7\pr & 15 & 22 & 146\dg 29.0\pr & S12\dg 25.8\pr & 15 & 23 & 161\dg 29.0\pr & S12\dg 25.0\pr \\
16 & 0 & 176\dg 29.0\pr & S12\dg 24.1\pr & 16 & 1 & 191\dg 29.1\pr & S12\dg 23.2\pr & 16 & 2 & 206\dg 29.1\pr & S12\dg 22.4\pr \\
16 & 3 & 221\dg 29.1\pr & S12\dg 21.5\pr & 16 & 4 & 236\dg 29.2\pr & S12\dg 20.6\pr & 16 & 5 & 251\dg 29.2\pr & S12\dg 19.8\pr \\
16 & 6 & 266\dg 29.2\pr & S12\dg 18.9\pr & 16 & 7 & 281\dg 29.3\pr & S12\dg 18.0\pr & 16 & 8 & 296\dg 29.3\pr & S12\dg 17.2\pr \\
16 & 9 & 311\dg 29.4\pr & S12\dg 16.3\pr & 16 & 10 & 326\dg 29.4\pr & S12\dg 15.4\pr & 16 & 11 & 341\dg 29.4\pr & S12\dg 14.5\pr \\
16 & 12 & 356\dg 29.5\pr & S12\dg 13.7\pr & 16 & 13 & 11\dg 29.5\pr & S12\dg 12.8\pr & 16 & 14 & 26\dg 29.6\pr & S12\dg 11.9\pr \\
16 & 15 & 41\dg 29.6\pr & S12\dg 11.1\pr & 16 & 16 & 56\dg 29.6\pr & S12\dg 10.2\pr & 16 & 17 & 71\dg 29.7\pr & S12\dg 09.3\pr \\
16 & 18 & 86\dg 29.7\pr & S12\dg 08.5\pr & 16 & 19 & 101\dg 29.8\pr & S12\dg 07.6\pr & 16 & 20 & 116\dg 29.8\pr & S12\dg 06.7\pr \\
16 & 21 & 131\dg 29.8\pr & S12\dg 05.8\pr & 16 & 22 & 146\dg 29.9\pr & S12\dg 05.0\pr & 16 & 23 & 161\dg 29.9\pr & S12\dg 04.1\pr \\
17 & 0 & 176\dg 30.0\pr & S12\dg 03.2\pr & 17 & 1 & 191\dg 30.0\pr & S12\dg 02.3\pr & 17 & 2 & 206\dg 30.1\pr & S12\dg 01.5\pr \\
17 & 3 & 221\dg 30.1\pr & S12\dg 00.6\pr & 17 & 4 & 236\dg 30.1\pr & S11\dg 59.7\pr & 17 & 5 & 251\dg 30.2\pr & S11\dg 58.9\pr \\
17 & 6 & 266\dg 30.2\pr & S11\dg 58.0\pr & 17 & 7 & 281\dg 30.3\pr & S11\dg 57.1\pr & 17 & 8 & 296\dg 30.3\pr & S11\dg 56.2\pr \\
17 & 9 & 311\dg 30.4\pr & S11\dg 55.3\pr & 17 & 10 & 326\dg 30.4\pr & S11\dg 54.5\pr & 17 & 11 & 341\dg 30.5\pr & S11\dg 53.6\pr \\
17 & 12 & 356\dg 30.5\pr & S11\dg 52.7\pr & 17 & 13 & 11\dg 30.5\pr & S11\dg 51.8\pr & 17 & 14 & 26\dg 30.6\pr & S11\dg 51.0\pr \\
17 & 15 & 41\dg 30.6\pr & S11\dg 50.1\pr & 17 & 16 & 56\dg 30.7\pr & S11\dg 49.2\pr & 17 & 17 & 71\dg 30.7\pr & S11\dg 48.3\pr \\
17 & 18 & 86\dg 30.8\pr & S11\dg 47.4\pr & 17 & 19 & 101\dg 30.8\pr & S11\dg 46.6\pr & 17 & 20 & 116\dg 30.9\pr & S11\dg 45.7\pr \\
17 & 21 & 131\dg 30.9\pr & S11\dg 44.8\pr & 17 & 22 & 146\dg 31.0\pr & S11\dg 43.9\pr & 17 & 23 & 161\dg 31.0\pr & S11\dg 43.0\pr \\
18 & 0 & 176\dg 31.1\pr & S11\dg 42.2\pr & 18 & 1 & 191\dg 31.1\pr & S11\dg 41.3\pr & 18 & 2 & 206\dg 31.2\pr & S11\dg 40.4\pr \\
18 & 3 & 221\dg 31.2\pr & S11\dg 39.5\pr & 18 & 4 & 236\dg 31.3\pr & S11\dg 38.6\pr & 18 & 5 & 251\dg 31.3\pr & S11\dg 37.8\pr \\
18 & 6 & 266\dg 31.4\pr & S11\dg 36.9\pr & 18 & 7 & 281\dg 31.4\pr & S11\dg 36.0\pr & 18 & 8 & 296\dg 31.5\pr & S11\dg 35.1\pr \\
18 & 9 & 311\dg 31.5\pr & S11\dg 34.2\pr & 18 & 10 & 326\dg 31.6\pr & S11\dg 33.3\pr & 18 & 11 & 341\dg 31.7\pr & S11\dg 32.4\pr \\
18 & 12 & 356\dg 31.7\pr & S11\dg 31.6\pr & 18 & 13 & 11\dg 31.8\pr & S11\dg 30.7\pr & 18 & 14 & 26\dg 31.8\pr & S11\dg 29.8\pr \\
18 & 15 & 41\dg 31.9\pr & S11\dg 28.9\pr & 18 & 16 & 56\dg 31.9\pr & S11\dg 28.0\pr & 18 & 17 & 71\dg 32.0\pr & S11\dg 27.1\pr \\
18 & 18 & 86\dg 32.0\pr & S11\dg 26.2\pr & 18 & 19 & 101\dg 32.1\pr & S11\dg 25.4\pr & 18 & 20 & 116\dg 32.1\pr & S11\dg 24.5\pr \\
18 & 21 & 131\dg 32.2\pr & S11\dg 23.6\pr & 18 & 22 & 146\dg 32.3\pr & S11\dg 22.7\pr & 18 & 23 & 161\dg 32.3\pr & S11\dg 21.8\pr \\
19 & 0 & 176\dg 32.4\pr & S11\dg 20.9\pr & 19 & 1 & 191\dg 32.4\pr & S11\dg 20.0\pr & 19 & 2 & 206\dg 32.5\pr & S11\dg 19.1\pr \\
19 & 3 & 221\dg 32.5\pr & S11\dg 18.3\pr & 19 & 4 & 236\dg 32.6\pr & S11\dg 17.4\pr & 19 & 5 & 251\dg 32.7\pr & S11\dg 16.5\pr \\
19 & 6 & 266\dg 32.7\pr & S11\dg 15.6\pr & 19 & 7 & 281\dg 32.8\pr & S11\dg 14.7\pr & 19 & 8 & 296\dg 32.8\pr & S11\dg 13.8\pr \\
19 & 9 & 311\dg 32.9\pr & S11\dg 12.9\pr & 19 & 10 & 326\dg 33.0\pr & S11\dg 12.0\pr & 19 & 11 & 341\dg 33.0\pr & S11\dg 11.1\pr \\
19 & 12 & 356\dg 33.1\pr & S11\dg 10.2\pr & 19 & 13 & 11\dg 33.1\pr & S11\dg 09.3\pr & 19 & 14 & 26\dg 33.2\pr & S11\dg 08.4\pr \\
19 & 15 & 41\dg 33.3\pr & S11\dg 07.6\pr & 19 & 16 & 56\dg 33.3\pr & S11\dg 06.7\pr & 19 & 17 & 71\dg 33.4\pr & S11\dg 05.8\pr \\
19 & 18 & 86\dg 33.4\pr & S11\dg 04.9\pr & 19 & 19 & 101\dg 33.5\pr & S11\dg 04.0\pr & 19 & 20 & 116\dg 33.6\pr & S11\dg 03.1\pr \\
19 & 21 & 131\dg 33.6\pr & S11\dg 02.2\pr & 19 & 22 & 146\dg 33.7\pr & S11\dg 01.3\pr & 19 & 23 & 161\dg 33.8\pr & S11\dg 00.4\pr \\
20 & 0 & 176\dg 33.8\pr & S10\dg 59.5\pr & 20 & 1 & 191\dg 33.9\pr & S10\dg 58.6\pr & 20 & 2 & 206\dg 34.0\pr & S10\dg 57.7\pr \\
20 & 3 & 221\dg 34.0\pr & S10\dg 56.8\pr & 20 & 4 & 236\dg 34.1\pr & S10\dg 55.9\pr & 20 & 5 & 251\dg 34.2\pr & S10\dg 55.0\pr \\
20 & 6 & 266\dg 34.2\pr & S10\dg 54.1\pr & 20 & 7 & 281\dg 34.3\pr & S10\dg 53.2\pr & 20 & 8 & 296\dg 34.3\pr & S10\dg 52.3\pr \\
20 & 9 & 311\dg 34.4\pr & S10\dg 51.4\pr & 20 & 10 & 326\dg 34.5\pr & S10\dg 50.5\pr & 20 & 11 & 341\dg 34.5\pr & S10\dg 49.6\pr \\
20 & 12 & 356\dg 34.6\pr & S10\dg 48.7\pr & 20 & 13 & 11\dg 34.7\pr & S10\dg 47.8\pr & 20 & 14 & 26\dg 34.8\pr & S10\dg 46.9\pr \\
20 & 15 & 41\dg 34.8\pr & S10\dg 46.0\pr & 20 & 16 & 56\dg 34.9\pr & S10\dg 45.1\pr & 20 & 17 & 71\dg 35.0\pr & S10\dg 44.2\pr \\
20 & 18 & 86\dg 35.0\pr & S10\dg 43.3\pr & 20 & 19 & 101\dg 35.1\pr & S10\dg 42.4\pr & 20 & 20 & 116\dg 35.2\pr & S10\dg 41.5\pr \\
20 & 21 & 131\dg 35.2\pr & S10\dg 40.6\pr & 20 & 22 & 146\dg 35.3\pr & S10\dg 39.7\pr & 20 & 23 & 161\dg 35.4\pr & S10\dg 38.8\pr \\
21 & 0 & 176\dg 35.4\pr & S10\dg 37.9\pr & 21 & 1 & 191\dg 35.5\pr & S10\dg 37.0\pr & 21 & 2 & 206\dg 35.6\pr & S10\dg 36.1\pr \\
21 & 3 & 221\dg 35.7\pr & S10\dg 35.2\pr & 21 & 4 & 236\dg 35.7\pr & S10\dg 34.3\pr & 21 & 5 & 251\dg 35.8\pr & S10\dg 33.4\pr \\
21 & 6 & 266\dg 35.9\pr & S10\dg 32.5\pr & 21 & 7 & 281\dg 36.0\pr & S10\dg 31.6\pr & 21 & 8 & 296\dg 36.0\pr & S10\dg 30.7\pr \\
21 & 9 & 311\dg 36.1\pr & S10\dg 29.8\pr & 21 & 10 & 326\dg 36.2\pr & S10\dg 28.9\pr & 21 & 11 & 341\dg 36.2\pr & S10\dg 28.0\pr \\
21 & 12 & 356\dg 36.3\pr & S10\dg 27.0\pr & 21 & 13 & 11\dg 36.4\pr & S10\dg 26.1\pr & 21 & 14 & 26\dg 36.5\pr & S10\dg 25.2\pr \\
21 & 15 & 41\dg 36.5\pr & S10\dg 24.3\pr & 21 & 16 & 56\dg 36.6\pr & S10\dg 23.4\pr & 21 & 17 & 71\dg 36.7\pr & S10\dg 22.5\pr \\
21 & 18 & 86\dg 36.8\pr & S10\dg 21.6\pr & 21 & 19 & 101\dg 36.8\pr & S10\dg 20.7\pr & 21 & 20 & 116\dg 36.9\pr & S10\dg 19.8\pr \\
21 & 21 & 131\dg 37.0\pr & S10\dg 18.9\pr & 21 & 22 & 146\dg 37.1\pr & S10\dg 18.0\pr & 21 & 23 & 161\dg 37.2\pr & S10\dg 17.1\pr \\
22 & 0 & 176\dg 37.2\pr & S10\dg 16.1\pr & 22 & 1 & 191\dg 37.3\pr & S10\dg 15.2\pr & 22 & 2 & 206\dg 37.4\pr & S10\dg 14.3\pr \\
22 & 3 & 221\dg 37.5\pr & S10\dg 13.4\pr & 22 & 4 & 236\dg 37.5\pr & S10\dg 12.5\pr & 22 & 5 & 251\dg 37.6\pr & S10\dg 11.6\pr \\
22 & 6 & 266\dg 37.7\pr & S10\dg 10.7\pr & 22 & 7 & 281\dg 37.8\pr & S10\dg 09.8\pr & 22 & 8 & 296\dg 37.9\pr & S10\dg 08.9\pr \\
22 & 9 & 311\dg 37.9\pr & S10\dg 07.9\pr & 22 & 10 & 326\dg 38.0\pr & S10\dg 07.0\pr & 22 & 11 & 341\dg 38.1\pr & S10\dg 06.1\pr \\
22 & 12 & 356\dg 38.2\pr & S10\dg 05.2\pr & 22 & 13 & 11\dg 38.3\pr & S10\dg 04.3\pr & 22 & 14 & 26\dg 38.3\pr & S10\dg 03.4\pr \\
22 & 15 & 41\dg 38.4\pr & S10\dg 02.5\pr & 22 & 16 & 56\dg 38.5\pr & S10\dg 01.6\pr & 22 & 17 & 71\dg 38.6\pr & S10\dg 00.6\pr \\
22 & 18 & 86\dg 38.7\pr & S9\dg 59.7\pr & 22 & 19 & 101\dg 38.8\pr & S9\dg 58.8\pr & 22 & 20 & 116\dg 38.8\pr & S9\dg 57.9\pr \\
22 & 21 & 131\dg 38.9\pr & S9\dg 57.0\pr & 22 & 22 & 146\dg 39.0\pr & S9\dg 56.1\pr & 22 & 23 & 161\dg 39.1\pr & S9\dg 55.1\pr \\
23 & 0 & 176\dg 39.2\pr & S9\dg 54.2\pr & 23 & 1 & 191\dg 39.3\pr & S9\dg 53.3\pr & 23 & 2 & 206\dg 39.3\pr & S9\dg 52.4\pr \\
23 & 3 & 221\dg 39.4\pr & S9\dg 51.5\pr & 23 & 4 & 236\dg 39.5\pr & S9\dg 50.6\pr & 23 & 5 & 251\dg 39.6\pr & S9\dg 49.6\pr \\
23 & 6 & 266\dg 39.7\pr & S9\dg 48.7\pr & 23 & 7 & 281\dg 39.8\pr & S9\dg 47.8\pr & 23 & 8 & 296\dg 39.9\pr & S9\dg 46.9\pr \\
23 & 9 & 311\dg 39.9\pr & S9\dg 46.0\pr & 23 & 10 & 326\dg 40.0\pr & S9\dg 45.1\pr & 23 & 11 & 341\dg 40.1\pr & S9\dg 44.1\pr \\
23 & 12 & 356\dg 40.2\pr & S9\dg 43.2\pr & 23 & 13 & 11\dg 40.3\pr & S9\dg 42.3\pr & 23 & 14 & 26\dg 40.4\pr & S9\dg 41.4\pr \\
23 & 15 & 41\dg 40.5\pr & S9\dg 40.5\pr & 23 & 16 & 56\dg 40.5\pr & S9\dg 39.5\pr & 23 & 17 & 71\dg 40.6\pr & S9\dg 38.6\pr \\
23 & 18 & 86\dg 40.7\pr & S9\dg 37.7\pr & 23 & 19 & 101\dg 40.8\pr & S9\dg 36.8\pr & 23 & 20 & 116\dg 40.9\pr & S9\dg 35.8\pr \\
23 & 21 & 131\dg 41.0\pr & S9\dg 34.9\pr & 23 & 22 & 146\dg 41.1\pr & S9\dg 34.0\pr & 23 & 23 & 161\dg 41.2\pr & S9\dg 33.1\pr \\
24 & 0 & 176\dg 41.3\pr & S9\dg 32.2\pr & 24 & 1 & 191\dg 41.4\pr & S9\dg 31.2\pr & 24 & 2 & 206\dg 41.4\pr & S9\dg 30.3\pr \\
24 & 3 & 221\dg 41.5\pr & S9\dg 29.4\pr & 24 & 4 & 236\dg 41.6\pr & S9\dg 28.5\pr & 24 & 5 & 251\dg 41.7\pr & S9\dg 27.5\pr \\
24 & 6 & 266\dg 41.8\pr & S9\dg 26.6\pr & 24 & 7 & 281\dg 41.9\pr & S9\dg 25.7\pr & 24 & 8 & 296\dg 42.0\pr & S9\dg 24.8\pr \\
24 & 9 & 311\dg 42.1\pr & S9\dg 23.8\pr & 24 & 10 & 326\dg 42.2\pr & S9\dg 22.9\pr & 24 & 11 & 341\dg 42.3\pr & S9\dg 22.0\pr \\
24 & 12 & 356\dg 42.4\pr & S9\dg 21.1\pr & 24 & 13 & 11\dg 42.5\pr & S9\dg 20.1\pr & 24 & 14 & 26\dg 42.6\pr & S9\dg 19.2\pr \\
24 & 15 & 41\dg 42.6\pr & S9\dg 18.3\pr & 24 & 16 & 56\dg 42.7\pr & S9\dg 17.4\pr & 24 & 17 & 71\dg 42.8\pr & S9\dg 16.4\pr \\
24 & 18 & 86\dg 42.9\pr & S9\dg 15.5\pr & 24 & 19 & 101\dg 43.0\pr & S9\dg 14.6\pr & 24 & 20 & 116\dg 43.1\pr & S9\dg 13.7\pr \\
24 & 21 & 131\dg 43.2\pr & S9\dg 12.7\pr & 24 & 22 & 146\dg 43.3\pr & S9\dg 11.8\pr & 24 & 23 & 161\dg 43.4\pr & S9\dg 10.9\pr \\
25 & 0 & 176\dg 43.5\pr & S9\dg 10.0\pr & 25 & 1 & 191\dg 43.6\pr & S9\dg 09.0\pr & 25 & 2 & 206\dg 43.7\pr & S9\dg 08.1\pr \\
25 & 3 & 221\dg 43.8\pr & S9\dg 07.2\pr & 25 & 4 & 236\dg 43.9\pr & S9\dg 06.2\pr & 25 & 5 & 251\dg 44.0\pr & S9\dg 05.3\pr \\
25 & 6 & 266\dg 44.1\pr & S9\dg 04.4\pr & 25 & 7 & 281\dg 44.2\pr & S9\dg 03.4\pr & 25 & 8 & 296\dg 44.3\pr & S9\dg 02.5\pr \\
25 & 9 & 311\dg 44.4\pr & S9\dg 01.6\pr & 25 & 10 & 326\dg 44.5\pr & S9\dg 00.7\pr & 25 & 11 & 341\dg 44.6\pr & S8\dg 59.7\pr \\
25 & 12 & 356\dg 44.7\pr & S8\dg 58.8\pr & 25 & 13 & 11\dg 44.8\pr & S8\dg 57.9\pr & 25 & 14 & 26\dg 44.9\pr & S8\dg 56.9\pr \\
25 & 15 & 41\dg 45.0\pr & S8\dg 56.0\pr & 25 & 16 & 56\dg 45.1\pr & S8\dg 55.1\pr & 25 & 17 & 71\dg 45.2\pr & S8\dg 54.1\pr \\
25 & 18 & 86\dg 45.3\pr & S8\dg 53.2\pr & 25 & 19 & 101\dg 45.4\pr & S8\dg 52.3\pr & 25 & 20 & 116\dg 45.5\pr & S8\dg 51.3\pr \\
25 & 21 & 131\dg 45.6\pr & S8\dg 50.4\pr & 25 & 22 & 146\dg 45.7\pr & S8\dg 49.5\pr & 25 & 23 & 161\dg 45.8\pr & S8\dg 48.5\pr \\
26 & 0 & 176\dg 45.9\pr & S8\dg 47.6\pr & 26 & 1 & 191\dg 46.0\pr & S8\dg 46.7\pr & 26 & 2 & 206\dg 46.1\pr & S8\dg 45.7\pr \\
26 & 3 & 221\dg 46.2\pr & S8\dg 44.8\pr & 26 & 4 & 236\dg 46.3\pr & S8\dg 43.9\pr & 26 & 5 & 251\dg 46.4\pr & S8\dg 42.9\pr \\
26 & 6 & 266\dg 46.5\pr & S8\dg 42.0\pr & 26 & 7 & 281\dg 46.6\pr & S8\dg 41.1\pr & 26 & 8 & 296\dg 46.7\pr & S8\dg 40.1\pr \\
26 & 9 & 311\dg 46.8\pr & S8\dg 39.2\pr & 26 & 10 & 326\dg 46.9\pr & S8\dg 38.3\pr & 26 & 11 & 341\dg 47.0\pr & S8\dg 37.3\pr \\
26 & 12 & 356\dg 47.1\pr & S8\dg 36.4\pr & 26 & 13 & 11\dg 47.2\pr & S8\dg 35.4\pr & 26 & 14 & 26\dg 47.4\pr & S8\dg 34.5\pr \\
26 & 15 & 41\dg 47.5\pr & S8\dg 33.6\pr & 26 & 16 & 56\dg 47.6\pr & S8\dg 32.6\pr & 26 & 17 & 71\dg 47.7\pr & S8\dg 31.7\pr \\
26 & 18 & 86\dg 47.8\pr & S8\dg 30.8\pr & 26 & 19 & 101\dg 47.9\pr & S8\dg 29.8\pr & 26 & 20 & 116\dg 48.0\pr & S8\dg 28.9\pr \\
26 & 21 & 131\dg 48.1\pr & S8\dg 27.9\pr & 26 & 22 & 146\dg 48.2\pr & S8\dg 27.0\pr & 26 & 23 & 161\dg 48.3\pr & S8\dg 26.1\pr \\
27 & 0 & 176\dg 48.4\pr & S8\dg 25.1\pr & 27 & 1 & 191\dg 48.5\pr & S8\dg 24.2\pr & 27 & 2 & 206\dg 48.6\pr & S8\dg 23.2\pr \\
27 & 3 & 221\dg 48.7\pr & S8\dg 22.3\pr & 27 & 4 & 236\dg 48.9\pr & S8\dg 21.4\pr & 27 & 5 & 251\dg 49.0\pr & S8\dg 20.4\pr \\
27 & 6 & 266\dg 49.1\pr & S8\dg 19.5\pr & 27 & 7 & 281\dg 49.2\pr & S8\dg 18.5\pr & 27 & 8 & 296\dg 49.3\pr & S8\dg 17.6\pr \\
27 & 9 & 311\dg 49.4\pr & S8\dg 16.7\pr & 27 & 10 & 326\dg 49.5\pr & S8\dg 15.7\pr & 27 & 11 & 341\dg 49.6\pr & S8\dg 14.8\pr \\
27 & 12 & 356\dg 49.7\pr & S8\dg 13.8\pr & 27 & 13 & 11\dg 49.8\pr & S8\dg 12.9\pr & 27 & 14 & 26\dg 50.0\pr & S8\dg 12.0\pr \\
27 & 15 & 41\dg 50.1\pr & S8\dg 11.0\pr & 27 & 16 & 56\dg 50.2\pr & S8\dg 10.1\pr & 27 & 17 & 71\dg 50.3\pr & S8\dg 09.1\pr \\
27 & 18 & 86\dg 50.4\pr & S8\dg 08.2\pr & 27 & 19 & 101\dg 50.5\pr & S8\dg 07.2\pr & 27 & 20 & 116\dg 50.6\pr & S8\dg 06.3\pr \\
27 & 21 & 131\dg 50.7\pr & S8\dg 05.4\pr & 27 & 22 & 146\dg 50.9\pr & S8\dg 04.4\pr & 27 & 23 & 161\dg 51.0\pr & S8\dg 03.5\pr \\
28 & 0 & 176\dg 51.1\pr & S8\dg 02.5\pr & 28 & 1 & 191\dg 51.2\pr & S8\dg 01.6\pr & 28 & 2 & 206\dg 51.3\pr & S8\dg 00.6\pr \\
28 & 3 & 221\dg 51.4\pr & S7\dg 59.7\pr & 28 & 4 & 236\dg 51.5\pr & S7\dg 58.7\pr & 28 & 5 & 251\dg 51.7\pr & S7\dg 57.8\pr \\
28 & 6 & 266\dg 51.8\pr & S7\dg 56.9\pr & 28 & 7 & 281\dg 51.9\pr & S7\dg 55.9\pr & 28 & 8 & 296\dg 52.0\pr & S7\dg 55.0\pr \\
28 & 9 & 311\dg 52.1\pr & S7\dg 54.0\pr & 28 & 10 & 326\dg 52.2\pr & S7\dg 53.1\pr & 28 & 11 & 341\dg 52.4\pr & S7\dg 52.1\pr \\
28 & 12 & 356\dg 52.5\pr & S7\dg 51.2\pr & 28 & 13 & 11\dg 52.6\pr & S7\dg 50.2\pr & 28 & 14 & 26\dg 52.7\pr & S7\dg 49.3\pr \\
28 & 15 & 41\dg 52.8\pr & S7\dg 48.3\pr & 28 & 16 & 56\dg 52.9\pr & S7\dg 47.4\pr & 28 & 17 & 71\dg 53.1\pr & S7\dg 46.4\pr \\
28 & 18 & 86\dg 53.2\pr & S7\dg 45.5\pr & 28 & 19 & 101\dg 53.3\pr & S7\dg 44.5\pr & 28 & 20 & 116\dg 53.4\pr & S7\dg 43.6\pr \\
28 & 21 & 131\dg 53.5\pr & S7\dg 42.6\pr & 28 & 22 & 146\dg 53.6\pr & S7\dg 41.7\pr & 28 & 23 & 161\dg 53.8\pr & S7\dg 40.7\pr \\
\end{longtable}}

{\footnotesize
\begin{longtable}{rrrr|rrrr|rrrr}
\caption{Sun --- GHA and Declination, March 2026 (UT)}\label{sun03}\\
\toprule
\textbf{d} & \textbf{h} & \textbf{GHA} & \textbf{Dec} & \textbf{d} & \textbf{h} & \textbf{GHA} & \textbf{Dec} & \textbf{d} & \textbf{h} & \textbf{GHA} & \textbf{Dec} \\
\midrule
\endfirsthead
\multicolumn{12}{c}{\textit{Sun --- GHA and Declination, March 2026 (UT) (continued)}}\\
\toprule
\textbf{d} & \textbf{h} & \textbf{GHA} & \textbf{Dec} & \textbf{d} & \textbf{h} & \textbf{GHA} & \textbf{Dec} & \textbf{d} & \textbf{h} & \textbf{GHA} & \textbf{Dec} \\
\midrule
\endhead
\midrule\multicolumn{12}{r}{\textit{continued on next page}}\\
\endfoot
\bottomrule
\endlastfoot
1 & 0 & 176\dg 53.9\pr & S7\dg 39.8\pr & 1 & 1 & 191\dg 54.0\pr & S7\dg 38.9\pr & 1 & 2 & 206\dg 54.1\pr & S7\dg 37.9\pr \\
1 & 3 & 221\dg 54.2\pr & S7\dg 37.0\pr & 1 & 4 & 236\dg 54.4\pr & S7\dg 36.0\pr & 1 & 5 & 251\dg 54.5\pr & S7\dg 35.1\pr \\
1 & 6 & 266\dg 54.6\pr & S7\dg 34.1\pr & 1 & 7 & 281\dg 54.7\pr & S7\dg 33.2\pr & 1 & 8 & 296\dg 54.8\pr & S7\dg 32.2\pr \\
1 & 9 & 311\dg 55.0\pr & S7\dg 31.3\pr & 1 & 10 & 326\dg 55.1\pr & S7\dg 30.3\pr & 1 & 11 & 341\dg 55.2\pr & S7\dg 29.3\pr \\
1 & 12 & 356\dg 55.3\pr & S7\dg 28.4\pr & 1 & 13 & 11\dg 55.5\pr & S7\dg 27.4\pr & 1 & 14 & 26\dg 55.6\pr & S7\dg 26.5\pr \\
1 & 15 & 41\dg 55.7\pr & S7\dg 25.5\pr & 1 & 16 & 56\dg 55.8\pr & S7\dg 24.6\pr & 1 & 17 & 71\dg 55.9\pr & S7\dg 23.6\pr \\
1 & 18 & 86\dg 56.1\pr & S7\dg 22.7\pr & 1 & 19 & 101\dg 56.2\pr & S7\dg 21.7\pr & 1 & 20 & 116\dg 56.3\pr & S7\dg 20.8\pr \\
1 & 21 & 131\dg 56.4\pr & S7\dg 19.8\pr & 1 & 22 & 146\dg 56.6\pr & S7\dg 18.9\pr & 1 & 23 & 161\dg 56.7\pr & S7\dg 17.9\pr \\
2 & 0 & 176\dg 56.8\pr & S7\dg 17.0\pr & 2 & 1 & 191\dg 56.9\pr & S7\dg 16.0\pr & 2 & 2 & 206\dg 57.1\pr & S7\dg 15.1\pr \\
2 & 3 & 221\dg 57.2\pr & S7\dg 14.1\pr & 2 & 4 & 236\dg 57.3\pr & S7\dg 13.2\pr & 2 & 5 & 251\dg 57.4\pr & S7\dg 12.2\pr \\
2 & 6 & 266\dg 57.6\pr & S7\dg 11.2\pr & 2 & 7 & 281\dg 57.7\pr & S7\dg 10.3\pr & 2 & 8 & 296\dg 57.8\pr & S7\dg 09.3\pr \\
2 & 9 & 311\dg 57.9\pr & S7\dg 08.4\pr & 2 & 10 & 326\dg 58.1\pr & S7\dg 07.4\pr & 2 & 11 & 341\dg 58.2\pr & S7\dg 06.5\pr \\
2 & 12 & 356\dg 58.3\pr & S7\dg 05.5\pr & 2 & 13 & 11\dg 58.4\pr & S7\dg 04.6\pr & 2 & 14 & 26\dg 58.6\pr & S7\dg 03.6\pr \\
2 & 15 & 41\dg 58.7\pr & S7\dg 02.6\pr & 2 & 16 & 56\dg 58.8\pr & S7\dg 01.7\pr & 2 & 17 & 71\dg 59.0\pr & S7\dg 00.7\pr \\
2 & 18 & 86\dg 59.1\pr & S6\dg 59.8\pr & 2 & 19 & 101\dg 59.2\pr & S6\dg 58.8\pr & 2 & 20 & 116\dg 59.3\pr & S6\dg 57.9\pr \\
2 & 21 & 131\dg 59.5\pr & S6\dg 56.9\pr & 2 & 22 & 146\dg 59.6\pr & S6\dg 55.9\pr & 2 & 23 & 161\dg 59.7\pr & S6\dg 55.0\pr \\
3 & 0 & 176\dg 59.9\pr & S6\dg 54.0\pr & 3 & 1 & 192\dg 00.0\pr & S6\dg 53.1\pr & 3 & 2 & 207\dg 00.1\pr & S6\dg 52.1\pr \\
3 & 3 & 222\dg 00.2\pr & S6\dg 51.2\pr & 3 & 4 & 237\dg 00.4\pr & S6\dg 50.2\pr & 3 & 5 & 252\dg 00.5\pr & S6\dg 49.2\pr \\
3 & 6 & 267\dg 00.6\pr & S6\dg 48.3\pr & 3 & 7 & 282\dg 00.8\pr & S6\dg 47.3\pr & 3 & 8 & 297\dg 00.9\pr & S6\dg 46.4\pr \\
3 & 9 & 312\dg 01.0\pr & S6\dg 45.4\pr & 3 & 10 & 327\dg 01.2\pr & S6\dg 44.4\pr & 3 & 11 & 342\dg 01.3\pr & S6\dg 43.5\pr \\
3 & 12 & 357\dg 01.4\pr & S6\dg 42.5\pr & 3 & 13 & 12\dg 01.6\pr & S6\dg 41.6\pr & 3 & 14 & 27\dg 01.7\pr & S6\dg 40.6\pr \\
3 & 15 & 42\dg 01.8\pr & S6\dg 39.6\pr & 3 & 16 & 57\dg 02.0\pr & S6\dg 38.7\pr & 3 & 17 & 72\dg 02.1\pr & S6\dg 37.7\pr \\
3 & 18 & 87\dg 02.2\pr & S6\dg 36.8\pr & 3 & 19 & 102\dg 02.4\pr & S6\dg 35.8\pr & 3 & 20 & 117\dg 02.5\pr & S6\dg 34.8\pr \\
3 & 21 & 132\dg 02.6\pr & S6\dg 33.9\pr & 3 & 22 & 147\dg 02.8\pr & S6\dg 32.9\pr & 3 & 23 & 162\dg 02.9\pr & S6\dg 32.0\pr \\
4 & 0 & 177\dg 03.0\pr & S6\dg 31.0\pr & 4 & 1 & 192\dg 03.2\pr & S6\dg 30.0\pr & 4 & 2 & 207\dg 03.3\pr & S6\dg 29.1\pr \\
4 & 3 & 222\dg 03.4\pr & S6\dg 28.1\pr & 4 & 4 & 237\dg 03.6\pr & S6\dg 27.2\pr & 4 & 5 & 252\dg 03.7\pr & S6\dg 26.2\pr \\
4 & 6 & 267\dg 03.8\pr & S6\dg 25.2\pr & 4 & 7 & 282\dg 04.0\pr & S6\dg 24.3\pr & 4 & 8 & 297\dg 04.1\pr & S6\dg 23.3\pr \\
4 & 9 & 312\dg 04.2\pr & S6\dg 22.3\pr & 4 & 10 & 327\dg 04.4\pr & S6\dg 21.4\pr & 4 & 11 & 342\dg 04.5\pr & S6\dg 20.4\pr \\
4 & 12 & 357\dg 04.7\pr & S6\dg 19.5\pr & 4 & 13 & 12\dg 04.8\pr & S6\dg 18.5\pr & 4 & 14 & 27\dg 04.9\pr & S6\dg 17.5\pr \\
4 & 15 & 42\dg 05.1\pr & S6\dg 16.6\pr & 4 & 16 & 57\dg 05.2\pr & S6\dg 15.6\pr & 4 & 17 & 72\dg 05.3\pr & S6\dg 14.6\pr \\
4 & 18 & 87\dg 05.5\pr & S6\dg 13.7\pr & 4 & 19 & 102\dg 05.6\pr & S6\dg 12.7\pr & 4 & 20 & 117\dg 05.8\pr & S6\dg 11.7\pr \\
4 & 21 & 132\dg 05.9\pr & S6\dg 10.8\pr & 4 & 22 & 147\dg 06.0\pr & S6\dg 09.8\pr & 4 & 23 & 162\dg 06.2\pr & S6\dg 08.8\pr \\
5 & 0 & 177\dg 06.3\pr & S6\dg 07.9\pr & 5 & 1 & 192\dg 06.4\pr & S6\dg 06.9\pr & 5 & 2 & 207\dg 06.6\pr & S6\dg 05.9\pr \\
5 & 3 & 222\dg 06.7\pr & S6\dg 05.0\pr & 5 & 4 & 237\dg 06.9\pr & S6\dg 04.0\pr & 5 & 5 & 252\dg 07.0\pr & S6\dg 03.0\pr \\
5 & 6 & 267\dg 07.1\pr & S6\dg 02.1\pr & 5 & 7 & 282\dg 07.3\pr & S6\dg 01.1\pr & 5 & 8 & 297\dg 07.4\pr & S6\dg 00.2\pr \\
5 & 9 & 312\dg 07.6\pr & S5\dg 59.2\pr & 5 & 10 & 327\dg 07.7\pr & S5\dg 58.2\pr & 5 & 11 & 342\dg 07.8\pr & S5\dg 57.3\pr \\
5 & 12 & 357\dg 08.0\pr & S5\dg 56.3\pr & 5 & 13 & 12\dg 08.1\pr & S5\dg 55.3\pr & 5 & 14 & 27\dg 08.3\pr & S5\dg 54.3\pr \\
5 & 15 & 42\dg 08.4\pr & S5\dg 53.4\pr & 5 & 16 & 57\dg 08.6\pr & S5\dg 52.4\pr & 5 & 17 & 72\dg 08.7\pr & S5\dg 51.4\pr \\
5 & 18 & 87\dg 08.8\pr & S5\dg 50.5\pr & 5 & 19 & 102\dg 09.0\pr & S5\dg 49.5\pr & 5 & 20 & 117\dg 09.1\pr & S5\dg 48.5\pr \\
5 & 21 & 132\dg 09.3\pr & S5\dg 47.6\pr & 5 & 22 & 147\dg 09.4\pr & S5\dg 46.6\pr & 5 & 23 & 162\dg 09.6\pr & S5\dg 45.6\pr \\
6 & 0 & 177\dg 09.7\pr & S5\dg 44.7\pr & 6 & 1 & 192\dg 09.8\pr & S5\dg 43.7\pr & 6 & 2 & 207\dg 10.0\pr & S5\dg 42.7\pr \\
6 & 3 & 222\dg 10.1\pr & S5\dg 41.8\pr & 6 & 4 & 237\dg 10.3\pr & S5\dg 40.8\pr & 6 & 5 & 252\dg 10.4\pr & S5\dg 39.8\pr \\
6 & 6 & 267\dg 10.6\pr & S5\dg 38.9\pr & 6 & 7 & 282\dg 10.7\pr & S5\dg 37.9\pr & 6 & 8 & 297\dg 10.8\pr & S5\dg 36.9\pr \\
6 & 9 & 312\dg 11.0\pr & S5\dg 35.9\pr & 6 & 10 & 327\dg 11.1\pr & S5\dg 35.0\pr & 6 & 11 & 342\dg 11.3\pr & S5\dg 34.0\pr \\
6 & 12 & 357\dg 11.4\pr & S5\dg 33.0\pr & 6 & 13 & 12\dg 11.6\pr & S5\dg 32.1\pr & 6 & 14 & 27\dg 11.7\pr & S5\dg 31.1\pr \\
6 & 15 & 42\dg 11.9\pr & S5\dg 30.1\pr & 6 & 16 & 57\dg 12.0\pr & S5\dg 29.2\pr & 6 & 17 & 72\dg 12.2\pr & S5\dg 28.2\pr \\
6 & 18 & 87\dg 12.3\pr & S5\dg 27.2\pr & 6 & 19 & 102\dg 12.5\pr & S5\dg 26.2\pr & 6 & 20 & 117\dg 12.6\pr & S5\dg 25.3\pr \\
6 & 21 & 132\dg 12.7\pr & S5\dg 24.3\pr & 6 & 22 & 147\dg 12.9\pr & S5\dg 23.3\pr & 6 & 23 & 162\dg 13.0\pr & S5\dg 22.4\pr \\
7 & 0 & 177\dg 13.2\pr & S5\dg 21.4\pr & 7 & 1 & 192\dg 13.3\pr & S5\dg 20.4\pr & 7 & 2 & 207\dg 13.5\pr & S5\dg 19.4\pr \\
7 & 3 & 222\dg 13.6\pr & S5\dg 18.5\pr & 7 & 4 & 237\dg 13.8\pr & S5\dg 17.5\pr & 7 & 5 & 252\dg 13.9\pr & S5\dg 16.5\pr \\
7 & 6 & 267\dg 14.1\pr & S5\dg 15.6\pr & 7 & 7 & 282\dg 14.2\pr & S5\dg 14.6\pr & 7 & 8 & 297\dg 14.4\pr & S5\dg 13.6\pr \\
7 & 9 & 312\dg 14.5\pr & S5\dg 12.6\pr & 7 & 10 & 327\dg 14.7\pr & S5\dg 11.7\pr & 7 & 11 & 342\dg 14.8\pr & S5\dg 10.7\pr \\
7 & 12 & 357\dg 15.0\pr & S5\dg 09.7\pr & 7 & 13 & 12\dg 15.1\pr & S5\dg 08.7\pr & 7 & 14 & 27\dg 15.3\pr & S5\dg 07.8\pr \\
7 & 15 & 42\dg 15.4\pr & S5\dg 06.8\pr & 7 & 16 & 57\dg 15.6\pr & S5\dg 05.8\pr & 7 & 17 & 72\dg 15.7\pr & S5\dg 04.8\pr \\
7 & 18 & 87\dg 15.9\pr & S5\dg 03.9\pr & 7 & 19 & 102\dg 16.0\pr & S5\dg 02.9\pr & 7 & 20 & 117\dg 16.2\pr & S5\dg 01.9\pr \\
7 & 21 & 132\dg 16.3\pr & S5\dg 00.9\pr & 7 & 22 & 147\dg 16.5\pr & S5\dg 00.0\pr & 7 & 23 & 162\dg 16.6\pr & S4\dg 59.0\pr \\
8 & 0 & 177\dg 16.8\pr & S4\dg 58.0\pr & 8 & 1 & 192\dg 16.9\pr & S4\dg 57.1\pr & 8 & 2 & 207\dg 17.1\pr & S4\dg 56.1\pr \\
8 & 3 & 222\dg 17.2\pr & S4\dg 55.1\pr & 8 & 4 & 237\dg 17.4\pr & S4\dg 54.1\pr & 8 & 5 & 252\dg 17.5\pr & S4\dg 53.2\pr \\
8 & 6 & 267\dg 17.7\pr & S4\dg 52.2\pr & 8 & 7 & 282\dg 17.8\pr & S4\dg 51.2\pr & 8 & 8 & 297\dg 18.0\pr & S4\dg 50.2\pr \\
8 & 9 & 312\dg 18.2\pr & S4\dg 49.3\pr & 8 & 10 & 327\dg 18.3\pr & S4\dg 48.3\pr & 8 & 11 & 342\dg 18.5\pr & S4\dg 47.3\pr \\
8 & 12 & 357\dg 18.6\pr & S4\dg 46.3\pr & 8 & 13 & 12\dg 18.8\pr & S4\dg 45.3\pr & 8 & 14 & 27\dg 18.9\pr & S4\dg 44.4\pr \\
8 & 15 & 42\dg 19.1\pr & S4\dg 43.4\pr & 8 & 16 & 57\dg 19.2\pr & S4\dg 42.4\pr & 8 & 17 & 72\dg 19.4\pr & S4\dg 41.4\pr \\
8 & 18 & 87\dg 19.5\pr & S4\dg 40.5\pr & 8 & 19 & 102\dg 19.7\pr & S4\dg 39.5\pr & 8 & 20 & 117\dg 19.9\pr & S4\dg 38.5\pr \\
8 & 21 & 132\dg 20.0\pr & S4\dg 37.5\pr & 8 & 22 & 147\dg 20.2\pr & S4\dg 36.6\pr & 8 & 23 & 162\dg 20.3\pr & S4\dg 35.6\pr \\
9 & 0 & 177\dg 20.5\pr & S4\dg 34.6\pr & 9 & 1 & 192\dg 20.6\pr & S4\dg 33.6\pr & 9 & 2 & 207\dg 20.8\pr & S4\dg 32.6\pr \\
9 & 3 & 222\dg 20.9\pr & S4\dg 31.7\pr & 9 & 4 & 237\dg 21.1\pr & S4\dg 30.7\pr & 9 & 5 & 252\dg 21.3\pr & S4\dg 29.7\pr \\
9 & 6 & 267\dg 21.4\pr & S4\dg 28.7\pr & 9 & 7 & 282\dg 21.6\pr & S4\dg 27.8\pr & 9 & 8 & 297\dg 21.7\pr & S4\dg 26.8\pr \\
9 & 9 & 312\dg 21.9\pr & S4\dg 25.8\pr & 9 & 10 & 327\dg 22.0\pr & S4\dg 24.8\pr & 9 & 11 & 342\dg 22.2\pr & S4\dg 23.8\pr \\
9 & 12 & 357\dg 22.3\pr & S4\dg 22.9\pr & 9 & 13 & 12\dg 22.5\pr & S4\dg 21.9\pr & 9 & 14 & 27\dg 22.7\pr & S4\dg 20.9\pr \\
9 & 15 & 42\dg 22.8\pr & S4\dg 19.9\pr & 9 & 16 & 57\dg 23.0\pr & S4\dg 19.0\pr & 9 & 17 & 72\dg 23.1\pr & S4\dg 18.0\pr \\
9 & 18 & 87\dg 23.3\pr & S4\dg 17.0\pr & 9 & 19 & 102\dg 23.5\pr & S4\dg 16.0\pr & 9 & 20 & 117\dg 23.6\pr & S4\dg 15.0\pr \\
9 & 21 & 132\dg 23.8\pr & S4\dg 14.1\pr & 9 & 22 & 147\dg 23.9\pr & S4\dg 13.1\pr & 9 & 23 & 162\dg 24.1\pr & S4\dg 12.1\pr \\
10 & 0 & 177\dg 24.2\pr & S4\dg 11.1\pr & 10 & 1 & 192\dg 24.4\pr & S4\dg 10.1\pr & 10 & 2 & 207\dg 24.6\pr & S4\dg 09.2\pr \\
10 & 3 & 222\dg 24.7\pr & S4\dg 08.2\pr & 10 & 4 & 237\dg 24.9\pr & S4\dg 07.2\pr & 10 & 5 & 252\dg 25.0\pr & S4\dg 06.2\pr \\
10 & 6 & 267\dg 25.2\pr & S4\dg 05.2\pr & 10 & 7 & 282\dg 25.4\pr & S4\dg 04.3\pr & 10 & 8 & 297\dg 25.5\pr & S4\dg 03.3\pr \\
10 & 9 & 312\dg 25.7\pr & S4\dg 02.3\pr & 10 & 10 & 327\dg 25.8\pr & S4\dg 01.3\pr & 10 & 11 & 342\dg 26.0\pr & S4\dg 00.3\pr \\
10 & 12 & 357\dg 26.2\pr & S3\dg 59.4\pr & 10 & 13 & 12\dg 26.3\pr & S3\dg 58.4\pr & 10 & 14 & 27\dg 26.5\pr & S3\dg 57.4\pr \\
10 & 15 & 42\dg 26.7\pr & S3\dg 56.4\pr & 10 & 16 & 57\dg 26.8\pr & S3\dg 55.4\pr & 10 & 17 & 72\dg 27.0\pr & S3\dg 54.5\pr \\
10 & 18 & 87\dg 27.1\pr & S3\dg 53.5\pr & 10 & 19 & 102\dg 27.3\pr & S3\dg 52.5\pr & 10 & 20 & 117\dg 27.5\pr & S3\dg 51.5\pr \\
10 & 21 & 132\dg 27.6\pr & S3\dg 50.5\pr & 10 & 22 & 147\dg 27.8\pr & S3\dg 49.6\pr & 10 & 23 & 162\dg 27.9\pr & S3\dg 48.6\pr \\
11 & 0 & 177\dg 28.1\pr & S3\dg 47.6\pr & 11 & 1 & 192\dg 28.3\pr & S3\dg 46.6\pr & 11 & 2 & 207\dg 28.4\pr & S3\dg 45.6\pr \\
11 & 3 & 222\dg 28.6\pr & S3\dg 44.6\pr & 11 & 4 & 237\dg 28.8\pr & S3\dg 43.7\pr & 11 & 5 & 252\dg 28.9\pr & S3\dg 42.7\pr \\
11 & 6 & 267\dg 29.1\pr & S3\dg 41.7\pr & 11 & 7 & 282\dg 29.3\pr & S3\dg 40.7\pr & 11 & 8 & 297\dg 29.4\pr & S3\dg 39.7\pr \\
11 & 9 & 312\dg 29.6\pr & S3\dg 38.8\pr & 11 & 10 & 327\dg 29.7\pr & S3\dg 37.8\pr & 11 & 11 & 342\dg 29.9\pr & S3\dg 36.8\pr \\
11 & 12 & 357\dg 30.1\pr & S3\dg 35.8\pr & 11 & 13 & 12\dg 30.2\pr & S3\dg 34.8\pr & 11 & 14 & 27\dg 30.4\pr & S3\dg 33.8\pr \\
11 & 15 & 42\dg 30.6\pr & S3\dg 32.9\pr & 11 & 16 & 57\dg 30.7\pr & S3\dg 31.9\pr & 11 & 17 & 72\dg 30.9\pr & S3\dg 30.9\pr \\
11 & 18 & 87\dg 31.1\pr & S3\dg 29.9\pr & 11 & 19 & 102\dg 31.2\pr & S3\dg 28.9\pr & 11 & 20 & 117\dg 31.4\pr & S3\dg 27.9\pr \\
11 & 21 & 132\dg 31.6\pr & S3\dg 27.0\pr & 11 & 22 & 147\dg 31.7\pr & S3\dg 26.0\pr & 11 & 23 & 162\dg 31.9\pr & S3\dg 25.0\pr \\
12 & 0 & 177\dg 32.1\pr & S3\dg 24.0\pr & 12 & 1 & 192\dg 32.2\pr & S3\dg 23.0\pr & 12 & 2 & 207\dg 32.4\pr & S3\dg 22.0\pr \\
12 & 3 & 222\dg 32.6\pr & S3\dg 21.1\pr & 12 & 4 & 237\dg 32.7\pr & S3\dg 20.1\pr & 12 & 5 & 252\dg 32.9\pr & S3\dg 19.1\pr \\
12 & 6 & 267\dg 33.1\pr & S3\dg 18.1\pr & 12 & 7 & 282\dg 33.2\pr & S3\dg 17.1\pr & 12 & 8 & 297\dg 33.4\pr & S3\dg 16.1\pr \\
12 & 9 & 312\dg 33.6\pr & S3\dg 15.2\pr & 12 & 10 & 327\dg 33.7\pr & S3\dg 14.2\pr & 12 & 11 & 342\dg 33.9\pr & S3\dg 13.2\pr \\
12 & 12 & 357\dg 34.1\pr & S3\dg 12.2\pr & 12 & 13 & 12\dg 34.2\pr & S3\dg 11.2\pr & 12 & 14 & 27\dg 34.4\pr & S3\dg 10.2\pr \\
12 & 15 & 42\dg 34.6\pr & S3\dg 09.3\pr & 12 & 16 & 57\dg 34.7\pr & S3\dg 08.3\pr & 12 & 17 & 72\dg 34.9\pr & S3\dg 07.3\pr \\
12 & 18 & 87\dg 35.1\pr & S3\dg 06.3\pr & 12 & 19 & 102\dg 35.2\pr & S3\dg 05.3\pr & 12 & 20 & 117\dg 35.4\pr & S3\dg 04.3\pr \\
12 & 21 & 132\dg 35.6\pr & S3\dg 03.3\pr & 12 & 22 & 147\dg 35.7\pr & S3\dg 02.4\pr & 12 & 23 & 162\dg 35.9\pr & S3\dg 01.4\pr \\
13 & 0 & 177\dg 36.1\pr & S3\dg 00.4\pr & 13 & 1 & 192\dg 36.2\pr & S2\dg 59.4\pr & 13 & 2 & 207\dg 36.4\pr & S2\dg 58.4\pr \\
13 & 3 & 222\dg 36.6\pr & S2\dg 57.4\pr & 13 & 4 & 237\dg 36.7\pr & S2\dg 56.4\pr & 13 & 5 & 252\dg 36.9\pr & S2\dg 55.5\pr \\
13 & 6 & 267\dg 37.1\pr & S2\dg 54.5\pr & 13 & 7 & 282\dg 37.3\pr & S2\dg 53.5\pr & 13 & 8 & 297\dg 37.4\pr & S2\dg 52.5\pr \\
13 & 9 & 312\dg 37.6\pr & S2\dg 51.5\pr & 13 & 10 & 327\dg 37.8\pr & S2\dg 50.5\pr & 13 & 11 & 342\dg 37.9\pr & S2\dg 49.6\pr \\
13 & 12 & 357\dg 38.1\pr & S2\dg 48.6\pr & 13 & 13 & 12\dg 38.3\pr & S2\dg 47.6\pr & 13 & 14 & 27\dg 38.4\pr & S2\dg 46.6\pr \\
13 & 15 & 42\dg 38.6\pr & S2\dg 45.6\pr & 13 & 16 & 57\dg 38.8\pr & S2\dg 44.6\pr & 13 & 17 & 72\dg 39.0\pr & S2\dg 43.6\pr \\
13 & 18 & 87\dg 39.1\pr & S2\dg 42.7\pr & 13 & 19 & 102\dg 39.3\pr & S2\dg 41.7\pr & 13 & 20 & 117\dg 39.5\pr & S2\dg 40.7\pr \\
13 & 21 & 132\dg 39.6\pr & S2\dg 39.7\pr & 13 & 22 & 147\dg 39.8\pr & S2\dg 38.7\pr & 13 & 23 & 162\dg 40.0\pr & S2\dg 37.7\pr \\
14 & 0 & 177\dg 40.2\pr & S2\dg 36.7\pr & 14 & 1 & 192\dg 40.3\pr & S2\dg 35.8\pr & 14 & 2 & 207\dg 40.5\pr & S2\dg 34.8\pr \\
14 & 3 & 222\dg 40.7\pr & S2\dg 33.8\pr & 14 & 4 & 237\dg 40.8\pr & S2\dg 32.8\pr & 14 & 5 & 252\dg 41.0\pr & S2\dg 31.8\pr \\
14 & 6 & 267\dg 41.2\pr & S2\dg 30.8\pr & 14 & 7 & 282\dg 41.4\pr & S2\dg 29.8\pr & 14 & 8 & 297\dg 41.5\pr & S2\dg 28.8\pr \\
14 & 9 & 312\dg 41.7\pr & S2\dg 27.9\pr & 14 & 10 & 327\dg 41.9\pr & S2\dg 26.9\pr & 14 & 11 & 342\dg 42.1\pr & S2\dg 25.9\pr \\
14 & 12 & 357\dg 42.2\pr & S2\dg 24.9\pr & 14 & 13 & 12\dg 42.4\pr & S2\dg 23.9\pr & 14 & 14 & 27\dg 42.6\pr & S2\dg 22.9\pr \\
14 & 15 & 42\dg 42.7\pr & S2\dg 21.9\pr & 14 & 16 & 57\dg 42.9\pr & S2\dg 21.0\pr & 14 & 17 & 72\dg 43.1\pr & S2\dg 20.0\pr \\
14 & 18 & 87\dg 43.3\pr & S2\dg 19.0\pr & 14 & 19 & 102\dg 43.4\pr & S2\dg 18.0\pr & 14 & 20 & 117\dg 43.6\pr & S2\dg 17.0\pr \\
14 & 21 & 132\dg 43.8\pr & S2\dg 16.0\pr & 14 & 22 & 147\dg 44.0\pr & S2\dg 15.0\pr & 14 & 23 & 162\dg 44.1\pr & S2\dg 14.0\pr \\
15 & 0 & 177\dg 44.3\pr & S2\dg 13.1\pr & 15 & 1 & 192\dg 44.5\pr & S2\dg 12.1\pr & 15 & 2 & 207\dg 44.7\pr & S2\dg 11.1\pr \\
15 & 3 & 222\dg 44.8\pr & S2\dg 10.1\pr & 15 & 4 & 237\dg 45.0\pr & S2\dg 09.1\pr & 15 & 5 & 252\dg 45.2\pr & S2\dg 08.1\pr \\
15 & 6 & 267\dg 45.4\pr & S2\dg 07.1\pr & 15 & 7 & 282\dg 45.5\pr & S2\dg 06.1\pr & 15 & 8 & 297\dg 45.7\pr & S2\dg 05.2\pr \\
15 & 9 & 312\dg 45.9\pr & S2\dg 04.2\pr & 15 & 10 & 327\dg 46.1\pr & S2\dg 03.2\pr & 15 & 11 & 342\dg 46.2\pr & S2\dg 02.2\pr \\
15 & 12 & 357\dg 46.4\pr & S2\dg 01.2\pr & 15 & 13 & 12\dg 46.6\pr & S2\dg 00.2\pr & 15 & 14 & 27\dg 46.8\pr & S1\dg 59.2\pr \\
15 & 15 & 42\dg 46.9\pr & S1\dg 58.2\pr & 15 & 16 & 57\dg 47.1\pr & S1\dg 57.3\pr & 15 & 17 & 72\dg 47.3\pr & S1\dg 56.3\pr \\
15 & 18 & 87\dg 47.5\pr & S1\dg 55.3\pr & 15 & 19 & 102\dg 47.6\pr & S1\dg 54.3\pr & 15 & 20 & 117\dg 47.8\pr & S1\dg 53.3\pr \\
15 & 21 & 132\dg 48.0\pr & S1\dg 52.3\pr & 15 & 22 & 147\dg 48.2\pr & S1\dg 51.3\pr & 15 & 23 & 162\dg 48.3\pr & S1\dg 50.3\pr \\
16 & 0 & 177\dg 48.5\pr & S1\dg 49.4\pr & 16 & 1 & 192\dg 48.7\pr & S1\dg 48.4\pr & 16 & 2 & 207\dg 48.9\pr & S1\dg 47.4\pr \\
16 & 3 & 222\dg 49.0\pr & S1\dg 46.4\pr & 16 & 4 & 237\dg 49.2\pr & S1\dg 45.4\pr & 16 & 5 & 252\dg 49.4\pr & S1\dg 44.4\pr \\
16 & 6 & 267\dg 49.6\pr & S1\dg 43.4\pr & 16 & 7 & 282\dg 49.8\pr & S1\dg 42.4\pr & 16 & 8 & 297\dg 49.9\pr & S1\dg 41.5\pr \\
16 & 9 & 312\dg 50.1\pr & S1\dg 40.5\pr & 16 & 10 & 327\dg 50.3\pr & S1\dg 39.5\pr & 16 & 11 & 342\dg 50.5\pr & S1\dg 38.5\pr \\
16 & 12 & 357\dg 50.6\pr & S1\dg 37.5\pr & 16 & 13 & 12\dg 50.8\pr & S1\dg 36.5\pr & 16 & 14 & 27\dg 51.0\pr & S1\dg 35.5\pr \\
16 & 15 & 42\dg 51.2\pr & S1\dg 34.5\pr & 16 & 16 & 57\dg 51.4\pr & S1\dg 33.5\pr & 16 & 17 & 72\dg 51.5\pr & S1\dg 32.6\pr \\
16 & 18 & 87\dg 51.7\pr & S1\dg 31.6\pr & 16 & 19 & 102\dg 51.9\pr & S1\dg 30.6\pr & 16 & 20 & 117\dg 52.1\pr & S1\dg 29.6\pr \\
16 & 21 & 132\dg 52.2\pr & S1\dg 28.6\pr & 16 & 22 & 147\dg 52.4\pr & S1\dg 27.6\pr & 16 & 23 & 162\dg 52.6\pr & S1\dg 26.6\pr \\
17 & 0 & 177\dg 52.8\pr & S1\dg 25.6\pr & 17 & 1 & 192\dg 53.0\pr & S1\dg 24.7\pr & 17 & 2 & 207\dg 53.1\pr & S1\dg 23.7\pr \\
17 & 3 & 222\dg 53.3\pr & S1\dg 22.7\pr & 17 & 4 & 237\dg 53.5\pr & S1\dg 21.7\pr & 17 & 5 & 252\dg 53.7\pr & S1\dg 20.7\pr \\
17 & 6 & 267\dg 53.9\pr & S1\dg 19.7\pr & 17 & 7 & 282\dg 54.0\pr & S1\dg 18.7\pr & 17 & 8 & 297\dg 54.2\pr & S1\dg 17.7\pr \\
17 & 9 & 312\dg 54.4\pr & S1\dg 16.7\pr & 17 & 10 & 327\dg 54.6\pr & S1\dg 15.8\pr & 17 & 11 & 342\dg 54.8\pr & S1\dg 14.8\pr \\
17 & 12 & 357\dg 54.9\pr & S1\dg 13.8\pr & 17 & 13 & 12\dg 55.1\pr & S1\dg 12.8\pr & 17 & 14 & 27\dg 55.3\pr & S1\dg 11.8\pr \\
17 & 15 & 42\dg 55.5\pr & S1\dg 10.8\pr & 17 & 16 & 57\dg 55.7\pr & S1\dg 09.8\pr & 17 & 17 & 72\dg 55.8\pr & S1\dg 08.8\pr \\
17 & 18 & 87\dg 56.0\pr & S1\dg 07.8\pr & 17 & 19 & 102\dg 56.2\pr & S1\dg 06.9\pr & 17 & 20 & 117\dg 56.4\pr & S1\dg 05.9\pr \\
17 & 21 & 132\dg 56.6\pr & S1\dg 04.9\pr & 17 & 22 & 147\dg 56.7\pr & S1\dg 03.9\pr & 17 & 23 & 162\dg 56.9\pr & S1\dg 02.9\pr \\
18 & 0 & 177\dg 57.1\pr & S1\dg 01.9\pr & 18 & 1 & 192\dg 57.3\pr & S1\dg 00.9\pr & 18 & 2 & 207\dg 57.5\pr & S0\dg 59.9\pr \\
18 & 3 & 222\dg 57.6\pr & S0\dg 58.9\pr & 18 & 4 & 237\dg 57.8\pr & S0\dg 58.0\pr & 18 & 5 & 252\dg 58.0\pr & S0\dg 57.0\pr \\
18 & 6 & 267\dg 58.2\pr & S0\dg 56.0\pr & 18 & 7 & 282\dg 58.4\pr & S0\dg 55.0\pr & 18 & 8 & 297\dg 58.5\pr & S0\dg 54.0\pr \\
18 & 9 & 312\dg 58.7\pr & S0\dg 53.0\pr & 18 & 10 & 327\dg 58.9\pr & S0\dg 52.0\pr & 18 & 11 & 342\dg 59.1\pr & S0\dg 51.0\pr \\
18 & 12 & 357\dg 59.3\pr & S0\dg 50.1\pr & 18 & 13 & 12\dg 59.5\pr & S0\dg 49.1\pr & 18 & 14 & 27\dg 59.6\pr & S0\dg 48.1\pr \\
18 & 15 & 42\dg 59.8\pr & S0\dg 47.1\pr & 18 & 16 & 58\dg 00.0\pr & S0\dg 46.1\pr & 18 & 17 & 73\dg 00.2\pr & S0\dg 45.1\pr \\
18 & 18 & 88\dg 00.4\pr & S0\dg 44.1\pr & 18 & 19 & 103\dg 00.5\pr & S0\dg 43.1\pr & 18 & 20 & 118\dg 00.7\pr & S0\dg 42.1\pr \\
18 & 21 & 133\dg 00.9\pr & S0\dg 41.2\pr & 18 & 22 & 148\dg 01.1\pr & S0\dg 40.2\pr & 18 & 23 & 163\dg 01.3\pr & S0\dg 39.2\pr \\
19 & 0 & 178\dg 01.5\pr & S0\dg 38.2\pr & 19 & 1 & 193\dg 01.6\pr & S0\dg 37.2\pr & 19 & 2 & 208\dg 01.8\pr & S0\dg 36.2\pr \\
19 & 3 & 223\dg 02.0\pr & S0\dg 35.2\pr & 19 & 4 & 238\dg 02.2\pr & S0\dg 34.2\pr & 19 & 5 & 253\dg 02.4\pr & S0\dg 33.2\pr \\
19 & 6 & 268\dg 02.5\pr & S0\dg 32.3\pr & 19 & 7 & 283\dg 02.7\pr & S0\dg 31.3\pr & 19 & 8 & 298\dg 02.9\pr & S0\dg 30.3\pr \\
19 & 9 & 313\dg 03.1\pr & S0\dg 29.3\pr & 19 & 10 & 328\dg 03.3\pr & S0\dg 28.3\pr & 19 & 11 & 343\dg 03.5\pr & S0\dg 27.3\pr \\
19 & 12 & 358\dg 03.6\pr & S0\dg 26.3\pr & 19 & 13 & 13\dg 03.8\pr & S0\dg 25.3\pr & 19 & 14 & 28\dg 04.0\pr & S0\dg 24.3\pr \\
19 & 15 & 43\dg 04.2\pr & S0\dg 23.4\pr & 19 & 16 & 58\dg 04.4\pr & S0\dg 22.4\pr & 19 & 17 & 73\dg 04.6\pr & S0\dg 21.4\pr \\
19 & 18 & 88\dg 04.7\pr & S0\dg 20.4\pr & 19 & 19 & 103\dg 04.9\pr & S0\dg 19.4\pr & 19 & 20 & 118\dg 05.1\pr & S0\dg 18.4\pr \\
19 & 21 & 133\dg 05.3\pr & S0\dg 17.4\pr & 19 & 22 & 148\dg 05.5\pr & S0\dg 16.4\pr & 19 & 23 & 163\dg 05.7\pr & S0\dg 15.5\pr \\
20 & 0 & 178\dg 05.9\pr & S0\dg 14.5\pr & 20 & 1 & 193\dg 06.0\pr & S0\dg 13.5\pr & 20 & 2 & 208\dg 06.2\pr & S0\dg 12.5\pr \\
20 & 3 & 223\dg 06.4\pr & S0\dg 11.5\pr & 20 & 4 & 238\dg 06.6\pr & S0\dg 10.5\pr & 20 & 5 & 253\dg 06.8\pr & S0\dg 09.5\pr \\
20 & 6 & 268\dg 07.0\pr & S0\dg 08.5\pr & 20 & 7 & 283\dg 07.1\pr & S0\dg 07.5\pr & 20 & 8 & 298\dg 07.3\pr & S0\dg 06.6\pr \\
20 & 9 & 313\dg 07.5\pr & S0\dg 05.6\pr & 20 & 10 & 328\dg 07.7\pr & S0\dg 04.6\pr & 20 & 11 & 343\dg 07.9\pr & S0\dg 03.6\pr \\
20 & 12 & 358\dg 08.1\pr & S0\dg 02.6\pr & 20 & 13 & 13\dg 08.2\pr & S0\dg 01.6\pr & 20 & 14 & 28\dg 08.4\pr & S0\dg 00.6\pr \\
20 & 15 & 43\dg 08.6\pr & N0\dg 00.4\pr & 20 & 16 & 58\dg 08.8\pr & N0\dg 01.3\pr & 20 & 17 & 73\dg 09.0\pr & N0\dg 02.3\pr \\
20 & 18 & 88\dg 09.2\pr & N0\dg 03.3\pr & 20 & 19 & 103\dg 09.4\pr & N0\dg 04.3\pr & 20 & 20 & 118\dg 09.5\pr & N0\dg 05.3\pr \\
20 & 21 & 133\dg 09.7\pr & N0\dg 06.3\pr & 20 & 22 & 148\dg 09.9\pr & N0\dg 07.3\pr & 20 & 23 & 163\dg 10.1\pr & N0\dg 08.3\pr \\
21 & 0 & 178\dg 10.3\pr & N0\dg 09.2\pr & 21 & 1 & 193\dg 10.5\pr & N0\dg 10.2\pr & 21 & 2 & 208\dg 10.7\pr & N0\dg 11.2\pr \\
21 & 3 & 223\dg 10.8\pr & N0\dg 12.2\pr & 21 & 4 & 238\dg 11.0\pr & N0\dg 13.2\pr & 21 & 5 & 253\dg 11.2\pr & N0\dg 14.2\pr \\
21 & 6 & 268\dg 11.4\pr & N0\dg 15.2\pr & 21 & 7 & 283\dg 11.6\pr & N0\dg 16.2\pr & 21 & 8 & 298\dg 11.8\pr & N0\dg 17.1\pr \\
21 & 9 & 313\dg 12.0\pr & N0\dg 18.1\pr & 21 & 10 & 328\dg 12.1\pr & N0\dg 19.1\pr & 21 & 11 & 343\dg 12.3\pr & N0\dg 20.1\pr \\
21 & 12 & 358\dg 12.5\pr & N0\dg 21.1\pr & 21 & 13 & 13\dg 12.7\pr & N0\dg 22.1\pr & 21 & 14 & 28\dg 12.9\pr & N0\dg 23.1\pr \\
21 & 15 & 43\dg 13.1\pr & N0\dg 24.1\pr & 21 & 16 & 58\dg 13.3\pr & N0\dg 25.0\pr & 21 & 17 & 73\dg 13.4\pr & N0\dg 26.0\pr \\
21 & 18 & 88\dg 13.6\pr & N0\dg 27.0\pr & 21 & 19 & 103\dg 13.8\pr & N0\dg 28.0\pr & 21 & 20 & 118\dg 14.0\pr & N0\dg 29.0\pr \\
21 & 21 & 133\dg 14.2\pr & N0\dg 30.0\pr & 21 & 22 & 148\dg 14.4\pr & N0\dg 31.0\pr & 21 & 23 & 163\dg 14.6\pr & N0\dg 32.0\pr \\
22 & 0 & 178\dg 14.7\pr & N0\dg 32.9\pr & 22 & 1 & 193\dg 14.9\pr & N0\dg 33.9\pr & 22 & 2 & 208\dg 15.1\pr & N0\dg 34.9\pr \\
22 & 3 & 223\dg 15.3\pr & N0\dg 35.9\pr & 22 & 4 & 238\dg 15.5\pr & N0\dg 36.9\pr & 22 & 5 & 253\dg 15.7\pr & N0\dg 37.9\pr \\
22 & 6 & 268\dg 15.9\pr & N0\dg 38.9\pr & 22 & 7 & 283\dg 16.0\pr & N0\dg 39.9\pr & 22 & 8 & 298\dg 16.2\pr & N0\dg 40.8\pr \\
22 & 9 & 313\dg 16.4\pr & N0\dg 41.8\pr & 22 & 10 & 328\dg 16.6\pr & N0\dg 42.8\pr & 22 & 11 & 343\dg 16.8\pr & N0\dg 43.8\pr \\
22 & 12 & 358\dg 17.0\pr & N0\dg 44.8\pr & 22 & 13 & 13\dg 17.2\pr & N0\dg 45.8\pr & 22 & 14 & 28\dg 17.4\pr & N0\dg 46.8\pr \\
22 & 15 & 43\dg 17.5\pr & N0\dg 47.7\pr & 22 & 16 & 58\dg 17.7\pr & N0\dg 48.7\pr & 22 & 17 & 73\dg 17.9\pr & N0\dg 49.7\pr \\
22 & 18 & 88\dg 18.1\pr & N0\dg 50.7\pr & 22 & 19 & 103\dg 18.3\pr & N0\dg 51.7\pr & 22 & 20 & 118\dg 18.5\pr & N0\dg 52.7\pr \\
22 & 21 & 133\dg 18.7\pr & N0\dg 53.7\pr & 22 & 22 & 148\dg 18.8\pr & N0\dg 54.7\pr & 22 & 23 & 163\dg 19.0\pr & N0\dg 55.6\pr \\
23 & 0 & 178\dg 19.2\pr & N0\dg 56.6\pr & 23 & 1 & 193\dg 19.4\pr & N0\dg 57.6\pr & 23 & 2 & 208\dg 19.6\pr & N0\dg 58.6\pr \\
23 & 3 & 223\dg 19.8\pr & N0\dg 59.6\pr & 23 & 4 & 238\dg 20.0\pr & N1\dg 00.6\pr & 23 & 5 & 253\dg 20.2\pr & N1\dg 01.6\pr \\
23 & 6 & 268\dg 20.3\pr & N1\dg 02.5\pr & 23 & 7 & 283\dg 20.5\pr & N1\dg 03.5\pr & 23 & 8 & 298\dg 20.7\pr & N1\dg 04.5\pr \\
23 & 9 & 313\dg 20.9\pr & N1\dg 05.5\pr & 23 & 10 & 328\dg 21.1\pr & N1\dg 06.5\pr & 23 & 11 & 343\dg 21.3\pr & N1\dg 07.5\pr \\
23 & 12 & 358\dg 21.5\pr & N1\dg 08.5\pr & 23 & 13 & 13\dg 21.7\pr & N1\dg 09.4\pr & 23 & 14 & 28\dg 21.8\pr & N1\dg 10.4\pr \\
23 & 15 & 43\dg 22.0\pr & N1\dg 11.4\pr & 23 & 16 & 58\dg 22.2\pr & N1\dg 12.4\pr & 23 & 17 & 73\dg 22.4\pr & N1\dg 13.4\pr \\
23 & 18 & 88\dg 22.6\pr & N1\dg 14.4\pr & 23 & 19 & 103\dg 22.8\pr & N1\dg 15.3\pr & 23 & 20 & 118\dg 23.0\pr & N1\dg 16.3\pr \\
23 & 21 & 133\dg 23.2\pr & N1\dg 17.3\pr & 23 & 22 & 148\dg 23.3\pr & N1\dg 18.3\pr & 23 & 23 & 163\dg 23.5\pr & N1\dg 19.3\pr \\
24 & 0 & 178\dg 23.7\pr & N1\dg 20.3\pr & 24 & 1 & 193\dg 23.9\pr & N1\dg 21.3\pr & 24 & 2 & 208\dg 24.1\pr & N1\dg 22.2\pr \\
24 & 3 & 223\dg 24.3\pr & N1\dg 23.2\pr & 24 & 4 & 238\dg 24.5\pr & N1\dg 24.2\pr & 24 & 5 & 253\dg 24.7\pr & N1\dg 25.2\pr \\
24 & 6 & 268\dg 24.9\pr & N1\dg 26.2\pr & 24 & 7 & 283\dg 25.0\pr & N1\dg 27.2\pr & 24 & 8 & 298\dg 25.2\pr & N1\dg 28.1\pr \\
24 & 9 & 313\dg 25.4\pr & N1\dg 29.1\pr & 24 & 10 & 328\dg 25.6\pr & N1\dg 30.1\pr & 24 & 11 & 343\dg 25.8\pr & N1\dg 31.1\pr \\
24 & 12 & 358\dg 26.0\pr & N1\dg 32.1\pr & 24 & 13 & 13\dg 26.2\pr & N1\dg 33.1\pr & 24 & 14 & 28\dg 26.4\pr & N1\dg 34.1\pr \\
24 & 15 & 43\dg 26.5\pr & N1\dg 35.0\pr & 24 & 16 & 58\dg 26.7\pr & N1\dg 36.0\pr & 24 & 17 & 73\dg 26.9\pr & N1\dg 37.0\pr \\
24 & 18 & 88\dg 27.1\pr & N1\dg 38.0\pr & 24 & 19 & 103\dg 27.3\pr & N1\dg 39.0\pr & 24 & 20 & 118\dg 27.5\pr & N1\dg 40.0\pr \\
24 & 21 & 133\dg 27.7\pr & N1\dg 40.9\pr & 24 & 22 & 148\dg 27.9\pr & N1\dg 41.9\pr & 24 & 23 & 163\dg 28.1\pr & N1\dg 42.9\pr \\
25 & 0 & 178\dg 28.2\pr & N1\dg 43.9\pr & 25 & 1 & 193\dg 28.4\pr & N1\dg 44.9\pr & 25 & 2 & 208\dg 28.6\pr & N1\dg 45.9\pr \\
25 & 3 & 223\dg 28.8\pr & N1\dg 46.8\pr & 25 & 4 & 238\dg 29.0\pr & N1\dg 47.8\pr & 25 & 5 & 253\dg 29.2\pr & N1\dg 48.8\pr \\
25 & 6 & 268\dg 29.4\pr & N1\dg 49.8\pr & 25 & 7 & 283\dg 29.6\pr & N1\dg 50.8\pr & 25 & 8 & 298\dg 29.7\pr & N1\dg 51.8\pr \\
25 & 9 & 313\dg 29.9\pr & N1\dg 52.7\pr & 25 & 10 & 328\dg 30.1\pr & N1\dg 53.7\pr & 25 & 11 & 343\dg 30.3\pr & N1\dg 54.7\pr \\
25 & 12 & 358\dg 30.5\pr & N1\dg 55.7\pr & 25 & 13 & 13\dg 30.7\pr & N1\dg 56.7\pr & 25 & 14 & 28\dg 30.9\pr & N1\dg 57.6\pr \\
25 & 15 & 43\dg 31.1\pr & N1\dg 58.6\pr & 25 & 16 & 58\dg 31.3\pr & N1\dg 59.6\pr & 25 & 17 & 73\dg 31.4\pr & N2\dg 00.6\pr \\
25 & 18 & 88\dg 31.6\pr & N2\dg 01.6\pr & 25 & 19 & 103\dg 31.8\pr & N2\dg 02.6\pr & 25 & 20 & 118\dg 32.0\pr & N2\dg 03.5\pr \\
25 & 21 & 133\dg 32.2\pr & N2\dg 04.5\pr & 25 & 22 & 148\dg 32.4\pr & N2\dg 05.5\pr & 25 & 23 & 163\dg 32.6\pr & N2\dg 06.5\pr \\
26 & 0 & 178\dg 32.8\pr & N2\dg 07.5\pr & 26 & 1 & 193\dg 33.0\pr & N2\dg 08.4\pr & 26 & 2 & 208\dg 33.1\pr & N2\dg 09.4\pr \\
26 & 3 & 223\dg 33.3\pr & N2\dg 10.4\pr & 26 & 4 & 238\dg 33.5\pr & N2\dg 11.4\pr & 26 & 5 & 253\dg 33.7\pr & N2\dg 12.4\pr \\
26 & 6 & 268\dg 33.9\pr & N2\dg 13.3\pr & 26 & 7 & 283\dg 34.1\pr & N2\dg 14.3\pr & 26 & 8 & 298\dg 34.3\pr & N2\dg 15.3\pr \\
26 & 9 & 313\dg 34.5\pr & N2\dg 16.3\pr & 26 & 10 & 328\dg 34.7\pr & N2\dg 17.3\pr & 26 & 11 & 343\dg 34.8\pr & N2\dg 18.3\pr \\
26 & 12 & 358\dg 35.0\pr & N2\dg 19.2\pr & 26 & 13 & 13\dg 35.2\pr & N2\dg 20.2\pr & 26 & 14 & 28\dg 35.4\pr & N2\dg 21.2\pr \\
26 & 15 & 43\dg 35.6\pr & N2\dg 22.2\pr & 26 & 16 & 58\dg 35.8\pr & N2\dg 23.2\pr & 26 & 17 & 73\dg 36.0\pr & N2\dg 24.1\pr \\
26 & 18 & 88\dg 36.2\pr & N2\dg 25.1\pr & 26 & 19 & 103\dg 36.3\pr & N2\dg 26.1\pr & 26 & 20 & 118\dg 36.5\pr & N2\dg 27.1\pr \\
26 & 21 & 133\dg 36.7\pr & N2\dg 28.1\pr & 26 & 22 & 148\dg 36.9\pr & N2\dg 29.0\pr & 26 & 23 & 163\dg 37.1\pr & N2\dg 30.0\pr \\
27 & 0 & 178\dg 37.3\pr & N2\dg 31.0\pr & 27 & 1 & 193\dg 37.5\pr & N2\dg 32.0\pr & 27 & 2 & 208\dg 37.7\pr & N2\dg 33.0\pr \\
27 & 3 & 223\dg 37.9\pr & N2\dg 33.9\pr & 27 & 4 & 238\dg 38.0\pr & N2\dg 34.9\pr & 27 & 5 & 253\dg 38.2\pr & N2\dg 35.9\pr \\
27 & 6 & 268\dg 38.4\pr & N2\dg 36.9\pr & 27 & 7 & 283\dg 38.6\pr & N2\dg 37.8\pr & 27 & 8 & 298\dg 38.8\pr & N2\dg 38.8\pr \\
27 & 9 & 313\dg 39.0\pr & N2\dg 39.8\pr & 27 & 10 & 328\dg 39.2\pr & N2\dg 40.8\pr & 27 & 11 & 343\dg 39.4\pr & N2\dg 41.8\pr \\
27 & 12 & 358\dg 39.6\pr & N2\dg 42.7\pr & 27 & 13 & 13\dg 39.7\pr & N2\dg 43.7\pr & 27 & 14 & 28\dg 39.9\pr & N2\dg 44.7\pr \\
27 & 15 & 43\dg 40.1\pr & N2\dg 45.7\pr & 27 & 16 & 58\dg 40.3\pr & N2\dg 46.7\pr & 27 & 17 & 73\dg 40.5\pr & N2\dg 47.6\pr \\
27 & 18 & 88\dg 40.7\pr & N2\dg 48.6\pr & 27 & 19 & 103\dg 40.9\pr & N2\dg 49.6\pr & 27 & 20 & 118\dg 41.1\pr & N2\dg 50.6\pr \\
27 & 21 & 133\dg 41.3\pr & N2\dg 51.5\pr & 27 & 22 & 148\dg 41.4\pr & N2\dg 52.5\pr & 27 & 23 & 163\dg 41.6\pr & N2\dg 53.5\pr \\
28 & 0 & 178\dg 41.8\pr & N2\dg 54.5\pr & 28 & 1 & 193\dg 42.0\pr & N2\dg 55.4\pr & 28 & 2 & 208\dg 42.2\pr & N2\dg 56.4\pr \\
28 & 3 & 223\dg 42.4\pr & N2\dg 57.4\pr & 28 & 4 & 238\dg 42.6\pr & N2\dg 58.4\pr & 28 & 5 & 253\dg 42.8\pr & N2\dg 59.4\pr \\
28 & 6 & 268\dg 42.9\pr & N3\dg 00.3\pr & 28 & 7 & 283\dg 43.1\pr & N3\dg 01.3\pr & 28 & 8 & 298\dg 43.3\pr & N3\dg 02.3\pr \\
28 & 9 & 313\dg 43.5\pr & N3\dg 03.3\pr & 28 & 10 & 328\dg 43.7\pr & N3\dg 04.2\pr & 28 & 11 & 343\dg 43.9\pr & N3\dg 05.2\pr \\
28 & 12 & 358\dg 44.1\pr & N3\dg 06.2\pr & 28 & 13 & 13\dg 44.3\pr & N3\dg 07.2\pr & 28 & 14 & 28\dg 44.5\pr & N3\dg 08.1\pr \\
28 & 15 & 43\dg 44.6\pr & N3\dg 09.1\pr & 28 & 16 & 58\dg 44.8\pr & N3\dg 10.1\pr & 28 & 17 & 73\dg 45.0\pr & N3\dg 11.1\pr \\
28 & 18 & 88\dg 45.2\pr & N3\dg 12.0\pr & 28 & 19 & 103\dg 45.4\pr & N3\dg 13.0\pr & 28 & 20 & 118\dg 45.6\pr & N3\dg 14.0\pr \\
28 & 21 & 133\dg 45.8\pr & N3\dg 15.0\pr & 28 & 22 & 148\dg 46.0\pr & N3\dg 15.9\pr & 28 & 23 & 163\dg 46.1\pr & N3\dg 16.9\pr \\
29 & 0 & 178\dg 46.3\pr & N3\dg 17.9\pr & 29 & 1 & 193\dg 46.5\pr & N3\dg 18.9\pr & 29 & 2 & 208\dg 46.7\pr & N3\dg 19.8\pr \\
29 & 3 & 223\dg 46.9\pr & N3\dg 20.8\pr & 29 & 4 & 238\dg 47.1\pr & N3\dg 21.8\pr & 29 & 5 & 253\dg 47.3\pr & N3\dg 22.8\pr \\
29 & 6 & 268\dg 47.5\pr & N3\dg 23.7\pr & 29 & 7 & 283\dg 47.7\pr & N3\dg 24.7\pr & 29 & 8 & 298\dg 47.8\pr & N3\dg 25.7\pr \\
29 & 9 & 313\dg 48.0\pr & N3\dg 26.7\pr & 29 & 10 & 328\dg 48.2\pr & N3\dg 27.6\pr & 29 & 11 & 343\dg 48.4\pr & N3\dg 28.6\pr \\
29 & 12 & 358\dg 48.6\pr & N3\dg 29.6\pr & 29 & 13 & 13\dg 48.8\pr & N3\dg 30.6\pr & 29 & 14 & 28\dg 49.0\pr & N3\dg 31.5\pr \\
29 & 15 & 43\dg 49.2\pr & N3\dg 32.5\pr & 29 & 16 & 58\dg 49.3\pr & N3\dg 33.5\pr & 29 & 17 & 73\dg 49.5\pr & N3\dg 34.4\pr \\
29 & 18 & 88\dg 49.7\pr & N3\dg 35.4\pr & 29 & 19 & 103\dg 49.9\pr & N3\dg 36.4\pr & 29 & 20 & 118\dg 50.1\pr & N3\dg 37.4\pr \\
29 & 21 & 133\dg 50.3\pr & N3\dg 38.3\pr & 29 & 22 & 148\dg 50.5\pr & N3\dg 39.3\pr & 29 & 23 & 163\dg 50.7\pr & N3\dg 40.3\pr \\
30 & 0 & 178\dg 50.8\pr & N3\dg 41.2\pr & 30 & 1 & 193\dg 51.0\pr & N3\dg 42.2\pr & 30 & 2 & 208\dg 51.2\pr & N3\dg 43.2\pr \\
30 & 3 & 223\dg 51.4\pr & N3\dg 44.2\pr & 30 & 4 & 238\dg 51.6\pr & N3\dg 45.1\pr & 30 & 5 & 253\dg 51.8\pr & N3\dg 46.1\pr \\
30 & 6 & 268\dg 52.0\pr & N3\dg 47.1\pr & 30 & 7 & 283\dg 52.2\pr & N3\dg 48.0\pr & 30 & 8 & 298\dg 52.3\pr & N3\dg 49.0\pr \\
30 & 9 & 313\dg 52.5\pr & N3\dg 50.0\pr & 30 & 10 & 328\dg 52.7\pr & N3\dg 51.0\pr & 30 & 11 & 343\dg 52.9\pr & N3\dg 51.9\pr \\
30 & 12 & 358\dg 53.1\pr & N3\dg 52.9\pr & 30 & 13 & 13\dg 53.3\pr & N3\dg 53.9\pr & 30 & 14 & 28\dg 53.5\pr & N3\dg 54.8\pr \\
30 & 15 & 43\dg 53.7\pr & N3\dg 55.8\pr & 30 & 16 & 58\dg 53.8\pr & N3\dg 56.8\pr & 30 & 17 & 73\dg 54.0\pr & N3\dg 57.8\pr \\
30 & 18 & 88\dg 54.2\pr & N3\dg 58.7\pr & 30 & 19 & 103\dg 54.4\pr & N3\dg 59.7\pr & 30 & 20 & 118\dg 54.6\pr & N4\dg 00.7\pr \\
30 & 21 & 133\dg 54.8\pr & N4\dg 01.6\pr & 30 & 22 & 148\dg 55.0\pr & N4\dg 02.6\pr & 30 & 23 & 163\dg 55.1\pr & N4\dg 03.6\pr \\
31 & 0 & 178\dg 55.3\pr & N4\dg 04.5\pr & 31 & 1 & 193\dg 55.5\pr & N4\dg 05.5\pr & 31 & 2 & 208\dg 55.7\pr & N4\dg 06.5\pr \\
31 & 3 & 223\dg 55.9\pr & N4\dg 07.4\pr & 31 & 4 & 238\dg 56.1\pr & N4\dg 08.4\pr & 31 & 5 & 253\dg 56.3\pr & N4\dg 09.4\pr \\
31 & 6 & 268\dg 56.5\pr & N4\dg 10.3\pr & 31 & 7 & 283\dg 56.6\pr & N4\dg 11.3\pr & 31 & 8 & 298\dg 56.8\pr & N4\dg 12.3\pr \\
31 & 9 & 313\dg 57.0\pr & N4\dg 13.3\pr & 31 & 10 & 328\dg 57.2\pr & N4\dg 14.2\pr & 31 & 11 & 343\dg 57.4\pr & N4\dg 15.2\pr \\
31 & 12 & 358\dg 57.6\pr & N4\dg 16.2\pr & 31 & 13 & 13\dg 57.8\pr & N4\dg 17.1\pr & 31 & 14 & 28\dg 57.9\pr & N4\dg 18.1\pr \\
31 & 15 & 43\dg 58.1\pr & N4\dg 19.1\pr & 31 & 16 & 58\dg 58.3\pr & N4\dg 20.0\pr & 31 & 17 & 73\dg 58.5\pr & N4\dg 21.0\pr \\
31 & 18 & 88\dg 58.7\pr & N4\dg 22.0\pr & 31 & 19 & 103\dg 58.9\pr & N4\dg 22.9\pr & 31 & 20 & 118\dg 59.1\pr & N4\dg 23.9\pr \\
31 & 21 & 133\dg 59.2\pr & N4\dg 24.9\pr & 31 & 22 & 148\dg 59.4\pr & N4\dg 25.8\pr & 31 & 23 & 163\dg 59.6\pr & N4\dg 26.8\pr \\
\end{longtable}}

{\footnotesize
\begin{longtable}{rrrr|rrrr|rrrr}
\caption{Sun --- GHA and Declination, April 2026 (UT)}\label{sun04}\\
\toprule
\textbf{d} & \textbf{h} & \textbf{GHA} & \textbf{Dec} & \textbf{d} & \textbf{h} & \textbf{GHA} & \textbf{Dec} & \textbf{d} & \textbf{h} & \textbf{GHA} & \textbf{Dec} \\
\midrule
\endfirsthead
\multicolumn{12}{c}{\textit{Sun --- GHA and Declination, April 2026 (UT) (continued)}}\\
\toprule
\textbf{d} & \textbf{h} & \textbf{GHA} & \textbf{Dec} & \textbf{d} & \textbf{h} & \textbf{GHA} & \textbf{Dec} & \textbf{d} & \textbf{h} & \textbf{GHA} & \textbf{Dec} \\
\midrule
\endhead
\midrule\multicolumn{12}{r}{\textit{continued on next page}}\\
\endfoot
\bottomrule
\endlastfoot
1 & 0 & 178\dg 59.8\pr & N4\dg 27.8\pr & 1 & 1 & 194\dg 00.0\pr & N4\dg 28.7\pr & 1 & 2 & 209\dg 00.2\pr & N4\dg 29.7\pr \\
1 & 3 & 224\dg 00.4\pr & N4\dg 30.6\pr & 1 & 4 & 239\dg 00.5\pr & N4\dg 31.6\pr & 1 & 5 & 254\dg 00.7\pr & N4\dg 32.6\pr \\
1 & 6 & 269\dg 00.9\pr & N4\dg 33.5\pr & 1 & 7 & 284\dg 01.1\pr & N4\dg 34.5\pr & 1 & 8 & 299\dg 01.3\pr & N4\dg 35.5\pr \\
1 & 9 & 314\dg 01.5\pr & N4\dg 36.4\pr & 1 & 10 & 329\dg 01.7\pr & N4\dg 37.4\pr & 1 & 11 & 344\dg 01.8\pr & N4\dg 38.4\pr \\
1 & 12 & 359\dg 02.0\pr & N4\dg 39.3\pr & 1 & 13 & 14\dg 02.2\pr & N4\dg 40.3\pr & 1 & 14 & 29\dg 02.4\pr & N4\dg 41.3\pr \\
1 & 15 & 44\dg 02.6\pr & N4\dg 42.2\pr & 1 & 16 & 59\dg 02.8\pr & N4\dg 43.2\pr & 1 & 17 & 74\dg 03.0\pr & N4\dg 44.1\pr \\
1 & 18 & 89\dg 03.1\pr & N4\dg 45.1\pr & 1 & 19 & 104\dg 03.3\pr & N4\dg 46.1\pr & 1 & 20 & 119\dg 03.5\pr & N4\dg 47.0\pr \\
1 & 21 & 134\dg 03.7\pr & N4\dg 48.0\pr & 1 & 22 & 149\dg 03.9\pr & N4\dg 49.0\pr & 1 & 23 & 164\dg 04.1\pr & N4\dg 49.9\pr \\
2 & 0 & 179\dg 04.2\pr & N4\dg 50.9\pr & 2 & 1 & 194\dg 04.4\pr & N4\dg 51.8\pr & 2 & 2 & 209\dg 04.6\pr & N4\dg 52.8\pr \\
2 & 3 & 224\dg 04.8\pr & N4\dg 53.8\pr & 2 & 4 & 239\dg 05.0\pr & N4\dg 54.7\pr & 2 & 5 & 254\dg 05.2\pr & N4\dg 55.7\pr \\
2 & 6 & 269\dg 05.4\pr & N4\dg 56.7\pr & 2 & 7 & 284\dg 05.5\pr & N4\dg 57.6\pr & 2 & 8 & 299\dg 05.7\pr & N4\dg 58.6\pr \\
2 & 9 & 314\dg 05.9\pr & N4\dg 59.5\pr & 2 & 10 & 329\dg 06.1\pr & N5\dg 00.5\pr & 2 & 11 & 344\dg 06.3\pr & N5\dg 01.5\pr \\
2 & 12 & 359\dg 06.5\pr & N5\dg 02.4\pr & 2 & 13 & 14\dg 06.6\pr & N5\dg 03.4\pr & 2 & 14 & 29\dg 06.8\pr & N5\dg 04.3\pr \\
2 & 15 & 44\dg 07.0\pr & N5\dg 05.3\pr & 2 & 16 & 59\dg 07.2\pr & N5\dg 06.3\pr & 2 & 17 & 74\dg 07.4\pr & N5\dg 07.2\pr \\
2 & 18 & 89\dg 07.6\pr & N5\dg 08.2\pr & 2 & 19 & 104\dg 07.7\pr & N5\dg 09.1\pr & 2 & 20 & 119\dg 07.9\pr & N5\dg 10.1\pr \\
2 & 21 & 134\dg 08.1\pr & N5\dg 11.1\pr & 2 & 22 & 149\dg 08.3\pr & N5\dg 12.0\pr & 2 & 23 & 164\dg 08.5\pr & N5\dg 13.0\pr \\
3 & 0 & 179\dg 08.7\pr & N5\dg 13.9\pr & 3 & 1 & 194\dg 08.8\pr & N5\dg 14.9\pr & 3 & 2 & 209\dg 09.0\pr & N5\dg 15.8\pr \\
3 & 3 & 224\dg 09.2\pr & N5\dg 16.8\pr & 3 & 4 & 239\dg 09.4\pr & N5\dg 17.8\pr & 3 & 5 & 254\dg 09.6\pr & N5\dg 18.7\pr \\
3 & 6 & 269\dg 09.8\pr & N5\dg 19.7\pr & 3 & 7 & 284\dg 09.9\pr & N5\dg 20.6\pr & 3 & 8 & 299\dg 10.1\pr & N5\dg 21.6\pr \\
3 & 9 & 314\dg 10.3\pr & N5\dg 22.5\pr & 3 & 10 & 329\dg 10.5\pr & N5\dg 23.5\pr & 3 & 11 & 344\dg 10.7\pr & N5\dg 24.5\pr \\
3 & 12 & 359\dg 10.9\pr & N5\dg 25.4\pr & 3 & 13 & 14\dg 11.0\pr & N5\dg 26.4\pr & 3 & 14 & 29\dg 11.2\pr & N5\dg 27.3\pr \\
3 & 15 & 44\dg 11.4\pr & N5\dg 28.3\pr & 3 & 16 & 59\dg 11.6\pr & N5\dg 29.2\pr & 3 & 17 & 74\dg 11.8\pr & N5\dg 30.2\pr \\
3 & 18 & 89\dg 11.9\pr & N5\dg 31.2\pr & 3 & 19 & 104\dg 12.1\pr & N5\dg 32.1\pr & 3 & 20 & 119\dg 12.3\pr & N5\dg 33.1\pr \\
3 & 21 & 134\dg 12.5\pr & N5\dg 34.0\pr & 3 & 22 & 149\dg 12.7\pr & N5\dg 35.0\pr & 3 & 23 & 164\dg 12.9\pr & N5\dg 35.9\pr \\
4 & 0 & 179\dg 13.0\pr & N5\dg 36.9\pr & 4 & 1 & 194\dg 13.2\pr & N5\dg 37.8\pr & 4 & 2 & 209\dg 13.4\pr & N5\dg 38.8\pr \\
4 & 3 & 224\dg 13.6\pr & N5\dg 39.7\pr & 4 & 4 & 239\dg 13.8\pr & N5\dg 40.7\pr & 4 & 5 & 254\dg 13.9\pr & N5\dg 41.7\pr \\
4 & 6 & 269\dg 14.1\pr & N5\dg 42.6\pr & 4 & 7 & 284\dg 14.3\pr & N5\dg 43.6\pr & 4 & 8 & 299\dg 14.5\pr & N5\dg 44.5\pr \\
4 & 9 & 314\dg 14.7\pr & N5\dg 45.5\pr & 4 & 10 & 329\dg 14.9\pr & N5\dg 46.4\pr & 4 & 11 & 344\dg 15.0\pr & N5\dg 47.4\pr \\
4 & 12 & 359\dg 15.2\pr & N5\dg 48.3\pr & 4 & 13 & 14\dg 15.4\pr & N5\dg 49.3\pr & 4 & 14 & 29\dg 15.6\pr & N5\dg 50.2\pr \\
4 & 15 & 44\dg 15.8\pr & N5\dg 51.2\pr & 4 & 16 & 59\dg 15.9\pr & N5\dg 52.1\pr & 4 & 17 & 74\dg 16.1\pr & N5\dg 53.1\pr \\
4 & 18 & 89\dg 16.3\pr & N5\dg 54.0\pr & 4 & 19 & 104\dg 16.5\pr & N5\dg 55.0\pr & 4 & 20 & 119\dg 16.7\pr & N5\dg 55.9\pr \\
4 & 21 & 134\dg 16.8\pr & N5\dg 56.9\pr & 4 & 22 & 149\dg 17.0\pr & N5\dg 57.8\pr & 4 & 23 & 164\dg 17.2\pr & N5\dg 58.8\pr \\
5 & 0 & 179\dg 17.4\pr & N5\dg 59.7\pr & 5 & 1 & 194\dg 17.6\pr & N6\dg 00.7\pr & 5 & 2 & 209\dg 17.7\pr & N6\dg 01.6\pr \\
5 & 3 & 224\dg 17.9\pr & N6\dg 02.6\pr & 5 & 4 & 239\dg 18.1\pr & N6\dg 03.5\pr & 5 & 5 & 254\dg 18.3\pr & N6\dg 04.5\pr \\
5 & 6 & 269\dg 18.5\pr & N6\dg 05.4\pr & 5 & 7 & 284\dg 18.6\pr & N6\dg 06.4\pr & 5 & 8 & 299\dg 18.8\pr & N6\dg 07.3\pr \\
5 & 9 & 314\dg 19.0\pr & N6\dg 08.3\pr & 5 & 10 & 329\dg 19.2\pr & N6\dg 09.2\pr & 5 & 11 & 344\dg 19.4\pr & N6\dg 10.2\pr \\
5 & 12 & 359\dg 19.5\pr & N6\dg 11.1\pr & 5 & 13 & 14\dg 19.7\pr & N6\dg 12.1\pr & 5 & 14 & 29\dg 19.9\pr & N6\dg 13.0\pr \\
5 & 15 & 44\dg 20.1\pr & N6\dg 14.0\pr & 5 & 16 & 59\dg 20.2\pr & N6\dg 14.9\pr & 5 & 17 & 74\dg 20.4\pr & N6\dg 15.9\pr \\
5 & 18 & 89\dg 20.6\pr & N6\dg 16.8\pr & 5 & 19 & 104\dg 20.8\pr & N6\dg 17.8\pr & 5 & 20 & 119\dg 21.0\pr & N6\dg 18.7\pr \\
5 & 21 & 134\dg 21.1\pr & N6\dg 19.7\pr & 5 & 22 & 149\dg 21.3\pr & N6\dg 20.6\pr & 5 & 23 & 164\dg 21.5\pr & N6\dg 21.5\pr \\
6 & 0 & 179\dg 21.7\pr & N6\dg 22.5\pr & 6 & 1 & 194\dg 21.8\pr & N6\dg 23.4\pr & 6 & 2 & 209\dg 22.0\pr & N6\dg 24.4\pr \\
6 & 3 & 224\dg 22.2\pr & N6\dg 25.3\pr & 6 & 4 & 239\dg 22.4\pr & N6\dg 26.3\pr & 6 & 5 & 254\dg 22.6\pr & N6\dg 27.2\pr \\
6 & 6 & 269\dg 22.7\pr & N6\dg 28.2\pr & 6 & 7 & 284\dg 22.9\pr & N6\dg 29.1\pr & 6 & 8 & 299\dg 23.1\pr & N6\dg 30.0\pr \\
6 & 9 & 314\dg 23.3\pr & N6\dg 31.0\pr & 6 & 10 & 329\dg 23.4\pr & N6\dg 31.9\pr & 6 & 11 & 344\dg 23.6\pr & N6\dg 32.9\pr \\
6 & 12 & 359\dg 23.8\pr & N6\dg 33.8\pr & 6 & 13 & 14\dg 24.0\pr & N6\dg 34.8\pr & 6 & 14 & 29\dg 24.2\pr & N6\dg 35.7\pr \\
6 & 15 & 44\dg 24.3\pr & N6\dg 36.7\pr & 6 & 16 & 59\dg 24.5\pr & N6\dg 37.6\pr & 6 & 17 & 74\dg 24.7\pr & N6\dg 38.5\pr \\
6 & 18 & 89\dg 24.9\pr & N6\dg 39.5\pr & 6 & 19 & 104\dg 25.0\pr & N6\dg 40.4\pr & 6 & 20 & 119\dg 25.2\pr & N6\dg 41.4\pr \\
6 & 21 & 134\dg 25.4\pr & N6\dg 42.3\pr & 6 & 22 & 149\dg 25.6\pr & N6\dg 43.2\pr & 6 & 23 & 164\dg 25.7\pr & N6\dg 44.2\pr \\
7 & 0 & 179\dg 25.9\pr & N6\dg 45.1\pr & 7 & 1 & 194\dg 26.1\pr & N6\dg 46.1\pr & 7 & 2 & 209\dg 26.3\pr & N6\dg 47.0\pr \\
7 & 3 & 224\dg 26.4\pr & N6\dg 48.0\pr & 7 & 4 & 239\dg 26.6\pr & N6\dg 48.9\pr & 7 & 5 & 254\dg 26.8\pr & N6\dg 49.8\pr \\
7 & 6 & 269\dg 27.0\pr & N6\dg 50.8\pr & 7 & 7 & 284\dg 27.1\pr & N6\dg 51.7\pr & 7 & 8 & 299\dg 27.3\pr & N6\dg 52.7\pr \\
7 & 9 & 314\dg 27.5\pr & N6\dg 53.6\pr & 7 & 10 & 329\dg 27.7\pr & N6\dg 54.5\pr & 7 & 11 & 344\dg 27.8\pr & N6\dg 55.5\pr \\
7 & 12 & 359\dg 28.0\pr & N6\dg 56.4\pr & 7 & 13 & 14\dg 28.2\pr & N6\dg 57.3\pr & 7 & 14 & 29\dg 28.4\pr & N6\dg 58.3\pr \\
7 & 15 & 44\dg 28.5\pr & N6\dg 59.2\pr & 7 & 16 & 59\dg 28.7\pr & N7\dg 00.2\pr & 7 & 17 & 74\dg 28.9\pr & N7\dg 01.1\pr \\
7 & 18 & 89\dg 29.1\pr & N7\dg 02.0\pr & 7 & 19 & 104\dg 29.2\pr & N7\dg 03.0\pr & 7 & 20 & 119\dg 29.4\pr & N7\dg 03.9\pr \\
7 & 21 & 134\dg 29.6\pr & N7\dg 04.8\pr & 7 & 22 & 149\dg 29.8\pr & N7\dg 05.8\pr & 7 & 23 & 164\dg 29.9\pr & N7\dg 06.7\pr \\
8 & 0 & 179\dg 30.1\pr & N7\dg 07.7\pr & 8 & 1 & 194\dg 30.3\pr & N7\dg 08.6\pr & 8 & 2 & 209\dg 30.4\pr & N7\dg 09.5\pr \\
8 & 3 & 224\dg 30.6\pr & N7\dg 10.5\pr & 8 & 4 & 239\dg 30.8\pr & N7\dg 11.4\pr & 8 & 5 & 254\dg 31.0\pr & N7\dg 12.3\pr \\
8 & 6 & 269\dg 31.1\pr & N7\dg 13.3\pr & 8 & 7 & 284\dg 31.3\pr & N7\dg 14.2\pr & 8 & 8 & 299\dg 31.5\pr & N7\dg 15.1\pr \\
8 & 9 & 314\dg 31.7\pr & N7\dg 16.1\pr & 8 & 10 & 329\dg 31.8\pr & N7\dg 17.0\pr & 8 & 11 & 344\dg 32.0\pr & N7\dg 17.9\pr \\
8 & 12 & 359\dg 32.2\pr & N7\dg 18.9\pr & 8 & 13 & 14\dg 32.3\pr & N7\dg 19.8\pr & 8 & 14 & 29\dg 32.5\pr & N7\dg 20.7\pr \\
8 & 15 & 44\dg 32.7\pr & N7\dg 21.7\pr & 8 & 16 & 59\dg 32.9\pr & N7\dg 22.6\pr & 8 & 17 & 74\dg 33.0\pr & N7\dg 23.5\pr \\
8 & 18 & 89\dg 33.2\pr & N7\dg 24.5\pr & 8 & 19 & 104\dg 33.4\pr & N7\dg 25.4\pr & 8 & 20 & 119\dg 33.5\pr & N7\dg 26.3\pr \\
8 & 21 & 134\dg 33.7\pr & N7\dg 27.3\pr & 8 & 22 & 149\dg 33.9\pr & N7\dg 28.2\pr & 8 & 23 & 164\dg 34.1\pr & N7\dg 29.1\pr \\
9 & 0 & 179\dg 34.2\pr & N7\dg 30.1\pr & 9 & 1 & 194\dg 34.4\pr & N7\dg 31.0\pr & 9 & 2 & 209\dg 34.6\pr & N7\dg 31.9\pr \\
9 & 3 & 224\dg 34.7\pr & N7\dg 32.9\pr & 9 & 4 & 239\dg 34.9\pr & N7\dg 33.8\pr & 9 & 5 & 254\dg 35.1\pr & N7\dg 34.7\pr \\
9 & 6 & 269\dg 35.3\pr & N7\dg 35.6\pr & 9 & 7 & 284\dg 35.4\pr & N7\dg 36.6\pr & 9 & 8 & 299\dg 35.6\pr & N7\dg 37.5\pr \\
9 & 9 & 314\dg 35.8\pr & N7\dg 38.4\pr & 9 & 10 & 329\dg 35.9\pr & N7\dg 39.4\pr & 9 & 11 & 344\dg 36.1\pr & N7\dg 40.3\pr \\
9 & 12 & 359\dg 36.3\pr & N7\dg 41.2\pr & 9 & 13 & 14\dg 36.4\pr & N7\dg 42.1\pr & 9 & 14 & 29\dg 36.6\pr & N7\dg 43.1\pr \\
9 & 15 & 44\dg 36.8\pr & N7\dg 44.0\pr & 9 & 16 & 59\dg 37.0\pr & N7\dg 44.9\pr & 9 & 17 & 74\dg 37.1\pr & N7\dg 45.9\pr \\
9 & 18 & 89\dg 37.3\pr & N7\dg 46.8\pr & 9 & 19 & 104\dg 37.5\pr & N7\dg 47.7\pr & 9 & 20 & 119\dg 37.6\pr & N7\dg 48.6\pr \\
9 & 21 & 134\dg 37.8\pr & N7\dg 49.6\pr & 9 & 22 & 149\dg 38.0\pr & N7\dg 50.5\pr & 9 & 23 & 164\dg 38.1\pr & N7\dg 51.4\pr \\
10 & 0 & 179\dg 38.3\pr & N7\dg 52.3\pr & 10 & 1 & 194\dg 38.5\pr & N7\dg 53.3\pr & 10 & 2 & 209\dg 38.6\pr & N7\dg 54.2\pr \\
10 & 3 & 224\dg 38.8\pr & N7\dg 55.1\pr & 10 & 4 & 239\dg 39.0\pr & N7\dg 56.0\pr & 10 & 5 & 254\dg 39.1\pr & N7\dg 57.0\pr \\
10 & 6 & 269\dg 39.3\pr & N7\dg 57.9\pr & 10 & 7 & 284\dg 39.5\pr & N7\dg 58.8\pr & 10 & 8 & 299\dg 39.6\pr & N7\dg 59.7\pr \\
10 & 9 & 314\dg 39.8\pr & N8\dg 00.7\pr & 10 & 10 & 329\dg 40.0\pr & N8\dg 01.6\pr & 10 & 11 & 344\dg 40.1\pr & N8\dg 02.5\pr \\
10 & 12 & 359\dg 40.3\pr & N8\dg 03.4\pr & 10 & 13 & 14\dg 40.5\pr & N8\dg 04.4\pr & 10 & 14 & 29\dg 40.6\pr & N8\dg 05.3\pr \\
10 & 15 & 44\dg 40.8\pr & N8\dg 06.2\pr & 10 & 16 & 59\dg 41.0\pr & N8\dg 07.1\pr & 10 & 17 & 74\dg 41.1\pr & N8\dg 08.0\pr \\
10 & 18 & 89\dg 41.3\pr & N8\dg 09.0\pr & 10 & 19 & 104\dg 41.5\pr & N8\dg 09.9\pr & 10 & 20 & 119\dg 41.6\pr & N8\dg 10.8\pr \\
10 & 21 & 134\dg 41.8\pr & N8\dg 11.7\pr & 10 & 22 & 149\dg 42.0\pr & N8\dg 12.6\pr & 10 & 23 & 164\dg 42.1\pr & N8\dg 13.6\pr \\
11 & 0 & 179\dg 42.3\pr & N8\dg 14.5\pr & 11 & 1 & 194\dg 42.5\pr & N8\dg 15.4\pr & 11 & 2 & 209\dg 42.6\pr & N8\dg 16.3\pr \\
11 & 3 & 224\dg 42.8\pr & N8\dg 17.2\pr & 11 & 4 & 239\dg 43.0\pr & N8\dg 18.2\pr & 11 & 5 & 254\dg 43.1\pr & N8\dg 19.1\pr \\
11 & 6 & 269\dg 43.3\pr & N8\dg 20.0\pr & 11 & 7 & 284\dg 43.5\pr & N8\dg 20.9\pr & 11 & 8 & 299\dg 43.6\pr & N8\dg 21.8\pr \\
11 & 9 & 314\dg 43.8\pr & N8\dg 22.8\pr & 11 & 10 & 329\dg 43.9\pr & N8\dg 23.7\pr & 11 & 11 & 344\dg 44.1\pr & N8\dg 24.6\pr \\
11 & 12 & 359\dg 44.3\pr & N8\dg 25.5\pr & 11 & 13 & 14\dg 44.4\pr & N8\dg 26.4\pr & 11 & 14 & 29\dg 44.6\pr & N8\dg 27.3\pr \\
11 & 15 & 44\dg 44.8\pr & N8\dg 28.3\pr & 11 & 16 & 59\dg 44.9\pr & N8\dg 29.2\pr & 11 & 17 & 74\dg 45.1\pr & N8\dg 30.1\pr \\
11 & 18 & 89\dg 45.3\pr & N8\dg 31.0\pr & 11 & 19 & 104\dg 45.4\pr & N8\dg 31.9\pr & 11 & 20 & 119\dg 45.6\pr & N8\dg 32.8\pr \\
11 & 21 & 134\dg 45.7\pr & N8\dg 33.8\pr & 11 & 22 & 149\dg 45.9\pr & N8\dg 34.7\pr & 11 & 23 & 164\dg 46.1\pr & N8\dg 35.6\pr \\
12 & 0 & 179\dg 46.2\pr & N8\dg 36.5\pr & 12 & 1 & 194\dg 46.4\pr & N8\dg 37.4\pr & 12 & 2 & 209\dg 46.6\pr & N8\dg 38.3\pr \\
12 & 3 & 224\dg 46.7\pr & N8\dg 39.2\pr & 12 & 4 & 239\dg 46.9\pr & N8\dg 40.1\pr & 12 & 5 & 254\dg 47.0\pr & N8\dg 41.1\pr \\
12 & 6 & 269\dg 47.2\pr & N8\dg 42.0\pr & 12 & 7 & 284\dg 47.4\pr & N8\dg 42.9\pr & 12 & 8 & 299\dg 47.5\pr & N8\dg 43.8\pr \\
12 & 9 & 314\dg 47.7\pr & N8\dg 44.7\pr & 12 & 10 & 329\dg 47.8\pr & N8\dg 45.6\pr & 12 & 11 & 344\dg 48.0\pr & N8\dg 46.5\pr \\
12 & 12 & 359\dg 48.2\pr & N8\dg 47.4\pr & 12 & 13 & 14\dg 48.3\pr & N8\dg 48.4\pr & 12 & 14 & 29\dg 48.5\pr & N8\dg 49.3\pr \\
12 & 15 & 44\dg 48.7\pr & N8\dg 50.2\pr & 12 & 16 & 59\dg 48.8\pr & N8\dg 51.1\pr & 12 & 17 & 74\dg 49.0\pr & N8\dg 52.0\pr \\
12 & 18 & 89\dg 49.1\pr & N8\dg 52.9\pr & 12 & 19 & 104\dg 49.3\pr & N8\dg 53.8\pr & 12 & 20 & 119\dg 49.5\pr & N8\dg 54.7\pr \\
12 & 21 & 134\dg 49.6\pr & N8\dg 55.6\pr & 12 & 22 & 149\dg 49.8\pr & N8\dg 56.5\pr & 12 & 23 & 164\dg 49.9\pr & N8\dg 57.5\pr \\
13 & 0 & 179\dg 50.1\pr & N8\dg 58.4\pr & 13 & 1 & 194\dg 50.2\pr & N8\dg 59.3\pr & 13 & 2 & 209\dg 50.4\pr & N9\dg 00.2\pr \\
13 & 3 & 224\dg 50.6\pr & N9\dg 01.1\pr & 13 & 4 & 239\dg 50.7\pr & N9\dg 02.0\pr & 13 & 5 & 254\dg 50.9\pr & N9\dg 02.9\pr \\
13 & 6 & 269\dg 51.0\pr & N9\dg 03.8\pr & 13 & 7 & 284\dg 51.2\pr & N9\dg 04.7\pr & 13 & 8 & 299\dg 51.4\pr & N9\dg 05.6\pr \\
13 & 9 & 314\dg 51.5\pr & N9\dg 06.5\pr & 13 & 10 & 329\dg 51.7\pr & N9\dg 07.4\pr & 13 & 11 & 344\dg 51.8\pr & N9\dg 08.3\pr \\
13 & 12 & 359\dg 52.0\pr & N9\dg 09.2\pr & 13 & 13 & 14\dg 52.1\pr & N9\dg 10.1\pr & 13 & 14 & 29\dg 52.3\pr & N9\dg 11.0\pr \\
13 & 15 & 44\dg 52.5\pr & N9\dg 11.9\pr & 13 & 16 & 59\dg 52.6\pr & N9\dg 12.9\pr & 13 & 17 & 74\dg 52.8\pr & N9\dg 13.8\pr \\
13 & 18 & 89\dg 52.9\pr & N9\dg 14.7\pr & 13 & 19 & 104\dg 53.1\pr & N9\dg 15.6\pr & 13 & 20 & 119\dg 53.2\pr & N9\dg 16.5\pr \\
13 & 21 & 134\dg 53.4\pr & N9\dg 17.4\pr & 13 & 22 & 149\dg 53.6\pr & N9\dg 18.3\pr & 13 & 23 & 164\dg 53.7\pr & N9\dg 19.2\pr \\
14 & 0 & 179\dg 53.9\pr & N9\dg 20.1\pr & 14 & 1 & 194\dg 54.0\pr & N9\dg 21.0\pr & 14 & 2 & 209\dg 54.2\pr & N9\dg 21.9\pr \\
14 & 3 & 224\dg 54.3\pr & N9\dg 22.8\pr & 14 & 4 & 239\dg 54.5\pr & N9\dg 23.7\pr & 14 & 5 & 254\dg 54.6\pr & N9\dg 24.6\pr \\
14 & 6 & 269\dg 54.8\pr & N9\dg 25.5\pr & 14 & 7 & 284\dg 54.9\pr & N9\dg 26.4\pr & 14 & 8 & 299\dg 55.1\pr & N9\dg 27.3\pr \\
14 & 9 & 314\dg 55.3\pr & N9\dg 28.2\pr & 14 & 10 & 329\dg 55.4\pr & N9\dg 29.1\pr & 14 & 11 & 344\dg 55.6\pr & N9\dg 30.0\pr \\
14 & 12 & 359\dg 55.7\pr & N9\dg 30.9\pr & 14 & 13 & 14\dg 55.9\pr & N9\dg 31.8\pr & 14 & 14 & 29\dg 56.0\pr & N9\dg 32.7\pr \\
14 & 15 & 44\dg 56.2\pr & N9\dg 33.6\pr & 14 & 16 & 59\dg 56.3\pr & N9\dg 34.5\pr & 14 & 17 & 74\dg 56.5\pr & N9\dg 35.4\pr \\
14 & 18 & 89\dg 56.6\pr & N9\dg 36.3\pr & 14 & 19 & 104\dg 56.8\pr & N9\dg 37.2\pr & 14 & 20 & 119\dg 56.9\pr & N9\dg 38.0\pr \\
14 & 21 & 134\dg 57.1\pr & N9\dg 38.9\pr & 14 & 22 & 149\dg 57.2\pr & N9\dg 39.8\pr & 14 & 23 & 164\dg 57.4\pr & N9\dg 40.7\pr \\
15 & 0 & 179\dg 57.6\pr & N9\dg 41.6\pr & 15 & 1 & 194\dg 57.7\pr & N9\dg 42.5\pr & 15 & 2 & 209\dg 57.9\pr & N9\dg 43.4\pr \\
15 & 3 & 224\dg 58.0\pr & N9\dg 44.3\pr & 15 & 4 & 239\dg 58.2\pr & N9\dg 45.2\pr & 15 & 5 & 254\dg 58.3\pr & N9\dg 46.1\pr \\
15 & 6 & 269\dg 58.5\pr & N9\dg 47.0\pr & 15 & 7 & 284\dg 58.6\pr & N9\dg 47.9\pr & 15 & 8 & 299\dg 58.8\pr & N9\dg 48.8\pr \\
15 & 9 & 314\dg 58.9\pr & N9\dg 49.7\pr & 15 & 10 & 329\dg 59.1\pr & N9\dg 50.6\pr & 15 & 11 & 344\dg 59.2\pr & N9\dg 51.5\pr \\
15 & 12 & 359\dg 59.4\pr & N9\dg 52.4\pr & 15 & 13 & 14\dg 59.5\pr & N9\dg 53.2\pr & 15 & 14 & 29\dg 59.7\pr & N9\dg 54.1\pr \\
15 & 15 & 44\dg 59.8\pr & N9\dg 55.0\pr & 15 & 16 & 60\dg 00.0\pr & N9\dg 55.9\pr & 15 & 17 & 75\dg 00.1\pr & N9\dg 56.8\pr \\
15 & 18 & 90\dg 00.3\pr & N9\dg 57.7\pr & 15 & 19 & 105\dg 00.4\pr & N9\dg 58.6\pr & 15 & 20 & 120\dg 00.6\pr & N9\dg 59.5\pr \\
15 & 21 & 135\dg 00.7\pr & N10\dg 00.4\pr & 15 & 22 & 150\dg 00.9\pr & N10\dg 01.3\pr & 15 & 23 & 165\dg 01.0\pr & N10\dg 02.1\pr \\
16 & 0 & 180\dg 01.2\pr & N10\dg 03.0\pr & 16 & 1 & 195\dg 01.3\pr & N10\dg 03.9\pr & 16 & 2 & 210\dg 01.5\pr & N10\dg 04.8\pr \\
16 & 3 & 225\dg 01.6\pr & N10\dg 05.7\pr & 16 & 4 & 240\dg 01.8\pr & N10\dg 06.6\pr & 16 & 5 & 255\dg 01.9\pr & N10\dg 07.5\pr \\
16 & 6 & 270\dg 02.0\pr & N10\dg 08.4\pr & 16 & 7 & 285\dg 02.2\pr & N10\dg 09.2\pr & 16 & 8 & 300\dg 02.3\pr & N10\dg 10.1\pr \\
16 & 9 & 315\dg 02.5\pr & N10\dg 11.0\pr & 16 & 10 & 330\dg 02.6\pr & N10\dg 11.9\pr & 16 & 11 & 345\dg 02.8\pr & N10\dg 12.8\pr \\
16 & 12 & 0\dg 02.9\pr & N10\dg 13.7\pr & 16 & 13 & 15\dg 03.1\pr & N10\dg 14.6\pr & 16 & 14 & 30\dg 03.2\pr & N10\dg 15.4\pr \\
16 & 15 & 45\dg 03.4\pr & N10\dg 16.3\pr & 16 & 16 & 60\dg 03.5\pr & N10\dg 17.2\pr & 16 & 17 & 75\dg 03.7\pr & N10\dg 18.1\pr \\
16 & 18 & 90\dg 03.8\pr & N10\dg 19.0\pr & 16 & 19 & 105\dg 03.9\pr & N10\dg 19.9\pr & 16 & 20 & 120\dg 04.1\pr & N10\dg 20.7\pr \\
16 & 21 & 135\dg 04.2\pr & N10\dg 21.6\pr & 16 & 22 & 150\dg 04.4\pr & N10\dg 22.5\pr & 16 & 23 & 165\dg 04.5\pr & N10\dg 23.4\pr \\
17 & 0 & 180\dg 04.7\pr & N10\dg 24.3\pr & 17 & 1 & 195\dg 04.8\pr & N10\dg 25.1\pr & 17 & 2 & 210\dg 05.0\pr & N10\dg 26.0\pr \\
17 & 3 & 225\dg 05.1\pr & N10\dg 26.9\pr & 17 & 4 & 240\dg 05.2\pr & N10\dg 27.8\pr & 17 & 5 & 255\dg 05.4\pr & N10\dg 28.7\pr \\
17 & 6 & 270\dg 05.5\pr & N10\dg 29.5\pr & 17 & 7 & 285\dg 05.7\pr & N10\dg 30.4\pr & 17 & 8 & 300\dg 05.8\pr & N10\dg 31.3\pr \\
17 & 9 & 315\dg 06.0\pr & N10\dg 32.2\pr & 17 & 10 & 330\dg 06.1\pr & N10\dg 33.1\pr & 17 & 11 & 345\dg 06.3\pr & N10\dg 33.9\pr \\
17 & 12 & 0\dg 06.4\pr & N10\dg 34.8\pr & 17 & 13 & 15\dg 06.5\pr & N10\dg 35.7\pr & 17 & 14 & 30\dg 06.7\pr & N10\dg 36.6\pr \\
17 & 15 & 45\dg 06.8\pr & N10\dg 37.4\pr & 17 & 16 & 60\dg 07.0\pr & N10\dg 38.3\pr & 17 & 17 & 75\dg 07.1\pr & N10\dg 39.2\pr \\
17 & 18 & 90\dg 07.2\pr & N10\dg 40.1\pr & 17 & 19 & 105\dg 07.4\pr & N10\dg 40.9\pr & 17 & 20 & 120\dg 07.5\pr & N10\dg 41.8\pr \\
17 & 21 & 135\dg 07.7\pr & N10\dg 42.7\pr & 17 & 22 & 150\dg 07.8\pr & N10\dg 43.6\pr & 17 & 23 & 165\dg 08.0\pr & N10\dg 44.4\pr \\
18 & 0 & 180\dg 08.1\pr & N10\dg 45.3\pr & 18 & 1 & 195\dg 08.2\pr & N10\dg 46.2\pr & 18 & 2 & 210\dg 08.4\pr & N10\dg 47.1\pr \\
18 & 3 & 225\dg 08.5\pr & N10\dg 47.9\pr & 18 & 4 & 240\dg 08.7\pr & N10\dg 48.8\pr & 18 & 5 & 255\dg 08.8\pr & N10\dg 49.7\pr \\
18 & 6 & 270\dg 08.9\pr & N10\dg 50.6\pr & 18 & 7 & 285\dg 09.1\pr & N10\dg 51.4\pr & 18 & 8 & 300\dg 09.2\pr & N10\dg 52.3\pr \\
18 & 9 & 315\dg 09.3\pr & N10\dg 53.2\pr & 18 & 10 & 330\dg 09.5\pr & N10\dg 54.0\pr & 18 & 11 & 345\dg 09.6\pr & N10\dg 54.9\pr \\
18 & 12 & 0\dg 09.8\pr & N10\dg 55.8\pr & 18 & 13 & 15\dg 09.9\pr & N10\dg 56.7\pr & 18 & 14 & 30\dg 10.0\pr & N10\dg 57.5\pr \\
18 & 15 & 45\dg 10.2\pr & N10\dg 58.4\pr & 18 & 16 & 60\dg 10.3\pr & N10\dg 59.3\pr & 18 & 17 & 75\dg 10.5\pr & N11\dg 00.1\pr \\
18 & 18 & 90\dg 10.6\pr & N11\dg 01.0\pr & 18 & 19 & 105\dg 10.7\pr & N11\dg 01.9\pr & 18 & 20 & 120\dg 10.9\pr & N11\dg 02.7\pr \\
18 & 21 & 135\dg 11.0\pr & N11\dg 03.6\pr & 18 & 22 & 150\dg 11.1\pr & N11\dg 04.5\pr & 18 & 23 & 165\dg 11.3\pr & N11\dg 05.3\pr \\
19 & 0 & 180\dg 11.4\pr & N11\dg 06.2\pr & 19 & 1 & 195\dg 11.5\pr & N11\dg 07.1\pr & 19 & 2 & 210\dg 11.7\pr & N11\dg 07.9\pr \\
19 & 3 & 225\dg 11.8\pr & N11\dg 08.8\pr & 19 & 4 & 240\dg 12.0\pr & N11\dg 09.7\pr & 19 & 5 & 255\dg 12.1\pr & N11\dg 10.5\pr \\
19 & 6 & 270\dg 12.2\pr & N11\dg 11.4\pr & 19 & 7 & 285\dg 12.4\pr & N11\dg 12.3\pr & 19 & 8 & 300\dg 12.5\pr & N11\dg 13.1\pr \\
19 & 9 & 315\dg 12.6\pr & N11\dg 14.0\pr & 19 & 10 & 330\dg 12.8\pr & N11\dg 14.9\pr & 19 & 11 & 345\dg 12.9\pr & N11\dg 15.7\pr \\
19 & 12 & 0\dg 13.0\pr & N11\dg 16.6\pr & 19 & 13 & 15\dg 13.2\pr & N11\dg 17.4\pr & 19 & 14 & 30\dg 13.3\pr & N11\dg 18.3\pr \\
19 & 15 & 45\dg 13.4\pr & N11\dg 19.2\pr & 19 & 16 & 60\dg 13.6\pr & N11\dg 20.0\pr & 19 & 17 & 75\dg 13.7\pr & N11\dg 20.9\pr \\
19 & 18 & 90\dg 13.8\pr & N11\dg 21.7\pr & 19 & 19 & 105\dg 14.0\pr & N11\dg 22.6\pr & 19 & 20 & 120\dg 14.1\pr & N11\dg 23.5\pr \\
19 & 21 & 135\dg 14.2\pr & N11\dg 24.3\pr & 19 & 22 & 150\dg 14.4\pr & N11\dg 25.2\pr & 19 & 23 & 165\dg 14.5\pr & N11\dg 26.0\pr \\
20 & 0 & 180\dg 14.6\pr & N11\dg 26.9\pr & 20 & 1 & 195\dg 14.8\pr & N11\dg 27.8\pr & 20 & 2 & 210\dg 14.9\pr & N11\dg 28.6\pr \\
20 & 3 & 225\dg 15.0\pr & N11\dg 29.5\pr & 20 & 4 & 240\dg 15.2\pr & N11\dg 30.3\pr & 20 & 5 & 255\dg 15.3\pr & N11\dg 31.2\pr \\
20 & 6 & 270\dg 15.4\pr & N11\dg 32.1\pr & 20 & 7 & 285\dg 15.5\pr & N11\dg 32.9\pr & 20 & 8 & 300\dg 15.7\pr & N11\dg 33.8\pr \\
20 & 9 & 315\dg 15.8\pr & N11\dg 34.6\pr & 20 & 10 & 330\dg 15.9\pr & N11\dg 35.5\pr & 20 & 11 & 345\dg 16.1\pr & N11\dg 36.3\pr \\
20 & 12 & 0\dg 16.2\pr & N11\dg 37.2\pr & 20 & 13 & 15\dg 16.3\pr & N11\dg 38.0\pr & 20 & 14 & 30\dg 16.5\pr & N11\dg 38.9\pr \\
20 & 15 & 45\dg 16.6\pr & N11\dg 39.7\pr & 20 & 16 & 60\dg 16.7\pr & N11\dg 40.6\pr & 20 & 17 & 75\dg 16.8\pr & N11\dg 41.5\pr \\
20 & 18 & 90\dg 17.0\pr & N11\dg 42.3\pr & 20 & 19 & 105\dg 17.1\pr & N11\dg 43.2\pr & 20 & 20 & 120\dg 17.2\pr & N11\dg 44.0\pr \\
20 & 21 & 135\dg 17.4\pr & N11\dg 44.9\pr & 20 & 22 & 150\dg 17.5\pr & N11\dg 45.7\pr & 20 & 23 & 165\dg 17.6\pr & N11\dg 46.6\pr \\
21 & 0 & 180\dg 17.7\pr & N11\dg 47.4\pr & 21 & 1 & 195\dg 17.9\pr & N11\dg 48.3\pr & 21 & 2 & 210\dg 18.0\pr & N11\dg 49.1\pr \\
21 & 3 & 225\dg 18.1\pr & N11\dg 50.0\pr & 21 & 4 & 240\dg 18.3\pr & N11\dg 50.8\pr & 21 & 5 & 255\dg 18.4\pr & N11\dg 51.7\pr \\
21 & 6 & 270\dg 18.5\pr & N11\dg 52.5\pr & 21 & 7 & 285\dg 18.6\pr & N11\dg 53.4\pr & 21 & 8 & 300\dg 18.8\pr & N11\dg 54.2\pr \\
21 & 9 & 315\dg 18.9\pr & N11\dg 55.1\pr & 21 & 10 & 330\dg 19.0\pr & N11\dg 55.9\pr & 21 & 11 & 345\dg 19.1\pr & N11\dg 56.8\pr \\
21 & 12 & 0\dg 19.3\pr & N11\dg 57.6\pr & 21 & 13 & 15\dg 19.4\pr & N11\dg 58.4\pr & 21 & 14 & 30\dg 19.5\pr & N11\dg 59.3\pr \\
21 & 15 & 45\dg 19.6\pr & N12\dg 00.1\pr & 21 & 16 & 60\dg 19.8\pr & N12\dg 01.0\pr & 21 & 17 & 75\dg 19.9\pr & N12\dg 01.8\pr \\
21 & 18 & 90\dg 20.0\pr & N12\dg 02.7\pr & 21 & 19 & 105\dg 20.1\pr & N12\dg 03.5\pr & 21 & 20 & 120\dg 20.3\pr & N12\dg 04.4\pr \\
21 & 21 & 135\dg 20.4\pr & N12\dg 05.2\pr & 21 & 22 & 150\dg 20.5\pr & N12\dg 06.0\pr & 21 & 23 & 165\dg 20.6\pr & N12\dg 06.9\pr \\
22 & 0 & 180\dg 20.7\pr & N12\dg 07.7\pr & 22 & 1 & 195\dg 20.9\pr & N12\dg 08.6\pr & 22 & 2 & 210\dg 21.0\pr & N12\dg 09.4\pr \\
22 & 3 & 225\dg 21.1\pr & N12\dg 10.3\pr & 22 & 4 & 240\dg 21.2\pr & N12\dg 11.1\pr & 22 & 5 & 255\dg 21.4\pr & N12\dg 11.9\pr \\
22 & 6 & 270\dg 21.5\pr & N12\dg 12.8\pr & 22 & 7 & 285\dg 21.6\pr & N12\dg 13.6\pr & 22 & 8 & 300\dg 21.7\pr & N12\dg 14.5\pr \\
22 & 9 & 315\dg 21.8\pr & N12\dg 15.3\pr & 22 & 10 & 330\dg 22.0\pr & N12\dg 16.1\pr & 22 & 11 & 345\dg 22.1\pr & N12\dg 17.0\pr \\
22 & 12 & 0\dg 22.2\pr & N12\dg 17.8\pr & 22 & 13 & 15\dg 22.3\pr & N12\dg 18.7\pr & 22 & 14 & 30\dg 22.5\pr & N12\dg 19.5\pr \\
22 & 15 & 45\dg 22.6\pr & N12\dg 20.3\pr & 22 & 16 & 60\dg 22.7\pr & N12\dg 21.2\pr & 22 & 17 & 75\dg 22.8\pr & N12\dg 22.0\pr \\
22 & 18 & 90\dg 22.9\pr & N12\dg 22.8\pr & 22 & 19 & 105\dg 23.0\pr & N12\dg 23.7\pr & 22 & 20 & 120\dg 23.2\pr & N12\dg 24.5\pr \\
22 & 21 & 135\dg 23.3\pr & N12\dg 25.3\pr & 22 & 22 & 150\dg 23.4\pr & N12\dg 26.2\pr & 22 & 23 & 165\dg 23.5\pr & N12\dg 27.0\pr \\
23 & 0 & 180\dg 23.6\pr & N12\dg 27.9\pr & 23 & 1 & 195\dg 23.8\pr & N12\dg 28.7\pr & 23 & 2 & 210\dg 23.9\pr & N12\dg 29.5\pr \\
23 & 3 & 225\dg 24.0\pr & N12\dg 30.4\pr & 23 & 4 & 240\dg 24.1\pr & N12\dg 31.2\pr & 23 & 5 & 255\dg 24.2\pr & N12\dg 32.0\pr \\
23 & 6 & 270\dg 24.3\pr & N12\dg 32.8\pr & 23 & 7 & 285\dg 24.5\pr & N12\dg 33.7\pr & 23 & 8 & 300\dg 24.6\pr & N12\dg 34.5\pr \\
23 & 9 & 315\dg 24.7\pr & N12\dg 35.3\pr & 23 & 10 & 330\dg 24.8\pr & N12\dg 36.2\pr & 23 & 11 & 345\dg 24.9\pr & N12\dg 37.0\pr \\
23 & 12 & 0\dg 25.0\pr & N12\dg 37.8\pr & 23 & 13 & 15\dg 25.2\pr & N12\dg 38.7\pr & 23 & 14 & 30\dg 25.3\pr & N12\dg 39.5\pr \\
23 & 15 & 45\dg 25.4\pr & N12\dg 40.3\pr & 23 & 16 & 60\dg 25.5\pr & N12\dg 41.1\pr & 23 & 17 & 75\dg 25.6\pr & N12\dg 42.0\pr \\
23 & 18 & 90\dg 25.7\pr & N12\dg 42.8\pr & 23 & 19 & 105\dg 25.9\pr & N12\dg 43.6\pr & 23 & 20 & 120\dg 26.0\pr & N12\dg 44.5\pr \\
23 & 21 & 135\dg 26.1\pr & N12\dg 45.3\pr & 23 & 22 & 150\dg 26.2\pr & N12\dg 46.1\pr & 23 & 23 & 165\dg 26.3\pr & N12\dg 46.9\pr \\
24 & 0 & 180\dg 26.4\pr & N12\dg 47.8\pr & 24 & 1 & 195\dg 26.5\pr & N12\dg 48.6\pr & 24 & 2 & 210\dg 26.6\pr & N12\dg 49.4\pr \\
24 & 3 & 225\dg 26.8\pr & N12\dg 50.2\pr & 24 & 4 & 240\dg 26.9\pr & N12\dg 51.1\pr & 24 & 5 & 255\dg 27.0\pr & N12\dg 51.9\pr \\
24 & 6 & 270\dg 27.1\pr & N12\dg 52.7\pr & 24 & 7 & 285\dg 27.2\pr & N12\dg 53.5\pr & 24 & 8 & 300\dg 27.3\pr & N12\dg 54.4\pr \\
24 & 9 & 315\dg 27.4\pr & N12\dg 55.2\pr & 24 & 10 & 330\dg 27.5\pr & N12\dg 56.0\pr & 24 & 11 & 345\dg 27.7\pr & N12\dg 56.8\pr \\
24 & 12 & 0\dg 27.8\pr & N12\dg 57.6\pr & 24 & 13 & 15\dg 27.9\pr & N12\dg 58.5\pr & 24 & 14 & 30\dg 28.0\pr & N12\dg 59.3\pr \\
24 & 15 & 45\dg 28.1\pr & N13\dg 00.1\pr & 24 & 16 & 60\dg 28.2\pr & N13\dg 00.9\pr & 24 & 17 & 75\dg 28.3\pr & N13\dg 01.7\pr \\
24 & 18 & 90\dg 28.4\pr & N13\dg 02.6\pr & 24 & 19 & 105\dg 28.5\pr & N13\dg 03.4\pr & 24 & 20 & 120\dg 28.6\pr & N13\dg 04.2\pr \\
24 & 21 & 135\dg 28.8\pr & N13\dg 05.0\pr & 24 & 22 & 150\dg 28.9\pr & N13\dg 05.8\pr & 24 & 23 & 165\dg 29.0\pr & N13\dg 06.6\pr \\
25 & 0 & 180\dg 29.1\pr & N13\dg 07.5\pr & 25 & 1 & 195\dg 29.2\pr & N13\dg 08.3\pr & 25 & 2 & 210\dg 29.3\pr & N13\dg 09.1\pr \\
25 & 3 & 225\dg 29.4\pr & N13\dg 09.9\pr & 25 & 4 & 240\dg 29.5\pr & N13\dg 10.7\pr & 25 & 5 & 255\dg 29.6\pr & N13\dg 11.5\pr \\
25 & 6 & 270\dg 29.7\pr & N13\dg 12.4\pr & 25 & 7 & 285\dg 29.8\pr & N13\dg 13.2\pr & 25 & 8 & 300\dg 29.9\pr & N13\dg 14.0\pr \\
25 & 9 & 315\dg 30.1\pr & N13\dg 14.8\pr & 25 & 10 & 330\dg 30.2\pr & N13\dg 15.6\pr & 25 & 11 & 345\dg 30.3\pr & N13\dg 16.4\pr \\
25 & 12 & 0\dg 30.4\pr & N13\dg 17.2\pr & 25 & 13 & 15\dg 30.5\pr & N13\dg 18.0\pr & 25 & 14 & 30\dg 30.6\pr & N13\dg 18.9\pr \\
25 & 15 & 45\dg 30.7\pr & N13\dg 19.7\pr & 25 & 16 & 60\dg 30.8\pr & N13\dg 20.5\pr & 25 & 17 & 75\dg 30.9\pr & N13\dg 21.3\pr \\
25 & 18 & 90\dg 31.0\pr & N13\dg 22.1\pr & 25 & 19 & 105\dg 31.1\pr & N13\dg 22.9\pr & 25 & 20 & 120\dg 31.2\pr & N13\dg 23.7\pr \\
25 & 21 & 135\dg 31.3\pr & N13\dg 24.5\pr & 25 & 22 & 150\dg 31.4\pr & N13\dg 25.3\pr & 25 & 23 & 165\dg 31.5\pr & N13\dg 26.1\pr \\
26 & 0 & 180\dg 31.6\pr & N13\dg 26.9\pr & 26 & 1 & 195\dg 31.7\pr & N13\dg 27.8\pr & 26 & 2 & 210\dg 31.8\pr & N13\dg 28.6\pr \\
26 & 3 & 225\dg 31.9\pr & N13\dg 29.4\pr & 26 & 4 & 240\dg 32.0\pr & N13\dg 30.2\pr & 26 & 5 & 255\dg 32.1\pr & N13\dg 31.0\pr \\
26 & 6 & 270\dg 32.2\pr & N13\dg 31.8\pr & 26 & 7 & 285\dg 32.3\pr & N13\dg 32.6\pr & 26 & 8 & 300\dg 32.4\pr & N13\dg 33.4\pr \\
26 & 9 & 315\dg 32.5\pr & N13\dg 34.2\pr & 26 & 10 & 330\dg 32.6\pr & N13\dg 35.0\pr & 26 & 11 & 345\dg 32.7\pr & N13\dg 35.8\pr \\
26 & 12 & 0\dg 32.9\pr & N13\dg 36.6\pr & 26 & 13 & 15\dg 33.0\pr & N13\dg 37.4\pr & 26 & 14 & 30\dg 33.1\pr & N13\dg 38.2\pr \\
26 & 15 & 45\dg 33.2\pr & N13\dg 39.0\pr & 26 & 16 & 60\dg 33.3\pr & N13\dg 39.8\pr & 26 & 17 & 75\dg 33.4\pr & N13\dg 40.6\pr \\
26 & 18 & 90\dg 33.5\pr & N13\dg 41.4\pr & 26 & 19 & 105\dg 33.6\pr & N13\dg 42.2\pr & 26 & 20 & 120\dg 33.7\pr & N13\dg 43.0\pr \\
26 & 21 & 135\dg 33.7\pr & N13\dg 43.8\pr & 26 & 22 & 150\dg 33.8\pr & N13\dg 44.6\pr & 26 & 23 & 165\dg 33.9\pr & N13\dg 45.4\pr \\
27 & 0 & 180\dg 34.0\pr & N13\dg 46.2\pr & 27 & 1 & 195\dg 34.1\pr & N13\dg 47.0\pr & 27 & 2 & 210\dg 34.2\pr & N13\dg 47.8\pr \\
27 & 3 & 225\dg 34.3\pr & N13\dg 48.6\pr & 27 & 4 & 240\dg 34.4\pr & N13\dg 49.4\pr & 27 & 5 & 255\dg 34.5\pr & N13\dg 50.2\pr \\
27 & 6 & 270\dg 34.6\pr & N13\dg 51.0\pr & 27 & 7 & 285\dg 34.7\pr & N13\dg 51.8\pr & 27 & 8 & 300\dg 34.8\pr & N13\dg 52.6\pr \\
27 & 9 & 315\dg 34.9\pr & N13\dg 53.4\pr & 27 & 10 & 330\dg 35.0\pr & N13\dg 54.2\pr & 27 & 11 & 345\dg 35.1\pr & N13\dg 55.0\pr \\
27 & 12 & 0\dg 35.2\pr & N13\dg 55.8\pr & 27 & 13 & 15\dg 35.3\pr & N13\dg 56.6\pr & 27 & 14 & 30\dg 35.4\pr & N13\dg 57.3\pr \\
27 & 15 & 45\dg 35.5\pr & N13\dg 58.1\pr & 27 & 16 & 60\dg 35.6\pr & N13\dg 58.9\pr & 27 & 17 & 75\dg 35.7\pr & N13\dg 59.7\pr \\
27 & 18 & 90\dg 35.8\pr & N14\dg 00.5\pr & 27 & 19 & 105\dg 35.9\pr & N14\dg 01.3\pr & 27 & 20 & 120\dg 36.0\pr & N14\dg 02.1\pr \\
27 & 21 & 135\dg 36.1\pr & N14\dg 02.9\pr & 27 & 22 & 150\dg 36.2\pr & N14\dg 03.7\pr & 27 & 23 & 165\dg 36.2\pr & N14\dg 04.5\pr \\
28 & 0 & 180\dg 36.3\pr & N14\dg 05.3\pr & 28 & 1 & 195\dg 36.4\pr & N14\dg 06.0\pr & 28 & 2 & 210\dg 36.5\pr & N14\dg 06.8\pr \\
28 & 3 & 225\dg 36.6\pr & N14\dg 07.6\pr & 28 & 4 & 240\dg 36.7\pr & N14\dg 08.4\pr & 28 & 5 & 255\dg 36.8\pr & N14\dg 09.2\pr \\
28 & 6 & 270\dg 36.9\pr & N14\dg 10.0\pr & 28 & 7 & 285\dg 37.0\pr & N14\dg 10.8\pr & 28 & 8 & 300\dg 37.1\pr & N14\dg 11.5\pr \\
28 & 9 & 315\dg 37.2\pr & N14\dg 12.3\pr & 28 & 10 & 330\dg 37.3\pr & N14\dg 13.1\pr & 28 & 11 & 345\dg 37.3\pr & N14\dg 13.9\pr \\
28 & 12 & 0\dg 37.4\pr & N14\dg 14.7\pr & 28 & 13 & 15\dg 37.5\pr & N14\dg 15.5\pr & 28 & 14 & 30\dg 37.6\pr & N14\dg 16.2\pr \\
28 & 15 & 45\dg 37.7\pr & N14\dg 17.0\pr & 28 & 16 & 60\dg 37.8\pr & N14\dg 17.8\pr & 28 & 17 & 75\dg 37.9\pr & N14\dg 18.6\pr \\
28 & 18 & 90\dg 38.0\pr & N14\dg 19.4\pr & 28 & 19 & 105\dg 38.1\pr & N14\dg 20.2\pr & 28 & 20 & 120\dg 38.2\pr & N14\dg 20.9\pr \\
28 & 21 & 135\dg 38.2\pr & N14\dg 21.7\pr & 28 & 22 & 150\dg 38.3\pr & N14\dg 22.5\pr & 28 & 23 & 165\dg 38.4\pr & N14\dg 23.3\pr \\
29 & 0 & 180\dg 38.5\pr & N14\dg 24.1\pr & 29 & 1 & 195\dg 38.6\pr & N14\dg 24.8\pr & 29 & 2 & 210\dg 38.7\pr & N14\dg 25.6\pr \\
29 & 3 & 225\dg 38.8\pr & N14\dg 26.4\pr & 29 & 4 & 240\dg 38.9\pr & N14\dg 27.2\pr & 29 & 5 & 255\dg 38.9\pr & N14\dg 27.9\pr \\
29 & 6 & 270\dg 39.0\pr & N14\dg 28.7\pr & 29 & 7 & 285\dg 39.1\pr & N14\dg 29.5\pr & 29 & 8 & 300\dg 39.2\pr & N14\dg 30.3\pr \\
29 & 9 & 315\dg 39.3\pr & N14\dg 31.0\pr & 29 & 10 & 330\dg 39.4\pr & N14\dg 31.8\pr & 29 & 11 & 345\dg 39.5\pr & N14\dg 32.6\pr \\
29 & 12 & 0\dg 39.5\pr & N14\dg 33.4\pr & 29 & 13 & 15\dg 39.6\pr & N14\dg 34.1\pr & 29 & 14 & 30\dg 39.7\pr & N14\dg 34.9\pr \\
29 & 15 & 45\dg 39.8\pr & N14\dg 35.7\pr & 29 & 16 & 60\dg 39.9\pr & N14\dg 36.5\pr & 29 & 17 & 75\dg 40.0\pr & N14\dg 37.2\pr \\
29 & 18 & 90\dg 40.0\pr & N14\dg 38.0\pr & 29 & 19 & 105\dg 40.1\pr & N14\dg 38.8\pr & 29 & 20 & 120\dg 40.2\pr & N14\dg 39.5\pr \\
29 & 21 & 135\dg 40.3\pr & N14\dg 40.3\pr & 29 & 22 & 150\dg 40.4\pr & N14\dg 41.1\pr & 29 & 23 & 165\dg 40.5\pr & N14\dg 41.9\pr \\
30 & 0 & 180\dg 40.5\pr & N14\dg 42.6\pr & 30 & 1 & 195\dg 40.6\pr & N14\dg 43.4\pr & 30 & 2 & 210\dg 40.7\pr & N14\dg 44.2\pr \\
30 & 3 & 225\dg 40.8\pr & N14\dg 44.9\pr & 30 & 4 & 240\dg 40.9\pr & N14\dg 45.7\pr & 30 & 5 & 255\dg 40.9\pr & N14\dg 46.5\pr \\
30 & 6 & 270\dg 41.0\pr & N14\dg 47.2\pr & 30 & 7 & 285\dg 41.1\pr & N14\dg 48.0\pr & 30 & 8 & 300\dg 41.2\pr & N14\dg 48.8\pr \\
30 & 9 & 315\dg 41.3\pr & N14\dg 49.5\pr & 30 & 10 & 330\dg 41.4\pr & N14\dg 50.3\pr & 30 & 11 & 345\dg 41.4\pr & N14\dg 51.1\pr \\
30 & 12 & 0\dg 41.5\pr & N14\dg 51.8\pr & 30 & 13 & 15\dg 41.6\pr & N14\dg 52.6\pr & 30 & 14 & 30\dg 41.7\pr & N14\dg 53.3\pr \\
30 & 15 & 45\dg 41.7\pr & N14\dg 54.1\pr & 30 & 16 & 60\dg 41.8\pr & N14\dg 54.9\pr & 30 & 17 & 75\dg 41.9\pr & N14\dg 55.6\pr \\
30 & 18 & 90\dg 42.0\pr & N14\dg 56.4\pr & 30 & 19 & 105\dg 42.1\pr & N14\dg 57.2\pr & 30 & 20 & 120\dg 42.1\pr & N14\dg 57.9\pr \\
30 & 21 & 135\dg 42.2\pr & N14\dg 58.7\pr & 30 & 22 & 150\dg 42.3\pr & N14\dg 59.4\pr & 30 & 23 & 165\dg 42.4\pr & N15\dg 00.2\pr \\
\end{longtable}}

{\footnotesize
\begin{longtable}{rrrr|rrrr|rrrr}
\caption{Sun --- GHA and Declination, May 2026 (UT)}\label{sun05}\\
\toprule
\textbf{d} & \textbf{h} & \textbf{GHA} & \textbf{Dec} & \textbf{d} & \textbf{h} & \textbf{GHA} & \textbf{Dec} & \textbf{d} & \textbf{h} & \textbf{GHA} & \textbf{Dec} \\
\midrule
\endfirsthead
\multicolumn{12}{c}{\textit{Sun --- GHA and Declination, May 2026 (UT) (continued)}}\\
\toprule
\textbf{d} & \textbf{h} & \textbf{GHA} & \textbf{Dec} & \textbf{d} & \textbf{h} & \textbf{GHA} & \textbf{Dec} & \textbf{d} & \textbf{h} & \textbf{GHA} & \textbf{Dec} \\
\midrule
\endhead
\midrule\multicolumn{12}{r}{\textit{continued on next page}}\\
\endfoot
\bottomrule
\endlastfoot
1 & 0 & 180\dg 42.4\pr & N15\dg 00.9\pr & 1 & 1 & 195\dg 42.5\pr & N15\dg 01.7\pr & 1 & 2 & 210\dg 42.6\pr & N15\dg 02.5\pr \\
1 & 3 & 225\dg 42.7\pr & N15\dg 03.2\pr & 1 & 4 & 240\dg 42.7\pr & N15\dg 04.0\pr & 1 & 5 & 255\dg 42.8\pr & N15\dg 04.7\pr \\
1 & 6 & 270\dg 42.9\pr & N15\dg 05.5\pr & 1 & 7 & 285\dg 43.0\pr & N15\dg 06.2\pr & 1 & 8 & 300\dg 43.1\pr & N15\dg 07.0\pr \\
1 & 9 & 315\dg 43.1\pr & N15\dg 07.8\pr & 1 & 10 & 330\dg 43.2\pr & N15\dg 08.5\pr & 1 & 11 & 345\dg 43.3\pr & N15\dg 09.3\pr \\
1 & 12 & 0\dg 43.3\pr & N15\dg 10.0\pr & 1 & 13 & 15\dg 43.4\pr & N15\dg 10.8\pr & 1 & 14 & 30\dg 43.5\pr & N15\dg 11.5\pr \\
1 & 15 & 45\dg 43.6\pr & N15\dg 12.3\pr & 1 & 16 & 60\dg 43.6\pr & N15\dg 13.0\pr & 1 & 17 & 75\dg 43.7\pr & N15\dg 13.8\pr \\
1 & 18 & 90\dg 43.8\pr & N15\dg 14.5\pr & 1 & 19 & 105\dg 43.9\pr & N15\dg 15.3\pr & 1 & 20 & 120\dg 43.9\pr & N15\dg 16.0\pr \\
1 & 21 & 135\dg 44.0\pr & N15\dg 16.8\pr & 1 & 22 & 150\dg 44.1\pr & N15\dg 17.5\pr & 1 & 23 & 165\dg 44.1\pr & N15\dg 18.3\pr \\
2 & 0 & 180\dg 44.2\pr & N15\dg 19.0\pr & 2 & 1 & 195\dg 44.3\pr & N15\dg 19.8\pr & 2 & 2 & 210\dg 44.4\pr & N15\dg 20.5\pr \\
2 & 3 & 225\dg 44.4\pr & N15\dg 21.3\pr & 2 & 4 & 240\dg 44.5\pr & N15\dg 22.0\pr & 2 & 5 & 255\dg 44.6\pr & N15\dg 22.8\pr \\
2 & 6 & 270\dg 44.6\pr & N15\dg 23.5\pr & 2 & 7 & 285\dg 44.7\pr & N15\dg 24.2\pr & 2 & 8 & 300\dg 44.8\pr & N15\dg 25.0\pr \\
2 & 9 & 315\dg 44.8\pr & N15\dg 25.7\pr & 2 & 10 & 330\dg 44.9\pr & N15\dg 26.5\pr & 2 & 11 & 345\dg 45.0\pr & N15\dg 27.2\pr \\
2 & 12 & 0\dg 45.1\pr & N15\dg 28.0\pr & 2 & 13 & 15\dg 45.1\pr & N15\dg 28.7\pr & 2 & 14 & 30\dg 45.2\pr & N15\dg 29.5\pr \\
2 & 15 & 45\dg 45.3\pr & N15\dg 30.2\pr & 2 & 16 & 60\dg 45.3\pr & N15\dg 30.9\pr & 2 & 17 & 75\dg 45.4\pr & N15\dg 31.7\pr \\
2 & 18 & 90\dg 45.5\pr & N15\dg 32.4\pr & 2 & 19 & 105\dg 45.5\pr & N15\dg 33.2\pr & 2 & 20 & 120\dg 45.6\pr & N15\dg 33.9\pr \\
2 & 21 & 135\dg 45.7\pr & N15\dg 34.6\pr & 2 & 22 & 150\dg 45.7\pr & N15\dg 35.4\pr & 2 & 23 & 165\dg 45.8\pr & N15\dg 36.1\pr \\
3 & 0 & 180\dg 45.9\pr & N15\dg 36.8\pr & 3 & 1 & 195\dg 45.9\pr & N15\dg 37.6\pr & 3 & 2 & 210\dg 46.0\pr & N15\dg 38.3\pr \\
3 & 3 & 225\dg 46.0\pr & N15\dg 39.1\pr & 3 & 4 & 240\dg 46.1\pr & N15\dg 39.8\pr & 3 & 5 & 255\dg 46.2\pr & N15\dg 40.5\pr \\
3 & 6 & 270\dg 46.2\pr & N15\dg 41.3\pr & 3 & 7 & 285\dg 46.3\pr & N15\dg 42.0\pr & 3 & 8 & 300\dg 46.4\pr & N15\dg 42.7\pr \\
3 & 9 & 315\dg 46.4\pr & N15\dg 43.5\pr & 3 & 10 & 330\dg 46.5\pr & N15\dg 44.2\pr & 3 & 11 & 345\dg 46.6\pr & N15\dg 44.9\pr \\
3 & 12 & 0\dg 46.6\pr & N15\dg 45.7\pr & 3 & 13 & 15\dg 46.7\pr & N15\dg 46.4\pr & 3 & 14 & 30\dg 46.7\pr & N15\dg 47.1\pr \\
3 & 15 & 45\dg 46.8\pr & N15\dg 47.9\pr & 3 & 16 & 60\dg 46.9\pr & N15\dg 48.6\pr & 3 & 17 & 75\dg 46.9\pr & N15\dg 49.3\pr \\
3 & 18 & 90\dg 47.0\pr & N15\dg 50.0\pr & 3 & 19 & 105\dg 47.0\pr & N15\dg 50.8\pr & 3 & 20 & 120\dg 47.1\pr & N15\dg 51.5\pr \\
3 & 21 & 135\dg 47.2\pr & N15\dg 52.2\pr & 3 & 22 & 150\dg 47.2\pr & N15\dg 53.0\pr & 3 & 23 & 165\dg 47.3\pr & N15\dg 53.7\pr \\
4 & 0 & 180\dg 47.4\pr & N15\dg 54.4\pr & 4 & 1 & 195\dg 47.4\pr & N15\dg 55.1\pr & 4 & 2 & 210\dg 47.5\pr & N15\dg 55.9\pr \\
4 & 3 & 225\dg 47.5\pr & N15\dg 56.6\pr & 4 & 4 & 240\dg 47.6\pr & N15\dg 57.3\pr & 4 & 5 & 255\dg 47.6\pr & N15\dg 58.0\pr \\
4 & 6 & 270\dg 47.7\pr & N15\dg 58.8\pr & 4 & 7 & 285\dg 47.8\pr & N15\dg 59.5\pr & 4 & 8 & 300\dg 47.8\pr & N16\dg 00.2\pr \\
4 & 9 & 315\dg 47.9\pr & N16\dg 00.9\pr & 4 & 10 & 330\dg 47.9\pr & N16\dg 01.7\pr & 4 & 11 & 345\dg 48.0\pr & N16\dg 02.4\pr \\
4 & 12 & 0\dg 48.0\pr & N16\dg 03.1\pr & 4 & 13 & 15\dg 48.1\pr & N16\dg 03.8\pr & 4 & 14 & 30\dg 48.2\pr & N16\dg 04.5\pr \\
4 & 15 & 45\dg 48.2\pr & N16\dg 05.3\pr & 4 & 16 & 60\dg 48.3\pr & N16\dg 06.0\pr & 4 & 17 & 75\dg 48.3\pr & N16\dg 06.7\pr \\
4 & 18 & 90\dg 48.4\pr & N16\dg 07.4\pr & 4 & 19 & 105\dg 48.4\pr & N16\dg 08.1\pr & 4 & 20 & 120\dg 48.5\pr & N16\dg 08.8\pr \\
4 & 21 & 135\dg 48.5\pr & N16\dg 09.6\pr & 4 & 22 & 150\dg 48.6\pr & N16\dg 10.3\pr & 4 & 23 & 165\dg 48.7\pr & N16\dg 11.0\pr \\
5 & 0 & 180\dg 48.7\pr & N16\dg 11.7\pr & 5 & 1 & 195\dg 48.8\pr & N16\dg 12.4\pr & 5 & 2 & 210\dg 48.8\pr & N16\dg 13.1\pr \\
5 & 3 & 225\dg 48.9\pr & N16\dg 13.9\pr & 5 & 4 & 240\dg 48.9\pr & N16\dg 14.6\pr & 5 & 5 & 255\dg 49.0\pr & N16\dg 15.3\pr \\
5 & 6 & 270\dg 49.0\pr & N16\dg 16.0\pr & 5 & 7 & 285\dg 49.1\pr & N16\dg 16.7\pr & 5 & 8 & 300\dg 49.1\pr & N16\dg 17.4\pr \\
5 & 9 & 315\dg 49.2\pr & N16\dg 18.1\pr & 5 & 10 & 330\dg 49.2\pr & N16\dg 18.8\pr & 5 & 11 & 345\dg 49.3\pr & N16\dg 19.5\pr \\
5 & 12 & 0\dg 49.3\pr & N16\dg 20.3\pr & 5 & 13 & 15\dg 49.4\pr & N16\dg 21.0\pr & 5 & 14 & 30\dg 49.4\pr & N16\dg 21.7\pr \\
5 & 15 & 45\dg 49.5\pr & N16\dg 22.4\pr & 5 & 16 & 60\dg 49.5\pr & N16\dg 23.1\pr & 5 & 17 & 75\dg 49.6\pr & N16\dg 23.8\pr \\
5 & 18 & 90\dg 49.6\pr & N16\dg 24.5\pr & 5 & 19 & 105\dg 49.7\pr & N16\dg 25.2\pr & 5 & 20 & 120\dg 49.7\pr & N16\dg 25.9\pr \\
5 & 21 & 135\dg 49.8\pr & N16\dg 26.6\pr & 5 & 22 & 150\dg 49.8\pr & N16\dg 27.3\pr & 5 & 23 & 165\dg 49.9\pr & N16\dg 28.0\pr \\
6 & 0 & 180\dg 49.9\pr & N16\dg 28.7\pr & 6 & 1 & 195\dg 50.0\pr & N16\dg 29.4\pr & 6 & 2 & 210\dg 50.0\pr & N16\dg 30.1\pr \\
6 & 3 & 225\dg 50.1\pr & N16\dg 30.9\pr & 6 & 4 & 240\dg 50.1\pr & N16\dg 31.6\pr & 6 & 5 & 255\dg 50.2\pr & N16\dg 32.3\pr \\
6 & 6 & 270\dg 50.2\pr & N16\dg 33.0\pr & 6 & 7 & 285\dg 50.3\pr & N16\dg 33.7\pr & 6 & 8 & 300\dg 50.3\pr & N16\dg 34.4\pr \\
6 & 9 & 315\dg 50.4\pr & N16\dg 35.1\pr & 6 & 10 & 330\dg 50.4\pr & N16\dg 35.8\pr & 6 & 11 & 345\dg 50.4\pr & N16\dg 36.5\pr \\
6 & 12 & 0\dg 50.5\pr & N16\dg 37.2\pr & 6 & 13 & 15\dg 50.5\pr & N16\dg 37.9\pr & 6 & 14 & 30\dg 50.6\pr & N16\dg 38.5\pr \\
6 & 15 & 45\dg 50.6\pr & N16\dg 39.2\pr & 6 & 16 & 60\dg 50.7\pr & N16\dg 39.9\pr & 6 & 17 & 75\dg 50.7\pr & N16\dg 40.6\pr \\
6 & 18 & 90\dg 50.8\pr & N16\dg 41.3\pr & 6 & 19 & 105\dg 50.8\pr & N16\dg 42.0\pr & 6 & 20 & 120\dg 50.8\pr & N16\dg 42.7\pr \\
6 & 21 & 135\dg 50.9\pr & N16\dg 43.4\pr & 6 & 22 & 150\dg 50.9\pr & N16\dg 44.1\pr & 6 & 23 & 165\dg 51.0\pr & N16\dg 44.8\pr \\
7 & 0 & 180\dg 51.0\pr & N16\dg 45.5\pr & 7 & 1 & 195\dg 51.1\pr & N16\dg 46.2\pr & 7 & 2 & 210\dg 51.1\pr & N16\dg 46.9\pr \\
7 & 3 & 225\dg 51.1\pr & N16\dg 47.6\pr & 7 & 4 & 240\dg 51.2\pr & N16\dg 48.3\pr & 7 & 5 & 255\dg 51.2\pr & N16\dg 49.0\pr \\
7 & 6 & 270\dg 51.3\pr & N16\dg 49.6\pr & 7 & 7 & 285\dg 51.3\pr & N16\dg 50.3\pr & 7 & 8 & 300\dg 51.3\pr & N16\dg 51.0\pr \\
7 & 9 & 315\dg 51.4\pr & N16\dg 51.7\pr & 7 & 10 & 330\dg 51.4\pr & N16\dg 52.4\pr & 7 & 11 & 345\dg 51.5\pr & N16\dg 53.1\pr \\
7 & 12 & 0\dg 51.5\pr & N16\dg 53.8\pr & 7 & 13 & 15\dg 51.5\pr & N16\dg 54.5\pr & 7 & 14 & 30\dg 51.6\pr & N16\dg 55.1\pr \\
7 & 15 & 45\dg 51.6\pr & N16\dg 55.8\pr & 7 & 16 & 60\dg 51.7\pr & N16\dg 56.5\pr & 7 & 17 & 75\dg 51.7\pr & N16\dg 57.2\pr \\
7 & 18 & 90\dg 51.7\pr & N16\dg 57.9\pr & 7 & 19 & 105\dg 51.8\pr & N16\dg 58.6\pr & 7 & 20 & 120\dg 51.8\pr & N16\dg 59.2\pr \\
7 & 21 & 135\dg 51.8\pr & N16\dg 59.9\pr & 7 & 22 & 150\dg 51.9\pr & N17\dg 00.6\pr & 7 & 23 & 165\dg 51.9\pr & N17\dg 01.3\pr \\
8 & 0 & 180\dg 51.9\pr & N17\dg 02.0\pr & 8 & 1 & 195\dg 52.0\pr & N17\dg 02.7\pr & 8 & 2 & 210\dg 52.0\pr & N17\dg 03.3\pr \\
8 & 3 & 225\dg 52.1\pr & N17\dg 04.0\pr & 8 & 4 & 240\dg 52.1\pr & N17\dg 04.7\pr & 8 & 5 & 255\dg 52.1\pr & N17\dg 05.4\pr \\
8 & 6 & 270\dg 52.2\pr & N17\dg 06.0\pr & 8 & 7 & 285\dg 52.2\pr & N17\dg 06.7\pr & 8 & 8 & 300\dg 52.2\pr & N17\dg 07.4\pr \\
8 & 9 & 315\dg 52.3\pr & N17\dg 08.1\pr & 8 & 10 & 330\dg 52.3\pr & N17\dg 08.8\pr & 8 & 11 & 345\dg 52.3\pr & N17\dg 09.4\pr \\
8 & 12 & 0\dg 52.4\pr & N17\dg 10.1\pr & 8 & 13 & 15\dg 52.4\pr & N17\dg 10.8\pr & 8 & 14 & 30\dg 52.4\pr & N17\dg 11.5\pr \\
8 & 15 & 45\dg 52.5\pr & N17\dg 12.1\pr & 8 & 16 & 60\dg 52.5\pr & N17\dg 12.8\pr & 8 & 17 & 75\dg 52.5\pr & N17\dg 13.5\pr \\
8 & 18 & 90\dg 52.6\pr & N17\dg 14.1\pr & 8 & 19 & 105\dg 52.6\pr & N17\dg 14.8\pr & 8 & 20 & 120\dg 52.6\pr & N17\dg 15.5\pr \\
8 & 21 & 135\dg 52.7\pr & N17\dg 16.2\pr & 8 & 22 & 150\dg 52.7\pr & N17\dg 16.8\pr & 8 & 23 & 165\dg 52.7\pr & N17\dg 17.5\pr \\
9 & 0 & 180\dg 52.7\pr & N17\dg 18.2\pr & 9 & 1 & 195\dg 52.8\pr & N17\dg 18.8\pr & 9 & 2 & 210\dg 52.8\pr & N17\dg 19.5\pr \\
9 & 3 & 225\dg 52.8\pr & N17\dg 20.2\pr & 9 & 4 & 240\dg 52.9\pr & N17\dg 20.8\pr & 9 & 5 & 255\dg 52.9\pr & N17\dg 21.5\pr \\
9 & 6 & 270\dg 52.9\pr & N17\dg 22.2\pr & 9 & 7 & 285\dg 53.0\pr & N17\dg 22.8\pr & 9 & 8 & 300\dg 53.0\pr & N17\dg 23.5\pr \\
9 & 9 & 315\dg 53.0\pr & N17\dg 24.2\pr & 9 & 10 & 330\dg 53.0\pr & N17\dg 24.8\pr & 9 & 11 & 345\dg 53.1\pr & N17\dg 25.5\pr \\
9 & 12 & 0\dg 53.1\pr & N17\dg 26.2\pr & 9 & 13 & 15\dg 53.1\pr & N17\dg 26.8\pr & 9 & 14 & 30\dg 53.1\pr & N17\dg 27.5\pr \\
9 & 15 & 45\dg 53.2\pr & N17\dg 28.1\pr & 9 & 16 & 60\dg 53.2\pr & N17\dg 28.8\pr & 9 & 17 & 75\dg 53.2\pr & N17\dg 29.5\pr \\
9 & 18 & 90\dg 53.2\pr & N17\dg 30.1\pr & 9 & 19 & 105\dg 53.3\pr & N17\dg 30.8\pr & 9 & 20 & 120\dg 53.3\pr & N17\dg 31.4\pr \\
9 & 21 & 135\dg 53.3\pr & N17\dg 32.1\pr & 9 & 22 & 150\dg 53.4\pr & N17\dg 32.8\pr & 9 & 23 & 165\dg 53.4\pr & N17\dg 33.4\pr \\
10 & 0 & 180\dg 53.4\pr & N17\dg 34.1\pr & 10 & 1 & 195\dg 53.4\pr & N17\dg 34.7\pr & 10 & 2 & 210\dg 53.4\pr & N17\dg 35.4\pr \\
10 & 3 & 225\dg 53.5\pr & N17\dg 36.0\pr & 10 & 4 & 240\dg 53.5\pr & N17\dg 36.7\pr & 10 & 5 & 255\dg 53.5\pr & N17\dg 37.3\pr \\
10 & 6 & 270\dg 53.5\pr & N17\dg 38.0\pr & 10 & 7 & 285\dg 53.6\pr & N17\dg 38.7\pr & 10 & 8 & 300\dg 53.6\pr & N17\dg 39.3\pr \\
10 & 9 & 315\dg 53.6\pr & N17\dg 40.0\pr & 10 & 10 & 330\dg 53.6\pr & N17\dg 40.6\pr & 10 & 11 & 345\dg 53.7\pr & N17\dg 41.3\pr \\
10 & 12 & 0\dg 53.7\pr & N17\dg 41.9\pr & 10 & 13 & 15\dg 53.7\pr & N17\dg 42.6\pr & 10 & 14 & 30\dg 53.7\pr & N17\dg 43.2\pr \\
10 & 15 & 45\dg 53.7\pr & N17\dg 43.9\pr & 10 & 16 & 60\dg 53.8\pr & N17\dg 44.5\pr & 10 & 17 & 75\dg 53.8\pr & N17\dg 45.2\pr \\
10 & 18 & 90\dg 53.8\pr & N17\dg 45.8\pr & 10 & 19 & 105\dg 53.8\pr & N17\dg 46.5\pr & 10 & 20 & 120\dg 53.8\pr & N17\dg 47.1\pr \\
10 & 21 & 135\dg 53.9\pr & N17\dg 47.7\pr & 10 & 22 & 150\dg 53.9\pr & N17\dg 48.4\pr & 10 & 23 & 165\dg 53.9\pr & N17\dg 49.0\pr \\
11 & 0 & 180\dg 53.9\pr & N17\dg 49.7\pr & 11 & 1 & 195\dg 53.9\pr & N17\dg 50.3\pr & 11 & 2 & 210\dg 53.9\pr & N17\dg 51.0\pr \\
11 & 3 & 225\dg 54.0\pr & N17\dg 51.6\pr & 11 & 4 & 240\dg 54.0\pr & N17\dg 52.3\pr & 11 & 5 & 255\dg 54.0\pr & N17\dg 52.9\pr \\
11 & 6 & 270\dg 54.0\pr & N17\dg 53.5\pr & 11 & 7 & 285\dg 54.0\pr & N17\dg 54.2\pr & 11 & 8 & 300\dg 54.0\pr & N17\dg 54.8\pr \\
11 & 9 & 315\dg 54.1\pr & N17\dg 55.5\pr & 11 & 10 & 330\dg 54.1\pr & N17\dg 56.1\pr & 11 & 11 & 345\dg 54.1\pr & N17\dg 56.7\pr \\
11 & 12 & 0\dg 54.1\pr & N17\dg 57.4\pr & 11 & 13 & 15\dg 54.1\pr & N17\dg 58.0\pr & 11 & 14 & 30\dg 54.1\pr & N17\dg 58.6\pr \\
11 & 15 & 45\dg 54.2\pr & N17\dg 59.3\pr & 11 & 16 & 60\dg 54.2\pr & N17\dg 59.9\pr & 11 & 17 & 75\dg 54.2\pr & N18\dg 00.6\pr \\
11 & 18 & 90\dg 54.2\pr & N18\dg 01.2\pr & 11 & 19 & 105\dg 54.2\pr & N18\dg 01.8\pr & 11 & 20 & 120\dg 54.2\pr & N18\dg 02.5\pr \\
11 & 21 & 135\dg 54.2\pr & N18\dg 03.1\pr & 11 & 22 & 150\dg 54.3\pr & N18\dg 03.7\pr & 11 & 23 & 165\dg 54.3\pr & N18\dg 04.4\pr \\
12 & 0 & 180\dg 54.3\pr & N18\dg 05.0\pr & 12 & 1 & 195\dg 54.3\pr & N18\dg 05.6\pr & 12 & 2 & 210\dg 54.3\pr & N18\dg 06.3\pr \\
12 & 3 & 225\dg 54.3\pr & N18\dg 06.9\pr & 12 & 4 & 240\dg 54.3\pr & N18\dg 07.5\pr & 12 & 5 & 255\dg 54.3\pr & N18\dg 08.1\pr \\
12 & 6 & 270\dg 54.3\pr & N18\dg 08.8\pr & 12 & 7 & 285\dg 54.4\pr & N18\dg 09.4\pr & 12 & 8 & 300\dg 54.4\pr & N18\dg 10.0\pr \\
12 & 9 & 315\dg 54.4\pr & N18\dg 10.7\pr & 12 & 10 & 330\dg 54.4\pr & N18\dg 11.3\pr & 12 & 11 & 345\dg 54.4\pr & N18\dg 11.9\pr \\
12 & 12 & 0\dg 54.4\pr & N18\dg 12.5\pr & 12 & 13 & 15\dg 54.4\pr & N18\dg 13.2\pr & 12 & 14 & 30\dg 54.4\pr & N18\dg 13.8\pr \\
12 & 15 & 45\dg 54.4\pr & N18\dg 14.4\pr & 12 & 16 & 60\dg 54.4\pr & N18\dg 15.0\pr & 12 & 17 & 75\dg 54.5\pr & N18\dg 15.7\pr \\
12 & 18 & 90\dg 54.5\pr & N18\dg 16.3\pr & 12 & 19 & 105\dg 54.5\pr & N18\dg 16.9\pr & 12 & 20 & 120\dg 54.5\pr & N18\dg 17.5\pr \\
12 & 21 & 135\dg 54.5\pr & N18\dg 18.1\pr & 12 & 22 & 150\dg 54.5\pr & N18\dg 18.8\pr & 12 & 23 & 165\dg 54.5\pr & N18\dg 19.4\pr \\
13 & 0 & 180\dg 54.5\pr & N18\dg 20.0\pr & 13 & 1 & 195\dg 54.5\pr & N18\dg 20.6\pr & 13 & 2 & 210\dg 54.5\pr & N18\dg 21.2\pr \\
13 & 3 & 225\dg 54.5\pr & N18\dg 21.9\pr & 13 & 4 & 240\dg 54.5\pr & N18\dg 22.5\pr & 13 & 5 & 255\dg 54.5\pr & N18\dg 23.1\pr \\
13 & 6 & 270\dg 54.5\pr & N18\dg 23.7\pr & 13 & 7 & 285\dg 54.5\pr & N18\dg 24.3\pr & 13 & 8 & 300\dg 54.5\pr & N18\dg 24.9\pr \\
13 & 9 & 315\dg 54.6\pr & N18\dg 25.5\pr & 13 & 10 & 330\dg 54.6\pr & N18\dg 26.2\pr & 13 & 11 & 345\dg 54.6\pr & N18\dg 26.8\pr \\
13 & 12 & 0\dg 54.6\pr & N18\dg 27.4\pr & 13 & 13 & 15\dg 54.6\pr & N18\dg 28.0\pr & 13 & 14 & 30\dg 54.6\pr & N18\dg 28.6\pr \\
13 & 15 & 45\dg 54.6\pr & N18\dg 29.2\pr & 13 & 16 & 60\dg 54.6\pr & N18\dg 29.8\pr & 13 & 17 & 75\dg 54.6\pr & N18\dg 30.4\pr \\
13 & 18 & 90\dg 54.6\pr & N18\dg 31.0\pr & 13 & 19 & 105\dg 54.6\pr & N18\dg 31.7\pr & 13 & 20 & 120\dg 54.6\pr & N18\dg 32.3\pr \\
13 & 21 & 135\dg 54.6\pr & N18\dg 32.9\pr & 13 & 22 & 150\dg 54.6\pr & N18\dg 33.5\pr & 13 & 23 & 165\dg 54.6\pr & N18\dg 34.1\pr \\
14 & 0 & 180\dg 54.6\pr & N18\dg 34.7\pr & 14 & 1 & 195\dg 54.6\pr & N18\dg 35.3\pr & 14 & 2 & 210\dg 54.6\pr & N18\dg 35.9\pr \\
14 & 3 & 225\dg 54.6\pr & N18\dg 36.5\pr & 14 & 4 & 240\dg 54.6\pr & N18\dg 37.1\pr & 14 & 5 & 255\dg 54.6\pr & N18\dg 37.7\pr \\
14 & 6 & 270\dg 54.6\pr & N18\dg 38.3\pr & 14 & 7 & 285\dg 54.6\pr & N18\dg 38.9\pr & 14 & 8 & 300\dg 54.6\pr & N18\dg 39.5\pr \\
14 & 9 & 315\dg 54.6\pr & N18\dg 40.1\pr & 14 & 10 & 330\dg 54.6\pr & N18\dg 40.7\pr & 14 & 11 & 345\dg 54.6\pr & N18\dg 41.3\pr \\
14 & 12 & 0\dg 54.6\pr & N18\dg 41.9\pr & 14 & 13 & 15\dg 54.6\pr & N18\dg 42.5\pr & 14 & 14 & 30\dg 54.6\pr & N18\dg 43.1\pr \\
14 & 15 & 45\dg 54.6\pr & N18\dg 43.7\pr & 14 & 16 & 60\dg 54.6\pr & N18\dg 44.3\pr & 14 & 17 & 75\dg 54.6\pr & N18\dg 44.9\pr \\
14 & 18 & 90\dg 54.6\pr & N18\dg 45.5\pr & 14 & 19 & 105\dg 54.5\pr & N18\dg 46.1\pr & 14 & 20 & 120\dg 54.5\pr & N18\dg 46.7\pr \\
14 & 21 & 135\dg 54.5\pr & N18\dg 47.3\pr & 14 & 22 & 150\dg 54.5\pr & N18\dg 47.9\pr & 14 & 23 & 165\dg 54.5\pr & N18\dg 48.5\pr \\
15 & 0 & 180\dg 54.5\pr & N18\dg 49.1\pr & 15 & 1 & 195\dg 54.5\pr & N18\dg 49.7\pr & 15 & 2 & 210\dg 54.5\pr & N18\dg 50.3\pr \\
15 & 3 & 225\dg 54.5\pr & N18\dg 50.9\pr & 15 & 4 & 240\dg 54.5\pr & N18\dg 51.4\pr & 15 & 5 & 255\dg 54.5\pr & N18\dg 52.0\pr \\
15 & 6 & 270\dg 54.5\pr & N18\dg 52.6\pr & 15 & 7 & 285\dg 54.5\pr & N18\dg 53.2\pr & 15 & 8 & 300\dg 54.5\pr & N18\dg 53.8\pr \\
15 & 9 & 315\dg 54.5\pr & N18\dg 54.4\pr & 15 & 10 & 330\dg 54.5\pr & N18\dg 55.0\pr & 15 & 11 & 345\dg 54.5\pr & N18\dg 55.6\pr \\
15 & 12 & 0\dg 54.4\pr & N18\dg 56.2\pr & 15 & 13 & 15\dg 54.4\pr & N18\dg 56.7\pr & 15 & 14 & 30\dg 54.4\pr & N18\dg 57.3\pr \\
15 & 15 & 45\dg 54.4\pr & N18\dg 57.9\pr & 15 & 16 & 60\dg 54.4\pr & N18\dg 58.5\pr & 15 & 17 & 75\dg 54.4\pr & N18\dg 59.1\pr \\
15 & 18 & 90\dg 54.4\pr & N18\dg 59.7\pr & 15 & 19 & 105\dg 54.4\pr & N19\dg 00.2\pr & 15 & 20 & 120\dg 54.4\pr & N19\dg 00.8\pr \\
15 & 21 & 135\dg 54.4\pr & N19\dg 01.4\pr & 15 & 22 & 150\dg 54.3\pr & N19\dg 02.0\pr & 15 & 23 & 165\dg 54.3\pr & N19\dg 02.6\pr \\
16 & 0 & 180\dg 54.3\pr & N19\dg 03.1\pr & 16 & 1 & 195\dg 54.3\pr & N19\dg 03.7\pr & 16 & 2 & 210\dg 54.3\pr & N19\dg 04.3\pr \\
16 & 3 & 225\dg 54.3\pr & N19\dg 04.9\pr & 16 & 4 & 240\dg 54.3\pr & N19\dg 05.5\pr & 16 & 5 & 255\dg 54.3\pr & N19\dg 06.0\pr \\
16 & 6 & 270\dg 54.3\pr & N19\dg 06.6\pr & 16 & 7 & 285\dg 54.2\pr & N19\dg 07.2\pr & 16 & 8 & 300\dg 54.2\pr & N19\dg 07.8\pr \\
16 & 9 & 315\dg 54.2\pr & N19\dg 08.3\pr & 16 & 10 & 330\dg 54.2\pr & N19\dg 08.9\pr & 16 & 11 & 345\dg 54.2\pr & N19\dg 09.5\pr \\
16 & 12 & 0\dg 54.2\pr & N19\dg 10.1\pr & 16 & 13 & 15\dg 54.2\pr & N19\dg 10.6\pr & 16 & 14 & 30\dg 54.1\pr & N19\dg 11.2\pr \\
16 & 15 & 45\dg 54.1\pr & N19\dg 11.8\pr & 16 & 16 & 60\dg 54.1\pr & N19\dg 12.3\pr & 16 & 17 & 75\dg 54.1\pr & N19\dg 12.9\pr \\
16 & 18 & 90\dg 54.1\pr & N19\dg 13.5\pr & 16 & 19 & 105\dg 54.1\pr & N19\dg 14.1\pr & 16 & 20 & 120\dg 54.0\pr & N19\dg 14.6\pr \\
16 & 21 & 135\dg 54.0\pr & N19\dg 15.2\pr & 16 & 22 & 150\dg 54.0\pr & N19\dg 15.8\pr & 16 & 23 & 165\dg 54.0\pr & N19\dg 16.3\pr \\
17 & 0 & 180\dg 54.0\pr & N19\dg 16.9\pr & 17 & 1 & 195\dg 54.0\pr & N19\dg 17.5\pr & 17 & 2 & 210\dg 53.9\pr & N19\dg 18.0\pr \\
17 & 3 & 225\dg 53.9\pr & N19\dg 18.6\pr & 17 & 4 & 240\dg 53.9\pr & N19\dg 19.1\pr & 17 & 5 & 255\dg 53.9\pr & N19\dg 19.7\pr \\
17 & 6 & 270\dg 53.9\pr & N19\dg 20.3\pr & 17 & 7 & 285\dg 53.9\pr & N19\dg 20.8\pr & 17 & 8 & 300\dg 53.8\pr & N19\dg 21.4\pr \\
17 & 9 & 315\dg 53.8\pr & N19\dg 22.0\pr & 17 & 10 & 330\dg 53.8\pr & N19\dg 22.5\pr & 17 & 11 & 345\dg 53.8\pr & N19\dg 23.1\pr \\
17 & 12 & 0\dg 53.8\pr & N19\dg 23.6\pr & 17 & 13 & 15\dg 53.7\pr & N19\dg 24.2\pr & 17 & 14 & 30\dg 53.7\pr & N19\dg 24.8\pr \\
17 & 15 & 45\dg 53.7\pr & N19\dg 25.3\pr & 17 & 16 & 60\dg 53.7\pr & N19\dg 25.9\pr & 17 & 17 & 75\dg 53.7\pr & N19\dg 26.4\pr \\
17 & 18 & 90\dg 53.6\pr & N19\dg 27.0\pr & 17 & 19 & 105\dg 53.6\pr & N19\dg 27.5\pr & 17 & 20 & 120\dg 53.6\pr & N19\dg 28.1\pr \\
17 & 21 & 135\dg 53.6\pr & N19\dg 28.6\pr & 17 & 22 & 150\dg 53.5\pr & N19\dg 29.2\pr & 17 & 23 & 165\dg 53.5\pr & N19\dg 29.8\pr \\
18 & 0 & 180\dg 53.5\pr & N19\dg 30.3\pr & 18 & 1 & 195\dg 53.5\pr & N19\dg 30.9\pr & 18 & 2 & 210\dg 53.5\pr & N19\dg 31.4\pr \\
18 & 3 & 225\dg 53.4\pr & N19\dg 32.0\pr & 18 & 4 & 240\dg 53.4\pr & N19\dg 32.5\pr & 18 & 5 & 255\dg 53.4\pr & N19\dg 33.1\pr \\
18 & 6 & 270\dg 53.4\pr & N19\dg 33.6\pr & 18 & 7 & 285\dg 53.3\pr & N19\dg 34.2\pr & 18 & 8 & 300\dg 53.3\pr & N19\dg 34.7\pr \\
18 & 9 & 315\dg 53.3\pr & N19\dg 35.2\pr & 18 & 10 & 330\dg 53.3\pr & N19\dg 35.8\pr & 18 & 11 & 345\dg 53.2\pr & N19\dg 36.3\pr \\
18 & 12 & 0\dg 53.2\pr & N19\dg 36.9\pr & 18 & 13 & 15\dg 53.2\pr & N19\dg 37.4\pr & 18 & 14 & 30\dg 53.2\pr & N19\dg 38.0\pr \\
18 & 15 & 45\dg 53.1\pr & N19\dg 38.5\pr & 18 & 16 & 60\dg 53.1\pr & N19\dg 39.1\pr & 18 & 17 & 75\dg 53.1\pr & N19\dg 39.6\pr \\
18 & 18 & 90\dg 53.0\pr & N19\dg 40.1\pr & 18 & 19 & 105\dg 53.0\pr & N19\dg 40.7\pr & 18 & 20 & 120\dg 53.0\pr & N19\dg 41.2\pr \\
18 & 21 & 135\dg 53.0\pr & N19\dg 41.8\pr & 18 & 22 & 150\dg 52.9\pr & N19\dg 42.3\pr & 18 & 23 & 165\dg 52.9\pr & N19\dg 42.8\pr \\
19 & 0 & 180\dg 52.9\pr & N19\dg 43.4\pr & 19 & 1 & 195\dg 52.8\pr & N19\dg 43.9\pr & 19 & 2 & 210\dg 52.8\pr & N19\dg 44.5\pr \\
19 & 3 & 225\dg 52.8\pr & N19\dg 45.0\pr & 19 & 4 & 240\dg 52.8\pr & N19\dg 45.5\pr & 19 & 5 & 255\dg 52.7\pr & N19\dg 46.1\pr \\
19 & 6 & 270\dg 52.7\pr & N19\dg 46.6\pr & 19 & 7 & 285\dg 52.7\pr & N19\dg 47.1\pr & 19 & 8 & 300\dg 52.6\pr & N19\dg 47.7\pr \\
19 & 9 & 315\dg 52.6\pr & N19\dg 48.2\pr & 19 & 10 & 330\dg 52.6\pr & N19\dg 48.7\pr & 19 & 11 & 345\dg 52.5\pr & N19\dg 49.3\pr \\
19 & 12 & 0\dg 52.5\pr & N19\dg 49.8\pr & 19 & 13 & 15\dg 52.5\pr & N19\dg 50.3\pr & 19 & 14 & 30\dg 52.5\pr & N19\dg 50.9\pr \\
19 & 15 & 45\dg 52.4\pr & N19\dg 51.4\pr & 19 & 16 & 60\dg 52.4\pr & N19\dg 51.9\pr & 19 & 17 & 75\dg 52.4\pr & N19\dg 52.5\pr \\
19 & 18 & 90\dg 52.3\pr & N19\dg 53.0\pr & 19 & 19 & 105\dg 52.3\pr & N19\dg 53.5\pr & 19 & 20 & 120\dg 52.3\pr & N19\dg 54.0\pr \\
19 & 21 & 135\dg 52.2\pr & N19\dg 54.6\pr & 19 & 22 & 150\dg 52.2\pr & N19\dg 55.1\pr & 19 & 23 & 165\dg 52.2\pr & N19\dg 55.6\pr \\
20 & 0 & 180\dg 52.1\pr & N19\dg 56.1\pr & 20 & 1 & 195\dg 52.1\pr & N19\dg 56.7\pr & 20 & 2 & 210\dg 52.0\pr & N19\dg 57.2\pr \\
20 & 3 & 225\dg 52.0\pr & N19\dg 57.7\pr & 20 & 4 & 240\dg 52.0\pr & N19\dg 58.2\pr & 20 & 5 & 255\dg 51.9\pr & N19\dg 58.7\pr \\
20 & 6 & 270\dg 51.9\pr & N19\dg 59.3\pr & 20 & 7 & 285\dg 51.9\pr & N19\dg 59.8\pr & 20 & 8 & 300\dg 51.8\pr & N20\dg 00.3\pr \\
20 & 9 & 315\dg 51.8\pr & N20\dg 00.8\pr & 20 & 10 & 330\dg 51.8\pr & N20\dg 01.3\pr & 20 & 11 & 345\dg 51.7\pr & N20\dg 01.9\pr \\
20 & 12 & 0\dg 51.7\pr & N20\dg 02.4\pr & 20 & 13 & 15\dg 51.7\pr & N20\dg 02.9\pr & 20 & 14 & 30\dg 51.6\pr & N20\dg 03.4\pr \\
20 & 15 & 45\dg 51.6\pr & N20\dg 03.9\pr & 20 & 16 & 60\dg 51.5\pr & N20\dg 04.4\pr & 20 & 17 & 75\dg 51.5\pr & N20\dg 05.0\pr \\
20 & 18 & 90\dg 51.5\pr & N20\dg 05.5\pr & 20 & 19 & 105\dg 51.4\pr & N20\dg 06.0\pr & 20 & 20 & 120\dg 51.4\pr & N20\dg 06.5\pr \\
20 & 21 & 135\dg 51.3\pr & N20\dg 07.0\pr & 20 & 22 & 150\dg 51.3\pr & N20\dg 07.5\pr & 20 & 23 & 165\dg 51.3\pr & N20\dg 08.0\pr \\
21 & 0 & 180\dg 51.2\pr & N20\dg 08.5\pr & 21 & 1 & 195\dg 51.2\pr & N20\dg 09.1\pr & 21 & 2 & 210\dg 51.1\pr & N20\dg 09.6\pr \\
21 & 3 & 225\dg 51.1\pr & N20\dg 10.1\pr & 21 & 4 & 240\dg 51.1\pr & N20\dg 10.6\pr & 21 & 5 & 255\dg 51.0\pr & N20\dg 11.1\pr \\
21 & 6 & 270\dg 51.0\pr & N20\dg 11.6\pr & 21 & 7 & 285\dg 50.9\pr & N20\dg 12.1\pr & 21 & 8 & 300\dg 50.9\pr & N20\dg 12.6\pr \\
21 & 9 & 315\dg 50.9\pr & N20\dg 13.1\pr & 21 & 10 & 330\dg 50.8\pr & N20\dg 13.6\pr & 21 & 11 & 345\dg 50.8\pr & N20\dg 14.1\pr \\
21 & 12 & 0\dg 50.7\pr & N20\dg 14.6\pr & 21 & 13 & 15\dg 50.7\pr & N20\dg 15.1\pr & 21 & 14 & 30\dg 50.6\pr & N20\dg 15.6\pr \\
21 & 15 & 45\dg 50.6\pr & N20\dg 16.1\pr & 21 & 16 & 60\dg 50.6\pr & N20\dg 16.6\pr & 21 & 17 & 75\dg 50.5\pr & N20\dg 17.1\pr \\
21 & 18 & 90\dg 50.5\pr & N20\dg 17.6\pr & 21 & 19 & 105\dg 50.4\pr & N20\dg 18.1\pr & 21 & 20 & 120\dg 50.4\pr & N20\dg 18.6\pr \\
21 & 21 & 135\dg 50.3\pr & N20\dg 19.1\pr & 21 & 22 & 150\dg 50.3\pr & N20\dg 19.6\pr & 21 & 23 & 165\dg 50.2\pr & N20\dg 20.1\pr \\
22 & 0 & 180\dg 50.2\pr & N20\dg 20.6\pr & 22 & 1 & 195\dg 50.2\pr & N20\dg 21.1\pr & 22 & 2 & 210\dg 50.1\pr & N20\dg 21.6\pr \\
22 & 3 & 225\dg 50.1\pr & N20\dg 22.1\pr & 22 & 4 & 240\dg 50.0\pr & N20\dg 22.6\pr & 22 & 5 & 255\dg 50.0\pr & N20\dg 23.1\pr \\
22 & 6 & 270\dg 49.9\pr & N20\dg 23.6\pr & 22 & 7 & 285\dg 49.9\pr & N20\dg 24.1\pr & 22 & 8 & 300\dg 49.8\pr & N20\dg 24.6\pr \\
22 & 9 & 315\dg 49.8\pr & N20\dg 25.0\pr & 22 & 10 & 330\dg 49.7\pr & N20\dg 25.5\pr & 22 & 11 & 345\dg 49.7\pr & N20\dg 26.0\pr \\
22 & 12 & 0\dg 49.6\pr & N20\dg 26.5\pr & 22 & 13 & 15\dg 49.6\pr & N20\dg 27.0\pr & 22 & 14 & 30\dg 49.5\pr & N20\dg 27.5\pr \\
22 & 15 & 45\dg 49.5\pr & N20\dg 28.0\pr & 22 & 16 & 60\dg 49.4\pr & N20\dg 28.5\pr & 22 & 17 & 75\dg 49.4\pr & N20\dg 28.9\pr \\
22 & 18 & 90\dg 49.3\pr & N20\dg 29.4\pr & 22 & 19 & 105\dg 49.3\pr & N20\dg 29.9\pr & 22 & 20 & 120\dg 49.2\pr & N20\dg 30.4\pr \\
22 & 21 & 135\dg 49.2\pr & N20\dg 30.9\pr & 22 & 22 & 150\dg 49.1\pr & N20\dg 31.4\pr & 22 & 23 & 165\dg 49.1\pr & N20\dg 31.8\pr \\
23 & 0 & 180\dg 49.0\pr & N20\dg 32.3\pr & 23 & 1 & 195\dg 49.0\pr & N20\dg 32.8\pr & 23 & 2 & 210\dg 48.9\pr & N20\dg 33.3\pr \\
23 & 3 & 225\dg 48.9\pr & N20\dg 33.8\pr & 23 & 4 & 240\dg 48.8\pr & N20\dg 34.2\pr & 23 & 5 & 255\dg 48.8\pr & N20\dg 34.7\pr \\
23 & 6 & 270\dg 48.7\pr & N20\dg 35.2\pr & 23 & 7 & 285\dg 48.7\pr & N20\dg 35.7\pr & 23 & 8 & 300\dg 48.6\pr & N20\dg 36.2\pr \\
23 & 9 & 315\dg 48.6\pr & N20\dg 36.6\pr & 23 & 10 & 330\dg 48.5\pr & N20\dg 37.1\pr & 23 & 11 & 345\dg 48.5\pr & N20\dg 37.6\pr \\
23 & 12 & 0\dg 48.4\pr & N20\dg 38.0\pr & 23 & 13 & 15\dg 48.4\pr & N20\dg 38.5\pr & 23 & 14 & 30\dg 48.3\pr & N20\dg 39.0\pr \\
23 & 15 & 45\dg 48.3\pr & N20\dg 39.5\pr & 23 & 16 & 60\dg 48.2\pr & N20\dg 39.9\pr & 23 & 17 & 75\dg 48.1\pr & N20\dg 40.4\pr \\
23 & 18 & 90\dg 48.1\pr & N20\dg 40.9\pr & 23 & 19 & 105\dg 48.0\pr & N20\dg 41.3\pr & 23 & 20 & 120\dg 48.0\pr & N20\dg 41.8\pr \\
23 & 21 & 135\dg 47.9\pr & N20\dg 42.3\pr & 23 & 22 & 150\dg 47.9\pr & N20\dg 42.8\pr & 23 & 23 & 165\dg 47.8\pr & N20\dg 43.2\pr \\
24 & 0 & 180\dg 47.8\pr & N20\dg 43.7\pr & 24 & 1 & 195\dg 47.7\pr & N20\dg 44.2\pr & 24 & 2 & 210\dg 47.6\pr & N20\dg 44.6\pr \\
24 & 3 & 225\dg 47.6\pr & N20\dg 45.1\pr & 24 & 4 & 240\dg 47.5\pr & N20\dg 45.5\pr & 24 & 5 & 255\dg 47.5\pr & N20\dg 46.0\pr \\
24 & 6 & 270\dg 47.4\pr & N20\dg 46.5\pr & 24 & 7 & 285\dg 47.4\pr & N20\dg 46.9\pr & 24 & 8 & 300\dg 47.3\pr & N20\dg 47.4\pr \\
24 & 9 & 315\dg 47.2\pr & N20\dg 47.9\pr & 24 & 10 & 330\dg 47.2\pr & N20\dg 48.3\pr & 24 & 11 & 345\dg 47.1\pr & N20\dg 48.8\pr \\
24 & 12 & 0\dg 47.1\pr & N20\dg 49.2\pr & 24 & 13 & 15\dg 47.0\pr & N20\dg 49.7\pr & 24 & 14 & 30\dg 46.9\pr & N20\dg 50.1\pr \\
24 & 15 & 45\dg 46.9\pr & N20\dg 50.6\pr & 24 & 16 & 60\dg 46.8\pr & N20\dg 51.1\pr & 24 & 17 & 75\dg 46.8\pr & N20\dg 51.5\pr \\
24 & 18 & 90\dg 46.7\pr & N20\dg 52.0\pr & 24 & 19 & 105\dg 46.6\pr & N20\dg 52.4\pr & 24 & 20 & 120\dg 46.6\pr & N20\dg 52.9\pr \\
24 & 21 & 135\dg 46.5\pr & N20\dg 53.3\pr & 24 & 22 & 150\dg 46.5\pr & N20\dg 53.8\pr & 24 & 23 & 165\dg 46.4\pr & N20\dg 54.2\pr \\
25 & 0 & 180\dg 46.3\pr & N20\dg 54.7\pr & 25 & 1 & 195\dg 46.3\pr & N20\dg 55.1\pr & 25 & 2 & 210\dg 46.2\pr & N20\dg 55.6\pr \\
25 & 3 & 225\dg 46.2\pr & N20\dg 56.0\pr & 25 & 4 & 240\dg 46.1\pr & N20\dg 56.5\pr & 25 & 5 & 255\dg 46.0\pr & N20\dg 56.9\pr \\
25 & 6 & 270\dg 46.0\pr & N20\dg 57.4\pr & 25 & 7 & 285\dg 45.9\pr & N20\dg 57.8\pr & 25 & 8 & 300\dg 45.8\pr & N20\dg 58.3\pr \\
25 & 9 & 315\dg 45.8\pr & N20\dg 58.7\pr & 25 & 10 & 330\dg 45.7\pr & N20\dg 59.2\pr & 25 & 11 & 345\dg 45.6\pr & N20\dg 59.6\pr \\
25 & 12 & 0\dg 45.6\pr & N21\dg 00.1\pr & 25 & 13 & 15\dg 45.5\pr & N21\dg 00.5\pr & 25 & 14 & 30\dg 45.5\pr & N21\dg 00.9\pr \\
25 & 15 & 45\dg 45.4\pr & N21\dg 01.4\pr & 25 & 16 & 60\dg 45.3\pr & N21\dg 01.8\pr & 25 & 17 & 75\dg 45.3\pr & N21\dg 02.3\pr \\
25 & 18 & 90\dg 45.2\pr & N21\dg 02.7\pr & 25 & 19 & 105\dg 45.1\pr & N21\dg 03.2\pr & 25 & 20 & 120\dg 45.1\pr & N21\dg 03.6\pr \\
25 & 21 & 135\dg 45.0\pr & N21\dg 04.0\pr & 25 & 22 & 150\dg 44.9\pr & N21\dg 04.5\pr & 25 & 23 & 165\dg 44.9\pr & N21\dg 04.9\pr \\
26 & 0 & 180\dg 44.8\pr & N21\dg 05.3\pr & 26 & 1 & 195\dg 44.7\pr & N21\dg 05.8\pr & 26 & 2 & 210\dg 44.7\pr & N21\dg 06.2\pr \\
26 & 3 & 225\dg 44.6\pr & N21\dg 06.6\pr & 26 & 4 & 240\dg 44.5\pr & N21\dg 07.1\pr & 26 & 5 & 255\dg 44.5\pr & N21\dg 07.5\pr \\
26 & 6 & 270\dg 44.4\pr & N21\dg 07.9\pr & 26 & 7 & 285\dg 44.3\pr & N21\dg 08.4\pr & 26 & 8 & 300\dg 44.3\pr & N21\dg 08.8\pr \\
26 & 9 & 315\dg 44.2\pr & N21\dg 09.2\pr & 26 & 10 & 330\dg 44.1\pr & N21\dg 09.7\pr & 26 & 11 & 345\dg 44.1\pr & N21\dg 10.1\pr \\
26 & 12 & 0\dg 44.0\pr & N21\dg 10.5\pr & 26 & 13 & 15\dg 43.9\pr & N21\dg 11.0\pr & 26 & 14 & 30\dg 43.8\pr & N21\dg 11.4\pr \\
26 & 15 & 45\dg 43.8\pr & N21\dg 11.8\pr & 26 & 16 & 60\dg 43.7\pr & N21\dg 12.2\pr & 26 & 17 & 75\dg 43.6\pr & N21\dg 12.7\pr \\
26 & 18 & 90\dg 43.6\pr & N21\dg 13.1\pr & 26 & 19 & 105\dg 43.5\pr & N21\dg 13.5\pr & 26 & 20 & 120\dg 43.4\pr & N21\dg 13.9\pr \\
26 & 21 & 135\dg 43.4\pr & N21\dg 14.4\pr & 26 & 22 & 150\dg 43.3\pr & N21\dg 14.8\pr & 26 & 23 & 165\dg 43.2\pr & N21\dg 15.2\pr \\
27 & 0 & 180\dg 43.1\pr & N21\dg 15.6\pr & 27 & 1 & 195\dg 43.1\pr & N21\dg 16.0\pr & 27 & 2 & 210\dg 43.0\pr & N21\dg 16.5\pr \\
27 & 3 & 225\dg 42.9\pr & N21\dg 16.9\pr & 27 & 4 & 240\dg 42.9\pr & N21\dg 17.3\pr & 27 & 5 & 255\dg 42.8\pr & N21\dg 17.7\pr \\
27 & 6 & 270\dg 42.7\pr & N21\dg 18.1\pr & 27 & 7 & 285\dg 42.6\pr & N21\dg 18.6\pr & 27 & 8 & 300\dg 42.6\pr & N21\dg 19.0\pr \\
27 & 9 & 315\dg 42.5\pr & N21\dg 19.4\pr & 27 & 10 & 330\dg 42.4\pr & N21\dg 19.8\pr & 27 & 11 & 345\dg 42.3\pr & N21\dg 20.2\pr \\
27 & 12 & 0\dg 42.3\pr & N21\dg 20.6\pr & 27 & 13 & 15\dg 42.2\pr & N21\dg 21.0\pr & 27 & 14 & 30\dg 42.1\pr & N21\dg 21.4\pr \\
27 & 15 & 45\dg 42.0\pr & N21\dg 21.9\pr & 27 & 16 & 60\dg 42.0\pr & N21\dg 22.3\pr & 27 & 17 & 75\dg 41.9\pr & N21\dg 22.7\pr \\
27 & 18 & 90\dg 41.8\pr & N21\dg 23.1\pr & 27 & 19 & 105\dg 41.7\pr & N21\dg 23.5\pr & 27 & 20 & 120\dg 41.7\pr & N21\dg 23.9\pr \\
27 & 21 & 135\dg 41.6\pr & N21\dg 24.3\pr & 27 & 22 & 150\dg 41.5\pr & N21\dg 24.7\pr & 27 & 23 & 165\dg 41.4\pr & N21\dg 25.1\pr \\
28 & 0 & 180\dg 41.4\pr & N21\dg 25.5\pr & 28 & 1 & 195\dg 41.3\pr & N21\dg 25.9\pr & 28 & 2 & 210\dg 41.2\pr & N21\dg 26.3\pr \\
28 & 3 & 225\dg 41.1\pr & N21\dg 26.7\pr & 28 & 4 & 240\dg 41.1\pr & N21\dg 27.2\pr & 28 & 5 & 255\dg 41.0\pr & N21\dg 27.6\pr \\
28 & 6 & 270\dg 40.9\pr & N21\dg 28.0\pr & 28 & 7 & 285\dg 40.8\pr & N21\dg 28.4\pr & 28 & 8 & 300\dg 40.7\pr & N21\dg 28.8\pr \\
28 & 9 & 315\dg 40.7\pr & N21\dg 29.2\pr & 28 & 10 & 330\dg 40.6\pr & N21\dg 29.6\pr & 28 & 11 & 345\dg 40.5\pr & N21\dg 30.0\pr \\
28 & 12 & 0\dg 40.4\pr & N21\dg 30.4\pr & 28 & 13 & 15\dg 40.4\pr & N21\dg 30.8\pr & 28 & 14 & 30\dg 40.3\pr & N21\dg 31.1\pr \\
28 & 15 & 45\dg 40.2\pr & N21\dg 31.5\pr & 28 & 16 & 60\dg 40.1\pr & N21\dg 31.9\pr & 28 & 17 & 75\dg 40.0\pr & N21\dg 32.3\pr \\
28 & 18 & 90\dg 40.0\pr & N21\dg 32.7\pr & 28 & 19 & 105\dg 39.9\pr & N21\dg 33.1\pr & 28 & 20 & 120\dg 39.8\pr & N21\dg 33.5\pr \\
28 & 21 & 135\dg 39.7\pr & N21\dg 33.9\pr & 28 & 22 & 150\dg 39.6\pr & N21\dg 34.3\pr & 28 & 23 & 165\dg 39.6\pr & N21\dg 34.7\pr \\
29 & 0 & 180\dg 39.5\pr & N21\dg 35.1\pr & 29 & 1 & 195\dg 39.4\pr & N21\dg 35.5\pr & 29 & 2 & 210\dg 39.3\pr & N21\dg 35.9\pr \\
29 & 3 & 225\dg 39.2\pr & N21\dg 36.2\pr & 29 & 4 & 240\dg 39.1\pr & N21\dg 36.6\pr & 29 & 5 & 255\dg 39.1\pr & N21\dg 37.0\pr \\
29 & 6 & 270\dg 39.0\pr & N21\dg 37.4\pr & 29 & 7 & 285\dg 38.9\pr & N21\dg 37.8\pr & 29 & 8 & 300\dg 38.8\pr & N21\dg 38.2\pr \\
29 & 9 & 315\dg 38.7\pr & N21\dg 38.6\pr & 29 & 10 & 330\dg 38.7\pr & N21\dg 38.9\pr & 29 & 11 & 345\dg 38.6\pr & N21\dg 39.3\pr \\
29 & 12 & 0\dg 38.5\pr & N21\dg 39.7\pr & 29 & 13 & 15\dg 38.4\pr & N21\dg 40.1\pr & 29 & 14 & 30\dg 38.3\pr & N21\dg 40.5\pr \\
29 & 15 & 45\dg 38.2\pr & N21\dg 40.9\pr & 29 & 16 & 60\dg 38.1\pr & N21\dg 41.2\pr & 29 & 17 & 75\dg 38.1\pr & N21\dg 41.6\pr \\
29 & 18 & 90\dg 38.0\pr & N21\dg 42.0\pr & 29 & 19 & 105\dg 37.9\pr & N21\dg 42.4\pr & 29 & 20 & 120\dg 37.8\pr & N21\dg 42.7\pr \\
29 & 21 & 135\dg 37.7\pr & N21\dg 43.1\pr & 29 & 22 & 150\dg 37.6\pr & N21\dg 43.5\pr & 29 & 23 & 165\dg 37.6\pr & N21\dg 43.9\pr \\
30 & 0 & 180\dg 37.5\pr & N21\dg 44.2\pr & 30 & 1 & 195\dg 37.4\pr & N21\dg 44.6\pr & 30 & 2 & 210\dg 37.3\pr & N21\dg 45.0\pr \\
30 & 3 & 225\dg 37.2\pr & N21\dg 45.4\pr & 30 & 4 & 240\dg 37.1\pr & N21\dg 45.7\pr & 30 & 5 & 255\dg 37.0\pr & N21\dg 46.1\pr \\
30 & 6 & 270\dg 37.0\pr & N21\dg 46.5\pr & 30 & 7 & 285\dg 36.9\pr & N21\dg 46.9\pr & 30 & 8 & 300\dg 36.8\pr & N21\dg 47.2\pr \\
30 & 9 & 315\dg 36.7\pr & N21\dg 47.6\pr & 30 & 10 & 330\dg 36.6\pr & N21\dg 48.0\pr & 30 & 11 & 345\dg 36.5\pr & N21\dg 48.3\pr \\
30 & 12 & 0\dg 36.4\pr & N21\dg 48.7\pr & 30 & 13 & 15\dg 36.3\pr & N21\dg 49.1\pr & 30 & 14 & 30\dg 36.3\pr & N21\dg 49.4\pr \\
30 & 15 & 45\dg 36.2\pr & N21\dg 49.8\pr & 30 & 16 & 60\dg 36.1\pr & N21\dg 50.2\pr & 30 & 17 & 75\dg 36.0\pr & N21\dg 50.5\pr \\
30 & 18 & 90\dg 35.9\pr & N21\dg 50.9\pr & 30 & 19 & 105\dg 35.8\pr & N21\dg 51.2\pr & 30 & 20 & 120\dg 35.7\pr & N21\dg 51.6\pr \\
30 & 21 & 135\dg 35.6\pr & N21\dg 52.0\pr & 30 & 22 & 150\dg 35.5\pr & N21\dg 52.3\pr & 30 & 23 & 165\dg 35.4\pr & N21\dg 52.7\pr \\
31 & 0 & 180\dg 35.4\pr & N21\dg 53.0\pr & 31 & 1 & 195\dg 35.3\pr & N21\dg 53.4\pr & 31 & 2 & 210\dg 35.2\pr & N21\dg 53.8\pr \\
31 & 3 & 225\dg 35.1\pr & N21\dg 54.1\pr & 31 & 4 & 240\dg 35.0\pr & N21\dg 54.5\pr & 31 & 5 & 255\dg 34.9\pr & N21\dg 54.8\pr \\
31 & 6 & 270\dg 34.8\pr & N21\dg 55.2\pr & 31 & 7 & 285\dg 34.7\pr & N21\dg 55.5\pr & 31 & 8 & 300\dg 34.6\pr & N21\dg 55.9\pr \\
31 & 9 & 315\dg 34.5\pr & N21\dg 56.2\pr & 31 & 10 & 330\dg 34.4\pr & N21\dg 56.6\pr & 31 & 11 & 345\dg 34.4\pr & N21\dg 56.9\pr \\
31 & 12 & 0\dg 34.3\pr & N21\dg 57.3\pr & 31 & 13 & 15\dg 34.2\pr & N21\dg 57.6\pr & 31 & 14 & 30\dg 34.1\pr & N21\dg 58.0\pr \\
31 & 15 & 45\dg 34.0\pr & N21\dg 58.3\pr & 31 & 16 & 60\dg 33.9\pr & N21\dg 58.7\pr & 31 & 17 & 75\dg 33.8\pr & N21\dg 59.0\pr \\
31 & 18 & 90\dg 33.7\pr & N21\dg 59.4\pr & 31 & 19 & 105\dg 33.6\pr & N21\dg 59.7\pr & 31 & 20 & 120\dg 33.5\pr & N22\dg 00.1\pr \\
31 & 21 & 135\dg 33.4\pr & N22\dg 00.4\pr & 31 & 22 & 150\dg 33.3\pr & N22\dg 00.8\pr & 31 & 23 & 165\dg 33.2\pr & N22\dg 01.1\pr \\
\end{longtable}}

\input{tables/sun_06.tex}
{\footnotesize
\begin{longtable}{rrrr|rrrr|rrrr}
\caption{Sun --- GHA and Declination, July 2026 (UT)}\label{sun07}\\
\toprule
\textbf{d} & \textbf{h} & \textbf{GHA} & \textbf{Dec} & \textbf{d} & \textbf{h} & \textbf{GHA} & \textbf{Dec} & \textbf{d} & \textbf{h} & \textbf{GHA} & \textbf{Dec} \\
\midrule
\endfirsthead
\multicolumn{12}{c}{\textit{Sun --- GHA and Declination, July 2026 (UT) (continued)}}\\
\toprule
\textbf{d} & \textbf{h} & \textbf{GHA} & \textbf{Dec} & \textbf{d} & \textbf{h} & \textbf{GHA} & \textbf{Dec} & \textbf{d} & \textbf{h} & \textbf{GHA} & \textbf{Dec} \\
\midrule
\endhead
\midrule\multicolumn{12}{r}{\textit{continued on next page}}\\
\endfoot
\bottomrule
\endlastfoot
1 & 0 & 179\dg 02.8\pr & N23\dg 06.9\pr & 1 & 1 & 194\dg 02.7\pr & N23\dg 06.8\pr & 1 & 2 & 209\dg 02.6\pr & N23\dg 06.6\pr \\
1 & 3 & 224\dg 02.5\pr & N23\dg 06.4\pr & 1 & 4 & 239\dg 02.4\pr & N23\dg 06.3\pr & 1 & 5 & 254\dg 02.2\pr & N23\dg 06.1\pr \\
1 & 6 & 269\dg 02.1\pr & N23\dg 05.9\pr & 1 & 7 & 284\dg 02.0\pr & N23\dg 05.8\pr & 1 & 8 & 299\dg 01.9\pr & N23\dg 05.6\pr \\
1 & 9 & 314\dg 01.8\pr & N23\dg 05.4\pr & 1 & 10 & 329\dg 01.6\pr & N23\dg 05.3\pr & 1 & 11 & 344\dg 01.5\pr & N23\dg 05.1\pr \\
1 & 12 & 359\dg 01.4\pr & N23\dg 04.9\pr & 1 & 13 & 14\dg 01.3\pr & N23\dg 04.7\pr & 1 & 14 & 29\dg 01.2\pr & N23\dg 04.6\pr \\
1 & 15 & 44\dg 01.0\pr & N23\dg 04.4\pr & 1 & 16 & 59\dg 00.9\pr & N23\dg 04.2\pr & 1 & 17 & 74\dg 00.8\pr & N23\dg 04.0\pr \\
1 & 18 & 89\dg 00.7\pr & N23\dg 03.9\pr & 1 & 19 & 104\dg 00.6\pr & N23\dg 03.7\pr & 1 & 20 & 119\dg 00.4\pr & N23\dg 03.5\pr \\
1 & 21 & 134\dg 00.3\pr & N23\dg 03.3\pr & 1 & 22 & 149\dg 00.2\pr & N23\dg 03.1\pr & 1 & 23 & 164\dg 00.1\pr & N23\dg 03.0\pr \\
2 & 0 & 179\dg 00.0\pr & N23\dg 02.8\pr & 2 & 1 & 193\dg 59.8\pr & N23\dg 02.6\pr & 2 & 2 & 208\dg 59.7\pr & N23\dg 02.4\pr \\
2 & 3 & 223\dg 59.6\pr & N23\dg 02.2\pr & 2 & 4 & 238\dg 59.5\pr & N23\dg 02.1\pr & 2 & 5 & 253\dg 59.4\pr & N23\dg 01.9\pr \\
2 & 6 & 268\dg 59.3\pr & N23\dg 01.7\pr & 2 & 7 & 283\dg 59.1\pr & N23\dg 01.5\pr & 2 & 8 & 298\dg 59.0\pr & N23\dg 01.3\pr \\
2 & 9 & 313\dg 58.9\pr & N23\dg 01.1\pr & 2 & 10 & 328\dg 58.8\pr & N23\dg 00.9\pr & 2 & 11 & 343\dg 58.7\pr & N23\dg 00.7\pr \\
2 & 12 & 358\dg 58.5\pr & N23\dg 00.6\pr & 2 & 13 & 13\dg 58.4\pr & N23\dg 00.4\pr & 2 & 14 & 28\dg 58.3\pr & N23\dg 00.2\pr \\
2 & 15 & 43\dg 58.2\pr & N23\dg 00.0\pr & 2 & 16 & 58\dg 58.1\pr & N22\dg 59.8\pr & 2 & 17 & 73\dg 58.0\pr & N22\dg 59.6\pr \\
2 & 18 & 88\dg 57.8\pr & N22\dg 59.4\pr & 2 & 19 & 103\dg 57.7\pr & N22\dg 59.2\pr & 2 & 20 & 118\dg 57.6\pr & N22\dg 59.0\pr \\
2 & 21 & 133\dg 57.5\pr & N22\dg 58.8\pr & 2 & 22 & 148\dg 57.4\pr & N22\dg 58.6\pr & 2 & 23 & 163\dg 57.3\pr & N22\dg 58.4\pr \\
3 & 0 & 178\dg 57.2\pr & N22\dg 58.2\pr & 3 & 1 & 193\dg 57.0\pr & N22\dg 58.0\pr & 3 & 2 & 208\dg 56.9\pr & N22\dg 57.8\pr \\
3 & 3 & 223\dg 56.8\pr & N22\dg 57.6\pr & 3 & 4 & 238\dg 56.7\pr & N22\dg 57.4\pr & 3 & 5 & 253\dg 56.6\pr & N22\dg 57.2\pr \\
3 & 6 & 268\dg 56.5\pr & N22\dg 57.0\pr & 3 & 7 & 283\dg 56.3\pr & N22\dg 56.8\pr & 3 & 8 & 298\dg 56.2\pr & N22\dg 56.6\pr \\
3 & 9 & 313\dg 56.1\pr & N22\dg 56.4\pr & 3 & 10 & 328\dg 56.0\pr & N22\dg 56.2\pr & 3 & 11 & 343\dg 55.9\pr & N22\dg 56.0\pr \\
3 & 12 & 358\dg 55.8\pr & N22\dg 55.8\pr & 3 & 13 & 13\dg 55.7\pr & N22\dg 55.6\pr & 3 & 14 & 28\dg 55.5\pr & N22\dg 55.4\pr \\
3 & 15 & 43\dg 55.4\pr & N22\dg 55.2\pr & 3 & 16 & 58\dg 55.3\pr & N22\dg 55.0\pr & 3 & 17 & 73\dg 55.2\pr & N22\dg 54.8\pr \\
3 & 18 & 88\dg 55.1\pr & N22\dg 54.5\pr & 3 & 19 & 103\dg 55.0\pr & N22\dg 54.3\pr & 3 & 20 & 118\dg 54.9\pr & N22\dg 54.1\pr \\
3 & 21 & 133\dg 54.7\pr & N22\dg 53.9\pr & 3 & 22 & 148\dg 54.6\pr & N22\dg 53.7\pr & 3 & 23 & 163\dg 54.5\pr & N22\dg 53.5\pr \\
4 & 0 & 178\dg 54.4\pr & N22\dg 53.3\pr & 4 & 1 & 193\dg 54.3\pr & N22\dg 53.0\pr & 4 & 2 & 208\dg 54.2\pr & N22\dg 52.8\pr \\
4 & 3 & 223\dg 54.1\pr & N22\dg 52.6\pr & 4 & 4 & 238\dg 54.0\pr & N22\dg 52.4\pr & 4 & 5 & 253\dg 53.9\pr & N22\dg 52.2\pr \\
4 & 6 & 268\dg 53.7\pr & N22\dg 52.0\pr & 4 & 7 & 283\dg 53.6\pr & N22\dg 51.7\pr & 4 & 8 & 298\dg 53.5\pr & N22\dg 51.5\pr \\
4 & 9 & 313\dg 53.4\pr & N22\dg 51.3\pr & 4 & 10 & 328\dg 53.3\pr & N22\dg 51.1\pr & 4 & 11 & 343\dg 53.2\pr & N22\dg 50.9\pr \\
4 & 12 & 358\dg 53.1\pr & N22\dg 50.6\pr & 4 & 13 & 13\dg 53.0\pr & N22\dg 50.4\pr & 4 & 14 & 28\dg 52.9\pr & N22\dg 50.2\pr \\
4 & 15 & 43\dg 52.7\pr & N22\dg 50.0\pr & 4 & 16 & 58\dg 52.6\pr & N22\dg 49.7\pr & 4 & 17 & 73\dg 52.5\pr & N22\dg 49.5\pr \\
4 & 18 & 88\dg 52.4\pr & N22\dg 49.3\pr & 4 & 19 & 103\dg 52.3\pr & N22\dg 49.1\pr & 4 & 20 & 118\dg 52.2\pr & N22\dg 48.8\pr \\
4 & 21 & 133\dg 52.1\pr & N22\dg 48.6\pr & 4 & 22 & 148\dg 52.0\pr & N22\dg 48.4\pr & 4 & 23 & 163\dg 51.9\pr & N22\dg 48.1\pr \\
5 & 0 & 178\dg 51.8\pr & N22\dg 47.9\pr & 5 & 1 & 193\dg 51.6\pr & N22\dg 47.7\pr & 5 & 2 & 208\dg 51.5\pr & N22\dg 47.4\pr \\
5 & 3 & 223\dg 51.4\pr & N22\dg 47.2\pr & 5 & 4 & 238\dg 51.3\pr & N22\dg 47.0\pr & 5 & 5 & 253\dg 51.2\pr & N22\dg 46.7\pr \\
5 & 6 & 268\dg 51.1\pr & N22\dg 46.5\pr & 5 & 7 & 283\dg 51.0\pr & N22\dg 46.3\pr & 5 & 8 & 298\dg 50.9\pr & N22\dg 46.0\pr \\
5 & 9 & 313\dg 50.8\pr & N22\dg 45.8\pr & 5 & 10 & 328\dg 50.7\pr & N22\dg 45.6\pr & 5 & 11 & 343\dg 50.6\pr & N22\dg 45.3\pr \\
5 & 12 & 358\dg 50.5\pr & N22\dg 45.1\pr & 5 & 13 & 13\dg 50.3\pr & N22\dg 44.8\pr & 5 & 14 & 28\dg 50.2\pr & N22\dg 44.6\pr \\
5 & 15 & 43\dg 50.1\pr & N22\dg 44.4\pr & 5 & 16 & 58\dg 50.0\pr & N22\dg 44.1\pr & 5 & 17 & 73\dg 49.9\pr & N22\dg 43.9\pr \\
5 & 18 & 88\dg 49.8\pr & N22\dg 43.6\pr & 5 & 19 & 103\dg 49.7\pr & N22\dg 43.4\pr & 5 & 20 & 118\dg 49.6\pr & N22\dg 43.1\pr \\
5 & 21 & 133\dg 49.5\pr & N22\dg 42.9\pr & 5 & 22 & 148\dg 49.4\pr & N22\dg 42.6\pr & 5 & 23 & 163\dg 49.3\pr & N22\dg 42.4\pr \\
6 & 0 & 178\dg 49.2\pr & N22\dg 42.1\pr & 6 & 1 & 193\dg 49.1\pr & N22\dg 41.9\pr & 6 & 2 & 208\dg 49.0\pr & N22\dg 41.6\pr \\
6 & 3 & 223\dg 48.9\pr & N22\dg 41.4\pr & 6 & 4 & 238\dg 48.8\pr & N22\dg 41.1\pr & 6 & 5 & 253\dg 48.6\pr & N22\dg 40.9\pr \\
6 & 6 & 268\dg 48.5\pr & N22\dg 40.6\pr & 6 & 7 & 283\dg 48.4\pr & N22\dg 40.4\pr & 6 & 8 & 298\dg 48.3\pr & N22\dg 40.1\pr \\
6 & 9 & 313\dg 48.2\pr & N22\dg 39.9\pr & 6 & 10 & 328\dg 48.1\pr & N22\dg 39.6\pr & 6 & 11 & 343\dg 48.0\pr & N22\dg 39.4\pr \\
6 & 12 & 358\dg 47.9\pr & N22\dg 39.1\pr & 6 & 13 & 13\dg 47.8\pr & N22\dg 38.9\pr & 6 & 14 & 28\dg 47.7\pr & N22\dg 38.6\pr \\
6 & 15 & 43\dg 47.6\pr & N22\dg 38.3\pr & 6 & 16 & 58\dg 47.5\pr & N22\dg 38.1\pr & 6 & 17 & 73\dg 47.4\pr & N22\dg 37.8\pr \\
6 & 18 & 88\dg 47.3\pr & N22\dg 37.6\pr & 6 & 19 & 103\dg 47.2\pr & N22\dg 37.3\pr & 6 & 20 & 118\dg 47.1\pr & N22\dg 37.0\pr \\
6 & 21 & 133\dg 47.0\pr & N22\dg 36.8\pr & 6 & 22 & 148\dg 46.9\pr & N22\dg 36.5\pr & 6 & 23 & 163\dg 46.8\pr & N22\dg 36.3\pr \\
7 & 0 & 178\dg 46.7\pr & N22\dg 36.0\pr & 7 & 1 & 193\dg 46.6\pr & N22\dg 35.7\pr & 7 & 2 & 208\dg 46.5\pr & N22\dg 35.5\pr \\
7 & 3 & 223\dg 46.4\pr & N22\dg 35.2\pr & 7 & 4 & 238\dg 46.3\pr & N22\dg 34.9\pr & 7 & 5 & 253\dg 46.2\pr & N22\dg 34.7\pr \\
7 & 6 & 268\dg 46.1\pr & N22\dg 34.4\pr & 7 & 7 & 283\dg 46.0\pr & N22\dg 34.1\pr & 7 & 8 & 298\dg 45.9\pr & N22\dg 33.9\pr \\
7 & 9 & 313\dg 45.8\pr & N22\dg 33.6\pr & 7 & 10 & 328\dg 45.7\pr & N22\dg 33.3\pr & 7 & 11 & 343\dg 45.6\pr & N22\dg 33.0\pr \\
7 & 12 & 358\dg 45.5\pr & N22\dg 32.8\pr & 7 & 13 & 13\dg 45.4\pr & N22\dg 32.5\pr & 7 & 14 & 28\dg 45.3\pr & N22\dg 32.2\pr \\
7 & 15 & 43\dg 45.2\pr & N22\dg 32.0\pr & 7 & 16 & 58\dg 45.1\pr & N22\dg 31.7\pr & 7 & 17 & 73\dg 45.0\pr & N22\dg 31.4\pr \\
7 & 18 & 88\dg 44.9\pr & N22\dg 31.1\pr & 7 & 19 & 103\dg 44.8\pr & N22\dg 30.8\pr & 7 & 20 & 118\dg 44.7\pr & N22\dg 30.6\pr \\
7 & 21 & 133\dg 44.6\pr & N22\dg 30.3\pr & 7 & 22 & 148\dg 44.5\pr & N22\dg 30.0\pr & 7 & 23 & 163\dg 44.4\pr & N22\dg 29.7\pr \\
8 & 0 & 178\dg 44.3\pr & N22\dg 29.5\pr & 8 & 1 & 193\dg 44.2\pr & N22\dg 29.2\pr & 8 & 2 & 208\dg 44.1\pr & N22\dg 28.9\pr \\
8 & 3 & 223\dg 44.0\pr & N22\dg 28.6\pr & 8 & 4 & 238\dg 43.9\pr & N22\dg 28.3\pr & 8 & 5 & 253\dg 43.8\pr & N22\dg 28.0\pr \\
8 & 6 & 268\dg 43.7\pr & N22\dg 27.8\pr & 8 & 7 & 283\dg 43.6\pr & N22\dg 27.5\pr & 8 & 8 & 298\dg 43.5\pr & N22\dg 27.2\pr \\
8 & 9 & 313\dg 43.4\pr & N22\dg 26.9\pr & 8 & 10 & 328\dg 43.3\pr & N22\dg 26.6\pr & 8 & 11 & 343\dg 43.2\pr & N22\dg 26.3\pr \\
8 & 12 & 358\dg 43.1\pr & N22\dg 26.0\pr & 8 & 13 & 13\dg 43.0\pr & N22\dg 25.7\pr & 8 & 14 & 28\dg 42.9\pr & N22\dg 25.5\pr \\
8 & 15 & 43\dg 42.8\pr & N22\dg 25.2\pr & 8 & 16 & 58\dg 42.8\pr & N22\dg 24.9\pr & 8 & 17 & 73\dg 42.7\pr & N22\dg 24.6\pr \\
8 & 18 & 88\dg 42.6\pr & N22\dg 24.3\pr & 8 & 19 & 103\dg 42.5\pr & N22\dg 24.0\pr & 8 & 20 & 118\dg 42.4\pr & N22\dg 23.7\pr \\
8 & 21 & 133\dg 42.3\pr & N22\dg 23.4\pr & 8 & 22 & 148\dg 42.2\pr & N22\dg 23.1\pr & 8 & 23 & 163\dg 42.1\pr & N22\dg 22.8\pr \\
9 & 0 & 178\dg 42.0\pr & N22\dg 22.5\pr & 9 & 1 & 193\dg 41.9\pr & N22\dg 22.2\pr & 9 & 2 & 208\dg 41.8\pr & N22\dg 21.9\pr \\
9 & 3 & 223\dg 41.7\pr & N22\dg 21.6\pr & 9 & 4 & 238\dg 41.6\pr & N22\dg 21.3\pr & 9 & 5 & 253\dg 41.5\pr & N22\dg 21.0\pr \\
9 & 6 & 268\dg 41.4\pr & N22\dg 20.7\pr & 9 & 7 & 283\dg 41.3\pr & N22\dg 20.4\pr & 9 & 8 & 298\dg 41.3\pr & N22\dg 20.1\pr \\
9 & 9 & 313\dg 41.2\pr & N22\dg 19.8\pr & 9 & 10 & 328\dg 41.1\pr & N22\dg 19.5\pr & 9 & 11 & 343\dg 41.0\pr & N22\dg 19.2\pr \\
9 & 12 & 358\dg 40.9\pr & N22\dg 18.9\pr & 9 & 13 & 13\dg 40.8\pr & N22\dg 18.6\pr & 9 & 14 & 28\dg 40.7\pr & N22\dg 18.3\pr \\
9 & 15 & 43\dg 40.6\pr & N22\dg 18.0\pr & 9 & 16 & 58\dg 40.5\pr & N22\dg 17.7\pr & 9 & 17 & 73\dg 40.4\pr & N22\dg 17.4\pr \\
9 & 18 & 88\dg 40.3\pr & N22\dg 17.1\pr & 9 & 19 & 103\dg 40.3\pr & N22\dg 16.8\pr & 9 & 20 & 118\dg 40.2\pr & N22\dg 16.5\pr \\
9 & 21 & 133\dg 40.1\pr & N22\dg 16.1\pr & 9 & 22 & 148\dg 40.0\pr & N22\dg 15.8\pr & 9 & 23 & 163\dg 39.9\pr & N22\dg 15.5\pr \\
10 & 0 & 178\dg 39.8\pr & N22\dg 15.2\pr & 10 & 1 & 193\dg 39.7\pr & N22\dg 14.9\pr & 10 & 2 & 208\dg 39.6\pr & N22\dg 14.6\pr \\
10 & 3 & 223\dg 39.5\pr & N22\dg 14.3\pr & 10 & 4 & 238\dg 39.5\pr & N22\dg 13.9\pr & 10 & 5 & 253\dg 39.4\pr & N22\dg 13.6\pr \\
10 & 6 & 268\dg 39.3\pr & N22\dg 13.3\pr & 10 & 7 & 283\dg 39.2\pr & N22\dg 13.0\pr & 10 & 8 & 298\dg 39.1\pr & N22\dg 12.7\pr \\
10 & 9 & 313\dg 39.0\pr & N22\dg 12.4\pr & 10 & 10 & 328\dg 38.9\pr & N22\dg 12.0\pr & 10 & 11 & 343\dg 38.8\pr & N22\dg 11.7\pr \\
10 & 12 & 358\dg 38.8\pr & N22\dg 11.4\pr & 10 & 13 & 13\dg 38.7\pr & N22\dg 11.1\pr & 10 & 14 & 28\dg 38.6\pr & N22\dg 10.8\pr \\
10 & 15 & 43\dg 38.5\pr & N22\dg 10.4\pr & 10 & 16 & 58\dg 38.4\pr & N22\dg 10.1\pr & 10 & 17 & 73\dg 38.3\pr & N22\dg 09.8\pr \\
10 & 18 & 88\dg 38.2\pr & N22\dg 09.5\pr & 10 & 19 & 103\dg 38.1\pr & N22\dg 09.1\pr & 10 & 20 & 118\dg 38.1\pr & N22\dg 08.8\pr \\
10 & 21 & 133\dg 38.0\pr & N22\dg 08.5\pr & 10 & 22 & 148\dg 37.9\pr & N22\dg 08.2\pr & 10 & 23 & 163\dg 37.8\pr & N22\dg 07.8\pr \\
11 & 0 & 178\dg 37.7\pr & N22\dg 07.5\pr & 11 & 1 & 193\dg 37.6\pr & N22\dg 07.2\pr & 11 & 2 & 208\dg 37.6\pr & N22\dg 06.8\pr \\
11 & 3 & 223\dg 37.5\pr & N22\dg 06.5\pr & 11 & 4 & 238\dg 37.4\pr & N22\dg 06.2\pr & 11 & 5 & 253\dg 37.3\pr & N22\dg 05.9\pr \\
11 & 6 & 268\dg 37.2\pr & N22\dg 05.5\pr & 11 & 7 & 283\dg 37.1\pr & N22\dg 05.2\pr & 11 & 8 & 298\dg 37.1\pr & N22\dg 04.9\pr \\
11 & 9 & 313\dg 37.0\pr & N22\dg 04.5\pr & 11 & 10 & 328\dg 36.9\pr & N22\dg 04.2\pr & 11 & 11 & 343\dg 36.8\pr & N22\dg 03.8\pr \\
11 & 12 & 358\dg 36.7\pr & N22\dg 03.5\pr & 11 & 13 & 13\dg 36.6\pr & N22\dg 03.2\pr & 11 & 14 & 28\dg 36.6\pr & N22\dg 02.8\pr \\
11 & 15 & 43\dg 36.5\pr & N22\dg 02.5\pr & 11 & 16 & 58\dg 36.4\pr & N22\dg 02.2\pr & 11 & 17 & 73\dg 36.3\pr & N22\dg 01.8\pr \\
11 & 18 & 88\dg 36.2\pr & N22\dg 01.5\pr & 11 & 19 & 103\dg 36.1\pr & N22\dg 01.1\pr & 11 & 20 & 118\dg 36.1\pr & N22\dg 00.8\pr \\
11 & 21 & 133\dg 36.0\pr & N22\dg 00.5\pr & 11 & 22 & 148\dg 35.9\pr & N22\dg 00.1\pr & 11 & 23 & 163\dg 35.8\pr & N21\dg 59.8\pr \\
12 & 0 & 178\dg 35.7\pr & N21\dg 59.4\pr & 12 & 1 & 193\dg 35.7\pr & N21\dg 59.1\pr & 12 & 2 & 208\dg 35.6\pr & N21\dg 58.7\pr \\
12 & 3 & 223\dg 35.5\pr & N21\dg 58.4\pr & 12 & 4 & 238\dg 35.4\pr & N21\dg 58.0\pr & 12 & 5 & 253\dg 35.4\pr & N21\dg 57.7\pr \\
12 & 6 & 268\dg 35.3\pr & N21\dg 57.3\pr & 12 & 7 & 283\dg 35.2\pr & N21\dg 57.0\pr & 12 & 8 & 298\dg 35.1\pr & N21\dg 56.6\pr \\
12 & 9 & 313\dg 35.0\pr & N21\dg 56.3\pr & 12 & 10 & 328\dg 35.0\pr & N21\dg 55.9\pr & 12 & 11 & 343\dg 34.9\pr & N21\dg 55.6\pr \\
12 & 12 & 358\dg 34.8\pr & N21\dg 55.2\pr & 12 & 13 & 13\dg 34.7\pr & N21\dg 54.9\pr & 12 & 14 & 28\dg 34.6\pr & N21\dg 54.5\pr \\
12 & 15 & 43\dg 34.6\pr & N21\dg 54.2\pr & 12 & 16 & 58\dg 34.5\pr & N21\dg 53.8\pr & 12 & 17 & 73\dg 34.4\pr & N21\dg 53.5\pr \\
12 & 18 & 88\dg 34.3\pr & N21\dg 53.1\pr & 12 & 19 & 103\dg 34.3\pr & N21\dg 52.8\pr & 12 & 20 & 118\dg 34.2\pr & N21\dg 52.4\pr \\
12 & 21 & 133\dg 34.1\pr & N21\dg 52.0\pr & 12 & 22 & 148\dg 34.0\pr & N21\dg 51.7\pr & 12 & 23 & 163\dg 34.0\pr & N21\dg 51.3\pr \\
13 & 0 & 178\dg 33.9\pr & N21\dg 51.0\pr & 13 & 1 & 193\dg 33.8\pr & N21\dg 50.6\pr & 13 & 2 & 208\dg 33.7\pr & N21\dg 50.2\pr \\
13 & 3 & 223\dg 33.7\pr & N21\dg 49.9\pr & 13 & 4 & 238\dg 33.6\pr & N21\dg 49.5\pr & 13 & 5 & 253\dg 33.5\pr & N21\dg 49.2\pr \\
13 & 6 & 268\dg 33.4\pr & N21\dg 48.8\pr & 13 & 7 & 283\dg 33.4\pr & N21\dg 48.4\pr & 13 & 8 & 298\dg 33.3\pr & N21\dg 48.1\pr \\
13 & 9 & 313\dg 33.2\pr & N21\dg 47.7\pr & 13 & 10 & 328\dg 33.1\pr & N21\dg 47.3\pr & 13 & 11 & 343\dg 33.1\pr & N21\dg 47.0\pr \\
13 & 12 & 358\dg 33.0\pr & N21\dg 46.6\pr & 13 & 13 & 13\dg 32.9\pr & N21\dg 46.2\pr & 13 & 14 & 28\dg 32.9\pr & N21\dg 45.9\pr \\
13 & 15 & 43\dg 32.8\pr & N21\dg 45.5\pr & 13 & 16 & 58\dg 32.7\pr & N21\dg 45.1\pr & 13 & 17 & 73\dg 32.6\pr & N21\dg 44.8\pr \\
13 & 18 & 88\dg 32.6\pr & N21\dg 44.4\pr & 13 & 19 & 103\dg 32.5\pr & N21\dg 44.0\pr & 13 & 20 & 118\dg 32.4\pr & N21\dg 43.6\pr \\
13 & 21 & 133\dg 32.4\pr & N21\dg 43.3\pr & 13 & 22 & 148\dg 32.3\pr & N21\dg 42.9\pr & 13 & 23 & 163\dg 32.2\pr & N21\dg 42.5\pr \\
14 & 0 & 178\dg 32.1\pr & N21\dg 42.1\pr & 14 & 1 & 193\dg 32.1\pr & N21\dg 41.8\pr & 14 & 2 & 208\dg 32.0\pr & N21\dg 41.4\pr \\
14 & 3 & 223\dg 31.9\pr & N21\dg 41.0\pr & 14 & 4 & 238\dg 31.9\pr & N21\dg 40.6\pr & 14 & 5 & 253\dg 31.8\pr & N21\dg 40.2\pr \\
14 & 6 & 268\dg 31.7\pr & N21\dg 39.9\pr & 14 & 7 & 283\dg 31.7\pr & N21\dg 39.5\pr & 14 & 8 & 298\dg 31.6\pr & N21\dg 39.1\pr \\
14 & 9 & 313\dg 31.5\pr & N21\dg 38.7\pr & 14 & 10 & 328\dg 31.5\pr & N21\dg 38.3\pr & 14 & 11 & 343\dg 31.4\pr & N21\dg 38.0\pr \\
14 & 12 & 358\dg 31.3\pr & N21\dg 37.6\pr & 14 & 13 & 13\dg 31.3\pr & N21\dg 37.2\pr & 14 & 14 & 28\dg 31.2\pr & N21\dg 36.8\pr \\
14 & 15 & 43\dg 31.1\pr & N21\dg 36.4\pr & 14 & 16 & 58\dg 31.1\pr & N21\dg 36.0\pr & 14 & 17 & 73\dg 31.0\pr & N21\dg 35.7\pr \\
14 & 18 & 88\dg 30.9\pr & N21\dg 35.3\pr & 14 & 19 & 103\dg 30.9\pr & N21\dg 34.9\pr & 14 & 20 & 118\dg 30.8\pr & N21\dg 34.5\pr \\
14 & 21 & 133\dg 30.7\pr & N21\dg 34.1\pr & 14 & 22 & 148\dg 30.7\pr & N21\dg 33.7\pr & 14 & 23 & 163\dg 30.6\pr & N21\dg 33.3\pr \\
15 & 0 & 178\dg 30.5\pr & N21\dg 32.9\pr & 15 & 1 & 193\dg 30.5\pr & N21\dg 32.5\pr & 15 & 2 & 208\dg 30.4\pr & N21\dg 32.2\pr \\
15 & 3 & 223\dg 30.3\pr & N21\dg 31.8\pr & 15 & 4 & 238\dg 30.3\pr & N21\dg 31.4\pr & 15 & 5 & 253\dg 30.2\pr & N21\dg 31.0\pr \\
15 & 6 & 268\dg 30.1\pr & N21\dg 30.6\pr & 15 & 7 & 283\dg 30.1\pr & N21\dg 30.2\pr & 15 & 8 & 298\dg 30.0\pr & N21\dg 29.8\pr \\
15 & 9 & 313\dg 30.0\pr & N21\dg 29.4\pr & 15 & 10 & 328\dg 29.9\pr & N21\dg 29.0\pr & 15 & 11 & 343\dg 29.8\pr & N21\dg 28.6\pr \\
15 & 12 & 358\dg 29.8\pr & N21\dg 28.2\pr & 15 & 13 & 13\dg 29.7\pr & N21\dg 27.8\pr & 15 & 14 & 28\dg 29.6\pr & N21\dg 27.4\pr \\
15 & 15 & 43\dg 29.6\pr & N21\dg 27.0\pr & 15 & 16 & 58\dg 29.5\pr & N21\dg 26.6\pr & 15 & 17 & 73\dg 29.5\pr & N21\dg 26.2\pr \\
15 & 18 & 88\dg 29.4\pr & N21\dg 25.8\pr & 15 & 19 & 103\dg 29.3\pr & N21\dg 25.4\pr & 15 & 20 & 118\dg 29.3\pr & N21\dg 25.0\pr \\
15 & 21 & 133\dg 29.2\pr & N21\dg 24.6\pr & 15 & 22 & 148\dg 29.2\pr & N21\dg 24.2\pr & 15 & 23 & 163\dg 29.1\pr & N21\dg 23.8\pr \\
16 & 0 & 178\dg 29.0\pr & N21\dg 23.4\pr & 16 & 1 & 193\dg 29.0\pr & N21\dg 23.0\pr & 16 & 2 & 208\dg 28.9\pr & N21\dg 22.6\pr \\
16 & 3 & 223\dg 28.9\pr & N21\dg 22.1\pr & 16 & 4 & 238\dg 28.8\pr & N21\dg 21.7\pr & 16 & 5 & 253\dg 28.7\pr & N21\dg 21.3\pr \\
16 & 6 & 268\dg 28.7\pr & N21\dg 20.9\pr & 16 & 7 & 283\dg 28.6\pr & N21\dg 20.5\pr & 16 & 8 & 298\dg 28.6\pr & N21\dg 20.1\pr \\
16 & 9 & 313\dg 28.5\pr & N21\dg 19.7\pr & 16 & 10 & 328\dg 28.4\pr & N21\dg 19.3\pr & 16 & 11 & 343\dg 28.4\pr & N21\dg 18.9\pr \\
16 & 12 & 358\dg 28.3\pr & N21\dg 18.4\pr & 16 & 13 & 13\dg 28.3\pr & N21\dg 18.0\pr & 16 & 14 & 28\dg 28.2\pr & N21\dg 17.6\pr \\
16 & 15 & 43\dg 28.2\pr & N21\dg 17.2\pr & 16 & 16 & 58\dg 28.1\pr & N21\dg 16.8\pr & 16 & 17 & 73\dg 28.1\pr & N21\dg 16.4\pr \\
16 & 18 & 88\dg 28.0\pr & N21\dg 16.0\pr & 16 & 19 & 103\dg 27.9\pr & N21\dg 15.5\pr & 16 & 20 & 118\dg 27.9\pr & N21\dg 15.1\pr \\
16 & 21 & 133\dg 27.8\pr & N21\dg 14.7\pr & 16 & 22 & 148\dg 27.8\pr & N21\dg 14.3\pr & 16 & 23 & 163\dg 27.7\pr & N21\dg 13.9\pr \\
17 & 0 & 178\dg 27.7\pr & N21\dg 13.4\pr & 17 & 1 & 193\dg 27.6\pr & N21\dg 13.0\pr & 17 & 2 & 208\dg 27.6\pr & N21\dg 12.6\pr \\
17 & 3 & 223\dg 27.5\pr & N21\dg 12.2\pr & 17 & 4 & 238\dg 27.5\pr & N21\dg 11.7\pr & 17 & 5 & 253\dg 27.4\pr & N21\dg 11.3\pr \\
17 & 6 & 268\dg 27.3\pr & N21\dg 10.9\pr & 17 & 7 & 283\dg 27.3\pr & N21\dg 10.5\pr & 17 & 8 & 298\dg 27.2\pr & N21\dg 10.0\pr \\
17 & 9 & 313\dg 27.2\pr & N21\dg 09.6\pr & 17 & 10 & 328\dg 27.1\pr & N21\dg 09.2\pr & 17 & 11 & 343\dg 27.1\pr & N21\dg 08.8\pr \\
17 & 12 & 358\dg 27.0\pr & N21\dg 08.3\pr & 17 & 13 & 13\dg 27.0\pr & N21\dg 07.9\pr & 17 & 14 & 28\dg 26.9\pr & N21\dg 07.5\pr \\
17 & 15 & 43\dg 26.9\pr & N21\dg 07.0\pr & 17 & 16 & 58\dg 26.8\pr & N21\dg 06.6\pr & 17 & 17 & 73\dg 26.8\pr & N21\dg 06.2\pr \\
17 & 18 & 88\dg 26.7\pr & N21\dg 05.7\pr & 17 & 19 & 103\dg 26.7\pr & N21\dg 05.3\pr & 17 & 20 & 118\dg 26.6\pr & N21\dg 04.9\pr \\
17 & 21 & 133\dg 26.6\pr & N21\dg 04.4\pr & 17 & 22 & 148\dg 26.5\pr & N21\dg 04.0\pr & 17 & 23 & 163\dg 26.5\pr & N21\dg 03.6\pr \\
18 & 0 & 178\dg 26.4\pr & N21\dg 03.1\pr & 18 & 1 & 193\dg 26.4\pr & N21\dg 02.7\pr & 18 & 2 & 208\dg 26.3\pr & N21\dg 02.3\pr \\
18 & 3 & 223\dg 26.3\pr & N21\dg 01.8\pr & 18 & 4 & 238\dg 26.2\pr & N21\dg 01.4\pr & 18 & 5 & 253\dg 26.2\pr & N21\dg 01.0\pr \\
18 & 6 & 268\dg 26.1\pr & N21\dg 00.5\pr & 18 & 7 & 283\dg 26.1\pr & N21\dg 00.1\pr & 18 & 8 & 298\dg 26.1\pr & N20\dg 59.6\pr \\
18 & 9 & 313\dg 26.0\pr & N20\dg 59.2\pr & 18 & 10 & 328\dg 26.0\pr & N20\dg 58.7\pr & 18 & 11 & 343\dg 25.9\pr & N20\dg 58.3\pr \\
18 & 12 & 358\dg 25.9\pr & N20\dg 57.9\pr & 18 & 13 & 13\dg 25.8\pr & N20\dg 57.4\pr & 18 & 14 & 28\dg 25.8\pr & N20\dg 57.0\pr \\
18 & 15 & 43\dg 25.7\pr & N20\dg 56.5\pr & 18 & 16 & 58\dg 25.7\pr & N20\dg 56.1\pr & 18 & 17 & 73\dg 25.6\pr & N20\dg 55.6\pr \\
18 & 18 & 88\dg 25.6\pr & N20\dg 55.2\pr & 18 & 19 & 103\dg 25.6\pr & N20\dg 54.7\pr & 18 & 20 & 118\dg 25.5\pr & N20\dg 54.3\pr \\
18 & 21 & 133\dg 25.5\pr & N20\dg 53.8\pr & 18 & 22 & 148\dg 25.4\pr & N20\dg 53.4\pr & 18 & 23 & 163\dg 25.4\pr & N20\dg 52.9\pr \\
19 & 0 & 178\dg 25.3\pr & N20\dg 52.5\pr & 19 & 1 & 193\dg 25.3\pr & N20\dg 52.0\pr & 19 & 2 & 208\dg 25.2\pr & N20\dg 51.6\pr \\
19 & 3 & 223\dg 25.2\pr & N20\dg 51.1\pr & 19 & 4 & 238\dg 25.2\pr & N20\dg 50.7\pr & 19 & 5 & 253\dg 25.1\pr & N20\dg 50.2\pr \\
19 & 6 & 268\dg 25.1\pr & N20\dg 49.8\pr & 19 & 7 & 283\dg 25.0\pr & N20\dg 49.3\pr & 19 & 8 & 298\dg 25.0\pr & N20\dg 48.9\pr \\
19 & 9 & 313\dg 25.0\pr & N20\dg 48.4\pr & 19 & 10 & 328\dg 24.9\pr & N20\dg 48.0\pr & 19 & 11 & 343\dg 24.9\pr & N20\dg 47.5\pr \\
19 & 12 & 358\dg 24.8\pr & N20\dg 47.0\pr & 19 & 13 & 13\dg 24.8\pr & N20\dg 46.6\pr & 19 & 14 & 28\dg 24.8\pr & N20\dg 46.1\pr \\
19 & 15 & 43\dg 24.7\pr & N20\dg 45.7\pr & 19 & 16 & 58\dg 24.7\pr & N20\dg 45.2\pr & 19 & 17 & 73\dg 24.6\pr & N20\dg 44.7\pr \\
19 & 18 & 88\dg 24.6\pr & N20\dg 44.3\pr & 19 & 19 & 103\dg 24.6\pr & N20\dg 43.8\pr & 19 & 20 & 118\dg 24.5\pr & N20\dg 43.4\pr \\
19 & 21 & 133\dg 24.5\pr & N20\dg 42.9\pr & 19 & 22 & 148\dg 24.4\pr & N20\dg 42.4\pr & 19 & 23 & 163\dg 24.4\pr & N20\dg 42.0\pr \\
20 & 0 & 178\dg 24.4\pr & N20\dg 41.5\pr & 20 & 1 & 193\dg 24.3\pr & N20\dg 41.0\pr & 20 & 2 & 208\dg 24.3\pr & N20\dg 40.6\pr \\
20 & 3 & 223\dg 24.3\pr & N20\dg 40.1\pr & 20 & 4 & 238\dg 24.2\pr & N20\dg 39.6\pr & 20 & 5 & 253\dg 24.2\pr & N20\dg 39.2\pr \\
20 & 6 & 268\dg 24.2\pr & N20\dg 38.7\pr & 20 & 7 & 283\dg 24.1\pr & N20\dg 38.2\pr & 20 & 8 & 298\dg 24.1\pr & N20\dg 37.8\pr \\
20 & 9 & 313\dg 24.0\pr & N20\dg 37.3\pr & 20 & 10 & 328\dg 24.0\pr & N20\dg 36.8\pr & 20 & 11 & 343\dg 24.0\pr & N20\dg 36.3\pr \\
20 & 12 & 358\dg 23.9\pr & N20\dg 35.9\pr & 20 & 13 & 13\dg 23.9\pr & N20\dg 35.4\pr & 20 & 14 & 28\dg 23.9\pr & N20\dg 34.9\pr \\
20 & 15 & 43\dg 23.8\pr & N20\dg 34.4\pr & 20 & 16 & 58\dg 23.8\pr & N20\dg 34.0\pr & 20 & 17 & 73\dg 23.8\pr & N20\dg 33.5\pr \\
20 & 18 & 88\dg 23.7\pr & N20\dg 33.0\pr & 20 & 19 & 103\dg 23.7\pr & N20\dg 32.5\pr & 20 & 20 & 118\dg 23.7\pr & N20\dg 32.1\pr \\
20 & 21 & 133\dg 23.6\pr & N20\dg 31.6\pr & 20 & 22 & 148\dg 23.6\pr & N20\dg 31.1\pr & 20 & 23 & 163\dg 23.6\pr & N20\dg 30.6\pr \\
21 & 0 & 178\dg 23.5\pr & N20\dg 30.1\pr & 21 & 1 & 193\dg 23.5\pr & N20\dg 29.7\pr & 21 & 2 & 208\dg 23.5\pr & N20\dg 29.2\pr \\
21 & 3 & 223\dg 23.5\pr & N20\dg 28.7\pr & 21 & 4 & 238\dg 23.4\pr & N20\dg 28.2\pr & 21 & 5 & 253\dg 23.4\pr & N20\dg 27.7\pr \\
21 & 6 & 268\dg 23.4\pr & N20\dg 27.3\pr & 21 & 7 & 283\dg 23.3\pr & N20\dg 26.8\pr & 21 & 8 & 298\dg 23.3\pr & N20\dg 26.3\pr \\
21 & 9 & 313\dg 23.3\pr & N20\dg 25.8\pr & 21 & 10 & 328\dg 23.2\pr & N20\dg 25.3\pr & 21 & 11 & 343\dg 23.2\pr & N20\dg 24.8\pr \\
21 & 12 & 358\dg 23.2\pr & N20\dg 24.3\pr & 21 & 13 & 13\dg 23.2\pr & N20\dg 23.9\pr & 21 & 14 & 28\dg 23.1\pr & N20\dg 23.4\pr \\
21 & 15 & 43\dg 23.1\pr & N20\dg 22.9\pr & 21 & 16 & 58\dg 23.1\pr & N20\dg 22.4\pr & 21 & 17 & 73\dg 23.1\pr & N20\dg 21.9\pr \\
21 & 18 & 88\dg 23.0\pr & N20\dg 21.4\pr & 21 & 19 & 103\dg 23.0\pr & N20\dg 20.9\pr & 21 & 20 & 118\dg 23.0\pr & N20\dg 20.4\pr \\
21 & 21 & 133\dg 22.9\pr & N20\dg 19.9\pr & 21 & 22 & 148\dg 22.9\pr & N20\dg 19.4\pr & 21 & 23 & 163\dg 22.9\pr & N20\dg 19.0\pr \\
22 & 0 & 178\dg 22.9\pr & N20\dg 18.5\pr & 22 & 1 & 193\dg 22.8\pr & N20\dg 18.0\pr & 22 & 2 & 208\dg 22.8\pr & N20\dg 17.5\pr \\
22 & 3 & 223\dg 22.8\pr & N20\dg 17.0\pr & 22 & 4 & 238\dg 22.8\pr & N20\dg 16.5\pr & 22 & 5 & 253\dg 22.7\pr & N20\dg 16.0\pr \\
22 & 6 & 268\dg 22.7\pr & N20\dg 15.5\pr & 22 & 7 & 283\dg 22.7\pr & N20\dg 15.0\pr & 22 & 8 & 298\dg 22.7\pr & N20\dg 14.5\pr \\
22 & 9 & 313\dg 22.6\pr & N20\dg 14.0\pr & 22 & 10 & 328\dg 22.6\pr & N20\dg 13.5\pr & 22 & 11 & 343\dg 22.6\pr & N20\dg 13.0\pr \\
22 & 12 & 358\dg 22.6\pr & N20\dg 12.5\pr & 22 & 13 & 13\dg 22.6\pr & N20\dg 12.0\pr & 22 & 14 & 28\dg 22.5\pr & N20\dg 11.5\pr \\
22 & 15 & 43\dg 22.5\pr & N20\dg 11.0\pr & 22 & 16 & 58\dg 22.5\pr & N20\dg 10.5\pr & 22 & 17 & 73\dg 22.5\pr & N20\dg 10.0\pr \\
22 & 18 & 88\dg 22.5\pr & N20\dg 09.5\pr & 22 & 19 & 103\dg 22.4\pr & N20\dg 09.0\pr & 22 & 20 & 118\dg 22.4\pr & N20\dg 08.5\pr \\
22 & 21 & 133\dg 22.4\pr & N20\dg 07.9\pr & 22 & 22 & 148\dg 22.4\pr & N20\dg 07.4\pr & 22 & 23 & 163\dg 22.3\pr & N20\dg 06.9\pr \\
23 & 0 & 178\dg 22.3\pr & N20\dg 06.4\pr & 23 & 1 & 193\dg 22.3\pr & N20\dg 05.9\pr & 23 & 2 & 208\dg 22.3\pr & N20\dg 05.4\pr \\
23 & 3 & 223\dg 22.3\pr & N20\dg 04.9\pr & 23 & 4 & 238\dg 22.3\pr & N20\dg 04.4\pr & 23 & 5 & 253\dg 22.2\pr & N20\dg 03.9\pr \\
23 & 6 & 268\dg 22.2\pr & N20\dg 03.4\pr & 23 & 7 & 283\dg 22.2\pr & N20\dg 02.9\pr & 23 & 8 & 298\dg 22.2\pr & N20\dg 02.3\pr \\
23 & 9 & 313\dg 22.2\pr & N20\dg 01.8\pr & 23 & 10 & 328\dg 22.1\pr & N20\dg 01.3\pr & 23 & 11 & 343\dg 22.1\pr & N20\dg 00.8\pr \\
23 & 12 & 358\dg 22.1\pr & N20\dg 00.3\pr & 23 & 13 & 13\dg 22.1\pr & N19\dg 59.8\pr & 23 & 14 & 28\dg 22.1\pr & N19\dg 59.3\pr \\
23 & 15 & 43\dg 22.1\pr & N19\dg 58.7\pr & 23 & 16 & 58\dg 22.1\pr & N19\dg 58.2\pr & 23 & 17 & 73\dg 22.0\pr & N19\dg 57.7\pr \\
23 & 18 & 88\dg 22.0\pr & N19\dg 57.2\pr & 23 & 19 & 103\dg 22.0\pr & N19\dg 56.7\pr & 23 & 20 & 118\dg 22.0\pr & N19\dg 56.1\pr \\
23 & 21 & 133\dg 22.0\pr & N19\dg 55.6\pr & 23 & 22 & 148\dg 22.0\pr & N19\dg 55.1\pr & 23 & 23 & 163\dg 21.9\pr & N19\dg 54.6\pr \\
24 & 0 & 178\dg 21.9\pr & N19\dg 54.1\pr & 24 & 1 & 193\dg 21.9\pr & N19\dg 53.5\pr & 24 & 2 & 208\dg 21.9\pr & N19\dg 53.0\pr \\
24 & 3 & 223\dg 21.9\pr & N19\dg 52.5\pr & 24 & 4 & 238\dg 21.9\pr & N19\dg 52.0\pr & 24 & 5 & 253\dg 21.9\pr & N19\dg 51.4\pr \\
24 & 6 & 268\dg 21.9\pr & N19\dg 50.9\pr & 24 & 7 & 283\dg 21.8\pr & N19\dg 50.4\pr & 24 & 8 & 298\dg 21.8\pr & N19\dg 49.9\pr \\
24 & 9 & 313\dg 21.8\pr & N19\dg 49.3\pr & 24 & 10 & 328\dg 21.8\pr & N19\dg 48.8\pr & 24 & 11 & 343\dg 21.8\pr & N19\dg 48.3\pr \\
24 & 12 & 358\dg 21.8\pr & N19\dg 47.8\pr & 24 & 13 & 13\dg 21.8\pr & N19\dg 47.2\pr & 24 & 14 & 28\dg 21.8\pr & N19\dg 46.7\pr \\
24 & 15 & 43\dg 21.8\pr & N19\dg 46.2\pr & 24 & 16 & 58\dg 21.8\pr & N19\dg 45.6\pr & 24 & 17 & 73\dg 21.7\pr & N19\dg 45.1\pr \\
24 & 18 & 88\dg 21.7\pr & N19\dg 44.6\pr & 24 & 19 & 103\dg 21.7\pr & N19\dg 44.0\pr & 24 & 20 & 118\dg 21.7\pr & N19\dg 43.5\pr \\
24 & 21 & 133\dg 21.7\pr & N19\dg 43.0\pr & 24 & 22 & 148\dg 21.7\pr & N19\dg 42.4\pr & 24 & 23 & 163\dg 21.7\pr & N19\dg 41.9\pr \\
25 & 0 & 178\dg 21.7\pr & N19\dg 41.4\pr & 25 & 1 & 193\dg 21.7\pr & N19\dg 40.8\pr & 25 & 2 & 208\dg 21.7\pr & N19\dg 40.3\pr \\
25 & 3 & 223\dg 21.7\pr & N19\dg 39.8\pr & 25 & 4 & 238\dg 21.7\pr & N19\dg 39.2\pr & 25 & 5 & 253\dg 21.7\pr & N19\dg 38.7\pr \\
25 & 6 & 268\dg 21.7\pr & N19\dg 38.1\pr & 25 & 7 & 283\dg 21.6\pr & N19\dg 37.6\pr & 25 & 8 & 298\dg 21.6\pr & N19\dg 37.1\pr \\
25 & 9 & 313\dg 21.6\pr & N19\dg 36.5\pr & 25 & 10 & 328\dg 21.6\pr & N19\dg 36.0\pr & 25 & 11 & 343\dg 21.6\pr & N19\dg 35.4\pr \\
25 & 12 & 358\dg 21.6\pr & N19\dg 34.9\pr & 25 & 13 & 13\dg 21.6\pr & N19\dg 34.3\pr & 25 & 14 & 28\dg 21.6\pr & N19\dg 33.8\pr \\
25 & 15 & 43\dg 21.6\pr & N19\dg 33.3\pr & 25 & 16 & 58\dg 21.6\pr & N19\dg 32.7\pr & 25 & 17 & 73\dg 21.6\pr & N19\dg 32.2\pr \\
25 & 18 & 88\dg 21.6\pr & N19\dg 31.6\pr & 25 & 19 & 103\dg 21.6\pr & N19\dg 31.1\pr & 25 & 20 & 118\dg 21.6\pr & N19\dg 30.5\pr \\
25 & 21 & 133\dg 21.6\pr & N19\dg 30.0\pr & 25 & 22 & 148\dg 21.6\pr & N19\dg 29.4\pr & 25 & 23 & 163\dg 21.6\pr & N19\dg 28.9\pr \\
26 & 0 & 178\dg 21.6\pr & N19\dg 28.3\pr & 26 & 1 & 193\dg 21.6\pr & N19\dg 27.8\pr & 26 & 2 & 208\dg 21.6\pr & N19\dg 27.2\pr \\
26 & 3 & 223\dg 21.6\pr & N19\dg 26.7\pr & 26 & 4 & 238\dg 21.6\pr & N19\dg 26.1\pr & 26 & 5 & 253\dg 21.6\pr & N19\dg 25.6\pr \\
26 & 6 & 268\dg 21.6\pr & N19\dg 25.0\pr & 26 & 7 & 283\dg 21.6\pr & N19\dg 24.5\pr & 26 & 8 & 298\dg 21.6\pr & N19\dg 23.9\pr \\
26 & 9 & 313\dg 21.6\pr & N19\dg 23.4\pr & 26 & 10 & 328\dg 21.6\pr & N19\dg 22.8\pr & 26 & 11 & 343\dg 21.6\pr & N19\dg 22.3\pr \\
26 & 12 & 358\dg 21.6\pr & N19\dg 21.7\pr & 26 & 13 & 13\dg 21.6\pr & N19\dg 21.1\pr & 26 & 14 & 28\dg 21.6\pr & N19\dg 20.6\pr \\
26 & 15 & 43\dg 21.6\pr & N19\dg 20.0\pr & 26 & 16 & 58\dg 21.6\pr & N19\dg 19.5\pr & 26 & 17 & 73\dg 21.6\pr & N19\dg 18.9\pr \\
26 & 18 & 88\dg 21.6\pr & N19\dg 18.4\pr & 26 & 19 & 103\dg 21.6\pr & N19\dg 17.8\pr & 26 & 20 & 118\dg 21.6\pr & N19\dg 17.2\pr \\
26 & 21 & 133\dg 21.6\pr & N19\dg 16.7\pr & 26 & 22 & 148\dg 21.6\pr & N19\dg 16.1\pr & 26 & 23 & 163\dg 21.6\pr & N19\dg 15.6\pr \\
27 & 0 & 178\dg 21.6\pr & N19\dg 15.0\pr & 27 & 1 & 193\dg 21.6\pr & N19\dg 14.4\pr & 27 & 2 & 208\dg 21.7\pr & N19\dg 13.9\pr \\
27 & 3 & 223\dg 21.7\pr & N19\dg 13.3\pr & 27 & 4 & 238\dg 21.7\pr & N19\dg 12.7\pr & 27 & 5 & 253\dg 21.7\pr & N19\dg 12.2\pr \\
27 & 6 & 268\dg 21.7\pr & N19\dg 11.6\pr & 27 & 7 & 283\dg 21.7\pr & N19\dg 11.0\pr & 27 & 8 & 298\dg 21.7\pr & N19\dg 10.5\pr \\
27 & 9 & 313\dg 21.7\pr & N19\dg 09.9\pr & 27 & 10 & 328\dg 21.7\pr & N19\dg 09.3\pr & 27 & 11 & 343\dg 21.7\pr & N19\dg 08.8\pr \\
27 & 12 & 358\dg 21.7\pr & N19\dg 08.2\pr & 27 & 13 & 13\dg 21.7\pr & N19\dg 07.6\pr & 27 & 14 & 28\dg 21.7\pr & N19\dg 07.1\pr \\
27 & 15 & 43\dg 21.7\pr & N19\dg 06.5\pr & 27 & 16 & 58\dg 21.8\pr & N19\dg 05.9\pr & 27 & 17 & 73\dg 21.8\pr & N19\dg 05.3\pr \\
27 & 18 & 88\dg 21.8\pr & N19\dg 04.8\pr & 27 & 19 & 103\dg 21.8\pr & N19\dg 04.2\pr & 27 & 20 & 118\dg 21.8\pr & N19\dg 03.6\pr \\
27 & 21 & 133\dg 21.8\pr & N19\dg 03.1\pr & 27 & 22 & 148\dg 21.8\pr & N19\dg 02.5\pr & 27 & 23 & 163\dg 21.8\pr & N19\dg 01.9\pr \\
28 & 0 & 178\dg 21.8\pr & N19\dg 01.3\pr & 28 & 1 & 193\dg 21.9\pr & N19\dg 00.7\pr & 28 & 2 & 208\dg 21.9\pr & N19\dg 00.2\pr \\
28 & 3 & 223\dg 21.9\pr & N18\dg 59.6\pr & 28 & 4 & 238\dg 21.9\pr & N18\dg 59.0\pr & 28 & 5 & 253\dg 21.9\pr & N18\dg 58.4\pr \\
28 & 6 & 268\dg 21.9\pr & N18\dg 57.9\pr & 28 & 7 & 283\dg 21.9\pr & N18\dg 57.3\pr & 28 & 8 & 298\dg 21.9\pr & N18\dg 56.7\pr \\
28 & 9 & 313\dg 22.0\pr & N18\dg 56.1\pr & 28 & 10 & 328\dg 22.0\pr & N18\dg 55.5\pr & 28 & 11 & 343\dg 22.0\pr & N18\dg 55.0\pr \\
28 & 12 & 358\dg 22.0\pr & N18\dg 54.4\pr & 28 & 13 & 13\dg 22.0\pr & N18\dg 53.8\pr & 28 & 14 & 28\dg 22.0\pr & N18\dg 53.2\pr \\
28 & 15 & 43\dg 22.0\pr & N18\dg 52.6\pr & 28 & 16 & 58\dg 22.1\pr & N18\dg 52.0\pr & 28 & 17 & 73\dg 22.1\pr & N18\dg 51.5\pr \\
28 & 18 & 88\dg 22.1\pr & N18\dg 50.9\pr & 28 & 19 & 103\dg 22.1\pr & N18\dg 50.3\pr & 28 & 20 & 118\dg 22.1\pr & N18\dg 49.7\pr \\
28 & 21 & 133\dg 22.1\pr & N18\dg 49.1\pr & 28 & 22 & 148\dg 22.2\pr & N18\dg 48.5\pr & 28 & 23 & 163\dg 22.2\pr & N18\dg 47.9\pr \\
29 & 0 & 178\dg 22.2\pr & N18\dg 47.3\pr & 29 & 1 & 193\dg 22.2\pr & N18\dg 46.8\pr & 29 & 2 & 208\dg 22.2\pr & N18\dg 46.2\pr \\
29 & 3 & 223\dg 22.2\pr & N18\dg 45.6\pr & 29 & 4 & 238\dg 22.3\pr & N18\dg 45.0\pr & 29 & 5 & 253\dg 22.3\pr & N18\dg 44.4\pr \\
29 & 6 & 268\dg 22.3\pr & N18\dg 43.8\pr & 29 & 7 & 283\dg 22.3\pr & N18\dg 43.2\pr & 29 & 8 & 298\dg 22.3\pr & N18\dg 42.6\pr \\
29 & 9 & 313\dg 22.4\pr & N18\dg 42.0\pr & 29 & 10 & 328\dg 22.4\pr & N18\dg 41.4\pr & 29 & 11 & 343\dg 22.4\pr & N18\dg 40.8\pr \\
29 & 12 & 358\dg 22.4\pr & N18\dg 40.2\pr & 29 & 13 & 13\dg 22.4\pr & N18\dg 39.6\pr & 29 & 14 & 28\dg 22.5\pr & N18\dg 39.0\pr \\
29 & 15 & 43\dg 22.5\pr & N18\dg 38.5\pr & 29 & 16 & 58\dg 22.5\pr & N18\dg 37.9\pr & 29 & 17 & 73\dg 22.5\pr & N18\dg 37.3\pr \\
29 & 18 & 88\dg 22.6\pr & N18\dg 36.7\pr & 29 & 19 & 103\dg 22.6\pr & N18\dg 36.1\pr & 29 & 20 & 118\dg 22.6\pr & N18\dg 35.5\pr \\
29 & 21 & 133\dg 22.6\pr & N18\dg 34.9\pr & 29 & 22 & 148\dg 22.6\pr & N18\dg 34.3\pr & 29 & 23 & 163\dg 22.7\pr & N18\dg 33.7\pr \\
30 & 0 & 178\dg 22.7\pr & N18\dg 33.1\pr & 30 & 1 & 193\dg 22.7\pr & N18\dg 32.5\pr & 30 & 2 & 208\dg 22.7\pr & N18\dg 31.9\pr \\
30 & 3 & 223\dg 22.8\pr & N18\dg 31.3\pr & 30 & 4 & 238\dg 22.8\pr & N18\dg 30.6\pr & 30 & 5 & 253\dg 22.8\pr & N18\dg 30.0\pr \\
30 & 6 & 268\dg 22.8\pr & N18\dg 29.4\pr & 30 & 7 & 283\dg 22.9\pr & N18\dg 28.8\pr & 30 & 8 & 298\dg 22.9\pr & N18\dg 28.2\pr \\
30 & 9 & 313\dg 22.9\pr & N18\dg 27.6\pr & 30 & 10 & 328\dg 23.0\pr & N18\dg 27.0\pr & 30 & 11 & 343\dg 23.0\pr & N18\dg 26.4\pr \\
30 & 12 & 358\dg 23.0\pr & N18\dg 25.8\pr & 30 & 13 & 13\dg 23.0\pr & N18\dg 25.2\pr & 30 & 14 & 28\dg 23.1\pr & N18\dg 24.6\pr \\
30 & 15 & 43\dg 23.1\pr & N18\dg 24.0\pr & 30 & 16 & 58\dg 23.1\pr & N18\dg 23.4\pr & 30 & 17 & 73\dg 23.1\pr & N18\dg 22.8\pr \\
30 & 18 & 88\dg 23.2\pr & N18\dg 22.1\pr & 30 & 19 & 103\dg 23.2\pr & N18\dg 21.5\pr & 30 & 20 & 118\dg 23.2\pr & N18\dg 20.9\pr \\
30 & 21 & 133\dg 23.3\pr & N18\dg 20.3\pr & 30 & 22 & 148\dg 23.3\pr & N18\dg 19.7\pr & 30 & 23 & 163\dg 23.3\pr & N18\dg 19.1\pr \\
31 & 0 & 178\dg 23.4\pr & N18\dg 18.5\pr & 31 & 1 & 193\dg 23.4\pr & N18\dg 17.9\pr & 31 & 2 & 208\dg 23.4\pr & N18\dg 17.2\pr \\
31 & 3 & 223\dg 23.4\pr & N18\dg 16.6\pr & 31 & 4 & 238\dg 23.5\pr & N18\dg 16.0\pr & 31 & 5 & 253\dg 23.5\pr & N18\dg 15.4\pr \\
31 & 6 & 268\dg 23.5\pr & N18\dg 14.8\pr & 31 & 7 & 283\dg 23.6\pr & N18\dg 14.2\pr & 31 & 8 & 298\dg 23.6\pr & N18\dg 13.5\pr \\
31 & 9 & 313\dg 23.6\pr & N18\dg 12.9\pr & 31 & 10 & 328\dg 23.7\pr & N18\dg 12.3\pr & 31 & 11 & 343\dg 23.7\pr & N18\dg 11.7\pr \\
31 & 12 & 358\dg 23.7\pr & N18\dg 11.1\pr & 31 & 13 & 13\dg 23.8\pr & N18\dg 10.4\pr & 31 & 14 & 28\dg 23.8\pr & N18\dg 09.8\pr \\
31 & 15 & 43\dg 23.8\pr & N18\dg 09.2\pr & 31 & 16 & 58\dg 23.9\pr & N18\dg 08.6\pr & 31 & 17 & 73\dg 23.9\pr & N18\dg 07.9\pr \\
31 & 18 & 88\dg 23.9\pr & N18\dg 07.3\pr & 31 & 19 & 103\dg 24.0\pr & N18\dg 06.7\pr & 31 & 20 & 118\dg 24.0\pr & N18\dg 06.1\pr \\
31 & 21 & 133\dg 24.1\pr & N18\dg 05.5\pr & 31 & 22 & 148\dg 24.1\pr & N18\dg 04.8\pr & 31 & 23 & 163\dg 24.1\pr & N18\dg 04.2\pr \\
\end{longtable}}

{\footnotesize
\begin{longtable}{rrrr|rrrr|rrrr}
\caption{Sun --- GHA and Declination, August 2026 (UT)}\label{sun08}\\
\toprule
\textbf{d} & \textbf{h} & \textbf{GHA} & \textbf{Dec} & \textbf{d} & \textbf{h} & \textbf{GHA} & \textbf{Dec} & \textbf{d} & \textbf{h} & \textbf{GHA} & \textbf{Dec} \\
\midrule
\endfirsthead
\multicolumn{12}{c}{\textit{Sun --- GHA and Declination, August 2026 (UT) (continued)}}\\
\toprule
\textbf{d} & \textbf{h} & \textbf{GHA} & \textbf{Dec} & \textbf{d} & \textbf{h} & \textbf{GHA} & \textbf{Dec} & \textbf{d} & \textbf{h} & \textbf{GHA} & \textbf{Dec} \\
\midrule
\endhead
\midrule\multicolumn{12}{r}{\textit{continued on next page}}\\
\endfoot
\bottomrule
\endlastfoot
1 & 0 & 178\dg 24.2\pr & N18\dg 03.6\pr & 1 & 1 & 193\dg 24.2\pr & N18\dg 02.9\pr & 1 & 2 & 208\dg 24.2\pr & N18\dg 02.3\pr \\
1 & 3 & 223\dg 24.3\pr & N18\dg 01.7\pr & 1 & 4 & 238\dg 24.3\pr & N18\dg 01.1\pr & 1 & 5 & 253\dg 24.3\pr & N18\dg 00.4\pr \\
1 & 6 & 268\dg 24.4\pr & N17\dg 59.8\pr & 1 & 7 & 283\dg 24.4\pr & N17\dg 59.2\pr & 1 & 8 & 298\dg 24.5\pr & N17\dg 58.5\pr \\
1 & 9 & 313\dg 24.5\pr & N17\dg 57.9\pr & 1 & 10 & 328\dg 24.5\pr & N17\dg 57.3\pr & 1 & 11 & 343\dg 24.6\pr & N17\dg 56.7\pr \\
1 & 12 & 358\dg 24.6\pr & N17\dg 56.0\pr & 1 & 13 & 13\dg 24.7\pr & N17\dg 55.4\pr & 1 & 14 & 28\dg 24.7\pr & N17\dg 54.8\pr \\
1 & 15 & 43\dg 24.7\pr & N17\dg 54.1\pr & 1 & 16 & 58\dg 24.8\pr & N17\dg 53.5\pr & 1 & 17 & 73\dg 24.8\pr & N17\dg 52.8\pr \\
1 & 18 & 88\dg 24.9\pr & N17\dg 52.2\pr & 1 & 19 & 103\dg 24.9\pr & N17\dg 51.6\pr & 1 & 20 & 118\dg 24.9\pr & N17\dg 50.9\pr \\
1 & 21 & 133\dg 25.0\pr & N17\dg 50.3\pr & 1 & 22 & 148\dg 25.0\pr & N17\dg 49.7\pr & 1 & 23 & 163\dg 25.1\pr & N17\dg 49.0\pr \\
2 & 0 & 178\dg 25.1\pr & N17\dg 48.4\pr & 2 & 1 & 193\dg 25.2\pr & N17\dg 47.8\pr & 2 & 2 & 208\dg 25.2\pr & N17\dg 47.1\pr \\
2 & 3 & 223\dg 25.2\pr & N17\dg 46.5\pr & 2 & 4 & 238\dg 25.3\pr & N17\dg 45.8\pr & 2 & 5 & 253\dg 25.3\pr & N17\dg 45.2\pr \\
2 & 6 & 268\dg 25.4\pr & N17\dg 44.5\pr & 2 & 7 & 283\dg 25.4\pr & N17\dg 43.9\pr & 2 & 8 & 298\dg 25.5\pr & N17\dg 43.3\pr \\
2 & 9 & 313\dg 25.5\pr & N17\dg 42.6\pr & 2 & 10 & 328\dg 25.6\pr & N17\dg 42.0\pr & 2 & 11 & 343\dg 25.6\pr & N17\dg 41.3\pr \\
2 & 12 & 358\dg 25.7\pr & N17\dg 40.7\pr & 2 & 13 & 13\dg 25.7\pr & N17\dg 40.0\pr & 2 & 14 & 28\dg 25.7\pr & N17\dg 39.4\pr \\
2 & 15 & 43\dg 25.8\pr & N17\dg 38.8\pr & 2 & 16 & 58\dg 25.8\pr & N17\dg 38.1\pr & 2 & 17 & 73\dg 25.9\pr & N17\dg 37.5\pr \\
2 & 18 & 88\dg 25.9\pr & N17\dg 36.8\pr & 2 & 19 & 103\dg 26.0\pr & N17\dg 36.2\pr & 2 & 20 & 118\dg 26.0\pr & N17\dg 35.5\pr \\
2 & 21 & 133\dg 26.1\pr & N17\dg 34.9\pr & 2 & 22 & 148\dg 26.1\pr & N17\dg 34.2\pr & 2 & 23 & 163\dg 26.2\pr & N17\dg 33.6\pr \\
3 & 0 & 178\dg 26.2\pr & N17\dg 32.9\pr & 3 & 1 & 193\dg 26.3\pr & N17\dg 32.3\pr & 3 & 2 & 208\dg 26.3\pr & N17\dg 31.6\pr \\
3 & 3 & 223\dg 26.4\pr & N17\dg 31.0\pr & 3 & 4 & 238\dg 26.4\pr & N17\dg 30.3\pr & 3 & 5 & 253\dg 26.5\pr & N17\dg 29.7\pr \\
3 & 6 & 268\dg 26.5\pr & N17\dg 29.0\pr & 3 & 7 & 283\dg 26.6\pr & N17\dg 28.3\pr & 3 & 8 & 298\dg 26.6\pr & N17\dg 27.7\pr \\
3 & 9 & 313\dg 26.7\pr & N17\dg 27.0\pr & 3 & 10 & 328\dg 26.7\pr & N17\dg 26.4\pr & 3 & 11 & 343\dg 26.8\pr & N17\dg 25.7\pr \\
3 & 12 & 358\dg 26.8\pr & N17\dg 25.1\pr & 3 & 13 & 13\dg 26.9\pr & N17\dg 24.4\pr & 3 & 14 & 28\dg 26.9\pr & N17\dg 23.8\pr \\
3 & 15 & 43\dg 27.0\pr & N17\dg 23.1\pr & 3 & 16 & 58\dg 27.1\pr & N17\dg 22.4\pr & 3 & 17 & 73\dg 27.1\pr & N17\dg 21.8\pr \\
3 & 18 & 88\dg 27.2\pr & N17\dg 21.1\pr & 3 & 19 & 103\dg 27.2\pr & N17\dg 20.5\pr & 3 & 20 & 118\dg 27.3\pr & N17\dg 19.8\pr \\
3 & 21 & 133\dg 27.3\pr & N17\dg 19.1\pr & 3 & 22 & 148\dg 27.4\pr & N17\dg 18.5\pr & 3 & 23 & 163\dg 27.4\pr & N17\dg 17.8\pr \\
4 & 0 & 178\dg 27.5\pr & N17\dg 17.2\pr & 4 & 1 & 193\dg 27.5\pr & N17\dg 16.5\pr & 4 & 2 & 208\dg 27.6\pr & N17\dg 15.8\pr \\
4 & 3 & 223\dg 27.7\pr & N17\dg 15.2\pr & 4 & 4 & 238\dg 27.7\pr & N17\dg 14.5\pr & 4 & 5 & 253\dg 27.8\pr & N17\dg 13.8\pr \\
4 & 6 & 268\dg 27.8\pr & N17\dg 13.2\pr & 4 & 7 & 283\dg 27.9\pr & N17\dg 12.5\pr & 4 & 8 & 298\dg 27.9\pr & N17\dg 11.8\pr \\
4 & 9 & 313\dg 28.0\pr & N17\dg 11.2\pr & 4 & 10 & 328\dg 28.1\pr & N17\dg 10.5\pr & 4 & 11 & 343\dg 28.1\pr & N17\dg 09.8\pr \\
4 & 12 & 358\dg 28.2\pr & N17\dg 09.2\pr & 4 & 13 & 13\dg 28.2\pr & N17\dg 08.5\pr & 4 & 14 & 28\dg 28.3\pr & N17\dg 07.8\pr \\
4 & 15 & 43\dg 28.4\pr & N17\dg 07.2\pr & 4 & 16 & 58\dg 28.4\pr & N17\dg 06.5\pr & 4 & 17 & 73\dg 28.5\pr & N17\dg 05.8\pr \\
4 & 18 & 88\dg 28.5\pr & N17\dg 05.1\pr & 4 & 19 & 103\dg 28.6\pr & N17\dg 04.5\pr & 4 & 20 & 118\dg 28.7\pr & N17\dg 03.8\pr \\
4 & 21 & 133\dg 28.7\pr & N17\dg 03.1\pr & 4 & 22 & 148\dg 28.8\pr & N17\dg 02.5\pr & 4 & 23 & 163\dg 28.8\pr & N17\dg 01.8\pr \\
5 & 0 & 178\dg 28.9\pr & N17\dg 01.1\pr & 5 & 1 & 193\dg 29.0\pr & N17\dg 00.4\pr & 5 & 2 & 208\dg 29.0\pr & N16\dg 59.8\pr \\
5 & 3 & 223\dg 29.1\pr & N16\dg 59.1\pr & 5 & 4 & 238\dg 29.2\pr & N16\dg 58.4\pr & 5 & 5 & 253\dg 29.2\pr & N16\dg 57.7\pr \\
5 & 6 & 268\dg 29.3\pr & N16\dg 57.1\pr & 5 & 7 & 283\dg 29.3\pr & N16\dg 56.4\pr & 5 & 8 & 298\dg 29.4\pr & N16\dg 55.7\pr \\
5 & 9 & 313\dg 29.5\pr & N16\dg 55.0\pr & 5 & 10 & 328\dg 29.5\pr & N16\dg 54.3\pr & 5 & 11 & 343\dg 29.6\pr & N16\dg 53.7\pr \\
5 & 12 & 358\dg 29.7\pr & N16\dg 53.0\pr & 5 & 13 & 13\dg 29.7\pr & N16\dg 52.3\pr & 5 & 14 & 28\dg 29.8\pr & N16\dg 51.6\pr \\
5 & 15 & 43\dg 29.9\pr & N16\dg 50.9\pr & 5 & 16 & 58\dg 29.9\pr & N16\dg 50.3\pr & 5 & 17 & 73\dg 30.0\pr & N16\dg 49.6\pr \\
5 & 18 & 88\dg 30.1\pr & N16\dg 48.9\pr & 5 & 19 & 103\dg 30.1\pr & N16\dg 48.2\pr & 5 & 20 & 118\dg 30.2\pr & N16\dg 47.5\pr \\
5 & 21 & 133\dg 30.3\pr & N16\dg 46.9\pr & 5 & 22 & 148\dg 30.3\pr & N16\dg 46.2\pr & 5 & 23 & 163\dg 30.4\pr & N16\dg 45.5\pr \\
6 & 0 & 178\dg 30.5\pr & N16\dg 44.8\pr & 6 & 1 & 193\dg 30.5\pr & N16\dg 44.1\pr & 6 & 2 & 208\dg 30.6\pr & N16\dg 43.4\pr \\
6 & 3 & 223\dg 30.7\pr & N16\dg 42.7\pr & 6 & 4 & 238\dg 30.7\pr & N16\dg 42.1\pr & 6 & 5 & 253\dg 30.8\pr & N16\dg 41.4\pr \\
6 & 6 & 268\dg 30.9\pr & N16\dg 40.7\pr & 6 & 7 & 283\dg 31.0\pr & N16\dg 40.0\pr & 6 & 8 & 298\dg 31.0\pr & N16\dg 39.3\pr \\
6 & 9 & 313\dg 31.1\pr & N16\dg 38.6\pr & 6 & 10 & 328\dg 31.2\pr & N16\dg 37.9\pr & 6 & 11 & 343\dg 31.2\pr & N16\dg 37.2\pr \\
6 & 12 & 358\dg 31.3\pr & N16\dg 36.5\pr & 6 & 13 & 13\dg 31.4\pr & N16\dg 35.8\pr & 6 & 14 & 28\dg 31.5\pr & N16\dg 35.2\pr \\
6 & 15 & 43\dg 31.5\pr & N16\dg 34.5\pr & 6 & 16 & 58\dg 31.6\pr & N16\dg 33.8\pr & 6 & 17 & 73\dg 31.7\pr & N16\dg 33.1\pr \\
6 & 18 & 88\dg 31.7\pr & N16\dg 32.4\pr & 6 & 19 & 103\dg 31.8\pr & N16\dg 31.7\pr & 6 & 20 & 118\dg 31.9\pr & N16\dg 31.0\pr \\
6 & 21 & 133\dg 32.0\pr & N16\dg 30.3\pr & 6 & 22 & 148\dg 32.0\pr & N16\dg 29.6\pr & 6 & 23 & 163\dg 32.1\pr & N16\dg 28.9\pr \\
7 & 0 & 178\dg 32.2\pr & N16\dg 28.2\pr & 7 & 1 & 193\dg 32.3\pr & N16\dg 27.5\pr & 7 & 2 & 208\dg 32.3\pr & N16\dg 26.8\pr \\
7 & 3 & 223\dg 32.4\pr & N16\dg 26.1\pr & 7 & 4 & 238\dg 32.5\pr & N16\dg 25.4\pr & 7 & 5 & 253\dg 32.6\pr & N16\dg 24.7\pr \\
7 & 6 & 268\dg 32.6\pr & N16\dg 24.0\pr & 7 & 7 & 283\dg 32.7\pr & N16\dg 23.3\pr & 7 & 8 & 298\dg 32.8\pr & N16\dg 22.6\pr \\
7 & 9 & 313\dg 32.9\pr & N16\dg 21.9\pr & 7 & 10 & 328\dg 32.9\pr & N16\dg 21.2\pr & 7 & 11 & 343\dg 33.0\pr & N16\dg 20.5\pr \\
7 & 12 & 358\dg 33.1\pr & N16\dg 19.8\pr & 7 & 13 & 13\dg 33.2\pr & N16\dg 19.1\pr & 7 & 14 & 28\dg 33.3\pr & N16\dg 18.4\pr \\
7 & 15 & 43\dg 33.3\pr & N16\dg 17.7\pr & 7 & 16 & 58\dg 33.4\pr & N16\dg 17.0\pr & 7 & 17 & 73\dg 33.5\pr & N16\dg 16.3\pr \\
7 & 18 & 88\dg 33.6\pr & N16\dg 15.6\pr & 7 & 19 & 103\dg 33.6\pr & N16\dg 14.9\pr & 7 & 20 & 118\dg 33.7\pr & N16\dg 14.2\pr \\
7 & 21 & 133\dg 33.8\pr & N16\dg 13.5\pr & 7 & 22 & 148\dg 33.9\pr & N16\dg 12.8\pr & 7 & 23 & 163\dg 34.0\pr & N16\dg 12.1\pr \\
8 & 0 & 178\dg 34.0\pr & N16\dg 11.4\pr & 8 & 1 & 193\dg 34.1\pr & N16\dg 10.7\pr & 8 & 2 & 208\dg 34.2\pr & N16\dg 09.9\pr \\
8 & 3 & 223\dg 34.3\pr & N16\dg 09.2\pr & 8 & 4 & 238\dg 34.4\pr & N16\dg 08.5\pr & 8 & 5 & 253\dg 34.5\pr & N16\dg 07.8\pr \\
8 & 6 & 268\dg 34.5\pr & N16\dg 07.1\pr & 8 & 7 & 283\dg 34.6\pr & N16\dg 06.4\pr & 8 & 8 & 298\dg 34.7\pr & N16\dg 05.7\pr \\
8 & 9 & 313\dg 34.8\pr & N16\dg 05.0\pr & 8 & 10 & 328\dg 34.9\pr & N16\dg 04.3\pr & 8 & 11 & 343\dg 34.9\pr & N16\dg 03.6\pr \\
8 & 12 & 358\dg 35.0\pr & N16\dg 02.8\pr & 8 & 13 & 13\dg 35.1\pr & N16\dg 02.1\pr & 8 & 14 & 28\dg 35.2\pr & N16\dg 01.4\pr \\
8 & 15 & 43\dg 35.3\pr & N16\dg 00.7\pr & 8 & 16 & 58\dg 35.4\pr & N16\dg 00.0\pr & 8 & 17 & 73\dg 35.5\pr & N15\dg 59.3\pr \\
8 & 18 & 88\dg 35.5\pr & N15\dg 58.6\pr & 8 & 19 & 103\dg 35.6\pr & N15\dg 57.8\pr & 8 & 20 & 118\dg 35.7\pr & N15\dg 57.1\pr \\
8 & 21 & 133\dg 35.8\pr & N15\dg 56.4\pr & 8 & 22 & 148\dg 35.9\pr & N15\dg 55.7\pr & 8 & 23 & 163\dg 36.0\pr & N15\dg 55.0\pr \\
9 & 0 & 178\dg 36.1\pr & N15\dg 54.3\pr & 9 & 1 & 193\dg 36.1\pr & N15\dg 53.5\pr & 9 & 2 & 208\dg 36.2\pr & N15\dg 52.8\pr \\
9 & 3 & 223\dg 36.3\pr & N15\dg 52.1\pr & 9 & 4 & 238\dg 36.4\pr & N15\dg 51.4\pr & 9 & 5 & 253\dg 36.5\pr & N15\dg 50.7\pr \\
9 & 6 & 268\dg 36.6\pr & N15\dg 49.9\pr & 9 & 7 & 283\dg 36.7\pr & N15\dg 49.2\pr & 9 & 8 & 298\dg 36.8\pr & N15\dg 48.5\pr \\
9 & 9 & 313\dg 36.8\pr & N15\dg 47.8\pr & 9 & 10 & 328\dg 36.9\pr & N15\dg 47.0\pr & 9 & 11 & 343\dg 37.0\pr & N15\dg 46.3\pr \\
9 & 12 & 358\dg 37.1\pr & N15\dg 45.6\pr & 9 & 13 & 13\dg 37.2\pr & N15\dg 44.9\pr & 9 & 14 & 28\dg 37.3\pr & N15\dg 44.2\pr \\
9 & 15 & 43\dg 37.4\pr & N15\dg 43.4\pr & 9 & 16 & 58\dg 37.5\pr & N15\dg 42.7\pr & 9 & 17 & 73\dg 37.6\pr & N15\dg 42.0\pr \\
9 & 18 & 88\dg 37.7\pr & N15\dg 41.3\pr & 9 & 19 & 103\dg 37.7\pr & N15\dg 40.5\pr & 9 & 20 & 118\dg 37.8\pr & N15\dg 39.8\pr \\
9 & 21 & 133\dg 37.9\pr & N15\dg 39.1\pr & 9 & 22 & 148\dg 38.0\pr & N15\dg 38.3\pr & 9 & 23 & 163\dg 38.1\pr & N15\dg 37.6\pr \\
10 & 0 & 178\dg 38.2\pr & N15\dg 36.9\pr & 10 & 1 & 193\dg 38.3\pr & N15\dg 36.2\pr & 10 & 2 & 208\dg 38.4\pr & N15\dg 35.4\pr \\
10 & 3 & 223\dg 38.5\pr & N15\dg 34.7\pr & 10 & 4 & 238\dg 38.6\pr & N15\dg 34.0\pr & 10 & 5 & 253\dg 38.7\pr & N15\dg 33.2\pr \\
10 & 6 & 268\dg 38.8\pr & N15\dg 32.5\pr & 10 & 7 & 283\dg 38.9\pr & N15\dg 31.8\pr & 10 & 8 & 298\dg 39.0\pr & N15\dg 31.0\pr \\
10 & 9 & 313\dg 39.1\pr & N15\dg 30.3\pr & 10 & 10 & 328\dg 39.1\pr & N15\dg 29.6\pr & 10 & 11 & 343\dg 39.2\pr & N15\dg 28.8\pr \\
10 & 12 & 358\dg 39.3\pr & N15\dg 28.1\pr & 10 & 13 & 13\dg 39.4\pr & N15\dg 27.4\pr & 10 & 14 & 28\dg 39.5\pr & N15\dg 26.6\pr \\
10 & 15 & 43\dg 39.6\pr & N15\dg 25.9\pr & 10 & 16 & 58\dg 39.7\pr & N15\dg 25.2\pr & 10 & 17 & 73\dg 39.8\pr & N15\dg 24.4\pr \\
10 & 18 & 88\dg 39.9\pr & N15\dg 23.7\pr & 10 & 19 & 103\dg 40.0\pr & N15\dg 23.0\pr & 10 & 20 & 118\dg 40.1\pr & N15\dg 22.2\pr \\
10 & 21 & 133\dg 40.2\pr & N15\dg 21.5\pr & 10 & 22 & 148\dg 40.3\pr & N15\dg 20.7\pr & 10 & 23 & 163\dg 40.4\pr & N15\dg 20.0\pr \\
11 & 0 & 178\dg 40.5\pr & N15\dg 19.3\pr & 11 & 1 & 193\dg 40.6\pr & N15\dg 18.5\pr & 11 & 2 & 208\dg 40.7\pr & N15\dg 17.8\pr \\
11 & 3 & 223\dg 40.8\pr & N15\dg 17.1\pr & 11 & 4 & 238\dg 40.9\pr & N15\dg 16.3\pr & 11 & 5 & 253\dg 41.0\pr & N15\dg 15.6\pr \\
11 & 6 & 268\dg 41.1\pr & N15\dg 14.8\pr & 11 & 7 & 283\dg 41.2\pr & N15\dg 14.1\pr & 11 & 8 & 298\dg 41.3\pr & N15\dg 13.3\pr \\
11 & 9 & 313\dg 41.4\pr & N15\dg 12.6\pr & 11 & 10 & 328\dg 41.5\pr & N15\dg 11.9\pr & 11 & 11 & 343\dg 41.6\pr & N15\dg 11.1\pr \\
11 & 12 & 358\dg 41.7\pr & N15\dg 10.4\pr & 11 & 13 & 13\dg 41.8\pr & N15\dg 09.6\pr & 11 & 14 & 28\dg 41.9\pr & N15\dg 08.9\pr \\
11 & 15 & 43\dg 42.0\pr & N15\dg 08.1\pr & 11 & 16 & 58\dg 42.1\pr & N15\dg 07.4\pr & 11 & 17 & 73\dg 42.2\pr & N15\dg 06.6\pr \\
11 & 18 & 88\dg 42.3\pr & N15\dg 05.9\pr & 11 & 19 & 103\dg 42.4\pr & N15\dg 05.2\pr & 11 & 20 & 118\dg 42.5\pr & N15\dg 04.4\pr \\
11 & 21 & 133\dg 42.6\pr & N15\dg 03.7\pr & 11 & 22 & 148\dg 42.7\pr & N15\dg 02.9\pr & 11 & 23 & 163\dg 42.8\pr & N15\dg 02.2\pr \\
12 & 0 & 178\dg 42.9\pr & N15\dg 01.4\pr & 12 & 1 & 193\dg 43.0\pr & N15\dg 00.7\pr & 12 & 2 & 208\dg 43.2\pr & N14\dg 59.9\pr \\
12 & 3 & 223\dg 43.3\pr & N14\dg 59.2\pr & 12 & 4 & 238\dg 43.4\pr & N14\dg 58.4\pr & 12 & 5 & 253\dg 43.5\pr & N14\dg 57.7\pr \\
12 & 6 & 268\dg 43.6\pr & N14\dg 56.9\pr & 12 & 7 & 283\dg 43.7\pr & N14\dg 56.2\pr & 12 & 8 & 298\dg 43.8\pr & N14\dg 55.4\pr \\
12 & 9 & 313\dg 43.9\pr & N14\dg 54.7\pr & 12 & 10 & 328\dg 44.0\pr & N14\dg 53.9\pr & 12 & 11 & 343\dg 44.1\pr & N14\dg 53.1\pr \\
12 & 12 & 358\dg 44.2\pr & N14\dg 52.4\pr & 12 & 13 & 13\dg 44.3\pr & N14\dg 51.6\pr & 12 & 14 & 28\dg 44.4\pr & N14\dg 50.9\pr \\
12 & 15 & 43\dg 44.5\pr & N14\dg 50.1\pr & 12 & 16 & 58\dg 44.6\pr & N14\dg 49.4\pr & 12 & 17 & 73\dg 44.8\pr & N14\dg 48.6\pr \\
12 & 18 & 88\dg 44.9\pr & N14\dg 47.9\pr & 12 & 19 & 103\dg 45.0\pr & N14\dg 47.1\pr & 12 & 20 & 118\dg 45.1\pr & N14\dg 46.3\pr \\
12 & 21 & 133\dg 45.2\pr & N14\dg 45.6\pr & 12 & 22 & 148\dg 45.3\pr & N14\dg 44.8\pr & 12 & 23 & 163\dg 45.4\pr & N14\dg 44.1\pr \\
13 & 0 & 178\dg 45.5\pr & N14\dg 43.3\pr & 13 & 1 & 193\dg 45.6\pr & N14\dg 42.6\pr & 13 & 2 & 208\dg 45.7\pr & N14\dg 41.8\pr \\
13 & 3 & 223\dg 45.9\pr & N14\dg 41.0\pr & 13 & 4 & 238\dg 46.0\pr & N14\dg 40.3\pr & 13 & 5 & 253\dg 46.1\pr & N14\dg 39.5\pr \\
13 & 6 & 268\dg 46.2\pr & N14\dg 38.8\pr & 13 & 7 & 283\dg 46.3\pr & N14\dg 38.0\pr & 13 & 8 & 298\dg 46.4\pr & N14\dg 37.2\pr \\
13 & 9 & 313\dg 46.5\pr & N14\dg 36.5\pr & 13 & 10 & 328\dg 46.6\pr & N14\dg 35.7\pr & 13 & 11 & 343\dg 46.8\pr & N14\dg 34.9\pr \\
13 & 12 & 358\dg 46.9\pr & N14\dg 34.2\pr & 13 & 13 & 13\dg 47.0\pr & N14\dg 33.4\pr & 13 & 14 & 28\dg 47.1\pr & N14\dg 32.6\pr \\
13 & 15 & 43\dg 47.2\pr & N14\dg 31.9\pr & 13 & 16 & 58\dg 47.3\pr & N14\dg 31.1\pr & 13 & 17 & 73\dg 47.4\pr & N14\dg 30.4\pr \\
13 & 18 & 88\dg 47.5\pr & N14\dg 29.6\pr & 13 & 19 & 103\dg 47.7\pr & N14\dg 28.8\pr & 13 & 20 & 118\dg 47.8\pr & N14\dg 28.1\pr \\
13 & 21 & 133\dg 47.9\pr & N14\dg 27.3\pr & 13 & 22 & 148\dg 48.0\pr & N14\dg 26.5\pr & 13 & 23 & 163\dg 48.1\pr & N14\dg 25.7\pr \\
14 & 0 & 178\dg 48.2\pr & N14\dg 25.0\pr & 14 & 1 & 193\dg 48.4\pr & N14\dg 24.2\pr & 14 & 2 & 208\dg 48.5\pr & N14\dg 23.4\pr \\
14 & 3 & 223\dg 48.6\pr & N14\dg 22.7\pr & 14 & 4 & 238\dg 48.7\pr & N14\dg 21.9\pr & 14 & 5 & 253\dg 48.8\pr & N14\dg 21.1\pr \\
14 & 6 & 268\dg 48.9\pr & N14\dg 20.4\pr & 14 & 7 & 283\dg 49.1\pr & N14\dg 19.6\pr & 14 & 8 & 298\dg 49.2\pr & N14\dg 18.8\pr \\
14 & 9 & 313\dg 49.3\pr & N14\dg 18.0\pr & 14 & 10 & 328\dg 49.4\pr & N14\dg 17.3\pr & 14 & 11 & 343\dg 49.5\pr & N14\dg 16.5\pr \\
14 & 12 & 358\dg 49.7\pr & N14\dg 15.7\pr & 14 & 13 & 13\dg 49.8\pr & N14\dg 15.0\pr & 14 & 14 & 28\dg 49.9\pr & N14\dg 14.2\pr \\
14 & 15 & 43\dg 50.0\pr & N14\dg 13.4\pr & 14 & 16 & 58\dg 50.1\pr & N14\dg 12.6\pr & 14 & 17 & 73\dg 50.2\pr & N14\dg 11.9\pr \\
14 & 18 & 88\dg 50.4\pr & N14\dg 11.1\pr & 14 & 19 & 103\dg 50.5\pr & N14\dg 10.3\pr & 14 & 20 & 118\dg 50.6\pr & N14\dg 09.5\pr \\
14 & 21 & 133\dg 50.7\pr & N14\dg 08.7\pr & 14 & 22 & 148\dg 50.9\pr & N14\dg 08.0\pr & 14 & 23 & 163\dg 51.0\pr & N14\dg 07.2\pr \\
15 & 0 & 178\dg 51.1\pr & N14\dg 06.4\pr & 15 & 1 & 193\dg 51.2\pr & N14\dg 05.6\pr & 15 & 2 & 208\dg 51.3\pr & N14\dg 04.9\pr \\
15 & 3 & 223\dg 51.5\pr & N14\dg 04.1\pr & 15 & 4 & 238\dg 51.6\pr & N14\dg 03.3\pr & 15 & 5 & 253\dg 51.7\pr & N14\dg 02.5\pr \\
15 & 6 & 268\dg 51.8\pr & N14\dg 01.7\pr & 15 & 7 & 283\dg 52.0\pr & N14\dg 01.0\pr & 15 & 8 & 298\dg 52.1\pr & N14\dg 00.2\pr \\
15 & 9 & 313\dg 52.2\pr & N13\dg 59.4\pr & 15 & 10 & 328\dg 52.3\pr & N13\dg 58.6\pr & 15 & 11 & 343\dg 52.4\pr & N13\dg 57.8\pr \\
15 & 12 & 358\dg 52.6\pr & N13\dg 57.1\pr & 15 & 13 & 13\dg 52.7\pr & N13\dg 56.3\pr & 15 & 14 & 28\dg 52.8\pr & N13\dg 55.5\pr \\
15 & 15 & 43\dg 52.9\pr & N13\dg 54.7\pr & 15 & 16 & 58\dg 53.1\pr & N13\dg 53.9\pr & 15 & 17 & 73\dg 53.2\pr & N13\dg 53.1\pr \\
15 & 18 & 88\dg 53.3\pr & N13\dg 52.3\pr & 15 & 19 & 103\dg 53.4\pr & N13\dg 51.6\pr & 15 & 20 & 118\dg 53.6\pr & N13\dg 50.8\pr \\
15 & 21 & 133\dg 53.7\pr & N13\dg 50.0\pr & 15 & 22 & 148\dg 53.8\pr & N13\dg 49.2\pr & 15 & 23 & 163\dg 54.0\pr & N13\dg 48.4\pr \\
16 & 0 & 178\dg 54.1\pr & N13\dg 47.6\pr & 16 & 1 & 193\dg 54.2\pr & N13\dg 46.8\pr & 16 & 2 & 208\dg 54.3\pr & N13\dg 46.1\pr \\
16 & 3 & 223\dg 54.5\pr & N13\dg 45.3\pr & 16 & 4 & 238\dg 54.6\pr & N13\dg 44.5\pr & 16 & 5 & 253\dg 54.7\pr & N13\dg 43.7\pr \\
16 & 6 & 268\dg 54.8\pr & N13\dg 42.9\pr & 16 & 7 & 283\dg 55.0\pr & N13\dg 42.1\pr & 16 & 8 & 298\dg 55.1\pr & N13\dg 41.3\pr \\
16 & 9 & 313\dg 55.2\pr & N13\dg 40.5\pr & 16 & 10 & 328\dg 55.4\pr & N13\dg 39.7\pr & 16 & 11 & 343\dg 55.5\pr & N13\dg 38.9\pr \\
16 & 12 & 358\dg 55.6\pr & N13\dg 38.2\pr & 16 & 13 & 13\dg 55.8\pr & N13\dg 37.4\pr & 16 & 14 & 28\dg 55.9\pr & N13\dg 36.6\pr \\
16 & 15 & 43\dg 56.0\pr & N13\dg 35.8\pr & 16 & 16 & 58\dg 56.1\pr & N13\dg 35.0\pr & 16 & 17 & 73\dg 56.3\pr & N13\dg 34.2\pr \\
16 & 18 & 88\dg 56.4\pr & N13\dg 33.4\pr & 16 & 19 & 103\dg 56.5\pr & N13\dg 32.6\pr & 16 & 20 & 118\dg 56.7\pr & N13\dg 31.8\pr \\
16 & 21 & 133\dg 56.8\pr & N13\dg 31.0\pr & 16 & 22 & 148\dg 56.9\pr & N13\dg 30.2\pr & 16 & 23 & 163\dg 57.1\pr & N13\dg 29.4\pr \\
17 & 0 & 178\dg 57.2\pr & N13\dg 28.6\pr & 17 & 1 & 193\dg 57.3\pr & N13\dg 27.8\pr & 17 & 2 & 208\dg 57.5\pr & N13\dg 27.0\pr \\
17 & 3 & 223\dg 57.6\pr & N13\dg 26.2\pr & 17 & 4 & 238\dg 57.7\pr & N13\dg 25.4\pr & 17 & 5 & 253\dg 57.9\pr & N13\dg 24.6\pr \\
17 & 6 & 268\dg 58.0\pr & N13\dg 23.8\pr & 17 & 7 & 283\dg 58.1\pr & N13\dg 23.0\pr & 17 & 8 & 298\dg 58.3\pr & N13\dg 22.2\pr \\
17 & 9 & 313\dg 58.4\pr & N13\dg 21.4\pr & 17 & 10 & 328\dg 58.5\pr & N13\dg 20.6\pr & 17 & 11 & 343\dg 58.7\pr & N13\dg 19.8\pr \\
17 & 12 & 358\dg 58.8\pr & N13\dg 19.0\pr & 17 & 13 & 13\dg 58.9\pr & N13\dg 18.2\pr & 17 & 14 & 28\dg 59.1\pr & N13\dg 17.4\pr \\
17 & 15 & 43\dg 59.2\pr & N13\dg 16.6\pr & 17 & 16 & 58\dg 59.4\pr & N13\dg 15.8\pr & 17 & 17 & 73\dg 59.5\pr & N13\dg 15.0\pr \\
17 & 18 & 88\dg 59.6\pr & N13\dg 14.2\pr & 17 & 19 & 103\dg 59.8\pr & N13\dg 13.4\pr & 17 & 20 & 118\dg 59.9\pr & N13\dg 12.6\pr \\
17 & 21 & 134\dg 00.0\pr & N13\dg 11.8\pr & 17 & 22 & 149\dg 00.2\pr & N13\dg 11.0\pr & 17 & 23 & 164\dg 00.3\pr & N13\dg 10.2\pr \\
18 & 0 & 179\dg 00.4\pr & N13\dg 09.4\pr & 18 & 1 & 194\dg 00.6\pr & N13\dg 08.6\pr & 18 & 2 & 209\dg 00.7\pr & N13\dg 07.8\pr \\
18 & 3 & 224\dg 00.9\pr & N13\dg 07.0\pr & 18 & 4 & 239\dg 01.0\pr & N13\dg 06.2\pr & 18 & 5 & 254\dg 01.1\pr & N13\dg 05.4\pr \\
18 & 6 & 269\dg 01.3\pr & N13\dg 04.6\pr & 18 & 7 & 284\dg 01.4\pr & N13\dg 03.8\pr & 18 & 8 & 299\dg 01.6\pr & N13\dg 03.0\pr \\
18 & 9 & 314\dg 01.7\pr & N13\dg 02.1\pr & 18 & 10 & 329\dg 01.8\pr & N13\dg 01.3\pr & 18 & 11 & 344\dg 02.0\pr & N13\dg 00.5\pr \\
18 & 12 & 359\dg 02.1\pr & N12\dg 59.7\pr & 18 & 13 & 14\dg 02.3\pr & N12\dg 58.9\pr & 18 & 14 & 29\dg 02.4\pr & N12\dg 58.1\pr \\
18 & 15 & 44\dg 02.5\pr & N12\dg 57.3\pr & 18 & 16 & 59\dg 02.7\pr & N12\dg 56.5\pr & 18 & 17 & 74\dg 02.8\pr & N12\dg 55.7\pr \\
18 & 18 & 89\dg 03.0\pr & N12\dg 54.9\pr & 18 & 19 & 104\dg 03.1\pr & N12\dg 54.0\pr & 18 & 20 & 119\dg 03.2\pr & N12\dg 53.2\pr \\
18 & 21 & 134\dg 03.4\pr & N12\dg 52.4\pr & 18 & 22 & 149\dg 03.5\pr & N12\dg 51.6\pr & 18 & 23 & 164\dg 03.7\pr & N12\dg 50.8\pr \\
19 & 0 & 179\dg 03.8\pr & N12\dg 50.0\pr & 19 & 1 & 194\dg 04.0\pr & N12\dg 49.2\pr & 19 & 2 & 209\dg 04.1\pr & N12\dg 48.3\pr \\
19 & 3 & 224\dg 04.2\pr & N12\dg 47.5\pr & 19 & 4 & 239\dg 04.4\pr & N12\dg 46.7\pr & 19 & 5 & 254\dg 04.5\pr & N12\dg 45.9\pr \\
19 & 6 & 269\dg 04.7\pr & N12\dg 45.1\pr & 19 & 7 & 284\dg 04.8\pr & N12\dg 44.3\pr & 19 & 8 & 299\dg 05.0\pr & N12\dg 43.5\pr \\
19 & 9 & 314\dg 05.1\pr & N12\dg 42.6\pr & 19 & 10 & 329\dg 05.3\pr & N12\dg 41.8\pr & 19 & 11 & 344\dg 05.4\pr & N12\dg 41.0\pr \\
19 & 12 & 359\dg 05.6\pr & N12\dg 40.2\pr & 19 & 13 & 14\dg 05.7\pr & N12\dg 39.4\pr & 19 & 14 & 29\dg 05.8\pr & N12\dg 38.6\pr \\
19 & 15 & 44\dg 06.0\pr & N12\dg 37.7\pr & 19 & 16 & 59\dg 06.1\pr & N12\dg 36.9\pr & 19 & 17 & 74\dg 06.3\pr & N12\dg 36.1\pr \\
19 & 18 & 89\dg 06.4\pr & N12\dg 35.3\pr & 19 & 19 & 104\dg 06.6\pr & N12\dg 34.5\pr & 19 & 20 & 119\dg 06.7\pr & N12\dg 33.6\pr \\
19 & 21 & 134\dg 06.9\pr & N12\dg 32.8\pr & 19 & 22 & 149\dg 07.0\pr & N12\dg 32.0\pr & 19 & 23 & 164\dg 07.2\pr & N12\dg 31.2\pr \\
20 & 0 & 179\dg 07.3\pr & N12\dg 30.3\pr & 20 & 1 & 194\dg 07.5\pr & N12\dg 29.5\pr & 20 & 2 & 209\dg 07.6\pr & N12\dg 28.7\pr \\
20 & 3 & 224\dg 07.8\pr & N12\dg 27.9\pr & 20 & 4 & 239\dg 07.9\pr & N12\dg 27.1\pr & 20 & 5 & 254\dg 08.1\pr & N12\dg 26.2\pr \\
20 & 6 & 269\dg 08.2\pr & N12\dg 25.4\pr & 20 & 7 & 284\dg 08.4\pr & N12\dg 24.6\pr & 20 & 8 & 299\dg 08.5\pr & N12\dg 23.8\pr \\
20 & 9 & 314\dg 08.7\pr & N12\dg 22.9\pr & 20 & 10 & 329\dg 08.8\pr & N12\dg 22.1\pr & 20 & 11 & 344\dg 09.0\pr & N12\dg 21.3\pr \\
20 & 12 & 359\dg 09.1\pr & N12\dg 20.5\pr & 20 & 13 & 14\dg 09.3\pr & N12\dg 19.6\pr & 20 & 14 & 29\dg 09.4\pr & N12\dg 18.8\pr \\
20 & 15 & 44\dg 09.6\pr & N12\dg 18.0\pr & 20 & 16 & 59\dg 09.7\pr & N12\dg 17.2\pr & 20 & 17 & 74\dg 09.9\pr & N12\dg 16.3\pr \\
20 & 18 & 89\dg 10.0\pr & N12\dg 15.5\pr & 20 & 19 & 104\dg 10.2\pr & N12\dg 14.7\pr & 20 & 20 & 119\dg 10.3\pr & N12\dg 13.8\pr \\
20 & 21 & 134\dg 10.5\pr & N12\dg 13.0\pr & 20 & 22 & 149\dg 10.6\pr & N12\dg 12.2\pr & 20 & 23 & 164\dg 10.8\pr & N12\dg 11.4\pr \\
21 & 0 & 179\dg 10.9\pr & N12\dg 10.5\pr & 21 & 1 & 194\dg 11.1\pr & N12\dg 09.7\pr & 21 & 2 & 209\dg 11.2\pr & N12\dg 08.9\pr \\
21 & 3 & 224\dg 11.4\pr & N12\dg 08.0\pr & 21 & 4 & 239\dg 11.5\pr & N12\dg 07.2\pr & 21 & 5 & 254\dg 11.7\pr & N12\dg 06.4\pr \\
21 & 6 & 269\dg 11.9\pr & N12\dg 05.5\pr & 21 & 7 & 284\dg 12.0\pr & N12\dg 04.7\pr & 21 & 8 & 299\dg 12.2\pr & N12\dg 03.9\pr \\
21 & 9 & 314\dg 12.3\pr & N12\dg 03.0\pr & 21 & 10 & 329\dg 12.5\pr & N12\dg 02.2\pr & 21 & 11 & 344\dg 12.6\pr & N12\dg 01.4\pr \\
21 & 12 & 359\dg 12.8\pr & N12\dg 00.5\pr & 21 & 13 & 14\dg 12.9\pr & N11\dg 59.7\pr & 21 & 14 & 29\dg 13.1\pr & N11\dg 58.9\pr \\
21 & 15 & 44\dg 13.2\pr & N11\dg 58.0\pr & 21 & 16 & 59\dg 13.4\pr & N11\dg 57.2\pr & 21 & 17 & 74\dg 13.6\pr & N11\dg 56.4\pr \\
21 & 18 & 89\dg 13.7\pr & N11\dg 55.5\pr & 21 & 19 & 104\dg 13.9\pr & N11\dg 54.7\pr & 21 & 20 & 119\dg 14.0\pr & N11\dg 53.8\pr \\
21 & 21 & 134\dg 14.2\pr & N11\dg 53.0\pr & 21 & 22 & 149\dg 14.3\pr & N11\dg 52.2\pr & 21 & 23 & 164\dg 14.5\pr & N11\dg 51.3\pr \\
22 & 0 & 179\dg 14.7\pr & N11\dg 50.5\pr & 22 & 1 & 194\dg 14.8\pr & N11\dg 49.7\pr & 22 & 2 & 209\dg 15.0\pr & N11\dg 48.8\pr \\
22 & 3 & 224\dg 15.1\pr & N11\dg 48.0\pr & 22 & 4 & 239\dg 15.3\pr & N11\dg 47.1\pr & 22 & 5 & 254\dg 15.5\pr & N11\dg 46.3\pr \\
22 & 6 & 269\dg 15.6\pr & N11\dg 45.5\pr & 22 & 7 & 284\dg 15.8\pr & N11\dg 44.6\pr & 22 & 8 & 299\dg 15.9\pr & N11\dg 43.8\pr \\
22 & 9 & 314\dg 16.1\pr & N11\dg 42.9\pr & 22 & 10 & 329\dg 16.2\pr & N11\dg 42.1\pr & 22 & 11 & 344\dg 16.4\pr & N11\dg 41.3\pr \\
22 & 12 & 359\dg 16.6\pr & N11\dg 40.4\pr & 22 & 13 & 14\dg 16.7\pr & N11\dg 39.6\pr & 22 & 14 & 29\dg 16.9\pr & N11\dg 38.7\pr \\
22 & 15 & 44\dg 17.0\pr & N11\dg 37.9\pr & 22 & 16 & 59\dg 17.2\pr & N11\dg 37.0\pr & 22 & 17 & 74\dg 17.4\pr & N11\dg 36.2\pr \\
22 & 18 & 89\dg 17.5\pr & N11\dg 35.4\pr & 22 & 19 & 104\dg 17.7\pr & N11\dg 34.5\pr & 22 & 20 & 119\dg 17.9\pr & N11\dg 33.7\pr \\
22 & 21 & 134\dg 18.0\pr & N11\dg 32.8\pr & 22 & 22 & 149\dg 18.2\pr & N11\dg 32.0\pr & 22 & 23 & 164\dg 18.3\pr & N11\dg 31.1\pr \\
23 & 0 & 179\dg 18.5\pr & N11\dg 30.3\pr & 23 & 1 & 194\dg 18.7\pr & N11\dg 29.4\pr & 23 & 2 & 209\dg 18.8\pr & N11\dg 28.6\pr \\
23 & 3 & 224\dg 19.0\pr & N11\dg 27.8\pr & 23 & 4 & 239\dg 19.2\pr & N11\dg 26.9\pr & 23 & 5 & 254\dg 19.3\pr & N11\dg 26.1\pr \\
23 & 6 & 269\dg 19.5\pr & N11\dg 25.2\pr & 23 & 7 & 284\dg 19.6\pr & N11\dg 24.4\pr & 23 & 8 & 299\dg 19.8\pr & N11\dg 23.5\pr \\
23 & 9 & 314\dg 20.0\pr & N11\dg 22.7\pr & 23 & 10 & 329\dg 20.1\pr & N11\dg 21.8\pr & 23 & 11 & 344\dg 20.3\pr & N11\dg 21.0\pr \\
23 & 12 & 359\dg 20.5\pr & N11\dg 20.1\pr & 23 & 13 & 14\dg 20.6\pr & N11\dg 19.3\pr & 23 & 14 & 29\dg 20.8\pr & N11\dg 18.4\pr \\
23 & 15 & 44\dg 21.0\pr & N11\dg 17.6\pr & 23 & 16 & 59\dg 21.1\pr & N11\dg 16.7\pr & 23 & 17 & 74\dg 21.3\pr & N11\dg 15.9\pr \\
23 & 18 & 89\dg 21.5\pr & N11\dg 15.0\pr & 23 & 19 & 104\dg 21.6\pr & N11\dg 14.2\pr & 23 & 20 & 119\dg 21.8\pr & N11\dg 13.3\pr \\
23 & 21 & 134\dg 22.0\pr & N11\dg 12.5\pr & 23 & 22 & 149\dg 22.1\pr & N11\dg 11.6\pr & 23 & 23 & 164\dg 22.3\pr & N11\dg 10.8\pr \\
24 & 0 & 179\dg 22.5\pr & N11\dg 09.9\pr & 24 & 1 & 194\dg 22.6\pr & N11\dg 09.0\pr & 24 & 2 & 209\dg 22.8\pr & N11\dg 08.2\pr \\
24 & 3 & 224\dg 23.0\pr & N11\dg 07.3\pr & 24 & 4 & 239\dg 23.1\pr & N11\dg 06.5\pr & 24 & 5 & 254\dg 23.3\pr & N11\dg 05.6\pr \\
24 & 6 & 269\dg 23.5\pr & N11\dg 04.8\pr & 24 & 7 & 284\dg 23.6\pr & N11\dg 03.9\pr & 24 & 8 & 299\dg 23.8\pr & N11\dg 03.1\pr \\
24 & 9 & 314\dg 24.0\pr & N11\dg 02.2\pr & 24 & 10 & 329\dg 24.1\pr & N11\dg 01.4\pr & 24 & 11 & 344\dg 24.3\pr & N11\dg 00.5\pr \\
24 & 12 & 359\dg 24.5\pr & N10\dg 59.6\pr & 24 & 13 & 14\dg 24.6\pr & N10\dg 58.8\pr & 24 & 14 & 29\dg 24.8\pr & N10\dg 57.9\pr \\
24 & 15 & 44\dg 25.0\pr & N10\dg 57.1\pr & 24 & 16 & 59\dg 25.2\pr & N10\dg 56.2\pr & 24 & 17 & 74\dg 25.3\pr & N10\dg 55.4\pr \\
24 & 18 & 89\dg 25.5\pr & N10\dg 54.5\pr & 24 & 19 & 104\dg 25.7\pr & N10\dg 53.6\pr & 24 & 20 & 119\dg 25.8\pr & N10\dg 52.8\pr \\
24 & 21 & 134\dg 26.0\pr & N10\dg 51.9\pr & 24 & 22 & 149\dg 26.2\pr & N10\dg 51.1\pr & 24 & 23 & 164\dg 26.3\pr & N10\dg 50.2\pr \\
25 & 0 & 179\dg 26.5\pr & N10\dg 49.3\pr & 25 & 1 & 194\dg 26.7\pr & N10\dg 48.5\pr & 25 & 2 & 209\dg 26.9\pr & N10\dg 47.6\pr \\
25 & 3 & 224\dg 27.0\pr & N10\dg 46.8\pr & 25 & 4 & 239\dg 27.2\pr & N10\dg 45.9\pr & 25 & 5 & 254\dg 27.4\pr & N10\dg 45.0\pr \\
25 & 6 & 269\dg 27.6\pr & N10\dg 44.2\pr & 25 & 7 & 284\dg 27.7\pr & N10\dg 43.3\pr & 25 & 8 & 299\dg 27.9\pr & N10\dg 42.4\pr \\
25 & 9 & 314\dg 28.1\pr & N10\dg 41.6\pr & 25 & 10 & 329\dg 28.2\pr & N10\dg 40.7\pr & 25 & 11 & 344\dg 28.4\pr & N10\dg 39.9\pr \\
25 & 12 & 359\dg 28.6\pr & N10\dg 39.0\pr & 25 & 13 & 14\dg 28.8\pr & N10\dg 38.1\pr & 25 & 14 & 29\dg 28.9\pr & N10\dg 37.3\pr \\
25 & 15 & 44\dg 29.1\pr & N10\dg 36.4\pr & 25 & 16 & 59\dg 29.3\pr & N10\dg 35.5\pr & 25 & 17 & 74\dg 29.5\pr & N10\dg 34.7\pr \\
25 & 18 & 89\dg 29.6\pr & N10\dg 33.8\pr & 25 & 19 & 104\dg 29.8\pr & N10\dg 32.9\pr & 25 & 20 & 119\dg 30.0\pr & N10\dg 32.1\pr \\
25 & 21 & 134\dg 30.2\pr & N10\dg 31.2\pr & 25 & 22 & 149\dg 30.3\pr & N10\dg 30.3\pr & 25 & 23 & 164\dg 30.5\pr & N10\dg 29.5\pr \\
26 & 0 & 179\dg 30.7\pr & N10\dg 28.6\pr & 26 & 1 & 194\dg 30.9\pr & N10\dg 27.7\pr & 26 & 2 & 209\dg 31.0\pr & N10\dg 26.9\pr \\
26 & 3 & 224\dg 31.2\pr & N10\dg 26.0\pr & 26 & 4 & 239\dg 31.4\pr & N10\dg 25.1\pr & 26 & 5 & 254\dg 31.6\pr & N10\dg 24.3\pr \\
26 & 6 & 269\dg 31.7\pr & N10\dg 23.4\pr & 26 & 7 & 284\dg 31.9\pr & N10\dg 22.5\pr & 26 & 8 & 299\dg 32.1\pr & N10\dg 21.7\pr \\
26 & 9 & 314\dg 32.3\pr & N10\dg 20.8\pr & 26 & 10 & 329\dg 32.4\pr & N10\dg 19.9\pr & 26 & 11 & 344\dg 32.6\pr & N10\dg 19.0\pr \\
26 & 12 & 359\dg 32.8\pr & N10\dg 18.2\pr & 26 & 13 & 14\dg 33.0\pr & N10\dg 17.3\pr & 26 & 14 & 29\dg 33.2\pr & N10\dg 16.4\pr \\
26 & 15 & 44\dg 33.3\pr & N10\dg 15.6\pr & 26 & 16 & 59\dg 33.5\pr & N10\dg 14.7\pr & 26 & 17 & 74\dg 33.7\pr & N10\dg 13.8\pr \\
26 & 18 & 89\dg 33.9\pr & N10\dg 12.9\pr & 26 & 19 & 104\dg 34.1\pr & N10\dg 12.1\pr & 26 & 20 & 119\dg 34.2\pr & N10\dg 11.2\pr \\
26 & 21 & 134\dg 34.4\pr & N10\dg 10.3\pr & 26 & 22 & 149\dg 34.6\pr & N10\dg 09.4\pr & 26 & 23 & 164\dg 34.8\pr & N10\dg 08.6\pr \\
27 & 0 & 179\dg 34.9\pr & N10\dg 07.7\pr & 27 & 1 & 194\dg 35.1\pr & N10\dg 06.8\pr & 27 & 2 & 209\dg 35.3\pr & N10\dg 05.9\pr \\
27 & 3 & 224\dg 35.5\pr & N10\dg 05.1\pr & 27 & 4 & 239\dg 35.7\pr & N10\dg 04.2\pr & 27 & 5 & 254\dg 35.8\pr & N10\dg 03.3\pr \\
27 & 6 & 269\dg 36.0\pr & N10\dg 02.4\pr & 27 & 7 & 284\dg 36.2\pr & N10\dg 01.6\pr & 27 & 8 & 299\dg 36.4\pr & N10\dg 00.7\pr \\
27 & 9 & 314\dg 36.6\pr & N9\dg 59.8\pr & 27 & 10 & 329\dg 36.8\pr & N9\dg 58.9\pr & 27 & 11 & 344\dg 36.9\pr & N9\dg 58.1\pr \\
27 & 12 & 359\dg 37.1\pr & N9\dg 57.2\pr & 27 & 13 & 14\dg 37.3\pr & N9\dg 56.3\pr & 27 & 14 & 29\dg 37.5\pr & N9\dg 55.4\pr \\
27 & 15 & 44\dg 37.7\pr & N9\dg 54.6\pr & 27 & 16 & 59\dg 37.8\pr & N9\dg 53.7\pr & 27 & 17 & 74\dg 38.0\pr & N9\dg 52.8\pr \\
27 & 18 & 89\dg 38.2\pr & N9\dg 51.9\pr & 27 & 19 & 104\dg 38.4\pr & N9\dg 51.0\pr & 27 & 20 & 119\dg 38.6\pr & N9\dg 50.2\pr \\
27 & 21 & 134\dg 38.8\pr & N9\dg 49.3\pr & 27 & 22 & 149\dg 38.9\pr & N9\dg 48.4\pr & 27 & 23 & 164\dg 39.1\pr & N9\dg 47.5\pr \\
28 & 0 & 179\dg 39.3\pr & N9\dg 46.6\pr & 28 & 1 & 194\dg 39.5\pr & N9\dg 45.8\pr & 28 & 2 & 209\dg 39.7\pr & N9\dg 44.9\pr \\
28 & 3 & 224\dg 39.9\pr & N9\dg 44.0\pr & 28 & 4 & 239\dg 40.0\pr & N9\dg 43.1\pr & 28 & 5 & 254\dg 40.2\pr & N9\dg 42.2\pr \\
28 & 6 & 269\dg 40.4\pr & N9\dg 41.3\pr & 28 & 7 & 284\dg 40.6\pr & N9\dg 40.5\pr & 28 & 8 & 299\dg 40.8\pr & N9\dg 39.6\pr \\
28 & 9 & 314\dg 41.0\pr & N9\dg 38.7\pr & 28 & 10 & 329\dg 41.2\pr & N9\dg 37.8\pr & 28 & 11 & 344\dg 41.3\pr & N9\dg 36.9\pr \\
28 & 12 & 359\dg 41.5\pr & N9\dg 36.0\pr & 28 & 13 & 14\dg 41.7\pr & N9\dg 35.2\pr & 28 & 14 & 29\dg 41.9\pr & N9\dg 34.3\pr \\
28 & 15 & 44\dg 42.1\pr & N9\dg 33.4\pr & 28 & 16 & 59\dg 42.3\pr & N9\dg 32.5\pr & 28 & 17 & 74\dg 42.5\pr & N9\dg 31.6\pr \\
28 & 18 & 89\dg 42.6\pr & N9\dg 30.7\pr & 28 & 19 & 104\dg 42.8\pr & N9\dg 29.8\pr & 28 & 20 & 119\dg 43.0\pr & N9\dg 29.0\pr \\
28 & 21 & 134\dg 43.2\pr & N9\dg 28.1\pr & 28 & 22 & 149\dg 43.4\pr & N9\dg 27.2\pr & 28 & 23 & 164\dg 43.6\pr & N9\dg 26.3\pr \\
29 & 0 & 179\dg 43.8\pr & N9\dg 25.4\pr & 29 & 1 & 194\dg 43.9\pr & N9\dg 24.5\pr & 29 & 2 & 209\dg 44.1\pr & N9\dg 23.6\pr \\
29 & 3 & 224\dg 44.3\pr & N9\dg 22.8\pr & 29 & 4 & 239\dg 44.5\pr & N9\dg 21.9\pr & 29 & 5 & 254\dg 44.7\pr & N9\dg 21.0\pr \\
29 & 6 & 269\dg 44.9\pr & N9\dg 20.1\pr & 29 & 7 & 284\dg 45.1\pr & N9\dg 19.2\pr & 29 & 8 & 299\dg 45.3\pr & N9\dg 18.3\pr \\
29 & 9 & 314\dg 45.4\pr & N9\dg 17.4\pr & 29 & 10 & 329\dg 45.6\pr & N9\dg 16.5\pr & 29 & 11 & 344\dg 45.8\pr & N9\dg 15.6\pr \\
29 & 12 & 359\dg 46.0\pr & N9\dg 14.7\pr & 29 & 13 & 14\dg 46.2\pr & N9\dg 13.9\pr & 29 & 14 & 29\dg 46.4\pr & N9\dg 13.0\pr \\
29 & 15 & 44\dg 46.6\pr & N9\dg 12.1\pr & 29 & 16 & 59\dg 46.8\pr & N9\dg 11.2\pr & 29 & 17 & 74\dg 47.0\pr & N9\dg 10.3\pr \\
29 & 18 & 89\dg 47.2\pr & N9\dg 09.4\pr & 29 & 19 & 104\dg 47.3\pr & N9\dg 08.5\pr & 29 & 20 & 119\dg 47.5\pr & N9\dg 07.6\pr \\
29 & 21 & 134\dg 47.7\pr & N9\dg 06.7\pr & 29 & 22 & 149\dg 47.9\pr & N9\dg 05.8\pr & 29 & 23 & 164\dg 48.1\pr & N9\dg 04.9\pr \\
30 & 0 & 179\dg 48.3\pr & N9\dg 04.0\pr & 30 & 1 & 194\dg 48.5\pr & N9\dg 03.2\pr & 30 & 2 & 209\dg 48.7\pr & N9\dg 02.3\pr \\
30 & 3 & 224\dg 48.9\pr & N9\dg 01.4\pr & 30 & 4 & 239\dg 49.1\pr & N9\dg 00.5\pr & 30 & 5 & 254\dg 49.3\pr & N8\dg 59.6\pr \\
30 & 6 & 269\dg 49.4\pr & N8\dg 58.7\pr & 30 & 7 & 284\dg 49.6\pr & N8\dg 57.8\pr & 30 & 8 & 299\dg 49.8\pr & N8\dg 56.9\pr \\
30 & 9 & 314\dg 50.0\pr & N8\dg 56.0\pr & 30 & 10 & 329\dg 50.2\pr & N8\dg 55.1\pr & 30 & 11 & 344\dg 50.4\pr & N8\dg 54.2\pr \\
30 & 12 & 359\dg 50.6\pr & N8\dg 53.3\pr & 30 & 13 & 14\dg 50.8\pr & N8\dg 52.4\pr & 30 & 14 & 29\dg 51.0\pr & N8\dg 51.5\pr \\
30 & 15 & 44\dg 51.2\pr & N8\dg 50.6\pr & 30 & 16 & 59\dg 51.4\pr & N8\dg 49.7\pr & 30 & 17 & 74\dg 51.6\pr & N8\dg 48.8\pr \\
30 & 18 & 89\dg 51.8\pr & N8\dg 47.9\pr & 30 & 19 & 104\dg 51.9\pr & N8\dg 47.0\pr & 30 & 20 & 119\dg 52.1\pr & N8\dg 46.1\pr \\
30 & 21 & 134\dg 52.3\pr & N8\dg 45.2\pr & 30 & 22 & 149\dg 52.5\pr & N8\dg 44.3\pr & 30 & 23 & 164\dg 52.7\pr & N8\dg 43.4\pr \\
31 & 0 & 179\dg 52.9\pr & N8\dg 42.5\pr & 31 & 1 & 194\dg 53.1\pr & N8\dg 41.6\pr & 31 & 2 & 209\dg 53.3\pr & N8\dg 40.7\pr \\
31 & 3 & 224\dg 53.5\pr & N8\dg 39.8\pr & 31 & 4 & 239\dg 53.7\pr & N8\dg 38.9\pr & 31 & 5 & 254\dg 53.9\pr & N8\dg 38.0\pr \\
31 & 6 & 269\dg 54.1\pr & N8\dg 37.1\pr & 31 & 7 & 284\dg 54.3\pr & N8\dg 36.2\pr & 31 & 8 & 299\dg 54.5\pr & N8\dg 35.3\pr \\
31 & 9 & 314\dg 54.7\pr & N8\dg 34.4\pr & 31 & 10 & 329\dg 54.9\pr & N8\dg 33.5\pr & 31 & 11 & 344\dg 55.1\pr & N8\dg 32.6\pr \\
31 & 12 & 359\dg 55.3\pr & N8\dg 31.7\pr & 31 & 13 & 14\dg 55.5\pr & N8\dg 30.8\pr & 31 & 14 & 29\dg 55.7\pr & N8\dg 29.9\pr \\
31 & 15 & 44\dg 55.8\pr & N8\dg 29.0\pr & 31 & 16 & 59\dg 56.0\pr & N8\dg 28.1\pr & 31 & 17 & 74\dg 56.2\pr & N8\dg 27.2\pr \\
31 & 18 & 89\dg 56.4\pr & N8\dg 26.3\pr & 31 & 19 & 104\dg 56.6\pr & N8\dg 25.4\pr & 31 & 20 & 119\dg 56.8\pr & N8\dg 24.5\pr \\
31 & 21 & 134\dg 57.0\pr & N8\dg 23.6\pr & 31 & 22 & 149\dg 57.2\pr & N8\dg 22.7\pr & 31 & 23 & 164\dg 57.4\pr & N8\dg 21.8\pr \\
\end{longtable}}

{\footnotesize
\begin{longtable}{rrrr|rrrr|rrrr}
\caption{Sun --- GHA and Declination, September 2026 (UT)}\label{sun09}\\
\toprule
\textbf{d} & \textbf{h} & \textbf{GHA} & \textbf{Dec} & \textbf{d} & \textbf{h} & \textbf{GHA} & \textbf{Dec} & \textbf{d} & \textbf{h} & \textbf{GHA} & \textbf{Dec} \\
\midrule
\endfirsthead
\multicolumn{12}{c}{\textit{Sun --- GHA and Declination, September 2026 (UT) (continued)}}\\
\toprule
\textbf{d} & \textbf{h} & \textbf{GHA} & \textbf{Dec} & \textbf{d} & \textbf{h} & \textbf{GHA} & \textbf{Dec} & \textbf{d} & \textbf{h} & \textbf{GHA} & \textbf{Dec} \\
\midrule
\endhead
\midrule\multicolumn{12}{r}{\textit{continued on next page}}\\
\endfoot
\bottomrule
\endlastfoot
1 & 0 & 179\dg 57.6\pr & N8\dg 20.9\pr & 1 & 1 & 194\dg 57.8\pr & N8\dg 20.0\pr & 1 & 2 & 209\dg 58.0\pr & N8\dg 19.1\pr \\
1 & 3 & 224\dg 58.2\pr & N8\dg 18.2\pr & 1 & 4 & 239\dg 58.4\pr & N8\dg 17.3\pr & 1 & 5 & 254\dg 58.6\pr & N8\dg 16.4\pr \\
1 & 6 & 269\dg 58.8\pr & N8\dg 15.4\pr & 1 & 7 & 284\dg 59.0\pr & N8\dg 14.5\pr & 1 & 8 & 299\dg 59.2\pr & N8\dg 13.6\pr \\
1 & 9 & 314\dg 59.4\pr & N8\dg 12.7\pr & 1 & 10 & 329\dg 59.6\pr & N8\dg 11.8\pr & 1 & 11 & 344\dg 59.8\pr & N8\dg 10.9\pr \\
1 & 12 & 0\dg 00.0\pr & N8\dg 10.0\pr & 1 & 13 & 15\dg 00.2\pr & N8\dg 09.1\pr & 1 & 14 & 30\dg 00.4\pr & N8\dg 08.2\pr \\
1 & 15 & 45\dg 00.6\pr & N8\dg 07.3\pr & 1 & 16 & 60\dg 00.8\pr & N8\dg 06.4\pr & 1 & 17 & 75\dg 01.0\pr & N8\dg 05.5\pr \\
1 & 18 & 90\dg 01.2\pr & N8\dg 04.6\pr & 1 & 19 & 105\dg 01.4\pr & N8\dg 03.6\pr & 1 & 20 & 120\dg 01.6\pr & N8\dg 02.7\pr \\
1 & 21 & 135\dg 01.8\pr & N8\dg 01.8\pr & 1 & 22 & 150\dg 02.0\pr & N8\dg 00.9\pr & 1 & 23 & 165\dg 02.2\pr & N8\dg 00.0\pr \\
2 & 0 & 180\dg 02.4\pr & N7\dg 59.1\pr & 2 & 1 & 195\dg 02.6\pr & N7\dg 58.2\pr & 2 & 2 & 210\dg 02.8\pr & N7\dg 57.3\pr \\
2 & 3 & 225\dg 03.0\pr & N7\dg 56.4\pr & 2 & 4 & 240\dg 03.2\pr & N7\dg 55.5\pr & 2 & 5 & 255\dg 03.4\pr & N7\dg 54.5\pr \\
2 & 6 & 270\dg 03.6\pr & N7\dg 53.6\pr & 2 & 7 & 285\dg 03.8\pr & N7\dg 52.7\pr & 2 & 8 & 300\dg 04.0\pr & N7\dg 51.8\pr \\
2 & 9 & 315\dg 04.2\pr & N7\dg 50.9\pr & 2 & 10 & 330\dg 04.4\pr & N7\dg 50.0\pr & 2 & 11 & 345\dg 04.6\pr & N7\dg 49.1\pr \\
2 & 12 & 0\dg 04.8\pr & N7\dg 48.2\pr & 2 & 13 & 15\dg 05.0\pr & N7\dg 47.2\pr & 2 & 14 & 30\dg 05.2\pr & N7\dg 46.3\pr \\
2 & 15 & 45\dg 05.4\pr & N7\dg 45.4\pr & 2 & 16 & 60\dg 05.6\pr & N7\dg 44.5\pr & 2 & 17 & 75\dg 05.8\pr & N7\dg 43.6\pr \\
2 & 18 & 90\dg 06.0\pr & N7\dg 42.7\pr & 2 & 19 & 105\dg 06.2\pr & N7\dg 41.8\pr & 2 & 20 & 120\dg 06.4\pr & N7\dg 40.8\pr \\
2 & 21 & 135\dg 06.6\pr & N7\dg 39.9\pr & 2 & 22 & 150\dg 06.8\pr & N7\dg 39.0\pr & 2 & 23 & 165\dg 07.0\pr & N7\dg 38.1\pr \\
3 & 0 & 180\dg 07.2\pr & N7\dg 37.2\pr & 3 & 1 & 195\dg 07.4\pr & N7\dg 36.3\pr & 3 & 2 & 210\dg 07.7\pr & N7\dg 35.3\pr \\
3 & 3 & 225\dg 07.9\pr & N7\dg 34.4\pr & 3 & 4 & 240\dg 08.1\pr & N7\dg 33.5\pr & 3 & 5 & 255\dg 08.3\pr & N7\dg 32.6\pr \\
3 & 6 & 270\dg 08.5\pr & N7\dg 31.7\pr & 3 & 7 & 285\dg 08.7\pr & N7\dg 30.8\pr & 3 & 8 & 300\dg 08.9\pr & N7\dg 29.8\pr \\
3 & 9 & 315\dg 09.1\pr & N7\dg 28.9\pr & 3 & 10 & 330\dg 09.3\pr & N7\dg 28.0\pr & 3 & 11 & 345\dg 09.5\pr & N7\dg 27.1\pr \\
3 & 12 & 0\dg 09.7\pr & N7\dg 26.2\pr & 3 & 13 & 15\dg 09.9\pr & N7\dg 25.3\pr & 3 & 14 & 30\dg 10.1\pr & N7\dg 24.3\pr \\
3 & 15 & 45\dg 10.3\pr & N7\dg 23.4\pr & 3 & 16 & 60\dg 10.5\pr & N7\dg 22.5\pr & 3 & 17 & 75\dg 10.7\pr & N7\dg 21.6\pr \\
3 & 18 & 90\dg 10.9\pr & N7\dg 20.7\pr & 3 & 19 & 105\dg 11.1\pr & N7\dg 19.7\pr & 3 & 20 & 120\dg 11.3\pr & N7\dg 18.8\pr \\
3 & 21 & 135\dg 11.5\pr & N7\dg 17.9\pr & 3 & 22 & 150\dg 11.7\pr & N7\dg 17.0\pr & 3 & 23 & 165\dg 12.0\pr & N7\dg 16.1\pr \\
4 & 0 & 180\dg 12.2\pr & N7\dg 15.1\pr & 4 & 1 & 195\dg 12.4\pr & N7\dg 14.2\pr & 4 & 2 & 210\dg 12.6\pr & N7\dg 13.3\pr \\
4 & 3 & 225\dg 12.8\pr & N7\dg 12.4\pr & 4 & 4 & 240\dg 13.0\pr & N7\dg 11.5\pr & 4 & 5 & 255\dg 13.2\pr & N7\dg 10.5\pr \\
4 & 6 & 270\dg 13.4\pr & N7\dg 09.6\pr & 4 & 7 & 285\dg 13.6\pr & N7\dg 08.7\pr & 4 & 8 & 300\dg 13.8\pr & N7\dg 07.8\pr \\
4 & 9 & 315\dg 14.0\pr & N7\dg 06.8\pr & 4 & 10 & 330\dg 14.2\pr & N7\dg 05.9\pr & 4 & 11 & 345\dg 14.4\pr & N7\dg 05.0\pr \\
4 & 12 & 0\dg 14.6\pr & N7\dg 04.1\pr & 4 & 13 & 15\dg 14.8\pr & N7\dg 03.2\pr & 4 & 14 & 30\dg 15.1\pr & N7\dg 02.2\pr \\
4 & 15 & 45\dg 15.3\pr & N7\dg 01.3\pr & 4 & 16 & 60\dg 15.5\pr & N7\dg 00.4\pr & 4 & 17 & 75\dg 15.7\pr & N6\dg 59.5\pr \\
4 & 18 & 90\dg 15.9\pr & N6\dg 58.5\pr & 4 & 19 & 105\dg 16.1\pr & N6\dg 57.6\pr & 4 & 20 & 120\dg 16.3\pr & N6\dg 56.7\pr \\
4 & 21 & 135\dg 16.5\pr & N6\dg 55.8\pr & 4 & 22 & 150\dg 16.7\pr & N6\dg 54.8\pr & 4 & 23 & 165\dg 16.9\pr & N6\dg 53.9\pr \\
5 & 0 & 180\dg 17.1\pr & N6\dg 53.0\pr & 5 & 1 & 195\dg 17.3\pr & N6\dg 52.1\pr & 5 & 2 & 210\dg 17.6\pr & N6\dg 51.1\pr \\
5 & 3 & 225\dg 17.8\pr & N6\dg 50.2\pr & 5 & 4 & 240\dg 18.0\pr & N6\dg 49.3\pr & 5 & 5 & 255\dg 18.2\pr & N6\dg 48.4\pr \\
5 & 6 & 270\dg 18.4\pr & N6\dg 47.4\pr & 5 & 7 & 285\dg 18.6\pr & N6\dg 46.5\pr & 5 & 8 & 300\dg 18.8\pr & N6\dg 45.6\pr \\
5 & 9 & 315\dg 19.0\pr & N6\dg 44.7\pr & 5 & 10 & 330\dg 19.2\pr & N6\dg 43.7\pr & 5 & 11 & 345\dg 19.4\pr & N6\dg 42.8\pr \\
5 & 12 & 0\dg 19.6\pr & N6\dg 41.9\pr & 5 & 13 & 15\dg 19.9\pr & N6\dg 40.9\pr & 5 & 14 & 30\dg 20.1\pr & N6\dg 40.0\pr \\
5 & 15 & 45\dg 20.3\pr & N6\dg 39.1\pr & 5 & 16 & 60\dg 20.5\pr & N6\dg 38.2\pr & 5 & 17 & 75\dg 20.7\pr & N6\dg 37.2\pr \\
5 & 18 & 90\dg 20.9\pr & N6\dg 36.3\pr & 5 & 19 & 105\dg 21.1\pr & N6\dg 35.4\pr & 5 & 20 & 120\dg 21.3\pr & N6\dg 34.4\pr \\
5 & 21 & 135\dg 21.5\pr & N6\dg 33.5\pr & 5 & 22 & 150\dg 21.8\pr & N6\dg 32.6\pr & 5 & 23 & 165\dg 22.0\pr & N6\dg 31.7\pr \\
6 & 0 & 180\dg 22.2\pr & N6\dg 30.7\pr & 6 & 1 & 195\dg 22.4\pr & N6\dg 29.8\pr & 6 & 2 & 210\dg 22.6\pr & N6\dg 28.9\pr \\
6 & 3 & 225\dg 22.8\pr & N6\dg 27.9\pr & 6 & 4 & 240\dg 23.0\pr & N6\dg 27.0\pr & 6 & 5 & 255\dg 23.2\pr & N6\dg 26.1\pr \\
6 & 6 & 270\dg 23.4\pr & N6\dg 25.1\pr & 6 & 7 & 285\dg 23.7\pr & N6\dg 24.2\pr & 6 & 8 & 300\dg 23.9\pr & N6\dg 23.3\pr \\
6 & 9 & 315\dg 24.1\pr & N6\dg 22.3\pr & 6 & 10 & 330\dg 24.3\pr & N6\dg 21.4\pr & 6 & 11 & 345\dg 24.5\pr & N6\dg 20.5\pr \\
6 & 12 & 0\dg 24.7\pr & N6\dg 19.5\pr & 6 & 13 & 15\dg 24.9\pr & N6\dg 18.6\pr & 6 & 14 & 30\dg 25.1\pr & N6\dg 17.7\pr \\
6 & 15 & 45\dg 25.3\pr & N6\dg 16.7\pr & 6 & 16 & 60\dg 25.6\pr & N6\dg 15.8\pr & 6 & 17 & 75\dg 25.8\pr & N6\dg 14.9\pr \\
6 & 18 & 90\dg 26.0\pr & N6\dg 13.9\pr & 6 & 19 & 105\dg 26.2\pr & N6\dg 13.0\pr & 6 & 20 & 120\dg 26.4\pr & N6\dg 12.1\pr \\
6 & 21 & 135\dg 26.6\pr & N6\dg 11.1\pr & 6 & 22 & 150\dg 26.8\pr & N6\dg 10.2\pr & 6 & 23 & 165\dg 27.1\pr & N6\dg 09.3\pr \\
7 & 0 & 180\dg 27.3\pr & N6\dg 08.3\pr & 7 & 1 & 195\dg 27.5\pr & N6\dg 07.4\pr & 7 & 2 & 210\dg 27.7\pr & N6\dg 06.5\pr \\
7 & 3 & 225\dg 27.9\pr & N6\dg 05.5\pr & 7 & 4 & 240\dg 28.1\pr & N6\dg 04.6\pr & 7 & 5 & 255\dg 28.3\pr & N6\dg 03.7\pr \\
7 & 6 & 270\dg 28.5\pr & N6\dg 02.7\pr & 7 & 7 & 285\dg 28.8\pr & N6\dg 01.8\pr & 7 & 8 & 300\dg 29.0\pr & N6\dg 00.9\pr \\
7 & 9 & 315\dg 29.2\pr & N5\dg 59.9\pr & 7 & 10 & 330\dg 29.4\pr & N5\dg 59.0\pr & 7 & 11 & 345\dg 29.6\pr & N5\dg 58.1\pr \\
7 & 12 & 0\dg 29.8\pr & N5\dg 57.1\pr & 7 & 13 & 15\dg 30.0\pr & N5\dg 56.2\pr & 7 & 14 & 30\dg 30.3\pr & N5\dg 55.2\pr \\
7 & 15 & 45\dg 30.5\pr & N5\dg 54.3\pr & 7 & 16 & 60\dg 30.7\pr & N5\dg 53.4\pr & 7 & 17 & 75\dg 30.9\pr & N5\dg 52.4\pr \\
7 & 18 & 90\dg 31.1\pr & N5\dg 51.5\pr & 7 & 19 & 105\dg 31.3\pr & N5\dg 50.6\pr & 7 & 20 & 120\dg 31.5\pr & N5\dg 49.6\pr \\
7 & 21 & 135\dg 31.8\pr & N5\dg 48.7\pr & 7 & 22 & 150\dg 32.0\pr & N5\dg 47.7\pr & 7 & 23 & 165\dg 32.2\pr & N5\dg 46.8\pr \\
8 & 0 & 180\dg 32.4\pr & N5\dg 45.9\pr & 8 & 1 & 195\dg 32.6\pr & N5\dg 44.9\pr & 8 & 2 & 210\dg 32.8\pr & N5\dg 44.0\pr \\
8 & 3 & 225\dg 33.0\pr & N5\dg 43.1\pr & 8 & 4 & 240\dg 33.3\pr & N5\dg 42.1\pr & 8 & 5 & 255\dg 33.5\pr & N5\dg 41.2\pr \\
8 & 6 & 270\dg 33.7\pr & N5\dg 40.2\pr & 8 & 7 & 285\dg 33.9\pr & N5\dg 39.3\pr & 8 & 8 & 300\dg 34.1\pr & N5\dg 38.4\pr \\
8 & 9 & 315\dg 34.3\pr & N5\dg 37.4\pr & 8 & 10 & 330\dg 34.6\pr & N5\dg 36.5\pr & 8 & 11 & 345\dg 34.8\pr & N5\dg 35.5\pr \\
8 & 12 & 0\dg 35.0\pr & N5\dg 34.6\pr & 8 & 13 & 15\dg 35.2\pr & N5\dg 33.7\pr & 8 & 14 & 30\dg 35.4\pr & N5\dg 32.7\pr \\
8 & 15 & 45\dg 35.6\pr & N5\dg 31.8\pr & 8 & 16 & 60\dg 35.9\pr & N5\dg 30.8\pr & 8 & 17 & 75\dg 36.1\pr & N5\dg 29.9\pr \\
8 & 18 & 90\dg 36.3\pr & N5\dg 29.0\pr & 8 & 19 & 105\dg 36.5\pr & N5\dg 28.0\pr & 8 & 20 & 120\dg 36.7\pr & N5\dg 27.1\pr \\
8 & 21 & 135\dg 36.9\pr & N5\dg 26.1\pr & 8 & 22 & 150\dg 37.1\pr & N5\dg 25.2\pr & 8 & 23 & 165\dg 37.4\pr & N5\dg 24.2\pr \\
9 & 0 & 180\dg 37.6\pr & N5\dg 23.3\pr & 9 & 1 & 195\dg 37.8\pr & N5\dg 22.4\pr & 9 & 2 & 210\dg 38.0\pr & N5\dg 21.4\pr \\
9 & 3 & 225\dg 38.2\pr & N5\dg 20.5\pr & 9 & 4 & 240\dg 38.5\pr & N5\dg 19.5\pr & 9 & 5 & 255\dg 38.7\pr & N5\dg 18.6\pr \\
9 & 6 & 270\dg 38.9\pr & N5\dg 17.6\pr & 9 & 7 & 285\dg 39.1\pr & N5\dg 16.7\pr & 9 & 8 & 300\dg 39.3\pr & N5\dg 15.8\pr \\
9 & 9 & 315\dg 39.5\pr & N5\dg 14.8\pr & 9 & 10 & 330\dg 39.8\pr & N5\dg 13.9\pr & 9 & 11 & 345\dg 40.0\pr & N5\dg 12.9\pr \\
9 & 12 & 0\dg 40.2\pr & N5\dg 12.0\pr & 9 & 13 & 15\dg 40.4\pr & N5\dg 11.0\pr & 9 & 14 & 30\dg 40.6\pr & N5\dg 10.1\pr \\
9 & 15 & 45\dg 40.8\pr & N5\dg 09.1\pr & 9 & 16 & 60\dg 41.1\pr & N5\dg 08.2\pr & 9 & 17 & 75\dg 41.3\pr & N5\dg 07.3\pr \\
9 & 18 & 90\dg 41.5\pr & N5\dg 06.3\pr & 9 & 19 & 105\dg 41.7\pr & N5\dg 05.4\pr & 9 & 20 & 120\dg 41.9\pr & N5\dg 04.4\pr \\
9 & 21 & 135\dg 42.2\pr & N5\dg 03.5\pr & 9 & 22 & 150\dg 42.4\pr & N5\dg 02.5\pr & 9 & 23 & 165\dg 42.6\pr & N5\dg 01.6\pr \\
10 & 0 & 180\dg 42.8\pr & N5\dg 00.6\pr & 10 & 1 & 195\dg 43.0\pr & N4\dg 59.7\pr & 10 & 2 & 210\dg 43.2\pr & N4\dg 58.7\pr \\
10 & 3 & 225\dg 43.5\pr & N4\dg 57.8\pr & 10 & 4 & 240\dg 43.7\pr & N4\dg 56.9\pr & 10 & 5 & 255\dg 43.9\pr & N4\dg 55.9\pr \\
10 & 6 & 270\dg 44.1\pr & N4\dg 55.0\pr & 10 & 7 & 285\dg 44.3\pr & N4\dg 54.0\pr & 10 & 8 & 300\dg 44.6\pr & N4\dg 53.1\pr \\
10 & 9 & 315\dg 44.8\pr & N4\dg 52.1\pr & 10 & 10 & 330\dg 45.0\pr & N4\dg 51.2\pr & 10 & 11 & 345\dg 45.2\pr & N4\dg 50.2\pr \\
10 & 12 & 0\dg 45.4\pr & N4\dg 49.3\pr & 10 & 13 & 15\dg 45.6\pr & N4\dg 48.3\pr & 10 & 14 & 30\dg 45.9\pr & N4\dg 47.4\pr \\
10 & 15 & 45\dg 46.1\pr & N4\dg 46.4\pr & 10 & 16 & 60\dg 46.3\pr & N4\dg 45.5\pr & 10 & 17 & 75\dg 46.5\pr & N4\dg 44.5\pr \\
10 & 18 & 90\dg 46.7\pr & N4\dg 43.6\pr & 10 & 19 & 105\dg 47.0\pr & N4\dg 42.6\pr & 10 & 20 & 120\dg 47.2\pr & N4\dg 41.7\pr \\
10 & 21 & 135\dg 47.4\pr & N4\dg 40.7\pr & 10 & 22 & 150\dg 47.6\pr & N4\dg 39.8\pr & 10 & 23 & 165\dg 47.8\pr & N4\dg 38.8\pr \\
11 & 0 & 180\dg 48.1\pr & N4\dg 37.9\pr & 11 & 1 & 195\dg 48.3\pr & N4\dg 36.9\pr & 11 & 2 & 210\dg 48.5\pr & N4\dg 36.0\pr \\
11 & 3 & 225\dg 48.7\pr & N4\dg 35.0\pr & 11 & 4 & 240\dg 48.9\pr & N4\dg 34.1\pr & 11 & 5 & 255\dg 49.2\pr & N4\dg 33.1\pr \\
11 & 6 & 270\dg 49.4\pr & N4\dg 32.2\pr & 11 & 7 & 285\dg 49.6\pr & N4\dg 31.2\pr & 11 & 8 & 300\dg 49.8\pr & N4\dg 30.3\pr \\
11 & 9 & 315\dg 50.0\pr & N4\dg 29.3\pr & 11 & 10 & 330\dg 50.3\pr & N4\dg 28.4\pr & 11 & 11 & 345\dg 50.5\pr & N4\dg 27.4\pr \\
11 & 12 & 0\dg 50.7\pr & N4\dg 26.5\pr & 11 & 13 & 15\dg 50.9\pr & N4\dg 25.5\pr & 11 & 14 & 30\dg 51.1\pr & N4\dg 24.6\pr \\
11 & 15 & 45\dg 51.4\pr & N4\dg 23.6\pr & 11 & 16 & 60\dg 51.6\pr & N4\dg 22.7\pr & 11 & 17 & 75\dg 51.8\pr & N4\dg 21.7\pr \\
11 & 18 & 90\dg 52.0\pr & N4\dg 20.8\pr & 11 & 19 & 105\dg 52.2\pr & N4\dg 19.8\pr & 11 & 20 & 120\dg 52.5\pr & N4\dg 18.9\pr \\
11 & 21 & 135\dg 52.7\pr & N4\dg 17.9\pr & 11 & 22 & 150\dg 52.9\pr & N4\dg 17.0\pr & 11 & 23 & 165\dg 53.1\pr & N4\dg 16.0\pr \\
12 & 0 & 180\dg 53.3\pr & N4\dg 15.1\pr & 12 & 1 & 195\dg 53.6\pr & N4\dg 14.1\pr & 12 & 2 & 210\dg 53.8\pr & N4\dg 13.2\pr \\
12 & 3 & 225\dg 54.0\pr & N4\dg 12.2\pr & 12 & 4 & 240\dg 54.2\pr & N4\dg 11.2\pr & 12 & 5 & 255\dg 54.5\pr & N4\dg 10.3\pr \\
12 & 6 & 270\dg 54.7\pr & N4\dg 09.3\pr & 12 & 7 & 285\dg 54.9\pr & N4\dg 08.4\pr & 12 & 8 & 300\dg 55.1\pr & N4\dg 07.4\pr \\
12 & 9 & 315\dg 55.3\pr & N4\dg 06.5\pr & 12 & 10 & 330\dg 55.6\pr & N4\dg 05.5\pr & 12 & 11 & 345\dg 55.8\pr & N4\dg 04.6\pr \\
12 & 12 & 0\dg 56.0\pr & N4\dg 03.6\pr & 12 & 13 & 15\dg 56.2\pr & N4\dg 02.7\pr & 12 & 14 & 30\dg 56.4\pr & N4\dg 01.7\pr \\
12 & 15 & 45\dg 56.7\pr & N4\dg 00.8\pr & 12 & 16 & 60\dg 56.9\pr & N3\dg 59.8\pr & 12 & 17 & 75\dg 57.1\pr & N3\dg 58.8\pr \\
12 & 18 & 90\dg 57.3\pr & N3\dg 57.9\pr & 12 & 19 & 105\dg 57.6\pr & N3\dg 56.9\pr & 12 & 20 & 120\dg 57.8\pr & N3\dg 56.0\pr \\
12 & 21 & 135\dg 58.0\pr & N3\dg 55.0\pr & 12 & 22 & 150\dg 58.2\pr & N3\dg 54.1\pr & 12 & 23 & 165\dg 58.4\pr & N3\dg 53.1\pr \\
13 & 0 & 180\dg 58.7\pr & N3\dg 52.2\pr & 13 & 1 & 195\dg 58.9\pr & N3\dg 51.2\pr & 13 & 2 & 210\dg 59.1\pr & N3\dg 50.2\pr \\
13 & 3 & 225\dg 59.3\pr & N3\dg 49.3\pr & 13 & 4 & 240\dg 59.5\pr & N3\dg 48.3\pr & 13 & 5 & 255\dg 59.8\pr & N3\dg 47.4\pr \\
13 & 6 & 271\dg 00.0\pr & N3\dg 46.4\pr & 13 & 7 & 286\dg 00.2\pr & N3\dg 45.5\pr & 13 & 8 & 301\dg 00.4\pr & N3\dg 44.5\pr \\
13 & 9 & 316\dg 00.7\pr & N3\dg 43.6\pr & 13 & 10 & 331\dg 00.9\pr & N3\dg 42.6\pr & 13 & 11 & 346\dg 01.1\pr & N3\dg 41.6\pr \\
13 & 12 & 1\dg 01.3\pr & N3\dg 40.7\pr & 13 & 13 & 16\dg 01.5\pr & N3\dg 39.7\pr & 13 & 14 & 31\dg 01.8\pr & N3\dg 38.8\pr \\
13 & 15 & 46\dg 02.0\pr & N3\dg 37.8\pr & 13 & 16 & 61\dg 02.2\pr & N3\dg 36.9\pr & 13 & 17 & 76\dg 02.4\pr & N3\dg 35.9\pr \\
13 & 18 & 91\dg 02.7\pr & N3\dg 34.9\pr & 13 & 19 & 106\dg 02.9\pr & N3\dg 34.0\pr & 13 & 20 & 121\dg 03.1\pr & N3\dg 33.0\pr \\
13 & 21 & 136\dg 03.3\pr & N3\dg 32.1\pr & 13 & 22 & 151\dg 03.5\pr & N3\dg 31.1\pr & 13 & 23 & 166\dg 03.8\pr & N3\dg 30.1\pr \\
14 & 0 & 181\dg 04.0\pr & N3\dg 29.2\pr & 14 & 1 & 196\dg 04.2\pr & N3\dg 28.2\pr & 14 & 2 & 211\dg 04.4\pr & N3\dg 27.3\pr \\
14 & 3 & 226\dg 04.7\pr & N3\dg 26.3\pr & 14 & 4 & 241\dg 04.9\pr & N3\dg 25.4\pr & 14 & 5 & 256\dg 05.1\pr & N3\dg 24.4\pr \\
14 & 6 & 271\dg 05.3\pr & N3\dg 23.4\pr & 14 & 7 & 286\dg 05.5\pr & N3\dg 22.5\pr & 14 & 8 & 301\dg 05.8\pr & N3\dg 21.5\pr \\
14 & 9 & 316\dg 06.0\pr & N3\dg 20.6\pr & 14 & 10 & 331\dg 06.2\pr & N3\dg 19.6\pr & 14 & 11 & 346\dg 06.4\pr & N3\dg 18.6\pr \\
14 & 12 & 1\dg 06.7\pr & N3\dg 17.7\pr & 14 & 13 & 16\dg 06.9\pr & N3\dg 16.7\pr & 14 & 14 & 31\dg 07.1\pr & N3\dg 15.8\pr \\
14 & 15 & 46\dg 07.3\pr & N3\dg 14.8\pr & 14 & 16 & 61\dg 07.6\pr & N3\dg 13.8\pr & 14 & 17 & 76\dg 07.8\pr & N3\dg 12.9\pr \\
14 & 18 & 91\dg 08.0\pr & N3\dg 11.9\pr & 14 & 19 & 106\dg 08.2\pr & N3\dg 11.0\pr & 14 & 20 & 121\dg 08.4\pr & N3\dg 10.0\pr \\
14 & 21 & 136\dg 08.7\pr & N3\dg 09.0\pr & 14 & 22 & 151\dg 08.9\pr & N3\dg 08.1\pr & 14 & 23 & 166\dg 09.1\pr & N3\dg 07.1\pr \\
15 & 0 & 181\dg 09.3\pr & N3\dg 06.2\pr & 15 & 1 & 196\dg 09.6\pr & N3\dg 05.2\pr & 15 & 2 & 211\dg 09.8\pr & N3\dg 04.2\pr \\
15 & 3 & 226\dg 10.0\pr & N3\dg 03.3\pr & 15 & 4 & 241\dg 10.2\pr & N3\dg 02.3\pr & 15 & 5 & 256\dg 10.5\pr & N3\dg 01.3\pr \\
15 & 6 & 271\dg 10.7\pr & N3\dg 00.4\pr & 15 & 7 & 286\dg 10.9\pr & N2\dg 59.4\pr & 15 & 8 & 301\dg 11.1\pr & N2\dg 58.5\pr \\
15 & 9 & 316\dg 11.3\pr & N2\dg 57.5\pr & 15 & 10 & 331\dg 11.6\pr & N2\dg 56.5\pr & 15 & 11 & 346\dg 11.8\pr & N2\dg 55.6\pr \\
15 & 12 & 1\dg 12.0\pr & N2\dg 54.6\pr & 15 & 13 & 16\dg 12.2\pr & N2\dg 53.7\pr & 15 & 14 & 31\dg 12.5\pr & N2\dg 52.7\pr \\
15 & 15 & 46\dg 12.7\pr & N2\dg 51.7\pr & 15 & 16 & 61\dg 12.9\pr & N2\dg 50.8\pr & 15 & 17 & 76\dg 13.1\pr & N2\dg 49.8\pr \\
15 & 18 & 91\dg 13.4\pr & N2\dg 48.8\pr & 15 & 19 & 106\dg 13.6\pr & N2\dg 47.9\pr & 15 & 20 & 121\dg 13.8\pr & N2\dg 46.9\pr \\
15 & 21 & 136\dg 14.0\pr & N2\dg 46.0\pr & 15 & 22 & 151\dg 14.2\pr & N2\dg 45.0\pr & 15 & 23 & 166\dg 14.5\pr & N2\dg 44.0\pr \\
16 & 0 & 181\dg 14.7\pr & N2\dg 43.1\pr & 16 & 1 & 196\dg 14.9\pr & N2\dg 42.1\pr & 16 & 2 & 211\dg 15.1\pr & N2\dg 41.1\pr \\
16 & 3 & 226\dg 15.4\pr & N2\dg 40.2\pr & 16 & 4 & 241\dg 15.6\pr & N2\dg 39.2\pr & 16 & 5 & 256\dg 15.8\pr & N2\dg 38.2\pr \\
16 & 6 & 271\dg 16.0\pr & N2\dg 37.3\pr & 16 & 7 & 286\dg 16.3\pr & N2\dg 36.3\pr & 16 & 8 & 301\dg 16.5\pr & N2\dg 35.4\pr \\
16 & 9 & 316\dg 16.7\pr & N2\dg 34.4\pr & 16 & 10 & 331\dg 16.9\pr & N2\dg 33.4\pr & 16 & 11 & 346\dg 17.2\pr & N2\dg 32.5\pr \\
16 & 12 & 1\dg 17.4\pr & N2\dg 31.5\pr & 16 & 13 & 16\dg 17.6\pr & N2\dg 30.5\pr & 16 & 14 & 31\dg 17.8\pr & N2\dg 29.6\pr \\
16 & 15 & 46\dg 18.0\pr & N2\dg 28.6\pr & 16 & 16 & 61\dg 18.3\pr & N2\dg 27.6\pr & 16 & 17 & 76\dg 18.5\pr & N2\dg 26.7\pr \\
16 & 18 & 91\dg 18.7\pr & N2\dg 25.7\pr & 16 & 19 & 106\dg 18.9\pr & N2\dg 24.7\pr & 16 & 20 & 121\dg 19.2\pr & N2\dg 23.8\pr \\
16 & 21 & 136\dg 19.4\pr & N2\dg 22.8\pr & 16 & 22 & 151\dg 19.6\pr & N2\dg 21.9\pr & 16 & 23 & 166\dg 19.8\pr & N2\dg 20.9\pr \\
17 & 0 & 181\dg 20.1\pr & N2\dg 19.9\pr & 17 & 1 & 196\dg 20.3\pr & N2\dg 19.0\pr & 17 & 2 & 211\dg 20.5\pr & N2\dg 18.0\pr \\
17 & 3 & 226\dg 20.7\pr & N2\dg 17.0\pr & 17 & 4 & 241\dg 21.0\pr & N2\dg 16.1\pr & 17 & 5 & 256\dg 21.2\pr & N2\dg 15.1\pr \\
17 & 6 & 271\dg 21.4\pr & N2\dg 14.1\pr & 17 & 7 & 286\dg 21.6\pr & N2\dg 13.2\pr & 17 & 8 & 301\dg 21.8\pr & N2\dg 12.2\pr \\
17 & 9 & 316\dg 22.1\pr & N2\dg 11.2\pr & 17 & 10 & 331\dg 22.3\pr & N2\dg 10.3\pr & 17 & 11 & 346\dg 22.5\pr & N2\dg 09.3\pr \\
17 & 12 & 1\dg 22.7\pr & N2\dg 08.3\pr & 17 & 13 & 16\dg 23.0\pr & N2\dg 07.4\pr & 17 & 14 & 31\dg 23.2\pr & N2\dg 06.4\pr \\
17 & 15 & 46\dg 23.4\pr & N2\dg 05.4\pr & 17 & 16 & 61\dg 23.6\pr & N2\dg 04.5\pr & 17 & 17 & 76\dg 23.9\pr & N2\dg 03.5\pr \\
17 & 18 & 91\dg 24.1\pr & N2\dg 02.5\pr & 17 & 19 & 106\dg 24.3\pr & N2\dg 01.6\pr & 17 & 20 & 121\dg 24.5\pr & N2\dg 00.6\pr \\
17 & 21 & 136\dg 24.8\pr & N1\dg 59.6\pr & 17 & 22 & 151\dg 25.0\pr & N1\dg 58.7\pr & 17 & 23 & 166\dg 25.2\pr & N1\dg 57.7\pr \\
18 & 0 & 181\dg 25.4\pr & N1\dg 56.7\pr & 18 & 1 & 196\dg 25.6\pr & N1\dg 55.8\pr & 18 & 2 & 211\dg 25.9\pr & N1\dg 54.8\pr \\
18 & 3 & 226\dg 26.1\pr & N1\dg 53.8\pr & 18 & 4 & 241\dg 26.3\pr & N1\dg 52.9\pr & 18 & 5 & 256\dg 26.5\pr & N1\dg 51.9\pr \\
18 & 6 & 271\dg 26.8\pr & N1\dg 50.9\pr & 18 & 7 & 286\dg 27.0\pr & N1\dg 50.0\pr & 18 & 8 & 301\dg 27.2\pr & N1\dg 49.0\pr \\
18 & 9 & 316\dg 27.4\pr & N1\dg 48.0\pr & 18 & 10 & 331\dg 27.7\pr & N1\dg 47.1\pr & 18 & 11 & 346\dg 27.9\pr & N1\dg 46.1\pr \\
18 & 12 & 1\dg 28.1\pr & N1\dg 45.1\pr & 18 & 13 & 16\dg 28.3\pr & N1\dg 44.2\pr & 18 & 14 & 31\dg 28.5\pr & N1\dg 43.2\pr \\
18 & 15 & 46\dg 28.8\pr & N1\dg 42.2\pr & 18 & 16 & 61\dg 29.0\pr & N1\dg 41.2\pr & 18 & 17 & 76\dg 29.2\pr & N1\dg 40.3\pr \\
18 & 18 & 91\dg 29.4\pr & N1\dg 39.3\pr & 18 & 19 & 106\dg 29.7\pr & N1\dg 38.3\pr & 18 & 20 & 121\dg 29.9\pr & N1\dg 37.4\pr \\
18 & 21 & 136\dg 30.1\pr & N1\dg 36.4\pr & 18 & 22 & 151\dg 30.3\pr & N1\dg 35.4\pr & 18 & 23 & 166\dg 30.6\pr & N1\dg 34.5\pr \\
19 & 0 & 181\dg 30.8\pr & N1\dg 33.5\pr & 19 & 1 & 196\dg 31.0\pr & N1\dg 32.5\pr & 19 & 2 & 211\dg 31.2\pr & N1\dg 31.6\pr \\
19 & 3 & 226\dg 31.5\pr & N1\dg 30.6\pr & 19 & 4 & 241\dg 31.7\pr & N1\dg 29.6\pr & 19 & 5 & 256\dg 31.9\pr & N1\dg 28.7\pr \\
19 & 6 & 271\dg 32.1\pr & N1\dg 27.7\pr & 19 & 7 & 286\dg 32.3\pr & N1\dg 26.7\pr & 19 & 8 & 301\dg 32.6\pr & N1\dg 25.7\pr \\
19 & 9 & 316\dg 32.8\pr & N1\dg 24.8\pr & 19 & 10 & 331\dg 33.0\pr & N1\dg 23.8\pr & 19 & 11 & 346\dg 33.2\pr & N1\dg 22.8\pr \\
19 & 12 & 1\dg 33.5\pr & N1\dg 21.9\pr & 19 & 13 & 16\dg 33.7\pr & N1\dg 20.9\pr & 19 & 14 & 31\dg 33.9\pr & N1\dg 19.9\pr \\
19 & 15 & 46\dg 34.1\pr & N1\dg 19.0\pr & 19 & 16 & 61\dg 34.3\pr & N1\dg 18.0\pr & 19 & 17 & 76\dg 34.6\pr & N1\dg 17.0\pr \\
19 & 18 & 91\dg 34.8\pr & N1\dg 16.1\pr & 19 & 19 & 106\dg 35.0\pr & N1\dg 15.1\pr & 19 & 20 & 121\dg 35.2\pr & N1\dg 14.1\pr \\
19 & 21 & 136\dg 35.5\pr & N1\dg 13.1\pr & 19 & 22 & 151\dg 35.7\pr & N1\dg 12.2\pr & 19 & 23 & 166\dg 35.9\pr & N1\dg 11.2\pr \\
20 & 0 & 181\dg 36.1\pr & N1\dg 10.2\pr & 20 & 1 & 196\dg 36.4\pr & N1\dg 09.3\pr & 20 & 2 & 211\dg 36.6\pr & N1\dg 08.3\pr \\
20 & 3 & 226\dg 36.8\pr & N1\dg 07.3\pr & 20 & 4 & 241\dg 37.0\pr & N1\dg 06.4\pr & 20 & 5 & 256\dg 37.2\pr & N1\dg 05.4\pr \\
20 & 6 & 271\dg 37.5\pr & N1\dg 04.4\pr & 20 & 7 & 286\dg 37.7\pr & N1\dg 03.4\pr & 20 & 8 & 301\dg 37.9\pr & N1\dg 02.5\pr \\
20 & 9 & 316\dg 38.1\pr & N1\dg 01.5\pr & 20 & 10 & 331\dg 38.4\pr & N1\dg 00.5\pr & 20 & 11 & 346\dg 38.6\pr & N0\dg 59.6\pr \\
20 & 12 & 1\dg 38.8\pr & N0\dg 58.6\pr & 20 & 13 & 16\dg 39.0\pr & N0\dg 57.6\pr & 20 & 14 & 31\dg 39.2\pr & N0\dg 56.6\pr \\
20 & 15 & 46\dg 39.5\pr & N0\dg 55.7\pr & 20 & 16 & 61\dg 39.7\pr & N0\dg 54.7\pr & 20 & 17 & 76\dg 39.9\pr & N0\dg 53.7\pr \\
20 & 18 & 91\dg 40.1\pr & N0\dg 52.8\pr & 20 & 19 & 106\dg 40.4\pr & N0\dg 51.8\pr & 20 & 20 & 121\dg 40.6\pr & N0\dg 50.8\pr \\
20 & 21 & 136\dg 40.8\pr & N0\dg 49.8\pr & 20 & 22 & 151\dg 41.0\pr & N0\dg 48.9\pr & 20 & 23 & 166\dg 41.2\pr & N0\dg 47.9\pr \\
21 & 0 & 181\dg 41.5\pr & N0\dg 46.9\pr & 21 & 1 & 196\dg 41.7\pr & N0\dg 46.0\pr & 21 & 2 & 211\dg 41.9\pr & N0\dg 45.0\pr \\
21 & 3 & 226\dg 42.1\pr & N0\dg 44.0\pr & 21 & 4 & 241\dg 42.4\pr & N0\dg 43.0\pr & 21 & 5 & 256\dg 42.6\pr & N0\dg 42.1\pr \\
21 & 6 & 271\dg 42.8\pr & N0\dg 41.1\pr & 21 & 7 & 286\dg 43.0\pr & N0\dg 40.1\pr & 21 & 8 & 301\dg 43.2\pr & N0\dg 39.2\pr \\
21 & 9 & 316\dg 43.5\pr & N0\dg 38.2\pr & 21 & 10 & 331\dg 43.7\pr & N0\dg 37.2\pr & 21 & 11 & 346\dg 43.9\pr & N0\dg 36.2\pr \\
21 & 12 & 1\dg 44.1\pr & N0\dg 35.3\pr & 21 & 13 & 16\dg 44.3\pr & N0\dg 34.3\pr & 21 & 14 & 31\dg 44.6\pr & N0\dg 33.3\pr \\
21 & 15 & 46\dg 44.8\pr & N0\dg 32.4\pr & 21 & 16 & 61\dg 45.0\pr & N0\dg 31.4\pr & 21 & 17 & 76\dg 45.2\pr & N0\dg 30.4\pr \\
21 & 18 & 91\dg 45.5\pr & N0\dg 29.4\pr & 21 & 19 & 106\dg 45.7\pr & N0\dg 28.5\pr & 21 & 20 & 121\dg 45.9\pr & N0\dg 27.5\pr \\
21 & 21 & 136\dg 46.1\pr & N0\dg 26.5\pr & 21 & 22 & 151\dg 46.3\pr & N0\dg 25.6\pr & 21 & 23 & 166\dg 46.6\pr & N0\dg 24.6\pr \\
22 & 0 & 181\dg 46.8\pr & N0\dg 23.6\pr & 22 & 1 & 196\dg 47.0\pr & N0\dg 22.6\pr & 22 & 2 & 211\dg 47.2\pr & N0\dg 21.7\pr \\
22 & 3 & 226\dg 47.4\pr & N0\dg 20.7\pr & 22 & 4 & 241\dg 47.7\pr & N0\dg 19.7\pr & 22 & 5 & 256\dg 47.9\pr & N0\dg 18.8\pr \\
22 & 6 & 271\dg 48.1\pr & N0\dg 17.8\pr & 22 & 7 & 286\dg 48.3\pr & N0\dg 16.8\pr & 22 & 8 & 301\dg 48.5\pr & N0\dg 15.8\pr \\
22 & 9 & 316\dg 48.8\pr & N0\dg 14.9\pr & 22 & 10 & 331\dg 49.0\pr & N0\dg 13.9\pr & 22 & 11 & 346\dg 49.2\pr & N0\dg 12.9\pr \\
22 & 12 & 1\dg 49.4\pr & N0\dg 11.9\pr & 22 & 13 & 16\dg 49.6\pr & N0\dg 11.0\pr & 22 & 14 & 31\dg 49.9\pr & N0\dg 10.0\pr \\
22 & 15 & 46\dg 50.1\pr & N0\dg 09.0\pr & 22 & 16 & 61\dg 50.3\pr & N0\dg 08.1\pr & 22 & 17 & 76\dg 50.5\pr & N0\dg 07.1\pr \\
22 & 18 & 91\dg 50.7\pr & N0\dg 06.1\pr & 22 & 19 & 106\dg 51.0\pr & N0\dg 05.1\pr & 22 & 20 & 121\dg 51.2\pr & N0\dg 04.2\pr \\
22 & 21 & 136\dg 51.4\pr & N0\dg 03.2\pr & 22 & 22 & 151\dg 51.6\pr & N0\dg 02.2\pr & 22 & 23 & 166\dg 51.8\pr & N0\dg 01.2\pr \\
23 & 0 & 181\dg 52.1\pr & N0\dg 00.3\pr & 23 & 1 & 196\dg 52.3\pr & S0\dg 00.7\pr & 23 & 2 & 211\dg 52.5\pr & S0\dg 01.7\pr \\
23 & 3 & 226\dg 52.7\pr & S0\dg 02.6\pr & 23 & 4 & 241\dg 52.9\pr & S0\dg 03.6\pr & 23 & 5 & 256\dg 53.2\pr & S0\dg 04.6\pr \\
23 & 6 & 271\dg 53.4\pr & S0\dg 05.6\pr & 23 & 7 & 286\dg 53.6\pr & S0\dg 06.5\pr & 23 & 8 & 301\dg 53.8\pr & S0\dg 07.5\pr \\
23 & 9 & 316\dg 54.0\pr & S0\dg 08.5\pr & 23 & 10 & 331\dg 54.3\pr & S0\dg 09.5\pr & 23 & 11 & 346\dg 54.5\pr & S0\dg 10.4\pr \\
23 & 12 & 1\dg 54.7\pr & S0\dg 11.4\pr & 23 & 13 & 16\dg 54.9\pr & S0\dg 12.4\pr & 23 & 14 & 31\dg 55.1\pr & S0\dg 13.3\pr \\
23 & 15 & 46\dg 55.4\pr & S0\dg 14.3\pr & 23 & 16 & 61\dg 55.6\pr & S0\dg 15.3\pr & 23 & 17 & 76\dg 55.8\pr & S0\dg 16.3\pr \\
23 & 18 & 91\dg 56.0\pr & S0\dg 17.2\pr & 23 & 19 & 106\dg 56.2\pr & S0\dg 18.2\pr & 23 & 20 & 121\dg 56.5\pr & S0\dg 19.2\pr \\
23 & 21 & 136\dg 56.7\pr & S0\dg 20.2\pr & 23 & 22 & 151\dg 56.9\pr & S0\dg 21.1\pr & 23 & 23 & 166\dg 57.1\pr & S0\dg 22.1\pr \\
24 & 0 & 181\dg 57.3\pr & S0\dg 23.1\pr & 24 & 1 & 196\dg 57.5\pr & S0\dg 24.1\pr & 24 & 2 & 211\dg 57.8\pr & S0\dg 25.0\pr \\
24 & 3 & 226\dg 58.0\pr & S0\dg 26.0\pr & 24 & 4 & 241\dg 58.2\pr & S0\dg 27.0\pr & 24 & 5 & 256\dg 58.4\pr & S0\dg 27.9\pr \\
24 & 6 & 271\dg 58.6\pr & S0\dg 28.9\pr & 24 & 7 & 286\dg 58.9\pr & S0\dg 29.9\pr & 24 & 8 & 301\dg 59.1\pr & S0\dg 30.9\pr \\
24 & 9 & 316\dg 59.3\pr & S0\dg 31.8\pr & 24 & 10 & 331\dg 59.5\pr & S0\dg 32.8\pr & 24 & 11 & 346\dg 59.7\pr & S0\dg 33.8\pr \\
24 & 12 & 1\dg 59.9\pr & S0\dg 34.8\pr & 24 & 13 & 17\dg 00.2\pr & S0\dg 35.7\pr & 24 & 14 & 32\dg 00.4\pr & S0\dg 36.7\pr \\
24 & 15 & 47\dg 00.6\pr & S0\dg 37.7\pr & 24 & 16 & 62\dg 00.8\pr & S0\dg 38.7\pr & 24 & 17 & 77\dg 01.0\pr & S0\dg 39.6\pr \\
24 & 18 & 92\dg 01.3\pr & S0\dg 40.6\pr & 24 & 19 & 107\dg 01.5\pr & S0\dg 41.6\pr & 24 & 20 & 122\dg 01.7\pr & S0\dg 42.5\pr \\
24 & 21 & 137\dg 01.9\pr & S0\dg 43.5\pr & 24 & 22 & 152\dg 02.1\pr & S0\dg 44.5\pr & 24 & 23 & 167\dg 02.3\pr & S0\dg 45.5\pr \\
25 & 0 & 182\dg 02.6\pr & S0\dg 46.4\pr & 25 & 1 & 197\dg 02.8\pr & S0\dg 47.4\pr & 25 & 2 & 212\dg 03.0\pr & S0\dg 48.4\pr \\
25 & 3 & 227\dg 03.2\pr & S0\dg 49.4\pr & 25 & 4 & 242\dg 03.4\pr & S0\dg 50.3\pr & 25 & 5 & 257\dg 03.6\pr & S0\dg 51.3\pr \\
25 & 6 & 272\dg 03.9\pr & S0\dg 52.3\pr & 25 & 7 & 287\dg 04.1\pr & S0\dg 53.3\pr & 25 & 8 & 302\dg 04.3\pr & S0\dg 54.2\pr \\
25 & 9 & 317\dg 04.5\pr & S0\dg 55.2\pr & 25 & 10 & 332\dg 04.7\pr & S0\dg 56.2\pr & 25 & 11 & 347\dg 04.9\pr & S0\dg 57.2\pr \\
25 & 12 & 2\dg 05.2\pr & S0\dg 58.1\pr & 25 & 13 & 17\dg 05.4\pr & S0\dg 59.1\pr & 25 & 14 & 32\dg 05.6\pr & S1\dg 00.1\pr \\
25 & 15 & 47\dg 05.8\pr & S1\dg 01.0\pr & 25 & 16 & 62\dg 06.0\pr & S1\dg 02.0\pr & 25 & 17 & 77\dg 06.2\pr & S1\dg 03.0\pr \\
25 & 18 & 92\dg 06.4\pr & S1\dg 04.0\pr & 25 & 19 & 107\dg 06.7\pr & S1\dg 04.9\pr & 25 & 20 & 122\dg 06.9\pr & S1\dg 05.9\pr \\
25 & 21 & 137\dg 07.1\pr & S1\dg 06.9\pr & 25 & 22 & 152\dg 07.3\pr & S1\dg 07.9\pr & 25 & 23 & 167\dg 07.5\pr & S1\dg 08.8\pr \\
26 & 0 & 182\dg 07.7\pr & S1\dg 09.8\pr & 26 & 1 & 197\dg 08.0\pr & S1\dg 10.8\pr & 26 & 2 & 212\dg 08.2\pr & S1\dg 11.8\pr \\
26 & 3 & 227\dg 08.4\pr & S1\dg 12.7\pr & 26 & 4 & 242\dg 08.6\pr & S1\dg 13.7\pr & 26 & 5 & 257\dg 08.8\pr & S1\dg 14.7\pr \\
26 & 6 & 272\dg 09.0\pr & S1\dg 15.6\pr & 26 & 7 & 287\dg 09.2\pr & S1\dg 16.6\pr & 26 & 8 & 302\dg 09.5\pr & S1\dg 17.6\pr \\
26 & 9 & 317\dg 09.7\pr & S1\dg 18.6\pr & 26 & 10 & 332\dg 09.9\pr & S1\dg 19.5\pr & 26 & 11 & 347\dg 10.1\pr & S1\dg 20.5\pr \\
26 & 12 & 2\dg 10.3\pr & S1\dg 21.5\pr & 26 & 13 & 17\dg 10.5\pr & S1\dg 22.5\pr & 26 & 14 & 32\dg 10.7\pr & S1\dg 23.4\pr \\
26 & 15 & 47\dg 11.0\pr & S1\dg 24.4\pr & 26 & 16 & 62\dg 11.2\pr & S1\dg 25.4\pr & 26 & 17 & 77\dg 11.4\pr & S1\dg 26.4\pr \\
26 & 18 & 92\dg 11.6\pr & S1\dg 27.3\pr & 26 & 19 & 107\dg 11.8\pr & S1\dg 28.3\pr & 26 & 20 & 122\dg 12.0\pr & S1\dg 29.3\pr \\
26 & 21 & 137\dg 12.2\pr & S1\dg 30.2\pr & 26 & 22 & 152\dg 12.5\pr & S1\dg 31.2\pr & 26 & 23 & 167\dg 12.7\pr & S1\dg 32.2\pr \\
27 & 0 & 182\dg 12.9\pr & S1\dg 33.2\pr & 27 & 1 & 197\dg 13.1\pr & S1\dg 34.1\pr & 27 & 2 & 212\dg 13.3\pr & S1\dg 35.1\pr \\
27 & 3 & 227\dg 13.5\pr & S1\dg 36.1\pr & 27 & 4 & 242\dg 13.7\pr & S1\dg 37.1\pr & 27 & 5 & 257\dg 13.9\pr & S1\dg 38.0\pr \\
27 & 6 & 272\dg 14.2\pr & S1\dg 39.0\pr & 27 & 7 & 287\dg 14.4\pr & S1\dg 40.0\pr & 27 & 8 & 302\dg 14.6\pr & S1\dg 40.9\pr \\
27 & 9 & 317\dg 14.8\pr & S1\dg 41.9\pr & 27 & 10 & 332\dg 15.0\pr & S1\dg 42.9\pr & 27 & 11 & 347\dg 15.2\pr & S1\dg 43.9\pr \\
27 & 12 & 2\dg 15.4\pr & S1\dg 44.8\pr & 27 & 13 & 17\dg 15.6\pr & S1\dg 45.8\pr & 27 & 14 & 32\dg 15.9\pr & S1\dg 46.8\pr \\
27 & 15 & 47\dg 16.1\pr & S1\dg 47.8\pr & 27 & 16 & 62\dg 16.3\pr & S1\dg 48.7\pr & 27 & 17 & 77\dg 16.5\pr & S1\dg 49.7\pr \\
27 & 18 & 92\dg 16.7\pr & S1\dg 50.7\pr & 27 & 19 & 107\dg 16.9\pr & S1\dg 51.6\pr & 27 & 20 & 122\dg 17.1\pr & S1\dg 52.6\pr \\
27 & 21 & 137\dg 17.3\pr & S1\dg 53.6\pr & 27 & 22 & 152\dg 17.5\pr & S1\dg 54.6\pr & 27 & 23 & 167\dg 17.8\pr & S1\dg 55.5\pr \\
28 & 0 & 182\dg 18.0\pr & S1\dg 56.5\pr & 28 & 1 & 197\dg 18.2\pr & S1\dg 57.5\pr & 28 & 2 & 212\dg 18.4\pr & S1\dg 58.5\pr \\
28 & 3 & 227\dg 18.6\pr & S1\dg 59.4\pr & 28 & 4 & 242\dg 18.8\pr & S2\dg 00.4\pr & 28 & 5 & 257\dg 19.0\pr & S2\dg 01.4\pr \\
28 & 6 & 272\dg 19.2\pr & S2\dg 02.3\pr & 28 & 7 & 287\dg 19.4\pr & S2\dg 03.3\pr & 28 & 8 & 302\dg 19.7\pr & S2\dg 04.3\pr \\
28 & 9 & 317\dg 19.9\pr & S2\dg 05.3\pr & 28 & 10 & 332\dg 20.1\pr & S2\dg 06.2\pr & 28 & 11 & 347\dg 20.3\pr & S2\dg 07.2\pr \\
28 & 12 & 2\dg 20.5\pr & S2\dg 08.2\pr & 28 & 13 & 17\dg 20.7\pr & S2\dg 09.2\pr & 28 & 14 & 32\dg 20.9\pr & S2\dg 10.1\pr \\
28 & 15 & 47\dg 21.1\pr & S2\dg 11.1\pr & 28 & 16 & 62\dg 21.3\pr & S2\dg 12.1\pr & 28 & 17 & 77\dg 21.5\pr & S2\dg 13.0\pr \\
28 & 18 & 92\dg 21.7\pr & S2\dg 14.0\pr & 28 & 19 & 107\dg 22.0\pr & S2\dg 15.0\pr & 28 & 20 & 122\dg 22.2\pr & S2\dg 16.0\pr \\
28 & 21 & 137\dg 22.4\pr & S2\dg 16.9\pr & 28 & 22 & 152\dg 22.6\pr & S2\dg 17.9\pr & 28 & 23 & 167\dg 22.8\pr & S2\dg 18.9\pr \\
29 & 0 & 182\dg 23.0\pr & S2\dg 19.8\pr & 29 & 1 & 197\dg 23.2\pr & S2\dg 20.8\pr & 29 & 2 & 212\dg 23.4\pr & S2\dg 21.8\pr \\
29 & 3 & 227\dg 23.6\pr & S2\dg 22.8\pr & 29 & 4 & 242\dg 23.8\pr & S2\dg 23.7\pr & 29 & 5 & 257\dg 24.0\pr & S2\dg 24.7\pr \\
29 & 6 & 272\dg 24.2\pr & S2\dg 25.7\pr & 29 & 7 & 287\dg 24.5\pr & S2\dg 26.6\pr & 29 & 8 & 302\dg 24.7\pr & S2\dg 27.6\pr \\
29 & 9 & 317\dg 24.9\pr & S2\dg 28.6\pr & 29 & 10 & 332\dg 25.1\pr & S2\dg 29.6\pr & 29 & 11 & 347\dg 25.3\pr & S2\dg 30.5\pr \\
29 & 12 & 2\dg 25.5\pr & S2\dg 31.5\pr & 29 & 13 & 17\dg 25.7\pr & S2\dg 32.5\pr & 29 & 14 & 32\dg 25.9\pr & S2\dg 33.4\pr \\
29 & 15 & 47\dg 26.1\pr & S2\dg 34.4\pr & 29 & 16 & 62\dg 26.3\pr & S2\dg 35.4\pr & 29 & 17 & 77\dg 26.5\pr & S2\dg 36.4\pr \\
29 & 18 & 92\dg 26.7\pr & S2\dg 37.3\pr & 29 & 19 & 107\dg 26.9\pr & S2\dg 38.3\pr & 29 & 20 & 122\dg 27.1\pr & S2\dg 39.3\pr \\
29 & 21 & 137\dg 27.4\pr & S2\dg 40.2\pr & 29 & 22 & 152\dg 27.6\pr & S2\dg 41.2\pr & 29 & 23 & 167\dg 27.8\pr & S2\dg 42.2\pr \\
30 & 0 & 182\dg 28.0\pr & S2\dg 43.2\pr & 30 & 1 & 197\dg 28.2\pr & S2\dg 44.1\pr & 30 & 2 & 212\dg 28.4\pr & S2\dg 45.1\pr \\
30 & 3 & 227\dg 28.6\pr & S2\dg 46.1\pr & 30 & 4 & 242\dg 28.8\pr & S2\dg 47.0\pr & 30 & 5 & 257\dg 29.0\pr & S2\dg 48.0\pr \\
30 & 6 & 272\dg 29.2\pr & S2\dg 49.0\pr & 30 & 7 & 287\dg 29.4\pr & S2\dg 50.0\pr & 30 & 8 & 302\dg 29.6\pr & S2\dg 50.9\pr \\
30 & 9 & 317\dg 29.8\pr & S2\dg 51.9\pr & 30 & 10 & 332\dg 30.0\pr & S2\dg 52.9\pr & 30 & 11 & 347\dg 30.2\pr & S2\dg 53.8\pr \\
30 & 12 & 2\dg 30.4\pr & S2\dg 54.8\pr & 30 & 13 & 17\dg 30.6\pr & S2\dg 55.8\pr & 30 & 14 & 32\dg 30.8\pr & S2\dg 56.7\pr \\
30 & 15 & 47\dg 31.0\pr & S2\dg 57.7\pr & 30 & 16 & 62\dg 31.3\pr & S2\dg 58.7\pr & 30 & 17 & 77\dg 31.5\pr & S2\dg 59.7\pr \\
30 & 18 & 92\dg 31.7\pr & S3\dg 00.6\pr & 30 & 19 & 107\dg 31.9\pr & S3\dg 01.6\pr & 30 & 20 & 122\dg 32.1\pr & S3\dg 02.6\pr \\
30 & 21 & 137\dg 32.3\pr & S3\dg 03.5\pr & 30 & 22 & 152\dg 32.5\pr & S3\dg 04.5\pr & 30 & 23 & 167\dg 32.7\pr & S3\dg 05.5\pr \\
\end{longtable}}

{\footnotesize
\begin{longtable}{rrrr|rrrr|rrrr}
\caption{Sun --- GHA and Declination, October 2026 (UT)}\label{sun10}\\
\toprule
\textbf{d} & \textbf{h} & \textbf{GHA} & \textbf{Dec} & \textbf{d} & \textbf{h} & \textbf{GHA} & \textbf{Dec} & \textbf{d} & \textbf{h} & \textbf{GHA} & \textbf{Dec} \\
\midrule
\endfirsthead
\multicolumn{12}{c}{\textit{Sun --- GHA and Declination, October 2026 (UT) (continued)}}\\
\toprule
\textbf{d} & \textbf{h} & \textbf{GHA} & \textbf{Dec} & \textbf{d} & \textbf{h} & \textbf{GHA} & \textbf{Dec} & \textbf{d} & \textbf{h} & \textbf{GHA} & \textbf{Dec} \\
\midrule
\endhead
\midrule\multicolumn{12}{r}{\textit{continued on next page}}\\
\endfoot
\bottomrule
\endlastfoot
1 & 0 & 182\dg 32.9\pr & S3\dg 06.4\pr & 1 & 1 & 197\dg 33.1\pr & S3\dg 07.4\pr & 1 & 2 & 212\dg 33.3\pr & S3\dg 08.4\pr \\
1 & 3 & 227\dg 33.5\pr & S3\dg 09.4\pr & 1 & 4 & 242\dg 33.7\pr & S3\dg 10.3\pr & 1 & 5 & 257\dg 33.9\pr & S3\dg 11.3\pr \\
1 & 6 & 272\dg 34.1\pr & S3\dg 12.3\pr & 1 & 7 & 287\dg 34.3\pr & S3\dg 13.2\pr & 1 & 8 & 302\dg 34.5\pr & S3\dg 14.2\pr \\
1 & 9 & 317\dg 34.7\pr & S3\dg 15.2\pr & 1 & 10 & 332\dg 34.9\pr & S3\dg 16.1\pr & 1 & 11 & 347\dg 35.1\pr & S3\dg 17.1\pr \\
1 & 12 & 2\dg 35.3\pr & S3\dg 18.1\pr & 1 & 13 & 17\dg 35.5\pr & S3\dg 19.0\pr & 1 & 14 & 32\dg 35.7\pr & S3\dg 20.0\pr \\
1 & 15 & 47\dg 35.9\pr & S3\dg 21.0\pr & 1 & 16 & 62\dg 36.1\pr & S3\dg 21.9\pr & 1 & 17 & 77\dg 36.3\pr & S3\dg 22.9\pr \\
1 & 18 & 92\dg 36.5\pr & S3\dg 23.9\pr & 1 & 19 & 107\dg 36.7\pr & S3\dg 24.9\pr & 1 & 20 & 122\dg 36.9\pr & S3\dg 25.8\pr \\
1 & 21 & 137\dg 37.1\pr & S3\dg 26.8\pr & 1 & 22 & 152\dg 37.3\pr & S3\dg 27.8\pr & 1 & 23 & 167\dg 37.5\pr & S3\dg 28.7\pr \\
2 & 0 & 182\dg 37.7\pr & S3\dg 29.7\pr & 2 & 1 & 197\dg 37.9\pr & S3\dg 30.7\pr & 2 & 2 & 212\dg 38.1\pr & S3\dg 31.6\pr \\
2 & 3 & 227\dg 38.3\pr & S3\dg 32.6\pr & 2 & 4 & 242\dg 38.5\pr & S3\dg 33.6\pr & 2 & 5 & 257\dg 38.7\pr & S3\dg 34.5\pr \\
2 & 6 & 272\dg 38.9\pr & S3\dg 35.5\pr & 2 & 7 & 287\dg 39.1\pr & S3\dg 36.5\pr & 2 & 8 & 302\dg 39.3\pr & S3\dg 37.4\pr \\
2 & 9 & 317\dg 39.5\pr & S3\dg 38.4\pr & 2 & 10 & 332\dg 39.7\pr & S3\dg 39.4\pr & 2 & 11 & 347\dg 39.9\pr & S3\dg 40.3\pr \\
2 & 12 & 2\dg 40.1\pr & S3\dg 41.3\pr & 2 & 13 & 17\dg 40.3\pr & S3\dg 42.3\pr & 2 & 14 & 32\dg 40.5\pr & S3\dg 43.2\pr \\
2 & 15 & 47\dg 40.7\pr & S3\dg 44.2\pr & 2 & 16 & 62\dg 40.9\pr & S3\dg 45.2\pr & 2 & 17 & 77\dg 41.1\pr & S3\dg 46.1\pr \\
2 & 18 & 92\dg 41.3\pr & S3\dg 47.1\pr & 2 & 19 & 107\dg 41.5\pr & S3\dg 48.1\pr & 2 & 20 & 122\dg 41.7\pr & S3\dg 49.0\pr \\
2 & 21 & 137\dg 41.9\pr & S3\dg 50.0\pr & 2 & 22 & 152\dg 42.1\pr & S3\dg 51.0\pr & 2 & 23 & 167\dg 42.3\pr & S3\dg 51.9\pr \\
3 & 0 & 182\dg 42.5\pr & S3\dg 52.9\pr & 3 & 1 & 197\dg 42.7\pr & S3\dg 53.9\pr & 3 & 2 & 212\dg 42.9\pr & S3\dg 54.8\pr \\
3 & 3 & 227\dg 43.1\pr & S3\dg 55.8\pr & 3 & 4 & 242\dg 43.3\pr & S3\dg 56.8\pr & 3 & 5 & 257\dg 43.5\pr & S3\dg 57.7\pr \\
3 & 6 & 272\dg 43.7\pr & S3\dg 58.7\pr & 3 & 7 & 287\dg 43.8\pr & S3\dg 59.7\pr & 3 & 8 & 302\dg 44.0\pr & S4\dg 00.6\pr \\
3 & 9 & 317\dg 44.2\pr & S4\dg 01.6\pr & 3 & 10 & 332\dg 44.4\pr & S4\dg 02.6\pr & 3 & 11 & 347\dg 44.6\pr & S4\dg 03.5\pr \\
3 & 12 & 2\dg 44.8\pr & S4\dg 04.5\pr & 3 & 13 & 17\dg 45.0\pr & S4\dg 05.5\pr & 3 & 14 & 32\dg 45.2\pr & S4\dg 06.4\pr \\
3 & 15 & 47\dg 45.4\pr & S4\dg 07.4\pr & 3 & 16 & 62\dg 45.6\pr & S4\dg 08.4\pr & 3 & 17 & 77\dg 45.8\pr & S4\dg 09.3\pr \\
3 & 18 & 92\dg 46.0\pr & S4\dg 10.3\pr & 3 & 19 & 107\dg 46.2\pr & S4\dg 11.3\pr & 3 & 20 & 122\dg 46.4\pr & S4\dg 12.2\pr \\
3 & 21 & 137\dg 46.6\pr & S4\dg 13.2\pr & 3 & 22 & 152\dg 46.8\pr & S4\dg 14.1\pr & 3 & 23 & 167\dg 47.0\pr & S4\dg 15.1\pr \\
4 & 0 & 182\dg 47.2\pr & S4\dg 16.1\pr & 4 & 1 & 197\dg 47.3\pr & S4\dg 17.0\pr & 4 & 2 & 212\dg 47.5\pr & S4\dg 18.0\pr \\
4 & 3 & 227\dg 47.7\pr & S4\dg 19.0\pr & 4 & 4 & 242\dg 47.9\pr & S4\dg 19.9\pr & 4 & 5 & 257\dg 48.1\pr & S4\dg 20.9\pr \\
4 & 6 & 272\dg 48.3\pr & S4\dg 21.9\pr & 4 & 7 & 287\dg 48.5\pr & S4\dg 22.8\pr & 4 & 8 & 302\dg 48.7\pr & S4\dg 23.8\pr \\
4 & 9 & 317\dg 48.9\pr & S4\dg 24.7\pr & 4 & 10 & 332\dg 49.1\pr & S4\dg 25.7\pr & 4 & 11 & 347\dg 49.3\pr & S4\dg 26.7\pr \\
4 & 12 & 2\dg 49.5\pr & S4\dg 27.6\pr & 4 & 13 & 17\dg 49.7\pr & S4\dg 28.6\pr & 4 & 14 & 32\dg 49.8\pr & S4\dg 29.6\pr \\
4 & 15 & 47\dg 50.0\pr & S4\dg 30.5\pr & 4 & 16 & 62\dg 50.2\pr & S4\dg 31.5\pr & 4 & 17 & 77\dg 50.4\pr & S4\dg 32.5\pr \\
4 & 18 & 92\dg 50.6\pr & S4\dg 33.4\pr & 4 & 19 & 107\dg 50.8\pr & S4\dg 34.4\pr & 4 & 20 & 122\dg 51.0\pr & S4\dg 35.3\pr \\
4 & 21 & 137\dg 51.2\pr & S4\dg 36.3\pr & 4 & 22 & 152\dg 51.4\pr & S4\dg 37.3\pr & 4 & 23 & 167\dg 51.6\pr & S4\dg 38.2\pr \\
5 & 0 & 182\dg 51.7\pr & S4\dg 39.2\pr & 5 & 1 & 197\dg 51.9\pr & S4\dg 40.2\pr & 5 & 2 & 212\dg 52.1\pr & S4\dg 41.1\pr \\
5 & 3 & 227\dg 52.3\pr & S4\dg 42.1\pr & 5 & 4 & 242\dg 52.5\pr & S4\dg 43.0\pr & 5 & 5 & 257\dg 52.7\pr & S4\dg 44.0\pr \\
5 & 6 & 272\dg 52.9\pr & S4\dg 45.0\pr & 5 & 7 & 287\dg 53.1\pr & S4\dg 45.9\pr & 5 & 8 & 302\dg 53.3\pr & S4\dg 46.9\pr \\
5 & 9 & 317\dg 53.4\pr & S4\dg 47.8\pr & 5 & 10 & 332\dg 53.6\pr & S4\dg 48.8\pr & 5 & 11 & 347\dg 53.8\pr & S4\dg 49.8\pr \\
5 & 12 & 2\dg 54.0\pr & S4\dg 50.7\pr & 5 & 13 & 17\dg 54.2\pr & S4\dg 51.7\pr & 5 & 14 & 32\dg 54.4\pr & S4\dg 52.6\pr \\
5 & 15 & 47\dg 54.6\pr & S4\dg 53.6\pr & 5 & 16 & 62\dg 54.8\pr & S4\dg 54.6\pr & 5 & 17 & 77\dg 54.9\pr & S4\dg 55.5\pr \\
5 & 18 & 92\dg 55.1\pr & S4\dg 56.5\pr & 5 & 19 & 107\dg 55.3\pr & S4\dg 57.5\pr & 5 & 20 & 122\dg 55.5\pr & S4\dg 58.4\pr \\
5 & 21 & 137\dg 55.7\pr & S4\dg 59.4\pr & 5 & 22 & 152\dg 55.9\pr & S5\dg 00.3\pr & 5 & 23 & 167\dg 56.1\pr & S5\dg 01.3\pr \\
6 & 0 & 182\dg 56.2\pr & S5\dg 02.3\pr & 6 & 1 & 197\dg 56.4\pr & S5\dg 03.2\pr & 6 & 2 & 212\dg 56.6\pr & S5\dg 04.2\pr \\
6 & 3 & 227\dg 56.8\pr & S5\dg 05.1\pr & 6 & 4 & 242\dg 57.0\pr & S5\dg 06.1\pr & 6 & 5 & 257\dg 57.2\pr & S5\dg 07.0\pr \\
6 & 6 & 272\dg 57.4\pr & S5\dg 08.0\pr & 6 & 7 & 287\dg 57.5\pr & S5\dg 09.0\pr & 6 & 8 & 302\dg 57.7\pr & S5\dg 09.9\pr \\
6 & 9 & 317\dg 57.9\pr & S5\dg 10.9\pr & 6 & 10 & 332\dg 58.1\pr & S5\dg 11.8\pr & 6 & 11 & 347\dg 58.3\pr & S5\dg 12.8\pr \\
6 & 12 & 2\dg 58.5\pr & S5\dg 13.8\pr & 6 & 13 & 17\dg 58.6\pr & S5\dg 14.7\pr & 6 & 14 & 32\dg 58.8\pr & S5\dg 15.7\pr \\
6 & 15 & 47\dg 59.0\pr & S5\dg 16.6\pr & 6 & 16 & 62\dg 59.2\pr & S5\dg 17.6\pr & 6 & 17 & 77\dg 59.4\pr & S5\dg 18.5\pr \\
6 & 18 & 92\dg 59.6\pr & S5\dg 19.5\pr & 6 & 19 & 107\dg 59.7\pr & S5\dg 20.5\pr & 6 & 20 & 122\dg 59.9\pr & S5\dg 21.4\pr \\
6 & 21 & 138\dg 00.1\pr & S5\dg 22.4\pr & 6 & 22 & 153\dg 00.3\pr & S5\dg 23.3\pr & 6 & 23 & 168\dg 00.5\pr & S5\dg 24.3\pr \\
7 & 0 & 183\dg 00.7\pr & S5\dg 25.2\pr & 7 & 1 & 198\dg 00.8\pr & S5\dg 26.2\pr & 7 & 2 & 213\dg 01.0\pr & S5\dg 27.2\pr \\
7 & 3 & 228\dg 01.2\pr & S5\dg 28.1\pr & 7 & 4 & 243\dg 01.4\pr & S5\dg 29.1\pr & 7 & 5 & 258\dg 01.6\pr & S5\dg 30.0\pr \\
7 & 6 & 273\dg 01.7\pr & S5\dg 31.0\pr & 7 & 7 & 288\dg 01.9\pr & S5\dg 31.9\pr & 7 & 8 & 303\dg 02.1\pr & S5\dg 32.9\pr \\
7 & 9 & 318\dg 02.3\pr & S5\dg 33.9\pr & 7 & 10 & 333\dg 02.5\pr & S5\dg 34.8\pr & 7 & 11 & 348\dg 02.6\pr & S5\dg 35.8\pr \\
7 & 12 & 3\dg 02.8\pr & S5\dg 36.7\pr & 7 & 13 & 18\dg 03.0\pr & S5\dg 37.7\pr & 7 & 14 & 33\dg 03.2\pr & S5\dg 38.6\pr \\
7 & 15 & 48\dg 03.4\pr & S5\dg 39.6\pr & 7 & 16 & 63\dg 03.5\pr & S5\dg 40.5\pr & 7 & 17 & 78\dg 03.7\pr & S5\dg 41.5\pr \\
7 & 18 & 93\dg 03.9\pr & S5\dg 42.5\pr & 7 & 19 & 108\dg 04.1\pr & S5\dg 43.4\pr & 7 & 20 & 123\dg 04.2\pr & S5\dg 44.4\pr \\
7 & 21 & 138\dg 04.4\pr & S5\dg 45.3\pr & 7 & 22 & 153\dg 04.6\pr & S5\dg 46.3\pr & 7 & 23 & 168\dg 04.8\pr & S5\dg 47.2\pr \\
8 & 0 & 183\dg 05.0\pr & S5\dg 48.2\pr & 8 & 1 & 198\dg 05.1\pr & S5\dg 49.1\pr & 8 & 2 & 213\dg 05.3\pr & S5\dg 50.1\pr \\
8 & 3 & 228\dg 05.5\pr & S5\dg 51.0\pr & 8 & 4 & 243\dg 05.7\pr & S5\dg 52.0\pr & 8 & 5 & 258\dg 05.8\pr & S5\dg 52.9\pr \\
8 & 6 & 273\dg 06.0\pr & S5\dg 53.9\pr & 8 & 7 & 288\dg 06.2\pr & S5\dg 54.9\pr & 8 & 8 & 303\dg 06.4\pr & S5\dg 55.8\pr \\
8 & 9 & 318\dg 06.5\pr & S5\dg 56.8\pr & 8 & 10 & 333\dg 06.7\pr & S5\dg 57.7\pr & 8 & 11 & 348\dg 06.9\pr & S5\dg 58.7\pr \\
8 & 12 & 3\dg 07.1\pr & S5\dg 59.6\pr & 8 & 13 & 18\dg 07.2\pr & S6\dg 00.6\pr & 8 & 14 & 33\dg 07.4\pr & S6\dg 01.5\pr \\
8 & 15 & 48\dg 07.6\pr & S6\dg 02.5\pr & 8 & 16 & 63\dg 07.8\pr & S6\dg 03.4\pr & 8 & 17 & 78\dg 07.9\pr & S6\dg 04.4\pr \\
8 & 18 & 93\dg 08.1\pr & S6\dg 05.3\pr & 8 & 19 & 108\dg 08.3\pr & S6\dg 06.3\pr & 8 & 20 & 123\dg 08.5\pr & S6\dg 07.2\pr \\
8 & 21 & 138\dg 08.6\pr & S6\dg 08.2\pr & 8 & 22 & 153\dg 08.8\pr & S6\dg 09.1\pr & 8 & 23 & 168\dg 09.0\pr & S6\dg 10.1\pr \\
9 & 0 & 183\dg 09.2\pr & S6\dg 11.0\pr & 9 & 1 & 198\dg 09.3\pr & S6\dg 12.0\pr & 9 & 2 & 213\dg 09.5\pr & S6\dg 12.9\pr \\
9 & 3 & 228\dg 09.7\pr & S6\dg 13.9\pr & 9 & 4 & 243\dg 09.8\pr & S6\dg 14.8\pr & 9 & 5 & 258\dg 10.0\pr & S6\dg 15.8\pr \\
9 & 6 & 273\dg 10.2\pr & S6\dg 16.7\pr & 9 & 7 & 288\dg 10.4\pr & S6\dg 17.7\pr & 9 & 8 & 303\dg 10.5\pr & S6\dg 18.6\pr \\
9 & 9 & 318\dg 10.7\pr & S6\dg 19.6\pr & 9 & 10 & 333\dg 10.9\pr & S6\dg 20.5\pr & 9 & 11 & 348\dg 11.0\pr & S6\dg 21.5\pr \\
9 & 12 & 3\dg 11.2\pr & S6\dg 22.4\pr & 9 & 13 & 18\dg 11.4\pr & S6\dg 23.4\pr & 9 & 14 & 33\dg 11.6\pr & S6\dg 24.3\pr \\
9 & 15 & 48\dg 11.7\pr & S6\dg 25.3\pr & 9 & 16 & 63\dg 11.9\pr & S6\dg 26.2\pr & 9 & 17 & 78\dg 12.1\pr & S6\dg 27.2\pr \\
9 & 18 & 93\dg 12.2\pr & S6\dg 28.1\pr & 9 & 19 & 108\dg 12.4\pr & S6\dg 29.1\pr & 9 & 20 & 123\dg 12.6\pr & S6\dg 30.0\pr \\
9 & 21 & 138\dg 12.7\pr & S6\dg 31.0\pr & 9 & 22 & 153\dg 12.9\pr & S6\dg 31.9\pr & 9 & 23 & 168\dg 13.1\pr & S6\dg 32.9\pr \\
10 & 0 & 183\dg 13.2\pr & S6\dg 33.8\pr & 10 & 1 & 198\dg 13.4\pr & S6\dg 34.7\pr & 10 & 2 & 213\dg 13.6\pr & S6\dg 35.7\pr \\
10 & 3 & 228\dg 13.7\pr & S6\dg 36.6\pr & 10 & 4 & 243\dg 13.9\pr & S6\dg 37.6\pr & 10 & 5 & 258\dg 14.1\pr & S6\dg 38.5\pr \\
10 & 6 & 273\dg 14.2\pr & S6\dg 39.5\pr & 10 & 7 & 288\dg 14.4\pr & S6\dg 40.4\pr & 10 & 8 & 303\dg 14.6\pr & S6\dg 41.4\pr \\
10 & 9 & 318\dg 14.7\pr & S6\dg 42.3\pr & 10 & 10 & 333\dg 14.9\pr & S6\dg 43.3\pr & 10 & 11 & 348\dg 15.1\pr & S6\dg 44.2\pr \\
10 & 12 & 3\dg 15.2\pr & S6\dg 45.2\pr & 10 & 13 & 18\dg 15.4\pr & S6\dg 46.1\pr & 10 & 14 & 33\dg 15.6\pr & S6\dg 47.0\pr \\
10 & 15 & 48\dg 15.7\pr & S6\dg 48.0\pr & 10 & 16 & 63\dg 15.9\pr & S6\dg 48.9\pr & 10 & 17 & 78\dg 16.1\pr & S6\dg 49.9\pr \\
10 & 18 & 93\dg 16.2\pr & S6\dg 50.8\pr & 10 & 19 & 108\dg 16.4\pr & S6\dg 51.8\pr & 10 & 20 & 123\dg 16.6\pr & S6\dg 52.7\pr \\
10 & 21 & 138\dg 16.7\pr & S6\dg 53.7\pr & 10 & 22 & 153\dg 16.9\pr & S6\dg 54.6\pr & 10 & 23 & 168\dg 17.0\pr & S6\dg 55.5\pr \\
11 & 0 & 183\dg 17.2\pr & S6\dg 56.5\pr & 11 & 1 & 198\dg 17.4\pr & S6\dg 57.4\pr & 11 & 2 & 213\dg 17.5\pr & S6\dg 58.4\pr \\
11 & 3 & 228\dg 17.7\pr & S6\dg 59.3\pr & 11 & 4 & 243\dg 17.9\pr & S7\dg 00.3\pr & 11 & 5 & 258\dg 18.0\pr & S7\dg 01.2\pr \\
11 & 6 & 273\dg 18.2\pr & S7\dg 02.1\pr & 11 & 7 & 288\dg 18.3\pr & S7\dg 03.1\pr & 11 & 8 & 303\dg 18.5\pr & S7\dg 04.0\pr \\
11 & 9 & 318\dg 18.7\pr & S7\dg 05.0\pr & 11 & 10 & 333\dg 18.8\pr & S7\dg 05.9\pr & 11 & 11 & 348\dg 19.0\pr & S7\dg 06.9\pr \\
11 & 12 & 3\dg 19.2\pr & S7\dg 07.8\pr & 11 & 13 & 18\dg 19.3\pr & S7\dg 08.7\pr & 11 & 14 & 33\dg 19.5\pr & S7\dg 09.7\pr \\
11 & 15 & 48\dg 19.6\pr & S7\dg 10.6\pr & 11 & 16 & 63\dg 19.8\pr & S7\dg 11.6\pr & 11 & 17 & 78\dg 20.0\pr & S7\dg 12.5\pr \\
11 & 18 & 93\dg 20.1\pr & S7\dg 13.4\pr & 11 & 19 & 108\dg 20.3\pr & S7\dg 14.4\pr & 11 & 20 & 123\dg 20.4\pr & S7\dg 15.3\pr \\
11 & 21 & 138\dg 20.6\pr & S7\dg 16.3\pr & 11 & 22 & 153\dg 20.7\pr & S7\dg 17.2\pr & 11 & 23 & 168\dg 20.9\pr & S7\dg 18.1\pr \\
12 & 0 & 183\dg 21.1\pr & S7\dg 19.1\pr & 12 & 1 & 198\dg 21.2\pr & S7\dg 20.0\pr & 12 & 2 & 213\dg 21.4\pr & S7\dg 21.0\pr \\
12 & 3 & 228\dg 21.5\pr & S7\dg 21.9\pr & 12 & 4 & 243\dg 21.7\pr & S7\dg 22.8\pr & 12 & 5 & 258\dg 21.9\pr & S7\dg 23.8\pr \\
12 & 6 & 273\dg 22.0\pr & S7\dg 24.7\pr & 12 & 7 & 288\dg 22.2\pr & S7\dg 25.7\pr & 12 & 8 & 303\dg 22.3\pr & S7\dg 26.6\pr \\
12 & 9 & 318\dg 22.5\pr & S7\dg 27.5\pr & 12 & 10 & 333\dg 22.6\pr & S7\dg 28.5\pr & 12 & 11 & 348\dg 22.8\pr & S7\dg 29.4\pr \\
12 & 12 & 3\dg 22.9\pr & S7\dg 30.3\pr & 12 & 13 & 18\dg 23.1\pr & S7\dg 31.3\pr & 12 & 14 & 33\dg 23.3\pr & S7\dg 32.2\pr \\
12 & 15 & 48\dg 23.4\pr & S7\dg 33.1\pr & 12 & 16 & 63\dg 23.6\pr & S7\dg 34.1\pr & 12 & 17 & 78\dg 23.7\pr & S7\dg 35.0\pr \\
12 & 18 & 93\dg 23.9\pr & S7\dg 36.0\pr & 12 & 19 & 108\dg 24.0\pr & S7\dg 36.9\pr & 12 & 20 & 123\dg 24.2\pr & S7\dg 37.8\pr \\
12 & 21 & 138\dg 24.3\pr & S7\dg 38.8\pr & 12 & 22 & 153\dg 24.5\pr & S7\dg 39.7\pr & 12 & 23 & 168\dg 24.6\pr & S7\dg 40.6\pr \\
13 & 0 & 183\dg 24.8\pr & S7\dg 41.6\pr & 13 & 1 & 198\dg 24.9\pr & S7\dg 42.5\pr & 13 & 2 & 213\dg 25.1\pr & S7\dg 43.4\pr \\
13 & 3 & 228\dg 25.2\pr & S7\dg 44.4\pr & 13 & 4 & 243\dg 25.4\pr & S7\dg 45.3\pr & 13 & 5 & 258\dg 25.6\pr & S7\dg 46.2\pr \\
13 & 6 & 273\dg 25.7\pr & S7\dg 47.2\pr & 13 & 7 & 288\dg 25.9\pr & S7\dg 48.1\pr & 13 & 8 & 303\dg 26.0\pr & S7\dg 49.0\pr \\
13 & 9 & 318\dg 26.2\pr & S7\dg 50.0\pr & 13 & 10 & 333\dg 26.3\pr & S7\dg 50.9\pr & 13 & 11 & 348\dg 26.5\pr & S7\dg 51.8\pr \\
13 & 12 & 3\dg 26.6\pr & S7\dg 52.8\pr & 13 & 13 & 18\dg 26.8\pr & S7\dg 53.7\pr & 13 & 14 & 33\dg 26.9\pr & S7\dg 54.6\pr \\
13 & 15 & 48\dg 27.1\pr & S7\dg 55.6\pr & 13 & 16 & 63\dg 27.2\pr & S7\dg 56.5\pr & 13 & 17 & 78\dg 27.4\pr & S7\dg 57.4\pr \\
13 & 18 & 93\dg 27.5\pr & S7\dg 58.4\pr & 13 & 19 & 108\dg 27.6\pr & S7\dg 59.3\pr & 13 & 20 & 123\dg 27.8\pr & S8\dg 00.2\pr \\
13 & 21 & 138\dg 27.9\pr & S8\dg 01.2\pr & 13 & 22 & 153\dg 28.1\pr & S8\dg 02.1\pr & 13 & 23 & 168\dg 28.2\pr & S8\dg 03.0\pr \\
14 & 0 & 183\dg 28.4\pr & S8\dg 04.0\pr & 14 & 1 & 198\dg 28.5\pr & S8\dg 04.9\pr & 14 & 2 & 213\dg 28.7\pr & S8\dg 05.8\pr \\
14 & 3 & 228\dg 28.8\pr & S8\dg 06.7\pr & 14 & 4 & 243\dg 29.0\pr & S8\dg 07.7\pr & 14 & 5 & 258\dg 29.1\pr & S8\dg 08.6\pr \\
14 & 6 & 273\dg 29.3\pr & S8\dg 09.5\pr & 14 & 7 & 288\dg 29.4\pr & S8\dg 10.5\pr & 14 & 8 & 303\dg 29.6\pr & S8\dg 11.4\pr \\
14 & 9 & 318\dg 29.7\pr & S8\dg 12.3\pr & 14 & 10 & 333\dg 29.8\pr & S8\dg 13.3\pr & 14 & 11 & 348\dg 30.0\pr & S8\dg 14.2\pr \\
14 & 12 & 3\dg 30.1\pr & S8\dg 15.1\pr & 14 & 13 & 18\dg 30.3\pr & S8\dg 16.0\pr & 14 & 14 & 33\dg 30.4\pr & S8\dg 17.0\pr \\
14 & 15 & 48\dg 30.6\pr & S8\dg 17.9\pr & 14 & 16 & 63\dg 30.7\pr & S8\dg 18.8\pr & 14 & 17 & 78\dg 30.9\pr & S8\dg 19.7\pr \\
14 & 18 & 93\dg 31.0\pr & S8\dg 20.7\pr & 14 & 19 & 108\dg 31.1\pr & S8\dg 21.6\pr & 14 & 20 & 123\dg 31.3\pr & S8\dg 22.5\pr \\
14 & 21 & 138\dg 31.4\pr & S8\dg 23.5\pr & 14 & 22 & 153\dg 31.6\pr & S8\dg 24.4\pr & 14 & 23 & 168\dg 31.7\pr & S8\dg 25.3\pr \\
15 & 0 & 183\dg 31.9\pr & S8\dg 26.2\pr & 15 & 1 & 198\dg 32.0\pr & S8\dg 27.2\pr & 15 & 2 & 213\dg 32.1\pr & S8\dg 28.1\pr \\
15 & 3 & 228\dg 32.3\pr & S8\dg 29.0\pr & 15 & 4 & 243\dg 32.4\pr & S8\dg 29.9\pr & 15 & 5 & 258\dg 32.6\pr & S8\dg 30.9\pr \\
15 & 6 & 273\dg 32.7\pr & S8\dg 31.8\pr & 15 & 7 & 288\dg 32.8\pr & S8\dg 32.7\pr & 15 & 8 & 303\dg 33.0\pr & S8\dg 33.6\pr \\
15 & 9 & 318\dg 33.1\pr & S8\dg 34.6\pr & 15 & 10 & 333\dg 33.3\pr & S8\dg 35.5\pr & 15 & 11 & 348\dg 33.4\pr & S8\dg 36.4\pr \\
15 & 12 & 3\dg 33.5\pr & S8\dg 37.3\pr & 15 & 13 & 18\dg 33.7\pr & S8\dg 38.2\pr & 15 & 14 & 33\dg 33.8\pr & S8\dg 39.2\pr \\
15 & 15 & 48\dg 33.9\pr & S8\dg 40.1\pr & 15 & 16 & 63\dg 34.1\pr & S8\dg 41.0\pr & 15 & 17 & 78\dg 34.2\pr & S8\dg 41.9\pr \\
15 & 18 & 93\dg 34.4\pr & S8\dg 42.9\pr & 15 & 19 & 108\dg 34.5\pr & S8\dg 43.8\pr & 15 & 20 & 123\dg 34.6\pr & S8\dg 44.7\pr \\
15 & 21 & 138\dg 34.8\pr & S8\dg 45.6\pr & 15 & 22 & 153\dg 34.9\pr & S8\dg 46.5\pr & 15 & 23 & 168\dg 35.0\pr & S8\dg 47.5\pr \\
16 & 0 & 183\dg 35.2\pr & S8\dg 48.4\pr & 16 & 1 & 198\dg 35.3\pr & S8\dg 49.3\pr & 16 & 2 & 213\dg 35.4\pr & S8\dg 50.2\pr \\
16 & 3 & 228\dg 35.6\pr & S8\dg 51.1\pr & 16 & 4 & 243\dg 35.7\pr & S8\dg 52.1\pr & 16 & 5 & 258\dg 35.8\pr & S8\dg 53.0\pr \\
16 & 6 & 273\dg 36.0\pr & S8\dg 53.9\pr & 16 & 7 & 288\dg 36.1\pr & S8\dg 54.8\pr & 16 & 8 & 303\dg 36.3\pr & S8\dg 55.7\pr \\
16 & 9 & 318\dg 36.4\pr & S8\dg 56.7\pr & 16 & 10 & 333\dg 36.5\pr & S8\dg 57.6\pr & 16 & 11 & 348\dg 36.7\pr & S8\dg 58.5\pr \\
16 & 12 & 3\dg 36.8\pr & S8\dg 59.4\pr & 16 & 13 & 18\dg 36.9\pr & S9\dg 00.3\pr & 16 & 14 & 33\dg 37.0\pr & S9\dg 01.3\pr \\
16 & 15 & 48\dg 37.2\pr & S9\dg 02.2\pr & 16 & 16 & 63\dg 37.3\pr & S9\dg 03.1\pr & 16 & 17 & 78\dg 37.4\pr & S9\dg 04.0\pr \\
16 & 18 & 93\dg 37.6\pr & S9\dg 04.9\pr & 16 & 19 & 108\dg 37.7\pr & S9\dg 05.8\pr & 16 & 20 & 123\dg 37.8\pr & S9\dg 06.8\pr \\
16 & 21 & 138\dg 38.0\pr & S9\dg 07.7\pr & 16 & 22 & 153\dg 38.1\pr & S9\dg 08.6\pr & 16 & 23 & 168\dg 38.2\pr & S9\dg 09.5\pr \\
17 & 0 & 183\dg 38.4\pr & S9\dg 10.4\pr & 17 & 1 & 198\dg 38.5\pr & S9\dg 11.3\pr & 17 & 2 & 213\dg 38.6\pr & S9\dg 12.2\pr \\
17 & 3 & 228\dg 38.7\pr & S9\dg 13.2\pr & 17 & 4 & 243\dg 38.9\pr & S9\dg 14.1\pr & 17 & 5 & 258\dg 39.0\pr & S9\dg 15.0\pr \\
17 & 6 & 273\dg 39.1\pr & S9\dg 15.9\pr & 17 & 7 & 288\dg 39.3\pr & S9\dg 16.8\pr & 17 & 8 & 303\dg 39.4\pr & S9\dg 17.7\pr \\
17 & 9 & 318\dg 39.5\pr & S9\dg 18.6\pr & 17 & 10 & 333\dg 39.6\pr & S9\dg 19.6\pr & 17 & 11 & 348\dg 39.8\pr & S9\dg 20.5\pr \\
17 & 12 & 3\dg 39.9\pr & S9\dg 21.4\pr & 17 & 13 & 18\dg 40.0\pr & S9\dg 22.3\pr & 17 & 14 & 33\dg 40.1\pr & S9\dg 23.2\pr \\
17 & 15 & 48\dg 40.3\pr & S9\dg 24.1\pr & 17 & 16 & 63\dg 40.4\pr & S9\dg 25.0\pr & 17 & 17 & 78\dg 40.5\pr & S9\dg 25.9\pr \\
17 & 18 & 93\dg 40.6\pr & S9\dg 26.9\pr & 17 & 19 & 108\dg 40.8\pr & S9\dg 27.8\pr & 17 & 20 & 123\dg 40.9\pr & S9\dg 28.7\pr \\
17 & 21 & 138\dg 41.0\pr & S9\dg 29.6\pr & 17 & 22 & 153\dg 41.1\pr & S9\dg 30.5\pr & 17 & 23 & 168\dg 41.3\pr & S9\dg 31.4\pr \\
18 & 0 & 183\dg 41.4\pr & S9\dg 32.3\pr & 18 & 1 & 198\dg 41.5\pr & S9\dg 33.2\pr & 18 & 2 & 213\dg 41.6\pr & S9\dg 34.1\pr \\
18 & 3 & 228\dg 41.8\pr & S9\dg 35.0\pr & 18 & 4 & 243\dg 41.9\pr & S9\dg 36.0\pr & 18 & 5 & 258\dg 42.0\pr & S9\dg 36.9\pr \\
18 & 6 & 273\dg 42.1\pr & S9\dg 37.8\pr & 18 & 7 & 288\dg 42.2\pr & S9\dg 38.7\pr & 18 & 8 & 303\dg 42.4\pr & S9\dg 39.6\pr \\
18 & 9 & 318\dg 42.5\pr & S9\dg 40.5\pr & 18 & 10 & 333\dg 42.6\pr & S9\dg 41.4\pr & 18 & 11 & 348\dg 42.7\pr & S9\dg 42.3\pr \\
18 & 12 & 3\dg 42.8\pr & S9\dg 43.2\pr & 18 & 13 & 18\dg 43.0\pr & S9\dg 44.1\pr & 18 & 14 & 33\dg 43.1\pr & S9\dg 45.0\pr \\
18 & 15 & 48\dg 43.2\pr & S9\dg 45.9\pr & 18 & 16 & 63\dg 43.3\pr & S9\dg 46.8\pr & 18 & 17 & 78\dg 43.4\pr & S9\dg 47.7\pr \\
18 & 18 & 93\dg 43.6\pr & S9\dg 48.7\pr & 18 & 19 & 108\dg 43.7\pr & S9\dg 49.6\pr & 18 & 20 & 123\dg 43.8\pr & S9\dg 50.5\pr \\
18 & 21 & 138\dg 43.9\pr & S9\dg 51.4\pr & 18 & 22 & 153\dg 44.0\pr & S9\dg 52.3\pr & 18 & 23 & 168\dg 44.2\pr & S9\dg 53.2\pr \\
19 & 0 & 183\dg 44.3\pr & S9\dg 54.1\pr & 19 & 1 & 198\dg 44.4\pr & S9\dg 55.0\pr & 19 & 2 & 213\dg 44.5\pr & S9\dg 55.9\pr \\
19 & 3 & 228\dg 44.6\pr & S9\dg 56.8\pr & 19 & 4 & 243\dg 44.7\pr & S9\dg 57.7\pr & 19 & 5 & 258\dg 44.9\pr & S9\dg 58.6\pr \\
19 & 6 & 273\dg 45.0\pr & S9\dg 59.5\pr & 19 & 7 & 288\dg 45.1\pr & S10\dg 00.4\pr & 19 & 8 & 303\dg 45.2\pr & S10\dg 01.3\pr \\
19 & 9 & 318\dg 45.3\pr & S10\dg 02.2\pr & 19 & 10 & 333\dg 45.4\pr & S10\dg 03.1\pr & 19 & 11 & 348\dg 45.5\pr & S10\dg 04.0\pr \\
19 & 12 & 3\dg 45.7\pr & S10\dg 04.9\pr & 19 & 13 & 18\dg 45.8\pr & S10\dg 05.8\pr & 19 & 14 & 33\dg 45.9\pr & S10\dg 06.7\pr \\
19 & 15 & 48\dg 46.0\pr & S10\dg 07.6\pr & 19 & 16 & 63\dg 46.1\pr & S10\dg 08.5\pr & 19 & 17 & 78\dg 46.2\pr & S10\dg 09.4\pr \\
19 & 18 & 93\dg 46.3\pr & S10\dg 10.3\pr & 19 & 19 & 108\dg 46.4\pr & S10\dg 11.2\pr & 19 & 20 & 123\dg 46.6\pr & S10\dg 12.1\pr \\
19 & 21 & 138\dg 46.7\pr & S10\dg 13.0\pr & 19 & 22 & 153\dg 46.8\pr & S10\dg 13.9\pr & 19 & 23 & 168\dg 46.9\pr & S10\dg 14.8\pr \\
20 & 0 & 183\dg 47.0\pr & S10\dg 15.7\pr & 20 & 1 & 198\dg 47.1\pr & S10\dg 16.6\pr & 20 & 2 & 213\dg 47.2\pr & S10\dg 17.5\pr \\
20 & 3 & 228\dg 47.3\pr & S10\dg 18.4\pr & 20 & 4 & 243\dg 47.4\pr & S10\dg 19.3\pr & 20 & 5 & 258\dg 47.5\pr & S10\dg 20.2\pr \\
20 & 6 & 273\dg 47.7\pr & S10\dg 21.1\pr & 20 & 7 & 288\dg 47.8\pr & S10\dg 22.0\pr & 20 & 8 & 303\dg 47.9\pr & S10\dg 22.9\pr \\
20 & 9 & 318\dg 48.0\pr & S10\dg 23.8\pr & 20 & 10 & 333\dg 48.1\pr & S10\dg 24.7\pr & 20 & 11 & 348\dg 48.2\pr & S10\dg 25.6\pr \\
20 & 12 & 3\dg 48.3\pr & S10\dg 26.4\pr & 20 & 13 & 18\dg 48.4\pr & S10\dg 27.3\pr & 20 & 14 & 33\dg 48.5\pr & S10\dg 28.2\pr \\
20 & 15 & 48\dg 48.6\pr & S10\dg 29.1\pr & 20 & 16 & 63\dg 48.7\pr & S10\dg 30.0\pr & 20 & 17 & 78\dg 48.8\pr & S10\dg 30.9\pr \\
20 & 18 & 93\dg 48.9\pr & S10\dg 31.8\pr & 20 & 19 & 108\dg 49.0\pr & S10\dg 32.7\pr & 20 & 20 & 123\dg 49.1\pr & S10\dg 33.6\pr \\
20 & 21 & 138\dg 49.3\pr & S10\dg 34.5\pr & 20 & 22 & 153\dg 49.4\pr & S10\dg 35.4\pr & 20 & 23 & 168\dg 49.5\pr & S10\dg 36.3\pr \\
21 & 0 & 183\dg 49.6\pr & S10\dg 37.2\pr & 21 & 1 & 198\dg 49.7\pr & S10\dg 38.1\pr & 21 & 2 & 213\dg 49.8\pr & S10\dg 38.9\pr \\
21 & 3 & 228\dg 49.9\pr & S10\dg 39.8\pr & 21 & 4 & 243\dg 50.0\pr & S10\dg 40.7\pr & 21 & 5 & 258\dg 50.1\pr & S10\dg 41.6\pr \\
21 & 6 & 273\dg 50.2\pr & S10\dg 42.5\pr & 21 & 7 & 288\dg 50.3\pr & S10\dg 43.4\pr & 21 & 8 & 303\dg 50.4\pr & S10\dg 44.3\pr \\
21 & 9 & 318\dg 50.5\pr & S10\dg 45.2\pr & 21 & 10 & 333\dg 50.6\pr & S10\dg 46.1\pr & 21 & 11 & 348\dg 50.7\pr & S10\dg 47.0\pr \\
21 & 12 & 3\dg 50.8\pr & S10\dg 47.8\pr & 21 & 13 & 18\dg 50.9\pr & S10\dg 48.7\pr & 21 & 14 & 33\dg 51.0\pr & S10\dg 49.6\pr \\
21 & 15 & 48\dg 51.1\pr & S10\dg 50.5\pr & 21 & 16 & 63\dg 51.2\pr & S10\dg 51.4\pr & 21 & 17 & 78\dg 51.3\pr & S10\dg 52.3\pr \\
21 & 18 & 93\dg 51.4\pr & S10\dg 53.2\pr & 21 & 19 & 108\dg 51.5\pr & S10\dg 54.0\pr & 21 & 20 & 123\dg 51.6\pr & S10\dg 54.9\pr \\
21 & 21 & 138\dg 51.7\pr & S10\dg 55.8\pr & 21 & 22 & 153\dg 51.8\pr & S10\dg 56.7\pr & 21 & 23 & 168\dg 51.9\pr & S10\dg 57.6\pr \\
22 & 0 & 183\dg 52.0\pr & S10\dg 58.5\pr & 22 & 1 & 198\dg 52.1\pr & S10\dg 59.4\pr & 22 & 2 & 213\dg 52.2\pr & S11\dg 00.2\pr \\
22 & 3 & 228\dg 52.3\pr & S11\dg 01.1\pr & 22 & 4 & 243\dg 52.3\pr & S11\dg 02.0\pr & 22 & 5 & 258\dg 52.4\pr & S11\dg 02.9\pr \\
22 & 6 & 273\dg 52.5\pr & S11\dg 03.8\pr & 22 & 7 & 288\dg 52.6\pr & S11\dg 04.7\pr & 22 & 8 & 303\dg 52.7\pr & S11\dg 05.5\pr \\
22 & 9 & 318\dg 52.8\pr & S11\dg 06.4\pr & 22 & 10 & 333\dg 52.9\pr & S11\dg 07.3\pr & 22 & 11 & 348\dg 53.0\pr & S11\dg 08.2\pr \\
22 & 12 & 3\dg 53.1\pr & S11\dg 09.1\pr & 22 & 13 & 18\dg 53.2\pr & S11\dg 09.9\pr & 22 & 14 & 33\dg 53.3\pr & S11\dg 10.8\pr \\
22 & 15 & 48\dg 53.4\pr & S11\dg 11.7\pr & 22 & 16 & 63\dg 53.5\pr & S11\dg 12.6\pr & 22 & 17 & 78\dg 53.6\pr & S11\dg 13.5\pr \\
22 & 18 & 93\dg 53.7\pr & S11\dg 14.3\pr & 22 & 19 & 108\dg 53.7\pr & S11\dg 15.2\pr & 22 & 20 & 123\dg 53.8\pr & S11\dg 16.1\pr \\
22 & 21 & 138\dg 53.9\pr & S11\dg 17.0\pr & 22 & 22 & 153\dg 54.0\pr & S11\dg 17.9\pr & 22 & 23 & 168\dg 54.1\pr & S11\dg 18.7\pr \\
23 & 0 & 183\dg 54.2\pr & S11\dg 19.6\pr & 23 & 1 & 198\dg 54.3\pr & S11\dg 20.5\pr & 23 & 2 & 213\dg 54.4\pr & S11\dg 21.4\pr \\
23 & 3 & 228\dg 54.5\pr & S11\dg 22.3\pr & 23 & 4 & 243\dg 54.6\pr & S11\dg 23.1\pr & 23 & 5 & 258\dg 54.6\pr & S11\dg 24.0\pr \\
23 & 6 & 273\dg 54.7\pr & S11\dg 24.9\pr & 23 & 7 & 288\dg 54.8\pr & S11\dg 25.8\pr & 23 & 8 & 303\dg 54.9\pr & S11\dg 26.6\pr \\
23 & 9 & 318\dg 55.0\pr & S11\dg 27.5\pr & 23 & 10 & 333\dg 55.1\pr & S11\dg 28.4\pr & 23 & 11 & 348\dg 55.2\pr & S11\dg 29.3\pr \\
23 & 12 & 3\dg 55.3\pr & S11\dg 30.1\pr & 23 & 13 & 18\dg 55.3\pr & S11\dg 31.0\pr & 23 & 14 & 33\dg 55.4\pr & S11\dg 31.9\pr \\
23 & 15 & 48\dg 55.5\pr & S11\dg 32.8\pr & 23 & 16 & 63\dg 55.6\pr & S11\dg 33.6\pr & 23 & 17 & 78\dg 55.7\pr & S11\dg 34.5\pr \\
23 & 18 & 93\dg 55.8\pr & S11\dg 35.4\pr & 23 & 19 & 108\dg 55.8\pr & S11\dg 36.2\pr & 23 & 20 & 123\dg 55.9\pr & S11\dg 37.1\pr \\
23 & 21 & 138\dg 56.0\pr & S11\dg 38.0\pr & 23 & 22 & 153\dg 56.1\pr & S11\dg 38.9\pr & 23 & 23 & 168\dg 56.2\pr & S11\dg 39.7\pr \\
24 & 0 & 183\dg 56.3\pr & S11\dg 40.6\pr & 24 & 1 & 198\dg 56.3\pr & S11\dg 41.5\pr & 24 & 2 & 213\dg 56.4\pr & S11\dg 42.3\pr \\
24 & 3 & 228\dg 56.5\pr & S11\dg 43.2\pr & 24 & 4 & 243\dg 56.6\pr & S11\dg 44.1\pr & 24 & 5 & 258\dg 56.7\pr & S11\dg 44.9\pr \\
24 & 6 & 273\dg 56.8\pr & S11\dg 45.8\pr & 24 & 7 & 288\dg 56.8\pr & S11\dg 46.7\pr & 24 & 8 & 303\dg 56.9\pr & S11\dg 47.5\pr \\
24 & 9 & 318\dg 57.0\pr & S11\dg 48.4\pr & 24 & 10 & 333\dg 57.1\pr & S11\dg 49.3\pr & 24 & 11 & 348\dg 57.1\pr & S11\dg 50.2\pr \\
24 & 12 & 3\dg 57.2\pr & S11\dg 51.0\pr & 24 & 13 & 18\dg 57.3\pr & S11\dg 51.9\pr & 24 & 14 & 33\dg 57.4\pr & S11\dg 52.8\pr \\
24 & 15 & 48\dg 57.5\pr & S11\dg 53.6\pr & 24 & 16 & 63\dg 57.5\pr & S11\dg 54.5\pr & 24 & 17 & 78\dg 57.6\pr & S11\dg 55.3\pr \\
24 & 18 & 93\dg 57.7\pr & S11\dg 56.2\pr & 24 & 19 & 108\dg 57.8\pr & S11\dg 57.1\pr & 24 & 20 & 123\dg 57.8\pr & S11\dg 57.9\pr \\
24 & 21 & 138\dg 57.9\pr & S11\dg 58.8\pr & 24 & 22 & 153\dg 58.0\pr & S11\dg 59.7\pr & 24 & 23 & 168\dg 58.1\pr & S12\dg 00.5\pr \\
25 & 0 & 183\dg 58.1\pr & S12\dg 01.4\pr & 25 & 1 & 198\dg 58.2\pr & S12\dg 02.3\pr & 25 & 2 & 213\dg 58.3\pr & S12\dg 03.1\pr \\
25 & 3 & 228\dg 58.4\pr & S12\dg 04.0\pr & 25 & 4 & 243\dg 58.4\pr & S12\dg 04.8\pr & 25 & 5 & 258\dg 58.5\pr & S12\dg 05.7\pr \\
25 & 6 & 273\dg 58.6\pr & S12\dg 06.6\pr & 25 & 7 & 288\dg 58.7\pr & S12\dg 07.4\pr & 25 & 8 & 303\dg 58.7\pr & S12\dg 08.3\pr \\
25 & 9 & 318\dg 58.8\pr & S12\dg 09.1\pr & 25 & 10 & 333\dg 58.9\pr & S12\dg 10.0\pr & 25 & 11 & 348\dg 59.0\pr & S12\dg 10.9\pr \\
25 & 12 & 3\dg 59.0\pr & S12\dg 11.7\pr & 25 & 13 & 18\dg 59.1\pr & S12\dg 12.6\pr & 25 & 14 & 33\dg 59.2\pr & S12\dg 13.4\pr \\
25 & 15 & 48\dg 59.2\pr & S12\dg 14.3\pr & 25 & 16 & 63\dg 59.3\pr & S12\dg 15.2\pr & 25 & 17 & 78\dg 59.4\pr & S12\dg 16.0\pr \\
25 & 18 & 93\dg 59.4\pr & S12\dg 16.9\pr & 25 & 19 & 108\dg 59.5\pr & S12\dg 17.7\pr & 25 & 20 & 123\dg 59.6\pr & S12\dg 18.6\pr \\
25 & 21 & 138\dg 59.7\pr & S12\dg 19.4\pr & 25 & 22 & 153\dg 59.7\pr & S12\dg 20.3\pr & 25 & 23 & 168\dg 59.8\pr & S12\dg 21.2\pr \\
26 & 0 & 183\dg 59.9\pr & S12\dg 22.0\pr & 26 & 1 & 198\dg 59.9\pr & S12\dg 22.9\pr & 26 & 2 & 214\dg 00.0\pr & S12\dg 23.7\pr \\
26 & 3 & 229\dg 00.1\pr & S12\dg 24.6\pr & 26 & 4 & 244\dg 00.1\pr & S12\dg 25.4\pr & 26 & 5 & 259\dg 00.2\pr & S12\dg 26.3\pr \\
26 & 6 & 274\dg 00.3\pr & S12\dg 27.1\pr & 26 & 7 & 289\dg 00.3\pr & S12\dg 28.0\pr & 26 & 8 & 304\dg 00.4\pr & S12\dg 28.8\pr \\
26 & 9 & 319\dg 00.5\pr & S12\dg 29.7\pr & 26 & 10 & 334\dg 00.5\pr & S12\dg 30.5\pr & 26 & 11 & 349\dg 00.6\pr & S12\dg 31.4\pr \\
26 & 12 & 4\dg 00.6\pr & S12\dg 32.2\pr & 26 & 13 & 19\dg 00.7\pr & S12\dg 33.1\pr & 26 & 14 & 34\dg 00.8\pr & S12\dg 33.9\pr \\
26 & 15 & 49\dg 00.8\pr & S12\dg 34.8\pr & 26 & 16 & 64\dg 00.9\pr & S12\dg 35.6\pr & 26 & 17 & 79\dg 01.0\pr & S12\dg 36.5\pr \\
26 & 18 & 94\dg 01.0\pr & S12\dg 37.3\pr & 26 & 19 & 109\dg 01.1\pr & S12\dg 38.2\pr & 26 & 20 & 124\dg 01.1\pr & S12\dg 39.0\pr \\
26 & 21 & 139\dg 01.2\pr & S12\dg 39.9\pr & 26 & 22 & 154\dg 01.3\pr & S12\dg 40.7\pr & 26 & 23 & 169\dg 01.3\pr & S12\dg 41.6\pr \\
27 & 0 & 184\dg 01.4\pr & S12\dg 42.4\pr & 27 & 1 & 199\dg 01.4\pr & S12\dg 43.3\pr & 27 & 2 & 214\dg 01.5\pr & S12\dg 44.1\pr \\
27 & 3 & 229\dg 01.6\pr & S12\dg 45.0\pr & 27 & 4 & 244\dg 01.6\pr & S12\dg 45.8\pr & 27 & 5 & 259\dg 01.7\pr & S12\dg 46.7\pr \\
27 & 6 & 274\dg 01.7\pr & S12\dg 47.5\pr & 27 & 7 & 289\dg 01.8\pr & S12\dg 48.4\pr & 27 & 8 & 304\dg 01.9\pr & S12\dg 49.2\pr \\
27 & 9 & 319\dg 01.9\pr & S12\dg 50.0\pr & 27 & 10 & 334\dg 02.0\pr & S12\dg 50.9\pr & 27 & 11 & 349\dg 02.0\pr & S12\dg 51.7\pr \\
27 & 12 & 4\dg 02.1\pr & S12\dg 52.6\pr & 27 & 13 & 19\dg 02.1\pr & S12\dg 53.4\pr & 27 & 14 & 34\dg 02.2\pr & S12\dg 54.3\pr \\
27 & 15 & 49\dg 02.2\pr & S12\dg 55.1\pr & 27 & 16 & 64\dg 02.3\pr & S12\dg 55.9\pr & 27 & 17 & 79\dg 02.4\pr & S12\dg 56.8\pr \\
27 & 18 & 94\dg 02.4\pr & S12\dg 57.6\pr & 27 & 19 & 109\dg 02.5\pr & S12\dg 58.5\pr & 27 & 20 & 124\dg 02.5\pr & S12\dg 59.3\pr \\
27 & 21 & 139\dg 02.6\pr & S13\dg 00.1\pr & 27 & 22 & 154\dg 02.6\pr & S13\dg 01.0\pr & 27 & 23 & 169\dg 02.7\pr & S13\dg 01.8\pr \\
28 & 0 & 184\dg 02.7\pr & S13\dg 02.7\pr & 28 & 1 & 199\dg 02.8\pr & S13\dg 03.5\pr & 28 & 2 & 214\dg 02.8\pr & S13\dg 04.3\pr \\
28 & 3 & 229\dg 02.9\pr & S13\dg 05.2\pr & 28 & 4 & 244\dg 02.9\pr & S13\dg 06.0\pr & 28 & 5 & 259\dg 03.0\pr & S13\dg 06.9\pr \\
28 & 6 & 274\dg 03.0\pr & S13\dg 07.7\pr & 28 & 7 & 289\dg 03.1\pr & S13\dg 08.5\pr & 28 & 8 & 304\dg 03.1\pr & S13\dg 09.4\pr \\
28 & 9 & 319\dg 03.2\pr & S13\dg 10.2\pr & 28 & 10 & 334\dg 03.2\pr & S13\dg 11.0\pr & 28 & 11 & 349\dg 03.3\pr & S13\dg 11.9\pr \\
28 & 12 & 4\dg 03.3\pr & S13\dg 12.7\pr & 28 & 13 & 19\dg 03.4\pr & S13\dg 13.5\pr & 28 & 14 & 34\dg 03.4\pr & S13\dg 14.4\pr \\
28 & 15 & 49\dg 03.5\pr & S13\dg 15.2\pr & 28 & 16 & 64\dg 03.5\pr & S13\dg 16.0\pr & 28 & 17 & 79\dg 03.6\pr & S13\dg 16.9\pr \\
28 & 18 & 94\dg 03.6\pr & S13\dg 17.7\pr & 28 & 19 & 109\dg 03.7\pr & S13\dg 18.5\pr & 28 & 20 & 124\dg 03.7\pr & S13\dg 19.4\pr \\
28 & 21 & 139\dg 03.7\pr & S13\dg 20.2\pr & 28 & 22 & 154\dg 03.8\pr & S13\dg 21.0\pr & 28 & 23 & 169\dg 03.8\pr & S13\dg 21.9\pr \\
29 & 0 & 184\dg 03.9\pr & S13\dg 22.7\pr & 29 & 1 & 199\dg 03.9\pr & S13\dg 23.5\pr & 29 & 2 & 214\dg 04.0\pr & S13\dg 24.3\pr \\
29 & 3 & 229\dg 04.0\pr & S13\dg 25.2\pr & 29 & 4 & 244\dg 04.1\pr & S13\dg 26.0\pr & 29 & 5 & 259\dg 04.1\pr & S13\dg 26.8\pr \\
29 & 6 & 274\dg 04.1\pr & S13\dg 27.7\pr & 29 & 7 & 289\dg 04.2\pr & S13\dg 28.5\pr & 29 & 8 & 304\dg 04.2\pr & S13\dg 29.3\pr \\
29 & 9 & 319\dg 04.3\pr & S13\dg 30.1\pr & 29 & 10 & 334\dg 04.3\pr & S13\dg 31.0\pr & 29 & 11 & 349\dg 04.3\pr & S13\dg 31.8\pr \\
29 & 12 & 4\dg 04.4\pr & S13\dg 32.6\pr & 29 & 13 & 19\dg 04.4\pr & S13\dg 33.4\pr & 29 & 14 & 34\dg 04.5\pr & S13\dg 34.3\pr \\
29 & 15 & 49\dg 04.5\pr & S13\dg 35.1\pr & 29 & 16 & 64\dg 04.5\pr & S13\dg 35.9\pr & 29 & 17 & 79\dg 04.6\pr & S13\dg 36.7\pr \\
29 & 18 & 94\dg 04.6\pr & S13\dg 37.6\pr & 29 & 19 & 109\dg 04.7\pr & S13\dg 38.4\pr & 29 & 20 & 124\dg 04.7\pr & S13\dg 39.2\pr \\
29 & 21 & 139\dg 04.7\pr & S13\dg 40.0\pr & 29 & 22 & 154\dg 04.8\pr & S13\dg 40.9\pr & 29 & 23 & 169\dg 04.8\pr & S13\dg 41.7\pr \\
30 & 0 & 184\dg 04.8\pr & S13\dg 42.5\pr & 30 & 1 & 199\dg 04.9\pr & S13\dg 43.3\pr & 30 & 2 & 214\dg 04.9\pr & S13\dg 44.1\pr \\
30 & 3 & 229\dg 04.9\pr & S13\dg 45.0\pr & 30 & 4 & 244\dg 05.0\pr & S13\dg 45.8\pr & 30 & 5 & 259\dg 05.0\pr & S13\dg 46.6\pr \\
30 & 6 & 274\dg 05.1\pr & S13\dg 47.4\pr & 30 & 7 & 289\dg 05.1\pr & S13\dg 48.2\pr & 30 & 8 & 304\dg 05.1\pr & S13\dg 49.1\pr \\
30 & 9 & 319\dg 05.2\pr & S13\dg 49.9\pr & 30 & 10 & 334\dg 05.2\pr & S13\dg 50.7\pr & 30 & 11 & 349\dg 05.2\pr & S13\dg 51.5\pr \\
30 & 12 & 4\dg 05.3\pr & S13\dg 52.3\pr & 30 & 13 & 19\dg 05.3\pr & S13\dg 53.1\pr & 30 & 14 & 34\dg 05.3\pr & S13\dg 54.0\pr \\
30 & 15 & 49\dg 05.3\pr & S13\dg 54.8\pr & 30 & 16 & 64\dg 05.4\pr & S13\dg 55.6\pr & 30 & 17 & 79\dg 05.4\pr & S13\dg 56.4\pr \\
30 & 18 & 94\dg 05.4\pr & S13\dg 57.2\pr & 30 & 19 & 109\dg 05.5\pr & S13\dg 58.0\pr & 30 & 20 & 124\dg 05.5\pr & S13\dg 58.9\pr \\
30 & 21 & 139\dg 05.5\pr & S13\dg 59.7\pr & 30 & 22 & 154\dg 05.6\pr & S14\dg 00.5\pr & 30 & 23 & 169\dg 05.6\pr & S14\dg 01.3\pr \\
31 & 0 & 184\dg 05.6\pr & S14\dg 02.1\pr & 31 & 1 & 199\dg 05.6\pr & S14\dg 02.9\pr & 31 & 2 & 214\dg 05.7\pr & S14\dg 03.7\pr \\
31 & 3 & 229\dg 05.7\pr & S14\dg 04.5\pr & 31 & 4 & 244\dg 05.7\pr & S14\dg 05.3\pr & 31 & 5 & 259\dg 05.7\pr & S14\dg 06.2\pr \\
31 & 6 & 274\dg 05.8\pr & S14\dg 07.0\pr & 31 & 7 & 289\dg 05.8\pr & S14\dg 07.8\pr & 31 & 8 & 304\dg 05.8\pr & S14\dg 08.6\pr \\
31 & 9 & 319\dg 05.8\pr & S14\dg 09.4\pr & 31 & 10 & 334\dg 05.9\pr & S14\dg 10.2\pr & 31 & 11 & 349\dg 05.9\pr & S14\dg 11.0\pr \\
31 & 12 & 4\dg 05.9\pr & S14\dg 11.8\pr & 31 & 13 & 19\dg 05.9\pr & S14\dg 12.6\pr & 31 & 14 & 34\dg 06.0\pr & S14\dg 13.4\pr \\
31 & 15 & 49\dg 06.0\pr & S14\dg 14.2\pr & 31 & 16 & 64\dg 06.0\pr & S14\dg 15.0\pr & 31 & 17 & 79\dg 06.0\pr & S14\dg 15.8\pr \\
31 & 18 & 94\dg 06.1\pr & S14\dg 16.7\pr & 31 & 19 & 109\dg 06.1\pr & S14\dg 17.5\pr & 31 & 20 & 124\dg 06.1\pr & S14\dg 18.3\pr \\
31 & 21 & 139\dg 06.1\pr & S14\dg 19.1\pr & 31 & 22 & 154\dg 06.1\pr & S14\dg 19.9\pr & 31 & 23 & 169\dg 06.2\pr & S14\dg 20.7\pr \\
\end{longtable}}

\input{tables/sun_11.tex}
{\footnotesize
\begin{longtable}{rrrr|rrrr|rrrr}
\caption{Sun --- GHA and Declination, December 2026 (UT)}\label{sun12}\\
\toprule
\textbf{d} & \textbf{h} & \textbf{GHA} & \textbf{Dec} & \textbf{d} & \textbf{h} & \textbf{GHA} & \textbf{Dec} & \textbf{d} & \textbf{h} & \textbf{GHA} & \textbf{Dec} \\
\midrule
\endfirsthead
\multicolumn{12}{c}{\textit{Sun --- GHA and Declination, December 2026 (UT) (continued)}}\\
\toprule
\textbf{d} & \textbf{h} & \textbf{GHA} & \textbf{Dec} & \textbf{d} & \textbf{h} & \textbf{GHA} & \textbf{Dec} & \textbf{d} & \textbf{h} & \textbf{GHA} & \textbf{Dec} \\
\midrule
\endhead
\midrule\multicolumn{12}{r}{\textit{continued on next page}}\\
\endfoot
\bottomrule
\endlastfoot
1 & 0 & 182\dg 47.6\pr & S21\dg 45.8\pr & 1 & 1 & 197\dg 47.4\pr & S21\dg 46.2\pr & 1 & 2 & 212\dg 47.1\pr & S21\dg 46.6\pr \\
1 & 3 & 227\dg 46.9\pr & S21\dg 47.0\pr & 1 & 4 & 242\dg 46.7\pr & S21\dg 47.4\pr & 1 & 5 & 257\dg 46.4\pr & S21\dg 47.7\pr \\
1 & 6 & 272\dg 46.2\pr & S21\dg 48.1\pr & 1 & 7 & 287\dg 46.0\pr & S21\dg 48.5\pr & 1 & 8 & 302\dg 45.7\pr & S21\dg 48.9\pr \\
1 & 9 & 317\dg 45.5\pr & S21\dg 49.3\pr & 1 & 10 & 332\dg 45.3\pr & S21\dg 49.7\pr & 1 & 11 & 347\dg 45.0\pr & S21\dg 50.1\pr \\
1 & 12 & 2\dg 44.8\pr & S21\dg 50.4\pr & 1 & 13 & 17\dg 44.6\pr & S21\dg 50.8\pr & 1 & 14 & 32\dg 44.3\pr & S21\dg 51.2\pr \\
1 & 15 & 47\dg 44.1\pr & S21\dg 51.6\pr & 1 & 16 & 62\dg 43.9\pr & S21\dg 52.0\pr & 1 & 17 & 77\dg 43.6\pr & S21\dg 52.4\pr \\
1 & 18 & 92\dg 43.4\pr & S21\dg 52.7\pr & 1 & 19 & 107\dg 43.1\pr & S21\dg 53.1\pr & 1 & 20 & 122\dg 42.9\pr & S21\dg 53.5\pr \\
1 & 21 & 137\dg 42.7\pr & S21\dg 53.9\pr & 1 & 22 & 152\dg 42.4\pr & S21\dg 54.2\pr & 1 & 23 & 167\dg 42.2\pr & S21\dg 54.6\pr \\
2 & 0 & 182\dg 42.0\pr & S21\dg 55.0\pr & 2 & 1 & 197\dg 41.7\pr & S21\dg 55.4\pr & 2 & 2 & 212\dg 41.5\pr & S21\dg 55.7\pr \\
2 & 3 & 227\dg 41.2\pr & S21\dg 56.1\pr & 2 & 4 & 242\dg 41.0\pr & S21\dg 56.5\pr & 2 & 5 & 257\dg 40.8\pr & S21\dg 56.9\pr \\
2 & 6 & 272\dg 40.5\pr & S21\dg 57.2\pr & 2 & 7 & 287\dg 40.3\pr & S21\dg 57.6\pr & 2 & 8 & 302\dg 40.1\pr & S21\dg 58.0\pr \\
2 & 9 & 317\dg 39.8\pr & S21\dg 58.3\pr & 2 & 10 & 332\dg 39.6\pr & S21\dg 58.7\pr & 2 & 11 & 347\dg 39.3\pr & S21\dg 59.1\pr \\
2 & 12 & 2\dg 39.1\pr & S21\dg 59.4\pr & 2 & 13 & 17\dg 38.9\pr & S21\dg 59.8\pr & 2 & 14 & 32\dg 38.6\pr & S22\dg 00.2\pr \\
2 & 15 & 47\dg 38.4\pr & S22\dg 00.5\pr & 2 & 16 & 62\dg 38.1\pr & S22\dg 00.9\pr & 2 & 17 & 77\dg 37.9\pr & S22\dg 01.3\pr \\
2 & 18 & 92\dg 37.6\pr & S22\dg 01.6\pr & 2 & 19 & 107\dg 37.4\pr & S22\dg 02.0\pr & 2 & 20 & 122\dg 37.2\pr & S22\dg 02.4\pr \\
2 & 21 & 137\dg 36.9\pr & S22\dg 02.7\pr & 2 & 22 & 152\dg 36.7\pr & S22\dg 03.1\pr & 2 & 23 & 167\dg 36.4\pr & S22\dg 03.4\pr \\
3 & 0 & 182\dg 36.2\pr & S22\dg 03.8\pr & 3 & 1 & 197\dg 35.9\pr & S22\dg 04.1\pr & 3 & 2 & 212\dg 35.7\pr & S22\dg 04.5\pr \\
3 & 3 & 227\dg 35.5\pr & S22\dg 04.9\pr & 3 & 4 & 242\dg 35.2\pr & S22\dg 05.2\pr & 3 & 5 & 257\dg 35.0\pr & S22\dg 05.6\pr \\
3 & 6 & 272\dg 34.7\pr & S22\dg 05.9\pr & 3 & 7 & 287\dg 34.5\pr & S22\dg 06.3\pr & 3 & 8 & 302\dg 34.2\pr & S22\dg 06.6\pr \\
3 & 9 & 317\dg 34.0\pr & S22\dg 07.0\pr & 3 & 10 & 332\dg 33.7\pr & S22\dg 07.3\pr & 3 & 11 & 347\dg 33.5\pr & S22\dg 07.7\pr \\
3 & 12 & 2\dg 33.2\pr & S22\dg 08.0\pr & 3 & 13 & 17\dg 33.0\pr & S22\dg 08.4\pr & 3 & 14 & 32\dg 32.7\pr & S22\dg 08.7\pr \\
3 & 15 & 47\dg 32.5\pr & S22\dg 09.1\pr & 3 & 16 & 62\dg 32.2\pr & S22\dg 09.4\pr & 3 & 17 & 77\dg 32.0\pr & S22\dg 09.8\pr \\
3 & 18 & 92\dg 31.8\pr & S22\dg 10.1\pr & 3 & 19 & 107\dg 31.5\pr & S22\dg 10.4\pr & 3 & 20 & 122\dg 31.3\pr & S22\dg 10.8\pr \\
3 & 21 & 137\dg 31.0\pr & S22\dg 11.1\pr & 3 & 22 & 152\dg 30.8\pr & S22\dg 11.5\pr & 3 & 23 & 167\dg 30.5\pr & S22\dg 11.8\pr \\
4 & 0 & 182\dg 30.3\pr & S22\dg 12.2\pr & 4 & 1 & 197\dg 30.0\pr & S22\dg 12.5\pr & 4 & 2 & 212\dg 29.8\pr & S22\dg 12.8\pr \\
4 & 3 & 227\dg 29.5\pr & S22\dg 13.2\pr & 4 & 4 & 242\dg 29.3\pr & S22\dg 13.5\pr & 4 & 5 & 257\dg 29.0\pr & S22\dg 13.8\pr \\
4 & 6 & 272\dg 28.8\pr & S22\dg 14.2\pr & 4 & 7 & 287\dg 28.5\pr & S22\dg 14.5\pr & 4 & 8 & 302\dg 28.2\pr & S22\dg 14.8\pr \\
4 & 9 & 317\dg 28.0\pr & S22\dg 15.2\pr & 4 & 10 & 332\dg 27.7\pr & S22\dg 15.5\pr & 4 & 11 & 347\dg 27.5\pr & S22\dg 15.8\pr \\
4 & 12 & 2\dg 27.2\pr & S22\dg 16.2\pr & 4 & 13 & 17\dg 27.0\pr & S22\dg 16.5\pr & 4 & 14 & 32\dg 26.7\pr & S22\dg 16.8\pr \\
4 & 15 & 47\dg 26.5\pr & S22\dg 17.2\pr & 4 & 16 & 62\dg 26.2\pr & S22\dg 17.5\pr & 4 & 17 & 77\dg 26.0\pr & S22\dg 17.8\pr \\
4 & 18 & 92\dg 25.7\pr & S22\dg 18.1\pr & 4 & 19 & 107\dg 25.5\pr & S22\dg 18.5\pr & 4 & 20 & 122\dg 25.2\pr & S22\dg 18.8\pr \\
4 & 21 & 137\dg 24.9\pr & S22\dg 19.1\pr & 4 & 22 & 152\dg 24.7\pr & S22\dg 19.4\pr & 4 & 23 & 167\dg 24.4\pr & S22\dg 19.8\pr \\
5 & 0 & 182\dg 24.2\pr & S22\dg 20.1\pr & 5 & 1 & 197\dg 23.9\pr & S22\dg 20.4\pr & 5 & 2 & 212\dg 23.7\pr & S22\dg 20.7\pr \\
5 & 3 & 227\dg 23.4\pr & S22\dg 21.1\pr & 5 & 4 & 242\dg 23.2\pr & S22\dg 21.4\pr & 5 & 5 & 257\dg 22.9\pr & S22\dg 21.7\pr \\
5 & 6 & 272\dg 22.6\pr & S22\dg 22.0\pr & 5 & 7 & 287\dg 22.4\pr & S22\dg 22.3\pr & 5 & 8 & 302\dg 22.1\pr & S22\dg 22.6\pr \\
5 & 9 & 317\dg 21.9\pr & S22\dg 23.0\pr & 5 & 10 & 332\dg 21.6\pr & S22\dg 23.3\pr & 5 & 11 & 347\dg 21.3\pr & S22\dg 23.6\pr \\
5 & 12 & 2\dg 21.1\pr & S22\dg 23.9\pr & 5 & 13 & 17\dg 20.8\pr & S22\dg 24.2\pr & 5 & 14 & 32\dg 20.6\pr & S22\dg 24.5\pr \\
5 & 15 & 47\dg 20.3\pr & S22\dg 24.8\pr & 5 & 16 & 62\dg 20.1\pr & S22\dg 25.1\pr & 5 & 17 & 77\dg 19.8\pr & S22\dg 25.5\pr \\
5 & 18 & 92\dg 19.5\pr & S22\dg 25.8\pr & 5 & 19 & 107\dg 19.3\pr & S22\dg 26.1\pr & 5 & 20 & 122\dg 19.0\pr & S22\dg 26.4\pr \\
5 & 21 & 137\dg 18.8\pr & S22\dg 26.7\pr & 5 & 22 & 152\dg 18.5\pr & S22\dg 27.0\pr & 5 & 23 & 167\dg 18.2\pr & S22\dg 27.3\pr \\
6 & 0 & 182\dg 18.0\pr & S22\dg 27.6\pr & 6 & 1 & 197\dg 17.7\pr & S22\dg 27.9\pr & 6 & 2 & 212\dg 17.4\pr & S22\dg 28.2\pr \\
6 & 3 & 227\dg 17.2\pr & S22\dg 28.5\pr & 6 & 4 & 242\dg 16.9\pr & S22\dg 28.8\pr & 6 & 5 & 257\dg 16.7\pr & S22\dg 29.1\pr \\
6 & 6 & 272\dg 16.4\pr & S22\dg 29.4\pr & 6 & 7 & 287\dg 16.1\pr & S22\dg 29.7\pr & 6 & 8 & 302\dg 15.9\pr & S22\dg 30.0\pr \\
6 & 9 & 317\dg 15.6\pr & S22\dg 30.3\pr & 6 & 10 & 332\dg 15.3\pr & S22\dg 30.6\pr & 6 & 11 & 347\dg 15.1\pr & S22\dg 30.9\pr \\
6 & 12 & 2\dg 14.8\pr & S22\dg 31.2\pr & 6 & 13 & 17\dg 14.5\pr & S22\dg 31.5\pr & 6 & 14 & 32\dg 14.3\pr & S22\dg 31.8\pr \\
6 & 15 & 47\dg 14.0\pr & S22\dg 32.1\pr & 6 & 16 & 62\dg 13.7\pr & S22\dg 32.4\pr & 6 & 17 & 77\dg 13.5\pr & S22\dg 32.6\pr \\
6 & 18 & 92\dg 13.2\pr & S22\dg 32.9\pr & 6 & 19 & 107\dg 13.0\pr & S22\dg 33.2\pr & 6 & 20 & 122\dg 12.7\pr & S22\dg 33.5\pr \\
6 & 21 & 137\dg 12.4\pr & S22\dg 33.8\pr & 6 & 22 & 152\dg 12.2\pr & S22\dg 34.1\pr & 6 & 23 & 167\dg 11.9\pr & S22\dg 34.4\pr \\
7 & 0 & 182\dg 11.6\pr & S22\dg 34.7\pr & 7 & 1 & 197\dg 11.4\pr & S22\dg 34.9\pr & 7 & 2 & 212\dg 11.1\pr & S22\dg 35.2\pr \\
7 & 3 & 227\dg 10.8\pr & S22\dg 35.5\pr & 7 & 4 & 242\dg 10.5\pr & S22\dg 35.8\pr & 7 & 5 & 257\dg 10.3\pr & S22\dg 36.1\pr \\
7 & 6 & 272\dg 10.0\pr & S22\dg 36.4\pr & 7 & 7 & 287\dg 09.7\pr & S22\dg 36.6\pr & 7 & 8 & 302\dg 09.5\pr & S22\dg 36.9\pr \\
7 & 9 & 317\dg 09.2\pr & S22\dg 37.2\pr & 7 & 10 & 332\dg 08.9\pr & S22\dg 37.5\pr & 7 & 11 & 347\dg 08.7\pr & S22\dg 37.8\pr \\
7 & 12 & 2\dg 08.4\pr & S22\dg 38.0\pr & 7 & 13 & 17\dg 08.1\pr & S22\dg 38.3\pr & 7 & 14 & 32\dg 07.9\pr & S22\dg 38.6\pr \\
7 & 15 & 47\dg 07.6\pr & S22\dg 38.9\pr & 7 & 16 & 62\dg 07.3\pr & S22\dg 39.1\pr & 7 & 17 & 77\dg 07.0\pr & S22\dg 39.4\pr \\
7 & 18 & 92\dg 06.8\pr & S22\dg 39.7\pr & 7 & 19 & 107\dg 06.5\pr & S22\dg 39.9\pr & 7 & 20 & 122\dg 06.2\pr & S22\dg 40.2\pr \\
7 & 21 & 137\dg 06.0\pr & S22\dg 40.5\pr & 7 & 22 & 152\dg 05.7\pr & S22\dg 40.7\pr & 7 & 23 & 167\dg 05.4\pr & S22\dg 41.0\pr \\
8 & 0 & 182\dg 05.1\pr & S22\dg 41.3\pr & 8 & 1 & 197\dg 04.9\pr & S22\dg 41.5\pr & 8 & 2 & 212\dg 04.6\pr & S22\dg 41.8\pr \\
8 & 3 & 227\dg 04.3\pr & S22\dg 42.1\pr & 8 & 4 & 242\dg 04.1\pr & S22\dg 42.3\pr & 8 & 5 & 257\dg 03.8\pr & S22\dg 42.6\pr \\
8 & 6 & 272\dg 03.5\pr & S22\dg 42.9\pr & 8 & 7 & 287\dg 03.2\pr & S22\dg 43.1\pr & 8 & 8 & 302\dg 03.0\pr & S22\dg 43.4\pr \\
8 & 9 & 317\dg 02.7\pr & S22\dg 43.7\pr & 8 & 10 & 332\dg 02.4\pr & S22\dg 43.9\pr & 8 & 11 & 347\dg 02.1\pr & S22\dg 44.2\pr \\
8 & 12 & 2\dg 01.9\pr & S22\dg 44.4\pr & 8 & 13 & 17\dg 01.6\pr & S22\dg 44.7\pr & 8 & 14 & 32\dg 01.3\pr & S22\dg 44.9\pr \\
8 & 15 & 47\dg 01.0\pr & S22\dg 45.2\pr & 8 & 16 & 62\dg 00.8\pr & S22\dg 45.5\pr & 8 & 17 & 77\dg 00.5\pr & S22\dg 45.7\pr \\
8 & 18 & 92\dg 00.2\pr & S22\dg 46.0\pr & 8 & 19 & 106\dg 59.9\pr & S22\dg 46.2\pr & 8 & 20 & 121\dg 59.7\pr & S22\dg 46.5\pr \\
8 & 21 & 136\dg 59.4\pr & S22\dg 46.7\pr & 8 & 22 & 151\dg 59.1\pr & S22\dg 47.0\pr & 8 & 23 & 166\dg 58.8\pr & S22\dg 47.2\pr \\
9 & 0 & 181\dg 58.5\pr & S22\dg 47.5\pr & 9 & 1 & 196\dg 58.3\pr & S22\dg 47.7\pr & 9 & 2 & 211\dg 58.0\pr & S22\dg 48.0\pr \\
9 & 3 & 226\dg 57.7\pr & S22\dg 48.2\pr & 9 & 4 & 241\dg 57.4\pr & S22\dg 48.4\pr & 9 & 5 & 256\dg 57.2\pr & S22\dg 48.7\pr \\
9 & 6 & 271\dg 56.9\pr & S22\dg 48.9\pr & 9 & 7 & 286\dg 56.6\pr & S22\dg 49.2\pr & 9 & 8 & 301\dg 56.3\pr & S22\dg 49.4\pr \\
9 & 9 & 316\dg 56.0\pr & S22\dg 49.7\pr & 9 & 10 & 331\dg 55.8\pr & S22\dg 49.9\pr & 9 & 11 & 346\dg 55.5\pr & S22\dg 50.1\pr \\
9 & 12 & 1\dg 55.2\pr & S22\dg 50.4\pr & 9 & 13 & 16\dg 54.9\pr & S22\dg 50.6\pr & 9 & 14 & 31\dg 54.6\pr & S22\dg 50.9\pr \\
9 & 15 & 46\dg 54.4\pr & S22\dg 51.1\pr & 9 & 16 & 61\dg 54.1\pr & S22\dg 51.3\pr & 9 & 17 & 76\dg 53.8\pr & S22\dg 51.6\pr \\
9 & 18 & 91\dg 53.5\pr & S22\dg 51.8\pr & 9 & 19 & 106\dg 53.2\pr & S22\dg 52.0\pr & 9 & 20 & 121\dg 53.0\pr & S22\dg 52.3\pr \\
9 & 21 & 136\dg 52.7\pr & S22\dg 52.5\pr & 9 & 22 & 151\dg 52.4\pr & S22\dg 52.7\pr & 9 & 23 & 166\dg 52.1\pr & S22\dg 53.0\pr \\
10 & 0 & 181\dg 51.8\pr & S22\dg 53.2\pr & 10 & 1 & 196\dg 51.6\pr & S22\dg 53.4\pr & 10 & 2 & 211\dg 51.3\pr & S22\dg 53.7\pr \\
10 & 3 & 226\dg 51.0\pr & S22\dg 53.9\pr & 10 & 4 & 241\dg 50.7\pr & S22\dg 54.1\pr & 10 & 5 & 256\dg 50.4\pr & S22\dg 54.3\pr \\
10 & 6 & 271\dg 50.1\pr & S22\dg 54.6\pr & 10 & 7 & 286\dg 49.9\pr & S22\dg 54.8\pr & 10 & 8 & 301\dg 49.6\pr & S22\dg 55.0\pr \\
10 & 9 & 316\dg 49.3\pr & S22\dg 55.2\pr & 10 & 10 & 331\dg 49.0\pr & S22\dg 55.4\pr & 10 & 11 & 346\dg 48.7\pr & S22\dg 55.7\pr \\
10 & 12 & 1\dg 48.4\pr & S22\dg 55.9\pr & 10 & 13 & 16\dg 48.2\pr & S22\dg 56.1\pr & 10 & 14 & 31\dg 47.9\pr & S22\dg 56.3\pr \\
10 & 15 & 46\dg 47.6\pr & S22\dg 56.5\pr & 10 & 16 & 61\dg 47.3\pr & S22\dg 56.8\pr & 10 & 17 & 76\dg 47.0\pr & S22\dg 57.0\pr \\
10 & 18 & 91\dg 46.7\pr & S22\dg 57.2\pr & 10 & 19 & 106\dg 46.5\pr & S22\dg 57.4\pr & 10 & 20 & 121\dg 46.2\pr & S22\dg 57.6\pr \\
10 & 21 & 136\dg 45.9\pr & S22\dg 57.8\pr & 10 & 22 & 151\dg 45.6\pr & S22\dg 58.1\pr & 10 & 23 & 166\dg 45.3\pr & S22\dg 58.3\pr \\
11 & 0 & 181\dg 45.0\pr & S22\dg 58.5\pr & 11 & 1 & 196\dg 44.7\pr & S22\dg 58.7\pr & 11 & 2 & 211\dg 44.5\pr & S22\dg 58.9\pr \\
11 & 3 & 226\dg 44.2\pr & S22\dg 59.1\pr & 11 & 4 & 241\dg 43.9\pr & S22\dg 59.3\pr & 11 & 5 & 256\dg 43.6\pr & S22\dg 59.5\pr \\
11 & 6 & 271\dg 43.3\pr & S22\dg 59.7\pr & 11 & 7 & 286\dg 43.0\pr & S22\dg 59.9\pr & 11 & 8 & 301\dg 42.7\pr & S23\dg 00.1\pr \\
11 & 9 & 316\dg 42.4\pr & S23\dg 00.3\pr & 11 & 10 & 331\dg 42.2\pr & S23\dg 00.5\pr & 11 & 11 & 346\dg 41.9\pr & S23\dg 00.7\pr \\
11 & 12 & 1\dg 41.6\pr & S23\dg 00.9\pr & 11 & 13 & 16\dg 41.3\pr & S23\dg 01.1\pr & 11 & 14 & 31\dg 41.0\pr & S23\dg 01.3\pr \\
11 & 15 & 46\dg 40.7\pr & S23\dg 01.5\pr & 11 & 16 & 61\dg 40.4\pr & S23\dg 01.7\pr & 11 & 17 & 76\dg 40.1\pr & S23\dg 01.9\pr \\
11 & 18 & 91\dg 39.9\pr & S23\dg 02.1\pr & 11 & 19 & 106\dg 39.6\pr & S23\dg 02.3\pr & 11 & 20 & 121\dg 39.3\pr & S23\dg 02.5\pr \\
11 & 21 & 136\dg 39.0\pr & S23\dg 02.7\pr & 11 & 22 & 151\dg 38.7\pr & S23\dg 02.9\pr & 11 & 23 & 166\dg 38.4\pr & S23\dg 03.1\pr \\
12 & 0 & 181\dg 38.1\pr & S23\dg 03.3\pr & 12 & 1 & 196\dg 37.8\pr & S23\dg 03.5\pr & 12 & 2 & 211\dg 37.5\pr & S23\dg 03.7\pr \\
12 & 3 & 226\dg 37.2\pr & S23\dg 03.9\pr & 12 & 4 & 241\dg 37.0\pr & S23\dg 04.1\pr & 12 & 5 & 256\dg 36.7\pr & S23\dg 04.2\pr \\
12 & 6 & 271\dg 36.4\pr & S23\dg 04.4\pr & 12 & 7 & 286\dg 36.1\pr & S23\dg 04.6\pr & 12 & 8 & 301\dg 35.8\pr & S23\dg 04.8\pr \\
12 & 9 & 316\dg 35.5\pr & S23\dg 05.0\pr & 12 & 10 & 331\dg 35.2\pr & S23\dg 05.2\pr & 12 & 11 & 346\dg 34.9\pr & S23\dg 05.4\pr \\
12 & 12 & 1\dg 34.6\pr & S23\dg 05.5\pr & 12 & 13 & 16\dg 34.3\pr & S23\dg 05.7\pr & 12 & 14 & 31\dg 34.0\pr & S23\dg 05.9\pr \\
12 & 15 & 46\dg 33.7\pr & S23\dg 06.1\pr & 12 & 16 & 61\dg 33.5\pr & S23\dg 06.3\pr & 12 & 17 & 76\dg 33.2\pr & S23\dg 06.4\pr \\
12 & 18 & 91\dg 32.9\pr & S23\dg 06.6\pr & 12 & 19 & 106\dg 32.6\pr & S23\dg 06.8\pr & 12 & 20 & 121\dg 32.3\pr & S23\dg 07.0\pr \\
12 & 21 & 136\dg 32.0\pr & S23\dg 07.1\pr & 12 & 22 & 151\dg 31.7\pr & S23\dg 07.3\pr & 12 & 23 & 166\dg 31.4\pr & S23\dg 07.5\pr \\
13 & 0 & 181\dg 31.1\pr & S23\dg 07.7\pr & 13 & 1 & 196\dg 30.8\pr & S23\dg 07.8\pr & 13 & 2 & 211\dg 30.5\pr & S23\dg 08.0\pr \\
13 & 3 & 226\dg 30.2\pr & S23\dg 08.2\pr & 13 & 4 & 241\dg 29.9\pr & S23\dg 08.4\pr & 13 & 5 & 256\dg 29.6\pr & S23\dg 08.5\pr \\
13 & 6 & 271\dg 29.3\pr & S23\dg 08.7\pr & 13 & 7 & 286\dg 29.1\pr & S23\dg 08.9\pr & 13 & 8 & 301\dg 28.8\pr & S23\dg 09.0\pr \\
13 & 9 & 316\dg 28.5\pr & S23\dg 09.2\pr & 13 & 10 & 331\dg 28.2\pr & S23\dg 09.4\pr & 13 & 11 & 346\dg 27.9\pr & S23\dg 09.5\pr \\
13 & 12 & 1\dg 27.6\pr & S23\dg 09.7\pr & 13 & 13 & 16\dg 27.3\pr & S23\dg 09.8\pr & 13 & 14 & 31\dg 27.0\pr & S23\dg 10.0\pr \\
13 & 15 & 46\dg 26.7\pr & S23\dg 10.2\pr & 13 & 16 & 61\dg 26.4\pr & S23\dg 10.3\pr & 13 & 17 & 76\dg 26.1\pr & S23\dg 10.5\pr \\
13 & 18 & 91\dg 25.8\pr & S23\dg 10.6\pr & 13 & 19 & 106\dg 25.5\pr & S23\dg 10.8\pr & 13 & 20 & 121\dg 25.2\pr & S23\dg 11.0\pr \\
13 & 21 & 136\dg 24.9\pr & S23\dg 11.1\pr & 13 & 22 & 151\dg 24.6\pr & S23\dg 11.3\pr & 13 & 23 & 166\dg 24.3\pr & S23\dg 11.4\pr \\
14 & 0 & 181\dg 24.0\pr & S23\dg 11.6\pr & 14 & 1 & 196\dg 23.7\pr & S23\dg 11.7\pr & 14 & 2 & 211\dg 23.4\pr & S23\dg 11.9\pr \\
14 & 3 & 226\dg 23.1\pr & S23\dg 12.0\pr & 14 & 4 & 241\dg 22.8\pr & S23\dg 12.2\pr & 14 & 5 & 256\dg 22.5\pr & S23\dg 12.3\pr \\
14 & 6 & 271\dg 22.2\pr & S23\dg 12.5\pr & 14 & 7 & 286\dg 21.9\pr & S23\dg 12.6\pr & 14 & 8 & 301\dg 21.6\pr & S23\dg 12.8\pr \\
14 & 9 & 316\dg 21.4\pr & S23\dg 12.9\pr & 14 & 10 & 331\dg 21.1\pr & S23\dg 13.1\pr & 14 & 11 & 346\dg 20.8\pr & S23\dg 13.2\pr \\
14 & 12 & 1\dg 20.5\pr & S23\dg 13.4\pr & 14 & 13 & 16\dg 20.2\pr & S23\dg 13.5\pr & 14 & 14 & 31\dg 19.9\pr & S23\dg 13.6\pr \\
14 & 15 & 46\dg 19.6\pr & S23\dg 13.8\pr & 14 & 16 & 61\dg 19.3\pr & S23\dg 13.9\pr & 14 & 17 & 76\dg 19.0\pr & S23\dg 14.1\pr \\
14 & 18 & 91\dg 18.7\pr & S23\dg 14.2\pr & 14 & 19 & 106\dg 18.4\pr & S23\dg 14.3\pr & 14 & 20 & 121\dg 18.1\pr & S23\dg 14.5\pr \\
14 & 21 & 136\dg 17.8\pr & S23\dg 14.6\pr & 14 & 22 & 151\dg 17.5\pr & S23\dg 14.8\pr & 14 & 23 & 166\dg 17.2\pr & S23\dg 14.9\pr \\
15 & 0 & 181\dg 16.9\pr & S23\dg 15.0\pr & 15 & 1 & 196\dg 16.6\pr & S23\dg 15.2\pr & 15 & 2 & 211\dg 16.3\pr & S23\dg 15.3\pr \\
15 & 3 & 226\dg 16.0\pr & S23\dg 15.4\pr & 15 & 4 & 241\dg 15.7\pr & S23\dg 15.6\pr & 15 & 5 & 256\dg 15.4\pr & S23\dg 15.7\pr \\
15 & 6 & 271\dg 15.1\pr & S23\dg 15.8\pr & 15 & 7 & 286\dg 14.8\pr & S23\dg 15.9\pr & 15 & 8 & 301\dg 14.5\pr & S23\dg 16.1\pr \\
15 & 9 & 316\dg 14.2\pr & S23\dg 16.2\pr & 15 & 10 & 331\dg 13.9\pr & S23\dg 16.3\pr & 15 & 11 & 346\dg 13.6\pr & S23\dg 16.5\pr \\
15 & 12 & 1\dg 13.3\pr & S23\dg 16.6\pr & 15 & 13 & 16\dg 13.0\pr & S23\dg 16.7\pr & 15 & 14 & 31\dg 12.7\pr & S23\dg 16.8\pr \\
15 & 15 & 46\dg 12.4\pr & S23\dg 16.9\pr & 15 & 16 & 61\dg 12.1\pr & S23\dg 17.1\pr & 15 & 17 & 76\dg 11.8\pr & S23\dg 17.2\pr \\
15 & 18 & 91\dg 11.5\pr & S23\dg 17.3\pr & 15 & 19 & 106\dg 11.2\pr & S23\dg 17.4\pr & 15 & 20 & 121\dg 10.8\pr & S23\dg 17.5\pr \\
15 & 21 & 136\dg 10.5\pr & S23\dg 17.7\pr & 15 & 22 & 151\dg 10.2\pr & S23\dg 17.8\pr & 15 & 23 & 166\dg 09.9\pr & S23\dg 17.9\pr \\
16 & 0 & 181\dg 09.6\pr & S23\dg 18.0\pr & 16 & 1 & 196\dg 09.3\pr & S23\dg 18.1\pr & 16 & 2 & 211\dg 09.0\pr & S23\dg 18.2\pr \\
16 & 3 & 226\dg 08.7\pr & S23\dg 18.3\pr & 16 & 4 & 241\dg 08.4\pr & S23\dg 18.5\pr & 16 & 5 & 256\dg 08.1\pr & S23\dg 18.6\pr \\
16 & 6 & 271\dg 07.8\pr & S23\dg 18.7\pr & 16 & 7 & 286\dg 07.5\pr & S23\dg 18.8\pr & 16 & 8 & 301\dg 07.2\pr & S23\dg 18.9\pr \\
16 & 9 & 316\dg 06.9\pr & S23\dg 19.0\pr & 16 & 10 & 331\dg 06.6\pr & S23\dg 19.1\pr & 16 & 11 & 346\dg 06.3\pr & S23\dg 19.2\pr \\
16 & 12 & 1\dg 06.0\pr & S23\dg 19.3\pr & 16 & 13 & 16\dg 05.7\pr & S23\dg 19.4\pr & 16 & 14 & 31\dg 05.4\pr & S23\dg 19.5\pr \\
16 & 15 & 46\dg 05.1\pr & S23\dg 19.6\pr & 16 & 16 & 61\dg 04.8\pr & S23\dg 19.7\pr & 16 & 17 & 76\dg 04.5\pr & S23\dg 19.8\pr \\
16 & 18 & 91\dg 04.2\pr & S23\dg 19.9\pr & 16 & 19 & 106\dg 03.9\pr & S23\dg 20.0\pr & 16 & 20 & 121\dg 03.6\pr & S23\dg 20.1\pr \\
16 & 21 & 136\dg 03.3\pr & S23\dg 20.2\pr & 16 & 22 & 151\dg 03.0\pr & S23\dg 20.3\pr & 16 & 23 & 166\dg 02.7\pr & S23\dg 20.4\pr \\
17 & 0 & 181\dg 02.4\pr & S23\dg 20.5\pr & 17 & 1 & 196\dg 02.1\pr & S23\dg 20.6\pr & 17 & 2 & 211\dg 01.7\pr & S23\dg 20.7\pr \\
17 & 3 & 226\dg 01.4\pr & S23\dg 20.8\pr & 17 & 4 & 241\dg 01.1\pr & S23\dg 20.9\pr & 17 & 5 & 256\dg 00.8\pr & S23\dg 21.0\pr \\
17 & 6 & 271\dg 00.5\pr & S23\dg 21.1\pr & 17 & 7 & 286\dg 00.2\pr & S23\dg 21.2\pr & 17 & 8 & 300\dg 59.9\pr & S23\dg 21.3\pr \\
17 & 9 & 315\dg 59.6\pr & S23\dg 21.4\pr & 17 & 10 & 330\dg 59.3\pr & S23\dg 21.4\pr & 17 & 11 & 345\dg 59.0\pr & S23\dg 21.5\pr \\
17 & 12 & 0\dg 58.7\pr & S23\dg 21.6\pr & 17 & 13 & 15\dg 58.4\pr & S23\dg 21.7\pr & 17 & 14 & 30\dg 58.1\pr & S23\dg 21.8\pr \\
17 & 15 & 45\dg 57.8\pr & S23\dg 21.9\pr & 17 & 16 & 60\dg 57.5\pr & S23\dg 21.9\pr & 17 & 17 & 75\dg 57.2\pr & S23\dg 22.0\pr \\
17 & 18 & 90\dg 56.9\pr & S23\dg 22.1\pr & 17 & 19 & 105\dg 56.6\pr & S23\dg 22.2\pr & 17 & 20 & 120\dg 56.2\pr & S23\dg 22.3\pr \\
17 & 21 & 135\dg 55.9\pr & S23\dg 22.3\pr & 17 & 22 & 150\dg 55.6\pr & S23\dg 22.4\pr & 17 & 23 & 165\dg 55.3\pr & S23\dg 22.5\pr \\
18 & 0 & 180\dg 55.0\pr & S23\dg 22.6\pr & 18 & 1 & 195\dg 54.7\pr & S23\dg 22.7\pr & 18 & 2 & 210\dg 54.4\pr & S23\dg 22.7\pr \\
18 & 3 & 225\dg 54.1\pr & S23\dg 22.8\pr & 18 & 4 & 240\dg 53.8\pr & S23\dg 22.9\pr & 18 & 5 & 255\dg 53.5\pr & S23\dg 22.9\pr \\
18 & 6 & 270\dg 53.2\pr & S23\dg 23.0\pr & 18 & 7 & 285\dg 52.9\pr & S23\dg 23.1\pr & 18 & 8 & 300\dg 52.6\pr & S23\dg 23.2\pr \\
18 & 9 & 315\dg 52.3\pr & S23\dg 23.2\pr & 18 & 10 & 330\dg 52.0\pr & S23\dg 23.3\pr & 18 & 11 & 345\dg 51.6\pr & S23\dg 23.4\pr \\
18 & 12 & 0\dg 51.3\pr & S23\dg 23.4\pr & 18 & 13 & 15\dg 51.0\pr & S23\dg 23.5\pr & 18 & 14 & 30\dg 50.7\pr & S23\dg 23.6\pr \\
18 & 15 & 45\dg 50.4\pr & S23\dg 23.6\pr & 18 & 16 & 60\dg 50.1\pr & S23\dg 23.7\pr & 18 & 17 & 75\dg 49.8\pr & S23\dg 23.8\pr \\
18 & 18 & 90\dg 49.5\pr & S23\dg 23.8\pr & 18 & 19 & 105\dg 49.2\pr & S23\dg 23.9\pr & 18 & 20 & 120\dg 48.9\pr & S23\dg 23.9\pr \\
18 & 21 & 135\dg 48.6\pr & S23\dg 24.0\pr & 18 & 22 & 150\dg 48.3\pr & S23\dg 24.1\pr & 18 & 23 & 165\dg 48.0\pr & S23\dg 24.1\pr \\
19 & 0 & 180\dg 47.6\pr & S23\dg 24.2\pr & 19 & 1 & 195\dg 47.3\pr & S23\dg 24.2\pr & 19 & 2 & 210\dg 47.0\pr & S23\dg 24.3\pr \\
19 & 3 & 225\dg 46.7\pr & S23\dg 24.3\pr & 19 & 4 & 240\dg 46.4\pr & S23\dg 24.4\pr & 19 & 5 & 255\dg 46.1\pr & S23\dg 24.4\pr \\
19 & 6 & 270\dg 45.8\pr & S23\dg 24.5\pr & 19 & 7 & 285\dg 45.5\pr & S23\dg 24.5\pr & 19 & 8 & 300\dg 45.2\pr & S23\dg 24.6\pr \\
19 & 9 & 315\dg 44.9\pr & S23\dg 24.6\pr & 19 & 10 & 330\dg 44.6\pr & S23\dg 24.7\pr & 19 & 11 & 345\dg 44.2\pr & S23\dg 24.7\pr \\
19 & 12 & 0\dg 43.9\pr & S23\dg 24.8\pr & 19 & 13 & 15\dg 43.6\pr & S23\dg 24.8\pr & 19 & 14 & 30\dg 43.3\pr & S23\dg 24.9\pr \\
19 & 15 & 45\dg 43.0\pr & S23\dg 24.9\pr & 19 & 16 & 60\dg 42.7\pr & S23\dg 25.0\pr & 19 & 17 & 75\dg 42.4\pr & S23\dg 25.0\pr \\
19 & 18 & 90\dg 42.1\pr & S23\dg 25.0\pr & 19 & 19 & 105\dg 41.8\pr & S23\dg 25.1\pr & 19 & 20 & 120\dg 41.5\pr & S23\dg 25.1\pr \\
19 & 21 & 135\dg 41.2\pr & S23\dg 25.2\pr & 19 & 22 & 150\dg 40.8\pr & S23\dg 25.2\pr & 19 & 23 & 165\dg 40.5\pr & S23\dg 25.2\pr \\
20 & 0 & 180\dg 40.2\pr & S23\dg 25.3\pr & 20 & 1 & 195\dg 39.9\pr & S23\dg 25.3\pr & 20 & 2 & 210\dg 39.6\pr & S23\dg 25.4\pr \\
20 & 3 & 225\dg 39.3\pr & S23\dg 25.4\pr & 20 & 4 & 240\dg 39.0\pr & S23\dg 25.4\pr & 20 & 5 & 255\dg 38.7\pr & S23\dg 25.5\pr \\
20 & 6 & 270\dg 38.4\pr & S23\dg 25.5\pr & 20 & 7 & 285\dg 38.1\pr & S23\dg 25.5\pr & 20 & 8 & 300\dg 37.8\pr & S23\dg 25.5\pr \\
20 & 9 & 315\dg 37.4\pr & S23\dg 25.6\pr & 20 & 10 & 330\dg 37.1\pr & S23\dg 25.6\pr & 20 & 11 & 345\dg 36.8\pr & S23\dg 25.6\pr \\
20 & 12 & 0\dg 36.5\pr & S23\dg 25.7\pr & 20 & 13 & 15\dg 36.2\pr & S23\dg 25.7\pr & 20 & 14 & 30\dg 35.9\pr & S23\dg 25.7\pr \\
20 & 15 & 45\dg 35.6\pr & S23\dg 25.7\pr & 20 & 16 & 60\dg 35.3\pr & S23\dg 25.8\pr & 20 & 17 & 75\dg 35.0\pr & S23\dg 25.8\pr \\
20 & 18 & 90\dg 34.7\pr & S23\dg 25.8\pr & 20 & 19 & 105\dg 34.3\pr & S23\dg 25.8\pr & 20 & 20 & 120\dg 34.0\pr & S23\dg 25.9\pr \\
20 & 21 & 135\dg 33.7\pr & S23\dg 25.9\pr & 20 & 22 & 150\dg 33.4\pr & S23\dg 25.9\pr & 20 & 23 & 165\dg 33.1\pr & S23\dg 25.9\pr \\
21 & 0 & 180\dg 32.8\pr & S23\dg 25.9\pr & 21 & 1 & 195\dg 32.5\pr & S23\dg 25.9\pr & 21 & 2 & 210\dg 32.2\pr & S23\dg 26.0\pr \\
21 & 3 & 225\dg 31.9\pr & S23\dg 26.0\pr & 21 & 4 & 240\dg 31.6\pr & S23\dg 26.0\pr & 21 & 5 & 255\dg 31.2\pr & S23\dg 26.0\pr \\
21 & 6 & 270\dg 30.9\pr & S23\dg 26.0\pr & 21 & 7 & 285\dg 30.6\pr & S23\dg 26.0\pr & 21 & 8 & 300\dg 30.3\pr & S23\dg 26.0\pr \\
21 & 9 & 315\dg 30.0\pr & S23\dg 26.0\pr & 21 & 10 & 330\dg 29.7\pr & S23\dg 26.1\pr & 21 & 11 & 345\dg 29.4\pr & S23\dg 26.1\pr \\
21 & 12 & 0\dg 29.1\pr & S23\dg 26.1\pr & 21 & 13 & 15\dg 28.8\pr & S23\dg 26.1\pr & 21 & 14 & 30\dg 28.4\pr & S23\dg 26.1\pr \\
21 & 15 & 45\dg 28.1\pr & S23\dg 26.1\pr & 21 & 16 & 60\dg 27.8\pr & S23\dg 26.1\pr & 21 & 17 & 75\dg 27.5\pr & S23\dg 26.1\pr \\
21 & 18 & 90\dg 27.2\pr & S23\dg 26.1\pr & 21 & 19 & 105\dg 26.9\pr & S23\dg 26.1\pr & 21 & 20 & 120\dg 26.6\pr & S23\dg 26.1\pr \\
21 & 21 & 135\dg 26.3\pr & S23\dg 26.1\pr & 21 & 22 & 150\dg 26.0\pr & S23\dg 26.1\pr & 21 & 23 & 165\dg 25.7\pr & S23\dg 26.1\pr \\
22 & 0 & 180\dg 25.3\pr & S23\dg 26.1\pr & 22 & 1 & 195\dg 25.0\pr & S23\dg 26.1\pr & 22 & 2 & 210\dg 24.7\pr & S23\dg 26.1\pr \\
22 & 3 & 225\dg 24.4\pr & S23\dg 26.1\pr & 22 & 4 & 240\dg 24.1\pr & S23\dg 26.1\pr & 22 & 5 & 255\dg 23.8\pr & S23\dg 26.1\pr \\
22 & 6 & 270\dg 23.5\pr & S23\dg 26.1\pr & 22 & 7 & 285\dg 23.2\pr & S23\dg 26.1\pr & 22 & 8 & 300\dg 22.9\pr & S23\dg 26.1\pr \\
22 & 9 & 315\dg 22.5\pr & S23\dg 26.0\pr & 22 & 10 & 330\dg 22.2\pr & S23\dg 26.0\pr & 22 & 11 & 345\dg 21.9\pr & S23\dg 26.0\pr \\
22 & 12 & 0\dg 21.6\pr & S23\dg 26.0\pr & 22 & 13 & 15\dg 21.3\pr & S23\dg 26.0\pr & 22 & 14 & 30\dg 21.0\pr & S23\dg 26.0\pr \\
22 & 15 & 45\dg 20.7\pr & S23\dg 26.0\pr & 22 & 16 & 60\dg 20.4\pr & S23\dg 26.0\pr & 22 & 17 & 75\dg 20.1\pr & S23\dg 25.9\pr \\
22 & 18 & 90\dg 19.7\pr & S23\dg 25.9\pr & 22 & 19 & 105\dg 19.4\pr & S23\dg 25.9\pr & 22 & 20 & 120\dg 19.1\pr & S23\dg 25.9\pr \\
22 & 21 & 135\dg 18.8\pr & S23\dg 25.9\pr & 22 & 22 & 150\dg 18.5\pr & S23\dg 25.8\pr & 22 & 23 & 165\dg 18.2\pr & S23\dg 25.8\pr \\
23 & 0 & 180\dg 17.9\pr & S23\dg 25.8\pr & 23 & 1 & 195\dg 17.6\pr & S23\dg 25.8\pr & 23 & 2 & 210\dg 17.3\pr & S23\dg 25.8\pr \\
23 & 3 & 225\dg 16.9\pr & S23\dg 25.7\pr & 23 & 4 & 240\dg 16.6\pr & S23\dg 25.7\pr & 23 & 5 & 255\dg 16.3\pr & S23\dg 25.7\pr \\
23 & 6 & 270\dg 16.0\pr & S23\dg 25.7\pr & 23 & 7 & 285\dg 15.7\pr & S23\dg 25.6\pr & 23 & 8 & 300\dg 15.4\pr & S23\dg 25.6\pr \\
23 & 9 & 315\dg 15.1\pr & S23\dg 25.6\pr & 23 & 10 & 330\dg 14.8\pr & S23\dg 25.5\pr & 23 & 11 & 345\dg 14.5\pr & S23\dg 25.5\pr \\
23 & 12 & 0\dg 14.2\pr & S23\dg 25.5\pr & 23 & 13 & 15\dg 13.8\pr & S23\dg 25.4\pr & 23 & 14 & 30\dg 13.5\pr & S23\dg 25.4\pr \\
23 & 15 & 45\dg 13.2\pr & S23\dg 25.4\pr & 23 & 16 & 60\dg 12.9\pr & S23\dg 25.3\pr & 23 & 17 & 75\dg 12.6\pr & S23\dg 25.3\pr \\
23 & 18 & 90\dg 12.3\pr & S23\dg 25.3\pr & 23 & 19 & 105\dg 12.0\pr & S23\dg 25.2\pr & 23 & 20 & 120\dg 11.7\pr & S23\dg 25.2\pr \\
23 & 21 & 135\dg 11.4\pr & S23\dg 25.2\pr & 23 & 22 & 150\dg 11.0\pr & S23\dg 25.1\pr & 23 & 23 & 165\dg 10.7\pr & S23\dg 25.1\pr \\
24 & 0 & 180\dg 10.4\pr & S23\dg 25.0\pr & 24 & 1 & 195\dg 10.1\pr & S23\dg 25.0\pr & 24 & 2 & 210\dg 09.8\pr & S23\dg 25.0\pr \\
24 & 3 & 225\dg 09.5\pr & S23\dg 24.9\pr & 24 & 4 & 240\dg 09.2\pr & S23\dg 24.9\pr & 24 & 5 & 255\dg 08.9\pr & S23\dg 24.8\pr \\
24 & 6 & 270\dg 08.6\pr & S23\dg 24.8\pr & 24 & 7 & 285\dg 08.2\pr & S23\dg 24.7\pr & 24 & 8 & 300\dg 07.9\pr & S23\dg 24.7\pr \\
24 & 9 & 315\dg 07.6\pr & S23\dg 24.6\pr & 24 & 10 & 330\dg 07.3\pr & S23\dg 24.6\pr & 24 & 11 & 345\dg 07.0\pr & S23\dg 24.5\pr \\
24 & 12 & 0\dg 06.7\pr & S23\dg 24.5\pr & 24 & 13 & 15\dg 06.4\pr & S23\dg 24.4\pr & 24 & 14 & 30\dg 06.1\pr & S23\dg 24.4\pr \\
24 & 15 & 45\dg 05.8\pr & S23\dg 24.3\pr & 24 & 16 & 60\dg 05.5\pr & S23\dg 24.3\pr & 24 & 17 & 75\dg 05.1\pr & S23\dg 24.2\pr \\
24 & 18 & 90\dg 04.8\pr & S23\dg 24.2\pr & 24 & 19 & 105\dg 04.5\pr & S23\dg 24.1\pr & 24 & 20 & 120\dg 04.2\pr & S23\dg 24.0\pr \\
24 & 21 & 135\dg 03.9\pr & S23\dg 24.0\pr & 24 & 22 & 150\dg 03.6\pr & S23\dg 23.9\pr & 24 & 23 & 165\dg 03.3\pr & S23\dg 23.9\pr \\
25 & 0 & 180\dg 03.0\pr & S23\dg 23.8\pr & 25 & 1 & 195\dg 02.7\pr & S23\dg 23.7\pr & 25 & 2 & 210\dg 02.4\pr & S23\dg 23.7\pr \\
25 & 3 & 225\dg 02.0\pr & S23\dg 23.6\pr & 25 & 4 & 240\dg 01.7\pr & S23\dg 23.5\pr & 25 & 5 & 255\dg 01.4\pr & S23\dg 23.5\pr \\
25 & 6 & 270\dg 01.1\pr & S23\dg 23.4\pr & 25 & 7 & 285\dg 00.8\pr & S23\dg 23.3\pr & 25 & 8 & 300\dg 00.5\pr & S23\dg 23.3\pr \\
25 & 9 & 315\dg 00.2\pr & S23\dg 23.2\pr & 25 & 10 & 329\dg 59.9\pr & S23\dg 23.1\pr & 25 & 11 & 344\dg 59.6\pr & S23\dg 23.1\pr \\
25 & 12 & 359\dg 59.3\pr & S23\dg 23.0\pr & 25 & 13 & 14\dg 58.9\pr & S23\dg 22.9\pr & 25 & 14 & 29\dg 58.6\pr & S23\dg 22.9\pr \\
25 & 15 & 44\dg 58.3\pr & S23\dg 22.8\pr & 25 & 16 & 59\dg 58.0\pr & S23\dg 22.7\pr & 25 & 17 & 74\dg 57.7\pr & S23\dg 22.6\pr \\
25 & 18 & 89\dg 57.4\pr & S23\dg 22.6\pr & 25 & 19 & 104\dg 57.1\pr & S23\dg 22.5\pr & 25 & 20 & 119\dg 56.8\pr & S23\dg 22.4\pr \\
25 & 21 & 134\dg 56.5\pr & S23\dg 22.3\pr & 25 & 22 & 149\dg 56.2\pr & S23\dg 22.3\pr & 25 & 23 & 164\dg 55.8\pr & S23\dg 22.2\pr \\
26 & 0 & 179\dg 55.5\pr & S23\dg 22.1\pr & 26 & 1 & 194\dg 55.2\pr & S23\dg 22.0\pr & 26 & 2 & 209\dg 54.9\pr & S23\dg 21.9\pr \\
26 & 3 & 224\dg 54.6\pr & S23\dg 21.8\pr & 26 & 4 & 239\dg 54.3\pr & S23\dg 21.8\pr & 26 & 5 & 254\dg 54.0\pr & S23\dg 21.7\pr \\
26 & 6 & 269\dg 53.7\pr & S23\dg 21.6\pr & 26 & 7 & 284\dg 53.4\pr & S23\dg 21.5\pr & 26 & 8 & 299\dg 53.1\pr & S23\dg 21.4\pr \\
26 & 9 & 314\dg 52.8\pr & S23\dg 21.3\pr & 26 & 10 & 329\dg 52.5\pr & S23\dg 21.2\pr & 26 & 11 & 344\dg 52.1\pr & S23\dg 21.2\pr \\
26 & 12 & 359\dg 51.8\pr & S23\dg 21.1\pr & 26 & 13 & 14\dg 51.5\pr & S23\dg 21.0\pr & 26 & 14 & 29\dg 51.2\pr & S23\dg 20.9\pr \\
26 & 15 & 44\dg 50.9\pr & S23\dg 20.8\pr & 26 & 16 & 59\dg 50.6\pr & S23\dg 20.7\pr & 26 & 17 & 74\dg 50.3\pr & S23\dg 20.6\pr \\
26 & 18 & 89\dg 50.0\pr & S23\dg 20.5\pr & 26 & 19 & 104\dg 49.7\pr & S23\dg 20.4\pr & 26 & 20 & 119\dg 49.4\pr & S23\dg 20.3\pr \\
26 & 21 & 134\dg 49.1\pr & S23\dg 20.2\pr & 26 & 22 & 149\dg 48.7\pr & S23\dg 20.1\pr & 26 & 23 & 164\dg 48.4\pr & S23\dg 20.0\pr \\
27 & 0 & 179\dg 48.1\pr & S23\dg 19.9\pr & 27 & 1 & 194\dg 47.8\pr & S23\dg 19.8\pr & 27 & 2 & 209\dg 47.5\pr & S23\dg 19.7\pr \\
27 & 3 & 224\dg 47.2\pr & S23\dg 19.6\pr & 27 & 4 & 239\dg 46.9\pr & S23\dg 19.5\pr & 27 & 5 & 254\dg 46.6\pr & S23\dg 19.4\pr \\
27 & 6 & 269\dg 46.3\pr & S23\dg 19.3\pr & 27 & 7 & 284\dg 46.0\pr & S23\dg 19.2\pr & 27 & 8 & 299\dg 45.7\pr & S23\dg 19.1\pr \\
27 & 9 & 314\dg 45.4\pr & S23\dg 19.0\pr & 27 & 10 & 329\dg 45.1\pr & S23\dg 18.9\pr & 27 & 11 & 344\dg 44.8\pr & S23\dg 18.8\pr \\
27 & 12 & 359\dg 44.4\pr & S23\dg 18.7\pr & 27 & 13 & 14\dg 44.1\pr & S23\dg 18.5\pr & 27 & 14 & 29\dg 43.8\pr & S23\dg 18.4\pr \\
27 & 15 & 44\dg 43.5\pr & S23\dg 18.3\pr & 27 & 16 & 59\dg 43.2\pr & S23\dg 18.2\pr & 27 & 17 & 74\dg 42.9\pr & S23\dg 18.1\pr \\
27 & 18 & 89\dg 42.6\pr & S23\dg 18.0\pr & 27 & 19 & 104\dg 42.3\pr & S23\dg 17.9\pr & 27 & 20 & 119\dg 42.0\pr & S23\dg 17.7\pr \\
27 & 21 & 134\dg 41.7\pr & S23\dg 17.6\pr & 27 & 22 & 149\dg 41.4\pr & S23\dg 17.5\pr & 27 & 23 & 164\dg 41.1\pr & S23\dg 17.4\pr \\
28 & 0 & 179\dg 40.8\pr & S23\dg 17.3\pr & 28 & 1 & 194\dg 40.5\pr & S23\dg 17.1\pr & 28 & 2 & 209\dg 40.2\pr & S23\dg 17.0\pr \\
28 & 3 & 224\dg 39.8\pr & S23\dg 16.9\pr & 28 & 4 & 239\dg 39.5\pr & S23\dg 16.8\pr & 28 & 5 & 254\dg 39.2\pr & S23\dg 16.7\pr \\
28 & 6 & 269\dg 38.9\pr & S23\dg 16.5\pr & 28 & 7 & 284\dg 38.6\pr & S23\dg 16.4\pr & 28 & 8 & 299\dg 38.3\pr & S23\dg 16.3\pr \\
28 & 9 & 314\dg 38.0\pr & S23\dg 16.2\pr & 28 & 10 & 329\dg 37.7\pr & S23\dg 16.0\pr & 28 & 11 & 344\dg 37.4\pr & S23\dg 15.9\pr \\
28 & 12 & 359\dg 37.1\pr & S23\dg 15.8\pr & 28 & 13 & 14\dg 36.8\pr & S23\dg 15.6\pr & 28 & 14 & 29\dg 36.5\pr & S23\dg 15.5\pr \\
28 & 15 & 44\dg 36.2\pr & S23\dg 15.4\pr & 28 & 16 & 59\dg 35.9\pr & S23\dg 15.2\pr & 28 & 17 & 74\dg 35.6\pr & S23\dg 15.1\pr \\
28 & 18 & 89\dg 35.3\pr & S23\dg 15.0\pr & 28 & 19 & 104\dg 35.0\pr & S23\dg 14.8\pr & 28 & 20 & 119\dg 34.7\pr & S23\dg 14.7\pr \\
28 & 21 & 134\dg 34.3\pr & S23\dg 14.6\pr & 28 & 22 & 149\dg 34.0\pr & S23\dg 14.4\pr & 28 & 23 & 164\dg 33.7\pr & S23\dg 14.3\pr \\
29 & 0 & 179\dg 33.4\pr & S23\dg 14.2\pr & 29 & 1 & 194\dg 33.1\pr & S23\dg 14.0\pr & 29 & 2 & 209\dg 32.8\pr & S23\dg 13.9\pr \\
29 & 3 & 224\dg 32.5\pr & S23\dg 13.7\pr & 29 & 4 & 239\dg 32.2\pr & S23\dg 13.6\pr & 29 & 5 & 254\dg 31.9\pr & S23\dg 13.4\pr \\
29 & 6 & 269\dg 31.6\pr & S23\dg 13.3\pr & 29 & 7 & 284\dg 31.3\pr & S23\dg 13.2\pr & 29 & 8 & 299\dg 31.0\pr & S23\dg 13.0\pr \\
29 & 9 & 314\dg 30.7\pr & S23\dg 12.9\pr & 29 & 10 & 329\dg 30.4\pr & S23\dg 12.7\pr & 29 & 11 & 344\dg 30.1\pr & S23\dg 12.6\pr \\
29 & 12 & 359\dg 29.8\pr & S23\dg 12.4\pr & 29 & 13 & 14\dg 29.5\pr & S23\dg 12.3\pr & 29 & 14 & 29\dg 29.2\pr & S23\dg 12.1\pr \\
29 & 15 & 44\dg 28.9\pr & S23\dg 12.0\pr & 29 & 16 & 59\dg 28.6\pr & S23\dg 11.8\pr & 29 & 17 & 74\dg 28.3\pr & S23\dg 11.7\pr \\
29 & 18 & 89\dg 28.0\pr & S23\dg 11.5\pr & 29 & 19 & 104\dg 27.7\pr & S23\dg 11.4\pr & 29 & 20 & 119\dg 27.4\pr & S23\dg 11.2\pr \\
29 & 21 & 134\dg 27.1\pr & S23\dg 11.1\pr & 29 & 22 & 149\dg 26.8\pr & S23\dg 10.9\pr & 29 & 23 & 164\dg 26.5\pr & S23\dg 10.7\pr \\
30 & 0 & 179\dg 26.2\pr & S23\dg 10.6\pr & 30 & 1 & 194\dg 25.9\pr & S23\dg 10.4\pr & 30 & 2 & 209\dg 25.6\pr & S23\dg 10.3\pr \\
30 & 3 & 224\dg 25.3\pr & S23\dg 10.1\pr & 30 & 4 & 239\dg 24.9\pr & S23\dg 09.9\pr & 30 & 5 & 254\dg 24.6\pr & S23\dg 09.8\pr \\
30 & 6 & 269\dg 24.3\pr & S23\dg 09.6\pr & 30 & 7 & 284\dg 24.0\pr & S23\dg 09.4\pr & 30 & 8 & 299\dg 23.7\pr & S23\dg 09.3\pr \\
30 & 9 & 314\dg 23.4\pr & S23\dg 09.1\pr & 30 & 10 & 329\dg 23.1\pr & S23\dg 08.9\pr & 30 & 11 & 344\dg 22.8\pr & S23\dg 08.8\pr \\
30 & 12 & 359\dg 22.5\pr & S23\dg 08.6\pr & 30 & 13 & 14\dg 22.2\pr & S23\dg 08.4\pr & 30 & 14 & 29\dg 21.9\pr & S23\dg 08.3\pr \\
30 & 15 & 44\dg 21.6\pr & S23\dg 08.1\pr & 30 & 16 & 59\dg 21.3\pr & S23\dg 07.9\pr & 30 & 17 & 74\dg 21.0\pr & S23\dg 07.8\pr \\
30 & 18 & 89\dg 20.7\pr & S23\dg 07.6\pr & 30 & 19 & 104\dg 20.4\pr & S23\dg 07.4\pr & 30 & 20 & 119\dg 20.1\pr & S23\dg 07.2\pr \\
30 & 21 & 134\dg 19.8\pr & S23\dg 07.1\pr & 30 & 22 & 149\dg 19.5\pr & S23\dg 06.9\pr & 30 & 23 & 164\dg 19.2\pr & S23\dg 06.7\pr \\
31 & 0 & 179\dg 18.9\pr & S23\dg 06.5\pr & 31 & 1 & 194\dg 18.6\pr & S23\dg 06.4\pr & 31 & 2 & 209\dg 18.3\pr & S23\dg 06.2\pr \\
31 & 3 & 224\dg 18.0\pr & S23\dg 06.0\pr & 31 & 4 & 239\dg 17.7\pr & S23\dg 05.8\pr & 31 & 5 & 254\dg 17.4\pr & S23\dg 05.6\pr \\
31 & 6 & 269\dg 17.1\pr & S23\dg 05.5\pr & 31 & 7 & 284\dg 16.8\pr & S23\dg 05.3\pr & 31 & 8 & 299\dg 16.5\pr & S23\dg 05.1\pr \\
31 & 9 & 314\dg 16.3\pr & S23\dg 04.9\pr & 31 & 10 & 329\dg 16.0\pr & S23\dg 04.7\pr & 31 & 11 & 344\dg 15.7\pr & S23\dg 04.5\pr \\
31 & 12 & 359\dg 15.4\pr & S23\dg 04.3\pr & 31 & 13 & 14\dg 15.1\pr & S23\dg 04.2\pr & 31 & 14 & 29\dg 14.8\pr & S23\dg 04.0\pr \\
31 & 15 & 44\dg 14.5\pr & S23\dg 03.8\pr & 31 & 16 & 59\dg 14.2\pr & S23\dg 03.6\pr & 31 & 17 & 74\dg 13.9\pr & S23\dg 03.4\pr \\
31 & 18 & 89\dg 13.6\pr & S23\dg 03.2\pr & 31 & 19 & 104\dg 13.3\pr & S23\dg 03.0\pr & 31 & 20 & 119\dg 13.0\pr & S23\dg 02.8\pr \\
31 & 21 & 134\dg 12.7\pr & S23\dg 02.6\pr & 31 & 22 & 149\dg 12.4\pr & S23\dg 02.4\pr & 31 & 23 & 164\dg 12.1\pr & S23\dg 02.2\pr \\
\end{longtable}}


% ============================================================
% PART III --- ARIES & STARS
% ============================================================
\part{GHA Aries \& the Navigational Stars, 2026}
\chapter{The 58 navigational stars}

For each star the table gives the \textbf{Sidereal Hour Angle} (SHA, the star's
angular distance west of the First Point of Aries) and \textbf{Declination} at
epoch 2026.5, together with the visual magnitude. These are \emph{mean places} for
epoch 2026.5; precession and proper motion shift them by up to about half an
arcminute toward the ends of the year --- within this almanac's $\sim\!1\pr$ budget,
but re-generate the tables from the source engine if you need another year. To find a star's
GHA at your moment of observation, add its SHA to the GHA of Aries from the next
chapter:
\[
  \text{GHA}_\star = \text{GHA Aries} + \text{SHA}\quad(\text{reduce modulo }360\dg).
\]

{\footnotesize
\begin{longtable}{l r r r || l r r r}
\caption{The 58 navigational stars --- SHA, Declination (epoch 2026.5), and magnitude}\label{stars}\\
\toprule
\textbf{Star} & \textbf{SHA} & \textbf{Dec} & \textbf{Mag} & \textbf{Star} & \textbf{SHA} & \textbf{Dec} & \textbf{Mag}\\
\midrule
\endfirsthead
\toprule
\textbf{Star} & \textbf{SHA} & \textbf{Dec} & \textbf{Mag} & \textbf{Star} & \textbf{SHA} & \textbf{Dec} & \textbf{Mag}\\
\midrule
\endhead
\bottomrule
\endlastfoot
Acamar & 315\dg 10.9\pr & S40\dg 12.0\pr & 2.9 & Fomalhaut & 15\dg 13.4\pr & S29\dg 28.9\pr & 1.2 \\
Achernar & 335\dg 19.5\pr & S57\dg 06.2\pr & 0.5 & Gacrux & 171\dg 50.2\pr & S57\dg 15.7\pr & 1.6 \\
Acrux & 172\dg 58.6\pr & S63\dg 14.7\pr & 0.8 & Gienah & 175\dg 42.4\pr & S17\dg 41.3\pr & 2.6 \\
Adhara & 255\dg 05.0\pr & S29\dg 00.6\pr & 1.5 & Hadar & 148\dg 34.2\pr & S60\dg 30.0\pr & 0.6 \\
Aldebaran & 290\dg 38.3\pr & N16\dg 33.6\pr & 0.9 & Hamal & 327\dg 49.9\pr & N23\dg 35.2\pr & 2.0 \\
Alioth & 166\dg 12.2\pr & N55\dg 49.0\pr & 1.8 & Kaus Australis & 83\dg 31.1\pr & S34\dg 22.2\pr & 1.8 \\
Alkaid & 152\dg 51.3\pr & N49\dg 10.9\pr & 1.9 & Kochab & 137\dg 20.2\pr & N74\dg 02.8\pr & 2.1 \\
Al Na'ir & 27\dg 31.7\pr & S46\dg 49.9\pr & 1.7 & Markab & 13\dg 28.7\pr & N15\dg 20.9\pr & 2.5 \\
Alnilam & 275\dg 36.6\pr & S1\dg 11.2\pr & 1.7 & Menkar & 314\dg 05.0\pr & N4\dg 11.5\pr & 2.5 \\
Alphard & 217\dg 46.7\pr & S8\dg 46.5\pr & 2.0 & Menkent & 147\dg 56.2\pr & S36\dg 29.9\pr & 2.1 \\
Alphecca & 126\dg 02.8\pr & N26\dg 37.6\pr & 2.2 & Miaplacidus & 221\dg 37.7\pr & S69\dg 49.6\pr & 1.7 \\
Alpheratz & 357\dg 33.5\pr & N29\dg 14.2\pr & 2.1 & Mirfak & 308\dg 26.6\pr & N49\dg 57.2\pr & 1.8 \\
Altair & 61\dg 58.9\pr & N8\dg 56.4\pr & 0.8 & Nunki & 75\dg 46.4\pr & S26\dg 15.7\pr & 2.0 \\
Ankaa & 353\dg 06.2\pr & S42\dg 09.7\pr & 2.4 & Peacock & 53\dg 04.1\pr & S56\dg 38.9\pr & 1.9 \\
Antares & 112\dg 14.4\pr & S26\dg 29.3\pr & 1.1 & Polaris & 313\dg 19.3\pr & N89\dg 22.4\pr & 2.0 \\
Arcturus & 145\dg 46.9\pr & N19\dg 02.7\pr & -0.1 & Pollux & 243\dg 16.0\pr & N27\dg 57.6\pr & 1.2 \\
Atria & 107\dg 07.6\pr & S69\dg 04.3\pr & 1.9 & Procyon & 244\dg 49.7\pr & N5\dg 09.3\pr & 0.4 \\
Avior & 234\dg 14.2\pr & S59\dg 35.7\pr & 1.9 & Rasalhague & 95\dg 57.5\pr & N12\dg 32.6\pr & 2.1 \\
Bellatrix & 278\dg 21.7\pr & N6\dg 22.3\pr & 1.6 & Regulus & 207\dg 33.3\pr & N11\dg 50.2\pr & 1.4 \\
Betelgeuse & 270\dg 50.9\pr & N7\dg 24.6\pr & 0.5 & Rigel & 281\dg 02.8\pr & S8\dg 10.4\pr & 0.1 \\
Canopus & 263\dg 51.9\pr & S52\dg 42.7\pr & -0.7 & Rigil Kentaurus & 139\dg 38.6\pr & S60\dg 56.6\pr & -0.3 \\
Capella & 280\dg 20.2\pr & N46\dg 01.3\pr & 0.1 & Sabik & 102\dg 01.5\pr & S15\dg 45.3\pr & 2.4 \\
Deneb & 49\dg 25.0\pr & N45\dg 22.6\pr & 1.2 & Schedar & 349\dg 29.6\pr & N56\dg 40.9\pr & 2.2 \\
Denebola & 182\dg 23.9\pr & N14\dg 25.4\pr & 2.1 & Shaula & 96\dg 08.8\pr & S37\dg 07.2\pr & 1.6 \\
Diphda & 348\dg 46.2\pr & S17\dg 50.5\pr & 2.0 & Sirius & 258\dg 25.3\pr & S16\dg 45.3\pr & -1.5 \\
Dubhe & 193\dg 39.9\pr & N61\dg 36.5\pr & 1.8 & Spica & 158\dg 21.1\pr & S11\dg 17.9\pr & 1.0 \\
Elnath & 278\dg 00.5\pr & N28\dg 37.6\pr & 1.7 & Suhail & 222\dg 45.4\pr & S43\dg 32.4\pr & 2.2 \\
Eltanin & 90\dg 41.7\pr & N51\dg 29.2\pr & 2.2 & Vega & 80\dg 32.5\pr & N38\dg 48.6\pr & 0.0 \\
Enif & 33\dg 37.7\pr & N9\dg 59.9\pr & 2.4 & Zubenelgenubi & 136\dg 54.8\pr & S16\dg 09.0\pr & 2.8 \\
\end{longtable}}


\chapter{Hourly GHA of Aries}

The Greenwich Hour Angle of the First Point of Aries ($\gamma$) for each day and
whole hour of UT. Add the star's SHA (previous chapter) for its GHA, and the Aries
increment from Part~IV for the odd minutes and seconds. Entries run left to right,
then top to bottom, in four columns of day~/~hour~/~GHA~$\gamma$.

{\footnotesize
\begin{longtable}{rrr|rrr|rrr|rrr}
\caption{Aries --- GHA of the First Point of Aries, January 2026 (UT)}\label{aries01}\\
\toprule
\textbf{d} & \textbf{h} & \textbf{GHA} $\gamma$ & \textbf{d} & \textbf{h} & \textbf{GHA} $\gamma$ & \textbf{d} & \textbf{h} & \textbf{GHA} $\gamma$ & \textbf{d} & \textbf{h} & \textbf{GHA} $\gamma$ \\
\midrule
\endfirsthead
\multicolumn{12}{c}{\textit{Aries --- GHA of the First Point of Aries, January 2026 (UT) (continued)}}\\
\toprule
\textbf{d} & \textbf{h} & \textbf{GHA} $\gamma$ & \textbf{d} & \textbf{h} & \textbf{GHA} $\gamma$ & \textbf{d} & \textbf{h} & \textbf{GHA} $\gamma$ & \textbf{d} & \textbf{h} & \textbf{GHA} $\gamma$ \\
\midrule
\endhead
\midrule\multicolumn{12}{r}{\textit{continued on next page}}\\
\endfoot
\bottomrule
\endlastfoot
1 & 0 & 100\dg 39.6\pr & 1 & 1 & 115\dg 42.1\pr & 1 & 2 & 130\dg 44.6\pr & 1 & 3 & 145\dg 47.0\pr \\
1 & 4 & 160\dg 49.5\pr & 1 & 5 & 175\dg 52.0\pr & 1 & 6 & 190\dg 54.4\pr & 1 & 7 & 205\dg 56.9\pr \\
1 & 8 & 220\dg 59.4\pr & 1 & 9 & 236\dg 01.8\pr & 1 & 10 & 251\dg 04.3\pr & 1 & 11 & 266\dg 06.8\pr \\
1 & 12 & 281\dg 09.2\pr & 1 & 13 & 296\dg 11.7\pr & 1 & 14 & 311\dg 14.1\pr & 1 & 15 & 326\dg 16.6\pr \\
1 & 16 & 341\dg 19.1\pr & 1 & 17 & 356\dg 21.5\pr & 1 & 18 & 11\dg 24.0\pr & 1 & 19 & 26\dg 26.5\pr \\
1 & 20 & 41\dg 28.9\pr & 1 & 21 & 56\dg 31.4\pr & 1 & 22 & 71\dg 33.9\pr & 1 & 23 & 86\dg 36.3\pr \\
2 & 0 & 101\dg 38.8\pr & 2 & 1 & 116\dg 41.3\pr & 2 & 2 & 131\dg 43.7\pr & 2 & 3 & 146\dg 46.2\pr \\
2 & 4 & 161\dg 48.6\pr & 2 & 5 & 176\dg 51.1\pr & 2 & 6 & 191\dg 53.6\pr & 2 & 7 & 206\dg 56.0\pr \\
2 & 8 & 221\dg 58.5\pr & 2 & 9 & 237\dg 01.0\pr & 2 & 10 & 252\dg 03.4\pr & 2 & 11 & 267\dg 05.9\pr \\
2 & 12 & 282\dg 08.4\pr & 2 & 13 & 297\dg 10.8\pr & 2 & 14 & 312\dg 13.3\pr & 2 & 15 & 327\dg 15.7\pr \\
2 & 16 & 342\dg 18.2\pr & 2 & 17 & 357\dg 20.7\pr & 2 & 18 & 12\dg 23.1\pr & 2 & 19 & 27\dg 25.6\pr \\
2 & 20 & 42\dg 28.1\pr & 2 & 21 & 57\dg 30.5\pr & 2 & 22 & 72\dg 33.0\pr & 2 & 23 & 87\dg 35.5\pr \\
3 & 0 & 102\dg 37.9\pr & 3 & 1 & 117\dg 40.4\pr & 3 & 2 & 132\dg 42.9\pr & 3 & 3 & 147\dg 45.3\pr \\
3 & 4 & 162\dg 47.8\pr & 3 & 5 & 177\dg 50.2\pr & 3 & 6 & 192\dg 52.7\pr & 3 & 7 & 207\dg 55.2\pr \\
3 & 8 & 222\dg 57.6\pr & 3 & 9 & 238\dg 00.1\pr & 3 & 10 & 253\dg 02.6\pr & 3 & 11 & 268\dg 05.0\pr \\
3 & 12 & 283\dg 07.5\pr & 3 & 13 & 298\dg 10.0\pr & 3 & 14 & 313\dg 12.4\pr & 3 & 15 & 328\dg 14.9\pr \\
3 & 16 & 343\dg 17.4\pr & 3 & 17 & 358\dg 19.8\pr & 3 & 18 & 13\dg 22.3\pr & 3 & 19 & 28\dg 24.7\pr \\
3 & 20 & 43\dg 27.2\pr & 3 & 21 & 58\dg 29.7\pr & 3 & 22 & 73\dg 32.1\pr & 3 & 23 & 88\dg 34.6\pr \\
4 & 0 & 103\dg 37.1\pr & 4 & 1 & 118\dg 39.5\pr & 4 & 2 & 133\dg 42.0\pr & 4 & 3 & 148\dg 44.5\pr \\
4 & 4 & 163\dg 46.9\pr & 4 & 5 & 178\dg 49.4\pr & 4 & 6 & 193\dg 51.9\pr & 4 & 7 & 208\dg 54.3\pr \\
4 & 8 & 223\dg 56.8\pr & 4 & 9 & 238\dg 59.2\pr & 4 & 10 & 254\dg 01.7\pr & 4 & 11 & 269\dg 04.2\pr \\
4 & 12 & 284\dg 06.6\pr & 4 & 13 & 299\dg 09.1\pr & 4 & 14 & 314\dg 11.6\pr & 4 & 15 & 329\dg 14.0\pr \\
4 & 16 & 344\dg 16.5\pr & 4 & 17 & 359\dg 19.0\pr & 4 & 18 & 14\dg 21.4\pr & 4 & 19 & 29\dg 23.9\pr \\
4 & 20 & 44\dg 26.3\pr & 4 & 21 & 59\dg 28.8\pr & 4 & 22 & 74\dg 31.3\pr & 4 & 23 & 89\dg 33.7\pr \\
5 & 0 & 104\dg 36.2\pr & 5 & 1 & 119\dg 38.7\pr & 5 & 2 & 134\dg 41.1\pr & 5 & 3 & 149\dg 43.6\pr \\
5 & 4 & 164\dg 46.1\pr & 5 & 5 & 179\dg 48.5\pr & 5 & 6 & 194\dg 51.0\pr & 5 & 7 & 209\dg 53.5\pr \\
5 & 8 & 224\dg 55.9\pr & 5 & 9 & 239\dg 58.4\pr & 5 & 10 & 255\dg 00.8\pr & 5 & 11 & 270\dg 03.3\pr \\
5 & 12 & 285\dg 05.8\pr & 5 & 13 & 300\dg 08.2\pr & 5 & 14 & 315\dg 10.7\pr & 5 & 15 & 330\dg 13.2\pr \\
5 & 16 & 345\dg 15.6\pr & 5 & 17 & 0\dg 18.1\pr & 5 & 18 & 15\dg 20.6\pr & 5 & 19 & 30\dg 23.0\pr \\
5 & 20 & 45\dg 25.5\pr & 5 & 21 & 60\dg 28.0\pr & 5 & 22 & 75\dg 30.4\pr & 5 & 23 & 90\dg 32.9\pr \\
6 & 0 & 105\dg 35.3\pr & 6 & 1 & 120\dg 37.8\pr & 6 & 2 & 135\dg 40.3\pr & 6 & 3 & 150\dg 42.7\pr \\
6 & 4 & 165\dg 45.2\pr & 6 & 5 & 180\dg 47.7\pr & 6 & 6 & 195\dg 50.1\pr & 6 & 7 & 210\dg 52.6\pr \\
6 & 8 & 225\dg 55.1\pr & 6 & 9 & 240\dg 57.5\pr & 6 & 10 & 256\dg 00.0\pr & 6 & 11 & 271\dg 02.4\pr \\
6 & 12 & 286\dg 04.9\pr & 6 & 13 & 301\dg 07.4\pr & 6 & 14 & 316\dg 09.8\pr & 6 & 15 & 331\dg 12.3\pr \\
6 & 16 & 346\dg 14.8\pr & 6 & 17 & 1\dg 17.2\pr & 6 & 18 & 16\dg 19.7\pr & 6 & 19 & 31\dg 22.2\pr \\
6 & 20 & 46\dg 24.6\pr & 6 & 21 & 61\dg 27.1\pr & 6 & 22 & 76\dg 29.6\pr & 6 & 23 & 91\dg 32.0\pr \\
7 & 0 & 106\dg 34.5\pr & 7 & 1 & 121\dg 36.9\pr & 7 & 2 & 136\dg 39.4\pr & 7 & 3 & 151\dg 41.9\pr \\
7 & 4 & 166\dg 44.3\pr & 7 & 5 & 181\dg 46.8\pr & 7 & 6 & 196\dg 49.3\pr & 7 & 7 & 211\dg 51.7\pr \\
7 & 8 & 226\dg 54.2\pr & 7 & 9 & 241\dg 56.7\pr & 7 & 10 & 256\dg 59.1\pr & 7 & 11 & 272\dg 01.6\pr \\
7 & 12 & 287\dg 04.1\pr & 7 & 13 & 302\dg 06.5\pr & 7 & 14 & 317\dg 09.0\pr & 7 & 15 & 332\dg 11.4\pr \\
7 & 16 & 347\dg 13.9\pr & 7 & 17 & 2\dg 16.4\pr & 7 & 18 & 17\dg 18.8\pr & 7 & 19 & 32\dg 21.3\pr \\
7 & 20 & 47\dg 23.8\pr & 7 & 21 & 62\dg 26.2\pr & 7 & 22 & 77\dg 28.7\pr & 7 & 23 & 92\dg 31.2\pr \\
8 & 0 & 107\dg 33.6\pr & 8 & 1 & 122\dg 36.1\pr & 8 & 2 & 137\dg 38.5\pr & 8 & 3 & 152\dg 41.0\pr \\
8 & 4 & 167\dg 43.5\pr & 8 & 5 & 182\dg 45.9\pr & 8 & 6 & 197\dg 48.4\pr & 8 & 7 & 212\dg 50.9\pr \\
8 & 8 & 227\dg 53.3\pr & 8 & 9 & 242\dg 55.8\pr & 8 & 10 & 257\dg 58.3\pr & 8 & 11 & 273\dg 00.7\pr \\
8 & 12 & 288\dg 03.2\pr & 8 & 13 & 303\dg 05.7\pr & 8 & 14 & 318\dg 08.1\pr & 8 & 15 & 333\dg 10.6\pr \\
8 & 16 & 348\dg 13.0\pr & 8 & 17 & 3\dg 15.5\pr & 8 & 18 & 18\dg 18.0\pr & 8 & 19 & 33\dg 20.4\pr \\
8 & 20 & 48\dg 22.9\pr & 8 & 21 & 63\dg 25.4\pr & 8 & 22 & 78\dg 27.8\pr & 8 & 23 & 93\dg 30.3\pr \\
9 & 0 & 108\dg 32.8\pr & 9 & 1 & 123\dg 35.2\pr & 9 & 2 & 138\dg 37.7\pr & 9 & 3 & 153\dg 40.2\pr \\
9 & 4 & 168\dg 42.6\pr & 9 & 5 & 183\dg 45.1\pr & 9 & 6 & 198\dg 47.5\pr & 9 & 7 & 213\dg 50.0\pr \\
9 & 8 & 228\dg 52.5\pr & 9 & 9 & 243\dg 54.9\pr & 9 & 10 & 258\dg 57.4\pr & 9 & 11 & 273\dg 59.9\pr \\
9 & 12 & 289\dg 02.3\pr & 9 & 13 & 304\dg 04.8\pr & 9 & 14 & 319\dg 07.3\pr & 9 & 15 & 334\dg 09.7\pr \\
9 & 16 & 349\dg 12.2\pr & 9 & 17 & 4\dg 14.6\pr & 9 & 18 & 19\dg 17.1\pr & 9 & 19 & 34\dg 19.6\pr \\
9 & 20 & 49\dg 22.0\pr & 9 & 21 & 64\dg 24.5\pr & 9 & 22 & 79\dg 27.0\pr & 9 & 23 & 94\dg 29.4\pr \\
10 & 0 & 109\dg 31.9\pr & 10 & 1 & 124\dg 34.4\pr & 10 & 2 & 139\dg 36.8\pr & 10 & 3 & 154\dg 39.3\pr \\
10 & 4 & 169\dg 41.8\pr & 10 & 5 & 184\dg 44.2\pr & 10 & 6 & 199\dg 46.7\pr & 10 & 7 & 214\dg 49.1\pr \\
10 & 8 & 229\dg 51.6\pr & 10 & 9 & 244\dg 54.1\pr & 10 & 10 & 259\dg 56.5\pr & 10 & 11 & 274\dg 59.0\pr \\
10 & 12 & 290\dg 01.5\pr & 10 & 13 & 305\dg 03.9\pr & 10 & 14 & 320\dg 06.4\pr & 10 & 15 & 335\dg 08.9\pr \\
10 & 16 & 350\dg 11.3\pr & 10 & 17 & 5\dg 13.8\pr & 10 & 18 & 20\dg 16.3\pr & 10 & 19 & 35\dg 18.7\pr \\
10 & 20 & 50\dg 21.2\pr & 10 & 21 & 65\dg 23.6\pr & 10 & 22 & 80\dg 26.1\pr & 10 & 23 & 95\dg 28.6\pr \\
11 & 0 & 110\dg 31.0\pr & 11 & 1 & 125\dg 33.5\pr & 11 & 2 & 140\dg 36.0\pr & 11 & 3 & 155\dg 38.4\pr \\
11 & 4 & 170\dg 40.9\pr & 11 & 5 & 185\dg 43.4\pr & 11 & 6 & 200\dg 45.8\pr & 11 & 7 & 215\dg 48.3\pr \\
11 & 8 & 230\dg 50.8\pr & 11 & 9 & 245\dg 53.2\pr & 11 & 10 & 260\dg 55.7\pr & 11 & 11 & 275\dg 58.1\pr \\
11 & 12 & 291\dg 00.6\pr & 11 & 13 & 306\dg 03.1\pr & 11 & 14 & 321\dg 05.5\pr & 11 & 15 & 336\dg 08.0\pr \\
11 & 16 & 351\dg 10.5\pr & 11 & 17 & 6\dg 12.9\pr & 11 & 18 & 21\dg 15.4\pr & 11 & 19 & 36\dg 17.9\pr \\
11 & 20 & 51\dg 20.3\pr & 11 & 21 & 66\dg 22.8\pr & 11 & 22 & 81\dg 25.2\pr & 11 & 23 & 96\dg 27.7\pr \\
12 & 0 & 111\dg 30.2\pr & 12 & 1 & 126\dg 32.6\pr & 12 & 2 & 141\dg 35.1\pr & 12 & 3 & 156\dg 37.6\pr \\
12 & 4 & 171\dg 40.0\pr & 12 & 5 & 186\dg 42.5\pr & 12 & 6 & 201\dg 45.0\pr & 12 & 7 & 216\dg 47.4\pr \\
12 & 8 & 231\dg 49.9\pr & 12 & 9 & 246\dg 52.4\pr & 12 & 10 & 261\dg 54.8\pr & 12 & 11 & 276\dg 57.3\pr \\
12 & 12 & 291\dg 59.7\pr & 12 & 13 & 307\dg 02.2\pr & 12 & 14 & 322\dg 04.7\pr & 12 & 15 & 337\dg 07.1\pr \\
12 & 16 & 352\dg 09.6\pr & 12 & 17 & 7\dg 12.1\pr & 12 & 18 & 22\dg 14.5\pr & 12 & 19 & 37\dg 17.0\pr \\
12 & 20 & 52\dg 19.5\pr & 12 & 21 & 67\dg 21.9\pr & 12 & 22 & 82\dg 24.4\pr & 12 & 23 & 97\dg 26.9\pr \\
13 & 0 & 112\dg 29.3\pr & 13 & 1 & 127\dg 31.8\pr & 13 & 2 & 142\dg 34.2\pr & 13 & 3 & 157\dg 36.7\pr \\
13 & 4 & 172\dg 39.2\pr & 13 & 5 & 187\dg 41.6\pr & 13 & 6 & 202\dg 44.1\pr & 13 & 7 & 217\dg 46.6\pr \\
13 & 8 & 232\dg 49.0\pr & 13 & 9 & 247\dg 51.5\pr & 13 & 10 & 262\dg 54.0\pr & 13 & 11 & 277\dg 56.4\pr \\
13 & 12 & 292\dg 58.9\pr & 13 & 13 & 308\dg 01.3\pr & 13 & 14 & 323\dg 03.8\pr & 13 & 15 & 338\dg 06.3\pr \\
13 & 16 & 353\dg 08.7\pr & 13 & 17 & 8\dg 11.2\pr & 13 & 18 & 23\dg 13.7\pr & 13 & 19 & 38\dg 16.1\pr \\
13 & 20 & 53\dg 18.6\pr & 13 & 21 & 68\dg 21.1\pr & 13 & 22 & 83\dg 23.5\pr & 13 & 23 & 98\dg 26.0\pr \\
14 & 0 & 113\dg 28.5\pr & 14 & 1 & 128\dg 30.9\pr & 14 & 2 & 143\dg 33.4\pr & 14 & 3 & 158\dg 35.8\pr \\
14 & 4 & 173\dg 38.3\pr & 14 & 5 & 188\dg 40.8\pr & 14 & 6 & 203\dg 43.2\pr & 14 & 7 & 218\dg 45.7\pr \\
14 & 8 & 233\dg 48.2\pr & 14 & 9 & 248\dg 50.6\pr & 14 & 10 & 263\dg 53.1\pr & 14 & 11 & 278\dg 55.6\pr \\
14 & 12 & 293\dg 58.0\pr & 14 & 13 & 309\dg 00.5\pr & 14 & 14 & 324\dg 03.0\pr & 14 & 15 & 339\dg 05.4\pr \\
14 & 16 & 354\dg 07.9\pr & 14 & 17 & 9\dg 10.3\pr & 14 & 18 & 24\dg 12.8\pr & 14 & 19 & 39\dg 15.3\pr \\
14 & 20 & 54\dg 17.7\pr & 14 & 21 & 69\dg 20.2\pr & 14 & 22 & 84\dg 22.7\pr & 14 & 23 & 99\dg 25.1\pr \\
15 & 0 & 114\dg 27.6\pr & 15 & 1 & 129\dg 30.1\pr & 15 & 2 & 144\dg 32.5\pr & 15 & 3 & 159\dg 35.0\pr \\
15 & 4 & 174\dg 37.4\pr & 15 & 5 & 189\dg 39.9\pr & 15 & 6 & 204\dg 42.4\pr & 15 & 7 & 219\dg 44.8\pr \\
15 & 8 & 234\dg 47.3\pr & 15 & 9 & 249\dg 49.8\pr & 15 & 10 & 264\dg 52.2\pr & 15 & 11 & 279\dg 54.7\pr \\
15 & 12 & 294\dg 57.2\pr & 15 & 13 & 309\dg 59.6\pr & 15 & 14 & 325\dg 02.1\pr & 15 & 15 & 340\dg 04.6\pr \\
15 & 16 & 355\dg 07.0\pr & 15 & 17 & 10\dg 09.5\pr & 15 & 18 & 25\dg 11.9\pr & 15 & 19 & 40\dg 14.4\pr \\
15 & 20 & 55\dg 16.9\pr & 15 & 21 & 70\dg 19.3\pr & 15 & 22 & 85\dg 21.8\pr & 15 & 23 & 100\dg 24.3\pr \\
16 & 0 & 115\dg 26.7\pr & 16 & 1 & 130\dg 29.2\pr & 16 & 2 & 145\dg 31.7\pr & 16 & 3 & 160\dg 34.1\pr \\
16 & 4 & 175\dg 36.6\pr & 16 & 5 & 190\dg 39.1\pr & 16 & 6 & 205\dg 41.5\pr & 16 & 7 & 220\dg 44.0\pr \\
16 & 8 & 235\dg 46.4\pr & 16 & 9 & 250\dg 48.9\pr & 16 & 10 & 265\dg 51.4\pr & 16 & 11 & 280\dg 53.8\pr \\
16 & 12 & 295\dg 56.3\pr & 16 & 13 & 310\dg 58.8\pr & 16 & 14 & 326\dg 01.2\pr & 16 & 15 & 341\dg 03.7\pr \\
16 & 16 & 356\dg 06.2\pr & 16 & 17 & 11\dg 08.6\pr & 16 & 18 & 26\dg 11.1\pr & 16 & 19 & 41\dg 13.6\pr \\
16 & 20 & 56\dg 16.0\pr & 16 & 21 & 71\dg 18.5\pr & 16 & 22 & 86\dg 20.9\pr & 16 & 23 & 101\dg 23.4\pr \\
17 & 0 & 116\dg 25.9\pr & 17 & 1 & 131\dg 28.3\pr & 17 & 2 & 146\dg 30.8\pr & 17 & 3 & 161\dg 33.3\pr \\
17 & 4 & 176\dg 35.7\pr & 17 & 5 & 191\dg 38.2\pr & 17 & 6 & 206\dg 40.7\pr & 17 & 7 & 221\dg 43.1\pr \\
17 & 8 & 236\dg 45.6\pr & 17 & 9 & 251\dg 48.0\pr & 17 & 10 & 266\dg 50.5\pr & 17 & 11 & 281\dg 53.0\pr \\
17 & 12 & 296\dg 55.4\pr & 17 & 13 & 311\dg 57.9\pr & 17 & 14 & 327\dg 00.4\pr & 17 & 15 & 342\dg 02.8\pr \\
17 & 16 & 357\dg 05.3\pr & 17 & 17 & 12\dg 07.8\pr & 17 & 18 & 27\dg 10.2\pr & 17 & 19 & 42\dg 12.7\pr \\
17 & 20 & 57\dg 15.2\pr & 17 & 21 & 72\dg 17.6\pr & 17 & 22 & 87\dg 20.1\pr & 17 & 23 & 102\dg 22.5\pr \\
18 & 0 & 117\dg 25.0\pr & 18 & 1 & 132\dg 27.5\pr & 18 & 2 & 147\dg 29.9\pr & 18 & 3 & 162\dg 32.4\pr \\
18 & 4 & 177\dg 34.9\pr & 18 & 5 & 192\dg 37.3\pr & 18 & 6 & 207\dg 39.8\pr & 18 & 7 & 222\dg 42.3\pr \\
18 & 8 & 237\dg 44.7\pr & 18 & 9 & 252\dg 47.2\pr & 18 & 10 & 267\dg 49.7\pr & 18 & 11 & 282\dg 52.1\pr \\
18 & 12 & 297\dg 54.6\pr & 18 & 13 & 312\dg 57.0\pr & 18 & 14 & 327\dg 59.5\pr & 18 & 15 & 343\dg 02.0\pr \\
18 & 16 & 358\dg 04.4\pr & 18 & 17 & 13\dg 06.9\pr & 18 & 18 & 28\dg 09.4\pr & 18 & 19 & 43\dg 11.8\pr \\
18 & 20 & 58\dg 14.3\pr & 18 & 21 & 73\dg 16.8\pr & 18 & 22 & 88\dg 19.2\pr & 18 & 23 & 103\dg 21.7\pr \\
19 & 0 & 118\dg 24.1\pr & 19 & 1 & 133\dg 26.6\pr & 19 & 2 & 148\dg 29.1\pr & 19 & 3 & 163\dg 31.5\pr \\
19 & 4 & 178\dg 34.0\pr & 19 & 5 & 193\dg 36.5\pr & 19 & 6 & 208\dg 38.9\pr & 19 & 7 & 223\dg 41.4\pr \\
19 & 8 & 238\dg 43.9\pr & 19 & 9 & 253\dg 46.3\pr & 19 & 10 & 268\dg 48.8\pr & 19 & 11 & 283\dg 51.3\pr \\
19 & 12 & 298\dg 53.7\pr & 19 & 13 & 313\dg 56.2\pr & 19 & 14 & 328\dg 58.6\pr & 19 & 15 & 344\dg 01.1\pr \\
19 & 16 & 359\dg 03.6\pr & 19 & 17 & 14\dg 06.0\pr & 19 & 18 & 29\dg 08.5\pr & 19 & 19 & 44\dg 11.0\pr \\
19 & 20 & 59\dg 13.4\pr & 19 & 21 & 74\dg 15.9\pr & 19 & 22 & 89\dg 18.4\pr & 19 & 23 & 104\dg 20.8\pr \\
20 & 0 & 119\dg 23.3\pr & 20 & 1 & 134\dg 25.8\pr & 20 & 2 & 149\dg 28.2\pr & 20 & 3 & 164\dg 30.7\pr \\
20 & 4 & 179\dg 33.1\pr & 20 & 5 & 194\dg 35.6\pr & 20 & 6 & 209\dg 38.1\pr & 20 & 7 & 224\dg 40.5\pr \\
20 & 8 & 239\dg 43.0\pr & 20 & 9 & 254\dg 45.5\pr & 20 & 10 & 269\dg 47.9\pr & 20 & 11 & 284\dg 50.4\pr \\
20 & 12 & 299\dg 52.9\pr & 20 & 13 & 314\dg 55.3\pr & 20 & 14 & 329\dg 57.8\pr & 20 & 15 & 345\dg 00.2\pr \\
20 & 16 & 0\dg 02.7\pr & 20 & 17 & 15\dg 05.2\pr & 20 & 18 & 30\dg 07.6\pr & 20 & 19 & 45\dg 10.1\pr \\
20 & 20 & 60\dg 12.6\pr & 20 & 21 & 75\dg 15.0\pr & 20 & 22 & 90\dg 17.5\pr & 20 & 23 & 105\dg 20.0\pr \\
21 & 0 & 120\dg 22.4\pr & 21 & 1 & 135\dg 24.9\pr & 21 & 2 & 150\dg 27.4\pr & 21 & 3 & 165\dg 29.8\pr \\
21 & 4 & 180\dg 32.3\pr & 21 & 5 & 195\dg 34.7\pr & 21 & 6 & 210\dg 37.2\pr & 21 & 7 & 225\dg 39.7\pr \\
21 & 8 & 240\dg 42.1\pr & 21 & 9 & 255\dg 44.6\pr & 21 & 10 & 270\dg 47.1\pr & 21 & 11 & 285\dg 49.5\pr \\
21 & 12 & 300\dg 52.0\pr & 21 & 13 & 315\dg 54.5\pr & 21 & 14 & 330\dg 56.9\pr & 21 & 15 & 345\dg 59.4\pr \\
21 & 16 & 1\dg 01.9\pr & 21 & 17 & 16\dg 04.3\pr & 21 & 18 & 31\dg 06.8\pr & 21 & 19 & 46\dg 09.2\pr \\
21 & 20 & 61\dg 11.7\pr & 21 & 21 & 76\dg 14.2\pr & 21 & 22 & 91\dg 16.6\pr & 21 & 23 & 106\dg 19.1\pr \\
22 & 0 & 121\dg 21.6\pr & 22 & 1 & 136\dg 24.0\pr & 22 & 2 & 151\dg 26.5\pr & 22 & 3 & 166\dg 29.0\pr \\
22 & 4 & 181\dg 31.4\pr & 22 & 5 & 196\dg 33.9\pr & 22 & 6 & 211\dg 36.3\pr & 22 & 7 & 226\dg 38.8\pr \\
22 & 8 & 241\dg 41.3\pr & 22 & 9 & 256\dg 43.7\pr & 22 & 10 & 271\dg 46.2\pr & 22 & 11 & 286\dg 48.7\pr \\
22 & 12 & 301\dg 51.1\pr & 22 & 13 & 316\dg 53.6\pr & 22 & 14 & 331\dg 56.1\pr & 22 & 15 & 346\dg 58.5\pr \\
22 & 16 & 2\dg 01.0\pr & 22 & 17 & 17\dg 03.5\pr & 22 & 18 & 32\dg 05.9\pr & 22 & 19 & 47\dg 08.4\pr \\
22 & 20 & 62\dg 10.8\pr & 22 & 21 & 77\dg 13.3\pr & 22 & 22 & 92\dg 15.8\pr & 22 & 23 & 107\dg 18.2\pr \\
23 & 0 & 122\dg 20.7\pr & 23 & 1 & 137\dg 23.2\pr & 23 & 2 & 152\dg 25.6\pr & 23 & 3 & 167\dg 28.1\pr \\
23 & 4 & 182\dg 30.6\pr & 23 & 5 & 197\dg 33.0\pr & 23 & 6 & 212\dg 35.5\pr & 23 & 7 & 227\dg 38.0\pr \\
23 & 8 & 242\dg 40.4\pr & 23 & 9 & 257\dg 42.9\pr & 23 & 10 & 272\dg 45.3\pr & 23 & 11 & 287\dg 47.8\pr \\
23 & 12 & 302\dg 50.3\pr & 23 & 13 & 317\dg 52.7\pr & 23 & 14 & 332\dg 55.2\pr & 23 & 15 & 347\dg 57.7\pr \\
23 & 16 & 3\dg 00.1\pr & 23 & 17 & 18\dg 02.6\pr & 23 & 18 & 33\dg 05.1\pr & 23 & 19 & 48\dg 07.5\pr \\
23 & 20 & 63\dg 10.0\pr & 23 & 21 & 78\dg 12.5\pr & 23 & 22 & 93\dg 14.9\pr & 23 & 23 & 108\dg 17.4\pr \\
24 & 0 & 123\dg 19.8\pr & 24 & 1 & 138\dg 22.3\pr & 24 & 2 & 153\dg 24.8\pr & 24 & 3 & 168\dg 27.2\pr \\
24 & 4 & 183\dg 29.7\pr & 24 & 5 & 198\dg 32.2\pr & 24 & 6 & 213\dg 34.6\pr & 24 & 7 & 228\dg 37.1\pr \\
24 & 8 & 243\dg 39.6\pr & 24 & 9 & 258\dg 42.0\pr & 24 & 10 & 273\dg 44.5\pr & 24 & 11 & 288\dg 46.9\pr \\
24 & 12 & 303\dg 49.4\pr & 24 & 13 & 318\dg 51.9\pr & 24 & 14 & 333\dg 54.3\pr & 24 & 15 & 348\dg 56.8\pr \\
24 & 16 & 3\dg 59.3\pr & 24 & 17 & 19\dg 01.7\pr & 24 & 18 & 34\dg 04.2\pr & 24 & 19 & 49\dg 06.7\pr \\
24 & 20 & 64\dg 09.1\pr & 24 & 21 & 79\dg 11.6\pr & 24 & 22 & 94\dg 14.1\pr & 24 & 23 & 109\dg 16.5\pr \\
25 & 0 & 124\dg 19.0\pr & 25 & 1 & 139\dg 21.4\pr & 25 & 2 & 154\dg 23.9\pr & 25 & 3 & 169\dg 26.4\pr \\
25 & 4 & 184\dg 28.8\pr & 25 & 5 & 199\dg 31.3\pr & 25 & 6 & 214\dg 33.8\pr & 25 & 7 & 229\dg 36.2\pr \\
25 & 8 & 244\dg 38.7\pr & 25 & 9 & 259\dg 41.2\pr & 25 & 10 & 274\dg 43.6\pr & 25 & 11 & 289\dg 46.1\pr \\
25 & 12 & 304\dg 48.6\pr & 25 & 13 & 319\dg 51.0\pr & 25 & 14 & 334\dg 53.5\pr & 25 & 15 & 349\dg 55.9\pr \\
25 & 16 & 4\dg 58.4\pr & 25 & 17 & 20\dg 00.9\pr & 25 & 18 & 35\dg 03.3\pr & 25 & 19 & 50\dg 05.8\pr \\
25 & 20 & 65\dg 08.3\pr & 25 & 21 & 80\dg 10.7\pr & 25 & 22 & 95\dg 13.2\pr & 25 & 23 & 110\dg 15.7\pr \\
26 & 0 & 125\dg 18.1\pr & 26 & 1 & 140\dg 20.6\pr & 26 & 2 & 155\dg 23.0\pr & 26 & 3 & 170\dg 25.5\pr \\
26 & 4 & 185\dg 28.0\pr & 26 & 5 & 200\dg 30.4\pr & 26 & 6 & 215\dg 32.9\pr & 26 & 7 & 230\dg 35.4\pr \\
26 & 8 & 245\dg 37.8\pr & 26 & 9 & 260\dg 40.3\pr & 26 & 10 & 275\dg 42.8\pr & 26 & 11 & 290\dg 45.2\pr \\
26 & 12 & 305\dg 47.7\pr & 26 & 13 & 320\dg 50.2\pr & 26 & 14 & 335\dg 52.6\pr & 26 & 15 & 350\dg 55.1\pr \\
26 & 16 & 5\dg 57.5\pr & 26 & 17 & 21\dg 00.0\pr & 26 & 18 & 36\dg 02.5\pr & 26 & 19 & 51\dg 04.9\pr \\
26 & 20 & 66\dg 07.4\pr & 26 & 21 & 81\dg 09.9\pr & 26 & 22 & 96\dg 12.3\pr & 26 & 23 & 111\dg 14.8\pr \\
27 & 0 & 126\dg 17.3\pr & 27 & 1 & 141\dg 19.7\pr & 27 & 2 & 156\dg 22.2\pr & 27 & 3 & 171\dg 24.7\pr \\
27 & 4 & 186\dg 27.1\pr & 27 & 5 & 201\dg 29.6\pr & 27 & 6 & 216\dg 32.0\pr & 27 & 7 & 231\dg 34.5\pr \\
27 & 8 & 246\dg 37.0\pr & 27 & 9 & 261\dg 39.4\pr & 27 & 10 & 276\dg 41.9\pr & 27 & 11 & 291\dg 44.4\pr \\
27 & 12 & 306\dg 46.8\pr & 27 & 13 & 321\dg 49.3\pr & 27 & 14 & 336\dg 51.8\pr & 27 & 15 & 351\dg 54.2\pr \\
27 & 16 & 6\dg 56.7\pr & 27 & 17 & 21\dg 59.1\pr & 27 & 18 & 37\dg 01.6\pr & 27 & 19 & 52\dg 04.1\pr \\
27 & 20 & 67\dg 06.5\pr & 27 & 21 & 82\dg 09.0\pr & 27 & 22 & 97\dg 11.5\pr & 27 & 23 & 112\dg 13.9\pr \\
28 & 0 & 127\dg 16.4\pr & 28 & 1 & 142\dg 18.9\pr & 28 & 2 & 157\dg 21.3\pr & 28 & 3 & 172\dg 23.8\pr \\
28 & 4 & 187\dg 26.3\pr & 28 & 5 & 202\dg 28.7\pr & 28 & 6 & 217\dg 31.2\pr & 28 & 7 & 232\dg 33.6\pr \\
28 & 8 & 247\dg 36.1\pr & 28 & 9 & 262\dg 38.6\pr & 28 & 10 & 277\dg 41.0\pr & 28 & 11 & 292\dg 43.5\pr \\
28 & 12 & 307\dg 46.0\pr & 28 & 13 & 322\dg 48.4\pr & 28 & 14 & 337\dg 50.9\pr & 28 & 15 & 352\dg 53.4\pr \\
28 & 16 & 7\dg 55.8\pr & 28 & 17 & 22\dg 58.3\pr & 28 & 18 & 38\dg 00.8\pr & 28 & 19 & 53\dg 03.2\pr \\
28 & 20 & 68\dg 05.7\pr & 28 & 21 & 83\dg 08.1\pr & 28 & 22 & 98\dg 10.6\pr & 28 & 23 & 113\dg 13.1\pr \\
29 & 0 & 128\dg 15.5\pr & 29 & 1 & 143\dg 18.0\pr & 29 & 2 & 158\dg 20.5\pr & 29 & 3 & 173\dg 22.9\pr \\
29 & 4 & 188\dg 25.4\pr & 29 & 5 & 203\dg 27.9\pr & 29 & 6 & 218\dg 30.3\pr & 29 & 7 & 233\dg 32.8\pr \\
29 & 8 & 248\dg 35.2\pr & 29 & 9 & 263\dg 37.7\pr & 29 & 10 & 278\dg 40.2\pr & 29 & 11 & 293\dg 42.6\pr \\
29 & 12 & 308\dg 45.1\pr & 29 & 13 & 323\dg 47.6\pr & 29 & 14 & 338\dg 50.0\pr & 29 & 15 & 353\dg 52.5\pr \\
29 & 16 & 8\dg 55.0\pr & 29 & 17 & 23\dg 57.4\pr & 29 & 18 & 38\dg 59.9\pr & 29 & 19 & 54\dg 02.4\pr \\
29 & 20 & 69\dg 04.8\pr & 29 & 21 & 84\dg 07.3\pr & 29 & 22 & 99\dg 09.7\pr & 29 & 23 & 114\dg 12.2\pr \\
30 & 0 & 129\dg 14.7\pr & 30 & 1 & 144\dg 17.1\pr & 30 & 2 & 159\dg 19.6\pr & 30 & 3 & 174\dg 22.1\pr \\
30 & 4 & 189\dg 24.5\pr & 30 & 5 & 204\dg 27.0\pr & 30 & 6 & 219\dg 29.5\pr & 30 & 7 & 234\dg 31.9\pr \\
30 & 8 & 249\dg 34.4\pr & 30 & 9 & 264\dg 36.9\pr & 30 & 10 & 279\dg 39.3\pr & 30 & 11 & 294\dg 41.8\pr \\
30 & 12 & 309\dg 44.2\pr & 30 & 13 & 324\dg 46.7\pr & 30 & 14 & 339\dg 49.2\pr & 30 & 15 & 354\dg 51.6\pr \\
30 & 16 & 9\dg 54.1\pr & 30 & 17 & 24\dg 56.6\pr & 30 & 18 & 39\dg 59.0\pr & 30 & 19 & 55\dg 01.5\pr \\
30 & 20 & 70\dg 04.0\pr & 30 & 21 & 85\dg 06.4\pr & 30 & 22 & 100\dg 08.9\pr & 30 & 23 & 115\dg 11.4\pr \\
31 & 0 & 130\dg 13.8\pr & 31 & 1 & 145\dg 16.3\pr & 31 & 2 & 160\dg 18.7\pr & 31 & 3 & 175\dg 21.2\pr \\
31 & 4 & 190\dg 23.7\pr & 31 & 5 & 205\dg 26.1\pr & 31 & 6 & 220\dg 28.6\pr & 31 & 7 & 235\dg 31.1\pr \\
31 & 8 & 250\dg 33.5\pr & 31 & 9 & 265\dg 36.0\pr & 31 & 10 & 280\dg 38.5\pr & 31 & 11 & 295\dg 40.9\pr \\
31 & 12 & 310\dg 43.4\pr & 31 & 13 & 325\dg 45.8\pr & 31 & 14 & 340\dg 48.3\pr & 31 & 15 & 355\dg 50.8\pr \\
31 & 16 & 10\dg 53.2\pr & 31 & 17 & 25\dg 55.7\pr & 31 & 18 & 40\dg 58.2\pr & 31 & 19 & 56\dg 00.6\pr \\
31 & 20 & 71\dg 03.1\pr & 31 & 21 & 86\dg 05.6\pr & 31 & 22 & 101\dg 08.0\pr & 31 & 23 & 116\dg 10.5\pr \\
\end{longtable}}

{\footnotesize
\begin{longtable}{rrr|rrr|rrr|rrr}
\caption{Aries --- GHA of the First Point of Aries, February 2026 (UT)}\label{aries02}\\
\toprule
\textbf{d} & \textbf{h} & \textbf{GHA} $\gamma$ & \textbf{d} & \textbf{h} & \textbf{GHA} $\gamma$ & \textbf{d} & \textbf{h} & \textbf{GHA} $\gamma$ & \textbf{d} & \textbf{h} & \textbf{GHA} $\gamma$ \\
\midrule
\endfirsthead
\multicolumn{12}{c}{\textit{Aries --- GHA of the First Point of Aries, February 2026 (UT) (continued)}}\\
\toprule
\textbf{d} & \textbf{h} & \textbf{GHA} $\gamma$ & \textbf{d} & \textbf{h} & \textbf{GHA} $\gamma$ & \textbf{d} & \textbf{h} & \textbf{GHA} $\gamma$ & \textbf{d} & \textbf{h} & \textbf{GHA} $\gamma$ \\
\midrule
\endhead
\midrule\multicolumn{12}{r}{\textit{continued on next page}}\\
\endfoot
\bottomrule
\endlastfoot
1 & 0 & 131\dg 13.0\pr & 1 & 1 & 146\dg 15.4\pr & 1 & 2 & 161\dg 17.9\pr & 1 & 3 & 176\dg 20.3\pr \\
1 & 4 & 191\dg 22.8\pr & 1 & 5 & 206\dg 25.3\pr & 1 & 6 & 221\dg 27.7\pr & 1 & 7 & 236\dg 30.2\pr \\
1 & 8 & 251\dg 32.7\pr & 1 & 9 & 266\dg 35.1\pr & 1 & 10 & 281\dg 37.6\pr & 1 & 11 & 296\dg 40.1\pr \\
1 & 12 & 311\dg 42.5\pr & 1 & 13 & 326\dg 45.0\pr & 1 & 14 & 341\dg 47.5\pr & 1 & 15 & 356\dg 49.9\pr \\
1 & 16 & 11\dg 52.4\pr & 1 & 17 & 26\dg 54.8\pr & 1 & 18 & 41\dg 57.3\pr & 1 & 19 & 56\dg 59.8\pr \\
1 & 20 & 72\dg 02.2\pr & 1 & 21 & 87\dg 04.7\pr & 1 & 22 & 102\dg 07.2\pr & 1 & 23 & 117\dg 09.6\pr \\
2 & 0 & 132\dg 12.1\pr & 2 & 1 & 147\dg 14.6\pr & 2 & 2 & 162\dg 17.0\pr & 2 & 3 & 177\dg 19.5\pr \\
2 & 4 & 192\dg 21.9\pr & 2 & 5 & 207\dg 24.4\pr & 2 & 6 & 222\dg 26.9\pr & 2 & 7 & 237\dg 29.3\pr \\
2 & 8 & 252\dg 31.8\pr & 2 & 9 & 267\dg 34.3\pr & 2 & 10 & 282\dg 36.7\pr & 2 & 11 & 297\dg 39.2\pr \\
2 & 12 & 312\dg 41.7\pr & 2 & 13 & 327\dg 44.1\pr & 2 & 14 & 342\dg 46.6\pr & 2 & 15 & 357\dg 49.1\pr \\
2 & 16 & 12\dg 51.5\pr & 2 & 17 & 27\dg 54.0\pr & 2 & 18 & 42\dg 56.4\pr & 2 & 19 & 57\dg 58.9\pr \\
2 & 20 & 73\dg 01.4\pr & 2 & 21 & 88\dg 03.8\pr & 2 & 22 & 103\dg 06.3\pr & 2 & 23 & 118\dg 08.8\pr \\
3 & 0 & 133\dg 11.2\pr & 3 & 1 & 148\dg 13.7\pr & 3 & 2 & 163\dg 16.2\pr & 3 & 3 & 178\dg 18.6\pr \\
3 & 4 & 193\dg 21.1\pr & 3 & 5 & 208\dg 23.6\pr & 3 & 6 & 223\dg 26.0\pr & 3 & 7 & 238\dg 28.5\pr \\
3 & 8 & 253\dg 30.9\pr & 3 & 9 & 268\dg 33.4\pr & 3 & 10 & 283\dg 35.9\pr & 3 & 11 & 298\dg 38.3\pr \\
3 & 12 & 313\dg 40.8\pr & 3 & 13 & 328\dg 43.3\pr & 3 & 14 & 343\dg 45.7\pr & 3 & 15 & 358\dg 48.2\pr \\
3 & 16 & 13\dg 50.7\pr & 3 & 17 & 28\dg 53.1\pr & 3 & 18 & 43\dg 55.6\pr & 3 & 19 & 58\dg 58.0\pr \\
3 & 20 & 74\dg 00.5\pr & 3 & 21 & 89\dg 03.0\pr & 3 & 22 & 104\dg 05.4\pr & 3 & 23 & 119\dg 07.9\pr \\
4 & 0 & 134\dg 10.4\pr & 4 & 1 & 149\dg 12.8\pr & 4 & 2 & 164\dg 15.3\pr & 4 & 3 & 179\dg 17.8\pr \\
4 & 4 & 194\dg 20.2\pr & 4 & 5 & 209\dg 22.7\pr & 4 & 6 & 224\dg 25.2\pr & 4 & 7 & 239\dg 27.6\pr \\
4 & 8 & 254\dg 30.1\pr & 4 & 9 & 269\dg 32.5\pr & 4 & 10 & 284\dg 35.0\pr & 4 & 11 & 299\dg 37.5\pr \\
4 & 12 & 314\dg 39.9\pr & 4 & 13 & 329\dg 42.4\pr & 4 & 14 & 344\dg 44.9\pr & 4 & 15 & 359\dg 47.3\pr \\
4 & 16 & 14\dg 49.8\pr & 4 & 17 & 29\dg 52.3\pr & 4 & 18 & 44\dg 54.7\pr & 4 & 19 & 59\dg 57.2\pr \\
4 & 20 & 74\dg 59.7\pr & 4 & 21 & 90\dg 02.1\pr & 4 & 22 & 105\dg 04.6\pr & 4 & 23 & 120\dg 07.0\pr \\
5 & 0 & 135\dg 09.5\pr & 5 & 1 & 150\dg 12.0\pr & 5 & 2 & 165\dg 14.4\pr & 5 & 3 & 180\dg 16.9\pr \\
5 & 4 & 195\dg 19.4\pr & 5 & 5 & 210\dg 21.8\pr & 5 & 6 & 225\dg 24.3\pr & 5 & 7 & 240\dg 26.8\pr \\
5 & 8 & 255\dg 29.2\pr & 5 & 9 & 270\dg 31.7\pr & 5 & 10 & 285\dg 34.1\pr & 5 & 11 & 300\dg 36.6\pr \\
5 & 12 & 315\dg 39.1\pr & 5 & 13 & 330\dg 41.5\pr & 5 & 14 & 345\dg 44.0\pr & 5 & 15 & 0\dg 46.5\pr \\
5 & 16 & 15\dg 48.9\pr & 5 & 17 & 30\dg 51.4\pr & 5 & 18 & 45\dg 53.9\pr & 5 & 19 & 60\dg 56.3\pr \\
5 & 20 & 75\dg 58.8\pr & 5 & 21 & 91\dg 01.3\pr & 5 & 22 & 106\dg 03.7\pr & 5 & 23 & 121\dg 06.2\pr \\
6 & 0 & 136\dg 08.6\pr & 6 & 1 & 151\dg 11.1\pr & 6 & 2 & 166\dg 13.6\pr & 6 & 3 & 181\dg 16.0\pr \\
6 & 4 & 196\dg 18.5\pr & 6 & 5 & 211\dg 21.0\pr & 6 & 6 & 226\dg 23.4\pr & 6 & 7 & 241\dg 25.9\pr \\
6 & 8 & 256\dg 28.4\pr & 6 & 9 & 271\dg 30.8\pr & 6 & 10 & 286\dg 33.3\pr & 6 & 11 & 301\dg 35.8\pr \\
6 & 12 & 316\dg 38.2\pr & 6 & 13 & 331\dg 40.7\pr & 6 & 14 & 346\dg 43.1\pr & 6 & 15 & 1\dg 45.6\pr \\
6 & 16 & 16\dg 48.1\pr & 6 & 17 & 31\dg 50.5\pr & 6 & 18 & 46\dg 53.0\pr & 6 & 19 & 61\dg 55.5\pr \\
6 & 20 & 76\dg 57.9\pr & 6 & 21 & 92\dg 00.4\pr & 6 & 22 & 107\dg 02.9\pr & 6 & 23 & 122\dg 05.3\pr \\
7 & 0 & 137\dg 07.8\pr & 7 & 1 & 152\dg 10.3\pr & 7 & 2 & 167\dg 12.7\pr & 7 & 3 & 182\dg 15.2\pr \\
7 & 4 & 197\dg 17.6\pr & 7 & 5 & 212\dg 20.1\pr & 7 & 6 & 227\dg 22.6\pr & 7 & 7 & 242\dg 25.0\pr \\
7 & 8 & 257\dg 27.5\pr & 7 & 9 & 272\dg 30.0\pr & 7 & 10 & 287\dg 32.4\pr & 7 & 11 & 302\dg 34.9\pr \\
7 & 12 & 317\dg 37.4\pr & 7 & 13 & 332\dg 39.8\pr & 7 & 14 & 347\dg 42.3\pr & 7 & 15 & 2\dg 44.7\pr \\
7 & 16 & 17\dg 47.2\pr & 7 & 17 & 32\dg 49.7\pr & 7 & 18 & 47\dg 52.1\pr & 7 & 19 & 62\dg 54.6\pr \\
7 & 20 & 77\dg 57.1\pr & 7 & 21 & 92\dg 59.5\pr & 7 & 22 & 108\dg 02.0\pr & 7 & 23 & 123\dg 04.5\pr \\
8 & 0 & 138\dg 06.9\pr & 8 & 1 & 153\dg 09.4\pr & 8 & 2 & 168\dg 11.9\pr & 8 & 3 & 183\dg 14.3\pr \\
8 & 4 & 198\dg 16.8\pr & 8 & 5 & 213\dg 19.2\pr & 8 & 6 & 228\dg 21.7\pr & 8 & 7 & 243\dg 24.2\pr \\
8 & 8 & 258\dg 26.6\pr & 8 & 9 & 273\dg 29.1\pr & 8 & 10 & 288\dg 31.6\pr & 8 & 11 & 303\dg 34.0\pr \\
8 & 12 & 318\dg 36.5\pr & 8 & 13 & 333\dg 39.0\pr & 8 & 14 & 348\dg 41.4\pr & 8 & 15 & 3\dg 43.9\pr \\
8 & 16 & 18\dg 46.4\pr & 8 & 17 & 33\dg 48.8\pr & 8 & 18 & 48\dg 51.3\pr & 8 & 19 & 63\dg 53.7\pr \\
8 & 20 & 78\dg 56.2\pr & 8 & 21 & 93\dg 58.7\pr & 8 & 22 & 109\dg 01.1\pr & 8 & 23 & 124\dg 03.6\pr \\
9 & 0 & 139\dg 06.1\pr & 9 & 1 & 154\dg 08.5\pr & 9 & 2 & 169\dg 11.0\pr & 9 & 3 & 184\dg 13.5\pr \\
9 & 4 & 199\dg 15.9\pr & 9 & 5 & 214\dg 18.4\pr & 9 & 6 & 229\dg 20.8\pr & 9 & 7 & 244\dg 23.3\pr \\
9 & 8 & 259\dg 25.8\pr & 9 & 9 & 274\dg 28.2\pr & 9 & 10 & 289\dg 30.7\pr & 9 & 11 & 304\dg 33.2\pr \\
9 & 12 & 319\dg 35.6\pr & 9 & 13 & 334\dg 38.1\pr & 9 & 14 & 349\dg 40.6\pr & 9 & 15 & 4\dg 43.0\pr \\
9 & 16 & 19\dg 45.5\pr & 9 & 17 & 34\dg 48.0\pr & 9 & 18 & 49\dg 50.4\pr & 9 & 19 & 64\dg 52.9\pr \\
9 & 20 & 79\dg 55.3\pr & 9 & 21 & 94\dg 57.8\pr & 9 & 22 & 110\dg 00.3\pr & 9 & 23 & 125\dg 02.7\pr \\
10 & 0 & 140\dg 05.2\pr & 10 & 1 & 155\dg 07.7\pr & 10 & 2 & 170\dg 10.1\pr & 10 & 3 & 185\dg 12.6\pr \\
10 & 4 & 200\dg 15.1\pr & 10 & 5 & 215\dg 17.5\pr & 10 & 6 & 230\dg 20.0\pr & 10 & 7 & 245\dg 22.5\pr \\
10 & 8 & 260\dg 24.9\pr & 10 & 9 & 275\dg 27.4\pr & 10 & 10 & 290\dg 29.8\pr & 10 & 11 & 305\dg 32.3\pr \\
10 & 12 & 320\dg 34.8\pr & 10 & 13 & 335\dg 37.2\pr & 10 & 14 & 350\dg 39.7\pr & 10 & 15 & 5\dg 42.2\pr \\
10 & 16 & 20\dg 44.6\pr & 10 & 17 & 35\dg 47.1\pr & 10 & 18 & 50\dg 49.6\pr & 10 & 19 & 65\dg 52.0\pr \\
10 & 20 & 80\dg 54.5\pr & 10 & 21 & 95\dg 56.9\pr & 10 & 22 & 110\dg 59.4\pr & 10 & 23 & 126\dg 01.9\pr \\
11 & 0 & 141\dg 04.3\pr & 11 & 1 & 156\dg 06.8\pr & 11 & 2 & 171\dg 09.3\pr & 11 & 3 & 186\dg 11.7\pr \\
11 & 4 & 201\dg 14.2\pr & 11 & 5 & 216\dg 16.7\pr & 11 & 6 & 231\dg 19.1\pr & 11 & 7 & 246\dg 21.6\pr \\
11 & 8 & 261\dg 24.1\pr & 11 & 9 & 276\dg 26.5\pr & 11 & 10 & 291\dg 29.0\pr & 11 & 11 & 306\dg 31.4\pr \\
11 & 12 & 321\dg 33.9\pr & 11 & 13 & 336\dg 36.4\pr & 11 & 14 & 351\dg 38.8\pr & 11 & 15 & 6\dg 41.3\pr \\
11 & 16 & 21\dg 43.8\pr & 11 & 17 & 36\dg 46.2\pr & 11 & 18 & 51\dg 48.7\pr & 11 & 19 & 66\dg 51.2\pr \\
11 & 20 & 81\dg 53.6\pr & 11 & 21 & 96\dg 56.1\pr & 11 & 22 & 111\dg 58.6\pr & 11 & 23 & 127\dg 01.0\pr \\
12 & 0 & 142\dg 03.5\pr & 12 & 1 & 157\dg 05.9\pr & 12 & 2 & 172\dg 08.4\pr & 12 & 3 & 187\dg 10.9\pr \\
12 & 4 & 202\dg 13.3\pr & 12 & 5 & 217\dg 15.8\pr & 12 & 6 & 232\dg 18.3\pr & 12 & 7 & 247\dg 20.7\pr \\
12 & 8 & 262\dg 23.2\pr & 12 & 9 & 277\dg 25.7\pr & 12 & 10 & 292\dg 28.1\pr & 12 & 11 & 307\dg 30.6\pr \\
12 & 12 & 322\dg 33.0\pr & 12 & 13 & 337\dg 35.5\pr & 12 & 14 & 352\dg 38.0\pr & 12 & 15 & 7\dg 40.4\pr \\
12 & 16 & 22\dg 42.9\pr & 12 & 17 & 37\dg 45.4\pr & 12 & 18 & 52\dg 47.8\pr & 12 & 19 & 67\dg 50.3\pr \\
12 & 20 & 82\dg 52.8\pr & 12 & 21 & 97\dg 55.2\pr & 12 & 22 & 112\dg 57.7\pr & 12 & 23 & 128\dg 00.2\pr \\
13 & 0 & 143\dg 02.6\pr & 13 & 1 & 158\dg 05.1\pr & 13 & 2 & 173\dg 07.5\pr & 13 & 3 & 188\dg 10.0\pr \\
13 & 4 & 203\dg 12.5\pr & 13 & 5 & 218\dg 14.9\pr & 13 & 6 & 233\dg 17.4\pr & 13 & 7 & 248\dg 19.9\pr \\
13 & 8 & 263\dg 22.3\pr & 13 & 9 & 278\dg 24.8\pr & 13 & 10 & 293\dg 27.3\pr & 13 & 11 & 308\dg 29.7\pr \\
13 & 12 & 323\dg 32.2\pr & 13 & 13 & 338\dg 34.7\pr & 13 & 14 & 353\dg 37.1\pr & 13 & 15 & 8\dg 39.6\pr \\
13 & 16 & 23\dg 42.0\pr & 13 & 17 & 38\dg 44.5\pr & 13 & 18 & 53\dg 47.0\pr & 13 & 19 & 68\dg 49.4\pr \\
13 & 20 & 83\dg 51.9\pr & 13 & 21 & 98\dg 54.4\pr & 13 & 22 & 113\dg 56.8\pr & 13 & 23 & 128\dg 59.3\pr \\
14 & 0 & 144\dg 01.8\pr & 14 & 1 & 159\dg 04.2\pr & 14 & 2 & 174\dg 06.7\pr & 14 & 3 & 189\dg 09.2\pr \\
14 & 4 & 204\dg 11.6\pr & 14 & 5 & 219\dg 14.1\pr & 14 & 6 & 234\dg 16.5\pr & 14 & 7 & 249\dg 19.0\pr \\
14 & 8 & 264\dg 21.5\pr & 14 & 9 & 279\dg 23.9\pr & 14 & 10 & 294\dg 26.4\pr & 14 & 11 & 309\dg 28.9\pr \\
14 & 12 & 324\dg 31.3\pr & 14 & 13 & 339\dg 33.8\pr & 14 & 14 & 354\dg 36.3\pr & 14 & 15 & 9\dg 38.7\pr \\
14 & 16 & 24\dg 41.2\pr & 14 & 17 & 39\dg 43.6\pr & 14 & 18 & 54\dg 46.1\pr & 14 & 19 & 69\dg 48.6\pr \\
14 & 20 & 84\dg 51.0\pr & 14 & 21 & 99\dg 53.5\pr & 14 & 22 & 114\dg 56.0\pr & 14 & 23 & 129\dg 58.4\pr \\
15 & 0 & 145\dg 00.9\pr & 15 & 1 & 160\dg 03.4\pr & 15 & 2 & 175\dg 05.8\pr & 15 & 3 & 190\dg 08.3\pr \\
15 & 4 & 205\dg 10.8\pr & 15 & 5 & 220\dg 13.2\pr & 15 & 6 & 235\dg 15.7\pr & 15 & 7 & 250\dg 18.1\pr \\
15 & 8 & 265\dg 20.6\pr & 15 & 9 & 280\dg 23.1\pr & 15 & 10 & 295\dg 25.5\pr & 15 & 11 & 310\dg 28.0\pr \\
15 & 12 & 325\dg 30.5\pr & 15 & 13 & 340\dg 32.9\pr & 15 & 14 & 355\dg 35.4\pr & 15 & 15 & 10\dg 37.9\pr \\
15 & 16 & 25\dg 40.3\pr & 15 & 17 & 40\dg 42.8\pr & 15 & 18 & 55\dg 45.3\pr & 15 & 19 & 70\dg 47.7\pr \\
15 & 20 & 85\dg 50.2\pr & 15 & 21 & 100\dg 52.6\pr & 15 & 22 & 115\dg 55.1\pr & 15 & 23 & 130\dg 57.6\pr \\
16 & 0 & 146\dg 00.0\pr & 16 & 1 & 161\dg 02.5\pr & 16 & 2 & 176\dg 05.0\pr & 16 & 3 & 191\dg 07.4\pr \\
16 & 4 & 206\dg 09.9\pr & 16 & 5 & 221\dg 12.4\pr & 16 & 6 & 236\dg 14.8\pr & 16 & 7 & 251\dg 17.3\pr \\
16 & 8 & 266\dg 19.7\pr & 16 & 9 & 281\dg 22.2\pr & 16 & 10 & 296\dg 24.7\pr & 16 & 11 & 311\dg 27.1\pr \\
16 & 12 & 326\dg 29.6\pr & 16 & 13 & 341\dg 32.1\pr & 16 & 14 & 356\dg 34.5\pr & 16 & 15 & 11\dg 37.0\pr \\
16 & 16 & 26\dg 39.5\pr & 16 & 17 & 41\dg 41.9\pr & 16 & 18 & 56\dg 44.4\pr & 16 & 19 & 71\dg 46.9\pr \\
16 & 20 & 86\dg 49.3\pr & 16 & 21 & 101\dg 51.8\pr & 16 & 22 & 116\dg 54.2\pr & 16 & 23 & 131\dg 56.7\pr \\
17 & 0 & 146\dg 59.2\pr & 17 & 1 & 162\dg 01.6\pr & 17 & 2 & 177\dg 04.1\pr & 17 & 3 & 192\dg 06.6\pr \\
17 & 4 & 207\dg 09.0\pr & 17 & 5 & 222\dg 11.5\pr & 17 & 6 & 237\dg 14.0\pr & 17 & 7 & 252\dg 16.4\pr \\
17 & 8 & 267\dg 18.9\pr & 17 & 9 & 282\dg 21.4\pr & 17 & 10 & 297\dg 23.8\pr & 17 & 11 & 312\dg 26.3\pr \\
17 & 12 & 327\dg 28.7\pr & 17 & 13 & 342\dg 31.2\pr & 17 & 14 & 357\dg 33.7\pr & 17 & 15 & 12\dg 36.1\pr \\
17 & 16 & 27\dg 38.6\pr & 17 & 17 & 42\dg 41.1\pr & 17 & 18 & 57\dg 43.5\pr & 17 & 19 & 72\dg 46.0\pr \\
17 & 20 & 87\dg 48.5\pr & 17 & 21 & 102\dg 50.9\pr & 17 & 22 & 117\dg 53.4\pr & 17 & 23 & 132\dg 55.8\pr \\
18 & 0 & 147\dg 58.3\pr & 18 & 1 & 163\dg 00.8\pr & 18 & 2 & 178\dg 03.2\pr & 18 & 3 & 193\dg 05.7\pr \\
18 & 4 & 208\dg 08.2\pr & 18 & 5 & 223\dg 10.6\pr & 18 & 6 & 238\dg 13.1\pr & 18 & 7 & 253\dg 15.6\pr \\
18 & 8 & 268\dg 18.0\pr & 18 & 9 & 283\dg 20.5\pr & 18 & 10 & 298\dg 23.0\pr & 18 & 11 & 313\dg 25.4\pr \\
18 & 12 & 328\dg 27.9\pr & 18 & 13 & 343\dg 30.3\pr & 18 & 14 & 358\dg 32.8\pr & 18 & 15 & 13\dg 35.3\pr \\
18 & 16 & 28\dg 37.7\pr & 18 & 17 & 43\dg 40.2\pr & 18 & 18 & 58\dg 42.7\pr & 18 & 19 & 73\dg 45.1\pr \\
18 & 20 & 88\dg 47.6\pr & 18 & 21 & 103\dg 50.1\pr & 18 & 22 & 118\dg 52.5\pr & 18 & 23 & 133\dg 55.0\pr \\
19 & 0 & 148\dg 57.5\pr & 19 & 1 & 163\dg 59.9\pr & 19 & 2 & 179\dg 02.4\pr & 19 & 3 & 194\dg 04.8\pr \\
19 & 4 & 209\dg 07.3\pr & 19 & 5 & 224\dg 09.8\pr & 19 & 6 & 239\dg 12.2\pr & 19 & 7 & 254\dg 14.7\pr \\
19 & 8 & 269\dg 17.2\pr & 19 & 9 & 284\dg 19.6\pr & 19 & 10 & 299\dg 22.1\pr & 19 & 11 & 314\dg 24.6\pr \\
19 & 12 & 329\dg 27.0\pr & 19 & 13 & 344\dg 29.5\pr & 19 & 14 & 359\dg 32.0\pr & 19 & 15 & 14\dg 34.4\pr \\
19 & 16 & 29\dg 36.9\pr & 19 & 17 & 44\dg 39.3\pr & 19 & 18 & 59\dg 41.8\pr & 19 & 19 & 74\dg 44.3\pr \\
19 & 20 & 89\dg 46.7\pr & 19 & 21 & 104\dg 49.2\pr & 19 & 22 & 119\dg 51.7\pr & 19 & 23 & 134\dg 54.1\pr \\
20 & 0 & 149\dg 56.6\pr & 20 & 1 & 164\dg 59.1\pr & 20 & 2 & 180\dg 01.5\pr & 20 & 3 & 195\dg 04.0\pr \\
20 & 4 & 210\dg 06.4\pr & 20 & 5 & 225\dg 08.9\pr & 20 & 6 & 240\dg 11.4\pr & 20 & 7 & 255\dg 13.8\pr \\
20 & 8 & 270\dg 16.3\pr & 20 & 9 & 285\dg 18.8\pr & 20 & 10 & 300\dg 21.2\pr & 20 & 11 & 315\dg 23.7\pr \\
20 & 12 & 330\dg 26.2\pr & 20 & 13 & 345\dg 28.6\pr & 20 & 14 & 0\dg 31.1\pr & 20 & 15 & 15\dg 33.6\pr \\
20 & 16 & 30\dg 36.0\pr & 20 & 17 & 45\dg 38.5\pr & 20 & 18 & 60\dg 40.9\pr & 20 & 19 & 75\dg 43.4\pr \\
20 & 20 & 90\dg 45.9\pr & 20 & 21 & 105\dg 48.3\pr & 20 & 22 & 120\dg 50.8\pr & 20 & 23 & 135\dg 53.3\pr \\
21 & 0 & 150\dg 55.7\pr & 21 & 1 & 165\dg 58.2\pr & 21 & 2 & 181\dg 00.7\pr & 21 & 3 & 196\dg 03.1\pr \\
21 & 4 & 211\dg 05.6\pr & 21 & 5 & 226\dg 08.1\pr & 21 & 6 & 241\dg 10.5\pr & 21 & 7 & 256\dg 13.0\pr \\
21 & 8 & 271\dg 15.4\pr & 21 & 9 & 286\dg 17.9\pr & 21 & 10 & 301\dg 20.4\pr & 21 & 11 & 316\dg 22.8\pr \\
21 & 12 & 331\dg 25.3\pr & 21 & 13 & 346\dg 27.8\pr & 21 & 14 & 1\dg 30.2\pr & 21 & 15 & 16\dg 32.7\pr \\
21 & 16 & 31\dg 35.2\pr & 21 & 17 & 46\dg 37.6\pr & 21 & 18 & 61\dg 40.1\pr & 21 & 19 & 76\dg 42.5\pr \\
21 & 20 & 91\dg 45.0\pr & 21 & 21 & 106\dg 47.5\pr & 21 & 22 & 121\dg 49.9\pr & 21 & 23 & 136\dg 52.4\pr \\
22 & 0 & 151\dg 54.9\pr & 22 & 1 & 166\dg 57.3\pr & 22 & 2 & 181\dg 59.8\pr & 22 & 3 & 197\dg 02.3\pr \\
22 & 4 & 212\dg 04.7\pr & 22 & 5 & 227\dg 07.2\pr & 22 & 6 & 242\dg 09.7\pr & 22 & 7 & 257\dg 12.1\pr \\
22 & 8 & 272\dg 14.6\pr & 22 & 9 & 287\dg 17.0\pr & 22 & 10 & 302\dg 19.5\pr & 22 & 11 & 317\dg 22.0\pr \\
22 & 12 & 332\dg 24.4\pr & 22 & 13 & 347\dg 26.9\pr & 22 & 14 & 2\dg 29.4\pr & 22 & 15 & 17\dg 31.8\pr \\
22 & 16 & 32\dg 34.3\pr & 22 & 17 & 47\dg 36.8\pr & 22 & 18 & 62\dg 39.2\pr & 22 & 19 & 77\dg 41.7\pr \\
22 & 20 & 92\dg 44.2\pr & 22 & 21 & 107\dg 46.6\pr & 22 & 22 & 122\dg 49.1\pr & 22 & 23 & 137\dg 51.5\pr \\
23 & 0 & 152\dg 54.0\pr & 23 & 1 & 167\dg 56.5\pr & 23 & 2 & 182\dg 58.9\pr & 23 & 3 & 198\dg 01.4\pr \\
23 & 4 & 213\dg 03.9\pr & 23 & 5 & 228\dg 06.3\pr & 23 & 6 & 243\dg 08.8\pr & 23 & 7 & 258\dg 11.3\pr \\
23 & 8 & 273\dg 13.7\pr & 23 & 9 & 288\dg 16.2\pr & 23 & 10 & 303\dg 18.6\pr & 23 & 11 & 318\dg 21.1\pr \\
23 & 12 & 333\dg 23.6\pr & 23 & 13 & 348\dg 26.0\pr & 23 & 14 & 3\dg 28.5\pr & 23 & 15 & 18\dg 31.0\pr \\
23 & 16 & 33\dg 33.4\pr & 23 & 17 & 48\dg 35.9\pr & 23 & 18 & 63\dg 38.4\pr & 23 & 19 & 78\dg 40.8\pr \\
23 & 20 & 93\dg 43.3\pr & 23 & 21 & 108\dg 45.8\pr & 23 & 22 & 123\dg 48.2\pr & 23 & 23 & 138\dg 50.7\pr \\
24 & 0 & 153\dg 53.1\pr & 24 & 1 & 168\dg 55.6\pr & 24 & 2 & 183\dg 58.1\pr & 24 & 3 & 199\dg 00.5\pr \\
24 & 4 & 214\dg 03.0\pr & 24 & 5 & 229\dg 05.5\pr & 24 & 6 & 244\dg 07.9\pr & 24 & 7 & 259\dg 10.4\pr \\
24 & 8 & 274\dg 12.9\pr & 24 & 9 & 289\dg 15.3\pr & 24 & 10 & 304\dg 17.8\pr & 24 & 11 & 319\dg 20.3\pr \\
24 & 12 & 334\dg 22.7\pr & 24 & 13 & 349\dg 25.2\pr & 24 & 14 & 4\dg 27.6\pr & 24 & 15 & 19\dg 30.1\pr \\
24 & 16 & 34\dg 32.6\pr & 24 & 17 & 49\dg 35.0\pr & 24 & 18 & 64\dg 37.5\pr & 24 & 19 & 79\dg 40.0\pr \\
24 & 20 & 94\dg 42.4\pr & 24 & 21 & 109\dg 44.9\pr & 24 & 22 & 124\dg 47.4\pr & 24 & 23 & 139\dg 49.8\pr \\
25 & 0 & 154\dg 52.3\pr & 25 & 1 & 169\dg 54.7\pr & 25 & 2 & 184\dg 57.2\pr & 25 & 3 & 199\dg 59.7\pr \\
25 & 4 & 215\dg 02.1\pr & 25 & 5 & 230\dg 04.6\pr & 25 & 6 & 245\dg 07.1\pr & 25 & 7 & 260\dg 09.5\pr \\
25 & 8 & 275\dg 12.0\pr & 25 & 9 & 290\dg 14.5\pr & 25 & 10 & 305\dg 16.9\pr & 25 & 11 & 320\dg 19.4\pr \\
25 & 12 & 335\dg 21.9\pr & 25 & 13 & 350\dg 24.3\pr & 25 & 14 & 5\dg 26.8\pr & 25 & 15 & 20\dg 29.2\pr \\
25 & 16 & 35\dg 31.7\pr & 25 & 17 & 50\dg 34.2\pr & 25 & 18 & 65\dg 36.6\pr & 25 & 19 & 80\dg 39.1\pr \\
25 & 20 & 95\dg 41.6\pr & 25 & 21 & 110\dg 44.0\pr & 25 & 22 & 125\dg 46.5\pr & 25 & 23 & 140\dg 49.0\pr \\
26 & 0 & 155\dg 51.4\pr & 26 & 1 & 170\dg 53.9\pr & 26 & 2 & 185\dg 56.4\pr & 26 & 3 & 200\dg 58.8\pr \\
26 & 4 & 216\dg 01.3\pr & 26 & 5 & 231\dg 03.7\pr & 26 & 6 & 246\dg 06.2\pr & 26 & 7 & 261\dg 08.7\pr \\
26 & 8 & 276\dg 11.1\pr & 26 & 9 & 291\dg 13.6\pr & 26 & 10 & 306\dg 16.1\pr & 26 & 11 & 321\dg 18.5\pr \\
26 & 12 & 336\dg 21.0\pr & 26 & 13 & 351\dg 23.5\pr & 26 & 14 & 6\dg 25.9\pr & 26 & 15 & 21\dg 28.4\pr \\
26 & 16 & 36\dg 30.9\pr & 26 & 17 & 51\dg 33.3\pr & 26 & 18 & 66\dg 35.8\pr & 26 & 19 & 81\dg 38.2\pr \\
26 & 20 & 96\dg 40.7\pr & 26 & 21 & 111\dg 43.2\pr & 26 & 22 & 126\dg 45.6\pr & 26 & 23 & 141\dg 48.1\pr \\
27 & 0 & 156\dg 50.6\pr & 27 & 1 & 171\dg 53.0\pr & 27 & 2 & 186\dg 55.5\pr & 27 & 3 & 201\dg 58.0\pr \\
27 & 4 & 217\dg 00.4\pr & 27 & 5 & 232\dg 02.9\pr & 27 & 6 & 247\dg 05.3\pr & 27 & 7 & 262\dg 07.8\pr \\
27 & 8 & 277\dg 10.3\pr & 27 & 9 & 292\dg 12.7\pr & 27 & 10 & 307\dg 15.2\pr & 27 & 11 & 322\dg 17.7\pr \\
27 & 12 & 337\dg 20.1\pr & 27 & 13 & 352\dg 22.6\pr & 27 & 14 & 7\dg 25.1\pr & 27 & 15 & 22\dg 27.5\pr \\
27 & 16 & 37\dg 30.0\pr & 27 & 17 & 52\dg 32.5\pr & 27 & 18 & 67\dg 34.9\pr & 27 & 19 & 82\dg 37.4\pr \\
27 & 20 & 97\dg 39.8\pr & 27 & 21 & 112\dg 42.3\pr & 27 & 22 & 127\dg 44.8\pr & 27 & 23 & 142\dg 47.2\pr \\
28 & 0 & 157\dg 49.7\pr & 28 & 1 & 172\dg 52.2\pr & 28 & 2 & 187\dg 54.6\pr & 28 & 3 & 202\dg 57.1\pr \\
28 & 4 & 217\dg 59.6\pr & 28 & 5 & 233\dg 02.0\pr & 28 & 6 & 248\dg 04.5\pr & 28 & 7 & 263\dg 07.0\pr \\
28 & 8 & 278\dg 09.4\pr & 28 & 9 & 293\dg 11.9\pr & 28 & 10 & 308\dg 14.3\pr & 28 & 11 & 323\dg 16.8\pr \\
28 & 12 & 338\dg 19.3\pr & 28 & 13 & 353\dg 21.7\pr & 28 & 14 & 8\dg 24.2\pr & 28 & 15 & 23\dg 26.7\pr \\
28 & 16 & 38\dg 29.1\pr & 28 & 17 & 53\dg 31.6\pr & 28 & 18 & 68\dg 34.1\pr & 28 & 19 & 83\dg 36.5\pr \\
28 & 20 & 98\dg 39.0\pr & 28 & 21 & 113\dg 41.4\pr & 28 & 22 & 128\dg 43.9\pr & 28 & 23 & 143\dg 46.4\pr \\
\end{longtable}}

{\footnotesize
\begin{longtable}{rrr|rrr|rrr|rrr}
\caption{Aries --- GHA of the First Point of Aries, March 2026 (UT)}\label{aries03}\\
\toprule
\textbf{d} & \textbf{h} & \textbf{GHA} $\gamma$ & \textbf{d} & \textbf{h} & \textbf{GHA} $\gamma$ & \textbf{d} & \textbf{h} & \textbf{GHA} $\gamma$ & \textbf{d} & \textbf{h} & \textbf{GHA} $\gamma$ \\
\midrule
\endfirsthead
\multicolumn{12}{c}{\textit{Aries --- GHA of the First Point of Aries, March 2026 (UT) (continued)}}\\
\toprule
\textbf{d} & \textbf{h} & \textbf{GHA} $\gamma$ & \textbf{d} & \textbf{h} & \textbf{GHA} $\gamma$ & \textbf{d} & \textbf{h} & \textbf{GHA} $\gamma$ & \textbf{d} & \textbf{h} & \textbf{GHA} $\gamma$ \\
\midrule
\endhead
\midrule\multicolumn{12}{r}{\textit{continued on next page}}\\
\endfoot
\bottomrule
\endlastfoot
1 & 0 & 158\dg 48.8\pr & 1 & 1 & 173\dg 51.3\pr & 1 & 2 & 188\dg 53.8\pr & 1 & 3 & 203\dg 56.2\pr \\
1 & 4 & 218\dg 58.7\pr & 1 & 5 & 234\dg 01.2\pr & 1 & 6 & 249\dg 03.6\pr & 1 & 7 & 264\dg 06.1\pr \\
1 & 8 & 279\dg 08.6\pr & 1 & 9 & 294\dg 11.0\pr & 1 & 10 & 309\dg 13.5\pr & 1 & 11 & 324\dg 15.9\pr \\
1 & 12 & 339\dg 18.4\pr & 1 & 13 & 354\dg 20.9\pr & 1 & 14 & 9\dg 23.3\pr & 1 & 15 & 24\dg 25.8\pr \\
1 & 16 & 39\dg 28.3\pr & 1 & 17 & 54\dg 30.7\pr & 1 & 18 & 69\dg 33.2\pr & 1 & 19 & 84\dg 35.7\pr \\
1 & 20 & 99\dg 38.1\pr & 1 & 21 & 114\dg 40.6\pr & 1 & 22 & 129\dg 43.1\pr & 1 & 23 & 144\dg 45.5\pr \\
2 & 0 & 159\dg 48.0\pr & 2 & 1 & 174\dg 50.4\pr & 2 & 2 & 189\dg 52.9\pr & 2 & 3 & 204\dg 55.4\pr \\
2 & 4 & 219\dg 57.8\pr & 2 & 5 & 235\dg 00.3\pr & 2 & 6 & 250\dg 02.8\pr & 2 & 7 & 265\dg 05.2\pr \\
2 & 8 & 280\dg 07.7\pr & 2 & 9 & 295\dg 10.2\pr & 2 & 10 & 310\dg 12.6\pr & 2 & 11 & 325\dg 15.1\pr \\
2 & 12 & 340\dg 17.5\pr & 2 & 13 & 355\dg 20.0\pr & 2 & 14 & 10\dg 22.5\pr & 2 & 15 & 25\dg 24.9\pr \\
2 & 16 & 40\dg 27.4\pr & 2 & 17 & 55\dg 29.9\pr & 2 & 18 & 70\dg 32.3\pr & 2 & 19 & 85\dg 34.8\pr \\
2 & 20 & 100\dg 37.3\pr & 2 & 21 & 115\dg 39.7\pr & 2 & 22 & 130\dg 42.2\pr & 2 & 23 & 145\dg 44.7\pr \\
3 & 0 & 160\dg 47.1\pr & 3 & 1 & 175\dg 49.6\pr & 3 & 2 & 190\dg 52.0\pr & 3 & 3 & 205\dg 54.5\pr \\
3 & 4 & 220\dg 57.0\pr & 3 & 5 & 235\dg 59.4\pr & 3 & 6 & 251\dg 01.9\pr & 3 & 7 & 266\dg 04.4\pr \\
3 & 8 & 281\dg 06.8\pr & 3 & 9 & 296\dg 09.3\pr & 3 & 10 & 311\dg 11.8\pr & 3 & 11 & 326\dg 14.2\pr \\
3 & 12 & 341\dg 16.7\pr & 3 & 13 & 356\dg 19.2\pr & 3 & 14 & 11\dg 21.6\pr & 3 & 15 & 26\dg 24.1\pr \\
3 & 16 & 41\dg 26.5\pr & 3 & 17 & 56\dg 29.0\pr & 3 & 18 & 71\dg 31.5\pr & 3 & 19 & 86\dg 33.9\pr \\
3 & 20 & 101\dg 36.4\pr & 3 & 21 & 116\dg 38.9\pr & 3 & 22 & 131\dg 41.3\pr & 3 & 23 & 146\dg 43.8\pr \\
4 & 0 & 161\dg 46.3\pr & 4 & 1 & 176\dg 48.7\pr & 4 & 2 & 191\dg 51.2\pr & 4 & 3 & 206\dg 53.6\pr \\
4 & 4 & 221\dg 56.1\pr & 4 & 5 & 236\dg 58.6\pr & 4 & 6 & 252\dg 01.0\pr & 4 & 7 & 267\dg 03.5\pr \\
4 & 8 & 282\dg 06.0\pr & 4 & 9 & 297\dg 08.4\pr & 4 & 10 & 312\dg 10.9\pr & 4 & 11 & 327\dg 13.4\pr \\
4 & 12 & 342\dg 15.8\pr & 4 & 13 & 357\dg 18.3\pr & 4 & 14 & 12\dg 20.8\pr & 4 & 15 & 27\dg 23.2\pr \\
4 & 16 & 42\dg 25.7\pr & 4 & 17 & 57\dg 28.1\pr & 4 & 18 & 72\dg 30.6\pr & 4 & 19 & 87\dg 33.1\pr \\
4 & 20 & 102\dg 35.5\pr & 4 & 21 & 117\dg 38.0\pr & 4 & 22 & 132\dg 40.5\pr & 4 & 23 & 147\dg 42.9\pr \\
5 & 0 & 162\dg 45.4\pr & 5 & 1 & 177\dg 47.9\pr & 5 & 2 & 192\dg 50.3\pr & 5 & 3 & 207\dg 52.8\pr \\
5 & 4 & 222\dg 55.3\pr & 5 & 5 & 237\dg 57.7\pr & 5 & 6 & 253\dg 00.2\pr & 5 & 7 & 268\dg 02.6\pr \\
5 & 8 & 283\dg 05.1\pr & 5 & 9 & 298\dg 07.6\pr & 5 & 10 & 313\dg 10.0\pr & 5 & 11 & 328\dg 12.5\pr \\
5 & 12 & 343\dg 15.0\pr & 5 & 13 & 358\dg 17.4\pr & 5 & 14 & 13\dg 19.9\pr & 5 & 15 & 28\dg 22.4\pr \\
5 & 16 & 43\dg 24.8\pr & 5 & 17 & 58\dg 27.3\pr & 5 & 18 & 73\dg 29.8\pr & 5 & 19 & 88\dg 32.2\pr \\
5 & 20 & 103\dg 34.7\pr & 5 & 21 & 118\dg 37.1\pr & 5 & 22 & 133\dg 39.6\pr & 5 & 23 & 148\dg 42.1\pr \\
6 & 0 & 163\dg 44.5\pr & 6 & 1 & 178\dg 47.0\pr & 6 & 2 & 193\dg 49.5\pr & 6 & 3 & 208\dg 51.9\pr \\
6 & 4 & 223\dg 54.4\pr & 6 & 5 & 238\dg 56.9\pr & 6 & 6 & 253\dg 59.3\pr & 6 & 7 & 269\dg 01.8\pr \\
6 & 8 & 284\dg 04.2\pr & 6 & 9 & 299\dg 06.7\pr & 6 & 10 & 314\dg 09.2\pr & 6 & 11 & 329\dg 11.6\pr \\
6 & 12 & 344\dg 14.1\pr & 6 & 13 & 359\dg 16.6\pr & 6 & 14 & 14\dg 19.0\pr & 6 & 15 & 29\dg 21.5\pr \\
6 & 16 & 44\dg 24.0\pr & 6 & 17 & 59\dg 26.4\pr & 6 & 18 & 74\dg 28.9\pr & 6 & 19 & 89\dg 31.4\pr \\
6 & 20 & 104\dg 33.8\pr & 6 & 21 & 119\dg 36.3\pr & 6 & 22 & 134\dg 38.7\pr & 6 & 23 & 149\dg 41.2\pr \\
7 & 0 & 164\dg 43.7\pr & 7 & 1 & 179\dg 46.1\pr & 7 & 2 & 194\dg 48.6\pr & 7 & 3 & 209\dg 51.1\pr \\
7 & 4 & 224\dg 53.5\pr & 7 & 5 & 239\dg 56.0\pr & 7 & 6 & 254\dg 58.5\pr & 7 & 7 & 270\dg 00.9\pr \\
7 & 8 & 285\dg 03.4\pr & 7 & 9 & 300\dg 05.9\pr & 7 & 10 & 315\dg 08.3\pr & 7 & 11 & 330\dg 10.8\pr \\
7 & 12 & 345\dg 13.2\pr & 7 & 13 & 0\dg 15.7\pr & 7 & 14 & 15\dg 18.2\pr & 7 & 15 & 30\dg 20.6\pr \\
7 & 16 & 45\dg 23.1\pr & 7 & 17 & 60\dg 25.6\pr & 7 & 18 & 75\dg 28.0\pr & 7 & 19 & 90\dg 30.5\pr \\
7 & 20 & 105\dg 33.0\pr & 7 & 21 & 120\dg 35.4\pr & 7 & 22 & 135\dg 37.9\pr & 7 & 23 & 150\dg 40.3\pr \\
8 & 0 & 165\dg 42.8\pr & 8 & 1 & 180\dg 45.3\pr & 8 & 2 & 195\dg 47.7\pr & 8 & 3 & 210\dg 50.2\pr \\
8 & 4 & 225\dg 52.7\pr & 8 & 5 & 240\dg 55.1\pr & 8 & 6 & 255\dg 57.6\pr & 8 & 7 & 271\dg 00.1\pr \\
8 & 8 & 286\dg 02.5\pr & 8 & 9 & 301\dg 05.0\pr & 8 & 10 & 316\dg 07.5\pr & 8 & 11 & 331\dg 09.9\pr \\
8 & 12 & 346\dg 12.4\pr & 8 & 13 & 1\dg 14.8\pr & 8 & 14 & 16\dg 17.3\pr & 8 & 15 & 31\dg 19.8\pr \\
8 & 16 & 46\dg 22.2\pr & 8 & 17 & 61\dg 24.7\pr & 8 & 18 & 76\dg 27.2\pr & 8 & 19 & 91\dg 29.6\pr \\
8 & 20 & 106\dg 32.1\pr & 8 & 21 & 121\dg 34.6\pr & 8 & 22 & 136\dg 37.0\pr & 8 & 23 & 151\dg 39.5\pr \\
9 & 0 & 166\dg 42.0\pr & 9 & 1 & 181\dg 44.4\pr & 9 & 2 & 196\dg 46.9\pr & 9 & 3 & 211\dg 49.3\pr \\
9 & 4 & 226\dg 51.8\pr & 9 & 5 & 241\dg 54.3\pr & 9 & 6 & 256\dg 56.7\pr & 9 & 7 & 271\dg 59.2\pr \\
9 & 8 & 287\dg 01.7\pr & 9 & 9 & 302\dg 04.1\pr & 9 & 10 & 317\dg 06.6\pr & 9 & 11 & 332\dg 09.1\pr \\
9 & 12 & 347\dg 11.5\pr & 9 & 13 & 2\dg 14.0\pr & 9 & 14 & 17\dg 16.4\pr & 9 & 15 & 32\dg 18.9\pr \\
9 & 16 & 47\dg 21.4\pr & 9 & 17 & 62\dg 23.8\pr & 9 & 18 & 77\dg 26.3\pr & 9 & 19 & 92\dg 28.8\pr \\
9 & 20 & 107\dg 31.2\pr & 9 & 21 & 122\dg 33.7\pr & 9 & 22 & 137\dg 36.2\pr & 9 & 23 & 152\dg 38.6\pr \\
10 & 0 & 167\dg 41.1\pr & 10 & 1 & 182\dg 43.6\pr & 10 & 2 & 197\dg 46.0\pr & 10 & 3 & 212\dg 48.5\pr \\
10 & 4 & 227\dg 50.9\pr & 10 & 5 & 242\dg 53.4\pr & 10 & 6 & 257\dg 55.9\pr & 10 & 7 & 272\dg 58.3\pr \\
10 & 8 & 288\dg 00.8\pr & 10 & 9 & 303\dg 03.3\pr & 10 & 10 & 318\dg 05.7\pr & 10 & 11 & 333\dg 08.2\pr \\
10 & 12 & 348\dg 10.7\pr & 10 & 13 & 3\dg 13.1\pr & 10 & 14 & 18\dg 15.6\pr & 10 & 15 & 33\dg 18.1\pr \\
10 & 16 & 48\dg 20.5\pr & 10 & 17 & 63\dg 23.0\pr & 10 & 18 & 78\dg 25.4\pr & 10 & 19 & 93\dg 27.9\pr \\
10 & 20 & 108\dg 30.4\pr & 10 & 21 & 123\dg 32.8\pr & 10 & 22 & 138\dg 35.3\pr & 10 & 23 & 153\dg 37.8\pr \\
11 & 0 & 168\dg 40.2\pr & 11 & 1 & 183\dg 42.7\pr & 11 & 2 & 198\dg 45.2\pr & 11 & 3 & 213\dg 47.6\pr \\
11 & 4 & 228\dg 50.1\pr & 11 & 5 & 243\dg 52.5\pr & 11 & 6 & 258\dg 55.0\pr & 11 & 7 & 273\dg 57.5\pr \\
11 & 8 & 288\dg 59.9\pr & 11 & 9 & 304\dg 02.4\pr & 11 & 10 & 319\dg 04.9\pr & 11 & 11 & 334\dg 07.3\pr \\
11 & 12 & 349\dg 09.8\pr & 11 & 13 & 4\dg 12.3\pr & 11 & 14 & 19\dg 14.7\pr & 11 & 15 & 34\dg 17.2\pr \\
11 & 16 & 49\dg 19.7\pr & 11 & 17 & 64\dg 22.1\pr & 11 & 18 & 79\dg 24.6\pr & 11 & 19 & 94\dg 27.0\pr \\
11 & 20 & 109\dg 29.5\pr & 11 & 21 & 124\dg 32.0\pr & 11 & 22 & 139\dg 34.4\pr & 11 & 23 & 154\dg 36.9\pr \\
12 & 0 & 169\dg 39.4\pr & 12 & 1 & 184\dg 41.8\pr & 12 & 2 & 199\dg 44.3\pr & 12 & 3 & 214\dg 46.8\pr \\
12 & 4 & 229\dg 49.2\pr & 12 & 5 & 244\dg 51.7\pr & 12 & 6 & 259\dg 54.2\pr & 12 & 7 & 274\dg 56.6\pr \\
12 & 8 & 289\dg 59.1\pr & 12 & 9 & 305\dg 01.5\pr & 12 & 10 & 320\dg 04.0\pr & 12 & 11 & 335\dg 06.5\pr \\
12 & 12 & 350\dg 08.9\pr & 12 & 13 & 5\dg 11.4\pr & 12 & 14 & 20\dg 13.9\pr & 12 & 15 & 35\dg 16.3\pr \\
12 & 16 & 50\dg 18.8\pr & 12 & 17 & 65\dg 21.3\pr & 12 & 18 & 80\dg 23.7\pr & 12 & 19 & 95\dg 26.2\pr \\
12 & 20 & 110\dg 28.7\pr & 12 & 21 & 125\dg 31.1\pr & 12 & 22 & 140\dg 33.6\pr & 12 & 23 & 155\dg 36.0\pr \\
13 & 0 & 170\dg 38.5\pr & 13 & 1 & 185\dg 41.0\pr & 13 & 2 & 200\dg 43.4\pr & 13 & 3 & 215\dg 45.9\pr \\
13 & 4 & 230\dg 48.4\pr & 13 & 5 & 245\dg 50.8\pr & 13 & 6 & 260\dg 53.3\pr & 13 & 7 & 275\dg 55.8\pr \\
13 & 8 & 290\dg 58.2\pr & 13 & 9 & 306\dg 00.7\pr & 13 & 10 & 321\dg 03.1\pr & 13 & 11 & 336\dg 05.6\pr \\
13 & 12 & 351\dg 08.1\pr & 13 & 13 & 6\dg 10.5\pr & 13 & 14 & 21\dg 13.0\pr & 13 & 15 & 36\dg 15.5\pr \\
13 & 16 & 51\dg 17.9\pr & 13 & 17 & 66\dg 20.4\pr & 13 & 18 & 81\dg 22.9\pr & 13 & 19 & 96\dg 25.3\pr \\
13 & 20 & 111\dg 27.8\pr & 13 & 21 & 126\dg 30.3\pr & 13 & 22 & 141\dg 32.7\pr & 13 & 23 & 156\dg 35.2\pr \\
14 & 0 & 171\dg 37.6\pr & 14 & 1 & 186\dg 40.1\pr & 14 & 2 & 201\dg 42.6\pr & 14 & 3 & 216\dg 45.0\pr \\
14 & 4 & 231\dg 47.5\pr & 14 & 5 & 246\dg 50.0\pr & 14 & 6 & 261\dg 52.4\pr & 14 & 7 & 276\dg 54.9\pr \\
14 & 8 & 291\dg 57.4\pr & 14 & 9 & 306\dg 59.8\pr & 14 & 10 & 322\dg 02.3\pr & 14 & 11 & 337\dg 04.8\pr \\
14 & 12 & 352\dg 07.2\pr & 14 & 13 & 7\dg 09.7\pr & 14 & 14 & 22\dg 12.1\pr & 14 & 15 & 37\dg 14.6\pr \\
14 & 16 & 52\dg 17.1\pr & 14 & 17 & 67\dg 19.5\pr & 14 & 18 & 82\dg 22.0\pr & 14 & 19 & 97\dg 24.5\pr \\
14 & 20 & 112\dg 26.9\pr & 14 & 21 & 127\dg 29.4\pr & 14 & 22 & 142\dg 31.9\pr & 14 & 23 & 157\dg 34.3\pr \\
15 & 0 & 172\dg 36.8\pr & 15 & 1 & 187\dg 39.2\pr & 15 & 2 & 202\dg 41.7\pr & 15 & 3 & 217\dg 44.2\pr \\
15 & 4 & 232\dg 46.6\pr & 15 & 5 & 247\dg 49.1\pr & 15 & 6 & 262\dg 51.6\pr & 15 & 7 & 277\dg 54.0\pr \\
15 & 8 & 292\dg 56.5\pr & 15 & 9 & 307\dg 59.0\pr & 15 & 10 & 323\dg 01.4\pr & 15 & 11 & 338\dg 03.9\pr \\
15 & 12 & 353\dg 06.4\pr & 15 & 13 & 8\dg 08.8\pr & 15 & 14 & 23\dg 11.3\pr & 15 & 15 & 38\dg 13.7\pr \\
15 & 16 & 53\dg 16.2\pr & 15 & 17 & 68\dg 18.7\pr & 15 & 18 & 83\dg 21.1\pr & 15 & 19 & 98\dg 23.6\pr \\
15 & 20 & 113\dg 26.1\pr & 15 & 21 & 128\dg 28.5\pr & 15 & 22 & 143\dg 31.0\pr & 15 & 23 & 158\dg 33.5\pr \\
16 & 0 & 173\dg 35.9\pr & 16 & 1 & 188\dg 38.4\pr & 16 & 2 & 203\dg 40.9\pr & 16 & 3 & 218\dg 43.3\pr \\
16 & 4 & 233\dg 45.8\pr & 16 & 5 & 248\dg 48.2\pr & 16 & 6 & 263\dg 50.7\pr & 16 & 7 & 278\dg 53.2\pr \\
16 & 8 & 293\dg 55.6\pr & 16 & 9 & 308\dg 58.1\pr & 16 & 10 & 324\dg 00.6\pr & 16 & 11 & 339\dg 03.0\pr \\
16 & 12 & 354\dg 05.5\pr & 16 & 13 & 9\dg 08.0\pr & 16 & 14 & 24\dg 10.4\pr & 16 & 15 & 39\dg 12.9\pr \\
16 & 16 & 54\dg 15.3\pr & 16 & 17 & 69\dg 17.8\pr & 16 & 18 & 84\dg 20.3\pr & 16 & 19 & 99\dg 22.7\pr \\
16 & 20 & 114\dg 25.2\pr & 16 & 21 & 129\dg 27.7\pr & 16 & 22 & 144\dg 30.1\pr & 16 & 23 & 159\dg 32.6\pr \\
17 & 0 & 174\dg 35.1\pr & 17 & 1 & 189\dg 37.5\pr & 17 & 2 & 204\dg 40.0\pr & 17 & 3 & 219\dg 42.5\pr \\
17 & 4 & 234\dg 44.9\pr & 17 & 5 & 249\dg 47.4\pr & 17 & 6 & 264\dg 49.8\pr & 17 & 7 & 279\dg 52.3\pr \\
17 & 8 & 294\dg 54.8\pr & 17 & 9 & 309\dg 57.2\pr & 17 & 10 & 324\dg 59.7\pr & 17 & 11 & 340\dg 02.2\pr \\
17 & 12 & 355\dg 04.6\pr & 17 & 13 & 10\dg 07.1\pr & 17 & 14 & 25\dg 09.6\pr & 17 & 15 & 40\dg 12.0\pr \\
17 & 16 & 55\dg 14.5\pr & 17 & 17 & 70\dg 17.0\pr & 17 & 18 & 85\dg 19.4\pr & 17 & 19 & 100\dg 21.9\pr \\
17 & 20 & 115\dg 24.3\pr & 17 & 21 & 130\dg 26.8\pr & 17 & 22 & 145\dg 29.3\pr & 17 & 23 & 160\dg 31.7\pr \\
18 & 0 & 175\dg 34.2\pr & 18 & 1 & 190\dg 36.7\pr & 18 & 2 & 205\dg 39.1\pr & 18 & 3 & 220\dg 41.6\pr \\
18 & 4 & 235\dg 44.1\pr & 18 & 5 & 250\dg 46.5\pr & 18 & 6 & 265\dg 49.0\pr & 18 & 7 & 280\dg 51.4\pr \\
18 & 8 & 295\dg 53.9\pr & 18 & 9 & 310\dg 56.4\pr & 18 & 10 & 325\dg 58.8\pr & 18 & 11 & 341\dg 01.3\pr \\
18 & 12 & 356\dg 03.8\pr & 18 & 13 & 11\dg 06.2\pr & 18 & 14 & 26\dg 08.7\pr & 18 & 15 & 41\dg 11.2\pr \\
18 & 16 & 56\dg 13.6\pr & 18 & 17 & 71\dg 16.1\pr & 18 & 18 & 86\dg 18.6\pr & 18 & 19 & 101\dg 21.0\pr \\
18 & 20 & 116\dg 23.5\pr & 18 & 21 & 131\dg 25.9\pr & 18 & 22 & 146\dg 28.4\pr & 18 & 23 & 161\dg 30.9\pr \\
19 & 0 & 176\dg 33.3\pr & 19 & 1 & 191\dg 35.8\pr & 19 & 2 & 206\dg 38.3\pr & 19 & 3 & 221\dg 40.7\pr \\
19 & 4 & 236\dg 43.2\pr & 19 & 5 & 251\dg 45.7\pr & 19 & 6 & 266\dg 48.1\pr & 19 & 7 & 281\dg 50.6\pr \\
19 & 8 & 296\dg 53.1\pr & 19 & 9 & 311\dg 55.5\pr & 19 & 10 & 326\dg 58.0\pr & 19 & 11 & 342\dg 00.4\pr \\
19 & 12 & 357\dg 02.9\pr & 19 & 13 & 12\dg 05.4\pr & 19 & 14 & 27\dg 07.8\pr & 19 & 15 & 42\dg 10.3\pr \\
19 & 16 & 57\dg 12.8\pr & 19 & 17 & 72\dg 15.2\pr & 19 & 18 & 87\dg 17.7\pr & 19 & 19 & 102\dg 20.2\pr \\
19 & 20 & 117\dg 22.6\pr & 19 & 21 & 132\dg 25.1\pr & 19 & 22 & 147\dg 27.6\pr & 19 & 23 & 162\dg 30.0\pr \\
20 & 0 & 177\dg 32.5\pr & 20 & 1 & 192\dg 34.9\pr & 20 & 2 & 207\dg 37.4\pr & 20 & 3 & 222\dg 39.9\pr \\
20 & 4 & 237\dg 42.3\pr & 20 & 5 & 252\dg 44.8\pr & 20 & 6 & 267\dg 47.3\pr & 20 & 7 & 282\dg 49.7\pr \\
20 & 8 & 297\dg 52.2\pr & 20 & 9 & 312\dg 54.7\pr & 20 & 10 & 327\dg 57.1\pr & 20 & 11 & 342\dg 59.6\pr \\
20 & 12 & 358\dg 02.0\pr & 20 & 13 & 13\dg 04.5\pr & 20 & 14 & 28\dg 07.0\pr & 20 & 15 & 43\dg 09.4\pr \\
20 & 16 & 58\dg 11.9\pr & 20 & 17 & 73\dg 14.4\pr & 20 & 18 & 88\dg 16.8\pr & 20 & 19 & 103\dg 19.3\pr \\
20 & 20 & 118\dg 21.8\pr & 20 & 21 & 133\dg 24.2\pr & 20 & 22 & 148\dg 26.7\pr & 20 & 23 & 163\dg 29.2\pr \\
21 & 0 & 178\dg 31.6\pr & 21 & 1 & 193\dg 34.1\pr & 21 & 2 & 208\dg 36.5\pr & 21 & 3 & 223\dg 39.0\pr \\
21 & 4 & 238\dg 41.5\pr & 21 & 5 & 253\dg 43.9\pr & 21 & 6 & 268\dg 46.4\pr & 21 & 7 & 283\dg 48.9\pr \\
21 & 8 & 298\dg 51.3\pr & 21 & 9 & 313\dg 53.8\pr & 21 & 10 & 328\dg 56.3\pr & 21 & 11 & 343\dg 58.7\pr \\
21 & 12 & 359\dg 01.2\pr & 21 & 13 & 14\dg 03.7\pr & 21 & 14 & 29\dg 06.1\pr & 21 & 15 & 44\dg 08.6\pr \\
21 & 16 & 59\dg 11.0\pr & 21 & 17 & 74\dg 13.5\pr & 21 & 18 & 89\dg 16.0\pr & 21 & 19 & 104\dg 18.4\pr \\
21 & 20 & 119\dg 20.9\pr & 21 & 21 & 134\dg 23.4\pr & 21 & 22 & 149\dg 25.8\pr & 21 & 23 & 164\dg 28.3\pr \\
22 & 0 & 179\dg 30.8\pr & 22 & 1 & 194\dg 33.2\pr & 22 & 2 & 209\dg 35.7\pr & 22 & 3 & 224\dg 38.1\pr \\
22 & 4 & 239\dg 40.6\pr & 22 & 5 & 254\dg 43.1\pr & 22 & 6 & 269\dg 45.5\pr & 22 & 7 & 284\dg 48.0\pr \\
22 & 8 & 299\dg 50.5\pr & 22 & 9 & 314\dg 52.9\pr & 22 & 10 & 329\dg 55.4\pr & 22 & 11 & 344\dg 57.9\pr \\
22 & 12 & 0\dg 00.3\pr & 22 & 13 & 15\dg 02.8\pr & 22 & 14 & 30\dg 05.3\pr & 22 & 15 & 45\dg 07.7\pr \\
22 & 16 & 60\dg 10.2\pr & 22 & 17 & 75\dg 12.6\pr & 22 & 18 & 90\dg 15.1\pr & 22 & 19 & 105\dg 17.6\pr \\
22 & 20 & 120\dg 20.0\pr & 22 & 21 & 135\dg 22.5\pr & 22 & 22 & 150\dg 25.0\pr & 22 & 23 & 165\dg 27.4\pr \\
23 & 0 & 180\dg 29.9\pr & 23 & 1 & 195\dg 32.4\pr & 23 & 2 & 210\dg 34.8\pr & 23 & 3 & 225\dg 37.3\pr \\
23 & 4 & 240\dg 39.8\pr & 23 & 5 & 255\dg 42.2\pr & 23 & 6 & 270\dg 44.7\pr & 23 & 7 & 285\dg 47.1\pr \\
23 & 8 & 300\dg 49.6\pr & 23 & 9 & 315\dg 52.1\pr & 23 & 10 & 330\dg 54.5\pr & 23 & 11 & 345\dg 57.0\pr \\
23 & 12 & 0\dg 59.5\pr & 23 & 13 & 16\dg 01.9\pr & 23 & 14 & 31\dg 04.4\pr & 23 & 15 & 46\dg 06.9\pr \\
23 & 16 & 61\dg 09.3\pr & 23 & 17 & 76\dg 11.8\pr & 23 & 18 & 91\dg 14.2\pr & 23 & 19 & 106\dg 16.7\pr \\
23 & 20 & 121\dg 19.2\pr & 23 & 21 & 136\dg 21.6\pr & 23 & 22 & 151\dg 24.1\pr & 23 & 23 & 166\dg 26.6\pr \\
24 & 0 & 181\dg 29.0\pr & 24 & 1 & 196\dg 31.5\pr & 24 & 2 & 211\dg 34.0\pr & 24 & 3 & 226\dg 36.4\pr \\
24 & 4 & 241\dg 38.9\pr & 24 & 5 & 256\dg 41.4\pr & 24 & 6 & 271\dg 43.8\pr & 24 & 7 & 286\dg 46.3\pr \\
24 & 8 & 301\dg 48.7\pr & 24 & 9 & 316\dg 51.2\pr & 24 & 10 & 331\dg 53.7\pr & 24 & 11 & 346\dg 56.1\pr \\
24 & 12 & 1\dg 58.6\pr & 24 & 13 & 17\dg 01.1\pr & 24 & 14 & 32\dg 03.5\pr & 24 & 15 & 47\dg 06.0\pr \\
24 & 16 & 62\dg 08.5\pr & 24 & 17 & 77\dg 10.9\pr & 24 & 18 & 92\dg 13.4\pr & 24 & 19 & 107\dg 15.9\pr \\
24 & 20 & 122\dg 18.3\pr & 24 & 21 & 137\dg 20.8\pr & 24 & 22 & 152\dg 23.2\pr & 24 & 23 & 167\dg 25.7\pr \\
25 & 0 & 182\dg 28.2\pr & 25 & 1 & 197\dg 30.6\pr & 25 & 2 & 212\dg 33.1\pr & 25 & 3 & 227\dg 35.6\pr \\
25 & 4 & 242\dg 38.0\pr & 25 & 5 & 257\dg 40.5\pr & 25 & 6 & 272\dg 43.0\pr & 25 & 7 & 287\dg 45.4\pr \\
25 & 8 & 302\dg 47.9\pr & 25 & 9 & 317\dg 50.4\pr & 25 & 10 & 332\dg 52.8\pr & 25 & 11 & 347\dg 55.3\pr \\
25 & 12 & 2\dg 57.7\pr & 25 & 13 & 18\dg 00.2\pr & 25 & 14 & 33\dg 02.7\pr & 25 & 15 & 48\dg 05.1\pr \\
25 & 16 & 63\dg 07.6\pr & 25 & 17 & 78\dg 10.1\pr & 25 & 18 & 93\dg 12.5\pr & 25 & 19 & 108\dg 15.0\pr \\
25 & 20 & 123\dg 17.5\pr & 25 & 21 & 138\dg 19.9\pr & 25 & 22 & 153\dg 22.4\pr & 25 & 23 & 168\dg 24.8\pr \\
26 & 0 & 183\dg 27.3\pr & 26 & 1 & 198\dg 29.8\pr & 26 & 2 & 213\dg 32.2\pr & 26 & 3 & 228\dg 34.7\pr \\
26 & 4 & 243\dg 37.2\pr & 26 & 5 & 258\dg 39.6\pr & 26 & 6 & 273\dg 42.1\pr & 26 & 7 & 288\dg 44.6\pr \\
26 & 8 & 303\dg 47.0\pr & 26 & 9 & 318\dg 49.5\pr & 26 & 10 & 333\dg 52.0\pr & 26 & 11 & 348\dg 54.4\pr \\
26 & 12 & 3\dg 56.9\pr & 26 & 13 & 18\dg 59.3\pr & 26 & 14 & 34\dg 01.8\pr & 26 & 15 & 49\dg 04.3\pr \\
26 & 16 & 64\dg 06.7\pr & 26 & 17 & 79\dg 09.2\pr & 26 & 18 & 94\dg 11.7\pr & 26 & 19 & 109\dg 14.1\pr \\
26 & 20 & 124\dg 16.6\pr & 26 & 21 & 139\dg 19.1\pr & 26 & 22 & 154\dg 21.5\pr & 26 & 23 & 169\dg 24.0\pr \\
27 & 0 & 184\dg 26.5\pr & 27 & 1 & 199\dg 28.9\pr & 27 & 2 & 214\dg 31.4\pr & 27 & 3 & 229\dg 33.8\pr \\
27 & 4 & 244\dg 36.3\pr & 27 & 5 & 259\dg 38.8\pr & 27 & 6 & 274\dg 41.2\pr & 27 & 7 & 289\dg 43.7\pr \\
27 & 8 & 304\dg 46.2\pr & 27 & 9 & 319\dg 48.6\pr & 27 & 10 & 334\dg 51.1\pr & 27 & 11 & 349\dg 53.6\pr \\
27 & 12 & 4\dg 56.0\pr & 27 & 13 & 19\dg 58.5\pr & 27 & 14 & 35\dg 00.9\pr & 27 & 15 & 50\dg 03.4\pr \\
27 & 16 & 65\dg 05.9\pr & 27 & 17 & 80\dg 08.3\pr & 27 & 18 & 95\dg 10.8\pr & 27 & 19 & 110\dg 13.3\pr \\
27 & 20 & 125\dg 15.7\pr & 27 & 21 & 140\dg 18.2\pr & 27 & 22 & 155\dg 20.7\pr & 27 & 23 & 170\dg 23.1\pr \\
28 & 0 & 185\dg 25.6\pr & 28 & 1 & 200\dg 28.1\pr & 28 & 2 & 215\dg 30.5\pr & 28 & 3 & 230\dg 33.0\pr \\
28 & 4 & 245\dg 35.4\pr & 28 & 5 & 260\dg 37.9\pr & 28 & 6 & 275\dg 40.4\pr & 28 & 7 & 290\dg 42.8\pr \\
28 & 8 & 305\dg 45.3\pr & 28 & 9 & 320\dg 47.8\pr & 28 & 10 & 335\dg 50.2\pr & 28 & 11 & 350\dg 52.7\pr \\
28 & 12 & 5\dg 55.2\pr & 28 & 13 & 20\dg 57.6\pr & 28 & 14 & 36\dg 00.1\pr & 28 & 15 & 51\dg 02.6\pr \\
28 & 16 & 66\dg 05.0\pr & 28 & 17 & 81\dg 07.5\pr & 28 & 18 & 96\dg 09.9\pr & 28 & 19 & 111\dg 12.4\pr \\
28 & 20 & 126\dg 14.9\pr & 28 & 21 & 141\dg 17.3\pr & 28 & 22 & 156\dg 19.8\pr & 28 & 23 & 171\dg 22.3\pr \\
29 & 0 & 186\dg 24.7\pr & 29 & 1 & 201\dg 27.2\pr & 29 & 2 & 216\dg 29.7\pr & 29 & 3 & 231\dg 32.1\pr \\
29 & 4 & 246\dg 34.6\pr & 29 & 5 & 261\dg 37.0\pr & 29 & 6 & 276\dg 39.5\pr & 29 & 7 & 291\dg 42.0\pr \\
29 & 8 & 306\dg 44.4\pr & 29 & 9 & 321\dg 46.9\pr & 29 & 10 & 336\dg 49.4\pr & 29 & 11 & 351\dg 51.8\pr \\
29 & 12 & 6\dg 54.3\pr & 29 & 13 & 21\dg 56.8\pr & 29 & 14 & 36\dg 59.2\pr & 29 & 15 & 52\dg 01.7\pr \\
29 & 16 & 67\dg 04.2\pr & 29 & 17 & 82\dg 06.6\pr & 29 & 18 & 97\dg 09.1\pr & 29 & 19 & 112\dg 11.5\pr \\
29 & 20 & 127\dg 14.0\pr & 29 & 21 & 142\dg 16.5\pr & 29 & 22 & 157\dg 18.9\pr & 29 & 23 & 172\dg 21.4\pr \\
30 & 0 & 187\dg 23.9\pr & 30 & 1 & 202\dg 26.3\pr & 30 & 2 & 217\dg 28.8\pr & 30 & 3 & 232\dg 31.3\pr \\
30 & 4 & 247\dg 33.7\pr & 30 & 5 & 262\dg 36.2\pr & 30 & 6 & 277\dg 38.7\pr & 30 & 7 & 292\dg 41.1\pr \\
30 & 8 & 307\dg 43.6\pr & 30 & 9 & 322\dg 46.0\pr & 30 & 10 & 337\dg 48.5\pr & 30 & 11 & 352\dg 51.0\pr \\
30 & 12 & 7\dg 53.4\pr & 30 & 13 & 22\dg 55.9\pr & 30 & 14 & 37\dg 58.4\pr & 30 & 15 & 53\dg 00.8\pr \\
30 & 16 & 68\dg 03.3\pr & 30 & 17 & 83\dg 05.8\pr & 30 & 18 & 98\dg 08.2\pr & 30 & 19 & 113\dg 10.7\pr \\
30 & 20 & 128\dg 13.1\pr & 30 & 21 & 143\dg 15.6\pr & 30 & 22 & 158\dg 18.1\pr & 30 & 23 & 173\dg 20.5\pr \\
31 & 0 & 188\dg 23.0\pr & 31 & 1 & 203\dg 25.5\pr & 31 & 2 & 218\dg 27.9\pr & 31 & 3 & 233\dg 30.4\pr \\
31 & 4 & 248\dg 32.9\pr & 31 & 5 & 263\dg 35.3\pr & 31 & 6 & 278\dg 37.8\pr & 31 & 7 & 293\dg 40.3\pr \\
31 & 8 & 308\dg 42.7\pr & 31 & 9 & 323\dg 45.2\pr & 31 & 10 & 338\dg 47.6\pr & 31 & 11 & 353\dg 50.1\pr \\
31 & 12 & 8\dg 52.6\pr & 31 & 13 & 23\dg 55.0\pr & 31 & 14 & 38\dg 57.5\pr & 31 & 15 & 54\dg 00.0\pr \\
31 & 16 & 69\dg 02.4\pr & 31 & 17 & 84\dg 04.9\pr & 31 & 18 & 99\dg 07.4\pr & 31 & 19 & 114\dg 09.8\pr \\
31 & 20 & 129\dg 12.3\pr & 31 & 21 & 144\dg 14.8\pr & 31 & 22 & 159\dg 17.2\pr & 31 & 23 & 174\dg 19.7\pr \\
\end{longtable}}

{\footnotesize
\begin{longtable}{rrr|rrr|rrr|rrr}
\caption{Aries --- GHA of the First Point of Aries, April 2026 (UT)}\label{aries04}\\
\toprule
\textbf{d} & \textbf{h} & \textbf{GHA} $\gamma$ & \textbf{d} & \textbf{h} & \textbf{GHA} $\gamma$ & \textbf{d} & \textbf{h} & \textbf{GHA} $\gamma$ & \textbf{d} & \textbf{h} & \textbf{GHA} $\gamma$ \\
\midrule
\endfirsthead
\multicolumn{12}{c}{\textit{Aries --- GHA of the First Point of Aries, April 2026 (UT) (continued)}}\\
\toprule
\textbf{d} & \textbf{h} & \textbf{GHA} $\gamma$ & \textbf{d} & \textbf{h} & \textbf{GHA} $\gamma$ & \textbf{d} & \textbf{h} & \textbf{GHA} $\gamma$ & \textbf{d} & \textbf{h} & \textbf{GHA} $\gamma$ \\
\midrule
\endhead
\midrule\multicolumn{12}{r}{\textit{continued on next page}}\\
\endfoot
\bottomrule
\endlastfoot
1 & 0 & 189\dg 22.1\pr & 1 & 1 & 204\dg 24.6\pr & 1 & 2 & 219\dg 27.1\pr & 1 & 3 & 234\dg 29.5\pr \\
1 & 4 & 249\dg 32.0\pr & 1 & 5 & 264\dg 34.5\pr & 1 & 6 & 279\dg 36.9\pr & 1 & 7 & 294\dg 39.4\pr \\
1 & 8 & 309\dg 41.9\pr & 1 & 9 & 324\dg 44.3\pr & 1 & 10 & 339\dg 46.8\pr & 1 & 11 & 354\dg 49.3\pr \\
1 & 12 & 9\dg 51.7\pr & 1 & 13 & 24\dg 54.2\pr & 1 & 14 & 39\dg 56.6\pr & 1 & 15 & 54\dg 59.1\pr \\
1 & 16 & 70\dg 01.6\pr & 1 & 17 & 85\dg 04.0\pr & 1 & 18 & 100\dg 06.5\pr & 1 & 19 & 115\dg 09.0\pr \\
1 & 20 & 130\dg 11.4\pr & 1 & 21 & 145\dg 13.9\pr & 1 & 22 & 160\dg 16.4\pr & 1 & 23 & 175\dg 18.8\pr \\
2 & 0 & 190\dg 21.3\pr & 2 & 1 & 205\dg 23.7\pr & 2 & 2 & 220\dg 26.2\pr & 2 & 3 & 235\dg 28.7\pr \\
2 & 4 & 250\dg 31.1\pr & 2 & 5 & 265\dg 33.6\pr & 2 & 6 & 280\dg 36.1\pr & 2 & 7 & 295\dg 38.5\pr \\
2 & 8 & 310\dg 41.0\pr & 2 & 9 & 325\dg 43.5\pr & 2 & 10 & 340\dg 45.9\pr & 2 & 11 & 355\dg 48.4\pr \\
2 & 12 & 10\dg 50.9\pr & 2 & 13 & 25\dg 53.3\pr & 2 & 14 & 40\dg 55.8\pr & 2 & 15 & 55\dg 58.2\pr \\
2 & 16 & 71\dg 00.7\pr & 2 & 17 & 86\dg 03.2\pr & 2 & 18 & 101\dg 05.6\pr & 2 & 19 & 116\dg 08.1\pr \\
2 & 20 & 131\dg 10.6\pr & 2 & 21 & 146\dg 13.0\pr & 2 & 22 & 161\dg 15.5\pr & 2 & 23 & 176\dg 18.0\pr \\
3 & 0 & 191\dg 20.4\pr & 3 & 1 & 206\dg 22.9\pr & 3 & 2 & 221\dg 25.4\pr & 3 & 3 & 236\dg 27.8\pr \\
3 & 4 & 251\dg 30.3\pr & 3 & 5 & 266\dg 32.7\pr & 3 & 6 & 281\dg 35.2\pr & 3 & 7 & 296\dg 37.7\pr \\
3 & 8 & 311\dg 40.1\pr & 3 & 9 & 326\dg 42.6\pr & 3 & 10 & 341\dg 45.1\pr & 3 & 11 & 356\dg 47.5\pr \\
3 & 12 & 11\dg 50.0\pr & 3 & 13 & 26\dg 52.5\pr & 3 & 14 & 41\dg 54.9\pr & 3 & 15 & 56\dg 57.4\pr \\
3 & 16 & 71\dg 59.8\pr & 3 & 17 & 87\dg 02.3\pr & 3 & 18 & 102\dg 04.8\pr & 3 & 19 & 117\dg 07.2\pr \\
3 & 20 & 132\dg 09.7\pr & 3 & 21 & 147\dg 12.2\pr & 3 & 22 & 162\dg 14.6\pr & 3 & 23 & 177\dg 17.1\pr \\
4 & 0 & 192\dg 19.6\pr & 4 & 1 & 207\dg 22.0\pr & 4 & 2 & 222\dg 24.5\pr & 4 & 3 & 237\dg 27.0\pr \\
4 & 4 & 252\dg 29.4\pr & 4 & 5 & 267\dg 31.9\pr & 4 & 6 & 282\dg 34.3\pr & 4 & 7 & 297\dg 36.8\pr \\
4 & 8 & 312\dg 39.3\pr & 4 & 9 & 327\dg 41.7\pr & 4 & 10 & 342\dg 44.2\pr & 4 & 11 & 357\dg 46.7\pr \\
4 & 12 & 12\dg 49.1\pr & 4 & 13 & 27\dg 51.6\pr & 4 & 14 & 42\dg 54.1\pr & 4 & 15 & 57\dg 56.5\pr \\
4 & 16 & 72\dg 59.0\pr & 4 & 17 & 88\dg 01.5\pr & 4 & 18 & 103\dg 03.9\pr & 4 & 19 & 118\dg 06.4\pr \\
4 & 20 & 133\dg 08.8\pr & 4 & 21 & 148\dg 11.3\pr & 4 & 22 & 163\dg 13.8\pr & 4 & 23 & 178\dg 16.2\pr \\
5 & 0 & 193\dg 18.7\pr & 5 & 1 & 208\dg 21.2\pr & 5 & 2 & 223\dg 23.6\pr & 5 & 3 & 238\dg 26.1\pr \\
5 & 4 & 253\dg 28.6\pr & 5 & 5 & 268\dg 31.0\pr & 5 & 6 & 283\dg 33.5\pr & 5 & 7 & 298\dg 35.9\pr \\
5 & 8 & 313\dg 38.4\pr & 5 & 9 & 328\dg 40.9\pr & 5 & 10 & 343\dg 43.3\pr & 5 & 11 & 358\dg 45.8\pr \\
5 & 12 & 13\dg 48.3\pr & 5 & 13 & 28\dg 50.7\pr & 5 & 14 & 43\dg 53.2\pr & 5 & 15 & 58\dg 55.7\pr \\
5 & 16 & 73\dg 58.1\pr & 5 & 17 & 89\dg 00.6\pr & 5 & 18 & 104\dg 03.1\pr & 5 & 19 & 119\dg 05.5\pr \\
5 & 20 & 134\dg 08.0\pr & 5 & 21 & 149\dg 10.4\pr & 5 & 22 & 164\dg 12.9\pr & 5 & 23 & 179\dg 15.4\pr \\
6 & 0 & 194\dg 17.8\pr & 6 & 1 & 209\dg 20.3\pr & 6 & 2 & 224\dg 22.8\pr & 6 & 3 & 239\dg 25.2\pr \\
6 & 4 & 254\dg 27.7\pr & 6 & 5 & 269\dg 30.2\pr & 6 & 6 & 284\dg 32.6\pr & 6 & 7 & 299\dg 35.1\pr \\
6 & 8 & 314\dg 37.6\pr & 6 & 9 & 329\dg 40.0\pr & 6 & 10 & 344\dg 42.5\pr & 6 & 11 & 359\dg 44.9\pr \\
6 & 12 & 14\dg 47.4\pr & 6 & 13 & 29\dg 49.9\pr & 6 & 14 & 44\dg 52.3\pr & 6 & 15 & 59\dg 54.8\pr \\
6 & 16 & 74\dg 57.3\pr & 6 & 17 & 89\dg 59.7\pr & 6 & 18 & 105\dg 02.2\pr & 6 & 19 & 120\dg 04.7\pr \\
6 & 20 & 135\dg 07.1\pr & 6 & 21 & 150\dg 09.6\pr & 6 & 22 & 165\dg 12.0\pr & 6 & 23 & 180\dg 14.5\pr \\
7 & 0 & 195\dg 17.0\pr & 7 & 1 & 210\dg 19.4\pr & 7 & 2 & 225\dg 21.9\pr & 7 & 3 & 240\dg 24.4\pr \\
7 & 4 & 255\dg 26.8\pr & 7 & 5 & 270\dg 29.3\pr & 7 & 6 & 285\dg 31.8\pr & 7 & 7 & 300\dg 34.2\pr \\
7 & 8 & 315\dg 36.7\pr & 7 & 9 & 330\dg 39.2\pr & 7 & 10 & 345\dg 41.6\pr & 7 & 11 & 0\dg 44.1\pr \\
7 & 12 & 15\dg 46.5\pr & 7 & 13 & 30\dg 49.0\pr & 7 & 14 & 45\dg 51.5\pr & 7 & 15 & 60\dg 53.9\pr \\
7 & 16 & 75\dg 56.4\pr & 7 & 17 & 90\dg 58.9\pr & 7 & 18 & 106\dg 01.3\pr & 7 & 19 & 121\dg 03.8\pr \\
7 & 20 & 136\dg 06.3\pr & 7 & 21 & 151\dg 08.7\pr & 7 & 22 & 166\dg 11.2\pr & 7 & 23 & 181\dg 13.7\pr \\
8 & 0 & 196\dg 16.1\pr & 8 & 1 & 211\dg 18.6\pr & 8 & 2 & 226\dg 21.0\pr & 8 & 3 & 241\dg 23.5\pr \\
8 & 4 & 256\dg 26.0\pr & 8 & 5 & 271\dg 28.4\pr & 8 & 6 & 286\dg 30.9\pr & 8 & 7 & 301\dg 33.4\pr \\
8 & 8 & 316\dg 35.8\pr & 8 & 9 & 331\dg 38.3\pr & 8 & 10 & 346\dg 40.8\pr & 8 & 11 & 1\dg 43.2\pr \\
8 & 12 & 16\dg 45.7\pr & 8 & 13 & 31\dg 48.2\pr & 8 & 14 & 46\dg 50.6\pr & 8 & 15 & 61\dg 53.1\pr \\
8 & 16 & 76\dg 55.5\pr & 8 & 17 & 91\dg 58.0\pr & 8 & 18 & 107\dg 00.5\pr & 8 & 19 & 122\dg 02.9\pr \\
8 & 20 & 137\dg 05.4\pr & 8 & 21 & 152\dg 07.9\pr & 8 & 22 & 167\dg 10.3\pr & 8 & 23 & 182\dg 12.8\pr \\
9 & 0 & 197\dg 15.3\pr & 9 & 1 & 212\dg 17.7\pr & 9 & 2 & 227\dg 20.2\pr & 9 & 3 & 242\dg 22.6\pr \\
9 & 4 & 257\dg 25.1\pr & 9 & 5 & 272\dg 27.6\pr & 9 & 6 & 287\dg 30.0\pr & 9 & 7 & 302\dg 32.5\pr \\
9 & 8 & 317\dg 35.0\pr & 9 & 9 & 332\dg 37.4\pr & 9 & 10 & 347\dg 39.9\pr & 9 & 11 & 2\dg 42.4\pr \\
9 & 12 & 17\dg 44.8\pr & 9 & 13 & 32\dg 47.3\pr & 9 & 14 & 47\dg 49.8\pr & 9 & 15 & 62\dg 52.2\pr \\
9 & 16 & 77\dg 54.7\pr & 9 & 17 & 92\dg 57.1\pr & 9 & 18 & 107\dg 59.6\pr & 9 & 19 & 123\dg 02.1\pr \\
9 & 20 & 138\dg 04.5\pr & 9 & 21 & 153\dg 07.0\pr & 9 & 22 & 168\dg 09.5\pr & 9 & 23 & 183\dg 11.9\pr \\
10 & 0 & 198\dg 14.4\pr & 10 & 1 & 213\dg 16.9\pr & 10 & 2 & 228\dg 19.3\pr & 10 & 3 & 243\dg 21.8\pr \\
10 & 4 & 258\dg 24.3\pr & 10 & 5 & 273\dg 26.7\pr & 10 & 6 & 288\dg 29.2\pr & 10 & 7 & 303\dg 31.6\pr \\
10 & 8 & 318\dg 34.1\pr & 10 & 9 & 333\dg 36.6\pr & 10 & 10 & 348\dg 39.0\pr & 10 & 11 & 3\dg 41.5\pr \\
10 & 12 & 18\dg 44.0\pr & 10 & 13 & 33\dg 46.4\pr & 10 & 14 & 48\dg 48.9\pr & 10 & 15 & 63\dg 51.4\pr \\
10 & 16 & 78\dg 53.8\pr & 10 & 17 & 93\dg 56.3\pr & 10 & 18 & 108\dg 58.7\pr & 10 & 19 & 124\dg 01.2\pr \\
10 & 20 & 139\dg 03.7\pr & 10 & 21 & 154\dg 06.1\pr & 10 & 22 & 169\dg 08.6\pr & 10 & 23 & 184\dg 11.1\pr \\
11 & 0 & 199\dg 13.5\pr & 11 & 1 & 214\dg 16.0\pr & 11 & 2 & 229\dg 18.5\pr & 11 & 3 & 244\dg 20.9\pr \\
11 & 4 & 259\dg 23.4\pr & 11 & 5 & 274\dg 25.9\pr & 11 & 6 & 289\dg 28.3\pr & 11 & 7 & 304\dg 30.8\pr \\
11 & 8 & 319\dg 33.2\pr & 11 & 9 & 334\dg 35.7\pr & 11 & 10 & 349\dg 38.2\pr & 11 & 11 & 4\dg 40.6\pr \\
11 & 12 & 19\dg 43.1\pr & 11 & 13 & 34\dg 45.6\pr & 11 & 14 & 49\dg 48.0\pr & 11 & 15 & 64\dg 50.5\pr \\
11 & 16 & 79\dg 53.0\pr & 11 & 17 & 94\dg 55.4\pr & 11 & 18 & 109\dg 57.9\pr & 11 & 19 & 125\dg 00.4\pr \\
11 & 20 & 140\dg 02.8\pr & 11 & 21 & 155\dg 05.3\pr & 11 & 22 & 170\dg 07.7\pr & 11 & 23 & 185\dg 10.2\pr \\
12 & 0 & 200\dg 12.7\pr & 12 & 1 & 215\dg 15.1\pr & 12 & 2 & 230\dg 17.6\pr & 12 & 3 & 245\dg 20.1\pr \\
12 & 4 & 260\dg 22.5\pr & 12 & 5 & 275\dg 25.0\pr & 12 & 6 & 290\dg 27.5\pr & 12 & 7 & 305\dg 29.9\pr \\
12 & 8 & 320\dg 32.4\pr & 12 & 9 & 335\dg 34.8\pr & 12 & 10 & 350\dg 37.3\pr & 12 & 11 & 5\dg 39.8\pr \\
12 & 12 & 20\dg 42.2\pr & 12 & 13 & 35\dg 44.7\pr & 12 & 14 & 50\dg 47.2\pr & 12 & 15 & 65\dg 49.6\pr \\
12 & 16 & 80\dg 52.1\pr & 12 & 17 & 95\dg 54.6\pr & 12 & 18 & 110\dg 57.0\pr & 12 & 19 & 125\dg 59.5\pr \\
12 & 20 & 141\dg 02.0\pr & 12 & 21 & 156\dg 04.4\pr & 12 & 22 & 171\dg 06.9\pr & 12 & 23 & 186\dg 09.3\pr \\
13 & 0 & 201\dg 11.8\pr & 13 & 1 & 216\dg 14.3\pr & 13 & 2 & 231\dg 16.7\pr & 13 & 3 & 246\dg 19.2\pr \\
13 & 4 & 261\dg 21.7\pr & 13 & 5 & 276\dg 24.1\pr & 13 & 6 & 291\dg 26.6\pr & 13 & 7 & 306\dg 29.1\pr \\
13 & 8 & 321\dg 31.5\pr & 13 & 9 & 336\dg 34.0\pr & 13 & 10 & 351\dg 36.5\pr & 13 & 11 & 6\dg 38.9\pr \\
13 & 12 & 21\dg 41.4\pr & 13 & 13 & 36\dg 43.8\pr & 13 & 14 & 51\dg 46.3\pr & 13 & 15 & 66\dg 48.8\pr \\
13 & 16 & 81\dg 51.2\pr & 13 & 17 & 96\dg 53.7\pr & 13 & 18 & 111\dg 56.2\pr & 13 & 19 & 126\dg 58.6\pr \\
13 & 20 & 142\dg 01.1\pr & 13 & 21 & 157\dg 03.6\pr & 13 & 22 & 172\dg 06.0\pr & 13 & 23 & 187\dg 08.5\pr \\
14 & 0 & 202\dg 10.9\pr & 14 & 1 & 217\dg 13.4\pr & 14 & 2 & 232\dg 15.9\pr & 14 & 3 & 247\dg 18.3\pr \\
14 & 4 & 262\dg 20.8\pr & 14 & 5 & 277\dg 23.3\pr & 14 & 6 & 292\dg 25.7\pr & 14 & 7 & 307\dg 28.2\pr \\
14 & 8 & 322\dg 30.7\pr & 14 & 9 & 337\dg 33.1\pr & 14 & 10 & 352\dg 35.6\pr & 14 & 11 & 7\dg 38.1\pr \\
14 & 12 & 22\dg 40.5\pr & 14 & 13 & 37\dg 43.0\pr & 14 & 14 & 52\dg 45.4\pr & 14 & 15 & 67\dg 47.9\pr \\
14 & 16 & 82\dg 50.4\pr & 14 & 17 & 97\dg 52.8\pr & 14 & 18 & 112\dg 55.3\pr & 14 & 19 & 127\dg 57.8\pr \\
14 & 20 & 143\dg 00.2\pr & 14 & 21 & 158\dg 02.7\pr & 14 & 22 & 173\dg 05.2\pr & 14 & 23 & 188\dg 07.6\pr \\
15 & 0 & 203\dg 10.1\pr & 15 & 1 & 218\dg 12.6\pr & 15 & 2 & 233\dg 15.0\pr & 15 & 3 & 248\dg 17.5\pr \\
15 & 4 & 263\dg 19.9\pr & 15 & 5 & 278\dg 22.4\pr & 15 & 6 & 293\dg 24.9\pr & 15 & 7 & 308\dg 27.3\pr \\
15 & 8 & 323\dg 29.8\pr & 15 & 9 & 338\dg 32.3\pr & 15 & 10 & 353\dg 34.7\pr & 15 & 11 & 8\dg 37.2\pr \\
15 & 12 & 23\dg 39.7\pr & 15 & 13 & 38\dg 42.1\pr & 15 & 14 & 53\dg 44.6\pr & 15 & 15 & 68\dg 47.1\pr \\
15 & 16 & 83\dg 49.5\pr & 15 & 17 & 98\dg 52.0\pr & 15 & 18 & 113\dg 54.4\pr & 15 & 19 & 128\dg 56.9\pr \\
15 & 20 & 143\dg 59.4\pr & 15 & 21 & 159\dg 01.8\pr & 15 & 22 & 174\dg 04.3\pr & 15 & 23 & 189\dg 06.8\pr \\
16 & 0 & 204\dg 09.2\pr & 16 & 1 & 219\dg 11.7\pr & 16 & 2 & 234\dg 14.2\pr & 16 & 3 & 249\dg 16.6\pr \\
16 & 4 & 264\dg 19.1\pr & 16 & 5 & 279\dg 21.5\pr & 16 & 6 & 294\dg 24.0\pr & 16 & 7 & 309\dg 26.5\pr \\
16 & 8 & 324\dg 28.9\pr & 16 & 9 & 339\dg 31.4\pr & 16 & 10 & 354\dg 33.9\pr & 16 & 11 & 9\dg 36.3\pr \\
16 & 12 & 24\dg 38.8\pr & 16 & 13 & 39\dg 41.3\pr & 16 & 14 & 54\dg 43.7\pr & 16 & 15 & 69\dg 46.2\pr \\
16 & 16 & 84\dg 48.7\pr & 16 & 17 & 99\dg 51.1\pr & 16 & 18 & 114\dg 53.6\pr & 16 & 19 & 129\dg 56.0\pr \\
16 & 20 & 144\dg 58.5\pr & 16 & 21 & 160\dg 01.0\pr & 16 & 22 & 175\dg 03.4\pr & 16 & 23 & 190\dg 05.9\pr \\
17 & 0 & 205\dg 08.4\pr & 17 & 1 & 220\dg 10.8\pr & 17 & 2 & 235\dg 13.3\pr & 17 & 3 & 250\dg 15.8\pr \\
17 & 4 & 265\dg 18.2\pr & 17 & 5 & 280\dg 20.7\pr & 17 & 6 & 295\dg 23.2\pr & 17 & 7 & 310\dg 25.6\pr \\
17 & 8 & 325\dg 28.1\pr & 17 & 9 & 340\dg 30.5\pr & 17 & 10 & 355\dg 33.0\pr & 17 & 11 & 10\dg 35.5\pr \\
17 & 12 & 25\dg 37.9\pr & 17 & 13 & 40\dg 40.4\pr & 17 & 14 & 55\dg 42.9\pr & 17 & 15 & 70\dg 45.3\pr \\
17 & 16 & 85\dg 47.8\pr & 17 & 17 & 100\dg 50.3\pr & 17 & 18 & 115\dg 52.7\pr & 17 & 19 & 130\dg 55.2\pr \\
17 & 20 & 145\dg 57.6\pr & 17 & 21 & 161\dg 00.1\pr & 17 & 22 & 176\dg 02.6\pr & 17 & 23 & 191\dg 05.0\pr \\
18 & 0 & 206\dg 07.5\pr & 18 & 1 & 221\dg 10.0\pr & 18 & 2 & 236\dg 12.4\pr & 18 & 3 & 251\dg 14.9\pr \\
18 & 4 & 266\dg 17.4\pr & 18 & 5 & 281\dg 19.8\pr & 18 & 6 & 296\dg 22.3\pr & 18 & 7 & 311\dg 24.8\pr \\
18 & 8 & 326\dg 27.2\pr & 18 & 9 & 341\dg 29.7\pr & 18 & 10 & 356\dg 32.1\pr & 18 & 11 & 11\dg 34.6\pr \\
18 & 12 & 26\dg 37.1\pr & 18 & 13 & 41\dg 39.5\pr & 18 & 14 & 56\dg 42.0\pr & 18 & 15 & 71\dg 44.5\pr \\
18 & 16 & 86\dg 46.9\pr & 18 & 17 & 101\dg 49.4\pr & 18 & 18 & 116\dg 51.9\pr & 18 & 19 & 131\dg 54.3\pr \\
18 & 20 & 146\dg 56.8\pr & 18 & 21 & 161\dg 59.3\pr & 18 & 22 & 177\dg 01.7\pr & 18 & 23 & 192\dg 04.2\pr \\
19 & 0 & 207\dg 06.6\pr & 19 & 1 & 222\dg 09.1\pr & 19 & 2 & 237\dg 11.6\pr & 19 & 3 & 252\dg 14.0\pr \\
19 & 4 & 267\dg 16.5\pr & 19 & 5 & 282\dg 19.0\pr & 19 & 6 & 297\dg 21.4\pr & 19 & 7 & 312\dg 23.9\pr \\
19 & 8 & 327\dg 26.4\pr & 19 & 9 & 342\dg 28.8\pr & 19 & 10 & 357\dg 31.3\pr & 19 & 11 & 12\dg 33.7\pr \\
19 & 12 & 27\dg 36.2\pr & 19 & 13 & 42\dg 38.7\pr & 19 & 14 & 57\dg 41.1\pr & 19 & 15 & 72\dg 43.6\pr \\
19 & 16 & 87\dg 46.1\pr & 19 & 17 & 102\dg 48.5\pr & 19 & 18 & 117\dg 51.0\pr & 19 & 19 & 132\dg 53.5\pr \\
19 & 20 & 147\dg 55.9\pr & 19 & 21 & 162\dg 58.4\pr & 19 & 22 & 178\dg 00.9\pr & 19 & 23 & 193\dg 03.3\pr \\
20 & 0 & 208\dg 05.8\pr & 20 & 1 & 223\dg 08.2\pr & 20 & 2 & 238\dg 10.7\pr & 20 & 3 & 253\dg 13.2\pr \\
20 & 4 & 268\dg 15.6\pr & 20 & 5 & 283\dg 18.1\pr & 20 & 6 & 298\dg 20.6\pr & 20 & 7 & 313\dg 23.0\pr \\
20 & 8 & 328\dg 25.5\pr & 20 & 9 & 343\dg 28.0\pr & 20 & 10 & 358\dg 30.4\pr & 20 & 11 & 13\dg 32.9\pr \\
20 & 12 & 28\dg 35.4\pr & 20 & 13 & 43\dg 37.8\pr & 20 & 14 & 58\dg 40.3\pr & 20 & 15 & 73\dg 42.7\pr \\
20 & 16 & 88\dg 45.2\pr & 20 & 17 & 103\dg 47.7\pr & 20 & 18 & 118\dg 50.1\pr & 20 & 19 & 133\dg 52.6\pr \\
20 & 20 & 148\dg 55.1\pr & 20 & 21 & 163\dg 57.5\pr & 20 & 22 & 179\dg 00.0\pr & 20 & 23 & 194\dg 02.5\pr \\
21 & 0 & 209\dg 04.9\pr & 21 & 1 & 224\dg 07.4\pr & 21 & 2 & 239\dg 09.8\pr & 21 & 3 & 254\dg 12.3\pr \\
21 & 4 & 269\dg 14.8\pr & 21 & 5 & 284\dg 17.2\pr & 21 & 6 & 299\dg 19.7\pr & 21 & 7 & 314\dg 22.2\pr \\
21 & 8 & 329\dg 24.6\pr & 21 & 9 & 344\dg 27.1\pr & 21 & 10 & 359\dg 29.6\pr & 21 & 11 & 14\dg 32.0\pr \\
21 & 12 & 29\dg 34.5\pr & 21 & 13 & 44\dg 37.0\pr & 21 & 14 & 59\dg 39.4\pr & 21 & 15 & 74\dg 41.9\pr \\
21 & 16 & 89\dg 44.3\pr & 21 & 17 & 104\dg 46.8\pr & 21 & 18 & 119\dg 49.3\pr & 21 & 19 & 134\dg 51.7\pr \\
21 & 20 & 149\dg 54.2\pr & 21 & 21 & 164\dg 56.7\pr & 21 & 22 & 179\dg 59.1\pr & 21 & 23 & 195\dg 01.6\pr \\
22 & 0 & 210\dg 04.1\pr & 22 & 1 & 225\dg 06.5\pr & 22 & 2 & 240\dg 09.0\pr & 22 & 3 & 255\dg 11.5\pr \\
22 & 4 & 270\dg 13.9\pr & 22 & 5 & 285\dg 16.4\pr & 22 & 6 & 300\dg 18.8\pr & 22 & 7 & 315\dg 21.3\pr \\
22 & 8 & 330\dg 23.8\pr & 22 & 9 & 345\dg 26.2\pr & 22 & 10 & 0\dg 28.7\pr & 22 & 11 & 15\dg 31.2\pr \\
22 & 12 & 30\dg 33.6\pr & 22 & 13 & 45\dg 36.1\pr & 22 & 14 & 60\dg 38.6\pr & 22 & 15 & 75\dg 41.0\pr \\
22 & 16 & 90\dg 43.5\pr & 22 & 17 & 105\dg 46.0\pr & 22 & 18 & 120\dg 48.4\pr & 22 & 19 & 135\dg 50.9\pr \\
22 & 20 & 150\dg 53.3\pr & 22 & 21 & 165\dg 55.8\pr & 22 & 22 & 180\dg 58.3\pr & 22 & 23 & 196\dg 00.7\pr \\
23 & 0 & 211\dg 03.2\pr & 23 & 1 & 226\dg 05.7\pr & 23 & 2 & 241\dg 08.1\pr & 23 & 3 & 256\dg 10.6\pr \\
23 & 4 & 271\dg 13.1\pr & 23 & 5 & 286\dg 15.5\pr & 23 & 6 & 301\dg 18.0\pr & 23 & 7 & 316\dg 20.4\pr \\
23 & 8 & 331\dg 22.9\pr & 23 & 9 & 346\dg 25.4\pr & 23 & 10 & 1\dg 27.8\pr & 23 & 11 & 16\dg 30.3\pr \\
23 & 12 & 31\dg 32.8\pr & 23 & 13 & 46\dg 35.2\pr & 23 & 14 & 61\dg 37.7\pr & 23 & 15 & 76\dg 40.2\pr \\
23 & 16 & 91\dg 42.6\pr & 23 & 17 & 106\dg 45.1\pr & 23 & 18 & 121\dg 47.6\pr & 23 & 19 & 136\dg 50.0\pr \\
23 & 20 & 151\dg 52.5\pr & 23 & 21 & 166\dg 54.9\pr & 23 & 22 & 181\dg 57.4\pr & 23 & 23 & 196\dg 59.9\pr \\
24 & 0 & 212\dg 02.3\pr & 24 & 1 & 227\dg 04.8\pr & 24 & 2 & 242\dg 07.3\pr & 24 & 3 & 257\dg 09.7\pr \\
24 & 4 & 272\dg 12.2\pr & 24 & 5 & 287\dg 14.7\pr & 24 & 6 & 302\dg 17.1\pr & 24 & 7 & 317\dg 19.6\pr \\
24 & 8 & 332\dg 22.1\pr & 24 & 9 & 347\dg 24.5\pr & 24 & 10 & 2\dg 27.0\pr & 24 & 11 & 17\dg 29.4\pr \\
24 & 12 & 32\dg 31.9\pr & 24 & 13 & 47\dg 34.4\pr & 24 & 14 & 62\dg 36.8\pr & 24 & 15 & 77\dg 39.3\pr \\
24 & 16 & 92\dg 41.8\pr & 24 & 17 & 107\dg 44.2\pr & 24 & 18 & 122\dg 46.7\pr & 24 & 19 & 137\dg 49.2\pr \\
24 & 20 & 152\dg 51.6\pr & 24 & 21 & 167\dg 54.1\pr & 24 & 22 & 182\dg 56.5\pr & 24 & 23 & 197\dg 59.0\pr \\
25 & 0 & 213\dg 01.5\pr & 25 & 1 & 228\dg 03.9\pr & 25 & 2 & 243\dg 06.4\pr & 25 & 3 & 258\dg 08.9\pr \\
25 & 4 & 273\dg 11.3\pr & 25 & 5 & 288\dg 13.8\pr & 25 & 6 & 303\dg 16.3\pr & 25 & 7 & 318\dg 18.7\pr \\
25 & 8 & 333\dg 21.2\pr & 25 & 9 & 348\dg 23.7\pr & 25 & 10 & 3\dg 26.1\pr & 25 & 11 & 18\dg 28.6\pr \\
25 & 12 & 33\dg 31.0\pr & 25 & 13 & 48\dg 33.5\pr & 25 & 14 & 63\dg 36.0\pr & 25 & 15 & 78\dg 38.4\pr \\
25 & 16 & 93\dg 40.9\pr & 25 & 17 & 108\dg 43.4\pr & 25 & 18 & 123\dg 45.8\pr & 25 & 19 & 138\dg 48.3\pr \\
25 & 20 & 153\dg 50.8\pr & 25 & 21 & 168\dg 53.2\pr & 25 & 22 & 183\dg 55.7\pr & 25 & 23 & 198\dg 58.2\pr \\
26 & 0 & 214\dg 00.6\pr & 26 & 1 & 229\dg 03.1\pr & 26 & 2 & 244\dg 05.5\pr & 26 & 3 & 259\dg 08.0\pr \\
26 & 4 & 274\dg 10.5\pr & 26 & 5 & 289\dg 12.9\pr & 26 & 6 & 304\dg 15.4\pr & 26 & 7 & 319\dg 17.9\pr \\
26 & 8 & 334\dg 20.3\pr & 26 & 9 & 349\dg 22.8\pr & 26 & 10 & 4\dg 25.3\pr & 26 & 11 & 19\dg 27.7\pr \\
26 & 12 & 34\dg 30.2\pr & 26 & 13 & 49\dg 32.6\pr & 26 & 14 & 64\dg 35.1\pr & 26 & 15 & 79\dg 37.6\pr \\
26 & 16 & 94\dg 40.0\pr & 26 & 17 & 109\dg 42.5\pr & 26 & 18 & 124\dg 45.0\pr & 26 & 19 & 139\dg 47.4\pr \\
26 & 20 & 154\dg 49.9\pr & 26 & 21 & 169\dg 52.4\pr & 26 & 22 & 184\dg 54.8\pr & 26 & 23 & 199\dg 57.3\pr \\
27 & 0 & 214\dg 59.8\pr & 27 & 1 & 230\dg 02.2\pr & 27 & 2 & 245\dg 04.7\pr & 27 & 3 & 260\dg 07.1\pr \\
27 & 4 & 275\dg 09.6\pr & 27 & 5 & 290\dg 12.1\pr & 27 & 6 & 305\dg 14.5\pr & 27 & 7 & 320\dg 17.0\pr \\
27 & 8 & 335\dg 19.5\pr & 27 & 9 & 350\dg 21.9\pr & 27 & 10 & 5\dg 24.4\pr & 27 & 11 & 20\dg 26.9\pr \\
27 & 12 & 35\dg 29.3\pr & 27 & 13 & 50\dg 31.8\pr & 27 & 14 & 65\dg 34.3\pr & 27 & 15 & 80\dg 36.7\pr \\
27 & 16 & 95\dg 39.2\pr & 27 & 17 & 110\dg 41.6\pr & 27 & 18 & 125\dg 44.1\pr & 27 & 19 & 140\dg 46.6\pr \\
27 & 20 & 155\dg 49.0\pr & 27 & 21 & 170\dg 51.5\pr & 27 & 22 & 185\dg 54.0\pr & 27 & 23 & 200\dg 56.4\pr \\
28 & 0 & 215\dg 58.9\pr & 28 & 1 & 231\dg 01.4\pr & 28 & 2 & 246\dg 03.8\pr & 28 & 3 & 261\dg 06.3\pr \\
28 & 4 & 276\dg 08.8\pr & 28 & 5 & 291\dg 11.2\pr & 28 & 6 & 306\dg 13.7\pr & 28 & 7 & 321\dg 16.1\pr \\
28 & 8 & 336\dg 18.6\pr & 28 & 9 & 351\dg 21.1\pr & 28 & 10 & 6\dg 23.5\pr & 28 & 11 & 21\dg 26.0\pr \\
28 & 12 & 36\dg 28.5\pr & 28 & 13 & 51\dg 30.9\pr & 28 & 14 & 66\dg 33.4\pr & 28 & 15 & 81\dg 35.9\pr \\
28 & 16 & 96\dg 38.3\pr & 28 & 17 & 111\dg 40.8\pr & 28 & 18 & 126\dg 43.2\pr & 28 & 19 & 141\dg 45.7\pr \\
28 & 20 & 156\dg 48.2\pr & 28 & 21 & 171\dg 50.6\pr & 28 & 22 & 186\dg 53.1\pr & 28 & 23 & 201\dg 55.6\pr \\
29 & 0 & 216\dg 58.0\pr & 29 & 1 & 232\dg 00.5\pr & 29 & 2 & 247\dg 03.0\pr & 29 & 3 & 262\dg 05.4\pr \\
29 & 4 & 277\dg 07.9\pr & 29 & 5 & 292\dg 10.4\pr & 29 & 6 & 307\dg 12.8\pr & 29 & 7 & 322\dg 15.3\pr \\
29 & 8 & 337\dg 17.7\pr & 29 & 9 & 352\dg 20.2\pr & 29 & 10 & 7\dg 22.7\pr & 29 & 11 & 22\dg 25.1\pr \\
29 & 12 & 37\dg 27.6\pr & 29 & 13 & 52\dg 30.1\pr & 29 & 14 & 67\dg 32.5\pr & 29 & 15 & 82\dg 35.0\pr \\
29 & 16 & 97\dg 37.5\pr & 29 & 17 & 112\dg 39.9\pr & 29 & 18 & 127\dg 42.4\pr & 29 & 19 & 142\dg 44.9\pr \\
29 & 20 & 157\dg 47.3\pr & 29 & 21 & 172\dg 49.8\pr & 29 & 22 & 187\dg 52.2\pr & 29 & 23 & 202\dg 54.7\pr \\
30 & 0 & 217\dg 57.2\pr & 30 & 1 & 232\dg 59.6\pr & 30 & 2 & 248\dg 02.1\pr & 30 & 3 & 263\dg 04.6\pr \\
30 & 4 & 278\dg 07.0\pr & 30 & 5 & 293\dg 09.5\pr & 30 & 6 & 308\dg 12.0\pr & 30 & 7 & 323\dg 14.4\pr \\
30 & 8 & 338\dg 16.9\pr & 30 & 9 & 353\dg 19.3\pr & 30 & 10 & 8\dg 21.8\pr & 30 & 11 & 23\dg 24.3\pr \\
30 & 12 & 38\dg 26.7\pr & 30 & 13 & 53\dg 29.2\pr & 30 & 14 & 68\dg 31.7\pr & 30 & 15 & 83\dg 34.1\pr \\
30 & 16 & 98\dg 36.6\pr & 30 & 17 & 113\dg 39.1\pr & 30 & 18 & 128\dg 41.5\pr & 30 & 19 & 143\dg 44.0\pr \\
30 & 20 & 158\dg 46.5\pr & 30 & 21 & 173\dg 48.9\pr & 30 & 22 & 188\dg 51.4\pr & 30 & 23 & 203\dg 53.8\pr \\
\end{longtable}}

{\footnotesize
\begin{longtable}{rrr|rrr|rrr|rrr}
\caption{Aries --- GHA of the First Point of Aries, May 2026 (UT)}\label{aries05}\\
\toprule
\textbf{d} & \textbf{h} & \textbf{GHA} $\gamma$ & \textbf{d} & \textbf{h} & \textbf{GHA} $\gamma$ & \textbf{d} & \textbf{h} & \textbf{GHA} $\gamma$ & \textbf{d} & \textbf{h} & \textbf{GHA} $\gamma$ \\
\midrule
\endfirsthead
\multicolumn{12}{c}{\textit{Aries --- GHA of the First Point of Aries, May 2026 (UT) (continued)}}\\
\toprule
\textbf{d} & \textbf{h} & \textbf{GHA} $\gamma$ & \textbf{d} & \textbf{h} & \textbf{GHA} $\gamma$ & \textbf{d} & \textbf{h} & \textbf{GHA} $\gamma$ & \textbf{d} & \textbf{h} & \textbf{GHA} $\gamma$ \\
\midrule
\endhead
\midrule\multicolumn{12}{r}{\textit{continued on next page}}\\
\endfoot
\bottomrule
\endlastfoot
1 & 0 & 218\dg 56.3\pr & 1 & 1 & 233\dg 58.8\pr & 1 & 2 & 249\dg 01.2\pr & 1 & 3 & 264\dg 03.7\pr \\
1 & 4 & 279\dg 06.2\pr & 1 & 5 & 294\dg 08.6\pr & 1 & 6 & 309\dg 11.1\pr & 1 & 7 & 324\dg 13.6\pr \\
1 & 8 & 339\dg 16.0\pr & 1 & 9 & 354\dg 18.5\pr & 1 & 10 & 9\dg 21.0\pr & 1 & 11 & 24\dg 23.4\pr \\
1 & 12 & 39\dg 25.9\pr & 1 & 13 & 54\dg 28.3\pr & 1 & 14 & 69\dg 30.8\pr & 1 & 15 & 84\dg 33.3\pr \\
1 & 16 & 99\dg 35.7\pr & 1 & 17 & 114\dg 38.2\pr & 1 & 18 & 129\dg 40.7\pr & 1 & 19 & 144\dg 43.1\pr \\
1 & 20 & 159\dg 45.6\pr & 1 & 21 & 174\dg 48.1\pr & 1 & 22 & 189\dg 50.5\pr & 1 & 23 & 204\dg 53.0\pr \\
2 & 0 & 219\dg 55.4\pr & 2 & 1 & 234\dg 57.9\pr & 2 & 2 & 250\dg 00.4\pr & 2 & 3 & 265\dg 02.8\pr \\
2 & 4 & 280\dg 05.3\pr & 2 & 5 & 295\dg 07.8\pr & 2 & 6 & 310\dg 10.2\pr & 2 & 7 & 325\dg 12.7\pr \\
2 & 8 & 340\dg 15.2\pr & 2 & 9 & 355\dg 17.6\pr & 2 & 10 & 10\dg 20.1\pr & 2 & 11 & 25\dg 22.6\pr \\
2 & 12 & 40\dg 25.0\pr & 2 & 13 & 55\dg 27.5\pr & 2 & 14 & 70\dg 29.9\pr & 2 & 15 & 85\dg 32.4\pr \\
2 & 16 & 100\dg 34.9\pr & 2 & 17 & 115\dg 37.3\pr & 2 & 18 & 130\dg 39.8\pr & 2 & 19 & 145\dg 42.3\pr \\
2 & 20 & 160\dg 44.7\pr & 2 & 21 & 175\dg 47.2\pr & 2 & 22 & 190\dg 49.7\pr & 2 & 23 & 205\dg 52.1\pr \\
3 & 0 & 220\dg 54.6\pr & 3 & 1 & 235\dg 57.1\pr & 3 & 2 & 250\dg 59.5\pr & 3 & 3 & 266\dg 02.0\pr \\
3 & 4 & 281\dg 04.4\pr & 3 & 5 & 296\dg 06.9\pr & 3 & 6 & 311\dg 09.4\pr & 3 & 7 & 326\dg 11.8\pr \\
3 & 8 & 341\dg 14.3\pr & 3 & 9 & 356\dg 16.8\pr & 3 & 10 & 11\dg 19.2\pr & 3 & 11 & 26\dg 21.7\pr \\
3 & 12 & 41\dg 24.2\pr & 3 & 13 & 56\dg 26.6\pr & 3 & 14 & 71\dg 29.1\pr & 3 & 15 & 86\dg 31.5\pr \\
3 & 16 & 101\dg 34.0\pr & 3 & 17 & 116\dg 36.5\pr & 3 & 18 & 131\dg 38.9\pr & 3 & 19 & 146\dg 41.4\pr \\
3 & 20 & 161\dg 43.9\pr & 3 & 21 & 176\dg 46.3\pr & 3 & 22 & 191\dg 48.8\pr & 3 & 23 & 206\dg 51.3\pr \\
4 & 0 & 221\dg 53.7\pr & 4 & 1 & 236\dg 56.2\pr & 4 & 2 & 251\dg 58.7\pr & 4 & 3 & 267\dg 01.1\pr \\
4 & 4 & 282\dg 03.6\pr & 4 & 5 & 297\dg 06.0\pr & 4 & 6 & 312\dg 08.5\pr & 4 & 7 & 327\dg 11.0\pr \\
4 & 8 & 342\dg 13.4\pr & 4 & 9 & 357\dg 15.9\pr & 4 & 10 & 12\dg 18.4\pr & 4 & 11 & 27\dg 20.8\pr \\
4 & 12 & 42\dg 23.3\pr & 4 & 13 & 57\dg 25.8\pr & 4 & 14 & 72\dg 28.2\pr & 4 & 15 & 87\dg 30.7\pr \\
4 & 16 & 102\dg 33.2\pr & 4 & 17 & 117\dg 35.6\pr & 4 & 18 & 132\dg 38.1\pr & 4 & 19 & 147\dg 40.5\pr \\
4 & 20 & 162\dg 43.0\pr & 4 & 21 & 177\dg 45.5\pr & 4 & 22 & 192\dg 47.9\pr & 4 & 23 & 207\dg 50.4\pr \\
5 & 0 & 222\dg 52.9\pr & 5 & 1 & 237\dg 55.3\pr & 5 & 2 & 252\dg 57.8\pr & 5 & 3 & 268\dg 00.3\pr \\
5 & 4 & 283\dg 02.7\pr & 5 & 5 & 298\dg 05.2\pr & 5 & 6 & 313\dg 07.7\pr & 5 & 7 & 328\dg 10.1\pr \\
5 & 8 & 343\dg 12.6\pr & 5 & 9 & 358\dg 15.0\pr & 5 & 10 & 13\dg 17.5\pr & 5 & 11 & 28\dg 20.0\pr \\
5 & 12 & 43\dg 22.4\pr & 5 & 13 & 58\dg 24.9\pr & 5 & 14 & 73\dg 27.4\pr & 5 & 15 & 88\dg 29.8\pr \\
5 & 16 & 103\dg 32.3\pr & 5 & 17 & 118\dg 34.8\pr & 5 & 18 & 133\dg 37.2\pr & 5 & 19 & 148\dg 39.7\pr \\
5 & 20 & 163\dg 42.1\pr & 5 & 21 & 178\dg 44.6\pr & 5 & 22 & 193\dg 47.1\pr & 5 & 23 & 208\dg 49.5\pr \\
6 & 0 & 223\dg 52.0\pr & 6 & 1 & 238\dg 54.5\pr & 6 & 2 & 253\dg 56.9\pr & 6 & 3 & 268\dg 59.4\pr \\
6 & 4 & 284\dg 01.9\pr & 6 & 5 & 299\dg 04.3\pr & 6 & 6 & 314\dg 06.8\pr & 6 & 7 & 329\dg 09.3\pr \\
6 & 8 & 344\dg 11.7\pr & 6 & 9 & 359\dg 14.2\pr & 6 & 10 & 14\dg 16.6\pr & 6 & 11 & 29\dg 19.1\pr \\
6 & 12 & 44\dg 21.6\pr & 6 & 13 & 59\dg 24.0\pr & 6 & 14 & 74\dg 26.5\pr & 6 & 15 & 89\dg 29.0\pr \\
6 & 16 & 104\dg 31.4\pr & 6 & 17 & 119\dg 33.9\pr & 6 & 18 & 134\dg 36.4\pr & 6 & 19 & 149\dg 38.8\pr \\
6 & 20 & 164\dg 41.3\pr & 6 & 21 & 179\dg 43.8\pr & 6 & 22 & 194\dg 46.2\pr & 6 & 23 & 209\dg 48.7\pr \\
7 & 0 & 224\dg 51.1\pr & 7 & 1 & 239\dg 53.6\pr & 7 & 2 & 254\dg 56.1\pr & 7 & 3 & 269\dg 58.5\pr \\
7 & 4 & 285\dg 01.0\pr & 7 & 5 & 300\dg 03.5\pr & 7 & 6 & 315\dg 05.9\pr & 7 & 7 & 330\dg 08.4\pr \\
7 & 8 & 345\dg 10.9\pr & 7 & 9 & 0\dg 13.3\pr & 7 & 10 & 15\dg 15.8\pr & 7 & 11 & 30\dg 18.2\pr \\
7 & 12 & 45\dg 20.7\pr & 7 & 13 & 60\dg 23.2\pr & 7 & 14 & 75\dg 25.6\pr & 7 & 15 & 90\dg 28.1\pr \\
7 & 16 & 105\dg 30.6\pr & 7 & 17 & 120\dg 33.0\pr & 7 & 18 & 135\dg 35.5\pr & 7 & 19 & 150\dg 38.0\pr \\
7 & 20 & 165\dg 40.4\pr & 7 & 21 & 180\dg 42.9\pr & 7 & 22 & 195\dg 45.4\pr & 7 & 23 & 210\dg 47.8\pr \\
8 & 0 & 225\dg 50.3\pr & 8 & 1 & 240\dg 52.7\pr & 8 & 2 & 255\dg 55.2\pr & 8 & 3 & 270\dg 57.7\pr \\
8 & 4 & 286\dg 00.1\pr & 8 & 5 & 301\dg 02.6\pr & 8 & 6 & 316\dg 05.1\pr & 8 & 7 & 331\dg 07.5\pr \\
8 & 8 & 346\dg 10.0\pr & 8 & 9 & 1\dg 12.5\pr & 8 & 10 & 16\dg 14.9\pr & 8 & 11 & 31\dg 17.4\pr \\
8 & 12 & 46\dg 19.9\pr & 8 & 13 & 61\dg 22.3\pr & 8 & 14 & 76\dg 24.8\pr & 8 & 15 & 91\dg 27.2\pr \\
8 & 16 & 106\dg 29.7\pr & 8 & 17 & 121\dg 32.2\pr & 8 & 18 & 136\dg 34.6\pr & 8 & 19 & 151\dg 37.1\pr \\
8 & 20 & 166\dg 39.6\pr & 8 & 21 & 181\dg 42.0\pr & 8 & 22 & 196\dg 44.5\pr & 8 & 23 & 211\dg 47.0\pr \\
9 & 0 & 226\dg 49.4\pr & 9 & 1 & 241\dg 51.9\pr & 9 & 2 & 256\dg 54.3\pr & 9 & 3 & 271\dg 56.8\pr \\
9 & 4 & 286\dg 59.3\pr & 9 & 5 & 302\dg 01.7\pr & 9 & 6 & 317\dg 04.2\pr & 9 & 7 & 332\dg 06.7\pr \\
9 & 8 & 347\dg 09.1\pr & 9 & 9 & 2\dg 11.6\pr & 9 & 10 & 17\dg 14.1\pr & 9 & 11 & 32\dg 16.5\pr \\
9 & 12 & 47\dg 19.0\pr & 9 & 13 & 62\dg 21.5\pr & 9 & 14 & 77\dg 23.9\pr & 9 & 15 & 92\dg 26.4\pr \\
9 & 16 & 107\dg 28.8\pr & 9 & 17 & 122\dg 31.3\pr & 9 & 18 & 137\dg 33.8\pr & 9 & 19 & 152\dg 36.2\pr \\
9 & 20 & 167\dg 38.7\pr & 9 & 21 & 182\dg 41.2\pr & 9 & 22 & 197\dg 43.6\pr & 9 & 23 & 212\dg 46.1\pr \\
10 & 0 & 227\dg 48.6\pr & 10 & 1 & 242\dg 51.0\pr & 10 & 2 & 257\dg 53.5\pr & 10 & 3 & 272\dg 56.0\pr \\
10 & 4 & 287\dg 58.4\pr & 10 & 5 & 303\dg 00.9\pr & 10 & 6 & 318\dg 03.3\pr & 10 & 7 & 333\dg 05.8\pr \\
10 & 8 & 348\dg 08.3\pr & 10 & 9 & 3\dg 10.7\pr & 10 & 10 & 18\dg 13.2\pr & 10 & 11 & 33\dg 15.7\pr \\
10 & 12 & 48\dg 18.1\pr & 10 & 13 & 63\dg 20.6\pr & 10 & 14 & 78\dg 23.1\pr & 10 & 15 & 93\dg 25.5\pr \\
10 & 16 & 108\dg 28.0\pr & 10 & 17 & 123\dg 30.4\pr & 10 & 18 & 138\dg 32.9\pr & 10 & 19 & 153\dg 35.4\pr \\
10 & 20 & 168\dg 37.8\pr & 10 & 21 & 183\dg 40.3\pr & 10 & 22 & 198\dg 42.8\pr & 10 & 23 & 213\dg 45.2\pr \\
11 & 0 & 228\dg 47.7\pr & 11 & 1 & 243\dg 50.2\pr & 11 & 2 & 258\dg 52.6\pr & 11 & 3 & 273\dg 55.1\pr \\
11 & 4 & 288\dg 57.6\pr & 11 & 5 & 304\dg 00.0\pr & 11 & 6 & 319\dg 02.5\pr & 11 & 7 & 334\dg 04.9\pr \\
11 & 8 & 349\dg 07.4\pr & 11 & 9 & 4\dg 09.9\pr & 11 & 10 & 19\dg 12.3\pr & 11 & 11 & 34\dg 14.8\pr \\
11 & 12 & 49\dg 17.3\pr & 11 & 13 & 64\dg 19.7\pr & 11 & 14 & 79\dg 22.2\pr & 11 & 15 & 94\dg 24.7\pr \\
11 & 16 & 109\dg 27.1\pr & 11 & 17 & 124\dg 29.6\pr & 11 & 18 & 139\dg 32.1\pr & 11 & 19 & 154\dg 34.5\pr \\
11 & 20 & 169\dg 37.0\pr & 11 & 21 & 184\dg 39.4\pr & 11 & 22 & 199\dg 41.9\pr & 11 & 23 & 214\dg 44.4\pr \\
12 & 0 & 229\dg 46.8\pr & 12 & 1 & 244\dg 49.3\pr & 12 & 2 & 259\dg 51.8\pr & 12 & 3 & 274\dg 54.2\pr \\
12 & 4 & 289\dg 56.7\pr & 12 & 5 & 304\dg 59.2\pr & 12 & 6 & 320\dg 01.6\pr & 12 & 7 & 335\dg 04.1\pr \\
12 & 8 & 350\dg 06.6\pr & 12 & 9 & 5\dg 09.0\pr & 12 & 10 & 20\dg 11.5\pr & 12 & 11 & 35\dg 13.9\pr \\
12 & 12 & 50\dg 16.4\pr & 12 & 13 & 65\dg 18.9\pr & 12 & 14 & 80\dg 21.3\pr & 12 & 15 & 95\dg 23.8\pr \\
12 & 16 & 110\dg 26.3\pr & 12 & 17 & 125\dg 28.7\pr & 12 & 18 & 140\dg 31.2\pr & 12 & 19 & 155\dg 33.7\pr \\
12 & 20 & 170\dg 36.1\pr & 12 & 21 & 185\dg 38.6\pr & 12 & 22 & 200\dg 41.0\pr & 12 & 23 & 215\dg 43.5\pr \\
13 & 0 & 230\dg 46.0\pr & 13 & 1 & 245\dg 48.4\pr & 13 & 2 & 260\dg 50.9\pr & 13 & 3 & 275\dg 53.4\pr \\
13 & 4 & 290\dg 55.8\pr & 13 & 5 & 305\dg 58.3\pr & 13 & 6 & 321\dg 00.8\pr & 13 & 7 & 336\dg 03.2\pr \\
13 & 8 & 351\dg 05.7\pr & 13 & 9 & 6\dg 08.2\pr & 13 & 10 & 21\dg 10.6\pr & 13 & 11 & 36\dg 13.1\pr \\
13 & 12 & 51\dg 15.5\pr & 13 & 13 & 66\dg 18.0\pr & 13 & 14 & 81\dg 20.5\pr & 13 & 15 & 96\dg 22.9\pr \\
13 & 16 & 111\dg 25.4\pr & 13 & 17 & 126\dg 27.9\pr & 13 & 18 & 141\dg 30.3\pr & 13 & 19 & 156\dg 32.8\pr \\
13 & 20 & 171\dg 35.3\pr & 13 & 21 & 186\dg 37.7\pr & 13 & 22 & 201\dg 40.2\pr & 13 & 23 & 216\dg 42.7\pr \\
14 & 0 & 231\dg 45.1\pr & 14 & 1 & 246\dg 47.6\pr & 14 & 2 & 261\dg 50.0\pr & 14 & 3 & 276\dg 52.5\pr \\
14 & 4 & 291\dg 55.0\pr & 14 & 5 & 306\dg 57.4\pr & 14 & 6 & 321\dg 59.9\pr & 14 & 7 & 337\dg 02.4\pr \\
14 & 8 & 352\dg 04.8\pr & 14 & 9 & 7\dg 07.3\pr & 14 & 10 & 22\dg 09.8\pr & 14 & 11 & 37\dg 12.2\pr \\
14 & 12 & 52\dg 14.7\pr & 14 & 13 & 67\dg 17.1\pr & 14 & 14 & 82\dg 19.6\pr & 14 & 15 & 97\dg 22.1\pr \\
14 & 16 & 112\dg 24.5\pr & 14 & 17 & 127\dg 27.0\pr & 14 & 18 & 142\dg 29.5\pr & 14 & 19 & 157\dg 31.9\pr \\
14 & 20 & 172\dg 34.4\pr & 14 & 21 & 187\dg 36.9\pr & 14 & 22 & 202\dg 39.3\pr & 14 & 23 & 217\dg 41.8\pr \\
15 & 0 & 232\dg 44.3\pr & 15 & 1 & 247\dg 46.7\pr & 15 & 2 & 262\dg 49.2\pr & 15 & 3 & 277\dg 51.6\pr \\
15 & 4 & 292\dg 54.1\pr & 15 & 5 & 307\dg 56.6\pr & 15 & 6 & 322\dg 59.0\pr & 15 & 7 & 338\dg 01.5\pr \\
15 & 8 & 353\dg 04.0\pr & 15 & 9 & 8\dg 06.4\pr & 15 & 10 & 23\dg 08.9\pr & 15 & 11 & 38\dg 11.4\pr \\
15 & 12 & 53\dg 13.8\pr & 15 & 13 & 68\dg 16.3\pr & 15 & 14 & 83\dg 18.8\pr & 15 & 15 & 98\dg 21.2\pr \\
15 & 16 & 113\dg 23.7\pr & 15 & 17 & 128\dg 26.1\pr & 15 & 18 & 143\dg 28.6\pr & 15 & 19 & 158\dg 31.1\pr \\
15 & 20 & 173\dg 33.5\pr & 15 & 21 & 188\dg 36.0\pr & 15 & 22 & 203\dg 38.5\pr & 15 & 23 & 218\dg 40.9\pr \\
16 & 0 & 233\dg 43.4\pr & 16 & 1 & 248\dg 45.9\pr & 16 & 2 & 263\dg 48.3\pr & 16 & 3 & 278\dg 50.8\pr \\
16 & 4 & 293\dg 53.2\pr & 16 & 5 & 308\dg 55.7\pr & 16 & 6 & 323\dg 58.2\pr & 16 & 7 & 339\dg 00.6\pr \\
16 & 8 & 354\dg 03.1\pr & 16 & 9 & 9\dg 05.6\pr & 16 & 10 & 24\dg 08.0\pr & 16 & 11 & 39\dg 10.5\pr \\
16 & 12 & 54\dg 13.0\pr & 16 & 13 & 69\dg 15.4\pr & 16 & 14 & 84\dg 17.9\pr & 16 & 15 & 99\dg 20.4\pr \\
16 & 16 & 114\dg 22.8\pr & 16 & 17 & 129\dg 25.3\pr & 16 & 18 & 144\dg 27.7\pr & 16 & 19 & 159\dg 30.2\pr \\
16 & 20 & 174\dg 32.7\pr & 16 & 21 & 189\dg 35.1\pr & 16 & 22 & 204\dg 37.6\pr & 16 & 23 & 219\dg 40.1\pr \\
17 & 0 & 234\dg 42.5\pr & 17 & 1 & 249\dg 45.0\pr & 17 & 2 & 264\dg 47.5\pr & 17 & 3 & 279\dg 49.9\pr \\
17 & 4 & 294\dg 52.4\pr & 17 & 5 & 309\dg 54.9\pr & 17 & 6 & 324\dg 57.3\pr & 17 & 7 & 339\dg 59.8\pr \\
17 & 8 & 355\dg 02.2\pr & 17 & 9 & 10\dg 04.7\pr & 17 & 10 & 25\dg 07.2\pr & 17 & 11 & 40\dg 09.6\pr \\
17 & 12 & 55\dg 12.1\pr & 17 & 13 & 70\dg 14.6\pr & 17 & 14 & 85\dg 17.0\pr & 17 & 15 & 100\dg 19.5\pr \\
17 & 16 & 115\dg 22.0\pr & 17 & 17 & 130\dg 24.4\pr & 17 & 18 & 145\dg 26.9\pr & 17 & 19 & 160\dg 29.3\pr \\
17 & 20 & 175\dg 31.8\pr & 17 & 21 & 190\dg 34.3\pr & 17 & 22 & 205\dg 36.7\pr & 17 & 23 & 220\dg 39.2\pr \\
18 & 0 & 235\dg 41.7\pr & 18 & 1 & 250\dg 44.1\pr & 18 & 2 & 265\dg 46.6\pr & 18 & 3 & 280\dg 49.1\pr \\
18 & 4 & 295\dg 51.5\pr & 18 & 5 & 310\dg 54.0\pr & 18 & 6 & 325\dg 56.5\pr & 18 & 7 & 340\dg 58.9\pr \\
18 & 8 & 356\dg 01.4\pr & 18 & 9 & 11\dg 03.8\pr & 18 & 10 & 26\dg 06.3\pr & 18 & 11 & 41\dg 08.8\pr \\
18 & 12 & 56\dg 11.2\pr & 18 & 13 & 71\dg 13.7\pr & 18 & 14 & 86\dg 16.2\pr & 18 & 15 & 101\dg 18.6\pr \\
18 & 16 & 116\dg 21.1\pr & 18 & 17 & 131\dg 23.6\pr & 18 & 18 & 146\dg 26.0\pr & 18 & 19 & 161\dg 28.5\pr \\
18 & 20 & 176\dg 31.0\pr & 18 & 21 & 191\dg 33.4\pr & 18 & 22 & 206\dg 35.9\pr & 18 & 23 & 221\dg 38.3\pr \\
19 & 0 & 236\dg 40.8\pr & 19 & 1 & 251\dg 43.3\pr & 19 & 2 & 266\dg 45.7\pr & 19 & 3 & 281\dg 48.2\pr \\
19 & 4 & 296\dg 50.7\pr & 19 & 5 & 311\dg 53.1\pr & 19 & 6 & 326\dg 55.6\pr & 19 & 7 & 341\dg 58.1\pr \\
19 & 8 & 357\dg 00.5\pr & 19 & 9 & 12\dg 03.0\pr & 19 & 10 & 27\dg 05.5\pr & 19 & 11 & 42\dg 07.9\pr \\
19 & 12 & 57\dg 10.4\pr & 19 & 13 & 72\dg 12.8\pr & 19 & 14 & 87\dg 15.3\pr & 19 & 15 & 102\dg 17.8\pr \\
19 & 16 & 117\dg 20.2\pr & 19 & 17 & 132\dg 22.7\pr & 19 & 18 & 147\dg 25.2\pr & 19 & 19 & 162\dg 27.6\pr \\
19 & 20 & 177\dg 30.1\pr & 19 & 21 & 192\dg 32.6\pr & 19 & 22 & 207\dg 35.0\pr & 19 & 23 & 222\dg 37.5\pr \\
20 & 0 & 237\dg 39.9\pr & 20 & 1 & 252\dg 42.4\pr & 20 & 2 & 267\dg 44.9\pr & 20 & 3 & 282\dg 47.3\pr \\
20 & 4 & 297\dg 49.8\pr & 20 & 5 & 312\dg 52.3\pr & 20 & 6 & 327\dg 54.7\pr & 20 & 7 & 342\dg 57.2\pr \\
20 & 8 & 357\dg 59.7\pr & 20 & 9 & 13\dg 02.1\pr & 20 & 10 & 28\dg 04.6\pr & 20 & 11 & 43\dg 07.1\pr \\
20 & 12 & 58\dg 09.5\pr & 20 & 13 & 73\dg 12.0\pr & 20 & 14 & 88\dg 14.4\pr & 20 & 15 & 103\dg 16.9\pr \\
20 & 16 & 118\dg 19.4\pr & 20 & 17 & 133\dg 21.8\pr & 20 & 18 & 148\dg 24.3\pr & 20 & 19 & 163\dg 26.8\pr \\
20 & 20 & 178\dg 29.2\pr & 20 & 21 & 193\dg 31.7\pr & 20 & 22 & 208\dg 34.2\pr & 20 & 23 & 223\dg 36.6\pr \\
21 & 0 & 238\dg 39.1\pr & 21 & 1 & 253\dg 41.6\pr & 21 & 2 & 268\dg 44.0\pr & 21 & 3 & 283\dg 46.5\pr \\
21 & 4 & 298\dg 48.9\pr & 21 & 5 & 313\dg 51.4\pr & 21 & 6 & 328\dg 53.9\pr & 21 & 7 & 343\dg 56.3\pr \\
21 & 8 & 358\dg 58.8\pr & 21 & 9 & 14\dg 01.3\pr & 21 & 10 & 29\dg 03.7\pr & 21 & 11 & 44\dg 06.2\pr \\
21 & 12 & 59\dg 08.7\pr & 21 & 13 & 74\dg 11.1\pr & 21 & 14 & 89\dg 13.6\pr & 21 & 15 & 104\dg 16.0\pr \\
21 & 16 & 119\dg 18.5\pr & 21 & 17 & 134\dg 21.0\pr & 21 & 18 & 149\dg 23.4\pr & 21 & 19 & 164\dg 25.9\pr \\
21 & 20 & 179\dg 28.4\pr & 21 & 21 & 194\dg 30.8\pr & 21 & 22 & 209\dg 33.3\pr & 21 & 23 & 224\dg 35.8\pr \\
22 & 0 & 239\dg 38.2\pr & 22 & 1 & 254\dg 40.7\pr & 22 & 2 & 269\dg 43.2\pr & 22 & 3 & 284\dg 45.6\pr \\
22 & 4 & 299\dg 48.1\pr & 22 & 5 & 314\dg 50.5\pr & 22 & 6 & 329\dg 53.0\pr & 22 & 7 & 344\dg 55.5\pr \\
22 & 8 & 359\dg 57.9\pr & 22 & 9 & 15\dg 00.4\pr & 22 & 10 & 30\dg 02.9\pr & 22 & 11 & 45\dg 05.3\pr \\
22 & 12 & 60\dg 07.8\pr & 22 & 13 & 75\dg 10.3\pr & 22 & 14 & 90\dg 12.7\pr & 22 & 15 & 105\dg 15.2\pr \\
22 & 16 & 120\dg 17.7\pr & 22 & 17 & 135\dg 20.1\pr & 22 & 18 & 150\dg 22.6\pr & 22 & 19 & 165\dg 25.0\pr \\
22 & 20 & 180\dg 27.5\pr & 22 & 21 & 195\dg 30.0\pr & 22 & 22 & 210\dg 32.4\pr & 22 & 23 & 225\dg 34.9\pr \\
23 & 0 & 240\dg 37.4\pr & 23 & 1 & 255\dg 39.8\pr & 23 & 2 & 270\dg 42.3\pr & 23 & 3 & 285\dg 44.8\pr \\
23 & 4 & 300\dg 47.2\pr & 23 & 5 & 315\dg 49.7\pr & 23 & 6 & 330\dg 52.1\pr & 23 & 7 & 345\dg 54.6\pr \\
23 & 8 & 0\dg 57.1\pr & 23 & 9 & 15\dg 59.5\pr & 23 & 10 & 31\dg 02.0\pr & 23 & 11 & 46\dg 04.5\pr \\
23 & 12 & 61\dg 06.9\pr & 23 & 13 & 76\dg 09.4\pr & 23 & 14 & 91\dg 11.9\pr & 23 & 15 & 106\dg 14.3\pr \\
23 & 16 & 121\dg 16.8\pr & 23 & 17 & 136\dg 19.3\pr & 23 & 18 & 151\dg 21.7\pr & 23 & 19 & 166\dg 24.2\pr \\
23 & 20 & 181\dg 26.6\pr & 23 & 21 & 196\dg 29.1\pr & 23 & 22 & 211\dg 31.6\pr & 23 & 23 & 226\dg 34.0\pr \\
24 & 0 & 241\dg 36.5\pr & 24 & 1 & 256\dg 39.0\pr & 24 & 2 & 271\dg 41.4\pr & 24 & 3 & 286\dg 43.9\pr \\
24 & 4 & 301\dg 46.4\pr & 24 & 5 & 316\dg 48.8\pr & 24 & 6 & 331\dg 51.3\pr & 24 & 7 & 346\dg 53.8\pr \\
24 & 8 & 1\dg 56.2\pr & 24 & 9 & 16\dg 58.7\pr & 24 & 10 & 32\dg 01.1\pr & 24 & 11 & 47\dg 03.6\pr \\
24 & 12 & 62\dg 06.1\pr & 24 & 13 & 77\dg 08.5\pr & 24 & 14 & 92\dg 11.0\pr & 24 & 15 & 107\dg 13.5\pr \\
24 & 16 & 122\dg 15.9\pr & 24 & 17 & 137\dg 18.4\pr & 24 & 18 & 152\dg 20.9\pr & 24 & 19 & 167\dg 23.3\pr \\
24 & 20 & 182\dg 25.8\pr & 24 & 21 & 197\dg 28.3\pr & 24 & 22 & 212\dg 30.7\pr & 24 & 23 & 227\dg 33.2\pr \\
25 & 0 & 242\dg 35.6\pr & 25 & 1 & 257\dg 38.1\pr & 25 & 2 & 272\dg 40.6\pr & 25 & 3 & 287\dg 43.0\pr \\
25 & 4 & 302\dg 45.5\pr & 25 & 5 & 317\dg 48.0\pr & 25 & 6 & 332\dg 50.4\pr & 25 & 7 & 347\dg 52.9\pr \\
25 & 8 & 2\dg 55.4\pr & 25 & 9 & 17\dg 57.8\pr & 25 & 10 & 33\dg 00.3\pr & 25 & 11 & 48\dg 02.7\pr \\
25 & 12 & 63\dg 05.2\pr & 25 & 13 & 78\dg 07.7\pr & 25 & 14 & 93\dg 10.1\pr & 25 & 15 & 108\dg 12.6\pr \\
25 & 16 & 123\dg 15.1\pr & 25 & 17 & 138\dg 17.5\pr & 25 & 18 & 153\dg 20.0\pr & 25 & 19 & 168\dg 22.5\pr \\
25 & 20 & 183\dg 24.9\pr & 25 & 21 & 198\dg 27.4\pr & 25 & 22 & 213\dg 29.9\pr & 25 & 23 & 228\dg 32.3\pr \\
26 & 0 & 243\dg 34.8\pr & 26 & 1 & 258\dg 37.2\pr & 26 & 2 & 273\dg 39.7\pr & 26 & 3 & 288\dg 42.2\pr \\
26 & 4 & 303\dg 44.6\pr & 26 & 5 & 318\dg 47.1\pr & 26 & 6 & 333\dg 49.6\pr & 26 & 7 & 348\dg 52.0\pr \\
26 & 8 & 3\dg 54.5\pr & 26 & 9 & 18\dg 57.0\pr & 26 & 10 & 33\dg 59.4\pr & 26 & 11 & 49\dg 01.9\pr \\
26 & 12 & 64\dg 04.4\pr & 26 & 13 & 79\dg 06.8\pr & 26 & 14 & 94\dg 09.3\pr & 26 & 15 & 109\dg 11.7\pr \\
26 & 16 & 124\dg 14.2\pr & 26 & 17 & 139\dg 16.7\pr & 26 & 18 & 154\dg 19.1\pr & 26 & 19 & 169\dg 21.6\pr \\
26 & 20 & 184\dg 24.1\pr & 26 & 21 & 199\dg 26.5\pr & 26 & 22 & 214\dg 29.0\pr & 26 & 23 & 229\dg 31.5\pr \\
27 & 0 & 244\dg 33.9\pr & 27 & 1 & 259\dg 36.4\pr & 27 & 2 & 274\dg 38.8\pr & 27 & 3 & 289\dg 41.3\pr \\
27 & 4 & 304\dg 43.8\pr & 27 & 5 & 319\dg 46.2\pr & 27 & 6 & 334\dg 48.7\pr & 27 & 7 & 349\dg 51.2\pr \\
27 & 8 & 4\dg 53.6\pr & 27 & 9 & 19\dg 56.1\pr & 27 & 10 & 34\dg 58.6\pr & 27 & 11 & 50\dg 01.0\pr \\
27 & 12 & 65\dg 03.5\pr & 27 & 13 & 80\dg 06.0\pr & 27 & 14 & 95\dg 08.4\pr & 27 & 15 & 110\dg 10.9\pr \\
27 & 16 & 125\dg 13.3\pr & 27 & 17 & 140\dg 15.8\pr & 27 & 18 & 155\dg 18.3\pr & 27 & 19 & 170\dg 20.7\pr \\
27 & 20 & 185\dg 23.2\pr & 27 & 21 & 200\dg 25.7\pr & 27 & 22 & 215\dg 28.1\pr & 27 & 23 & 230\dg 30.6\pr \\
28 & 0 & 245\dg 33.1\pr & 28 & 1 & 260\dg 35.5\pr & 28 & 2 & 275\dg 38.0\pr & 28 & 3 & 290\dg 40.5\pr \\
28 & 4 & 305\dg 42.9\pr & 28 & 5 & 320\dg 45.4\pr & 28 & 6 & 335\dg 47.8\pr & 28 & 7 & 350\dg 50.3\pr \\
28 & 8 & 5\dg 52.8\pr & 28 & 9 & 20\dg 55.2\pr & 28 & 10 & 35\dg 57.7\pr & 28 & 11 & 51\dg 00.2\pr \\
28 & 12 & 66\dg 02.6\pr & 28 & 13 & 81\dg 05.1\pr & 28 & 14 & 96\dg 07.6\pr & 28 & 15 & 111\dg 10.0\pr \\
28 & 16 & 126\dg 12.5\pr & 28 & 17 & 141\dg 14.9\pr & 28 & 18 & 156\dg 17.4\pr & 28 & 19 & 171\dg 19.9\pr \\
28 & 20 & 186\dg 22.3\pr & 28 & 21 & 201\dg 24.8\pr & 28 & 22 & 216\dg 27.3\pr & 28 & 23 & 231\dg 29.7\pr \\
29 & 0 & 246\dg 32.2\pr & 29 & 1 & 261\dg 34.7\pr & 29 & 2 & 276\dg 37.1\pr & 29 & 3 & 291\dg 39.6\pr \\
29 & 4 & 306\dg 42.1\pr & 29 & 5 & 321\dg 44.5\pr & 29 & 6 & 336\dg 47.0\pr & 29 & 7 & 351\dg 49.4\pr \\
29 & 8 & 6\dg 51.9\pr & 29 & 9 & 21\dg 54.4\pr & 29 & 10 & 36\dg 56.8\pr & 29 & 11 & 51\dg 59.3\pr \\
29 & 12 & 67\dg 01.8\pr & 29 & 13 & 82\dg 04.2\pr & 29 & 14 & 97\dg 06.7\pr & 29 & 15 & 112\dg 09.2\pr \\
29 & 16 & 127\dg 11.6\pr & 29 & 17 & 142\dg 14.1\pr & 29 & 18 & 157\dg 16.6\pr & 29 & 19 & 172\dg 19.0\pr \\
29 & 20 & 187\dg 21.5\pr & 29 & 21 & 202\dg 23.9\pr & 29 & 22 & 217\dg 26.4\pr & 29 & 23 & 232\dg 28.9\pr \\
30 & 0 & 247\dg 31.3\pr & 30 & 1 & 262\dg 33.8\pr & 30 & 2 & 277\dg 36.3\pr & 30 & 3 & 292\dg 38.7\pr \\
30 & 4 & 307\dg 41.2\pr & 30 & 5 & 322\dg 43.7\pr & 30 & 6 & 337\dg 46.1\pr & 30 & 7 & 352\dg 48.6\pr \\
30 & 8 & 7\dg 51.0\pr & 30 & 9 & 22\dg 53.5\pr & 30 & 10 & 37\dg 56.0\pr & 30 & 11 & 52\dg 58.4\pr \\
30 & 12 & 68\dg 00.9\pr & 30 & 13 & 83\dg 03.4\pr & 30 & 14 & 98\dg 05.8\pr & 30 & 15 & 113\dg 08.3\pr \\
30 & 16 & 128\dg 10.8\pr & 30 & 17 & 143\dg 13.2\pr & 30 & 18 & 158\dg 15.7\pr & 30 & 19 & 173\dg 18.2\pr \\
30 & 20 & 188\dg 20.6\pr & 30 & 21 & 203\dg 23.1\pr & 30 & 22 & 218\dg 25.5\pr & 30 & 23 & 233\dg 28.0\pr \\
31 & 0 & 248\dg 30.5\pr & 31 & 1 & 263\dg 32.9\pr & 31 & 2 & 278\dg 35.4\pr & 31 & 3 & 293\dg 37.9\pr \\
31 & 4 & 308\dg 40.3\pr & 31 & 5 & 323\dg 42.8\pr & 31 & 6 & 338\dg 45.3\pr & 31 & 7 & 353\dg 47.7\pr \\
31 & 8 & 8\dg 50.2\pr & 31 & 9 & 23\dg 52.7\pr & 31 & 10 & 38\dg 55.1\pr & 31 & 11 & 53\dg 57.6\pr \\
31 & 12 & 69\dg 00.0\pr & 31 & 13 & 84\dg 02.5\pr & 31 & 14 & 99\dg 05.0\pr & 31 & 15 & 114\dg 07.4\pr \\
31 & 16 & 129\dg 09.9\pr & 31 & 17 & 144\dg 12.4\pr & 31 & 18 & 159\dg 14.8\pr & 31 & 19 & 174\dg 17.3\pr \\
31 & 20 & 189\dg 19.8\pr & 31 & 21 & 204\dg 22.2\pr & 31 & 22 & 219\dg 24.7\pr & 31 & 23 & 234\dg 27.2\pr \\
\end{longtable}}

{\footnotesize
\begin{longtable}{rrr|rrr|rrr|rrr}
\caption{Aries --- GHA of the First Point of Aries, June 2026 (UT)}\label{aries06}\\
\toprule
\textbf{d} & \textbf{h} & \textbf{GHA} $\gamma$ & \textbf{d} & \textbf{h} & \textbf{GHA} $\gamma$ & \textbf{d} & \textbf{h} & \textbf{GHA} $\gamma$ & \textbf{d} & \textbf{h} & \textbf{GHA} $\gamma$ \\
\midrule
\endfirsthead
\multicolumn{12}{c}{\textit{Aries --- GHA of the First Point of Aries, June 2026 (UT) (continued)}}\\
\toprule
\textbf{d} & \textbf{h} & \textbf{GHA} $\gamma$ & \textbf{d} & \textbf{h} & \textbf{GHA} $\gamma$ & \textbf{d} & \textbf{h} & \textbf{GHA} $\gamma$ & \textbf{d} & \textbf{h} & \textbf{GHA} $\gamma$ \\
\midrule
\endhead
\midrule\multicolumn{12}{r}{\textit{continued on next page}}\\
\endfoot
\bottomrule
\endlastfoot
1 & 0 & 249\dg 29.6\pr & 1 & 1 & 264\dg 32.1\pr & 1 & 2 & 279\dg 34.5\pr & 1 & 3 & 294\dg 37.0\pr \\
1 & 4 & 309\dg 39.5\pr & 1 & 5 & 324\dg 41.9\pr & 1 & 6 & 339\dg 44.4\pr & 1 & 7 & 354\dg 46.9\pr \\
1 & 8 & 9\dg 49.3\pr & 1 & 9 & 24\dg 51.8\pr & 1 & 10 & 39\dg 54.3\pr & 1 & 11 & 54\dg 56.7\pr \\
1 & 12 & 69\dg 59.2\pr & 1 & 13 & 85\dg 01.6\pr & 1 & 14 & 100\dg 04.1\pr & 1 & 15 & 115\dg 06.6\pr \\
1 & 16 & 130\dg 09.0\pr & 1 & 17 & 145\dg 11.5\pr & 1 & 18 & 160\dg 14.0\pr & 1 & 19 & 175\dg 16.4\pr \\
1 & 20 & 190\dg 18.9\pr & 1 & 21 & 205\dg 21.4\pr & 1 & 22 & 220\dg 23.8\pr & 1 & 23 & 235\dg 26.3\pr \\
2 & 0 & 250\dg 28.8\pr & 2 & 1 & 265\dg 31.2\pr & 2 & 2 & 280\dg 33.7\pr & 2 & 3 & 295\dg 36.1\pr \\
2 & 4 & 310\dg 38.6\pr & 2 & 5 & 325\dg 41.1\pr & 2 & 6 & 340\dg 43.5\pr & 2 & 7 & 355\dg 46.0\pr \\
2 & 8 & 10\dg 48.5\pr & 2 & 9 & 25\dg 50.9\pr & 2 & 10 & 40\dg 53.4\pr & 2 & 11 & 55\dg 55.9\pr \\
2 & 12 & 70\dg 58.3\pr & 2 & 13 & 86\dg 00.8\pr & 2 & 14 & 101\dg 03.3\pr & 2 & 15 & 116\dg 05.7\pr \\
2 & 16 & 131\dg 08.2\pr & 2 & 17 & 146\dg 10.6\pr & 2 & 18 & 161\dg 13.1\pr & 2 & 19 & 176\dg 15.6\pr \\
2 & 20 & 191\dg 18.0\pr & 2 & 21 & 206\dg 20.5\pr & 2 & 22 & 221\dg 23.0\pr & 2 & 23 & 236\dg 25.4\pr \\
3 & 0 & 251\dg 27.9\pr & 3 & 1 & 266\dg 30.4\pr & 3 & 2 & 281\dg 32.8\pr & 3 & 3 & 296\dg 35.3\pr \\
3 & 4 & 311\dg 37.7\pr & 3 & 5 & 326\dg 40.2\pr & 3 & 6 & 341\dg 42.7\pr & 3 & 7 & 356\dg 45.1\pr \\
3 & 8 & 11\dg 47.6\pr & 3 & 9 & 26\dg 50.1\pr & 3 & 10 & 41\dg 52.5\pr & 3 & 11 & 56\dg 55.0\pr \\
3 & 12 & 71\dg 57.5\pr & 3 & 13 & 86\dg 59.9\pr & 3 & 14 & 102\dg 02.4\pr & 3 & 15 & 117\dg 04.9\pr \\
3 & 16 & 132\dg 07.3\pr & 3 & 17 & 147\dg 09.8\pr & 3 & 18 & 162\dg 12.2\pr & 3 & 19 & 177\dg 14.7\pr \\
3 & 20 & 192\dg 17.2\pr & 3 & 21 & 207\dg 19.6\pr & 3 & 22 & 222\dg 22.1\pr & 3 & 23 & 237\dg 24.6\pr \\
4 & 0 & 252\dg 27.0\pr & 4 & 1 & 267\dg 29.5\pr & 4 & 2 & 282\dg 32.0\pr & 4 & 3 & 297\dg 34.4\pr \\
4 & 4 & 312\dg 36.9\pr & 4 & 5 & 327\dg 39.4\pr & 4 & 6 & 342\dg 41.8\pr & 4 & 7 & 357\dg 44.3\pr \\
4 & 8 & 12\dg 46.7\pr & 4 & 9 & 27\dg 49.2\pr & 4 & 10 & 42\dg 51.7\pr & 4 & 11 & 57\dg 54.1\pr \\
4 & 12 & 72\dg 56.6\pr & 4 & 13 & 87\dg 59.1\pr & 4 & 14 & 103\dg 01.5\pr & 4 & 15 & 118\dg 04.0\pr \\
4 & 16 & 133\dg 06.5\pr & 4 & 17 & 148\dg 08.9\pr & 4 & 18 & 163\dg 11.4\pr & 4 & 19 & 178\dg 13.8\pr \\
4 & 20 & 193\dg 16.3\pr & 4 & 21 & 208\dg 18.8\pr & 4 & 22 & 223\dg 21.2\pr & 4 & 23 & 238\dg 23.7\pr \\
5 & 0 & 253\dg 26.2\pr & 5 & 1 & 268\dg 28.6\pr & 5 & 2 & 283\dg 31.1\pr & 5 & 3 & 298\dg 33.6\pr \\
5 & 4 & 313\dg 36.0\pr & 5 & 5 & 328\dg 38.5\pr & 5 & 6 & 343\dg 41.0\pr & 5 & 7 & 358\dg 43.4\pr \\
5 & 8 & 13\dg 45.9\pr & 5 & 9 & 28\dg 48.3\pr & 5 & 10 & 43\dg 50.8\pr & 5 & 11 & 58\dg 53.3\pr \\
5 & 12 & 73\dg 55.7\pr & 5 & 13 & 88\dg 58.2\pr & 5 & 14 & 104\dg 00.7\pr & 5 & 15 & 119\dg 03.1\pr \\
5 & 16 & 134\dg 05.6\pr & 5 & 17 & 149\dg 08.1\pr & 5 & 18 & 164\dg 10.5\pr & 5 & 19 & 179\dg 13.0\pr \\
5 & 20 & 194\dg 15.5\pr & 5 & 21 & 209\dg 17.9\pr & 5 & 22 & 224\dg 20.4\pr & 5 & 23 & 239\dg 22.8\pr \\
6 & 0 & 254\dg 25.3\pr & 6 & 1 & 269\dg 27.8\pr & 6 & 2 & 284\dg 30.2\pr & 6 & 3 & 299\dg 32.7\pr \\
6 & 4 & 314\dg 35.2\pr & 6 & 5 & 329\dg 37.6\pr & 6 & 6 & 344\dg 40.1\pr & 6 & 7 & 359\dg 42.6\pr \\
6 & 8 & 14\dg 45.0\pr & 6 & 9 & 29\dg 47.5\pr & 6 & 10 & 44\dg 49.9\pr & 6 & 11 & 59\dg 52.4\pr \\
6 & 12 & 74\dg 54.9\pr & 6 & 13 & 89\dg 57.3\pr & 6 & 14 & 104\dg 59.8\pr & 6 & 15 & 120\dg 02.3\pr \\
6 & 16 & 135\dg 04.7\pr & 6 & 17 & 150\dg 07.2\pr & 6 & 18 & 165\dg 09.7\pr & 6 & 19 & 180\dg 12.1\pr \\
6 & 20 & 195\dg 14.6\pr & 6 & 21 & 210\dg 17.1\pr & 6 & 22 & 225\dg 19.5\pr & 6 & 23 & 240\dg 22.0\pr \\
7 & 0 & 255\dg 24.4\pr & 7 & 1 & 270\dg 26.9\pr & 7 & 2 & 285\dg 29.4\pr & 7 & 3 & 300\dg 31.8\pr \\
7 & 4 & 315\dg 34.3\pr & 7 & 5 & 330\dg 36.8\pr & 7 & 6 & 345\dg 39.2\pr & 7 & 7 & 0\dg 41.7\pr \\
7 & 8 & 15\dg 44.2\pr & 7 & 9 & 30\dg 46.6\pr & 7 & 10 & 45\dg 49.1\pr & 7 & 11 & 60\dg 51.6\pr \\
7 & 12 & 75\dg 54.0\pr & 7 & 13 & 90\dg 56.5\pr & 7 & 14 & 105\dg 58.9\pr & 7 & 15 & 121\dg 01.4\pr \\
7 & 16 & 136\dg 03.9\pr & 7 & 17 & 151\dg 06.3\pr & 7 & 18 & 166\dg 08.8\pr & 7 & 19 & 181\dg 11.3\pr \\
7 & 20 & 196\dg 13.7\pr & 7 & 21 & 211\dg 16.2\pr & 7 & 22 & 226\dg 18.7\pr & 7 & 23 & 241\dg 21.1\pr \\
8 & 0 & 256\dg 23.6\pr & 8 & 1 & 271\dg 26.1\pr & 8 & 2 & 286\dg 28.5\pr & 8 & 3 & 301\dg 31.0\pr \\
8 & 4 & 316\dg 33.4\pr & 8 & 5 & 331\dg 35.9\pr & 8 & 6 & 346\dg 38.4\pr & 8 & 7 & 1\dg 40.8\pr \\
8 & 8 & 16\dg 43.3\pr & 8 & 9 & 31\dg 45.8\pr & 8 & 10 & 46\dg 48.2\pr & 8 & 11 & 61\dg 50.7\pr \\
8 & 12 & 76\dg 53.2\pr & 8 & 13 & 91\dg 55.6\pr & 8 & 14 & 106\dg 58.1\pr & 8 & 15 & 122\dg 00.5\pr \\
8 & 16 & 137\dg 03.0\pr & 8 & 17 & 152\dg 05.5\pr & 8 & 18 & 167\dg 07.9\pr & 8 & 19 & 182\dg 10.4\pr \\
8 & 20 & 197\dg 12.9\pr & 8 & 21 & 212\dg 15.3\pr & 8 & 22 & 227\dg 17.8\pr & 8 & 23 & 242\dg 20.3\pr \\
9 & 0 & 257\dg 22.7\pr & 9 & 1 & 272\dg 25.2\pr & 9 & 2 & 287\dg 27.7\pr & 9 & 3 & 302\dg 30.1\pr \\
9 & 4 & 317\dg 32.6\pr & 9 & 5 & 332\dg 35.0\pr & 9 & 6 & 347\dg 37.5\pr & 9 & 7 & 2\dg 40.0\pr \\
9 & 8 & 17\dg 42.4\pr & 9 & 9 & 32\dg 44.9\pr & 9 & 10 & 47\dg 47.4\pr & 9 & 11 & 62\dg 49.8\pr \\
9 & 12 & 77\dg 52.3\pr & 9 & 13 & 92\dg 54.8\pr & 9 & 14 & 107\dg 57.2\pr & 9 & 15 & 122\dg 59.7\pr \\
9 & 16 & 138\dg 02.2\pr & 9 & 17 & 153\dg 04.6\pr & 9 & 18 & 168\dg 07.1\pr & 9 & 19 & 183\dg 09.5\pr \\
9 & 20 & 198\dg 12.0\pr & 9 & 21 & 213\dg 14.5\pr & 9 & 22 & 228\dg 16.9\pr & 9 & 23 & 243\dg 19.4\pr \\
10 & 0 & 258\dg 21.9\pr & 10 & 1 & 273\dg 24.3\pr & 10 & 2 & 288\dg 26.8\pr & 10 & 3 & 303\dg 29.3\pr \\
10 & 4 & 318\dg 31.7\pr & 10 & 5 & 333\dg 34.2\pr & 10 & 6 & 348\dg 36.6\pr & 10 & 7 & 3\dg 39.1\pr \\
10 & 8 & 18\dg 41.6\pr & 10 & 9 & 33\dg 44.0\pr & 10 & 10 & 48\dg 46.5\pr & 10 & 11 & 63\dg 49.0\pr \\
10 & 12 & 78\dg 51.4\pr & 10 & 13 & 93\dg 53.9\pr & 10 & 14 & 108\dg 56.4\pr & 10 & 15 & 123\dg 58.8\pr \\
10 & 16 & 139\dg 01.3\pr & 10 & 17 & 154\dg 03.8\pr & 10 & 18 & 169\dg 06.2\pr & 10 & 19 & 184\dg 08.7\pr \\
10 & 20 & 199\dg 11.1\pr & 10 & 21 & 214\dg 13.6\pr & 10 & 22 & 229\dg 16.1\pr & 10 & 23 & 244\dg 18.5\pr \\
11 & 0 & 259\dg 21.0\pr & 11 & 1 & 274\dg 23.5\pr & 11 & 2 & 289\dg 25.9\pr & 11 & 3 & 304\dg 28.4\pr \\
11 & 4 & 319\dg 30.9\pr & 11 & 5 & 334\dg 33.3\pr & 11 & 6 & 349\dg 35.8\pr & 11 & 7 & 4\dg 38.3\pr \\
11 & 8 & 19\dg 40.7\pr & 11 & 9 & 34\dg 43.2\pr & 11 & 10 & 49\dg 45.6\pr & 11 & 11 & 64\dg 48.1\pr \\
11 & 12 & 79\dg 50.6\pr & 11 & 13 & 94\dg 53.0\pr & 11 & 14 & 109\dg 55.5\pr & 11 & 15 & 124\dg 58.0\pr \\
11 & 16 & 140\dg 00.4\pr & 11 & 17 & 155\dg 02.9\pr & 11 & 18 & 170\dg 05.4\pr & 11 & 19 & 185\dg 07.8\pr \\
11 & 20 & 200\dg 10.3\pr & 11 & 21 & 215\dg 12.7\pr & 11 & 22 & 230\dg 15.2\pr & 11 & 23 & 245\dg 17.7\pr \\
12 & 0 & 260\dg 20.1\pr & 12 & 1 & 275\dg 22.6\pr & 12 & 2 & 290\dg 25.1\pr & 12 & 3 & 305\dg 27.5\pr \\
12 & 4 & 320\dg 30.0\pr & 12 & 5 & 335\dg 32.5\pr & 12 & 6 & 350\dg 34.9\pr & 12 & 7 & 5\dg 37.4\pr \\
12 & 8 & 20\dg 39.9\pr & 12 & 9 & 35\dg 42.3\pr & 12 & 10 & 50\dg 44.8\pr & 12 & 11 & 65\dg 47.2\pr \\
12 & 12 & 80\dg 49.7\pr & 12 & 13 & 95\dg 52.2\pr & 12 & 14 & 110\dg 54.6\pr & 12 & 15 & 125\dg 57.1\pr \\
12 & 16 & 140\dg 59.6\pr & 12 & 17 & 156\dg 02.0\pr & 12 & 18 & 171\dg 04.5\pr & 12 & 19 & 186\dg 07.0\pr \\
12 & 20 & 201\dg 09.4\pr & 12 & 21 & 216\dg 11.9\pr & 12 & 22 & 231\dg 14.4\pr & 12 & 23 & 246\dg 16.8\pr \\
13 & 0 & 261\dg 19.3\pr & 13 & 1 & 276\dg 21.7\pr & 13 & 2 & 291\dg 24.2\pr & 13 & 3 & 306\dg 26.7\pr \\
13 & 4 & 321\dg 29.1\pr & 13 & 5 & 336\dg 31.6\pr & 13 & 6 & 351\dg 34.1\pr & 13 & 7 & 6\dg 36.5\pr \\
13 & 8 & 21\dg 39.0\pr & 13 & 9 & 36\dg 41.5\pr & 13 & 10 & 51\dg 43.9\pr & 13 & 11 & 66\dg 46.4\pr \\
13 & 12 & 81\dg 48.8\pr & 13 & 13 & 96\dg 51.3\pr & 13 & 14 & 111\dg 53.8\pr & 13 & 15 & 126\dg 56.2\pr \\
13 & 16 & 141\dg 58.7\pr & 13 & 17 & 157\dg 01.2\pr & 13 & 18 & 172\dg 03.6\pr & 13 & 19 & 187\dg 06.1\pr \\
13 & 20 & 202\dg 08.6\pr & 13 & 21 & 217\dg 11.0\pr & 13 & 22 & 232\dg 13.5\pr & 13 & 23 & 247\dg 16.0\pr \\
14 & 0 & 262\dg 18.4\pr & 14 & 1 & 277\dg 20.9\pr & 14 & 2 & 292\dg 23.3\pr & 14 & 3 & 307\dg 25.8\pr \\
14 & 4 & 322\dg 28.3\pr & 14 & 5 & 337\dg 30.7\pr & 14 & 6 & 352\dg 33.2\pr & 14 & 7 & 7\dg 35.7\pr \\
14 & 8 & 22\dg 38.1\pr & 14 & 9 & 37\dg 40.6\pr & 14 & 10 & 52\dg 43.1\pr & 14 & 11 & 67\dg 45.5\pr \\
14 & 12 & 82\dg 48.0\pr & 14 & 13 & 97\dg 50.5\pr & 14 & 14 & 112\dg 52.9\pr & 14 & 15 & 127\dg 55.4\pr \\
14 & 16 & 142\dg 57.8\pr & 14 & 17 & 158\dg 00.3\pr & 14 & 18 & 173\dg 02.8\pr & 14 & 19 & 188\dg 05.2\pr \\
14 & 20 & 203\dg 07.7\pr & 14 & 21 & 218\dg 10.2\pr & 14 & 22 & 233\dg 12.6\pr & 14 & 23 & 248\dg 15.1\pr \\
15 & 0 & 263\dg 17.6\pr & 15 & 1 & 278\dg 20.0\pr & 15 & 2 & 293\dg 22.5\pr & 15 & 3 & 308\dg 25.0\pr \\
15 & 4 & 323\dg 27.4\pr & 15 & 5 & 338\dg 29.9\pr & 15 & 6 & 353\dg 32.3\pr & 15 & 7 & 8\dg 34.8\pr \\
15 & 8 & 23\dg 37.3\pr & 15 & 9 & 38\dg 39.7\pr & 15 & 10 & 53\dg 42.2\pr & 15 & 11 & 68\dg 44.7\pr \\
15 & 12 & 83\dg 47.1\pr & 15 & 13 & 98\dg 49.6\pr & 15 & 14 & 113\dg 52.1\pr & 15 & 15 & 128\dg 54.5\pr \\
15 & 16 & 143\dg 57.0\pr & 15 & 17 & 158\dg 59.4\pr & 15 & 18 & 174\dg 01.9\pr & 15 & 19 & 189\dg 04.4\pr \\
15 & 20 & 204\dg 06.8\pr & 15 & 21 & 219\dg 09.3\pr & 15 & 22 & 234\dg 11.8\pr & 15 & 23 & 249\dg 14.2\pr \\
16 & 0 & 264\dg 16.7\pr & 16 & 1 & 279\dg 19.2\pr & 16 & 2 & 294\dg 21.6\pr & 16 & 3 & 309\dg 24.1\pr \\
16 & 4 & 324\dg 26.6\pr & 16 & 5 & 339\dg 29.0\pr & 16 & 6 & 354\dg 31.5\pr & 16 & 7 & 9\dg 33.9\pr \\
16 & 8 & 24\dg 36.4\pr & 16 & 9 & 39\dg 38.9\pr & 16 & 10 & 54\dg 41.3\pr & 16 & 11 & 69\dg 43.8\pr \\
16 & 12 & 84\dg 46.3\pr & 16 & 13 & 99\dg 48.7\pr & 16 & 14 & 114\dg 51.2\pr & 16 & 15 & 129\dg 53.7\pr \\
16 & 16 & 144\dg 56.1\pr & 16 & 17 & 159\dg 58.6\pr & 16 & 18 & 175\dg 01.1\pr & 16 & 19 & 190\dg 03.5\pr \\
16 & 20 & 205\dg 06.0\pr & 16 & 21 & 220\dg 08.4\pr & 16 & 22 & 235\dg 10.9\pr & 16 & 23 & 250\dg 13.4\pr \\
17 & 0 & 265\dg 15.8\pr & 17 & 1 & 280\dg 18.3\pr & 17 & 2 & 295\dg 20.8\pr & 17 & 3 & 310\dg 23.2\pr \\
17 & 4 & 325\dg 25.7\pr & 17 & 5 & 340\dg 28.2\pr & 17 & 6 & 355\dg 30.6\pr & 17 & 7 & 10\dg 33.1\pr \\
17 & 8 & 25\dg 35.5\pr & 17 & 9 & 40\dg 38.0\pr & 17 & 10 & 55\dg 40.5\pr & 17 & 11 & 70\dg 42.9\pr \\
17 & 12 & 85\dg 45.4\pr & 17 & 13 & 100\dg 47.9\pr & 17 & 14 & 115\dg 50.3\pr & 17 & 15 & 130\dg 52.8\pr \\
17 & 16 & 145\dg 55.3\pr & 17 & 17 & 160\dg 57.7\pr & 17 & 18 & 176\dg 00.2\pr & 17 & 19 & 191\dg 02.7\pr \\
17 & 20 & 206\dg 05.1\pr & 17 & 21 & 221\dg 07.6\pr & 17 & 22 & 236\dg 10.0\pr & 17 & 23 & 251\dg 12.5\pr \\
18 & 0 & 266\dg 15.0\pr & 18 & 1 & 281\dg 17.4\pr & 18 & 2 & 296\dg 19.9\pr & 18 & 3 & 311\dg 22.4\pr \\
18 & 4 & 326\dg 24.8\pr & 18 & 5 & 341\dg 27.3\pr & 18 & 6 & 356\dg 29.8\pr & 18 & 7 & 11\dg 32.2\pr \\
18 & 8 & 26\dg 34.7\pr & 18 & 9 & 41\dg 37.2\pr & 18 & 10 & 56\dg 39.6\pr & 18 & 11 & 71\dg 42.1\pr \\
18 & 12 & 86\dg 44.5\pr & 18 & 13 & 101\dg 47.0\pr & 18 & 14 & 116\dg 49.5\pr & 18 & 15 & 131\dg 51.9\pr \\
18 & 16 & 146\dg 54.4\pr & 18 & 17 & 161\dg 56.9\pr & 18 & 18 & 176\dg 59.3\pr & 18 & 19 & 192\dg 01.8\pr \\
18 & 20 & 207\dg 04.3\pr & 18 & 21 & 222\dg 06.7\pr & 18 & 22 & 237\dg 09.2\pr & 18 & 23 & 252\dg 11.6\pr \\
19 & 0 & 267\dg 14.1\pr & 19 & 1 & 282\dg 16.6\pr & 19 & 2 & 297\dg 19.0\pr & 19 & 3 & 312\dg 21.5\pr \\
19 & 4 & 327\dg 24.0\pr & 19 & 5 & 342\dg 26.4\pr & 19 & 6 & 357\dg 28.9\pr & 19 & 7 & 12\dg 31.4\pr \\
19 & 8 & 27\dg 33.8\pr & 19 & 9 & 42\dg 36.3\pr & 19 & 10 & 57\dg 38.8\pr & 19 & 11 & 72\dg 41.2\pr \\
19 & 12 & 87\dg 43.7\pr & 19 & 13 & 102\dg 46.1\pr & 19 & 14 & 117\dg 48.6\pr & 19 & 15 & 132\dg 51.1\pr \\
19 & 16 & 147\dg 53.5\pr & 19 & 17 & 162\dg 56.0\pr & 19 & 18 & 177\dg 58.5\pr & 19 & 19 & 193\dg 00.9\pr \\
19 & 20 & 208\dg 03.4\pr & 19 & 21 & 223\dg 05.9\pr & 19 & 22 & 238\dg 08.3\pr & 19 & 23 & 253\dg 10.8\pr \\
20 & 0 & 268\dg 13.3\pr & 20 & 1 & 283\dg 15.7\pr & 20 & 2 & 298\dg 18.2\pr & 20 & 3 & 313\dg 20.6\pr \\
20 & 4 & 328\dg 23.1\pr & 20 & 5 & 343\dg 25.6\pr & 20 & 6 & 358\dg 28.0\pr & 20 & 7 & 13\dg 30.5\pr \\
20 & 8 & 28\dg 33.0\pr & 20 & 9 & 43\dg 35.4\pr & 20 & 10 & 58\dg 37.9\pr & 20 & 11 & 73\dg 40.4\pr \\
20 & 12 & 88\dg 42.8\pr & 20 & 13 & 103\dg 45.3\pr & 20 & 14 & 118\dg 47.7\pr & 20 & 15 & 133\dg 50.2\pr \\
20 & 16 & 148\dg 52.7\pr & 20 & 17 & 163\dg 55.1\pr & 20 & 18 & 178\dg 57.6\pr & 20 & 19 & 194\dg 00.1\pr \\
20 & 20 & 209\dg 02.5\pr & 20 & 21 & 224\dg 05.0\pr & 20 & 22 & 239\dg 07.5\pr & 20 & 23 & 254\dg 09.9\pr \\
21 & 0 & 269\dg 12.4\pr & 21 & 1 & 284\dg 14.9\pr & 21 & 2 & 299\dg 17.3\pr & 21 & 3 & 314\dg 19.8\pr \\
21 & 4 & 329\dg 22.2\pr & 21 & 5 & 344\dg 24.7\pr & 21 & 6 & 359\dg 27.2\pr & 21 & 7 & 14\dg 29.6\pr \\
21 & 8 & 29\dg 32.1\pr & 21 & 9 & 44\dg 34.6\pr & 21 & 10 & 59\dg 37.0\pr & 21 & 11 & 74\dg 39.5\pr \\
21 & 12 & 89\dg 42.0\pr & 21 & 13 & 104\dg 44.4\pr & 21 & 14 & 119\dg 46.9\pr & 21 & 15 & 134\dg 49.4\pr \\
21 & 16 & 149\dg 51.8\pr & 21 & 17 & 164\dg 54.3\pr & 21 & 18 & 179\dg 56.7\pr & 21 & 19 & 194\dg 59.2\pr \\
21 & 20 & 210\dg 01.7\pr & 21 & 21 & 225\dg 04.1\pr & 21 & 22 & 240\dg 06.6\pr & 21 & 23 & 255\dg 09.1\pr \\
22 & 0 & 270\dg 11.5\pr & 22 & 1 & 285\dg 14.0\pr & 22 & 2 & 300\dg 16.5\pr & 22 & 3 & 315\dg 18.9\pr \\
22 & 4 & 330\dg 21.4\pr & 22 & 5 & 345\dg 23.9\pr & 22 & 6 & 0\dg 26.3\pr & 22 & 7 & 15\dg 28.8\pr \\
22 & 8 & 30\dg 31.2\pr & 22 & 9 & 45\dg 33.7\pr & 22 & 10 & 60\dg 36.2\pr & 22 & 11 & 75\dg 38.6\pr \\
22 & 12 & 90\dg 41.1\pr & 22 & 13 & 105\dg 43.6\pr & 22 & 14 & 120\dg 46.0\pr & 22 & 15 & 135\dg 48.5\pr \\
22 & 16 & 150\dg 51.0\pr & 22 & 17 & 165\dg 53.4\pr & 22 & 18 & 180\dg 55.9\pr & 22 & 19 & 195\dg 58.3\pr \\
22 & 20 & 211\dg 00.8\pr & 22 & 21 & 226\dg 03.3\pr & 22 & 22 & 241\dg 05.7\pr & 22 & 23 & 256\dg 08.2\pr \\
23 & 0 & 271\dg 10.7\pr & 23 & 1 & 286\dg 13.1\pr & 23 & 2 & 301\dg 15.6\pr & 23 & 3 & 316\dg 18.1\pr \\
23 & 4 & 331\dg 20.5\pr & 23 & 5 & 346\dg 23.0\pr & 23 & 6 & 1\dg 25.5\pr & 23 & 7 & 16\dg 27.9\pr \\
23 & 8 & 31\dg 30.4\pr & 23 & 9 & 46\dg 32.8\pr & 23 & 10 & 61\dg 35.3\pr & 23 & 11 & 76\dg 37.8\pr \\
23 & 12 & 91\dg 40.2\pr & 23 & 13 & 106\dg 42.7\pr & 23 & 14 & 121\dg 45.2\pr & 23 & 15 & 136\dg 47.6\pr \\
23 & 16 & 151\dg 50.1\pr & 23 & 17 & 166\dg 52.6\pr & 23 & 18 & 181\dg 55.0\pr & 23 & 19 & 196\dg 57.5\pr \\
23 & 20 & 212\dg 00.0\pr & 23 & 21 & 227\dg 02.4\pr & 23 & 22 & 242\dg 04.9\pr & 23 & 23 & 257\dg 07.3\pr \\
24 & 0 & 272\dg 09.8\pr & 24 & 1 & 287\dg 12.3\pr & 24 & 2 & 302\dg 14.7\pr & 24 & 3 & 317\dg 17.2\pr \\
24 & 4 & 332\dg 19.7\pr & 24 & 5 & 347\dg 22.1\pr & 24 & 6 & 2\dg 24.6\pr & 24 & 7 & 17\dg 27.1\pr \\
24 & 8 & 32\dg 29.5\pr & 24 & 9 & 47\dg 32.0\pr & 24 & 10 & 62\dg 34.4\pr & 24 & 11 & 77\dg 36.9\pr \\
24 & 12 & 92\dg 39.4\pr & 24 & 13 & 107\dg 41.8\pr & 24 & 14 & 122\dg 44.3\pr & 24 & 15 & 137\dg 46.8\pr \\
24 & 16 & 152\dg 49.2\pr & 24 & 17 & 167\dg 51.7\pr & 24 & 18 & 182\dg 54.2\pr & 24 & 19 & 197\dg 56.6\pr \\
24 & 20 & 212\dg 59.1\pr & 24 & 21 & 228\dg 01.6\pr & 24 & 22 & 243\dg 04.0\pr & 24 & 23 & 258\dg 06.5\pr \\
25 & 0 & 273\dg 08.9\pr & 25 & 1 & 288\dg 11.4\pr & 25 & 2 & 303\dg 13.9\pr & 25 & 3 & 318\dg 16.3\pr \\
25 & 4 & 333\dg 18.8\pr & 25 & 5 & 348\dg 21.3\pr & 25 & 6 & 3\dg 23.7\pr & 25 & 7 & 18\dg 26.2\pr \\
25 & 8 & 33\dg 28.7\pr & 25 & 9 & 48\dg 31.1\pr & 25 & 10 & 63\dg 33.6\pr & 25 & 11 & 78\dg 36.1\pr \\
25 & 12 & 93\dg 38.5\pr & 25 & 13 & 108\dg 41.0\pr & 25 & 14 & 123\dg 43.4\pr & 25 & 15 & 138\dg 45.9\pr \\
25 & 16 & 153\dg 48.4\pr & 25 & 17 & 168\dg 50.8\pr & 25 & 18 & 183\dg 53.3\pr & 25 & 19 & 198\dg 55.8\pr \\
25 & 20 & 213\dg 58.2\pr & 25 & 21 & 229\dg 00.7\pr & 25 & 22 & 244\dg 03.2\pr & 25 & 23 & 259\dg 05.6\pr \\
26 & 0 & 274\dg 08.1\pr & 26 & 1 & 289\dg 10.5\pr & 26 & 2 & 304\dg 13.0\pr & 26 & 3 & 319\dg 15.5\pr \\
26 & 4 & 334\dg 17.9\pr & 26 & 5 & 349\dg 20.4\pr & 26 & 6 & 4\dg 22.9\pr & 26 & 7 & 19\dg 25.3\pr \\
26 & 8 & 34\dg 27.8\pr & 26 & 9 & 49\dg 30.3\pr & 26 & 10 & 64\dg 32.7\pr & 26 & 11 & 79\dg 35.2\pr \\
26 & 12 & 94\dg 37.7\pr & 26 & 13 & 109\dg 40.1\pr & 26 & 14 & 124\dg 42.6\pr & 26 & 15 & 139\dg 45.0\pr \\
26 & 16 & 154\dg 47.5\pr & 26 & 17 & 169\dg 50.0\pr & 26 & 18 & 184\dg 52.4\pr & 26 & 19 & 199\dg 54.9\pr \\
26 & 20 & 214\dg 57.4\pr & 26 & 21 & 229\dg 59.8\pr & 26 & 22 & 245\dg 02.3\pr & 26 & 23 & 260\dg 04.8\pr \\
27 & 0 & 275\dg 07.2\pr & 27 & 1 & 290\dg 09.7\pr & 27 & 2 & 305\dg 12.2\pr & 27 & 3 & 320\dg 14.6\pr \\
27 & 4 & 335\dg 17.1\pr & 27 & 5 & 350\dg 19.5\pr & 27 & 6 & 5\dg 22.0\pr & 27 & 7 & 20\dg 24.5\pr \\
27 & 8 & 35\dg 26.9\pr & 27 & 9 & 50\dg 29.4\pr & 27 & 10 & 65\dg 31.9\pr & 27 & 11 & 80\dg 34.3\pr \\
27 & 12 & 95\dg 36.8\pr & 27 & 13 & 110\dg 39.3\pr & 27 & 14 & 125\dg 41.7\pr & 27 & 15 & 140\dg 44.2\pr \\
27 & 16 & 155\dg 46.7\pr & 27 & 17 & 170\dg 49.1\pr & 27 & 18 & 185\dg 51.6\pr & 27 & 19 & 200\dg 54.0\pr \\
27 & 20 & 215\dg 56.5\pr & 27 & 21 & 230\dg 59.0\pr & 27 & 22 & 246\dg 01.4\pr & 27 & 23 & 261\dg 03.9\pr \\
28 & 0 & 276\dg 06.4\pr & 28 & 1 & 291\dg 08.8\pr & 28 & 2 & 306\dg 11.3\pr & 28 & 3 & 321\dg 13.8\pr \\
28 & 4 & 336\dg 16.2\pr & 28 & 5 & 351\dg 18.7\pr & 28 & 6 & 6\dg 21.1\pr & 28 & 7 & 21\dg 23.6\pr \\
28 & 8 & 36\dg 26.1\pr & 28 & 9 & 51\dg 28.5\pr & 28 & 10 & 66\dg 31.0\pr & 28 & 11 & 81\dg 33.5\pr \\
28 & 12 & 96\dg 35.9\pr & 28 & 13 & 111\dg 38.4\pr & 28 & 14 & 126\dg 40.9\pr & 28 & 15 & 141\dg 43.3\pr \\
28 & 16 & 156\dg 45.8\pr & 28 & 17 & 171\dg 48.3\pr & 28 & 18 & 186\dg 50.7\pr & 28 & 19 & 201\dg 53.2\pr \\
28 & 20 & 216\dg 55.6\pr & 28 & 21 & 231\dg 58.1\pr & 28 & 22 & 247\dg 00.6\pr & 28 & 23 & 262\dg 03.0\pr \\
29 & 0 & 277\dg 05.5\pr & 29 & 1 & 292\dg 08.0\pr & 29 & 2 & 307\dg 10.4\pr & 29 & 3 & 322\dg 12.9\pr \\
29 & 4 & 337\dg 15.4\pr & 29 & 5 & 352\dg 17.8\pr & 29 & 6 & 7\dg 20.3\pr & 29 & 7 & 22\dg 22.8\pr \\
29 & 8 & 37\dg 25.2\pr & 29 & 9 & 52\dg 27.7\pr & 29 & 10 & 67\dg 30.1\pr & 29 & 11 & 82\dg 32.6\pr \\
29 & 12 & 97\dg 35.1\pr & 29 & 13 & 112\dg 37.5\pr & 29 & 14 & 127\dg 40.0\pr & 29 & 15 & 142\dg 42.5\pr \\
29 & 16 & 157\dg 44.9\pr & 29 & 17 & 172\dg 47.4\pr & 29 & 18 & 187\dg 49.9\pr & 29 & 19 & 202\dg 52.3\pr \\
29 & 20 & 217\dg 54.8\pr & 29 & 21 & 232\dg 57.2\pr & 29 & 22 & 247\dg 59.7\pr & 29 & 23 & 263\dg 02.2\pr \\
30 & 0 & 278\dg 04.6\pr & 30 & 1 & 293\dg 07.1\pr & 30 & 2 & 308\dg 09.6\pr & 30 & 3 & 323\dg 12.0\pr \\
30 & 4 & 338\dg 14.5\pr & 30 & 5 & 353\dg 17.0\pr & 30 & 6 & 8\dg 19.4\pr & 30 & 7 & 23\dg 21.9\pr \\
30 & 8 & 38\dg 24.4\pr & 30 & 9 & 53\dg 26.8\pr & 30 & 10 & 68\dg 29.3\pr & 30 & 11 & 83\dg 31.7\pr \\
30 & 12 & 98\dg 34.2\pr & 30 & 13 & 113\dg 36.7\pr & 30 & 14 & 128\dg 39.1\pr & 30 & 15 & 143\dg 41.6\pr \\
30 & 16 & 158\dg 44.1\pr & 30 & 17 & 173\dg 46.5\pr & 30 & 18 & 188\dg 49.0\pr & 30 & 19 & 203\dg 51.5\pr \\
30 & 20 & 218\dg 53.9\pr & 30 & 21 & 233\dg 56.4\pr & 30 & 22 & 248\dg 58.9\pr & 30 & 23 & 264\dg 01.3\pr \\
\end{longtable}}

{\footnotesize
\begin{longtable}{rrr|rrr|rrr|rrr}
\caption{Aries --- GHA of the First Point of Aries, July 2026 (UT)}\label{aries07}\\
\toprule
\textbf{d} & \textbf{h} & \textbf{GHA} $\gamma$ & \textbf{d} & \textbf{h} & \textbf{GHA} $\gamma$ & \textbf{d} & \textbf{h} & \textbf{GHA} $\gamma$ & \textbf{d} & \textbf{h} & \textbf{GHA} $\gamma$ \\
\midrule
\endfirsthead
\multicolumn{12}{c}{\textit{Aries --- GHA of the First Point of Aries, July 2026 (UT) (continued)}}\\
\toprule
\textbf{d} & \textbf{h} & \textbf{GHA} $\gamma$ & \textbf{d} & \textbf{h} & \textbf{GHA} $\gamma$ & \textbf{d} & \textbf{h} & \textbf{GHA} $\gamma$ & \textbf{d} & \textbf{h} & \textbf{GHA} $\gamma$ \\
\midrule
\endhead
\midrule\multicolumn{12}{r}{\textit{continued on next page}}\\
\endfoot
\bottomrule
\endlastfoot
1 & 0 & 279\dg 03.8\pr & 1 & 1 & 294\dg 06.2\pr & 1 & 2 & 309\dg 08.7\pr & 1 & 3 & 324\dg 11.2\pr \\
1 & 4 & 339\dg 13.6\pr & 1 & 5 & 354\dg 16.1\pr & 1 & 6 & 9\dg 18.6\pr & 1 & 7 & 24\dg 21.0\pr \\
1 & 8 & 39\dg 23.5\pr & 1 & 9 & 54\dg 26.0\pr & 1 & 10 & 69\dg 28.4\pr & 1 & 11 & 84\dg 30.9\pr \\
1 & 12 & 99\dg 33.3\pr & 1 & 13 & 114\dg 35.8\pr & 1 & 14 & 129\dg 38.3\pr & 1 & 15 & 144\dg 40.7\pr \\
1 & 16 & 159\dg 43.2\pr & 1 & 17 & 174\dg 45.7\pr & 1 & 18 & 189\dg 48.1\pr & 1 & 19 & 204\dg 50.6\pr \\
1 & 20 & 219\dg 53.1\pr & 1 & 21 & 234\dg 55.5\pr & 1 & 22 & 249\dg 58.0\pr & 1 & 23 & 265\dg 00.5\pr \\
2 & 0 & 280\dg 02.9\pr & 2 & 1 & 295\dg 05.4\pr & 2 & 2 & 310\dg 07.8\pr & 2 & 3 & 325\dg 10.3\pr \\
2 & 4 & 340\dg 12.8\pr & 2 & 5 & 355\dg 15.2\pr & 2 & 6 & 10\dg 17.7\pr & 2 & 7 & 25\dg 20.2\pr \\
2 & 8 & 40\dg 22.6\pr & 2 & 9 & 55\dg 25.1\pr & 2 & 10 & 70\dg 27.6\pr & 2 & 11 & 85\dg 30.0\pr \\
2 & 12 & 100\dg 32.5\pr & 2 & 13 & 115\dg 35.0\pr & 2 & 14 & 130\dg 37.4\pr & 2 & 15 & 145\dg 39.9\pr \\
2 & 16 & 160\dg 42.3\pr & 2 & 17 & 175\dg 44.8\pr & 2 & 18 & 190\dg 47.3\pr & 2 & 19 & 205\dg 49.7\pr \\
2 & 20 & 220\dg 52.2\pr & 2 & 21 & 235\dg 54.7\pr & 2 & 22 & 250\dg 57.1\pr & 2 & 23 & 265\dg 59.6\pr \\
3 & 0 & 281\dg 02.1\pr & 3 & 1 & 296\dg 04.5\pr & 3 & 2 & 311\dg 07.0\pr & 3 & 3 & 326\dg 09.4\pr \\
3 & 4 & 341\dg 11.9\pr & 3 & 5 & 356\dg 14.4\pr & 3 & 6 & 11\dg 16.8\pr & 3 & 7 & 26\dg 19.3\pr \\
3 & 8 & 41\dg 21.8\pr & 3 & 9 & 56\dg 24.2\pr & 3 & 10 & 71\dg 26.7\pr & 3 & 11 & 86\dg 29.2\pr \\
3 & 12 & 101\dg 31.6\pr & 3 & 13 & 116\dg 34.1\pr & 3 & 14 & 131\dg 36.6\pr & 3 & 15 & 146\dg 39.0\pr \\
3 & 16 & 161\dg 41.5\pr & 3 & 17 & 176\dg 43.9\pr & 3 & 18 & 191\dg 46.4\pr & 3 & 19 & 206\dg 48.9\pr \\
3 & 20 & 221\dg 51.3\pr & 3 & 21 & 236\dg 53.8\pr & 3 & 22 & 251\dg 56.3\pr & 3 & 23 & 266\dg 58.7\pr \\
4 & 0 & 282\dg 01.2\pr & 4 & 1 & 297\dg 03.7\pr & 4 & 2 & 312\dg 06.1\pr & 4 & 3 & 327\dg 08.6\pr \\
4 & 4 & 342\dg 11.1\pr & 4 & 5 & 357\dg 13.5\pr & 4 & 6 & 12\dg 16.0\pr & 4 & 7 & 27\dg 18.4\pr \\
4 & 8 & 42\dg 20.9\pr & 4 & 9 & 57\dg 23.4\pr & 4 & 10 & 72\dg 25.8\pr & 4 & 11 & 87\dg 28.3\pr \\
4 & 12 & 102\dg 30.8\pr & 4 & 13 & 117\dg 33.2\pr & 4 & 14 & 132\dg 35.7\pr & 4 & 15 & 147\dg 38.2\pr \\
4 & 16 & 162\dg 40.6\pr & 4 & 17 & 177\dg 43.1\pr & 4 & 18 & 192\dg 45.6\pr & 4 & 19 & 207\dg 48.0\pr \\
4 & 20 & 222\dg 50.5\pr & 4 & 21 & 237\dg 52.9\pr & 4 & 22 & 252\dg 55.4\pr & 4 & 23 & 267\dg 57.9\pr \\
5 & 0 & 283\dg 00.3\pr & 5 & 1 & 298\dg 02.8\pr & 5 & 2 & 313\dg 05.3\pr & 5 & 3 & 328\dg 07.7\pr \\
5 & 4 & 343\dg 10.2\pr & 5 & 5 & 358\dg 12.7\pr & 5 & 6 & 13\dg 15.1\pr & 5 & 7 & 28\dg 17.6\pr \\
5 & 8 & 43\dg 20.0\pr & 5 & 9 & 58\dg 22.5\pr & 5 & 10 & 73\dg 25.0\pr & 5 & 11 & 88\dg 27.4\pr \\
5 & 12 & 103\dg 29.9\pr & 5 & 13 & 118\dg 32.4\pr & 5 & 14 & 133\dg 34.8\pr & 5 & 15 & 148\dg 37.3\pr \\
5 & 16 & 163\dg 39.8\pr & 5 & 17 & 178\dg 42.2\pr & 5 & 18 & 193\dg 44.7\pr & 5 & 19 & 208\dg 47.2\pr \\
5 & 20 & 223\dg 49.6\pr & 5 & 21 & 238\dg 52.1\pr & 5 & 22 & 253\dg 54.5\pr & 5 & 23 & 268\dg 57.0\pr \\
6 & 0 & 283\dg 59.5\pr & 6 & 1 & 299\dg 01.9\pr & 6 & 2 & 314\dg 04.4\pr & 6 & 3 & 329\dg 06.9\pr \\
6 & 4 & 344\dg 09.3\pr & 6 & 5 & 359\dg 11.8\pr & 6 & 6 & 14\dg 14.3\pr & 6 & 7 & 29\dg 16.7\pr \\
6 & 8 & 44\dg 19.2\pr & 6 & 9 & 59\dg 21.7\pr & 6 & 10 & 74\dg 24.1\pr & 6 & 11 & 89\dg 26.6\pr \\
6 & 12 & 104\dg 29.0\pr & 6 & 13 & 119\dg 31.5\pr & 6 & 14 & 134\dg 34.0\pr & 6 & 15 & 149\dg 36.4\pr \\
6 & 16 & 164\dg 38.9\pr & 6 & 17 & 179\dg 41.4\pr & 6 & 18 & 194\dg 43.8\pr & 6 & 19 & 209\dg 46.3\pr \\
6 & 20 & 224\dg 48.8\pr & 6 & 21 & 239\dg 51.2\pr & 6 & 22 & 254\dg 53.7\pr & 6 & 23 & 269\dg 56.1\pr \\
7 & 0 & 284\dg 58.6\pr & 7 & 1 & 300\dg 01.1\pr & 7 & 2 & 315\dg 03.5\pr & 7 & 3 & 330\dg 06.0\pr \\
7 & 4 & 345\dg 08.5\pr & 7 & 5 & 0\dg 10.9\pr & 7 & 6 & 15\dg 13.4\pr & 7 & 7 & 30\dg 15.9\pr \\
7 & 8 & 45\dg 18.3\pr & 7 & 9 & 60\dg 20.8\pr & 7 & 10 & 75\dg 23.3\pr & 7 & 11 & 90\dg 25.7\pr \\
7 & 12 & 105\dg 28.2\pr & 7 & 13 & 120\dg 30.6\pr & 7 & 14 & 135\dg 33.1\pr & 7 & 15 & 150\dg 35.6\pr \\
7 & 16 & 165\dg 38.0\pr & 7 & 17 & 180\dg 40.5\pr & 7 & 18 & 195\dg 43.0\pr & 7 & 19 & 210\dg 45.4\pr \\
7 & 20 & 225\dg 47.9\pr & 7 & 21 & 240\dg 50.4\pr & 7 & 22 & 255\dg 52.8\pr & 7 & 23 & 270\dg 55.3\pr \\
8 & 0 & 285\dg 57.8\pr & 8 & 1 & 301\dg 00.2\pr & 8 & 2 & 316\dg 02.7\pr & 8 & 3 & 331\dg 05.1\pr \\
8 & 4 & 346\dg 07.6\pr & 8 & 5 & 1\dg 10.1\pr & 8 & 6 & 16\dg 12.5\pr & 8 & 7 & 31\dg 15.0\pr \\
8 & 8 & 46\dg 17.5\pr & 8 & 9 & 61\dg 19.9\pr & 8 & 10 & 76\dg 22.4\pr & 8 & 11 & 91\dg 24.9\pr \\
8 & 12 & 106\dg 27.3\pr & 8 & 13 & 121\dg 29.8\pr & 8 & 14 & 136\dg 32.2\pr & 8 & 15 & 151\dg 34.7\pr \\
8 & 16 & 166\dg 37.2\pr & 8 & 17 & 181\dg 39.6\pr & 8 & 18 & 196\dg 42.1\pr & 8 & 19 & 211\dg 44.6\pr \\
8 & 20 & 226\dg 47.0\pr & 8 & 21 & 241\dg 49.5\pr & 8 & 22 & 256\dg 52.0\pr & 8 & 23 & 271\dg 54.4\pr \\
9 & 0 & 286\dg 56.9\pr & 9 & 1 & 301\dg 59.4\pr & 9 & 2 & 317\dg 01.8\pr & 9 & 3 & 332\dg 04.3\pr \\
9 & 4 & 347\dg 06.7\pr & 9 & 5 & 2\dg 09.2\pr & 9 & 6 & 17\dg 11.7\pr & 9 & 7 & 32\dg 14.1\pr \\
9 & 8 & 47\dg 16.6\pr & 9 & 9 & 62\dg 19.1\pr & 9 & 10 & 77\dg 21.5\pr & 9 & 11 & 92\dg 24.0\pr \\
9 & 12 & 107\dg 26.5\pr & 9 & 13 & 122\dg 28.9\pr & 9 & 14 & 137\dg 31.4\pr & 9 & 15 & 152\dg 33.9\pr \\
9 & 16 & 167\dg 36.3\pr & 9 & 17 & 182\dg 38.8\pr & 9 & 18 & 197\dg 41.2\pr & 9 & 19 & 212\dg 43.7\pr \\
9 & 20 & 227\dg 46.2\pr & 9 & 21 & 242\dg 48.6\pr & 9 & 22 & 257\dg 51.1\pr & 9 & 23 & 272\dg 53.6\pr \\
10 & 0 & 287\dg 56.0\pr & 10 & 1 & 302\dg 58.5\pr & 10 & 2 & 318\dg 01.0\pr & 10 & 3 & 333\dg 03.4\pr \\
10 & 4 & 348\dg 05.9\pr & 10 & 5 & 3\dg 08.3\pr & 10 & 6 & 18\dg 10.8\pr & 10 & 7 & 33\dg 13.3\pr \\
10 & 8 & 48\dg 15.7\pr & 10 & 9 & 63\dg 18.2\pr & 10 & 10 & 78\dg 20.7\pr & 10 & 11 & 93\dg 23.1\pr \\
10 & 12 & 108\dg 25.6\pr & 10 & 13 & 123\dg 28.1\pr & 10 & 14 & 138\dg 30.5\pr & 10 & 15 & 153\dg 33.0\pr \\
10 & 16 & 168\dg 35.5\pr & 10 & 17 & 183\dg 37.9\pr & 10 & 18 & 198\dg 40.4\pr & 10 & 19 & 213\dg 42.8\pr \\
10 & 20 & 228\dg 45.3\pr & 10 & 21 & 243\dg 47.8\pr & 10 & 22 & 258\dg 50.2\pr & 10 & 23 & 273\dg 52.7\pr \\
11 & 0 & 288\dg 55.2\pr & 11 & 1 & 303\dg 57.6\pr & 11 & 2 & 319\dg 00.1\pr & 11 & 3 & 334\dg 02.6\pr \\
11 & 4 & 349\dg 05.0\pr & 11 & 5 & 4\dg 07.5\pr & 11 & 6 & 19\dg 10.0\pr & 11 & 7 & 34\dg 12.4\pr \\
11 & 8 & 49\dg 14.9\pr & 11 & 9 & 64\dg 17.3\pr & 11 & 10 & 79\dg 19.8\pr & 11 & 11 & 94\dg 22.3\pr \\
11 & 12 & 109\dg 24.7\pr & 11 & 13 & 124\dg 27.2\pr & 11 & 14 & 139\dg 29.7\pr & 11 & 15 & 154\dg 32.1\pr \\
11 & 16 & 169\dg 34.6\pr & 11 & 17 & 184\dg 37.1\pr & 11 & 18 & 199\dg 39.5\pr & 11 & 19 & 214\dg 42.0\pr \\
11 & 20 & 229\dg 44.5\pr & 11 & 21 & 244\dg 46.9\pr & 11 & 22 & 259\dg 49.4\pr & 11 & 23 & 274\dg 51.8\pr \\
12 & 0 & 289\dg 54.3\pr & 12 & 1 & 304\dg 56.8\pr & 12 & 2 & 319\dg 59.2\pr & 12 & 3 & 335\dg 01.7\pr \\
12 & 4 & 350\dg 04.2\pr & 12 & 5 & 5\dg 06.6\pr & 12 & 6 & 20\dg 09.1\pr & 12 & 7 & 35\dg 11.6\pr \\
12 & 8 & 50\dg 14.0\pr & 12 & 9 & 65\dg 16.5\pr & 12 & 10 & 80\dg 18.9\pr & 12 & 11 & 95\dg 21.4\pr \\
12 & 12 & 110\dg 23.9\pr & 12 & 13 & 125\dg 26.3\pr & 12 & 14 & 140\dg 28.8\pr & 12 & 15 & 155\dg 31.3\pr \\
12 & 16 & 170\dg 33.7\pr & 12 & 17 & 185\dg 36.2\pr & 12 & 18 & 200\dg 38.7\pr & 12 & 19 & 215\dg 41.1\pr \\
12 & 20 & 230\dg 43.6\pr & 12 & 21 & 245\dg 46.1\pr & 12 & 22 & 260\dg 48.5\pr & 12 & 23 & 275\dg 51.0\pr \\
13 & 0 & 290\dg 53.4\pr & 13 & 1 & 305\dg 55.9\pr & 13 & 2 & 320\dg 58.4\pr & 13 & 3 & 336\dg 00.8\pr \\
13 & 4 & 351\dg 03.3\pr & 13 & 5 & 6\dg 05.8\pr & 13 & 6 & 21\dg 08.2\pr & 13 & 7 & 36\dg 10.7\pr \\
13 & 8 & 51\dg 13.2\pr & 13 & 9 & 66\dg 15.6\pr & 13 & 10 & 81\dg 18.1\pr & 13 & 11 & 96\dg 20.6\pr \\
13 & 12 & 111\dg 23.0\pr & 13 & 13 & 126\dg 25.5\pr & 13 & 14 & 141\dg 27.9\pr & 13 & 15 & 156\dg 30.4\pr \\
13 & 16 & 171\dg 32.9\pr & 13 & 17 & 186\dg 35.3\pr & 13 & 18 & 201\dg 37.8\pr & 13 & 19 & 216\dg 40.3\pr \\
13 & 20 & 231\dg 42.7\pr & 13 & 21 & 246\dg 45.2\pr & 13 & 22 & 261\dg 47.7\pr & 13 & 23 & 276\dg 50.1\pr \\
14 & 0 & 291\dg 52.6\pr & 14 & 1 & 306\dg 55.0\pr & 14 & 2 & 321\dg 57.5\pr & 14 & 3 & 337\dg 00.0\pr \\
14 & 4 & 352\dg 02.4\pr & 14 & 5 & 7\dg 04.9\pr & 14 & 6 & 22\dg 07.4\pr & 14 & 7 & 37\dg 09.8\pr \\
14 & 8 & 52\dg 12.3\pr & 14 & 9 & 67\dg 14.8\pr & 14 & 10 & 82\dg 17.2\pr & 14 & 11 & 97\dg 19.7\pr \\
14 & 12 & 112\dg 22.2\pr & 14 & 13 & 127\dg 24.6\pr & 14 & 14 & 142\dg 27.1\pr & 14 & 15 & 157\dg 29.5\pr \\
14 & 16 & 172\dg 32.0\pr & 14 & 17 & 187\dg 34.5\pr & 14 & 18 & 202\dg 36.9\pr & 14 & 19 & 217\dg 39.4\pr \\
14 & 20 & 232\dg 41.9\pr & 14 & 21 & 247\dg 44.3\pr & 14 & 22 & 262\dg 46.8\pr & 14 & 23 & 277\dg 49.3\pr \\
15 & 0 & 292\dg 51.7\pr & 15 & 1 & 307\dg 54.2\pr & 15 & 2 & 322\dg 56.7\pr & 15 & 3 & 337\dg 59.1\pr \\
15 & 4 & 353\dg 01.6\pr & 15 & 5 & 8\dg 04.0\pr & 15 & 6 & 23\dg 06.5\pr & 15 & 7 & 38\dg 09.0\pr \\
15 & 8 & 53\dg 11.4\pr & 15 & 9 & 68\dg 13.9\pr & 15 & 10 & 83\dg 16.4\pr & 15 & 11 & 98\dg 18.8\pr \\
15 & 12 & 113\dg 21.3\pr & 15 & 13 & 128\dg 23.8\pr & 15 & 14 & 143\dg 26.2\pr & 15 & 15 & 158\dg 28.7\pr \\
15 & 16 & 173\dg 31.1\pr & 15 & 17 & 188\dg 33.6\pr & 15 & 18 & 203\dg 36.1\pr & 15 & 19 & 218\dg 38.5\pr \\
15 & 20 & 233\dg 41.0\pr & 15 & 21 & 248\dg 43.5\pr & 15 & 22 & 263\dg 45.9\pr & 15 & 23 & 278\dg 48.4\pr \\
16 & 0 & 293\dg 50.9\pr & 16 & 1 & 308\dg 53.3\pr & 16 & 2 & 323\dg 55.8\pr & 16 & 3 & 338\dg 58.3\pr \\
16 & 4 & 354\dg 00.7\pr & 16 & 5 & 9\dg 03.2\pr & 16 & 6 & 24\dg 05.6\pr & 16 & 7 & 39\dg 08.1\pr \\
16 & 8 & 54\dg 10.6\pr & 16 & 9 & 69\dg 13.0\pr & 16 & 10 & 84\dg 15.5\pr & 16 & 11 & 99\dg 18.0\pr \\
16 & 12 & 114\dg 20.4\pr & 16 & 13 & 129\dg 22.9\pr & 16 & 14 & 144\dg 25.4\pr & 16 & 15 & 159\dg 27.8\pr \\
16 & 16 & 174\dg 30.3\pr & 16 & 17 & 189\dg 32.8\pr & 16 & 18 & 204\dg 35.2\pr & 16 & 19 & 219\dg 37.7\pr \\
16 & 20 & 234\dg 40.1\pr & 16 & 21 & 249\dg 42.6\pr & 16 & 22 & 264\dg 45.1\pr & 16 & 23 & 279\dg 47.5\pr \\
17 & 0 & 294\dg 50.0\pr & 17 & 1 & 309\dg 52.5\pr & 17 & 2 & 324\dg 54.9\pr & 17 & 3 & 339\dg 57.4\pr \\
17 & 4 & 354\dg 59.9\pr & 17 & 5 & 10\dg 02.3\pr & 17 & 6 & 25\dg 04.8\pr & 17 & 7 & 40\dg 07.2\pr \\
17 & 8 & 55\dg 09.7\pr & 17 & 9 & 70\dg 12.2\pr & 17 & 10 & 85\dg 14.6\pr & 17 & 11 & 100\dg 17.1\pr \\
17 & 12 & 115\dg 19.6\pr & 17 & 13 & 130\dg 22.0\pr & 17 & 14 & 145\dg 24.5\pr & 17 & 15 & 160\dg 27.0\pr \\
17 & 16 & 175\dg 29.4\pr & 17 & 17 & 190\dg 31.9\pr & 17 & 18 & 205\dg 34.4\pr & 17 & 19 & 220\dg 36.8\pr \\
17 & 20 & 235\dg 39.3\pr & 17 & 21 & 250\dg 41.7\pr & 17 & 22 & 265\dg 44.2\pr & 17 & 23 & 280\dg 46.7\pr \\
18 & 0 & 295\dg 49.1\pr & 18 & 1 & 310\dg 51.6\pr & 18 & 2 & 325\dg 54.1\pr & 18 & 3 & 340\dg 56.5\pr \\
18 & 4 & 355\dg 59.0\pr & 18 & 5 & 11\dg 01.5\pr & 18 & 6 & 26\dg 03.9\pr & 18 & 7 & 41\dg 06.4\pr \\
18 & 8 & 56\dg 08.9\pr & 18 & 9 & 71\dg 11.3\pr & 18 & 10 & 86\dg 13.8\pr & 18 & 11 & 101\dg 16.2\pr \\
18 & 12 & 116\dg 18.7\pr & 18 & 13 & 131\dg 21.2\pr & 18 & 14 & 146\dg 23.6\pr & 18 & 15 & 161\dg 26.1\pr \\
18 & 16 & 176\dg 28.6\pr & 18 & 17 & 191\dg 31.0\pr & 18 & 18 & 206\dg 33.5\pr & 18 & 19 & 221\dg 36.0\pr \\
18 & 20 & 236\dg 38.4\pr & 18 & 21 & 251\dg 40.9\pr & 18 & 22 & 266\dg 43.4\pr & 18 & 23 & 281\dg 45.8\pr \\
19 & 0 & 296\dg 48.3\pr & 19 & 1 & 311\dg 50.7\pr & 19 & 2 & 326\dg 53.2\pr & 19 & 3 & 341\dg 55.7\pr \\
19 & 4 & 356\dg 58.1\pr & 19 & 5 & 12\dg 00.6\pr & 19 & 6 & 27\dg 03.1\pr & 19 & 7 & 42\dg 05.5\pr \\
19 & 8 & 57\dg 08.0\pr & 19 & 9 & 72\dg 10.5\pr & 19 & 10 & 87\dg 12.9\pr & 19 & 11 & 102\dg 15.4\pr \\
19 & 12 & 117\dg 17.8\pr & 19 & 13 & 132\dg 20.3\pr & 19 & 14 & 147\dg 22.8\pr & 19 & 15 & 162\dg 25.2\pr \\
19 & 16 & 177\dg 27.7\pr & 19 & 17 & 192\dg 30.2\pr & 19 & 18 & 207\dg 32.6\pr & 19 & 19 & 222\dg 35.1\pr \\
19 & 20 & 237\dg 37.6\pr & 19 & 21 & 252\dg 40.0\pr & 19 & 22 & 267\dg 42.5\pr & 19 & 23 & 282\dg 45.0\pr \\
20 & 0 & 297\dg 47.4\pr & 20 & 1 & 312\dg 49.9\pr & 20 & 2 & 327\dg 52.3\pr & 20 & 3 & 342\dg 54.8\pr \\
20 & 4 & 357\dg 57.3\pr & 20 & 5 & 12\dg 59.7\pr & 20 & 6 & 28\dg 02.2\pr & 20 & 7 & 43\dg 04.7\pr \\
20 & 8 & 58\dg 07.1\pr & 20 & 9 & 73\dg 09.6\pr & 20 & 10 & 88\dg 12.1\pr & 20 & 11 & 103\dg 14.5\pr \\
20 & 12 & 118\dg 17.0\pr & 20 & 13 & 133\dg 19.5\pr & 20 & 14 & 148\dg 21.9\pr & 20 & 15 & 163\dg 24.4\pr \\
20 & 16 & 178\dg 26.8\pr & 20 & 17 & 193\dg 29.3\pr & 20 & 18 & 208\dg 31.8\pr & 20 & 19 & 223\dg 34.2\pr \\
20 & 20 & 238\dg 36.7\pr & 20 & 21 & 253\dg 39.2\pr & 20 & 22 & 268\dg 41.6\pr & 20 & 23 & 283\dg 44.1\pr \\
21 & 0 & 298\dg 46.6\pr & 21 & 1 & 313\dg 49.0\pr & 21 & 2 & 328\dg 51.5\pr & 21 & 3 & 343\dg 53.9\pr \\
21 & 4 & 358\dg 56.4\pr & 21 & 5 & 13\dg 58.9\pr & 21 & 6 & 29\dg 01.3\pr & 21 & 7 & 44\dg 03.8\pr \\
21 & 8 & 59\dg 06.3\pr & 21 & 9 & 74\dg 08.7\pr & 21 & 10 & 89\dg 11.2\pr & 21 & 11 & 104\dg 13.7\pr \\
21 & 12 & 119\dg 16.1\pr & 21 & 13 & 134\dg 18.6\pr & 21 & 14 & 149\dg 21.1\pr & 21 & 15 & 164\dg 23.5\pr \\
21 & 16 & 179\dg 26.0\pr & 21 & 17 & 194\dg 28.4\pr & 21 & 18 & 209\dg 30.9\pr & 21 & 19 & 224\dg 33.4\pr \\
21 & 20 & 239\dg 35.8\pr & 21 & 21 & 254\dg 38.3\pr & 21 & 22 & 269\dg 40.8\pr & 21 & 23 & 284\dg 43.2\pr \\
22 & 0 & 299\dg 45.7\pr & 22 & 1 & 314\dg 48.2\pr & 22 & 2 & 329\dg 50.6\pr & 22 & 3 & 344\dg 53.1\pr \\
22 & 4 & 359\dg 55.6\pr & 22 & 5 & 14\dg 58.0\pr & 22 & 6 & 30\dg 00.5\pr & 22 & 7 & 45\dg 02.9\pr \\
22 & 8 & 60\dg 05.4\pr & 22 & 9 & 75\dg 07.9\pr & 22 & 10 & 90\dg 10.3\pr & 22 & 11 & 105\dg 12.8\pr \\
22 & 12 & 120\dg 15.3\pr & 22 & 13 & 135\dg 17.7\pr & 22 & 14 & 150\dg 20.2\pr & 22 & 15 & 165\dg 22.7\pr \\
22 & 16 & 180\dg 25.1\pr & 22 & 17 & 195\dg 27.6\pr & 22 & 18 & 210\dg 30.0\pr & 22 & 19 & 225\dg 32.5\pr \\
22 & 20 & 240\dg 35.0\pr & 22 & 21 & 255\dg 37.4\pr & 22 & 22 & 270\dg 39.9\pr & 22 & 23 & 285\dg 42.4\pr \\
23 & 0 & 300\dg 44.8\pr & 23 & 1 & 315\dg 47.3\pr & 23 & 2 & 330\dg 49.8\pr & 23 & 3 & 345\dg 52.2\pr \\
23 & 4 & 0\dg 54.7\pr & 23 & 5 & 15\dg 57.2\pr & 23 & 6 & 30\dg 59.6\pr & 23 & 7 & 46\dg 02.1\pr \\
23 & 8 & 61\dg 04.5\pr & 23 & 9 & 76\dg 07.0\pr & 23 & 10 & 91\dg 09.5\pr & 23 & 11 & 106\dg 11.9\pr \\
23 & 12 & 121\dg 14.4\pr & 23 & 13 & 136\dg 16.9\pr & 23 & 14 & 151\dg 19.3\pr & 23 & 15 & 166\dg 21.8\pr \\
23 & 16 & 181\dg 24.3\pr & 23 & 17 & 196\dg 26.7\pr & 23 & 18 & 211\dg 29.2\pr & 23 & 19 & 226\dg 31.7\pr \\
23 & 20 & 241\dg 34.1\pr & 23 & 21 & 256\dg 36.6\pr & 23 & 22 & 271\dg 39.0\pr & 23 & 23 & 286\dg 41.5\pr \\
24 & 0 & 301\dg 44.0\pr & 24 & 1 & 316\dg 46.4\pr & 24 & 2 & 331\dg 48.9\pr & 24 & 3 & 346\dg 51.4\pr \\
24 & 4 & 1\dg 53.8\pr & 24 & 5 & 16\dg 56.3\pr & 24 & 6 & 31\dg 58.8\pr & 24 & 7 & 47\dg 01.2\pr \\
24 & 8 & 62\dg 03.7\pr & 24 & 9 & 77\dg 06.1\pr & 24 & 10 & 92\dg 08.6\pr & 24 & 11 & 107\dg 11.1\pr \\
24 & 12 & 122\dg 13.5\pr & 24 & 13 & 137\dg 16.0\pr & 24 & 14 & 152\dg 18.5\pr & 24 & 15 & 167\dg 20.9\pr \\
24 & 16 & 182\dg 23.4\pr & 24 & 17 & 197\dg 25.9\pr & 24 & 18 & 212\dg 28.3\pr & 24 & 19 & 227\dg 30.8\pr \\
24 & 20 & 242\dg 33.3\pr & 24 & 21 & 257\dg 35.7\pr & 24 & 22 & 272\dg 38.2\pr & 24 & 23 & 287\dg 40.6\pr \\
25 & 0 & 302\dg 43.1\pr & 25 & 1 & 317\dg 45.6\pr & 25 & 2 & 332\dg 48.0\pr & 25 & 3 & 347\dg 50.5\pr \\
25 & 4 & 2\dg 53.0\pr & 25 & 5 & 17\dg 55.4\pr & 25 & 6 & 32\dg 57.9\pr & 25 & 7 & 48\dg 00.4\pr \\
25 & 8 & 63\dg 02.8\pr & 25 & 9 & 78\dg 05.3\pr & 25 & 10 & 93\dg 07.8\pr & 25 & 11 & 108\dg 10.2\pr \\
25 & 12 & 123\dg 12.7\pr & 25 & 13 & 138\dg 15.1\pr & 25 & 14 & 153\dg 17.6\pr & 25 & 15 & 168\dg 20.1\pr \\
25 & 16 & 183\dg 22.5\pr & 25 & 17 & 198\dg 25.0\pr & 25 & 18 & 213\dg 27.5\pr & 25 & 19 & 228\dg 29.9\pr \\
25 & 20 & 243\dg 32.4\pr & 25 & 21 & 258\dg 34.9\pr & 25 & 22 & 273\dg 37.3\pr & 25 & 23 & 288\dg 39.8\pr \\
26 & 0 & 303\dg 42.3\pr & 26 & 1 & 318\dg 44.7\pr & 26 & 2 & 333\dg 47.2\pr & 26 & 3 & 348\dg 49.6\pr \\
26 & 4 & 3\dg 52.1\pr & 26 & 5 & 18\dg 54.6\pr & 26 & 6 & 33\dg 57.0\pr & 26 & 7 & 48\dg 59.5\pr \\
26 & 8 & 64\dg 02.0\pr & 26 & 9 & 79\dg 04.4\pr & 26 & 10 & 94\dg 06.9\pr & 26 & 11 & 109\dg 09.4\pr \\
26 & 12 & 124\dg 11.8\pr & 26 & 13 & 139\dg 14.3\pr & 26 & 14 & 154\dg 16.7\pr & 26 & 15 & 169\dg 19.2\pr \\
26 & 16 & 184\dg 21.7\pr & 26 & 17 & 199\dg 24.1\pr & 26 & 18 & 214\dg 26.6\pr & 26 & 19 & 229\dg 29.1\pr \\
26 & 20 & 244\dg 31.5\pr & 26 & 21 & 259\dg 34.0\pr & 26 & 22 & 274\dg 36.5\pr & 26 & 23 & 289\dg 38.9\pr \\
27 & 0 & 304\dg 41.4\pr & 27 & 1 & 319\dg 43.9\pr & 27 & 2 & 334\dg 46.3\pr & 27 & 3 & 349\dg 48.8\pr \\
27 & 4 & 4\dg 51.2\pr & 27 & 5 & 19\dg 53.7\pr & 27 & 6 & 34\dg 56.2\pr & 27 & 7 & 49\dg 58.6\pr \\
27 & 8 & 65\dg 01.1\pr & 27 & 9 & 80\dg 03.6\pr & 27 & 10 & 95\dg 06.0\pr & 27 & 11 & 110\dg 08.5\pr \\
27 & 12 & 125\dg 11.0\pr & 27 & 13 & 140\dg 13.4\pr & 27 & 14 & 155\dg 15.9\pr & 27 & 15 & 170\dg 18.4\pr \\
27 & 16 & 185\dg 20.8\pr & 27 & 17 & 200\dg 23.3\pr & 27 & 18 & 215\dg 25.7\pr & 27 & 19 & 230\dg 28.2\pr \\
27 & 20 & 245\dg 30.7\pr & 27 & 21 & 260\dg 33.1\pr & 27 & 22 & 275\dg 35.6\pr & 27 & 23 & 290\dg 38.1\pr \\
28 & 0 & 305\dg 40.5\pr & 28 & 1 & 320\dg 43.0\pr & 28 & 2 & 335\dg 45.5\pr & 28 & 3 & 350\dg 47.9\pr \\
28 & 4 & 5\dg 50.4\pr & 28 & 5 & 20\dg 52.8\pr & 28 & 6 & 35\dg 55.3\pr & 28 & 7 & 50\dg 57.8\pr \\
28 & 8 & 66\dg 00.2\pr & 28 & 9 & 81\dg 02.7\pr & 28 & 10 & 96\dg 05.2\pr & 28 & 11 & 111\dg 07.6\pr \\
28 & 12 & 126\dg 10.1\pr & 28 & 13 & 141\dg 12.6\pr & 28 & 14 & 156\dg 15.0\pr & 28 & 15 & 171\dg 17.5\pr \\
28 & 16 & 186\dg 20.0\pr & 28 & 17 & 201\dg 22.4\pr & 28 & 18 & 216\dg 24.9\pr & 28 & 19 & 231\dg 27.3\pr \\
28 & 20 & 246\dg 29.8\pr & 28 & 21 & 261\dg 32.3\pr & 28 & 22 & 276\dg 34.7\pr & 28 & 23 & 291\dg 37.2\pr \\
29 & 0 & 306\dg 39.7\pr & 29 & 1 & 321\dg 42.1\pr & 29 & 2 & 336\dg 44.6\pr & 29 & 3 & 351\dg 47.1\pr \\
29 & 4 & 6\dg 49.5\pr & 29 & 5 & 21\dg 52.0\pr & 29 & 6 & 36\dg 54.5\pr & 29 & 7 & 51\dg 56.9\pr \\
29 & 8 & 66\dg 59.4\pr & 29 & 9 & 82\dg 01.8\pr & 29 & 10 & 97\dg 04.3\pr & 29 & 11 & 112\dg 06.8\pr \\
29 & 12 & 127\dg 09.2\pr & 29 & 13 & 142\dg 11.7\pr & 29 & 14 & 157\dg 14.2\pr & 29 & 15 & 172\dg 16.6\pr \\
29 & 16 & 187\dg 19.1\pr & 29 & 17 & 202\dg 21.6\pr & 29 & 18 & 217\dg 24.0\pr & 29 & 19 & 232\dg 26.5\pr \\
29 & 20 & 247\dg 28.9\pr & 29 & 21 & 262\dg 31.4\pr & 29 & 22 & 277\dg 33.9\pr & 29 & 23 & 292\dg 36.3\pr \\
30 & 0 & 307\dg 38.8\pr & 30 & 1 & 322\dg 41.3\pr & 30 & 2 & 337\dg 43.7\pr & 30 & 3 & 352\dg 46.2\pr \\
30 & 4 & 7\dg 48.7\pr & 30 & 5 & 22\dg 51.1\pr & 30 & 6 & 37\dg 53.6\pr & 30 & 7 & 52\dg 56.1\pr \\
30 & 8 & 67\dg 58.5\pr & 30 & 9 & 83\dg 01.0\pr & 30 & 10 & 98\dg 03.4\pr & 30 & 11 & 113\dg 05.9\pr \\
30 & 12 & 128\dg 08.4\pr & 30 & 13 & 143\dg 10.8\pr & 30 & 14 & 158\dg 13.3\pr & 30 & 15 & 173\dg 15.8\pr \\
30 & 16 & 188\dg 18.2\pr & 30 & 17 & 203\dg 20.7\pr & 30 & 18 & 218\dg 23.2\pr & 30 & 19 & 233\dg 25.6\pr \\
30 & 20 & 248\dg 28.1\pr & 30 & 21 & 263\dg 30.6\pr & 30 & 22 & 278\dg 33.0\pr & 30 & 23 & 293\dg 35.5\pr \\
31 & 0 & 308\dg 37.9\pr & 31 & 1 & 323\dg 40.4\pr & 31 & 2 & 338\dg 42.9\pr & 31 & 3 & 353\dg 45.3\pr \\
31 & 4 & 8\dg 47.8\pr & 31 & 5 & 23\dg 50.3\pr & 31 & 6 & 38\dg 52.7\pr & 31 & 7 & 53\dg 55.2\pr \\
31 & 8 & 68\dg 57.7\pr & 31 & 9 & 84\dg 00.1\pr & 31 & 10 & 99\dg 02.6\pr & 31 & 11 & 114\dg 05.1\pr \\
31 & 12 & 129\dg 07.5\pr & 31 & 13 & 144\dg 10.0\pr & 31 & 14 & 159\dg 12.4\pr & 31 & 15 & 174\dg 14.9\pr \\
31 & 16 & 189\dg 17.4\pr & 31 & 17 & 204\dg 19.8\pr & 31 & 18 & 219\dg 22.3\pr & 31 & 19 & 234\dg 24.8\pr \\
31 & 20 & 249\dg 27.2\pr & 31 & 21 & 264\dg 29.7\pr & 31 & 22 & 279\dg 32.2\pr & 31 & 23 & 294\dg 34.6\pr \\
\end{longtable}}

{\footnotesize
\begin{longtable}{rrr|rrr|rrr|rrr}
\caption{Aries --- GHA of the First Point of Aries, August 2026 (UT)}\label{aries08}\\
\toprule
\textbf{d} & \textbf{h} & \textbf{GHA} $\gamma$ & \textbf{d} & \textbf{h} & \textbf{GHA} $\gamma$ & \textbf{d} & \textbf{h} & \textbf{GHA} $\gamma$ & \textbf{d} & \textbf{h} & \textbf{GHA} $\gamma$ \\
\midrule
\endfirsthead
\multicolumn{12}{c}{\textit{Aries --- GHA of the First Point of Aries, August 2026 (UT) (continued)}}\\
\toprule
\textbf{d} & \textbf{h} & \textbf{GHA} $\gamma$ & \textbf{d} & \textbf{h} & \textbf{GHA} $\gamma$ & \textbf{d} & \textbf{h} & \textbf{GHA} $\gamma$ & \textbf{d} & \textbf{h} & \textbf{GHA} $\gamma$ \\
\midrule
\endhead
\midrule\multicolumn{12}{r}{\textit{continued on next page}}\\
\endfoot
\bottomrule
\endlastfoot
1 & 0 & 309\dg 37.1\pr & 1 & 1 & 324\dg 39.5\pr & 1 & 2 & 339\dg 42.0\pr & 1 & 3 & 354\dg 44.5\pr \\
1 & 4 & 9\dg 46.9\pr & 1 & 5 & 24\dg 49.4\pr & 1 & 6 & 39\dg 51.9\pr & 1 & 7 & 54\dg 54.3\pr \\
1 & 8 & 69\dg 56.8\pr & 1 & 9 & 84\dg 59.3\pr & 1 & 10 & 100\dg 01.7\pr & 1 & 11 & 115\dg 04.2\pr \\
1 & 12 & 130\dg 06.7\pr & 1 & 13 & 145\dg 09.1\pr & 1 & 14 & 160\dg 11.6\pr & 1 & 15 & 175\dg 14.0\pr \\
1 & 16 & 190\dg 16.5\pr & 1 & 17 & 205\dg 19.0\pr & 1 & 18 & 220\dg 21.4\pr & 1 & 19 & 235\dg 23.9\pr \\
1 & 20 & 250\dg 26.4\pr & 1 & 21 & 265\dg 28.8\pr & 1 & 22 & 280\dg 31.3\pr & 1 & 23 & 295\dg 33.8\pr \\
2 & 0 & 310\dg 36.2\pr & 2 & 1 & 325\dg 38.7\pr & 2 & 2 & 340\dg 41.2\pr & 2 & 3 & 355\dg 43.6\pr \\
2 & 4 & 10\dg 46.1\pr & 2 & 5 & 25\dg 48.5\pr & 2 & 6 & 40\dg 51.0\pr & 2 & 7 & 55\dg 53.5\pr \\
2 & 8 & 70\dg 55.9\pr & 2 & 9 & 85\dg 58.4\pr & 2 & 10 & 101\dg 00.9\pr & 2 & 11 & 116\dg 03.3\pr \\
2 & 12 & 131\dg 05.8\pr & 2 & 13 & 146\dg 08.3\pr & 2 & 14 & 161\dg 10.7\pr & 2 & 15 & 176\dg 13.2\pr \\
2 & 16 & 191\dg 15.6\pr & 2 & 17 & 206\dg 18.1\pr & 2 & 18 & 221\dg 20.6\pr & 2 & 19 & 236\dg 23.0\pr \\
2 & 20 & 251\dg 25.5\pr & 2 & 21 & 266\dg 28.0\pr & 2 & 22 & 281\dg 30.4\pr & 2 & 23 & 296\dg 32.9\pr \\
3 & 0 & 311\dg 35.4\pr & 3 & 1 & 326\dg 37.8\pr & 3 & 2 & 341\dg 40.3\pr & 3 & 3 & 356\dg 42.8\pr \\
3 & 4 & 11\dg 45.2\pr & 3 & 5 & 26\dg 47.7\pr & 3 & 6 & 41\dg 50.1\pr & 3 & 7 & 56\dg 52.6\pr \\
3 & 8 & 71\dg 55.1\pr & 3 & 9 & 86\dg 57.5\pr & 3 & 10 & 102\dg 00.0\pr & 3 & 11 & 117\dg 02.5\pr \\
3 & 12 & 132\dg 04.9\pr & 3 & 13 & 147\dg 07.4\pr & 3 & 14 & 162\dg 09.9\pr & 3 & 15 & 177\dg 12.3\pr \\
3 & 16 & 192\dg 14.8\pr & 3 & 17 & 207\dg 17.3\pr & 3 & 18 & 222\dg 19.7\pr & 3 & 19 & 237\dg 22.2\pr \\
3 & 20 & 252\dg 24.6\pr & 3 & 21 & 267\dg 27.1\pr & 3 & 22 & 282\dg 29.6\pr & 3 & 23 & 297\dg 32.0\pr \\
4 & 0 & 312\dg 34.5\pr & 4 & 1 & 327\dg 37.0\pr & 4 & 2 & 342\dg 39.4\pr & 4 & 3 & 357\dg 41.9\pr \\
4 & 4 & 12\dg 44.4\pr & 4 & 5 & 27\dg 46.8\pr & 4 & 6 & 42\dg 49.3\pr & 4 & 7 & 57\dg 51.7\pr \\
4 & 8 & 72\dg 54.2\pr & 4 & 9 & 87\dg 56.7\pr & 4 & 10 & 102\dg 59.1\pr & 4 & 11 & 118\dg 01.6\pr \\
4 & 12 & 133\dg 04.1\pr & 4 & 13 & 148\dg 06.5\pr & 4 & 14 & 163\dg 09.0\pr & 4 & 15 & 178\dg 11.5\pr \\
4 & 16 & 193\dg 13.9\pr & 4 & 17 & 208\dg 16.4\pr & 4 & 18 & 223\dg 18.9\pr & 4 & 19 & 238\dg 21.3\pr \\
4 & 20 & 253\dg 23.8\pr & 4 & 21 & 268\dg 26.2\pr & 4 & 22 & 283\dg 28.7\pr & 4 & 23 & 298\dg 31.2\pr \\
5 & 0 & 313\dg 33.6\pr & 5 & 1 & 328\dg 36.1\pr & 5 & 2 & 343\dg 38.6\pr & 5 & 3 & 358\dg 41.0\pr \\
5 & 4 & 13\dg 43.5\pr & 5 & 5 & 28\dg 46.0\pr & 5 & 6 & 43\dg 48.4\pr & 5 & 7 & 58\dg 50.9\pr \\
5 & 8 & 73\dg 53.4\pr & 5 & 9 & 88\dg 55.8\pr & 5 & 10 & 103\dg 58.3\pr & 5 & 11 & 119\dg 00.7\pr \\
5 & 12 & 134\dg 03.2\pr & 5 & 13 & 149\dg 05.7\pr & 5 & 14 & 164\dg 08.1\pr & 5 & 15 & 179\dg 10.6\pr \\
5 & 16 & 194\dg 13.1\pr & 5 & 17 & 209\dg 15.5\pr & 5 & 18 & 224\dg 18.0\pr & 5 & 19 & 239\dg 20.5\pr \\
5 & 20 & 254\dg 22.9\pr & 5 & 21 & 269\dg 25.4\pr & 5 & 22 & 284\dg 27.8\pr & 5 & 23 & 299\dg 30.3\pr \\
6 & 0 & 314\dg 32.8\pr & 6 & 1 & 329\dg 35.2\pr & 6 & 2 & 344\dg 37.7\pr & 6 & 3 & 359\dg 40.2\pr \\
6 & 4 & 14\dg 42.6\pr & 6 & 5 & 29\dg 45.1\pr & 6 & 6 & 44\dg 47.6\pr & 6 & 7 & 59\dg 50.0\pr \\
6 & 8 & 74\dg 52.5\pr & 6 & 9 & 89\dg 55.0\pr & 6 & 10 & 104\dg 57.4\pr & 6 & 11 & 119\dg 59.9\pr \\
6 & 12 & 135\dg 02.3\pr & 6 & 13 & 150\dg 04.8\pr & 6 & 14 & 165\dg 07.3\pr & 6 & 15 & 180\dg 09.7\pr \\
6 & 16 & 195\dg 12.2\pr & 6 & 17 & 210\dg 14.7\pr & 6 & 18 & 225\dg 17.1\pr & 6 & 19 & 240\dg 19.6\pr \\
6 & 20 & 255\dg 22.1\pr & 6 & 21 & 270\dg 24.5\pr & 6 & 22 & 285\dg 27.0\pr & 6 & 23 & 300\dg 29.5\pr \\
7 & 0 & 315\dg 31.9\pr & 7 & 1 & 330\dg 34.4\pr & 7 & 2 & 345\dg 36.8\pr & 7 & 3 & 0\dg 39.3\pr \\
7 & 4 & 15\dg 41.8\pr & 7 & 5 & 30\dg 44.2\pr & 7 & 6 & 45\dg 46.7\pr & 7 & 7 & 60\dg 49.2\pr \\
7 & 8 & 75\dg 51.6\pr & 7 & 9 & 90\dg 54.1\pr & 7 & 10 & 105\dg 56.6\pr & 7 & 11 & 120\dg 59.0\pr \\
7 & 12 & 136\dg 01.5\pr & 7 & 13 & 151\dg 04.0\pr & 7 & 14 & 166\dg 06.4\pr & 7 & 15 & 181\dg 08.9\pr \\
7 & 16 & 196\dg 11.3\pr & 7 & 17 & 211\dg 13.8\pr & 7 & 18 & 226\dg 16.3\pr & 7 & 19 & 241\dg 18.7\pr \\
7 & 20 & 256\dg 21.2\pr & 7 & 21 & 271\dg 23.7\pr & 7 & 22 & 286\dg 26.1\pr & 7 & 23 & 301\dg 28.6\pr \\
8 & 0 & 316\dg 31.1\pr & 8 & 1 & 331\dg 33.5\pr & 8 & 2 & 346\dg 36.0\pr & 8 & 3 & 1\dg 38.4\pr \\
8 & 4 & 16\dg 40.9\pr & 8 & 5 & 31\dg 43.4\pr & 8 & 6 & 46\dg 45.8\pr & 8 & 7 & 61\dg 48.3\pr \\
8 & 8 & 76\dg 50.8\pr & 8 & 9 & 91\dg 53.2\pr & 8 & 10 & 106\dg 55.7\pr & 8 & 11 & 121\dg 58.2\pr \\
8 & 12 & 137\dg 00.6\pr & 8 & 13 & 152\dg 03.1\pr & 8 & 14 & 167\dg 05.6\pr & 8 & 15 & 182\dg 08.0\pr \\
8 & 16 & 197\dg 10.5\pr & 8 & 17 & 212\dg 12.9\pr & 8 & 18 & 227\dg 15.4\pr & 8 & 19 & 242\dg 17.9\pr \\
8 & 20 & 257\dg 20.3\pr & 8 & 21 & 272\dg 22.8\pr & 8 & 22 & 287\dg 25.3\pr & 8 & 23 & 302\dg 27.7\pr \\
9 & 0 & 317\dg 30.2\pr & 9 & 1 & 332\dg 32.7\pr & 9 & 2 & 347\dg 35.1\pr & 9 & 3 & 2\dg 37.6\pr \\
9 & 4 & 17\dg 40.1\pr & 9 & 5 & 32\dg 42.5\pr & 9 & 6 & 47\dg 45.0\pr & 9 & 7 & 62\dg 47.4\pr \\
9 & 8 & 77\dg 49.9\pr & 9 & 9 & 92\dg 52.4\pr & 9 & 10 & 107\dg 54.8\pr & 9 & 11 & 122\dg 57.3\pr \\
9 & 12 & 137\dg 59.8\pr & 9 & 13 & 153\dg 02.2\pr & 9 & 14 & 168\dg 04.7\pr & 9 & 15 & 183\dg 07.2\pr \\
9 & 16 & 198\dg 09.6\pr & 9 & 17 & 213\dg 12.1\pr & 9 & 18 & 228\dg 14.5\pr & 9 & 19 & 243\dg 17.0\pr \\
9 & 20 & 258\dg 19.5\pr & 9 & 21 & 273\dg 21.9\pr & 9 & 22 & 288\dg 24.4\pr & 9 & 23 & 303\dg 26.9\pr \\
10 & 0 & 318\dg 29.3\pr & 10 & 1 & 333\dg 31.8\pr & 10 & 2 & 348\dg 34.3\pr & 10 & 3 & 3\dg 36.7\pr \\
10 & 4 & 18\dg 39.2\pr & 10 & 5 & 33\dg 41.7\pr & 10 & 6 & 48\dg 44.1\pr & 10 & 7 & 63\dg 46.6\pr \\
10 & 8 & 78\dg 49.0\pr & 10 & 9 & 93\dg 51.5\pr & 10 & 10 & 108\dg 54.0\pr & 10 & 11 & 123\dg 56.4\pr \\
10 & 12 & 138\dg 58.9\pr & 10 & 13 & 154\dg 01.4\pr & 10 & 14 & 169\dg 03.8\pr & 10 & 15 & 184\dg 06.3\pr \\
10 & 16 & 199\dg 08.8\pr & 10 & 17 & 214\dg 11.2\pr & 10 & 18 & 229\dg 13.7\pr & 10 & 19 & 244\dg 16.2\pr \\
10 & 20 & 259\dg 18.6\pr & 10 & 21 & 274\dg 21.1\pr & 10 & 22 & 289\dg 23.5\pr & 10 & 23 & 304\dg 26.0\pr \\
11 & 0 & 319\dg 28.5\pr & 11 & 1 & 334\dg 30.9\pr & 11 & 2 & 349\dg 33.4\pr & 11 & 3 & 4\dg 35.9\pr \\
11 & 4 & 19\dg 38.3\pr & 11 & 5 & 34\dg 40.8\pr & 11 & 6 & 49\dg 43.3\pr & 11 & 7 & 64\dg 45.7\pr \\
11 & 8 & 79\dg 48.2\pr & 11 & 9 & 94\dg 50.6\pr & 11 & 10 & 109\dg 53.1\pr & 11 & 11 & 124\dg 55.6\pr \\
11 & 12 & 139\dg 58.0\pr & 11 & 13 & 155\dg 00.5\pr & 11 & 14 & 170\dg 03.0\pr & 11 & 15 & 185\dg 05.4\pr \\
11 & 16 & 200\dg 07.9\pr & 11 & 17 & 215\dg 10.4\pr & 11 & 18 & 230\dg 12.8\pr & 11 & 19 & 245\dg 15.3\pr \\
11 & 20 & 260\dg 17.8\pr & 11 & 21 & 275\dg 20.2\pr & 11 & 22 & 290\dg 22.7\pr & 11 & 23 & 305\dg 25.1\pr \\
12 & 0 & 320\dg 27.6\pr & 12 & 1 & 335\dg 30.1\pr & 12 & 2 & 350\dg 32.5\pr & 12 & 3 & 5\dg 35.0\pr \\
12 & 4 & 20\dg 37.5\pr & 12 & 5 & 35\dg 39.9\pr & 12 & 6 & 50\dg 42.4\pr & 12 & 7 & 65\dg 44.9\pr \\
12 & 8 & 80\dg 47.3\pr & 12 & 9 & 95\dg 49.8\pr & 12 & 10 & 110\dg 52.3\pr & 12 & 11 & 125\dg 54.7\pr \\
12 & 12 & 140\dg 57.2\pr & 12 & 13 & 155\dg 59.6\pr & 12 & 14 & 171\dg 02.1\pr & 12 & 15 & 186\dg 04.6\pr \\
12 & 16 & 201\dg 07.0\pr & 12 & 17 & 216\dg 09.5\pr & 12 & 18 & 231\dg 12.0\pr & 12 & 19 & 246\dg 14.4\pr \\
12 & 20 & 261\dg 16.9\pr & 12 & 21 & 276\dg 19.4\pr & 12 & 22 & 291\dg 21.8\pr & 12 & 23 & 306\dg 24.3\pr \\
13 & 0 & 321\dg 26.7\pr & 13 & 1 & 336\dg 29.2\pr & 13 & 2 & 351\dg 31.7\pr & 13 & 3 & 6\dg 34.1\pr \\
13 & 4 & 21\dg 36.6\pr & 13 & 5 & 36\dg 39.1\pr & 13 & 6 & 51\dg 41.5\pr & 13 & 7 & 66\dg 44.0\pr \\
13 & 8 & 81\dg 46.5\pr & 13 & 9 & 96\dg 48.9\pr & 13 & 10 & 111\dg 51.4\pr & 13 & 11 & 126\dg 53.9\pr \\
13 & 12 & 141\dg 56.3\pr & 13 & 13 & 156\dg 58.8\pr & 13 & 14 & 172\dg 01.2\pr & 13 & 15 & 187\dg 03.7\pr \\
13 & 16 & 202\dg 06.2\pr & 13 & 17 & 217\dg 08.6\pr & 13 & 18 & 232\dg 11.1\pr & 13 & 19 & 247\dg 13.6\pr \\
13 & 20 & 262\dg 16.0\pr & 13 & 21 & 277\dg 18.5\pr & 13 & 22 & 292\dg 21.0\pr & 13 & 23 & 307\dg 23.4\pr \\
14 & 0 & 322\dg 25.9\pr & 14 & 1 & 337\dg 28.4\pr & 14 & 2 & 352\dg 30.8\pr & 14 & 3 & 7\dg 33.3\pr \\
14 & 4 & 22\dg 35.7\pr & 14 & 5 & 37\dg 38.2\pr & 14 & 6 & 52\dg 40.7\pr & 14 & 7 & 67\dg 43.1\pr \\
14 & 8 & 82\dg 45.6\pr & 14 & 9 & 97\dg 48.1\pr & 14 & 10 & 112\dg 50.5\pr & 14 & 11 & 127\dg 53.0\pr \\
14 & 12 & 142\dg 55.5\pr & 14 & 13 & 157\dg 57.9\pr & 14 & 14 & 173\dg 00.4\pr & 14 & 15 & 188\dg 02.9\pr \\
14 & 16 & 203\dg 05.3\pr & 14 & 17 & 218\dg 07.8\pr & 14 & 18 & 233\dg 10.2\pr & 14 & 19 & 248\dg 12.7\pr \\
14 & 20 & 263\dg 15.2\pr & 14 & 21 & 278\dg 17.6\pr & 14 & 22 & 293\dg 20.1\pr & 14 & 23 & 308\dg 22.6\pr \\
15 & 0 & 323\dg 25.0\pr & 15 & 1 & 338\dg 27.5\pr & 15 & 2 & 353\dg 30.0\pr & 15 & 3 & 8\dg 32.4\pr \\
15 & 4 & 23\dg 34.9\pr & 15 & 5 & 38\dg 37.3\pr & 15 & 6 & 53\dg 39.8\pr & 15 & 7 & 68\dg 42.3\pr \\
15 & 8 & 83\dg 44.7\pr & 15 & 9 & 98\dg 47.2\pr & 15 & 10 & 113\dg 49.7\pr & 15 & 11 & 128\dg 52.1\pr \\
15 & 12 & 143\dg 54.6\pr & 15 & 13 & 158\dg 57.1\pr & 15 & 14 & 173\dg 59.5\pr & 15 & 15 & 189\dg 02.0\pr \\
15 & 16 & 204\dg 04.5\pr & 15 & 17 & 219\dg 06.9\pr & 15 & 18 & 234\dg 09.4\pr & 15 & 19 & 249\dg 11.8\pr \\
15 & 20 & 264\dg 14.3\pr & 15 & 21 & 279\dg 16.8\pr & 15 & 22 & 294\dg 19.2\pr & 15 & 23 & 309\dg 21.7\pr \\
16 & 0 & 324\dg 24.2\pr & 16 & 1 & 339\dg 26.6\pr & 16 & 2 & 354\dg 29.1\pr & 16 & 3 & 9\dg 31.6\pr \\
16 & 4 & 24\dg 34.0\pr & 16 & 5 & 39\dg 36.5\pr & 16 & 6 & 54\dg 39.0\pr & 16 & 7 & 69\dg 41.4\pr \\
16 & 8 & 84\dg 43.9\pr & 16 & 9 & 99\dg 46.3\pr & 16 & 10 & 114\dg 48.8\pr & 16 & 11 & 129\dg 51.3\pr \\
16 & 12 & 144\dg 53.7\pr & 16 & 13 & 159\dg 56.2\pr & 16 & 14 & 174\dg 58.7\pr & 16 & 15 & 190\dg 01.1\pr \\
16 & 16 & 205\dg 03.6\pr & 16 & 17 & 220\dg 06.1\pr & 16 & 18 & 235\dg 08.5\pr & 16 & 19 & 250\dg 11.0\pr \\
16 & 20 & 265\dg 13.4\pr & 16 & 21 & 280\dg 15.9\pr & 16 & 22 & 295\dg 18.4\pr & 16 & 23 & 310\dg 20.8\pr \\
17 & 0 & 325\dg 23.3\pr & 17 & 1 & 340\dg 25.8\pr & 17 & 2 & 355\dg 28.2\pr & 17 & 3 & 10\dg 30.7\pr \\
17 & 4 & 25\dg 33.2\pr & 17 & 5 & 40\dg 35.6\pr & 17 & 6 & 55\dg 38.1\pr & 17 & 7 & 70\dg 40.6\pr \\
17 & 8 & 85\dg 43.0\pr & 17 & 9 & 100\dg 45.5\pr & 17 & 10 & 115\dg 47.9\pr & 17 & 11 & 130\dg 50.4\pr \\
17 & 12 & 145\dg 52.9\pr & 17 & 13 & 160\dg 55.3\pr & 17 & 14 & 175\dg 57.8\pr & 17 & 15 & 191\dg 00.3\pr \\
17 & 16 & 206\dg 02.7\pr & 17 & 17 & 221\dg 05.2\pr & 17 & 18 & 236\dg 07.7\pr & 17 & 19 & 251\dg 10.1\pr \\
17 & 20 & 266\dg 12.6\pr & 17 & 21 & 281\dg 15.1\pr & 17 & 22 & 296\dg 17.5\pr & 17 & 23 & 311\dg 20.0\pr \\
18 & 0 & 326\dg 22.4\pr & 18 & 1 & 341\dg 24.9\pr & 18 & 2 & 356\dg 27.4\pr & 18 & 3 & 11\dg 29.8\pr \\
18 & 4 & 26\dg 32.3\pr & 18 & 5 & 41\dg 34.8\pr & 18 & 6 & 56\dg 37.2\pr & 18 & 7 & 71\dg 39.7\pr \\
18 & 8 & 86\dg 42.2\pr & 18 & 9 & 101\dg 44.6\pr & 18 & 10 & 116\dg 47.1\pr & 18 & 11 & 131\dg 49.5\pr \\
18 & 12 & 146\dg 52.0\pr & 18 & 13 & 161\dg 54.5\pr & 18 & 14 & 176\dg 56.9\pr & 18 & 15 & 191\dg 59.4\pr \\
18 & 16 & 207\dg 01.9\pr & 18 & 17 & 222\dg 04.3\pr & 18 & 18 & 237\dg 06.8\pr & 18 & 19 & 252\dg 09.3\pr \\
18 & 20 & 267\dg 11.7\pr & 18 & 21 & 282\dg 14.2\pr & 18 & 22 & 297\dg 16.7\pr & 18 & 23 & 312\dg 19.1\pr \\
19 & 0 & 327\dg 21.6\pr & 19 & 1 & 342\dg 24.0\pr & 19 & 2 & 357\dg 26.5\pr & 19 & 3 & 12\dg 29.0\pr \\
19 & 4 & 27\dg 31.4\pr & 19 & 5 & 42\dg 33.9\pr & 19 & 6 & 57\dg 36.4\pr & 19 & 7 & 72\dg 38.8\pr \\
19 & 8 & 87\dg 41.3\pr & 19 & 9 & 102\dg 43.8\pr & 19 & 10 & 117\dg 46.2\pr & 19 & 11 & 132\dg 48.7\pr \\
19 & 12 & 147\dg 51.2\pr & 19 & 13 & 162\dg 53.6\pr & 19 & 14 & 177\dg 56.1\pr & 19 & 15 & 192\dg 58.5\pr \\
19 & 16 & 208\dg 01.0\pr & 19 & 17 & 223\dg 03.5\pr & 19 & 18 & 238\dg 05.9\pr & 19 & 19 & 253\dg 08.4\pr \\
19 & 20 & 268\dg 10.9\pr & 19 & 21 & 283\dg 13.3\pr & 19 & 22 & 298\dg 15.8\pr & 19 & 23 & 313\dg 18.3\pr \\
20 & 0 & 328\dg 20.7\pr & 20 & 1 & 343\dg 23.2\pr & 20 & 2 & 358\dg 25.6\pr & 20 & 3 & 13\dg 28.1\pr \\
20 & 4 & 28\dg 30.6\pr & 20 & 5 & 43\dg 33.0\pr & 20 & 6 & 58\dg 35.5\pr & 20 & 7 & 73\dg 38.0\pr \\
20 & 8 & 88\dg 40.4\pr & 20 & 9 & 103\dg 42.9\pr & 20 & 10 & 118\dg 45.4\pr & 20 & 11 & 133\dg 47.8\pr \\
20 & 12 & 148\dg 50.3\pr & 20 & 13 & 163\dg 52.8\pr & 20 & 14 & 178\dg 55.2\pr & 20 & 15 & 193\dg 57.7\pr \\
20 & 16 & 209\dg 00.1\pr & 20 & 17 & 224\dg 02.6\pr & 20 & 18 & 239\dg 05.1\pr & 20 & 19 & 254\dg 07.5\pr \\
20 & 20 & 269\dg 10.0\pr & 20 & 21 & 284\dg 12.5\pr & 20 & 22 & 299\dg 14.9\pr & 20 & 23 & 314\dg 17.4\pr \\
21 & 0 & 329\dg 19.9\pr & 21 & 1 & 344\dg 22.3\pr & 21 & 2 & 359\dg 24.8\pr & 21 & 3 & 14\dg 27.3\pr \\
21 & 4 & 29\dg 29.7\pr & 21 & 5 & 44\dg 32.2\pr & 21 & 6 & 59\dg 34.6\pr & 21 & 7 & 74\dg 37.1\pr \\
21 & 8 & 89\dg 39.6\pr & 21 & 9 & 104\dg 42.0\pr & 21 & 10 & 119\dg 44.5\pr & 21 & 11 & 134\dg 47.0\pr \\
21 & 12 & 149\dg 49.4\pr & 21 & 13 & 164\dg 51.9\pr & 21 & 14 & 179\dg 54.4\pr & 21 & 15 & 194\dg 56.8\pr \\
21 & 16 & 209\dg 59.3\pr & 21 & 17 & 225\dg 01.8\pr & 21 & 18 & 240\dg 04.2\pr & 21 & 19 & 255\dg 06.7\pr \\
21 & 20 & 270\dg 09.1\pr & 21 & 21 & 285\dg 11.6\pr & 21 & 22 & 300\dg 14.1\pr & 21 & 23 & 315\dg 16.5\pr \\
22 & 0 & 330\dg 19.0\pr & 22 & 1 & 345\dg 21.5\pr & 22 & 2 & 0\dg 23.9\pr & 22 & 3 & 15\dg 26.4\pr \\
22 & 4 & 30\dg 28.9\pr & 22 & 5 & 45\dg 31.3\pr & 22 & 6 & 60\dg 33.8\pr & 22 & 7 & 75\dg 36.2\pr \\
22 & 8 & 90\dg 38.7\pr & 22 & 9 & 105\dg 41.2\pr & 22 & 10 & 120\dg 43.6\pr & 22 & 11 & 135\dg 46.1\pr \\
22 & 12 & 150\dg 48.6\pr & 22 & 13 & 165\dg 51.0\pr & 22 & 14 & 180\dg 53.5\pr & 22 & 15 & 195\dg 56.0\pr \\
22 & 16 & 210\dg 58.4\pr & 22 & 17 & 226\dg 00.9\pr & 22 & 18 & 241\dg 03.4\pr & 22 & 19 & 256\dg 05.8\pr \\
22 & 20 & 271\dg 08.3\pr & 22 & 21 & 286\dg 10.7\pr & 22 & 22 & 301\dg 13.2\pr & 22 & 23 & 316\dg 15.7\pr \\
23 & 0 & 331\dg 18.1\pr & 23 & 1 & 346\dg 20.6\pr & 23 & 2 & 1\dg 23.1\pr & 23 & 3 & 16\dg 25.5\pr \\
23 & 4 & 31\dg 28.0\pr & 23 & 5 & 46\dg 30.5\pr & 23 & 6 & 61\dg 32.9\pr & 23 & 7 & 76\dg 35.4\pr \\
23 & 8 & 91\dg 37.9\pr & 23 & 9 & 106\dg 40.3\pr & 23 & 10 & 121\dg 42.8\pr & 23 & 11 & 136\dg 45.2\pr \\
23 & 12 & 151\dg 47.7\pr & 23 & 13 & 166\dg 50.2\pr & 23 & 14 & 181\dg 52.6\pr & 23 & 15 & 196\dg 55.1\pr \\
23 & 16 & 211\dg 57.6\pr & 23 & 17 & 227\dg 00.0\pr & 23 & 18 & 242\dg 02.5\pr & 23 & 19 & 257\dg 05.0\pr \\
23 & 20 & 272\dg 07.4\pr & 23 & 21 & 287\dg 09.9\pr & 23 & 22 & 302\dg 12.3\pr & 23 & 23 & 317\dg 14.8\pr \\
24 & 0 & 332\dg 17.3\pr & 24 & 1 & 347\dg 19.7\pr & 24 & 2 & 2\dg 22.2\pr & 24 & 3 & 17\dg 24.7\pr \\
24 & 4 & 32\dg 27.1\pr & 24 & 5 & 47\dg 29.6\pr & 24 & 6 & 62\dg 32.1\pr & 24 & 7 & 77\dg 34.5\pr \\
24 & 8 & 92\dg 37.0\pr & 24 & 9 & 107\dg 39.5\pr & 24 & 10 & 122\dg 41.9\pr & 24 & 11 & 137\dg 44.4\pr \\
24 & 12 & 152\dg 46.8\pr & 24 & 13 & 167\dg 49.3\pr & 24 & 14 & 182\dg 51.8\pr & 24 & 15 & 197\dg 54.2\pr \\
24 & 16 & 212\dg 56.7\pr & 24 & 17 & 227\dg 59.2\pr & 24 & 18 & 243\dg 01.6\pr & 24 & 19 & 258\dg 04.1\pr \\
24 & 20 & 273\dg 06.6\pr & 24 & 21 & 288\dg 09.0\pr & 24 & 22 & 303\dg 11.5\pr & 24 & 23 & 318\dg 14.0\pr \\
25 & 0 & 333\dg 16.4\pr & 25 & 1 & 348\dg 18.9\pr & 25 & 2 & 3\dg 21.3\pr & 25 & 3 & 18\dg 23.8\pr \\
25 & 4 & 33\dg 26.3\pr & 25 & 5 & 48\dg 28.7\pr & 25 & 6 & 63\dg 31.2\pr & 25 & 7 & 78\dg 33.7\pr \\
25 & 8 & 93\dg 36.1\pr & 25 & 9 & 108\dg 38.6\pr & 25 & 10 & 123\dg 41.1\pr & 25 & 11 & 138\dg 43.5\pr \\
25 & 12 & 153\dg 46.0\pr & 25 & 13 & 168\dg 48.4\pr & 25 & 14 & 183\dg 50.9\pr & 25 & 15 & 198\dg 53.4\pr \\
25 & 16 & 213\dg 55.8\pr & 25 & 17 & 228\dg 58.3\pr & 25 & 18 & 244\dg 00.8\pr & 25 & 19 & 259\dg 03.2\pr \\
25 & 20 & 274\dg 05.7\pr & 25 & 21 & 289\dg 08.2\pr & 25 & 22 & 304\dg 10.6\pr & 25 & 23 & 319\dg 13.1\pr \\
26 & 0 & 334\dg 15.6\pr & 26 & 1 & 349\dg 18.0\pr & 26 & 2 & 4\dg 20.5\pr & 26 & 3 & 19\dg 22.9\pr \\
26 & 4 & 34\dg 25.4\pr & 26 & 5 & 49\dg 27.9\pr & 26 & 6 & 64\dg 30.3\pr & 26 & 7 & 79\dg 32.8\pr \\
26 & 8 & 94\dg 35.3\pr & 26 & 9 & 109\dg 37.7\pr & 26 & 10 & 124\dg 40.2\pr & 26 & 11 & 139\dg 42.7\pr \\
26 & 12 & 154\dg 45.1\pr & 26 & 13 & 169\dg 47.6\pr & 26 & 14 & 184\dg 50.1\pr & 26 & 15 & 199\dg 52.5\pr \\
26 & 16 & 214\dg 55.0\pr & 26 & 17 & 229\dg 57.4\pr & 26 & 18 & 244\dg 59.9\pr & 26 & 19 & 260\dg 02.4\pr \\
26 & 20 & 275\dg 04.8\pr & 26 & 21 & 290\dg 07.3\pr & 26 & 22 & 305\dg 09.8\pr & 26 & 23 & 320\dg 12.2\pr \\
27 & 0 & 335\dg 14.7\pr & 27 & 1 & 350\dg 17.2\pr & 27 & 2 & 5\dg 19.6\pr & 27 & 3 & 20\dg 22.1\pr \\
27 & 4 & 35\dg 24.5\pr & 27 & 5 & 50\dg 27.0\pr & 27 & 6 & 65\dg 29.5\pr & 27 & 7 & 80\dg 31.9\pr \\
27 & 8 & 95\dg 34.4\pr & 27 & 9 & 110\dg 36.9\pr & 27 & 10 & 125\dg 39.3\pr & 27 & 11 & 140\dg 41.8\pr \\
27 & 12 & 155\dg 44.3\pr & 27 & 13 & 170\dg 46.7\pr & 27 & 14 & 185\dg 49.2\pr & 27 & 15 & 200\dg 51.7\pr \\
27 & 16 & 215\dg 54.1\pr & 27 & 17 & 230\dg 56.6\pr & 27 & 18 & 245\dg 59.0\pr & 27 & 19 & 261\dg 01.5\pr \\
27 & 20 & 276\dg 04.0\pr & 27 & 21 & 291\dg 06.4\pr & 27 & 22 & 306\dg 08.9\pr & 27 & 23 & 321\dg 11.4\pr \\
28 & 0 & 336\dg 13.8\pr & 28 & 1 & 351\dg 16.3\pr & 28 & 2 & 6\dg 18.8\pr & 28 & 3 & 21\dg 21.2\pr \\
28 & 4 & 36\dg 23.7\pr & 28 & 5 & 51\dg 26.2\pr & 28 & 6 & 66\dg 28.6\pr & 28 & 7 & 81\dg 31.1\pr \\
28 & 8 & 96\dg 33.5\pr & 28 & 9 & 111\dg 36.0\pr & 28 & 10 & 126\dg 38.5\pr & 28 & 11 & 141\dg 40.9\pr \\
28 & 12 & 156\dg 43.4\pr & 28 & 13 & 171\dg 45.9\pr & 28 & 14 & 186\dg 48.3\pr & 28 & 15 & 201\dg 50.8\pr \\
28 & 16 & 216\dg 53.3\pr & 28 & 17 & 231\dg 55.7\pr & 28 & 18 & 246\dg 58.2\pr & 28 & 19 & 262\dg 00.7\pr \\
28 & 20 & 277\dg 03.1\pr & 28 & 21 & 292\dg 05.6\pr & 28 & 22 & 307\dg 08.0\pr & 28 & 23 & 322\dg 10.5\pr \\
29 & 0 & 337\dg 13.0\pr & 29 & 1 & 352\dg 15.4\pr & 29 & 2 & 7\dg 17.9\pr & 29 & 3 & 22\dg 20.4\pr \\
29 & 4 & 37\dg 22.8\pr & 29 & 5 & 52\dg 25.3\pr & 29 & 6 & 67\dg 27.8\pr & 29 & 7 & 82\dg 30.2\pr \\
29 & 8 & 97\dg 32.7\pr & 29 & 9 & 112\dg 35.1\pr & 29 & 10 & 127\dg 37.6\pr & 29 & 11 & 142\dg 40.1\pr \\
29 & 12 & 157\dg 42.5\pr & 29 & 13 & 172\dg 45.0\pr & 29 & 14 & 187\dg 47.5\pr & 29 & 15 & 202\dg 49.9\pr \\
29 & 16 & 217\dg 52.4\pr & 29 & 17 & 232\dg 54.9\pr & 29 & 18 & 247\dg 57.3\pr & 29 & 19 & 262\dg 59.8\pr \\
29 & 20 & 278\dg 02.3\pr & 29 & 21 & 293\dg 04.7\pr & 29 & 22 & 308\dg 07.2\pr & 29 & 23 & 323\dg 09.6\pr \\
30 & 0 & 338\dg 12.1\pr & 30 & 1 & 353\dg 14.6\pr & 30 & 2 & 8\dg 17.0\pr & 30 & 3 & 23\dg 19.5\pr \\
30 & 4 & 38\dg 22.0\pr & 30 & 5 & 53\dg 24.4\pr & 30 & 6 & 68\dg 26.9\pr & 30 & 7 & 83\dg 29.4\pr \\
30 & 8 & 98\dg 31.8\pr & 30 & 9 & 113\dg 34.3\pr & 30 & 10 & 128\dg 36.8\pr & 30 & 11 & 143\dg 39.2\pr \\
30 & 12 & 158\dg 41.7\pr & 30 & 13 & 173\dg 44.1\pr & 30 & 14 & 188\dg 46.6\pr & 30 & 15 & 203\dg 49.1\pr \\
30 & 16 & 218\dg 51.5\pr & 30 & 17 & 233\dg 54.0\pr & 30 & 18 & 248\dg 56.5\pr & 30 & 19 & 263\dg 58.9\pr \\
30 & 20 & 279\dg 01.4\pr & 30 & 21 & 294\dg 03.9\pr & 30 & 22 & 309\dg 06.3\pr & 30 & 23 & 324\dg 08.8\pr \\
31 & 0 & 339\dg 11.2\pr & 31 & 1 & 354\dg 13.7\pr & 31 & 2 & 9\dg 16.2\pr & 31 & 3 & 24\dg 18.6\pr \\
31 & 4 & 39\dg 21.1\pr & 31 & 5 & 54\dg 23.6\pr & 31 & 6 & 69\dg 26.0\pr & 31 & 7 & 84\dg 28.5\pr \\
31 & 8 & 99\dg 31.0\pr & 31 & 9 & 114\dg 33.4\pr & 31 & 10 & 129\dg 35.9\pr & 31 & 11 & 144\dg 38.4\pr \\
31 & 12 & 159\dg 40.8\pr & 31 & 13 & 174\dg 43.3\pr & 31 & 14 & 189\dg 45.7\pr & 31 & 15 & 204\dg 48.2\pr \\
31 & 16 & 219\dg 50.7\pr & 31 & 17 & 234\dg 53.1\pr & 31 & 18 & 249\dg 55.6\pr & 31 & 19 & 264\dg 58.1\pr \\
31 & 20 & 280\dg 00.5\pr & 31 & 21 & 295\dg 03.0\pr & 31 & 22 & 310\dg 05.5\pr & 31 & 23 & 325\dg 07.9\pr \\
\end{longtable}}

{\footnotesize
\begin{longtable}{rrr|rrr|rrr|rrr}
\caption{Aries --- GHA of the First Point of Aries, September 2026 (UT)}\label{aries09}\\
\toprule
\textbf{d} & \textbf{h} & \textbf{GHA} $\gamma$ & \textbf{d} & \textbf{h} & \textbf{GHA} $\gamma$ & \textbf{d} & \textbf{h} & \textbf{GHA} $\gamma$ & \textbf{d} & \textbf{h} & \textbf{GHA} $\gamma$ \\
\midrule
\endfirsthead
\multicolumn{12}{c}{\textit{Aries --- GHA of the First Point of Aries, September 2026 (UT) (continued)}}\\
\toprule
\textbf{d} & \textbf{h} & \textbf{GHA} $\gamma$ & \textbf{d} & \textbf{h} & \textbf{GHA} $\gamma$ & \textbf{d} & \textbf{h} & \textbf{GHA} $\gamma$ & \textbf{d} & \textbf{h} & \textbf{GHA} $\gamma$ \\
\midrule
\endhead
\midrule\multicolumn{12}{r}{\textit{continued on next page}}\\
\endfoot
\bottomrule
\endlastfoot
1 & 0 & 340\dg 10.4\pr & 1 & 1 & 355\dg 12.9\pr & 1 & 2 & 10\dg 15.3\pr & 1 & 3 & 25\dg 17.8\pr \\
1 & 4 & 40\dg 20.2\pr & 1 & 5 & 55\dg 22.7\pr & 1 & 6 & 70\dg 25.2\pr & 1 & 7 & 85\dg 27.6\pr \\
1 & 8 & 100\dg 30.1\pr & 1 & 9 & 115\dg 32.6\pr & 1 & 10 & 130\dg 35.0\pr & 1 & 11 & 145\dg 37.5\pr \\
1 & 12 & 160\dg 40.0\pr & 1 & 13 & 175\dg 42.4\pr & 1 & 14 & 190\dg 44.9\pr & 1 & 15 & 205\dg 47.3\pr \\
1 & 16 & 220\dg 49.8\pr & 1 & 17 & 235\dg 52.3\pr & 1 & 18 & 250\dg 54.7\pr & 1 & 19 & 265\dg 57.2\pr \\
1 & 20 & 280\dg 59.7\pr & 1 & 21 & 296\dg 02.1\pr & 1 & 22 & 311\dg 04.6\pr & 1 & 23 & 326\dg 07.1\pr \\
2 & 0 & 341\dg 09.5\pr & 2 & 1 & 356\dg 12.0\pr & 2 & 2 & 11\dg 14.5\pr & 2 & 3 & 26\dg 16.9\pr \\
2 & 4 & 41\dg 19.4\pr & 2 & 5 & 56\dg 21.8\pr & 2 & 6 & 71\dg 24.3\pr & 2 & 7 & 86\dg 26.8\pr \\
2 & 8 & 101\dg 29.2\pr & 2 & 9 & 116\dg 31.7\pr & 2 & 10 & 131\dg 34.2\pr & 2 & 11 & 146\dg 36.6\pr \\
2 & 12 & 161\dg 39.1\pr & 2 & 13 & 176\dg 41.6\pr & 2 & 14 & 191\dg 44.0\pr & 2 & 15 & 206\dg 46.5\pr \\
2 & 16 & 221\dg 49.0\pr & 2 & 17 & 236\dg 51.4\pr & 2 & 18 & 251\dg 53.9\pr & 2 & 19 & 266\dg 56.3\pr \\
2 & 20 & 281\dg 58.8\pr & 2 & 21 & 297\dg 01.3\pr & 2 & 22 & 312\dg 03.7\pr & 2 & 23 & 327\dg 06.2\pr \\
3 & 0 & 342\dg 08.7\pr & 3 & 1 & 357\dg 11.1\pr & 3 & 2 & 12\dg 13.6\pr & 3 & 3 & 27\dg 16.1\pr \\
3 & 4 & 42\dg 18.5\pr & 3 & 5 & 57\dg 21.0\pr & 3 & 6 & 72\dg 23.5\pr & 3 & 7 & 87\dg 25.9\pr \\
3 & 8 & 102\dg 28.4\pr & 3 & 9 & 117\dg 30.8\pr & 3 & 10 & 132\dg 33.3\pr & 3 & 11 & 147\dg 35.8\pr \\
3 & 12 & 162\dg 38.2\pr & 3 & 13 & 177\dg 40.7\pr & 3 & 14 & 192\dg 43.2\pr & 3 & 15 & 207\dg 45.6\pr \\
3 & 16 & 222\dg 48.1\pr & 3 & 17 & 237\dg 50.6\pr & 3 & 18 & 252\dg 53.0\pr & 3 & 19 & 267\dg 55.5\pr \\
3 & 20 & 282\dg 57.9\pr & 3 & 21 & 298\dg 00.4\pr & 3 & 22 & 313\dg 02.9\pr & 3 & 23 & 328\dg 05.3\pr \\
4 & 0 & 343\dg 07.8\pr & 4 & 1 & 358\dg 10.3\pr & 4 & 2 & 13\dg 12.7\pr & 4 & 3 & 28\dg 15.2\pr \\
4 & 4 & 43\dg 17.7\pr & 4 & 5 & 58\dg 20.1\pr & 4 & 6 & 73\dg 22.6\pr & 4 & 7 & 88\dg 25.1\pr \\
4 & 8 & 103\dg 27.5\pr & 4 & 9 & 118\dg 30.0\pr & 4 & 10 & 133\dg 32.4\pr & 4 & 11 & 148\dg 34.9\pr \\
4 & 12 & 163\dg 37.4\pr & 4 & 13 & 178\dg 39.8\pr & 4 & 14 & 193\dg 42.3\pr & 4 & 15 & 208\dg 44.8\pr \\
4 & 16 & 223\dg 47.2\pr & 4 & 17 & 238\dg 49.7\pr & 4 & 18 & 253\dg 52.2\pr & 4 & 19 & 268\dg 54.6\pr \\
4 & 20 & 283\dg 57.1\pr & 4 & 21 & 298\dg 59.6\pr & 4 & 22 & 314\dg 02.0\pr & 4 & 23 & 329\dg 04.5\pr \\
5 & 0 & 344\dg 06.9\pr & 5 & 1 & 359\dg 09.4\pr & 5 & 2 & 14\dg 11.9\pr & 5 & 3 & 29\dg 14.3\pr \\
5 & 4 & 44\dg 16.8\pr & 5 & 5 & 59\dg 19.3\pr & 5 & 6 & 74\dg 21.7\pr & 5 & 7 & 89\dg 24.2\pr \\
5 & 8 & 104\dg 26.7\pr & 5 & 9 & 119\dg 29.1\pr & 5 & 10 & 134\dg 31.6\pr & 5 & 11 & 149\dg 34.0\pr \\
5 & 12 & 164\dg 36.5\pr & 5 & 13 & 179\dg 39.0\pr & 5 & 14 & 194\dg 41.4\pr & 5 & 15 & 209\dg 43.9\pr \\
5 & 16 & 224\dg 46.4\pr & 5 & 17 & 239\dg 48.8\pr & 5 & 18 & 254\dg 51.3\pr & 5 & 19 & 269\dg 53.8\pr \\
5 & 20 & 284\dg 56.2\pr & 5 & 21 & 299\dg 58.7\pr & 5 & 22 & 315\dg 01.2\pr & 5 & 23 & 330\dg 03.6\pr \\
6 & 0 & 345\dg 06.1\pr & 6 & 1 & 0\dg 08.5\pr & 6 & 2 & 15\dg 11.0\pr & 6 & 3 & 30\dg 13.5\pr \\
6 & 4 & 45\dg 15.9\pr & 6 & 5 & 60\dg 18.4\pr & 6 & 6 & 75\dg 20.9\pr & 6 & 7 & 90\dg 23.3\pr \\
6 & 8 & 105\dg 25.8\pr & 6 & 9 & 120\dg 28.3\pr & 6 & 10 & 135\dg 30.7\pr & 6 & 11 & 150\dg 33.2\pr \\
6 & 12 & 165\dg 35.7\pr & 6 & 13 & 180\dg 38.1\pr & 6 & 14 & 195\dg 40.6\pr & 6 & 15 & 210\dg 43.0\pr \\
6 & 16 & 225\dg 45.5\pr & 6 & 17 & 240\dg 48.0\pr & 6 & 18 & 255\dg 50.4\pr & 6 & 19 & 270\dg 52.9\pr \\
6 & 20 & 285\dg 55.4\pr & 6 & 21 & 300\dg 57.8\pr & 6 & 22 & 316\dg 00.3\pr & 6 & 23 & 331\dg 02.8\pr \\
7 & 0 & 346\dg 05.2\pr & 7 & 1 & 1\dg 07.7\pr & 7 & 2 & 16\dg 10.1\pr & 7 & 3 & 31\dg 12.6\pr \\
7 & 4 & 46\dg 15.1\pr & 7 & 5 & 61\dg 17.5\pr & 7 & 6 & 76\dg 20.0\pr & 7 & 7 & 91\dg 22.5\pr \\
7 & 8 & 106\dg 24.9\pr & 7 & 9 & 121\dg 27.4\pr & 7 & 10 & 136\dg 29.9\pr & 7 & 11 & 151\dg 32.3\pr \\
7 & 12 & 166\dg 34.8\pr & 7 & 13 & 181\dg 37.3\pr & 7 & 14 & 196\dg 39.7\pr & 7 & 15 & 211\dg 42.2\pr \\
7 & 16 & 226\dg 44.6\pr & 7 & 17 & 241\dg 47.1\pr & 7 & 18 & 256\dg 49.6\pr & 7 & 19 & 271\dg 52.0\pr \\
7 & 20 & 286\dg 54.5\pr & 7 & 21 & 301\dg 57.0\pr & 7 & 22 & 316\dg 59.4\pr & 7 & 23 & 332\dg 01.9\pr \\
8 & 0 & 347\dg 04.4\pr & 8 & 1 & 2\dg 06.8\pr & 8 & 2 & 17\dg 09.3\pr & 8 & 3 & 32\dg 11.8\pr \\
8 & 4 & 47\dg 14.2\pr & 8 & 5 & 62\dg 16.7\pr & 8 & 6 & 77\dg 19.1\pr & 8 & 7 & 92\dg 21.6\pr \\
8 & 8 & 107\dg 24.1\pr & 8 & 9 & 122\dg 26.5\pr & 8 & 10 & 137\dg 29.0\pr & 8 & 11 & 152\dg 31.5\pr \\
8 & 12 & 167\dg 33.9\pr & 8 & 13 & 182\dg 36.4\pr & 8 & 14 & 197\dg 38.9\pr & 8 & 15 & 212\dg 41.3\pr \\
8 & 16 & 227\dg 43.8\pr & 8 & 17 & 242\dg 46.2\pr & 8 & 18 & 257\dg 48.7\pr & 8 & 19 & 272\dg 51.2\pr \\
8 & 20 & 287\dg 53.6\pr & 8 & 21 & 302\dg 56.1\pr & 8 & 22 & 317\dg 58.6\pr & 8 & 23 & 333\dg 01.0\pr \\
9 & 0 & 348\dg 03.5\pr & 9 & 1 & 3\dg 06.0\pr & 9 & 2 & 18\dg 08.4\pr & 9 & 3 & 33\dg 10.9\pr \\
9 & 4 & 48\dg 13.4\pr & 9 & 5 & 63\dg 15.8\pr & 9 & 6 & 78\dg 18.3\pr & 9 & 7 & 93\dg 20.7\pr \\
9 & 8 & 108\dg 23.2\pr & 9 & 9 & 123\dg 25.7\pr & 9 & 10 & 138\dg 28.1\pr & 9 & 11 & 153\dg 30.6\pr \\
9 & 12 & 168\dg 33.1\pr & 9 & 13 & 183\dg 35.5\pr & 9 & 14 & 198\dg 38.0\pr & 9 & 15 & 213\dg 40.5\pr \\
9 & 16 & 228\dg 42.9\pr & 9 & 17 & 243\dg 45.4\pr & 9 & 18 & 258\dg 47.9\pr & 9 & 19 & 273\dg 50.3\pr \\
9 & 20 & 288\dg 52.8\pr & 9 & 21 & 303\dg 55.2\pr & 9 & 22 & 318\dg 57.7\pr & 9 & 23 & 334\dg 00.2\pr \\
10 & 0 & 349\dg 02.6\pr & 10 & 1 & 4\dg 05.1\pr & 10 & 2 & 19\dg 07.6\pr & 10 & 3 & 34\dg 10.0\pr \\
10 & 4 & 49\dg 12.5\pr & 10 & 5 & 64\dg 15.0\pr & 10 & 6 & 79\dg 17.4\pr & 10 & 7 & 94\dg 19.9\pr \\
10 & 8 & 109\dg 22.4\pr & 10 & 9 & 124\dg 24.8\pr & 10 & 10 & 139\dg 27.3\pr & 10 & 11 & 154\dg 29.7\pr \\
10 & 12 & 169\dg 32.2\pr & 10 & 13 & 184\dg 34.7\pr & 10 & 14 & 199\dg 37.1\pr & 10 & 15 & 214\dg 39.6\pr \\
10 & 16 & 229\dg 42.1\pr & 10 & 17 & 244\dg 44.5\pr & 10 & 18 & 259\dg 47.0\pr & 10 & 19 & 274\dg 49.5\pr \\
10 & 20 & 289\dg 51.9\pr & 10 & 21 & 304\dg 54.4\pr & 10 & 22 & 319\dg 56.8\pr & 10 & 23 & 334\dg 59.3\pr \\
11 & 0 & 350\dg 01.8\pr & 11 & 1 & 5\dg 04.2\pr & 11 & 2 & 20\dg 06.7\pr & 11 & 3 & 35\dg 09.2\pr \\
11 & 4 & 50\dg 11.6\pr & 11 & 5 & 65\dg 14.1\pr & 11 & 6 & 80\dg 16.6\pr & 11 & 7 & 95\dg 19.0\pr \\
11 & 8 & 110\dg 21.5\pr & 11 & 9 & 125\dg 24.0\pr & 11 & 10 & 140\dg 26.4\pr & 11 & 11 & 155\dg 28.9\pr \\
11 & 12 & 170\dg 31.3\pr & 11 & 13 & 185\dg 33.8\pr & 11 & 14 & 200\dg 36.3\pr & 11 & 15 & 215\dg 38.7\pr \\
11 & 16 & 230\dg 41.2\pr & 11 & 17 & 245\dg 43.7\pr & 11 & 18 & 260\dg 46.1\pr & 11 & 19 & 275\dg 48.6\pr \\
11 & 20 & 290\dg 51.1\pr & 11 & 21 & 305\dg 53.5\pr & 11 & 22 & 320\dg 56.0\pr & 11 & 23 & 335\dg 58.5\pr \\
12 & 0 & 351\dg 00.9\pr & 12 & 1 & 6\dg 03.4\pr & 12 & 2 & 21\dg 05.8\pr & 12 & 3 & 36\dg 08.3\pr \\
12 & 4 & 51\dg 10.8\pr & 12 & 5 & 66\dg 13.2\pr & 12 & 6 & 81\dg 15.7\pr & 12 & 7 & 96\dg 18.2\pr \\
12 & 8 & 111\dg 20.6\pr & 12 & 9 & 126\dg 23.1\pr & 12 & 10 & 141\dg 25.6\pr & 12 & 11 & 156\dg 28.0\pr \\
12 & 12 & 171\dg 30.5\pr & 12 & 13 & 186\dg 32.9\pr & 12 & 14 & 201\dg 35.4\pr & 12 & 15 & 216\dg 37.9\pr \\
12 & 16 & 231\dg 40.3\pr & 12 & 17 & 246\dg 42.8\pr & 12 & 18 & 261\dg 45.3\pr & 12 & 19 & 276\dg 47.7\pr \\
12 & 20 & 291\dg 50.2\pr & 12 & 21 & 306\dg 52.7\pr & 12 & 22 & 321\dg 55.1\pr & 12 & 23 & 336\dg 57.6\pr \\
13 & 0 & 352\dg 00.1\pr & 13 & 1 & 7\dg 02.5\pr & 13 & 2 & 22\dg 05.0\pr & 13 & 3 & 37\dg 07.4\pr \\
13 & 4 & 52\dg 09.9\pr & 13 & 5 & 67\dg 12.4\pr & 13 & 6 & 82\dg 14.8\pr & 13 & 7 & 97\dg 17.3\pr \\
13 & 8 & 112\dg 19.8\pr & 13 & 9 & 127\dg 22.2\pr & 13 & 10 & 142\dg 24.7\pr & 13 & 11 & 157\dg 27.2\pr \\
13 & 12 & 172\dg 29.6\pr & 13 & 13 & 187\dg 32.1\pr & 13 & 14 & 202\dg 34.6\pr & 13 & 15 & 217\dg 37.0\pr \\
13 & 16 & 232\dg 39.5\pr & 13 & 17 & 247\dg 41.9\pr & 13 & 18 & 262\dg 44.4\pr & 13 & 19 & 277\dg 46.9\pr \\
13 & 20 & 292\dg 49.3\pr & 13 & 21 & 307\dg 51.8\pr & 13 & 22 & 322\dg 54.3\pr & 13 & 23 & 337\dg 56.7\pr \\
14 & 0 & 352\dg 59.2\pr & 14 & 1 & 8\dg 01.7\pr & 14 & 2 & 23\dg 04.1\pr & 14 & 3 & 38\dg 06.6\pr \\
14 & 4 & 53\dg 09.0\pr & 14 & 5 & 68\dg 11.5\pr & 14 & 6 & 83\dg 14.0\pr & 14 & 7 & 98\dg 16.4\pr \\
14 & 8 & 113\dg 18.9\pr & 14 & 9 & 128\dg 21.4\pr & 14 & 10 & 143\dg 23.8\pr & 14 & 11 & 158\dg 26.3\pr \\
14 & 12 & 173\dg 28.8\pr & 14 & 13 & 188\dg 31.2\pr & 14 & 14 & 203\dg 33.7\pr & 14 & 15 & 218\dg 36.2\pr \\
14 & 16 & 233\dg 38.6\pr & 14 & 17 & 248\dg 41.1\pr & 14 & 18 & 263\dg 43.5\pr & 14 & 19 & 278\dg 46.0\pr \\
14 & 20 & 293\dg 48.5\pr & 14 & 21 & 308\dg 50.9\pr & 14 & 22 & 323\dg 53.4\pr & 14 & 23 & 338\dg 55.9\pr \\
15 & 0 & 353\dg 58.3\pr & 15 & 1 & 9\dg 00.8\pr & 15 & 2 & 24\dg 03.3\pr & 15 & 3 & 39\dg 05.7\pr \\
15 & 4 & 54\dg 08.2\pr & 15 & 5 & 69\dg 10.7\pr & 15 & 6 & 84\dg 13.1\pr & 15 & 7 & 99\dg 15.6\pr \\
15 & 8 & 114\dg 18.0\pr & 15 & 9 & 129\dg 20.5\pr & 15 & 10 & 144\dg 23.0\pr & 15 & 11 & 159\dg 25.4\pr \\
15 & 12 & 174\dg 27.9\pr & 15 & 13 & 189\dg 30.4\pr & 15 & 14 & 204\dg 32.8\pr & 15 & 15 & 219\dg 35.3\pr \\
15 & 16 & 234\dg 37.8\pr & 15 & 17 & 249\dg 40.2\pr & 15 & 18 & 264\dg 42.7\pr & 15 & 19 & 279\dg 45.1\pr \\
15 & 20 & 294\dg 47.6\pr & 15 & 21 & 309\dg 50.1\pr & 15 & 22 & 324\dg 52.5\pr & 15 & 23 & 339\dg 55.0\pr \\
16 & 0 & 354\dg 57.5\pr & 16 & 1 & 9\dg 59.9\pr & 16 & 2 & 25\dg 02.4\pr & 16 & 3 & 40\dg 04.9\pr \\
16 & 4 & 55\dg 07.3\pr & 16 & 5 & 70\dg 09.8\pr & 16 & 6 & 85\dg 12.3\pr & 16 & 7 & 100\dg 14.7\pr \\
16 & 8 & 115\dg 17.2\pr & 16 & 9 & 130\dg 19.6\pr & 16 & 10 & 145\dg 22.1\pr & 16 & 11 & 160\dg 24.6\pr \\
16 & 12 & 175\dg 27.0\pr & 16 & 13 & 190\dg 29.5\pr & 16 & 14 & 205\dg 32.0\pr & 16 & 15 & 220\dg 34.4\pr \\
16 & 16 & 235\dg 36.9\pr & 16 & 17 & 250\dg 39.4\pr & 16 & 18 & 265\dg 41.8\pr & 16 & 19 & 280\dg 44.3\pr \\
16 & 20 & 295\dg 46.8\pr & 16 & 21 & 310\dg 49.2\pr & 16 & 22 & 325\dg 51.7\pr & 16 & 23 & 340\dg 54.1\pr \\
17 & 0 & 355\dg 56.6\pr & 17 & 1 & 10\dg 59.1\pr & 17 & 2 & 26\dg 01.5\pr & 17 & 3 & 41\dg 04.0\pr \\
17 & 4 & 56\dg 06.5\pr & 17 & 5 & 71\dg 08.9\pr & 17 & 6 & 86\dg 11.4\pr & 17 & 7 & 101\dg 13.9\pr \\
17 & 8 & 116\dg 16.3\pr & 17 & 9 & 131\dg 18.8\pr & 17 & 10 & 146\dg 21.3\pr & 17 & 11 & 161\dg 23.7\pr \\
17 & 12 & 176\dg 26.2\pr & 17 & 13 & 191\dg 28.6\pr & 17 & 14 & 206\dg 31.1\pr & 17 & 15 & 221\dg 33.6\pr \\
17 & 16 & 236\dg 36.0\pr & 17 & 17 & 251\dg 38.5\pr & 17 & 18 & 266\dg 41.0\pr & 17 & 19 & 281\dg 43.4\pr \\
17 & 20 & 296\dg 45.9\pr & 17 & 21 & 311\dg 48.4\pr & 17 & 22 & 326\dg 50.8\pr & 17 & 23 & 341\dg 53.3\pr \\
18 & 0 & 356\dg 55.7\pr & 18 & 1 & 11\dg 58.2\pr & 18 & 2 & 27\dg 00.7\pr & 18 & 3 & 42\dg 03.1\pr \\
18 & 4 & 57\dg 05.6\pr & 18 & 5 & 72\dg 08.1\pr & 18 & 6 & 87\dg 10.5\pr & 18 & 7 & 102\dg 13.0\pr \\
18 & 8 & 117\dg 15.5\pr & 18 & 9 & 132\dg 17.9\pr & 18 & 10 & 147\dg 20.4\pr & 18 & 11 & 162\dg 22.9\pr \\
18 & 12 & 177\dg 25.3\pr & 18 & 13 & 192\dg 27.8\pr & 18 & 14 & 207\dg 30.2\pr & 18 & 15 & 222\dg 32.7\pr \\
18 & 16 & 237\dg 35.2\pr & 18 & 17 & 252\dg 37.6\pr & 18 & 18 & 267\dg 40.1\pr & 18 & 19 & 282\dg 42.6\pr \\
18 & 20 & 297\dg 45.0\pr & 18 & 21 & 312\dg 47.5\pr & 18 & 22 & 327\dg 50.0\pr & 18 & 23 & 342\dg 52.4\pr \\
19 & 0 & 357\dg 54.9\pr & 19 & 1 & 12\dg 57.4\pr & 19 & 2 & 27\dg 59.8\pr & 19 & 3 & 43\dg 02.3\pr \\
19 & 4 & 58\dg 04.7\pr & 19 & 5 & 73\dg 07.2\pr & 19 & 6 & 88\dg 09.7\pr & 19 & 7 & 103\dg 12.1\pr \\
19 & 8 & 118\dg 14.6\pr & 19 & 9 & 133\dg 17.1\pr & 19 & 10 & 148\dg 19.5\pr & 19 & 11 & 163\dg 22.0\pr \\
19 & 12 & 178\dg 24.5\pr & 19 & 13 & 193\dg 26.9\pr & 19 & 14 & 208\dg 29.4\pr & 19 & 15 & 223\dg 31.8\pr \\
19 & 16 & 238\dg 34.3\pr & 19 & 17 & 253\dg 36.8\pr & 19 & 18 & 268\dg 39.2\pr & 19 & 19 & 283\dg 41.7\pr \\
19 & 20 & 298\dg 44.2\pr & 19 & 21 & 313\dg 46.6\pr & 19 & 22 & 328\dg 49.1\pr & 19 & 23 & 343\dg 51.6\pr \\
20 & 0 & 358\dg 54.0\pr & 20 & 1 & 13\dg 56.5\pr & 20 & 2 & 28\dg 59.0\pr & 20 & 3 & 44\dg 01.4\pr \\
20 & 4 & 59\dg 03.9\pr & 20 & 5 & 74\dg 06.3\pr & 20 & 6 & 89\dg 08.8\pr & 20 & 7 & 104\dg 11.3\pr \\
20 & 8 & 119\dg 13.7\pr & 20 & 9 & 134\dg 16.2\pr & 20 & 10 & 149\dg 18.7\pr & 20 & 11 & 164\dg 21.1\pr \\
20 & 12 & 179\dg 23.6\pr & 20 & 13 & 194\dg 26.1\pr & 20 & 14 & 209\dg 28.5\pr & 20 & 15 & 224\dg 31.0\pr \\
20 & 16 & 239\dg 33.5\pr & 20 & 17 & 254\dg 35.9\pr & 20 & 18 & 269\dg 38.4\pr & 20 & 19 & 284\dg 40.8\pr \\
20 & 20 & 299\dg 43.3\pr & 20 & 21 & 314\dg 45.8\pr & 20 & 22 & 329\dg 48.2\pr & 20 & 23 & 344\dg 50.7\pr \\
21 & 0 & 359\dg 53.2\pr & 21 & 1 & 14\dg 55.6\pr & 21 & 2 & 29\dg 58.1\pr & 21 & 3 & 45\dg 00.6\pr \\
21 & 4 & 60\dg 03.0\pr & 21 & 5 & 75\dg 05.5\pr & 21 & 6 & 90\dg 07.9\pr & 21 & 7 & 105\dg 10.4\pr \\
21 & 8 & 120\dg 12.9\pr & 21 & 9 & 135\dg 15.3\pr & 21 & 10 & 150\dg 17.8\pr & 21 & 11 & 165\dg 20.3\pr \\
21 & 12 & 180\dg 22.7\pr & 21 & 13 & 195\dg 25.2\pr & 21 & 14 & 210\dg 27.7\pr & 21 & 15 & 225\dg 30.1\pr \\
21 & 16 & 240\dg 32.6\pr & 21 & 17 & 255\dg 35.1\pr & 21 & 18 & 270\dg 37.5\pr & 21 & 19 & 285\dg 40.0\pr \\
21 & 20 & 300\dg 42.4\pr & 21 & 21 & 315\dg 44.9\pr & 21 & 22 & 330\dg 47.4\pr & 21 & 23 & 345\dg 49.8\pr \\
22 & 0 & 0\dg 52.3\pr & 22 & 1 & 15\dg 54.8\pr & 22 & 2 & 30\dg 57.2\pr & 22 & 3 & 45\dg 59.7\pr \\
22 & 4 & 61\dg 02.2\pr & 22 & 5 & 76\dg 04.6\pr & 22 & 6 & 91\dg 07.1\pr & 22 & 7 & 106\dg 09.6\pr \\
22 & 8 & 121\dg 12.0\pr & 22 & 9 & 136\dg 14.5\pr & 22 & 10 & 151\dg 16.9\pr & 22 & 11 & 166\dg 19.4\pr \\
22 & 12 & 181\dg 21.9\pr & 22 & 13 & 196\dg 24.3\pr & 22 & 14 & 211\dg 26.8\pr & 22 & 15 & 226\dg 29.3\pr \\
22 & 16 & 241\dg 31.7\pr & 22 & 17 & 256\dg 34.2\pr & 22 & 18 & 271\dg 36.7\pr & 22 & 19 & 286\dg 39.1\pr \\
22 & 20 & 301\dg 41.6\pr & 22 & 21 & 316\dg 44.0\pr & 22 & 22 & 331\dg 46.5\pr & 22 & 23 & 346\dg 49.0\pr \\
23 & 0 & 1\dg 51.4\pr & 23 & 1 & 16\dg 53.9\pr & 23 & 2 & 31\dg 56.4\pr & 23 & 3 & 46\dg 58.8\pr \\
23 & 4 & 62\dg 01.3\pr & 23 & 5 & 77\dg 03.8\pr & 23 & 6 & 92\dg 06.2\pr & 23 & 7 & 107\dg 08.7\pr \\
23 & 8 & 122\dg 11.2\pr & 23 & 9 & 137\dg 13.6\pr & 23 & 10 & 152\dg 16.1\pr & 23 & 11 & 167\dg 18.5\pr \\
23 & 12 & 182\dg 21.0\pr & 23 & 13 & 197\dg 23.5\pr & 23 & 14 & 212\dg 25.9\pr & 23 & 15 & 227\dg 28.4\pr \\
23 & 16 & 242\dg 30.9\pr & 23 & 17 & 257\dg 33.3\pr & 23 & 18 & 272\dg 35.8\pr & 23 & 19 & 287\dg 38.3\pr \\
23 & 20 & 302\dg 40.7\pr & 23 & 21 & 317\dg 43.2\pr & 23 & 22 & 332\dg 45.7\pr & 23 & 23 & 347\dg 48.1\pr \\
24 & 0 & 2\dg 50.6\pr & 24 & 1 & 17\dg 53.0\pr & 24 & 2 & 32\dg 55.5\pr & 24 & 3 & 47\dg 58.0\pr \\
24 & 4 & 63\dg 00.4\pr & 24 & 5 & 78\dg 02.9\pr & 24 & 6 & 93\dg 05.4\pr & 24 & 7 & 108\dg 07.8\pr \\
24 & 8 & 123\dg 10.3\pr & 24 & 9 & 138\dg 12.8\pr & 24 & 10 & 153\dg 15.2\pr & 24 & 11 & 168\dg 17.7\pr \\
24 & 12 & 183\dg 20.2\pr & 24 & 13 & 198\dg 22.6\pr & 24 & 14 & 213\dg 25.1\pr & 24 & 15 & 228\dg 27.5\pr \\
24 & 16 & 243\dg 30.0\pr & 24 & 17 & 258\dg 32.5\pr & 24 & 18 & 273\dg 34.9\pr & 24 & 19 & 288\dg 37.4\pr \\
24 & 20 & 303\dg 39.9\pr & 24 & 21 & 318\dg 42.3\pr & 24 & 22 & 333\dg 44.8\pr & 24 & 23 & 348\dg 47.3\pr \\
25 & 0 & 3\dg 49.7\pr & 25 & 1 & 18\dg 52.2\pr & 25 & 2 & 33\dg 54.6\pr & 25 & 3 & 48\dg 57.1\pr \\
25 & 4 & 63\dg 59.6\pr & 25 & 5 & 79\dg 02.0\pr & 25 & 6 & 94\dg 04.5\pr & 25 & 7 & 109\dg 07.0\pr \\
25 & 8 & 124\dg 09.4\pr & 25 & 9 & 139\dg 11.9\pr & 25 & 10 & 154\dg 14.4\pr & 25 & 11 & 169\dg 16.8\pr \\
25 & 12 & 184\dg 19.3\pr & 25 & 13 & 199\dg 21.8\pr & 25 & 14 & 214\dg 24.2\pr & 25 & 15 & 229\dg 26.7\pr \\
25 & 16 & 244\dg 29.1\pr & 25 & 17 & 259\dg 31.6\pr & 25 & 18 & 274\dg 34.1\pr & 25 & 19 & 289\dg 36.5\pr \\
25 & 20 & 304\dg 39.0\pr & 25 & 21 & 319\dg 41.5\pr & 25 & 22 & 334\dg 43.9\pr & 25 & 23 & 349\dg 46.4\pr \\
26 & 0 & 4\dg 48.9\pr & 26 & 1 & 19\dg 51.3\pr & 26 & 2 & 34\dg 53.8\pr & 26 & 3 & 49\dg 56.3\pr \\
26 & 4 & 64\dg 58.7\pr & 26 & 5 & 80\dg 01.2\pr & 26 & 6 & 95\dg 03.6\pr & 26 & 7 & 110\dg 06.1\pr \\
26 & 8 & 125\dg 08.6\pr & 26 & 9 & 140\dg 11.0\pr & 26 & 10 & 155\dg 13.5\pr & 26 & 11 & 170\dg 16.0\pr \\
26 & 12 & 185\dg 18.4\pr & 26 & 13 & 200\dg 20.9\pr & 26 & 14 & 215\dg 23.4\pr & 26 & 15 & 230\dg 25.8\pr \\
26 & 16 & 245\dg 28.3\pr & 26 & 17 & 260\dg 30.7\pr & 26 & 18 & 275\dg 33.2\pr & 26 & 19 & 290\dg 35.7\pr \\
26 & 20 & 305\dg 38.1\pr & 26 & 21 & 320\dg 40.6\pr & 26 & 22 & 335\dg 43.1\pr & 26 & 23 & 350\dg 45.5\pr \\
27 & 0 & 5\dg 48.0\pr & 27 & 1 & 20\dg 50.5\pr & 27 & 2 & 35\dg 52.9\pr & 27 & 3 & 50\dg 55.4\pr \\
27 & 4 & 65\dg 57.9\pr & 27 & 5 & 81\dg 00.3\pr & 27 & 6 & 96\dg 02.8\pr & 27 & 7 & 111\dg 05.2\pr \\
27 & 8 & 126\dg 07.7\pr & 27 & 9 & 141\dg 10.2\pr & 27 & 10 & 156\dg 12.6\pr & 27 & 11 & 171\dg 15.1\pr \\
27 & 12 & 186\dg 17.6\pr & 27 & 13 & 201\dg 20.0\pr & 27 & 14 & 216\dg 22.5\pr & 27 & 15 & 231\dg 25.0\pr \\
27 & 16 & 246\dg 27.4\pr & 27 & 17 & 261\dg 29.9\pr & 27 & 18 & 276\dg 32.4\pr & 27 & 19 & 291\dg 34.8\pr \\
27 & 20 & 306\dg 37.3\pr & 27 & 21 & 321\dg 39.7\pr & 27 & 22 & 336\dg 42.2\pr & 27 & 23 & 351\dg 44.7\pr \\
28 & 0 & 6\dg 47.1\pr & 28 & 1 & 21\dg 49.6\pr & 28 & 2 & 36\dg 52.1\pr & 28 & 3 & 51\dg 54.5\pr \\
28 & 4 & 66\dg 57.0\pr & 28 & 5 & 81\dg 59.5\pr & 28 & 6 & 97\dg 01.9\pr & 28 & 7 & 112\dg 04.4\pr \\
28 & 8 & 127\dg 06.8\pr & 28 & 9 & 142\dg 09.3\pr & 28 & 10 & 157\dg 11.8\pr & 28 & 11 & 172\dg 14.2\pr \\
28 & 12 & 187\dg 16.7\pr & 28 & 13 & 202\dg 19.2\pr & 28 & 14 & 217\dg 21.6\pr & 28 & 15 & 232\dg 24.1\pr \\
28 & 16 & 247\dg 26.6\pr & 28 & 17 & 262\dg 29.0\pr & 28 & 18 & 277\dg 31.5\pr & 28 & 19 & 292\dg 34.0\pr \\
28 & 20 & 307\dg 36.4\pr & 28 & 21 & 322\dg 38.9\pr & 28 & 22 & 337\dg 41.3\pr & 28 & 23 & 352\dg 43.8\pr \\
29 & 0 & 7\dg 46.3\pr & 29 & 1 & 22\dg 48.7\pr & 29 & 2 & 37\dg 51.2\pr & 29 & 3 & 52\dg 53.7\pr \\
29 & 4 & 67\dg 56.1\pr & 29 & 5 & 82\dg 58.6\pr & 29 & 6 & 98\dg 01.1\pr & 29 & 7 & 113\dg 03.5\pr \\
29 & 8 & 128\dg 06.0\pr & 29 & 9 & 143\dg 08.5\pr & 29 & 10 & 158\dg 10.9\pr & 29 & 11 & 173\dg 13.4\pr \\
29 & 12 & 188\dg 15.8\pr & 29 & 13 & 203\dg 18.3\pr & 29 & 14 & 218\dg 20.8\pr & 29 & 15 & 233\dg 23.2\pr \\
29 & 16 & 248\dg 25.7\pr & 29 & 17 & 263\dg 28.2\pr & 29 & 18 & 278\dg 30.6\pr & 29 & 19 & 293\dg 33.1\pr \\
29 & 20 & 308\dg 35.6\pr & 29 & 21 & 323\dg 38.0\pr & 29 & 22 & 338\dg 40.5\pr & 29 & 23 & 353\dg 43.0\pr \\
30 & 0 & 8\dg 45.4\pr & 30 & 1 & 23\dg 47.9\pr & 30 & 2 & 38\dg 50.3\pr & 30 & 3 & 53\dg 52.8\pr \\
30 & 4 & 68\dg 55.3\pr & 30 & 5 & 83\dg 57.7\pr & 30 & 6 & 99\dg 00.2\pr & 30 & 7 & 114\dg 02.7\pr \\
30 & 8 & 129\dg 05.1\pr & 30 & 9 & 144\dg 07.6\pr & 30 & 10 & 159\dg 10.1\pr & 30 & 11 & 174\dg 12.5\pr \\
30 & 12 & 189\dg 15.0\pr & 30 & 13 & 204\dg 17.4\pr & 30 & 14 & 219\dg 19.9\pr & 30 & 15 & 234\dg 22.4\pr \\
30 & 16 & 249\dg 24.8\pr & 30 & 17 & 264\dg 27.3\pr & 30 & 18 & 279\dg 29.8\pr & 30 & 19 & 294\dg 32.2\pr \\
30 & 20 & 309\dg 34.7\pr & 30 & 21 & 324\dg 37.2\pr & 30 & 22 & 339\dg 39.6\pr & 30 & 23 & 354\dg 42.1\pr \\
\end{longtable}}

{\footnotesize
\begin{longtable}{rrr|rrr|rrr|rrr}
\caption{Aries --- GHA of the First Point of Aries, October 2026 (UT)}\label{aries10}\\
\toprule
\textbf{d} & \textbf{h} & \textbf{GHA} $\gamma$ & \textbf{d} & \textbf{h} & \textbf{GHA} $\gamma$ & \textbf{d} & \textbf{h} & \textbf{GHA} $\gamma$ & \textbf{d} & \textbf{h} & \textbf{GHA} $\gamma$ \\
\midrule
\endfirsthead
\multicolumn{12}{c}{\textit{Aries --- GHA of the First Point of Aries, October 2026 (UT) (continued)}}\\
\toprule
\textbf{d} & \textbf{h} & \textbf{GHA} $\gamma$ & \textbf{d} & \textbf{h} & \textbf{GHA} $\gamma$ & \textbf{d} & \textbf{h} & \textbf{GHA} $\gamma$ & \textbf{d} & \textbf{h} & \textbf{GHA} $\gamma$ \\
\midrule
\endhead
\midrule\multicolumn{12}{r}{\textit{continued on next page}}\\
\endfoot
\bottomrule
\endlastfoot
1 & 0 & 9\dg 44.6\pr & 1 & 1 & 24\dg 47.0\pr & 1 & 2 & 39\dg 49.5\pr & 1 & 3 & 54\dg 51.9\pr \\
1 & 4 & 69\dg 54.4\pr & 1 & 5 & 84\dg 56.9\pr & 1 & 6 & 99\dg 59.3\pr & 1 & 7 & 115\dg 01.8\pr \\
1 & 8 & 130\dg 04.3\pr & 1 & 9 & 145\dg 06.7\pr & 1 & 10 & 160\dg 09.2\pr & 1 & 11 & 175\dg 11.7\pr \\
1 & 12 & 190\dg 14.1\pr & 1 & 13 & 205\dg 16.6\pr & 1 & 14 & 220\dg 19.1\pr & 1 & 15 & 235\dg 21.5\pr \\
1 & 16 & 250\dg 24.0\pr & 1 & 17 & 265\dg 26.4\pr & 1 & 18 & 280\dg 28.9\pr & 1 & 19 & 295\dg 31.4\pr \\
1 & 20 & 310\dg 33.8\pr & 1 & 21 & 325\dg 36.3\pr & 1 & 22 & 340\dg 38.8\pr & 1 & 23 & 355\dg 41.2\pr \\
2 & 0 & 10\dg 43.7\pr & 2 & 1 & 25\dg 46.2\pr & 2 & 2 & 40\dg 48.6\pr & 2 & 3 & 55\dg 51.1\pr \\
2 & 4 & 70\dg 53.5\pr & 2 & 5 & 85\dg 56.0\pr & 2 & 6 & 100\dg 58.5\pr & 2 & 7 & 116\dg 00.9\pr \\
2 & 8 & 131\dg 03.4\pr & 2 & 9 & 146\dg 05.9\pr & 2 & 10 & 161\dg 08.3\pr & 2 & 11 & 176\dg 10.8\pr \\
2 & 12 & 191\dg 13.3\pr & 2 & 13 & 206\dg 15.7\pr & 2 & 14 & 221\dg 18.2\pr & 2 & 15 & 236\dg 20.7\pr \\
2 & 16 & 251\dg 23.1\pr & 2 & 17 & 266\dg 25.6\pr & 2 & 18 & 281\dg 28.0\pr & 2 & 19 & 296\dg 30.5\pr \\
2 & 20 & 311\dg 33.0\pr & 2 & 21 & 326\dg 35.4\pr & 2 & 22 & 341\dg 37.9\pr & 2 & 23 & 356\dg 40.4\pr \\
3 & 0 & 11\dg 42.8\pr & 3 & 1 & 26\dg 45.3\pr & 3 & 2 & 41\dg 47.8\pr & 3 & 3 & 56\dg 50.2\pr \\
3 & 4 & 71\dg 52.7\pr & 3 & 5 & 86\dg 55.2\pr & 3 & 6 & 101\dg 57.6\pr & 3 & 7 & 117\dg 00.1\pr \\
3 & 8 & 132\dg 02.5\pr & 3 & 9 & 147\dg 05.0\pr & 3 & 10 & 162\dg 07.5\pr & 3 & 11 & 177\dg 09.9\pr \\
3 & 12 & 192\dg 12.4\pr & 3 & 13 & 207\dg 14.9\pr & 3 & 14 & 222\dg 17.3\pr & 3 & 15 & 237\dg 19.8\pr \\
3 & 16 & 252\dg 22.3\pr & 3 & 17 & 267\dg 24.7\pr & 3 & 18 & 282\dg 27.2\pr & 3 & 19 & 297\dg 29.6\pr \\
3 & 20 & 312\dg 32.1\pr & 3 & 21 & 327\dg 34.6\pr & 3 & 22 & 342\dg 37.0\pr & 3 & 23 & 357\dg 39.5\pr \\
4 & 0 & 12\dg 42.0\pr & 4 & 1 & 27\dg 44.4\pr & 4 & 2 & 42\dg 46.9\pr & 4 & 3 & 57\dg 49.4\pr \\
4 & 4 & 72\dg 51.8\pr & 4 & 5 & 87\dg 54.3\pr & 4 & 6 & 102\dg 56.8\pr & 4 & 7 & 117\dg 59.2\pr \\
4 & 8 & 133\dg 01.7\pr & 4 & 9 & 148\dg 04.1\pr & 4 & 10 & 163\dg 06.6\pr & 4 & 11 & 178\dg 09.1\pr \\
4 & 12 & 193\dg 11.5\pr & 4 & 13 & 208\dg 14.0\pr & 4 & 14 & 223\dg 16.5\pr & 4 & 15 & 238\dg 18.9\pr \\
4 & 16 & 253\dg 21.4\pr & 4 & 17 & 268\dg 23.9\pr & 4 & 18 & 283\dg 26.3\pr & 4 & 19 & 298\dg 28.8\pr \\
4 & 20 & 313\dg 31.3\pr & 4 & 21 & 328\dg 33.7\pr & 4 & 22 & 343\dg 36.2\pr & 4 & 23 & 358\dg 38.6\pr \\
5 & 0 & 13\dg 41.1\pr & 5 & 1 & 28\dg 43.6\pr & 5 & 2 & 43\dg 46.0\pr & 5 & 3 & 58\dg 48.5\pr \\
5 & 4 & 73\dg 51.0\pr & 5 & 5 & 88\dg 53.4\pr & 5 & 6 & 103\dg 55.9\pr & 5 & 7 & 118\dg 58.4\pr \\
5 & 8 & 134\dg 00.8\pr & 5 & 9 & 149\dg 03.3\pr & 5 & 10 & 164\dg 05.7\pr & 5 & 11 & 179\dg 08.2\pr \\
5 & 12 & 194\dg 10.7\pr & 5 & 13 & 209\dg 13.1\pr & 5 & 14 & 224\dg 15.6\pr & 5 & 15 & 239\dg 18.1\pr \\
5 & 16 & 254\dg 20.5\pr & 5 & 17 & 269\dg 23.0\pr & 5 & 18 & 284\dg 25.5\pr & 5 & 19 & 299\dg 27.9\pr \\
5 & 20 & 314\dg 30.4\pr & 5 & 21 & 329\dg 32.9\pr & 5 & 22 & 344\dg 35.3\pr & 5 & 23 & 359\dg 37.8\pr \\
6 & 0 & 14\dg 40.2\pr & 6 & 1 & 29\dg 42.7\pr & 6 & 2 & 44\dg 45.2\pr & 6 & 3 & 59\dg 47.6\pr \\
6 & 4 & 74\dg 50.1\pr & 6 & 5 & 89\dg 52.6\pr & 6 & 6 & 104\dg 55.0\pr & 6 & 7 & 119\dg 57.5\pr \\
6 & 8 & 135\dg 00.0\pr & 6 & 9 & 150\dg 02.4\pr & 6 & 10 & 165\dg 04.9\pr & 6 & 11 & 180\dg 07.4\pr \\
6 & 12 & 195\dg 09.8\pr & 6 & 13 & 210\dg 12.3\pr & 6 & 14 & 225\dg 14.7\pr & 6 & 15 & 240\dg 17.2\pr \\
6 & 16 & 255\dg 19.7\pr & 6 & 17 & 270\dg 22.1\pr & 6 & 18 & 285\dg 24.6\pr & 6 & 19 & 300\dg 27.1\pr \\
6 & 20 & 315\dg 29.5\pr & 6 & 21 & 330\dg 32.0\pr & 6 & 22 & 345\dg 34.5\pr & 6 & 23 & 0\dg 36.9\pr \\
7 & 0 & 15\dg 39.4\pr & 7 & 1 & 30\dg 41.9\pr & 7 & 2 & 45\dg 44.3\pr & 7 & 3 & 60\dg 46.8\pr \\
7 & 4 & 75\dg 49.2\pr & 7 & 5 & 90\dg 51.7\pr & 7 & 6 & 105\dg 54.2\pr & 7 & 7 & 120\dg 56.6\pr \\
7 & 8 & 135\dg 59.1\pr & 7 & 9 & 151\dg 01.6\pr & 7 & 10 & 166\dg 04.0\pr & 7 & 11 & 181\dg 06.5\pr \\
7 & 12 & 196\dg 09.0\pr & 7 & 13 & 211\dg 11.4\pr & 7 & 14 & 226\dg 13.9\pr & 7 & 15 & 241\dg 16.3\pr \\
7 & 16 & 256\dg 18.8\pr & 7 & 17 & 271\dg 21.3\pr & 7 & 18 & 286\dg 23.7\pr & 7 & 19 & 301\dg 26.2\pr \\
7 & 20 & 316\dg 28.7\pr & 7 & 21 & 331\dg 31.1\pr & 7 & 22 & 346\dg 33.6\pr & 7 & 23 & 1\dg 36.1\pr \\
8 & 0 & 16\dg 38.5\pr & 8 & 1 & 31\dg 41.0\pr & 8 & 2 & 46\dg 43.5\pr & 8 & 3 & 61\dg 45.9\pr \\
8 & 4 & 76\dg 48.4\pr & 8 & 5 & 91\dg 50.8\pr & 8 & 6 & 106\dg 53.3\pr & 8 & 7 & 121\dg 55.8\pr \\
8 & 8 & 136\dg 58.2\pr & 8 & 9 & 152\dg 00.7\pr & 8 & 10 & 167\dg 03.2\pr & 8 & 11 & 182\dg 05.6\pr \\
8 & 12 & 197\dg 08.1\pr & 8 & 13 & 212\dg 10.6\pr & 8 & 14 & 227\dg 13.0\pr & 8 & 15 & 242\dg 15.5\pr \\
8 & 16 & 257\dg 18.0\pr & 8 & 17 & 272\dg 20.4\pr & 8 & 18 & 287\dg 22.9\pr & 8 & 19 & 302\dg 25.3\pr \\
8 & 20 & 317\dg 27.8\pr & 8 & 21 & 332\dg 30.3\pr & 8 & 22 & 347\dg 32.7\pr & 8 & 23 & 2\dg 35.2\pr \\
9 & 0 & 17\dg 37.7\pr & 9 & 1 & 32\dg 40.1\pr & 9 & 2 & 47\dg 42.6\pr & 9 & 3 & 62\dg 45.1\pr \\
9 & 4 & 77\dg 47.5\pr & 9 & 5 & 92\dg 50.0\pr & 9 & 6 & 107\dg 52.4\pr & 9 & 7 & 122\dg 54.9\pr \\
9 & 8 & 137\dg 57.4\pr & 9 & 9 & 152\dg 59.8\pr & 9 & 10 & 168\dg 02.3\pr & 9 & 11 & 183\dg 04.8\pr \\
9 & 12 & 198\dg 07.2\pr & 9 & 13 & 213\dg 09.7\pr & 9 & 14 & 228\dg 12.2\pr & 9 & 15 & 243\dg 14.6\pr \\
9 & 16 & 258\dg 17.1\pr & 9 & 17 & 273\dg 19.6\pr & 9 & 18 & 288\dg 22.0\pr & 9 & 19 & 303\dg 24.5\pr \\
9 & 20 & 318\dg 26.9\pr & 9 & 21 & 333\dg 29.4\pr & 9 & 22 & 348\dg 31.9\pr & 9 & 23 & 3\dg 34.3\pr \\
10 & 0 & 18\dg 36.8\pr & 10 & 1 & 33\dg 39.3\pr & 10 & 2 & 48\dg 41.7\pr & 10 & 3 & 63\dg 44.2\pr \\
10 & 4 & 78\dg 46.7\pr & 10 & 5 & 93\dg 49.1\pr & 10 & 6 & 108\dg 51.6\pr & 10 & 7 & 123\dg 54.1\pr \\
10 & 8 & 138\dg 56.5\pr & 10 & 9 & 153\dg 59.0\pr & 10 & 10 & 169\dg 01.4\pr & 10 & 11 & 184\dg 03.9\pr \\
10 & 12 & 199\dg 06.4\pr & 10 & 13 & 214\dg 08.8\pr & 10 & 14 & 229\dg 11.3\pr & 10 & 15 & 244\dg 13.8\pr \\
10 & 16 & 259\dg 16.2\pr & 10 & 17 & 274\dg 18.7\pr & 10 & 18 & 289\dg 21.2\pr & 10 & 19 & 304\dg 23.6\pr \\
10 & 20 & 319\dg 26.1\pr & 10 & 21 & 334\dg 28.5\pr & 10 & 22 & 349\dg 31.0\pr & 10 & 23 & 4\dg 33.5\pr \\
11 & 0 & 19\dg 35.9\pr & 11 & 1 & 34\dg 38.4\pr & 11 & 2 & 49\dg 40.9\pr & 11 & 3 & 64\dg 43.3\pr \\
11 & 4 & 79\dg 45.8\pr & 11 & 5 & 94\dg 48.3\pr & 11 & 6 & 109\dg 50.7\pr & 11 & 7 & 124\dg 53.2\pr \\
11 & 8 & 139\dg 55.7\pr & 11 & 9 & 154\dg 58.1\pr & 11 & 10 & 170\dg 00.6\pr & 11 & 11 & 185\dg 03.0\pr \\
11 & 12 & 200\dg 05.5\pr & 11 & 13 & 215\dg 08.0\pr & 11 & 14 & 230\dg 10.4\pr & 11 & 15 & 245\dg 12.9\pr \\
11 & 16 & 260\dg 15.4\pr & 11 & 17 & 275\dg 17.8\pr & 11 & 18 & 290\dg 20.3\pr & 11 & 19 & 305\dg 22.8\pr \\
11 & 20 & 320\dg 25.2\pr & 11 & 21 & 335\dg 27.7\pr & 11 & 22 & 350\dg 30.2\pr & 11 & 23 & 5\dg 32.6\pr \\
12 & 0 & 20\dg 35.1\pr & 12 & 1 & 35\dg 37.5\pr & 12 & 2 & 50\dg 40.0\pr & 12 & 3 & 65\dg 42.5\pr \\
12 & 4 & 80\dg 44.9\pr & 12 & 5 & 95\dg 47.4\pr & 12 & 6 & 110\dg 49.9\pr & 12 & 7 & 125\dg 52.3\pr \\
12 & 8 & 140\dg 54.8\pr & 12 & 9 & 155\dg 57.3\pr & 12 & 10 & 170\dg 59.7\pr & 12 & 11 & 186\dg 02.2\pr \\
12 & 12 & 201\dg 04.6\pr & 12 & 13 & 216\dg 07.1\pr & 12 & 14 & 231\dg 09.6\pr & 12 & 15 & 246\dg 12.0\pr \\
12 & 16 & 261\dg 14.5\pr & 12 & 17 & 276\dg 17.0\pr & 12 & 18 & 291\dg 19.4\pr & 12 & 19 & 306\dg 21.9\pr \\
12 & 20 & 321\dg 24.4\pr & 12 & 21 & 336\dg 26.8\pr & 12 & 22 & 351\dg 29.3\pr & 12 & 23 & 6\dg 31.8\pr \\
13 & 0 & 21\dg 34.2\pr & 13 & 1 & 36\dg 36.7\pr & 13 & 2 & 51\dg 39.1\pr & 13 & 3 & 66\dg 41.6\pr \\
13 & 4 & 81\dg 44.1\pr & 13 & 5 & 96\dg 46.5\pr & 13 & 6 & 111\dg 49.0\pr & 13 & 7 & 126\dg 51.5\pr \\
13 & 8 & 141\dg 53.9\pr & 13 & 9 & 156\dg 56.4\pr & 13 & 10 & 171\dg 58.9\pr & 13 & 11 & 187\dg 01.3\pr \\
13 & 12 & 202\dg 03.8\pr & 13 & 13 & 217\dg 06.3\pr & 13 & 14 & 232\dg 08.7\pr & 13 & 15 & 247\dg 11.2\pr \\
13 & 16 & 262\dg 13.6\pr & 13 & 17 & 277\dg 16.1\pr & 13 & 18 & 292\dg 18.6\pr & 13 & 19 & 307\dg 21.0\pr \\
13 & 20 & 322\dg 23.5\pr & 13 & 21 & 337\dg 26.0\pr & 13 & 22 & 352\dg 28.4\pr & 13 & 23 & 7\dg 30.9\pr \\
14 & 0 & 22\dg 33.4\pr & 14 & 1 & 37\dg 35.8\pr & 14 & 2 & 52\dg 38.3\pr & 14 & 3 & 67\dg 40.8\pr \\
14 & 4 & 82\dg 43.2\pr & 14 & 5 & 97\dg 45.7\pr & 14 & 6 & 112\dg 48.1\pr & 14 & 7 & 127\dg 50.6\pr \\
14 & 8 & 142\dg 53.1\pr & 14 & 9 & 157\dg 55.5\pr & 14 & 10 & 172\dg 58.0\pr & 14 & 11 & 188\dg 00.5\pr \\
14 & 12 & 203\dg 02.9\pr & 14 & 13 & 218\dg 05.4\pr & 14 & 14 & 233\dg 07.9\pr & 14 & 15 & 248\dg 10.3\pr \\
14 & 16 & 263\dg 12.8\pr & 14 & 17 & 278\dg 15.2\pr & 14 & 18 & 293\dg 17.7\pr & 14 & 19 & 308\dg 20.2\pr \\
14 & 20 & 323\dg 22.6\pr & 14 & 21 & 338\dg 25.1\pr & 14 & 22 & 353\dg 27.6\pr & 14 & 23 & 8\dg 30.0\pr \\
15 & 0 & 23\dg 32.5\pr & 15 & 1 & 38\dg 35.0\pr & 15 & 2 & 53\dg 37.4\pr & 15 & 3 & 68\dg 39.9\pr \\
15 & 4 & 83\dg 42.4\pr & 15 & 5 & 98\dg 44.8\pr & 15 & 6 & 113\dg 47.3\pr & 15 & 7 & 128\dg 49.7\pr \\
15 & 8 & 143\dg 52.2\pr & 15 & 9 & 158\dg 54.7\pr & 15 & 10 & 173\dg 57.1\pr & 15 & 11 & 188\dg 59.6\pr \\
15 & 12 & 204\dg 02.1\pr & 15 & 13 & 219\dg 04.5\pr & 15 & 14 & 234\dg 07.0\pr & 15 & 15 & 249\dg 09.5\pr \\
15 & 16 & 264\dg 11.9\pr & 15 & 17 & 279\dg 14.4\pr & 15 & 18 & 294\dg 16.9\pr & 15 & 19 & 309\dg 19.3\pr \\
15 & 20 & 324\dg 21.8\pr & 15 & 21 & 339\dg 24.2\pr & 15 & 22 & 354\dg 26.7\pr & 15 & 23 & 9\dg 29.2\pr \\
16 & 0 & 24\dg 31.6\pr & 16 & 1 & 39\dg 34.1\pr & 16 & 2 & 54\dg 36.6\pr & 16 & 3 & 69\dg 39.0\pr \\
16 & 4 & 84\dg 41.5\pr & 16 & 5 & 99\dg 44.0\pr & 16 & 6 & 114\dg 46.4\pr & 16 & 7 & 129\dg 48.9\pr \\
16 & 8 & 144\dg 51.3\pr & 16 & 9 & 159\dg 53.8\pr & 16 & 10 & 174\dg 56.3\pr & 16 & 11 & 189\dg 58.7\pr \\
16 & 12 & 205\dg 01.2\pr & 16 & 13 & 220\dg 03.7\pr & 16 & 14 & 235\dg 06.1\pr & 16 & 15 & 250\dg 08.6\pr \\
16 & 16 & 265\dg 11.1\pr & 16 & 17 & 280\dg 13.5\pr & 16 & 18 & 295\dg 16.0\pr & 16 & 19 & 310\dg 18.5\pr \\
16 & 20 & 325\dg 20.9\pr & 16 & 21 & 340\dg 23.4\pr & 16 & 22 & 355\dg 25.8\pr & 16 & 23 & 10\dg 28.3\pr \\
17 & 0 & 25\dg 30.8\pr & 17 & 1 & 40\dg 33.2\pr & 17 & 2 & 55\dg 35.7\pr & 17 & 3 & 70\dg 38.2\pr \\
17 & 4 & 85\dg 40.6\pr & 17 & 5 & 100\dg 43.1\pr & 17 & 6 & 115\dg 45.6\pr & 17 & 7 & 130\dg 48.0\pr \\
17 & 8 & 145\dg 50.5\pr & 17 & 9 & 160\dg 53.0\pr & 17 & 10 & 175\dg 55.4\pr & 17 & 11 & 190\dg 57.9\pr \\
17 & 12 & 206\dg 00.3\pr & 17 & 13 & 221\dg 02.8\pr & 17 & 14 & 236\dg 05.3\pr & 17 & 15 & 251\dg 07.7\pr \\
17 & 16 & 266\dg 10.2\pr & 17 & 17 & 281\dg 12.7\pr & 17 & 18 & 296\dg 15.1\pr & 17 & 19 & 311\dg 17.6\pr \\
17 & 20 & 326\dg 20.1\pr & 17 & 21 & 341\dg 22.5\pr & 17 & 22 & 356\dg 25.0\pr & 17 & 23 & 11\dg 27.4\pr \\
18 & 0 & 26\dg 29.9\pr & 18 & 1 & 41\dg 32.4\pr & 18 & 2 & 56\dg 34.8\pr & 18 & 3 & 71\dg 37.3\pr \\
18 & 4 & 86\dg 39.8\pr & 18 & 5 & 101\dg 42.2\pr & 18 & 6 & 116\dg 44.7\pr & 18 & 7 & 131\dg 47.2\pr \\
18 & 8 & 146\dg 49.6\pr & 18 & 9 & 161\dg 52.1\pr & 18 & 10 & 176\dg 54.6\pr & 18 & 11 & 191\dg 57.0\pr \\
18 & 12 & 206\dg 59.5\pr & 18 & 13 & 222\dg 01.9\pr & 18 & 14 & 237\dg 04.4\pr & 18 & 15 & 252\dg 06.9\pr \\
18 & 16 & 267\dg 09.3\pr & 18 & 17 & 282\dg 11.8\pr & 18 & 18 & 297\dg 14.3\pr & 18 & 19 & 312\dg 16.7\pr \\
18 & 20 & 327\dg 19.2\pr & 18 & 21 & 342\dg 21.7\pr & 18 & 22 & 357\dg 24.1\pr & 18 & 23 & 12\dg 26.6\pr \\
19 & 0 & 27\dg 29.1\pr & 19 & 1 & 42\dg 31.5\pr & 19 & 2 & 57\dg 34.0\pr & 19 & 3 & 72\dg 36.4\pr \\
19 & 4 & 87\dg 38.9\pr & 19 & 5 & 102\dg 41.4\pr & 19 & 6 & 117\dg 43.8\pr & 19 & 7 & 132\dg 46.3\pr \\
19 & 8 & 147\dg 48.8\pr & 19 & 9 & 162\dg 51.2\pr & 19 & 10 & 177\dg 53.7\pr & 19 & 11 & 192\dg 56.2\pr \\
19 & 12 & 207\dg 58.6\pr & 19 & 13 & 223\dg 01.1\pr & 19 & 14 & 238\dg 03.5\pr & 19 & 15 & 253\dg 06.0\pr \\
19 & 16 & 268\dg 08.5\pr & 19 & 17 & 283\dg 10.9\pr & 19 & 18 & 298\dg 13.4\pr & 19 & 19 & 313\dg 15.9\pr \\
19 & 20 & 328\dg 18.3\pr & 19 & 21 & 343\dg 20.8\pr & 19 & 22 & 358\dg 23.3\pr & 19 & 23 & 13\dg 25.7\pr \\
20 & 0 & 28\dg 28.2\pr & 20 & 1 & 43\dg 30.7\pr & 20 & 2 & 58\dg 33.1\pr & 20 & 3 & 73\dg 35.6\pr \\
20 & 4 & 88\dg 38.0\pr & 20 & 5 & 103\dg 40.5\pr & 20 & 6 & 118\dg 43.0\pr & 20 & 7 & 133\dg 45.4\pr \\
20 & 8 & 148\dg 47.9\pr & 20 & 9 & 163\dg 50.4\pr & 20 & 10 & 178\dg 52.8\pr & 20 & 11 & 193\dg 55.3\pr \\
20 & 12 & 208\dg 57.8\pr & 20 & 13 & 224\dg 00.2\pr & 20 & 14 & 239\dg 02.7\pr & 20 & 15 & 254\dg 05.2\pr \\
20 & 16 & 269\dg 07.6\pr & 20 & 17 & 284\dg 10.1\pr & 20 & 18 & 299\dg 12.5\pr & 20 & 19 & 314\dg 15.0\pr \\
20 & 20 & 329\dg 17.5\pr & 20 & 21 & 344\dg 19.9\pr & 20 & 22 & 359\dg 22.4\pr & 20 & 23 & 14\dg 24.9\pr \\
21 & 0 & 29\dg 27.3\pr & 21 & 1 & 44\dg 29.8\pr & 21 & 2 & 59\dg 32.3\pr & 21 & 3 & 74\dg 34.7\pr \\
21 & 4 & 89\dg 37.2\pr & 21 & 5 & 104\dg 39.7\pr & 21 & 6 & 119\dg 42.1\pr & 21 & 7 & 134\dg 44.6\pr \\
21 & 8 & 149\dg 47.0\pr & 21 & 9 & 164\dg 49.5\pr & 21 & 10 & 179\dg 52.0\pr & 21 & 11 & 194\dg 54.4\pr \\
21 & 12 & 209\dg 56.9\pr & 21 & 13 & 224\dg 59.4\pr & 21 & 14 & 240\dg 01.8\pr & 21 & 15 & 255\dg 04.3\pr \\
21 & 16 & 270\dg 06.8\pr & 21 & 17 & 285\dg 09.2\pr & 21 & 18 & 300\dg 11.7\pr & 21 & 19 & 315\dg 14.1\pr \\
21 & 20 & 330\dg 16.6\pr & 21 & 21 & 345\dg 19.1\pr & 21 & 22 & 0\dg 21.5\pr & 21 & 23 & 15\dg 24.0\pr \\
22 & 0 & 30\dg 26.5\pr & 22 & 1 & 45\dg 28.9\pr & 22 & 2 & 60\dg 31.4\pr & 22 & 3 & 75\dg 33.9\pr \\
22 & 4 & 90\dg 36.3\pr & 22 & 5 & 105\dg 38.8\pr & 22 & 6 & 120\dg 41.3\pr & 22 & 7 & 135\dg 43.7\pr \\
22 & 8 & 150\dg 46.2\pr & 22 & 9 & 165\dg 48.6\pr & 22 & 10 & 180\dg 51.1\pr & 22 & 11 & 195\dg 53.6\pr \\
22 & 12 & 210\dg 56.0\pr & 22 & 13 & 225\dg 58.5\pr & 22 & 14 & 241\dg 01.0\pr & 22 & 15 & 256\dg 03.4\pr \\
22 & 16 & 271\dg 05.9\pr & 22 & 17 & 286\dg 08.4\pr & 22 & 18 & 301\dg 10.8\pr & 22 & 19 & 316\dg 13.3\pr \\
22 & 20 & 331\dg 15.8\pr & 22 & 21 & 346\dg 18.2\pr & 22 & 22 & 1\dg 20.7\pr & 22 & 23 & 16\dg 23.1\pr \\
23 & 0 & 31\dg 25.6\pr & 23 & 1 & 46\dg 28.1\pr & 23 & 2 & 61\dg 30.5\pr & 23 & 3 & 76\dg 33.0\pr \\
23 & 4 & 91\dg 35.5\pr & 23 & 5 & 106\dg 37.9\pr & 23 & 6 & 121\dg 40.4\pr & 23 & 7 & 136\dg 42.9\pr \\
23 & 8 & 151\dg 45.3\pr & 23 & 9 & 166\dg 47.8\pr & 23 & 10 & 181\dg 50.2\pr & 23 & 11 & 196\dg 52.7\pr \\
23 & 12 & 211\dg 55.2\pr & 23 & 13 & 226\dg 57.6\pr & 23 & 14 & 242\dg 00.1\pr & 23 & 15 & 257\dg 02.6\pr \\
23 & 16 & 272\dg 05.0\pr & 23 & 17 & 287\dg 07.5\pr & 23 & 18 & 302\dg 10.0\pr & 23 & 19 & 317\dg 12.4\pr \\
23 & 20 & 332\dg 14.9\pr & 23 & 21 & 347\dg 17.4\pr & 23 & 22 & 2\dg 19.8\pr & 23 & 23 & 17\dg 22.3\pr \\
24 & 0 & 32\dg 24.7\pr & 24 & 1 & 47\dg 27.2\pr & 24 & 2 & 62\dg 29.7\pr & 24 & 3 & 77\dg 32.1\pr \\
24 & 4 & 92\dg 34.6\pr & 24 & 5 & 107\dg 37.1\pr & 24 & 6 & 122\dg 39.5\pr & 24 & 7 & 137\dg 42.0\pr \\
24 & 8 & 152\dg 44.5\pr & 24 & 9 & 167\dg 46.9\pr & 24 & 10 & 182\dg 49.4\pr & 24 & 11 & 197\dg 51.9\pr \\
24 & 12 & 212\dg 54.3\pr & 24 & 13 & 227\dg 56.8\pr & 24 & 14 & 242\dg 59.2\pr & 24 & 15 & 258\dg 01.7\pr \\
24 & 16 & 273\dg 04.2\pr & 24 & 17 & 288\dg 06.6\pr & 24 & 18 & 303\dg 09.1\pr & 24 & 19 & 318\dg 11.6\pr \\
24 & 20 & 333\dg 14.0\pr & 24 & 21 & 348\dg 16.5\pr & 24 & 22 & 3\dg 19.0\pr & 24 & 23 & 18\dg 21.4\pr \\
25 & 0 & 33\dg 23.9\pr & 25 & 1 & 48\dg 26.3\pr & 25 & 2 & 63\dg 28.8\pr & 25 & 3 & 78\dg 31.3\pr \\
25 & 4 & 93\dg 33.7\pr & 25 & 5 & 108\dg 36.2\pr & 25 & 6 & 123\dg 38.7\pr & 25 & 7 & 138\dg 41.1\pr \\
25 & 8 & 153\dg 43.6\pr & 25 & 9 & 168\dg 46.1\pr & 25 & 10 & 183\dg 48.5\pr & 25 & 11 & 198\dg 51.0\pr \\
25 & 12 & 213\dg 53.5\pr & 25 & 13 & 228\dg 55.9\pr & 25 & 14 & 243\dg 58.4\pr & 25 & 15 & 259\dg 00.8\pr \\
25 & 16 & 274\dg 03.3\pr & 25 & 17 & 289\dg 05.8\pr & 25 & 18 & 304\dg 08.2\pr & 25 & 19 & 319\dg 10.7\pr \\
25 & 20 & 334\dg 13.2\pr & 25 & 21 & 349\dg 15.6\pr & 25 & 22 & 4\dg 18.1\pr & 25 & 23 & 19\dg 20.6\pr \\
26 & 0 & 34\dg 23.0\pr & 26 & 1 & 49\dg 25.5\pr & 26 & 2 & 64\dg 28.0\pr & 26 & 3 & 79\dg 30.4\pr \\
26 & 4 & 94\dg 32.9\pr & 26 & 5 & 109\dg 35.3\pr & 26 & 6 & 124\dg 37.8\pr & 26 & 7 & 139\dg 40.3\pr \\
26 & 8 & 154\dg 42.7\pr & 26 & 9 & 169\dg 45.2\pr & 26 & 10 & 184\dg 47.7\pr & 26 & 11 & 199\dg 50.1\pr \\
26 & 12 & 214\dg 52.6\pr & 26 & 13 & 229\dg 55.1\pr & 26 & 14 & 244\dg 57.5\pr & 26 & 15 & 260\dg 00.0\pr \\
26 & 16 & 275\dg 02.4\pr & 26 & 17 & 290\dg 04.9\pr & 26 & 18 & 305\dg 07.4\pr & 26 & 19 & 320\dg 09.8\pr \\
26 & 20 & 335\dg 12.3\pr & 26 & 21 & 350\dg 14.8\pr & 26 & 22 & 5\dg 17.2\pr & 26 & 23 & 20\dg 19.7\pr \\
27 & 0 & 35\dg 22.2\pr & 27 & 1 & 50\dg 24.6\pr & 27 & 2 & 65\dg 27.1\pr & 27 & 3 & 80\dg 29.6\pr \\
27 & 4 & 95\dg 32.0\pr & 27 & 5 & 110\dg 34.5\pr & 27 & 6 & 125\dg 36.9\pr & 27 & 7 & 140\dg 39.4\pr \\
27 & 8 & 155\dg 41.9\pr & 27 & 9 & 170\dg 44.3\pr & 27 & 10 & 185\dg 46.8\pr & 27 & 11 & 200\dg 49.3\pr \\
27 & 12 & 215\dg 51.7\pr & 27 & 13 & 230\dg 54.2\pr & 27 & 14 & 245\dg 56.7\pr & 27 & 15 & 260\dg 59.1\pr \\
27 & 16 & 276\dg 01.6\pr & 27 & 17 & 291\dg 04.1\pr & 27 & 18 & 306\dg 06.5\pr & 27 & 19 & 321\dg 09.0\pr \\
27 & 20 & 336\dg 11.4\pr & 27 & 21 & 351\dg 13.9\pr & 27 & 22 & 6\dg 16.4\pr & 27 & 23 & 21\dg 18.8\pr \\
28 & 0 & 36\dg 21.3\pr & 28 & 1 & 51\dg 23.8\pr & 28 & 2 & 66\dg 26.2\pr & 28 & 3 & 81\dg 28.7\pr \\
28 & 4 & 96\dg 31.2\pr & 28 & 5 & 111\dg 33.6\pr & 28 & 6 & 126\dg 36.1\pr & 28 & 7 & 141\dg 38.6\pr \\
28 & 8 & 156\dg 41.0\pr & 28 & 9 & 171\dg 43.5\pr & 28 & 10 & 186\dg 45.9\pr & 28 & 11 & 201\dg 48.4\pr \\
28 & 12 & 216\dg 50.9\pr & 28 & 13 & 231\dg 53.3\pr & 28 & 14 & 246\dg 55.8\pr & 28 & 15 & 261\dg 58.3\pr \\
28 & 16 & 277\dg 00.7\pr & 28 & 17 & 292\dg 03.2\pr & 28 & 18 & 307\dg 05.7\pr & 28 & 19 & 322\dg 08.1\pr \\
28 & 20 & 337\dg 10.6\pr & 28 & 21 & 352\dg 13.0\pr & 28 & 22 & 7\dg 15.5\pr & 28 & 23 & 22\dg 18.0\pr \\
29 & 0 & 37\dg 20.4\pr & 29 & 1 & 52\dg 22.9\pr & 29 & 2 & 67\dg 25.4\pr & 29 & 3 & 82\dg 27.8\pr \\
29 & 4 & 97\dg 30.3\pr & 29 & 5 & 112\dg 32.8\pr & 29 & 6 & 127\dg 35.2\pr & 29 & 7 & 142\dg 37.7\pr \\
29 & 8 & 157\dg 40.2\pr & 29 & 9 & 172\dg 42.6\pr & 29 & 10 & 187\dg 45.1\pr & 29 & 11 & 202\dg 47.5\pr \\
29 & 12 & 217\dg 50.0\pr & 29 & 13 & 232\dg 52.5\pr & 29 & 14 & 247\dg 54.9\pr & 29 & 15 & 262\dg 57.4\pr \\
29 & 16 & 277\dg 59.9\pr & 29 & 17 & 293\dg 02.3\pr & 29 & 18 & 308\dg 04.8\pr & 29 & 19 & 323\dg 07.3\pr \\
29 & 20 & 338\dg 09.7\pr & 29 & 21 & 353\dg 12.2\pr & 29 & 22 & 8\dg 14.7\pr & 29 & 23 & 23\dg 17.1\pr \\
30 & 0 & 38\dg 19.6\pr & 30 & 1 & 53\dg 22.0\pr & 30 & 2 & 68\dg 24.5\pr & 30 & 3 & 83\dg 27.0\pr \\
30 & 4 & 98\dg 29.4\pr & 30 & 5 & 113\dg 31.9\pr & 30 & 6 & 128\dg 34.4\pr & 30 & 7 & 143\dg 36.8\pr \\
30 & 8 & 158\dg 39.3\pr & 30 & 9 & 173\dg 41.8\pr & 30 & 10 & 188\dg 44.2\pr & 30 & 11 & 203\dg 46.7\pr \\
30 & 12 & 218\dg 49.1\pr & 30 & 13 & 233\dg 51.6\pr & 30 & 14 & 248\dg 54.1\pr & 30 & 15 & 263\dg 56.5\pr \\
30 & 16 & 278\dg 59.0\pr & 30 & 17 & 294\dg 01.5\pr & 30 & 18 & 309\dg 03.9\pr & 30 & 19 & 324\dg 06.4\pr \\
30 & 20 & 339\dg 08.9\pr & 30 & 21 & 354\dg 11.3\pr & 30 & 22 & 9\dg 13.8\pr & 30 & 23 & 24\dg 16.3\pr \\
31 & 0 & 39\dg 18.7\pr & 31 & 1 & 54\dg 21.2\pr & 31 & 2 & 69\dg 23.6\pr & 31 & 3 & 84\dg 26.1\pr \\
31 & 4 & 99\dg 28.6\pr & 31 & 5 & 114\dg 31.0\pr & 31 & 6 & 129\dg 33.5\pr & 31 & 7 & 144\dg 36.0\pr \\
31 & 8 & 159\dg 38.4\pr & 31 & 9 & 174\dg 40.9\pr & 31 & 10 & 189\dg 43.4\pr & 31 & 11 & 204\dg 45.8\pr \\
31 & 12 & 219\dg 48.3\pr & 31 & 13 & 234\dg 50.8\pr & 31 & 14 & 249\dg 53.2\pr & 31 & 15 & 264\dg 55.7\pr \\
31 & 16 & 279\dg 58.1\pr & 31 & 17 & 295\dg 00.6\pr & 31 & 18 & 310\dg 03.1\pr & 31 & 19 & 325\dg 05.5\pr \\
31 & 20 & 340\dg 08.0\pr & 31 & 21 & 355\dg 10.5\pr & 31 & 22 & 10\dg 12.9\pr & 31 & 23 & 25\dg 15.4\pr \\
\end{longtable}}

{\footnotesize
\begin{longtable}{rrr|rrr|rrr|rrr}
\caption{Aries --- GHA of the First Point of Aries, November 2026 (UT)}\label{aries11}\\
\toprule
\textbf{d} & \textbf{h} & \textbf{GHA} $\gamma$ & \textbf{d} & \textbf{h} & \textbf{GHA} $\gamma$ & \textbf{d} & \textbf{h} & \textbf{GHA} $\gamma$ & \textbf{d} & \textbf{h} & \textbf{GHA} $\gamma$ \\
\midrule
\endfirsthead
\multicolumn{12}{c}{\textit{Aries --- GHA of the First Point of Aries, November 2026 (UT) (continued)}}\\
\toprule
\textbf{d} & \textbf{h} & \textbf{GHA} $\gamma$ & \textbf{d} & \textbf{h} & \textbf{GHA} $\gamma$ & \textbf{d} & \textbf{h} & \textbf{GHA} $\gamma$ & \textbf{d} & \textbf{h} & \textbf{GHA} $\gamma$ \\
\midrule
\endhead
\midrule\multicolumn{12}{r}{\textit{continued on next page}}\\
\endfoot
\bottomrule
\endlastfoot
1 & 0 & 40\dg 17.9\pr & 1 & 1 & 55\dg 20.3\pr & 1 & 2 & 70\dg 22.8\pr & 1 & 3 & 85\dg 25.2\pr \\
1 & 4 & 100\dg 27.7\pr & 1 & 5 & 115\dg 30.2\pr & 1 & 6 & 130\dg 32.6\pr & 1 & 7 & 145\dg 35.1\pr \\
1 & 8 & 160\dg 37.6\pr & 1 & 9 & 175\dg 40.0\pr & 1 & 10 & 190\dg 42.5\pr & 1 & 11 & 205\dg 45.0\pr \\
1 & 12 & 220\dg 47.4\pr & 1 & 13 & 235\dg 49.9\pr & 1 & 14 & 250\dg 52.4\pr & 1 & 15 & 265\dg 54.8\pr \\
1 & 16 & 280\dg 57.3\pr & 1 & 17 & 295\dg 59.7\pr & 1 & 18 & 311\dg 02.2\pr & 1 & 19 & 326\dg 04.7\pr \\
1 & 20 & 341\dg 07.1\pr & 1 & 21 & 356\dg 09.6\pr & 1 & 22 & 11\dg 12.1\pr & 1 & 23 & 26\dg 14.5\pr \\
2 & 0 & 41\dg 17.0\pr & 2 & 1 & 56\dg 19.5\pr & 2 & 2 & 71\dg 21.9\pr & 2 & 3 & 86\dg 24.4\pr \\
2 & 4 & 101\dg 26.9\pr & 2 & 5 & 116\dg 29.3\pr & 2 & 6 & 131\dg 31.8\pr & 2 & 7 & 146\dg 34.2\pr \\
2 & 8 & 161\dg 36.7\pr & 2 & 9 & 176\dg 39.2\pr & 2 & 10 & 191\dg 41.6\pr & 2 & 11 & 206\dg 44.1\pr \\
2 & 12 & 221\dg 46.6\pr & 2 & 13 & 236\dg 49.0\pr & 2 & 14 & 251\dg 51.5\pr & 2 & 15 & 266\dg 54.0\pr \\
2 & 16 & 281\dg 56.4\pr & 2 & 17 & 296\dg 58.9\pr & 2 & 18 & 312\dg 01.4\pr & 2 & 19 & 327\dg 03.8\pr \\
2 & 20 & 342\dg 06.3\pr & 2 & 21 & 357\dg 08.7\pr & 2 & 22 & 12\dg 11.2\pr & 2 & 23 & 27\dg 13.7\pr \\
3 & 0 & 42\dg 16.1\pr & 3 & 1 & 57\dg 18.6\pr & 3 & 2 & 72\dg 21.1\pr & 3 & 3 & 87\dg 23.5\pr \\
3 & 4 & 102\dg 26.0\pr & 3 & 5 & 117\dg 28.5\pr & 3 & 6 & 132\dg 30.9\pr & 3 & 7 & 147\dg 33.4\pr \\
3 & 8 & 162\dg 35.8\pr & 3 & 9 & 177\dg 38.3\pr & 3 & 10 & 192\dg 40.8\pr & 3 & 11 & 207\dg 43.2\pr \\
3 & 12 & 222\dg 45.7\pr & 3 & 13 & 237\dg 48.2\pr & 3 & 14 & 252\dg 50.6\pr & 3 & 15 & 267\dg 53.1\pr \\
3 & 16 & 282\dg 55.6\pr & 3 & 17 & 297\dg 58.0\pr & 3 & 18 & 313\dg 00.5\pr & 3 & 19 & 328\dg 03.0\pr \\
3 & 20 & 343\dg 05.4\pr & 3 & 21 & 358\dg 07.9\pr & 3 & 22 & 13\dg 10.3\pr & 3 & 23 & 28\dg 12.8\pr \\
4 & 0 & 43\dg 15.3\pr & 4 & 1 & 58\dg 17.7\pr & 4 & 2 & 73\dg 20.2\pr & 4 & 3 & 88\dg 22.7\pr \\
4 & 4 & 103\dg 25.1\pr & 4 & 5 & 118\dg 27.6\pr & 4 & 6 & 133\dg 30.1\pr & 4 & 7 & 148\dg 32.5\pr \\
4 & 8 & 163\dg 35.0\pr & 4 & 9 & 178\dg 37.5\pr & 4 & 10 & 193\dg 39.9\pr & 4 & 11 & 208\dg 42.4\pr \\
4 & 12 & 223\dg 44.8\pr & 4 & 13 & 238\dg 47.3\pr & 4 & 14 & 253\dg 49.8\pr & 4 & 15 & 268\dg 52.2\pr \\
4 & 16 & 283\dg 54.7\pr & 4 & 17 & 298\dg 57.2\pr & 4 & 18 & 313\dg 59.6\pr & 4 & 19 & 329\dg 02.1\pr \\
4 & 20 & 344\dg 04.6\pr & 4 & 21 & 359\dg 07.0\pr & 4 & 22 & 14\dg 09.5\pr & 4 & 23 & 29\dg 11.9\pr \\
5 & 0 & 44\dg 14.4\pr & 5 & 1 & 59\dg 16.9\pr & 5 & 2 & 74\dg 19.3\pr & 5 & 3 & 89\dg 21.8\pr \\
5 & 4 & 104\dg 24.3\pr & 5 & 5 & 119\dg 26.7\pr & 5 & 6 & 134\dg 29.2\pr & 5 & 7 & 149\dg 31.7\pr \\
5 & 8 & 164\dg 34.1\pr & 5 & 9 & 179\dg 36.6\pr & 5 & 10 & 194\dg 39.1\pr & 5 & 11 & 209\dg 41.5\pr \\
5 & 12 & 224\dg 44.0\pr & 5 & 13 & 239\dg 46.4\pr & 5 & 14 & 254\dg 48.9\pr & 5 & 15 & 269\dg 51.4\pr \\
5 & 16 & 284\dg 53.8\pr & 5 & 17 & 299\dg 56.3\pr & 5 & 18 & 314\dg 58.8\pr & 5 & 19 & 330\dg 01.2\pr \\
5 & 20 & 345\dg 03.7\pr & 5 & 21 & 0\dg 06.2\pr & 5 & 22 & 15\dg 08.6\pr & 5 & 23 & 30\dg 11.1\pr \\
6 & 0 & 45\dg 13.6\pr & 6 & 1 & 60\dg 16.0\pr & 6 & 2 & 75\dg 18.5\pr & 6 & 3 & 90\dg 20.9\pr \\
6 & 4 & 105\dg 23.4\pr & 6 & 5 & 120\dg 25.9\pr & 6 & 6 & 135\dg 28.3\pr & 6 & 7 & 150\dg 30.8\pr \\
6 & 8 & 165\dg 33.3\pr & 6 & 9 & 180\dg 35.7\pr & 6 & 10 & 195\dg 38.2\pr & 6 & 11 & 210\dg 40.7\pr \\
6 & 12 & 225\dg 43.1\pr & 6 & 13 & 240\dg 45.6\pr & 6 & 14 & 255\dg 48.0\pr & 6 & 15 & 270\dg 50.5\pr \\
6 & 16 & 285\dg 53.0\pr & 6 & 17 & 300\dg 55.4\pr & 6 & 18 & 315\dg 57.9\pr & 6 & 19 & 331\dg 00.4\pr \\
6 & 20 & 346\dg 02.8\pr & 6 & 21 & 1\dg 05.3\pr & 6 & 22 & 16\dg 07.8\pr & 6 & 23 & 31\dg 10.2\pr \\
7 & 0 & 46\dg 12.7\pr & 7 & 1 & 61\dg 15.2\pr & 7 & 2 & 76\dg 17.6\pr & 7 & 3 & 91\dg 20.1\pr \\
7 & 4 & 106\dg 22.5\pr & 7 & 5 & 121\dg 25.0\pr & 7 & 6 & 136\dg 27.5\pr & 7 & 7 & 151\dg 29.9\pr \\
7 & 8 & 166\dg 32.4\pr & 7 & 9 & 181\dg 34.9\pr & 7 & 10 & 196\dg 37.3\pr & 7 & 11 & 211\dg 39.8\pr \\
7 & 12 & 226\dg 42.3\pr & 7 & 13 & 241\dg 44.7\pr & 7 & 14 & 256\dg 47.2\pr & 7 & 15 & 271\dg 49.7\pr \\
7 & 16 & 286\dg 52.1\pr & 7 & 17 & 301\dg 54.6\pr & 7 & 18 & 316\dg 57.0\pr & 7 & 19 & 331\dg 59.5\pr \\
7 & 20 & 347\dg 02.0\pr & 7 & 21 & 2\dg 04.4\pr & 7 & 22 & 17\dg 06.9\pr & 7 & 23 & 32\dg 09.4\pr \\
8 & 0 & 47\dg 11.8\pr & 8 & 1 & 62\dg 14.3\pr & 8 & 2 & 77\dg 16.8\pr & 8 & 3 & 92\dg 19.2\pr \\
8 & 4 & 107\dg 21.7\pr & 8 & 5 & 122\dg 24.1\pr & 8 & 6 & 137\dg 26.6\pr & 8 & 7 & 152\dg 29.1\pr \\
8 & 8 & 167\dg 31.5\pr & 8 & 9 & 182\dg 34.0\pr & 8 & 10 & 197\dg 36.5\pr & 8 & 11 & 212\dg 38.9\pr \\
8 & 12 & 227\dg 41.4\pr & 8 & 13 & 242\dg 43.9\pr & 8 & 14 & 257\dg 46.3\pr & 8 & 15 & 272\dg 48.8\pr \\
8 & 16 & 287\dg 51.3\pr & 8 & 17 & 302\dg 53.7\pr & 8 & 18 & 317\dg 56.2\pr & 8 & 19 & 332\dg 58.6\pr \\
8 & 20 & 348\dg 01.1\pr & 8 & 21 & 3\dg 03.6\pr & 8 & 22 & 18\dg 06.0\pr & 8 & 23 & 33\dg 08.5\pr \\
9 & 0 & 48\dg 11.0\pr & 9 & 1 & 63\dg 13.4\pr & 9 & 2 & 78\dg 15.9\pr & 9 & 3 & 93\dg 18.4\pr \\
9 & 4 & 108\dg 20.8\pr & 9 & 5 & 123\dg 23.3\pr & 9 & 6 & 138\dg 25.8\pr & 9 & 7 & 153\dg 28.2\pr \\
9 & 8 & 168\dg 30.7\pr & 9 & 9 & 183\dg 33.1\pr & 9 & 10 & 198\dg 35.6\pr & 9 & 11 & 213\dg 38.1\pr \\
9 & 12 & 228\dg 40.5\pr & 9 & 13 & 243\dg 43.0\pr & 9 & 14 & 258\dg 45.5\pr & 9 & 15 & 273\dg 47.9\pr \\
9 & 16 & 288\dg 50.4\pr & 9 & 17 & 303\dg 52.9\pr & 9 & 18 & 318\dg 55.3\pr & 9 & 19 & 333\dg 57.8\pr \\
9 & 20 & 349\dg 00.3\pr & 9 & 21 & 4\dg 02.7\pr & 9 & 22 & 19\dg 05.2\pr & 9 & 23 & 34\dg 07.6\pr \\
10 & 0 & 49\dg 10.1\pr & 10 & 1 & 64\dg 12.6\pr & 10 & 2 & 79\dg 15.0\pr & 10 & 3 & 94\dg 17.5\pr \\
10 & 4 & 109\dg 20.0\pr & 10 & 5 & 124\dg 22.4\pr & 10 & 6 & 139\dg 24.9\pr & 10 & 7 & 154\dg 27.4\pr \\
10 & 8 & 169\dg 29.8\pr & 10 & 9 & 184\dg 32.3\pr & 10 & 10 & 199\dg 34.7\pr & 10 & 11 & 214\dg 37.2\pr \\
10 & 12 & 229\dg 39.7\pr & 10 & 13 & 244\dg 42.1\pr & 10 & 14 & 259\dg 44.6\pr & 10 & 15 & 274\dg 47.1\pr \\
10 & 16 & 289\dg 49.5\pr & 10 & 17 & 304\dg 52.0\pr & 10 & 18 & 319\dg 54.5\pr & 10 & 19 & 334\dg 56.9\pr \\
10 & 20 & 349\dg 59.4\pr & 10 & 21 & 5\dg 01.9\pr & 10 & 22 & 20\dg 04.3\pr & 10 & 23 & 35\dg 06.8\pr \\
11 & 0 & 50\dg 09.2\pr & 11 & 1 & 65\dg 11.7\pr & 11 & 2 & 80\dg 14.2\pr & 11 & 3 & 95\dg 16.6\pr \\
11 & 4 & 110\dg 19.1\pr & 11 & 5 & 125\dg 21.6\pr & 11 & 6 & 140\dg 24.0\pr & 11 & 7 & 155\dg 26.5\pr \\
11 & 8 & 170\dg 29.0\pr & 11 & 9 & 185\dg 31.4\pr & 11 & 10 & 200\dg 33.9\pr & 11 & 11 & 215\dg 36.4\pr \\
11 & 12 & 230\dg 38.8\pr & 11 & 13 & 245\dg 41.3\pr & 11 & 14 & 260\dg 43.7\pr & 11 & 15 & 275\dg 46.2\pr \\
11 & 16 & 290\dg 48.7\pr & 11 & 17 & 305\dg 51.1\pr & 11 & 18 & 320\dg 53.6\pr & 11 & 19 & 335\dg 56.1\pr \\
11 & 20 & 350\dg 58.5\pr & 11 & 21 & 6\dg 01.0\pr & 11 & 22 & 21\dg 03.5\pr & 11 & 23 & 36\dg 05.9\pr \\
12 & 0 & 51\dg 08.4\pr & 12 & 1 & 66\dg 10.8\pr & 12 & 2 & 81\dg 13.3\pr & 12 & 3 & 96\dg 15.8\pr \\
12 & 4 & 111\dg 18.2\pr & 12 & 5 & 126\dg 20.7\pr & 12 & 6 & 141\dg 23.2\pr & 12 & 7 & 156\dg 25.6\pr \\
12 & 8 & 171\dg 28.1\pr & 12 & 9 & 186\dg 30.6\pr & 12 & 10 & 201\dg 33.0\pr & 12 & 11 & 216\dg 35.5\pr \\
12 & 12 & 231\dg 38.0\pr & 12 & 13 & 246\dg 40.4\pr & 12 & 14 & 261\dg 42.9\pr & 12 & 15 & 276\dg 45.3\pr \\
12 & 16 & 291\dg 47.8\pr & 12 & 17 & 306\dg 50.3\pr & 12 & 18 & 321\dg 52.7\pr & 12 & 19 & 336\dg 55.2\pr \\
12 & 20 & 351\dg 57.7\pr & 12 & 21 & 7\dg 00.1\pr & 12 & 22 & 22\dg 02.6\pr & 12 & 23 & 37\dg 05.1\pr \\
13 & 0 & 52\dg 07.5\pr & 13 & 1 & 67\dg 10.0\pr & 13 & 2 & 82\dg 12.5\pr & 13 & 3 & 97\dg 14.9\pr \\
13 & 4 & 112\dg 17.4\pr & 13 & 5 & 127\dg 19.8\pr & 13 & 6 & 142\dg 22.3\pr & 13 & 7 & 157\dg 24.8\pr \\
13 & 8 & 172\dg 27.2\pr & 13 & 9 & 187\dg 29.7\pr & 13 & 10 & 202\dg 32.2\pr & 13 & 11 & 217\dg 34.6\pr \\
13 & 12 & 232\dg 37.1\pr & 13 & 13 & 247\dg 39.6\pr & 13 & 14 & 262\dg 42.0\pr & 13 & 15 & 277\dg 44.5\pr \\
13 & 16 & 292\dg 46.9\pr & 13 & 17 & 307\dg 49.4\pr & 13 & 18 & 322\dg 51.9\pr & 13 & 19 & 337\dg 54.3\pr \\
13 & 20 & 352\dg 56.8\pr & 13 & 21 & 7\dg 59.3\pr & 13 & 22 & 23\dg 01.7\pr & 13 & 23 & 38\dg 04.2\pr \\
14 & 0 & 53\dg 06.7\pr & 14 & 1 & 68\dg 09.1\pr & 14 & 2 & 83\dg 11.6\pr & 14 & 3 & 98\dg 14.1\pr \\
14 & 4 & 113\dg 16.5\pr & 14 & 5 & 128\dg 19.0\pr & 14 & 6 & 143\dg 21.4\pr & 14 & 7 & 158\dg 23.9\pr \\
14 & 8 & 173\dg 26.4\pr & 14 & 9 & 188\dg 28.8\pr & 14 & 10 & 203\dg 31.3\pr & 14 & 11 & 218\dg 33.8\pr \\
14 & 12 & 233\dg 36.2\pr & 14 & 13 & 248\dg 38.7\pr & 14 & 14 & 263\dg 41.2\pr & 14 & 15 & 278\dg 43.6\pr \\
14 & 16 & 293\dg 46.1\pr & 14 & 17 & 308\dg 48.6\pr & 14 & 18 & 323\dg 51.0\pr & 14 & 19 & 338\dg 53.5\pr \\
14 & 20 & 353\dg 55.9\pr & 14 & 21 & 8\dg 58.4\pr & 14 & 22 & 24\dg 00.9\pr & 14 & 23 & 39\dg 03.3\pr \\
15 & 0 & 54\dg 05.8\pr & 15 & 1 & 69\dg 08.3\pr & 15 & 2 & 84\dg 10.7\pr & 15 & 3 & 99\dg 13.2\pr \\
15 & 4 & 114\dg 15.7\pr & 15 & 5 & 129\dg 18.1\pr & 15 & 6 & 144\dg 20.6\pr & 15 & 7 & 159\dg 23.0\pr \\
15 & 8 & 174\dg 25.5\pr & 15 & 9 & 189\dg 28.0\pr & 15 & 10 & 204\dg 30.4\pr & 15 & 11 & 219\dg 32.9\pr \\
15 & 12 & 234\dg 35.4\pr & 15 & 13 & 249\dg 37.8\pr & 15 & 14 & 264\dg 40.3\pr & 15 & 15 & 279\dg 42.8\pr \\
15 & 16 & 294\dg 45.2\pr & 15 & 17 & 309\dg 47.7\pr & 15 & 18 & 324\dg 50.2\pr & 15 & 19 & 339\dg 52.6\pr \\
15 & 20 & 354\dg 55.1\pr & 15 & 21 & 9\dg 57.5\pr & 15 & 22 & 25\dg 00.0\pr & 15 & 23 & 40\dg 02.5\pr \\
16 & 0 & 55\dg 04.9\pr & 16 & 1 & 70\dg 07.4\pr & 16 & 2 & 85\dg 09.9\pr & 16 & 3 & 100\dg 12.3\pr \\
16 & 4 & 115\dg 14.8\pr & 16 & 5 & 130\dg 17.3\pr & 16 & 6 & 145\dg 19.7\pr & 16 & 7 & 160\dg 22.2\pr \\
16 & 8 & 175\dg 24.7\pr & 16 & 9 & 190\dg 27.1\pr & 16 & 10 & 205\dg 29.6\pr & 16 & 11 & 220\dg 32.0\pr \\
16 & 12 & 235\dg 34.5\pr & 16 & 13 & 250\dg 37.0\pr & 16 & 14 & 265\dg 39.4\pr & 16 & 15 & 280\dg 41.9\pr \\
16 & 16 & 295\dg 44.4\pr & 16 & 17 & 310\dg 46.8\pr & 16 & 18 & 325\dg 49.3\pr & 16 & 19 & 340\dg 51.8\pr \\
16 & 20 & 355\dg 54.2\pr & 16 & 21 & 10\dg 56.7\pr & 16 & 22 & 25\dg 59.2\pr & 16 & 23 & 41\dg 01.6\pr \\
17 & 0 & 56\dg 04.1\pr & 17 & 1 & 71\dg 06.5\pr & 17 & 2 & 86\dg 09.0\pr & 17 & 3 & 101\dg 11.5\pr \\
17 & 4 & 116\dg 13.9\pr & 17 & 5 & 131\dg 16.4\pr & 17 & 6 & 146\dg 18.9\pr & 17 & 7 & 161\dg 21.3\pr \\
17 & 8 & 176\dg 23.8\pr & 17 & 9 & 191\dg 26.3\pr & 17 & 10 & 206\dg 28.7\pr & 17 & 11 & 221\dg 31.2\pr \\
17 & 12 & 236\dg 33.6\pr & 17 & 13 & 251\dg 36.1\pr & 17 & 14 & 266\dg 38.6\pr & 17 & 15 & 281\dg 41.0\pr \\
17 & 16 & 296\dg 43.5\pr & 17 & 17 & 311\dg 46.0\pr & 17 & 18 & 326\dg 48.4\pr & 17 & 19 & 341\dg 50.9\pr \\
17 & 20 & 356\dg 53.4\pr & 17 & 21 & 11\dg 55.8\pr & 17 & 22 & 26\dg 58.3\pr & 17 & 23 & 42\dg 00.8\pr \\
18 & 0 & 57\dg 03.2\pr & 18 & 1 & 72\dg 05.7\pr & 18 & 2 & 87\dg 08.1\pr & 18 & 3 & 102\dg 10.6\pr \\
18 & 4 & 117\dg 13.1\pr & 18 & 5 & 132\dg 15.5\pr & 18 & 6 & 147\dg 18.0\pr & 18 & 7 & 162\dg 20.5\pr \\
18 & 8 & 177\dg 22.9\pr & 18 & 9 & 192\dg 25.4\pr & 18 & 10 & 207\dg 27.9\pr & 18 & 11 & 222\dg 30.3\pr \\
18 & 12 & 237\dg 32.8\pr & 18 & 13 & 252\dg 35.3\pr & 18 & 14 & 267\dg 37.7\pr & 18 & 15 & 282\dg 40.2\pr \\
18 & 16 & 297\dg 42.6\pr & 18 & 17 & 312\dg 45.1\pr & 18 & 18 & 327\dg 47.6\pr & 18 & 19 & 342\dg 50.0\pr \\
18 & 20 & 357\dg 52.5\pr & 18 & 21 & 12\dg 55.0\pr & 18 & 22 & 27\dg 57.4\pr & 18 & 23 & 42\dg 59.9\pr \\
19 & 0 & 58\dg 02.4\pr & 19 & 1 & 73\dg 04.8\pr & 19 & 2 & 88\dg 07.3\pr & 19 & 3 & 103\dg 09.7\pr \\
19 & 4 & 118\dg 12.2\pr & 19 & 5 & 133\dg 14.7\pr & 19 & 6 & 148\dg 17.1\pr & 19 & 7 & 163\dg 19.6\pr \\
19 & 8 & 178\dg 22.1\pr & 19 & 9 & 193\dg 24.5\pr & 19 & 10 & 208\dg 27.0\pr & 19 & 11 & 223\dg 29.5\pr \\
19 & 12 & 238\dg 31.9\pr & 19 & 13 & 253\dg 34.4\pr & 19 & 14 & 268\dg 36.9\pr & 19 & 15 & 283\dg 39.3\pr \\
19 & 16 & 298\dg 41.8\pr & 19 & 17 & 313\dg 44.2\pr & 19 & 18 & 328\dg 46.7\pr & 19 & 19 & 343\dg 49.2\pr \\
19 & 20 & 358\dg 51.6\pr & 19 & 21 & 13\dg 54.1\pr & 19 & 22 & 28\dg 56.6\pr & 19 & 23 & 43\dg 59.0\pr \\
20 & 0 & 59\dg 01.5\pr & 20 & 1 & 74\dg 04.0\pr & 20 & 2 & 89\dg 06.4\pr & 20 & 3 & 104\dg 08.9\pr \\
20 & 4 & 119\dg 11.4\pr & 20 & 5 & 134\dg 13.8\pr & 20 & 6 & 149\dg 16.3\pr & 20 & 7 & 164\dg 18.7\pr \\
20 & 8 & 179\dg 21.2\pr & 20 & 9 & 194\dg 23.7\pr & 20 & 10 & 209\dg 26.1\pr & 20 & 11 & 224\dg 28.6\pr \\
20 & 12 & 239\dg 31.1\pr & 20 & 13 & 254\dg 33.5\pr & 20 & 14 & 269\dg 36.0\pr & 20 & 15 & 284\dg 38.5\pr \\
20 & 16 & 299\dg 40.9\pr & 20 & 17 & 314\dg 43.4\pr & 20 & 18 & 329\dg 45.8\pr & 20 & 19 & 344\dg 48.3\pr \\
20 & 20 & 359\dg 50.8\pr & 20 & 21 & 14\dg 53.2\pr & 20 & 22 & 29\dg 55.7\pr & 20 & 23 & 44\dg 58.2\pr \\
21 & 0 & 60\dg 00.6\pr & 21 & 1 & 75\dg 03.1\pr & 21 & 2 & 90\dg 05.6\pr & 21 & 3 & 105\dg 08.0\pr \\
21 & 4 & 120\dg 10.5\pr & 21 & 5 & 135\dg 13.0\pr & 21 & 6 & 150\dg 15.4\pr & 21 & 7 & 165\dg 17.9\pr \\
21 & 8 & 180\dg 20.3\pr & 21 & 9 & 195\dg 22.8\pr & 21 & 10 & 210\dg 25.3\pr & 21 & 11 & 225\dg 27.7\pr \\
21 & 12 & 240\dg 30.2\pr & 21 & 13 & 255\dg 32.7\pr & 21 & 14 & 270\dg 35.1\pr & 21 & 15 & 285\dg 37.6\pr \\
21 & 16 & 300\dg 40.1\pr & 21 & 17 & 315\dg 42.5\pr & 21 & 18 & 330\dg 45.0\pr & 21 & 19 & 345\dg 47.5\pr \\
21 & 20 & 0\dg 49.9\pr & 21 & 21 & 15\dg 52.4\pr & 21 & 22 & 30\dg 54.8\pr & 21 & 23 & 45\dg 57.3\pr \\
22 & 0 & 60\dg 59.8\pr & 22 & 1 & 76\dg 02.2\pr & 22 & 2 & 91\dg 04.7\pr & 22 & 3 & 106\dg 07.2\pr \\
22 & 4 & 121\dg 09.6\pr & 22 & 5 & 136\dg 12.1\pr & 22 & 6 & 151\dg 14.6\pr & 22 & 7 & 166\dg 17.0\pr \\
22 & 8 & 181\dg 19.5\pr & 22 & 9 & 196\dg 21.9\pr & 22 & 10 & 211\dg 24.4\pr & 22 & 11 & 226\dg 26.9\pr \\
22 & 12 & 241\dg 29.3\pr & 22 & 13 & 256\dg 31.8\pr & 22 & 14 & 271\dg 34.3\pr & 22 & 15 & 286\dg 36.7\pr \\
22 & 16 & 301\dg 39.2\pr & 22 & 17 & 316\dg 41.7\pr & 22 & 18 & 331\dg 44.1\pr & 22 & 19 & 346\dg 46.6\pr \\
22 & 20 & 1\dg 49.1\pr & 22 & 21 & 16\dg 51.5\pr & 22 & 22 & 31\dg 54.0\pr & 22 & 23 & 46\dg 56.4\pr \\
23 & 0 & 61\dg 58.9\pr & 23 & 1 & 77\dg 01.4\pr & 23 & 2 & 92\dg 03.8\pr & 23 & 3 & 107\dg 06.3\pr \\
23 & 4 & 122\dg 08.8\pr & 23 & 5 & 137\dg 11.2\pr & 23 & 6 & 152\dg 13.7\pr & 23 & 7 & 167\dg 16.2\pr \\
23 & 8 & 182\dg 18.6\pr & 23 & 9 & 197\dg 21.1\pr & 23 & 10 & 212\dg 23.6\pr & 23 & 11 & 227\dg 26.0\pr \\
23 & 12 & 242\dg 28.5\pr & 23 & 13 & 257\dg 30.9\pr & 23 & 14 & 272\dg 33.4\pr & 23 & 15 & 287\dg 35.9\pr \\
23 & 16 & 302\dg 38.3\pr & 23 & 17 & 317\dg 40.8\pr & 23 & 18 & 332\dg 43.3\pr & 23 & 19 & 347\dg 45.7\pr \\
23 & 20 & 2\dg 48.2\pr & 23 & 21 & 17\dg 50.7\pr & 23 & 22 & 32\dg 53.1\pr & 23 & 23 & 47\dg 55.6\pr \\
24 & 0 & 62\dg 58.1\pr & 24 & 1 & 78\dg 00.5\pr & 24 & 2 & 93\dg 03.0\pr & 24 & 3 & 108\dg 05.4\pr \\
24 & 4 & 123\dg 07.9\pr & 24 & 5 & 138\dg 10.4\pr & 24 & 6 & 153\dg 12.8\pr & 24 & 7 & 168\dg 15.3\pr \\
24 & 8 & 183\dg 17.8\pr & 24 & 9 & 198\dg 20.2\pr & 24 & 10 & 213\dg 22.7\pr & 24 & 11 & 228\dg 25.2\pr \\
24 & 12 & 243\dg 27.6\pr & 24 & 13 & 258\dg 30.1\pr & 24 & 14 & 273\dg 32.5\pr & 24 & 15 & 288\dg 35.0\pr \\
24 & 16 & 303\dg 37.5\pr & 24 & 17 & 318\dg 39.9\pr & 24 & 18 & 333\dg 42.4\pr & 24 & 19 & 348\dg 44.9\pr \\
24 & 20 & 3\dg 47.3\pr & 24 & 21 & 18\dg 49.8\pr & 24 & 22 & 33\dg 52.3\pr & 24 & 23 & 48\dg 54.7\pr \\
25 & 0 & 63\dg 57.2\pr & 25 & 1 & 78\dg 59.7\pr & 25 & 2 & 94\dg 02.1\pr & 25 & 3 & 109\dg 04.6\pr \\
25 & 4 & 124\dg 07.0\pr & 25 & 5 & 139\dg 09.5\pr & 25 & 6 & 154\dg 12.0\pr & 25 & 7 & 169\dg 14.4\pr \\
25 & 8 & 184\dg 16.9\pr & 25 & 9 & 199\dg 19.4\pr & 25 & 10 & 214\dg 21.8\pr & 25 & 11 & 229\dg 24.3\pr \\
25 & 12 & 244\dg 26.8\pr & 25 & 13 & 259\dg 29.2\pr & 25 & 14 & 274\dg 31.7\pr & 25 & 15 & 289\dg 34.2\pr \\
25 & 16 & 304\dg 36.6\pr & 25 & 17 & 319\dg 39.1\pr & 25 & 18 & 334\dg 41.5\pr & 25 & 19 & 349\dg 44.0\pr \\
25 & 20 & 4\dg 46.5\pr & 25 & 21 & 19\dg 48.9\pr & 25 & 22 & 34\dg 51.4\pr & 25 & 23 & 49\dg 53.9\pr \\
26 & 0 & 64\dg 56.3\pr & 26 & 1 & 79\dg 58.8\pr & 26 & 2 & 95\dg 01.3\pr & 26 & 3 & 110\dg 03.7\pr \\
26 & 4 & 125\dg 06.2\pr & 26 & 5 & 140\dg 08.6\pr & 26 & 6 & 155\dg 11.1\pr & 26 & 7 & 170\dg 13.6\pr \\
26 & 8 & 185\dg 16.0\pr & 26 & 9 & 200\dg 18.5\pr & 26 & 10 & 215\dg 21.0\pr & 26 & 11 & 230\dg 23.4\pr \\
26 & 12 & 245\dg 25.9\pr & 26 & 13 & 260\dg 28.4\pr & 26 & 14 & 275\dg 30.8\pr & 26 & 15 & 290\dg 33.3\pr \\
26 & 16 & 305\dg 35.8\pr & 26 & 17 & 320\dg 38.2\pr & 26 & 18 & 335\dg 40.7\pr & 26 & 19 & 350\dg 43.1\pr \\
26 & 20 & 5\dg 45.6\pr & 26 & 21 & 20\dg 48.1\pr & 26 & 22 & 35\dg 50.5\pr & 26 & 23 & 50\dg 53.0\pr \\
27 & 0 & 65\dg 55.5\pr & 27 & 1 & 80\dg 57.9\pr & 27 & 2 & 96\dg 00.4\pr & 27 & 3 & 111\dg 02.9\pr \\
27 & 4 & 126\dg 05.3\pr & 27 & 5 & 141\dg 07.8\pr & 27 & 6 & 156\dg 10.3\pr & 27 & 7 & 171\dg 12.7\pr \\
27 & 8 & 186\dg 15.2\pr & 27 & 9 & 201\dg 17.6\pr & 27 & 10 & 216\dg 20.1\pr & 27 & 11 & 231\dg 22.6\pr \\
27 & 12 & 246\dg 25.0\pr & 27 & 13 & 261\dg 27.5\pr & 27 & 14 & 276\dg 30.0\pr & 27 & 15 & 291\dg 32.4\pr \\
27 & 16 & 306\dg 34.9\pr & 27 & 17 & 321\dg 37.4\pr & 27 & 18 & 336\dg 39.8\pr & 27 & 19 & 351\dg 42.3\pr \\
27 & 20 & 6\dg 44.7\pr & 27 & 21 & 21\dg 47.2\pr & 27 & 22 & 36\dg 49.7\pr & 27 & 23 & 51\dg 52.1\pr \\
28 & 0 & 66\dg 54.6\pr & 28 & 1 & 81\dg 57.1\pr & 28 & 2 & 96\dg 59.5\pr & 28 & 3 & 112\dg 02.0\pr \\
28 & 4 & 127\dg 04.5\pr & 28 & 5 & 142\dg 06.9\pr & 28 & 6 & 157\dg 09.4\pr & 28 & 7 & 172\dg 11.9\pr \\
28 & 8 & 187\dg 14.3\pr & 28 & 9 & 202\dg 16.8\pr & 28 & 10 & 217\dg 19.2\pr & 28 & 11 & 232\dg 21.7\pr \\
28 & 12 & 247\dg 24.2\pr & 28 & 13 & 262\dg 26.6\pr & 28 & 14 & 277\dg 29.1\pr & 28 & 15 & 292\dg 31.6\pr \\
28 & 16 & 307\dg 34.0\pr & 28 & 17 & 322\dg 36.5\pr & 28 & 18 & 337\dg 39.0\pr & 28 & 19 & 352\dg 41.4\pr \\
28 & 20 & 7\dg 43.9\pr & 28 & 21 & 22\dg 46.4\pr & 28 & 22 & 37\dg 48.8\pr & 28 & 23 & 52\dg 51.3\pr \\
29 & 0 & 67\dg 53.7\pr & 29 & 1 & 82\dg 56.2\pr & 29 & 2 & 97\dg 58.7\pr & 29 & 3 & 113\dg 01.1\pr \\
29 & 4 & 128\dg 03.6\pr & 29 & 5 & 143\dg 06.1\pr & 29 & 6 & 158\dg 08.5\pr & 29 & 7 & 173\dg 11.0\pr \\
29 & 8 & 188\dg 13.5\pr & 29 & 9 & 203\dg 15.9\pr & 29 & 10 & 218\dg 18.4\pr & 29 & 11 & 233\dg 20.8\pr \\
29 & 12 & 248\dg 23.3\pr & 29 & 13 & 263\dg 25.8\pr & 29 & 14 & 278\dg 28.2\pr & 29 & 15 & 293\dg 30.7\pr \\
29 & 16 & 308\dg 33.2\pr & 29 & 17 & 323\dg 35.6\pr & 29 & 18 & 338\dg 38.1\pr & 29 & 19 & 353\dg 40.6\pr \\
29 & 20 & 8\dg 43.0\pr & 29 & 21 & 23\dg 45.5\pr & 29 & 22 & 38\dg 48.0\pr & 29 & 23 & 53\dg 50.4\pr \\
30 & 0 & 68\dg 52.9\pr & 30 & 1 & 83\dg 55.3\pr & 30 & 2 & 98\dg 57.8\pr & 30 & 3 & 114\dg 00.3\pr \\
30 & 4 & 129\dg 02.7\pr & 30 & 5 & 144\dg 05.2\pr & 30 & 6 & 159\dg 07.7\pr & 30 & 7 & 174\dg 10.1\pr \\
30 & 8 & 189\dg 12.6\pr & 30 & 9 & 204\dg 15.1\pr & 30 & 10 & 219\dg 17.5\pr & 30 & 11 & 234\dg 20.0\pr \\
30 & 12 & 249\dg 22.5\pr & 30 & 13 & 264\dg 24.9\pr & 30 & 14 & 279\dg 27.4\pr & 30 & 15 & 294\dg 29.8\pr \\
30 & 16 & 309\dg 32.3\pr & 30 & 17 & 324\dg 34.8\pr & 30 & 18 & 339\dg 37.2\pr & 30 & 19 & 354\dg 39.7\pr \\
30 & 20 & 9\dg 42.2\pr & 30 & 21 & 24\dg 44.6\pr & 30 & 22 & 39\dg 47.1\pr & 30 & 23 & 54\dg 49.6\pr \\
\end{longtable}}

{\footnotesize
\begin{longtable}{rrr|rrr|rrr|rrr}
\caption{Aries --- GHA of the First Point of Aries, December 2026 (UT)}\label{aries12}\\
\toprule
\textbf{d} & \textbf{h} & \textbf{GHA} $\gamma$ & \textbf{d} & \textbf{h} & \textbf{GHA} $\gamma$ & \textbf{d} & \textbf{h} & \textbf{GHA} $\gamma$ & \textbf{d} & \textbf{h} & \textbf{GHA} $\gamma$ \\
\midrule
\endfirsthead
\multicolumn{12}{c}{\textit{Aries --- GHA of the First Point of Aries, December 2026 (UT) (continued)}}\\
\toprule
\textbf{d} & \textbf{h} & \textbf{GHA} $\gamma$ & \textbf{d} & \textbf{h} & \textbf{GHA} $\gamma$ & \textbf{d} & \textbf{h} & \textbf{GHA} $\gamma$ & \textbf{d} & \textbf{h} & \textbf{GHA} $\gamma$ \\
\midrule
\endhead
\midrule\multicolumn{12}{r}{\textit{continued on next page}}\\
\endfoot
\bottomrule
\endlastfoot
1 & 0 & 69\dg 52.0\pr & 1 & 1 & 84\dg 54.5\pr & 1 & 2 & 99\dg 57.0\pr & 1 & 3 & 114\dg 59.4\pr \\
1 & 4 & 130\dg 01.9\pr & 1 & 5 & 145\dg 04.3\pr & 1 & 6 & 160\dg 06.8\pr & 1 & 7 & 175\dg 09.3\pr \\
1 & 8 & 190\dg 11.7\pr & 1 & 9 & 205\dg 14.2\pr & 1 & 10 & 220\dg 16.7\pr & 1 & 11 & 235\dg 19.1\pr \\
1 & 12 & 250\dg 21.6\pr & 1 & 13 & 265\dg 24.1\pr & 1 & 14 & 280\dg 26.5\pr & 1 & 15 & 295\dg 29.0\pr \\
1 & 16 & 310\dg 31.4\pr & 1 & 17 & 325\dg 33.9\pr & 1 & 18 & 340\dg 36.4\pr & 1 & 19 & 355\dg 38.8\pr \\
1 & 20 & 10\dg 41.3\pr & 1 & 21 & 25\dg 43.8\pr & 1 & 22 & 40\dg 46.2\pr & 1 & 23 & 55\dg 48.7\pr \\
2 & 0 & 70\dg 51.2\pr & 2 & 1 & 85\dg 53.6\pr & 2 & 2 & 100\dg 56.1\pr & 2 & 3 & 115\dg 58.6\pr \\
2 & 4 & 131\dg 01.0\pr & 2 & 5 & 146\dg 03.5\pr & 2 & 6 & 161\dg 05.9\pr & 2 & 7 & 176\dg 08.4\pr \\
2 & 8 & 191\dg 10.9\pr & 2 & 9 & 206\dg 13.3\pr & 2 & 10 & 221\dg 15.8\pr & 2 & 11 & 236\dg 18.3\pr \\
2 & 12 & 251\dg 20.7\pr & 2 & 13 & 266\dg 23.2\pr & 2 & 14 & 281\dg 25.7\pr & 2 & 15 & 296\dg 28.1\pr \\
2 & 16 & 311\dg 30.6\pr & 2 & 17 & 326\dg 33.1\pr & 2 & 18 & 341\dg 35.5\pr & 2 & 19 & 356\dg 38.0\pr \\
2 & 20 & 11\dg 40.4\pr & 2 & 21 & 26\dg 42.9\pr & 2 & 22 & 41\dg 45.4\pr & 2 & 23 & 56\dg 47.8\pr \\
3 & 0 & 71\dg 50.3\pr & 3 & 1 & 86\dg 52.8\pr & 3 & 2 & 101\dg 55.2\pr & 3 & 3 & 116\dg 57.7\pr \\
3 & 4 & 132\dg 00.2\pr & 3 & 5 & 147\dg 02.6\pr & 3 & 6 & 162\dg 05.1\pr & 3 & 7 & 177\dg 07.5\pr \\
3 & 8 & 192\dg 10.0\pr & 3 & 9 & 207\dg 12.5\pr & 3 & 10 & 222\dg 14.9\pr & 3 & 11 & 237\dg 17.4\pr \\
3 & 12 & 252\dg 19.9\pr & 3 & 13 & 267\dg 22.3\pr & 3 & 14 & 282\dg 24.8\pr & 3 & 15 & 297\dg 27.3\pr \\
3 & 16 & 312\dg 29.7\pr & 3 & 17 & 327\dg 32.2\pr & 3 & 18 & 342\dg 34.7\pr & 3 & 19 & 357\dg 37.1\pr \\
3 & 20 & 12\dg 39.6\pr & 3 & 21 & 27\dg 42.0\pr & 3 & 22 & 42\dg 44.5\pr & 3 & 23 & 57\dg 47.0\pr \\
4 & 0 & 72\dg 49.4\pr & 4 & 1 & 87\dg 51.9\pr & 4 & 2 & 102\dg 54.4\pr & 4 & 3 & 117\dg 56.8\pr \\
4 & 4 & 132\dg 59.3\pr & 4 & 5 & 148\dg 01.8\pr & 4 & 6 & 163\dg 04.2\pr & 4 & 7 & 178\dg 06.7\pr \\
4 & 8 & 193\dg 09.2\pr & 4 & 9 & 208\dg 11.6\pr & 4 & 10 & 223\dg 14.1\pr & 4 & 11 & 238\dg 16.5\pr \\
4 & 12 & 253\dg 19.0\pr & 4 & 13 & 268\dg 21.5\pr & 4 & 14 & 283\dg 23.9\pr & 4 & 15 & 298\dg 26.4\pr \\
4 & 16 & 313\dg 28.9\pr & 4 & 17 & 328\dg 31.3\pr & 4 & 18 & 343\dg 33.8\pr & 4 & 19 & 358\dg 36.3\pr \\
4 & 20 & 13\dg 38.7\pr & 4 & 21 & 28\dg 41.2\pr & 4 & 22 & 43\dg 43.6\pr & 4 & 23 & 58\dg 46.1\pr \\
5 & 0 & 73\dg 48.6\pr & 5 & 1 & 88\dg 51.0\pr & 5 & 2 & 103\dg 53.5\pr & 5 & 3 & 118\dg 56.0\pr \\
5 & 4 & 133\dg 58.4\pr & 5 & 5 & 149\dg 00.9\pr & 5 & 6 & 164\dg 03.4\pr & 5 & 7 & 179\dg 05.8\pr \\
5 & 8 & 194\dg 08.3\pr & 5 & 9 & 209\dg 10.8\pr & 5 & 10 & 224\dg 13.2\pr & 5 & 11 & 239\dg 15.7\pr \\
5 & 12 & 254\dg 18.1\pr & 5 & 13 & 269\dg 20.6\pr & 5 & 14 & 284\dg 23.1\pr & 5 & 15 & 299\dg 25.5\pr \\
5 & 16 & 314\dg 28.0\pr & 5 & 17 & 329\dg 30.5\pr & 5 & 18 & 344\dg 32.9\pr & 5 & 19 & 359\dg 35.4\pr \\
5 & 20 & 14\dg 37.9\pr & 5 & 21 & 29\dg 40.3\pr & 5 & 22 & 44\dg 42.8\pr & 5 & 23 & 59\dg 45.3\pr \\
6 & 0 & 74\dg 47.7\pr & 6 & 1 & 89\dg 50.2\pr & 6 & 2 & 104\dg 52.6\pr & 6 & 3 & 119\dg 55.1\pr \\
6 & 4 & 134\dg 57.6\pr & 6 & 5 & 150\dg 00.0\pr & 6 & 6 & 165\dg 02.5\pr & 6 & 7 & 180\dg 05.0\pr \\
6 & 8 & 195\dg 07.4\pr & 6 & 9 & 210\dg 09.9\pr & 6 & 10 & 225\dg 12.4\pr & 6 & 11 & 240\dg 14.8\pr \\
6 & 12 & 255\dg 17.3\pr & 6 & 13 & 270\dg 19.8\pr & 6 & 14 & 285\dg 22.2\pr & 6 & 15 & 300\dg 24.7\pr \\
6 & 16 & 315\dg 27.1\pr & 6 & 17 & 330\dg 29.6\pr & 6 & 18 & 345\dg 32.1\pr & 6 & 19 & 0\dg 34.5\pr \\
6 & 20 & 15\dg 37.0\pr & 6 & 21 & 30\dg 39.5\pr & 6 & 22 & 45\dg 41.9\pr & 6 & 23 & 60\dg 44.4\pr \\
7 & 0 & 75\dg 46.9\pr & 7 & 1 & 90\dg 49.3\pr & 7 & 2 & 105\dg 51.8\pr & 7 & 3 & 120\dg 54.2\pr \\
7 & 4 & 135\dg 56.7\pr & 7 & 5 & 150\dg 59.2\pr & 7 & 6 & 166\dg 01.6\pr & 7 & 7 & 181\dg 04.1\pr \\
7 & 8 & 196\dg 06.6\pr & 7 & 9 & 211\dg 09.0\pr & 7 & 10 & 226\dg 11.5\pr & 7 & 11 & 241\dg 14.0\pr \\
7 & 12 & 256\dg 16.4\pr & 7 & 13 & 271\dg 18.9\pr & 7 & 14 & 286\dg 21.4\pr & 7 & 15 & 301\dg 23.8\pr \\
7 & 16 & 316\dg 26.3\pr & 7 & 17 & 331\dg 28.7\pr & 7 & 18 & 346\dg 31.2\pr & 7 & 19 & 1\dg 33.7\pr \\
7 & 20 & 16\dg 36.1\pr & 7 & 21 & 31\dg 38.6\pr & 7 & 22 & 46\dg 41.1\pr & 7 & 23 & 61\dg 43.5\pr \\
8 & 0 & 76\dg 46.0\pr & 8 & 1 & 91\dg 48.5\pr & 8 & 2 & 106\dg 50.9\pr & 8 & 3 & 121\dg 53.4\pr \\
8 & 4 & 136\dg 55.9\pr & 8 & 5 & 151\dg 58.3\pr & 8 & 6 & 167\dg 00.8\pr & 8 & 7 & 182\dg 03.2\pr \\
8 & 8 & 197\dg 05.7\pr & 8 & 9 & 212\dg 08.2\pr & 8 & 10 & 227\dg 10.6\pr & 8 & 11 & 242\dg 13.1\pr \\
8 & 12 & 257\dg 15.6\pr & 8 & 13 & 272\dg 18.0\pr & 8 & 14 & 287\dg 20.5\pr & 8 & 15 & 302\dg 23.0\pr \\
8 & 16 & 317\dg 25.4\pr & 8 & 17 & 332\dg 27.9\pr & 8 & 18 & 347\dg 30.3\pr & 8 & 19 & 2\dg 32.8\pr \\
8 & 20 & 17\dg 35.3\pr & 8 & 21 & 32\dg 37.7\pr & 8 & 22 & 47\dg 40.2\pr & 8 & 23 & 62\dg 42.7\pr \\
9 & 0 & 77\dg 45.1\pr & 9 & 1 & 92\dg 47.6\pr & 9 & 2 & 107\dg 50.1\pr & 9 & 3 & 122\dg 52.5\pr \\
9 & 4 & 137\dg 55.0\pr & 9 & 5 & 152\dg 57.5\pr & 9 & 6 & 167\dg 59.9\pr & 9 & 7 & 183\dg 02.4\pr \\
9 & 8 & 198\dg 04.8\pr & 9 & 9 & 213\dg 07.3\pr & 9 & 10 & 228\dg 09.8\pr & 9 & 11 & 243\dg 12.2\pr \\
9 & 12 & 258\dg 14.7\pr & 9 & 13 & 273\dg 17.2\pr & 9 & 14 & 288\dg 19.6\pr & 9 & 15 & 303\dg 22.1\pr \\
9 & 16 & 318\dg 24.6\pr & 9 & 17 & 333\dg 27.0\pr & 9 & 18 & 348\dg 29.5\pr & 9 & 19 & 3\dg 32.0\pr \\
9 & 20 & 18\dg 34.4\pr & 9 & 21 & 33\dg 36.9\pr & 9 & 22 & 48\dg 39.3\pr & 9 & 23 & 63\dg 41.8\pr \\
10 & 0 & 78\dg 44.3\pr & 10 & 1 & 93\dg 46.7\pr & 10 & 2 & 108\dg 49.2\pr & 10 & 3 & 123\dg 51.7\pr \\
10 & 4 & 138\dg 54.1\pr & 10 & 5 & 153\dg 56.6\pr & 10 & 6 & 168\dg 59.1\pr & 10 & 7 & 184\dg 01.5\pr \\
10 & 8 & 199\dg 04.0\pr & 10 & 9 & 214\dg 06.4\pr & 10 & 10 & 229\dg 08.9\pr & 10 & 11 & 244\dg 11.4\pr \\
10 & 12 & 259\dg 13.8\pr & 10 & 13 & 274\dg 16.3\pr & 10 & 14 & 289\dg 18.8\pr & 10 & 15 & 304\dg 21.2\pr \\
10 & 16 & 319\dg 23.7\pr & 10 & 17 & 334\dg 26.2\pr & 10 & 18 & 349\dg 28.6\pr & 10 & 19 & 4\dg 31.1\pr \\
10 & 20 & 19\dg 33.6\pr & 10 & 21 & 34\dg 36.0\pr & 10 & 22 & 49\dg 38.5\pr & 10 & 23 & 64\dg 40.9\pr \\
11 & 0 & 79\dg 43.4\pr & 11 & 1 & 94\dg 45.9\pr & 11 & 2 & 109\dg 48.3\pr & 11 & 3 & 124\dg 50.8\pr \\
11 & 4 & 139\dg 53.3\pr & 11 & 5 & 154\dg 55.7\pr & 11 & 6 & 169\dg 58.2\pr & 11 & 7 & 185\dg 00.7\pr \\
11 & 8 & 200\dg 03.1\pr & 11 & 9 & 215\dg 05.6\pr & 11 & 10 & 230\dg 08.1\pr & 11 & 11 & 245\dg 10.5\pr \\
11 & 12 & 260\dg 13.0\pr & 11 & 13 & 275\dg 15.4\pr & 11 & 14 & 290\dg 17.9\pr & 11 & 15 & 305\dg 20.4\pr \\
11 & 16 & 320\dg 22.8\pr & 11 & 17 & 335\dg 25.3\pr & 11 & 18 & 350\dg 27.8\pr & 11 & 19 & 5\dg 30.2\pr \\
11 & 20 & 20\dg 32.7\pr & 11 & 21 & 35\dg 35.2\pr & 11 & 22 & 50\dg 37.6\pr & 11 & 23 & 65\dg 40.1\pr \\
12 & 0 & 80\dg 42.5\pr & 12 & 1 & 95\dg 45.0\pr & 12 & 2 & 110\dg 47.5\pr & 12 & 3 & 125\dg 49.9\pr \\
12 & 4 & 140\dg 52.4\pr & 12 & 5 & 155\dg 54.9\pr & 12 & 6 & 170\dg 57.3\pr & 12 & 7 & 185\dg 59.8\pr \\
12 & 8 & 201\dg 02.3\pr & 12 & 9 & 216\dg 04.7\pr & 12 & 10 & 231\dg 07.2\pr & 12 & 11 & 246\dg 09.7\pr \\
12 & 12 & 261\dg 12.1\pr & 12 & 13 & 276\dg 14.6\pr & 12 & 14 & 291\dg 17.0\pr & 12 & 15 & 306\dg 19.5\pr \\
12 & 16 & 321\dg 22.0\pr & 12 & 17 & 336\dg 24.4\pr & 12 & 18 & 351\dg 26.9\pr & 12 & 19 & 6\dg 29.4\pr \\
12 & 20 & 21\dg 31.8\pr & 12 & 21 & 36\dg 34.3\pr & 12 & 22 & 51\dg 36.8\pr & 12 & 23 & 66\dg 39.2\pr \\
13 & 0 & 81\dg 41.7\pr & 13 & 1 & 96\dg 44.2\pr & 13 & 2 & 111\dg 46.6\pr & 13 & 3 & 126\dg 49.1\pr \\
13 & 4 & 141\dg 51.5\pr & 13 & 5 & 156\dg 54.0\pr & 13 & 6 & 171\dg 56.5\pr & 13 & 7 & 186\dg 58.9\pr \\
13 & 8 & 202\dg 01.4\pr & 13 & 9 & 217\dg 03.9\pr & 13 & 10 & 232\dg 06.3\pr & 13 & 11 & 247\dg 08.8\pr \\
13 & 12 & 262\dg 11.3\pr & 13 & 13 & 277\dg 13.7\pr & 13 & 14 & 292\dg 16.2\pr & 13 & 15 & 307\dg 18.7\pr \\
13 & 16 & 322\dg 21.1\pr & 13 & 17 & 337\dg 23.6\pr & 13 & 18 & 352\dg 26.0\pr & 13 & 19 & 7\dg 28.5\pr \\
13 & 20 & 22\dg 31.0\pr & 13 & 21 & 37\dg 33.4\pr & 13 & 22 & 52\dg 35.9\pr & 13 & 23 & 67\dg 38.4\pr \\
14 & 0 & 82\dg 40.8\pr & 14 & 1 & 97\dg 43.3\pr & 14 & 2 & 112\dg 45.8\pr & 14 & 3 & 127\dg 48.2\pr \\
14 & 4 & 142\dg 50.7\pr & 14 & 5 & 157\dg 53.1\pr & 14 & 6 & 172\dg 55.6\pr & 14 & 7 & 187\dg 58.1\pr \\
14 & 8 & 203\dg 00.5\pr & 14 & 9 & 218\dg 03.0\pr & 14 & 10 & 233\dg 05.5\pr & 14 & 11 & 248\dg 07.9\pr \\
14 & 12 & 263\dg 10.4\pr & 14 & 13 & 278\dg 12.9\pr & 14 & 14 & 293\dg 15.3\pr & 14 & 15 & 308\dg 17.8\pr \\
14 & 16 & 323\dg 20.3\pr & 14 & 17 & 338\dg 22.7\pr & 14 & 18 & 353\dg 25.2\pr & 14 & 19 & 8\dg 27.6\pr \\
14 & 20 & 23\dg 30.1\pr & 14 & 21 & 38\dg 32.6\pr & 14 & 22 & 53\dg 35.0\pr & 14 & 23 & 68\dg 37.5\pr \\
15 & 0 & 83\dg 40.0\pr & 15 & 1 & 98\dg 42.4\pr & 15 & 2 & 113\dg 44.9\pr & 15 & 3 & 128\dg 47.4\pr \\
15 & 4 & 143\dg 49.8\pr & 15 & 5 & 158\dg 52.3\pr & 15 & 6 & 173\dg 54.8\pr & 15 & 7 & 188\dg 57.2\pr \\
15 & 8 & 203\dg 59.7\pr & 15 & 9 & 219\dg 02.1\pr & 15 & 10 & 234\dg 04.6\pr & 15 & 11 & 249\dg 07.1\pr \\
15 & 12 & 264\dg 09.5\pr & 15 & 13 & 279\dg 12.0\pr & 15 & 14 & 294\dg 14.5\pr & 15 & 15 & 309\dg 16.9\pr \\
15 & 16 & 324\dg 19.4\pr & 15 & 17 & 339\dg 21.9\pr & 15 & 18 & 354\dg 24.3\pr & 15 & 19 & 9\dg 26.8\pr \\
15 & 20 & 24\dg 29.2\pr & 15 & 21 & 39\dg 31.7\pr & 15 & 22 & 54\dg 34.2\pr & 15 & 23 & 69\dg 36.6\pr \\
16 & 0 & 84\dg 39.1\pr & 16 & 1 & 99\dg 41.6\pr & 16 & 2 & 114\dg 44.0\pr & 16 & 3 & 129\dg 46.5\pr \\
16 & 4 & 144\dg 49.0\pr & 16 & 5 & 159\dg 51.4\pr & 16 & 6 & 174\dg 53.9\pr & 16 & 7 & 189\dg 56.4\pr \\
16 & 8 & 204\dg 58.8\pr & 16 & 9 & 220\dg 01.3\pr & 16 & 10 & 235\dg 03.7\pr & 16 & 11 & 250\dg 06.2\pr \\
16 & 12 & 265\dg 08.7\pr & 16 & 13 & 280\dg 11.1\pr & 16 & 14 & 295\dg 13.6\pr & 16 & 15 & 310\dg 16.1\pr \\
16 & 16 & 325\dg 18.5\pr & 16 & 17 & 340\dg 21.0\pr & 16 & 18 & 355\dg 23.5\pr & 16 & 19 & 10\dg 25.9\pr \\
16 & 20 & 25\dg 28.4\pr & 16 & 21 & 40\dg 30.9\pr & 16 & 22 & 55\dg 33.3\pr & 16 & 23 & 70\dg 35.8\pr \\
17 & 0 & 85\dg 38.2\pr & 17 & 1 & 100\dg 40.7\pr & 17 & 2 & 115\dg 43.2\pr & 17 & 3 & 130\dg 45.6\pr \\
17 & 4 & 145\dg 48.1\pr & 17 & 5 & 160\dg 50.6\pr & 17 & 6 & 175\dg 53.0\pr & 17 & 7 & 190\dg 55.5\pr \\
17 & 8 & 205\dg 58.0\pr & 17 & 9 & 221\dg 00.4\pr & 17 & 10 & 236\dg 02.9\pr & 17 & 11 & 251\dg 05.3\pr \\
17 & 12 & 266\dg 07.8\pr & 17 & 13 & 281\dg 10.3\pr & 17 & 14 & 296\dg 12.7\pr & 17 & 15 & 311\dg 15.2\pr \\
17 & 16 & 326\dg 17.7\pr & 17 & 17 & 341\dg 20.1\pr & 17 & 18 & 356\dg 22.6\pr & 17 & 19 & 11\dg 25.1\pr \\
17 & 20 & 26\dg 27.5\pr & 17 & 21 & 41\dg 30.0\pr & 17 & 22 & 56\dg 32.5\pr & 17 & 23 & 71\dg 34.9\pr \\
18 & 0 & 86\dg 37.4\pr & 18 & 1 & 101\dg 39.8\pr & 18 & 2 & 116\dg 42.3\pr & 18 & 3 & 131\dg 44.8\pr \\
18 & 4 & 146\dg 47.2\pr & 18 & 5 & 161\dg 49.7\pr & 18 & 6 & 176\dg 52.2\pr & 18 & 7 & 191\dg 54.6\pr \\
18 & 8 & 206\dg 57.1\pr & 18 & 9 & 221\dg 59.6\pr & 18 & 10 & 237\dg 02.0\pr & 18 & 11 & 252\dg 04.5\pr \\
18 & 12 & 267\dg 07.0\pr & 18 & 13 & 282\dg 09.4\pr & 18 & 14 & 297\dg 11.9\pr & 18 & 15 & 312\dg 14.3\pr \\
18 & 16 & 327\dg 16.8\pr & 18 & 17 & 342\dg 19.3\pr & 18 & 18 & 357\dg 21.7\pr & 18 & 19 & 12\dg 24.2\pr \\
18 & 20 & 27\dg 26.7\pr & 18 & 21 & 42\dg 29.1\pr & 18 & 22 & 57\dg 31.6\pr & 18 & 23 & 72\dg 34.1\pr \\
19 & 0 & 87\dg 36.5\pr & 19 & 1 & 102\dg 39.0\pr & 19 & 2 & 117\dg 41.4\pr & 19 & 3 & 132\dg 43.9\pr \\
19 & 4 & 147\dg 46.4\pr & 19 & 5 & 162\dg 48.8\pr & 19 & 6 & 177\dg 51.3\pr & 19 & 7 & 192\dg 53.8\pr \\
19 & 8 & 207\dg 56.2\pr & 19 & 9 & 222\dg 58.7\pr & 19 & 10 & 238\dg 01.2\pr & 19 & 11 & 253\dg 03.6\pr \\
19 & 12 & 268\dg 06.1\pr & 19 & 13 & 283\dg 08.6\pr & 19 & 14 & 298\dg 11.0\pr & 19 & 15 & 313\dg 13.5\pr \\
19 & 16 & 328\dg 15.9\pr & 19 & 17 & 343\dg 18.4\pr & 19 & 18 & 358\dg 20.9\pr & 19 & 19 & 13\dg 23.3\pr \\
19 & 20 & 28\dg 25.8\pr & 19 & 21 & 43\dg 28.3\pr & 19 & 22 & 58\dg 30.7\pr & 19 & 23 & 73\dg 33.2\pr \\
20 & 0 & 88\dg 35.7\pr & 20 & 1 & 103\dg 38.1\pr & 20 & 2 & 118\dg 40.6\pr & 20 & 3 & 133\dg 43.1\pr \\
20 & 4 & 148\dg 45.5\pr & 20 & 5 & 163\dg 48.0\pr & 20 & 6 & 178\dg 50.4\pr & 20 & 7 & 193\dg 52.9\pr \\
20 & 8 & 208\dg 55.4\pr & 20 & 9 & 223\dg 57.8\pr & 20 & 10 & 239\dg 00.3\pr & 20 & 11 & 254\dg 02.8\pr \\
20 & 12 & 269\dg 05.2\pr & 20 & 13 & 284\dg 07.7\pr & 20 & 14 & 299\dg 10.2\pr & 20 & 15 & 314\dg 12.6\pr \\
20 & 16 & 329\dg 15.1\pr & 20 & 17 & 344\dg 17.6\pr & 20 & 18 & 359\dg 20.0\pr & 20 & 19 & 14\dg 22.5\pr \\
20 & 20 & 29\dg 24.9\pr & 20 & 21 & 44\dg 27.4\pr & 20 & 22 & 59\dg 29.9\pr & 20 & 23 & 74\dg 32.3\pr \\
21 & 0 & 89\dg 34.8\pr & 21 & 1 & 104\dg 37.3\pr & 21 & 2 & 119\dg 39.7\pr & 21 & 3 & 134\dg 42.2\pr \\
21 & 4 & 149\dg 44.7\pr & 21 & 5 & 164\dg 47.1\pr & 21 & 6 & 179\dg 49.6\pr & 21 & 7 & 194\dg 52.0\pr \\
21 & 8 & 209\dg 54.5\pr & 21 & 9 & 224\dg 57.0\pr & 21 & 10 & 239\dg 59.4\pr & 21 & 11 & 255\dg 01.9\pr \\
21 & 12 & 270\dg 04.4\pr & 21 & 13 & 285\dg 06.8\pr & 21 & 14 & 300\dg 09.3\pr & 21 & 15 & 315\dg 11.8\pr \\
21 & 16 & 330\dg 14.2\pr & 21 & 17 & 345\dg 16.7\pr & 21 & 18 & 0\dg 19.2\pr & 21 & 19 & 15\dg 21.6\pr \\
21 & 20 & 30\dg 24.1\pr & 21 & 21 & 45\dg 26.5\pr & 21 & 22 & 60\dg 29.0\pr & 21 & 23 & 75\dg 31.5\pr \\
22 & 0 & 90\dg 33.9\pr & 22 & 1 & 105\dg 36.4\pr & 22 & 2 & 120\dg 38.9\pr & 22 & 3 & 135\dg 41.3\pr \\
22 & 4 & 150\dg 43.8\pr & 22 & 5 & 165\dg 46.3\pr & 22 & 6 & 180\dg 48.7\pr & 22 & 7 & 195\dg 51.2\pr \\
22 & 8 & 210\dg 53.7\pr & 22 & 9 & 225\dg 56.1\pr & 22 & 10 & 240\dg 58.6\pr & 22 & 11 & 256\dg 01.0\pr \\
22 & 12 & 271\dg 03.5\pr & 22 & 13 & 286\dg 06.0\pr & 22 & 14 & 301\dg 08.4\pr & 22 & 15 & 316\dg 10.9\pr \\
22 & 16 & 331\dg 13.4\pr & 22 & 17 & 346\dg 15.8\pr & 22 & 18 & 1\dg 18.3\pr & 22 & 19 & 16\dg 20.8\pr \\
22 & 20 & 31\dg 23.2\pr & 22 & 21 & 46\dg 25.7\pr & 22 & 22 & 61\dg 28.1\pr & 22 & 23 & 76\dg 30.6\pr \\
23 & 0 & 91\dg 33.1\pr & 23 & 1 & 106\dg 35.5\pr & 23 & 2 & 121\dg 38.0\pr & 23 & 3 & 136\dg 40.5\pr \\
23 & 4 & 151\dg 42.9\pr & 23 & 5 & 166\dg 45.4\pr & 23 & 6 & 181\dg 47.9\pr & 23 & 7 & 196\dg 50.3\pr \\
23 & 8 & 211\dg 52.8\pr & 23 & 9 & 226\dg 55.3\pr & 23 & 10 & 241\dg 57.7\pr & 23 & 11 & 257\dg 00.2\pr \\
23 & 12 & 272\dg 02.6\pr & 23 & 13 & 287\dg 05.1\pr & 23 & 14 & 302\dg 07.6\pr & 23 & 15 & 317\dg 10.0\pr \\
23 & 16 & 332\dg 12.5\pr & 23 & 17 & 347\dg 15.0\pr & 23 & 18 & 2\dg 17.4\pr & 23 & 19 & 17\dg 19.9\pr \\
23 & 20 & 32\dg 22.4\pr & 23 & 21 & 47\dg 24.8\pr & 23 & 22 & 62\dg 27.3\pr & 23 & 23 & 77\dg 29.8\pr \\
24 & 0 & 92\dg 32.2\pr & 24 & 1 & 107\dg 34.7\pr & 24 & 2 & 122\dg 37.1\pr & 24 & 3 & 137\dg 39.6\pr \\
24 & 4 & 152\dg 42.1\pr & 24 & 5 & 167\dg 44.5\pr & 24 & 6 & 182\dg 47.0\pr & 24 & 7 & 197\dg 49.5\pr \\
24 & 8 & 212\dg 51.9\pr & 24 & 9 & 227\dg 54.4\pr & 24 & 10 & 242\dg 56.9\pr & 24 & 11 & 257\dg 59.3\pr \\
24 & 12 & 273\dg 01.8\pr & 24 & 13 & 288\dg 04.2\pr & 24 & 14 & 303\dg 06.7\pr & 24 & 15 & 318\dg 09.2\pr \\
24 & 16 & 333\dg 11.6\pr & 24 & 17 & 348\dg 14.1\pr & 24 & 18 & 3\dg 16.6\pr & 24 & 19 & 18\dg 19.0\pr \\
24 & 20 & 33\dg 21.5\pr & 24 & 21 & 48\dg 24.0\pr & 24 & 22 & 63\dg 26.4\pr & 24 & 23 & 78\dg 28.9\pr \\
25 & 0 & 93\dg 31.4\pr & 25 & 1 & 108\dg 33.8\pr & 25 & 2 & 123\dg 36.3\pr & 25 & 3 & 138\dg 38.7\pr \\
25 & 4 & 153\dg 41.2\pr & 25 & 5 & 168\dg 43.7\pr & 25 & 6 & 183\dg 46.1\pr & 25 & 7 & 198\dg 48.6\pr \\
25 & 8 & 213\dg 51.1\pr & 25 & 9 & 228\dg 53.5\pr & 25 & 10 & 243\dg 56.0\pr & 25 & 11 & 258\dg 58.5\pr \\
25 & 12 & 274\dg 00.9\pr & 25 & 13 & 289\dg 03.4\pr & 25 & 14 & 304\dg 05.9\pr & 25 & 15 & 319\dg 08.3\pr \\
25 & 16 & 334\dg 10.8\pr & 25 & 17 & 349\dg 13.2\pr & 25 & 18 & 4\dg 15.7\pr & 25 & 19 & 19\dg 18.2\pr \\
25 & 20 & 34\dg 20.6\pr & 25 & 21 & 49\dg 23.1\pr & 25 & 22 & 64\dg 25.6\pr & 25 & 23 & 79\dg 28.0\pr \\
26 & 0 & 94\dg 30.5\pr & 26 & 1 & 109\dg 33.0\pr & 26 & 2 & 124\dg 35.4\pr & 26 & 3 & 139\dg 37.9\pr \\
26 & 4 & 154\dg 40.3\pr & 26 & 5 & 169\dg 42.8\pr & 26 & 6 & 184\dg 45.3\pr & 26 & 7 & 199\dg 47.7\pr \\
26 & 8 & 214\dg 50.2\pr & 26 & 9 & 229\dg 52.7\pr & 26 & 10 & 244\dg 55.1\pr & 26 & 11 & 259\dg 57.6\pr \\
26 & 12 & 275\dg 00.1\pr & 26 & 13 & 290\dg 02.5\pr & 26 & 14 & 305\dg 05.0\pr & 26 & 15 & 320\dg 07.5\pr \\
26 & 16 & 335\dg 09.9\pr & 26 & 17 & 350\dg 12.4\pr & 26 & 18 & 5\dg 14.8\pr & 26 & 19 & 20\dg 17.3\pr \\
26 & 20 & 35\dg 19.8\pr & 26 & 21 & 50\dg 22.2\pr & 26 & 22 & 65\dg 24.7\pr & 26 & 23 & 80\dg 27.2\pr \\
27 & 0 & 95\dg 29.6\pr & 27 & 1 & 110\dg 32.1\pr & 27 & 2 & 125\dg 34.6\pr & 27 & 3 & 140\dg 37.0\pr \\
27 & 4 & 155\dg 39.5\pr & 27 & 5 & 170\dg 42.0\pr & 27 & 6 & 185\dg 44.4\pr & 27 & 7 & 200\dg 46.9\pr \\
27 & 8 & 215\dg 49.3\pr & 27 & 9 & 230\dg 51.8\pr & 27 & 10 & 245\dg 54.3\pr & 27 & 11 & 260\dg 56.7\pr \\
27 & 12 & 275\dg 59.2\pr & 27 & 13 & 291\dg 01.7\pr & 27 & 14 & 306\dg 04.1\pr & 27 & 15 & 321\dg 06.6\pr \\
27 & 16 & 336\dg 09.1\pr & 27 & 17 & 351\dg 11.5\pr & 27 & 18 & 6\dg 14.0\pr & 27 & 19 & 21\dg 16.5\pr \\
27 & 20 & 36\dg 18.9\pr & 27 & 21 & 51\dg 21.4\pr & 27 & 22 & 66\dg 23.8\pr & 27 & 23 & 81\dg 26.3\pr \\
28 & 0 & 96\dg 28.8\pr & 28 & 1 & 111\dg 31.2\pr & 28 & 2 & 126\dg 33.7\pr & 28 & 3 & 141\dg 36.2\pr \\
28 & 4 & 156\dg 38.6\pr & 28 & 5 & 171\dg 41.1\pr & 28 & 6 & 186\dg 43.6\pr & 28 & 7 & 201\dg 46.0\pr \\
28 & 8 & 216\dg 48.5\pr & 28 & 9 & 231\dg 50.9\pr & 28 & 10 & 246\dg 53.4\pr & 28 & 11 & 261\dg 55.9\pr \\
28 & 12 & 276\dg 58.3\pr & 28 & 13 & 292\dg 00.8\pr & 28 & 14 & 307\dg 03.3\pr & 28 & 15 & 322\dg 05.7\pr \\
28 & 16 & 337\dg 08.2\pr & 28 & 17 & 352\dg 10.7\pr & 28 & 18 & 7\dg 13.1\pr & 28 & 19 & 22\dg 15.6\pr \\
28 & 20 & 37\dg 18.1\pr & 28 & 21 & 52\dg 20.5\pr & 28 & 22 & 67\dg 23.0\pr & 28 & 23 & 82\dg 25.4\pr \\
29 & 0 & 97\dg 27.9\pr & 29 & 1 & 112\dg 30.4\pr & 29 & 2 & 127\dg 32.8\pr & 29 & 3 & 142\dg 35.3\pr \\
29 & 4 & 157\dg 37.8\pr & 29 & 5 & 172\dg 40.2\pr & 29 & 6 & 187\dg 42.7\pr & 29 & 7 & 202\dg 45.2\pr \\
29 & 8 & 217\dg 47.6\pr & 29 & 9 & 232\dg 50.1\pr & 29 & 10 & 247\dg 52.6\pr & 29 & 11 & 262\dg 55.0\pr \\
29 & 12 & 277\dg 57.5\pr & 29 & 13 & 292\dg 59.9\pr & 29 & 14 & 308\dg 02.4\pr & 29 & 15 & 323\dg 04.9\pr \\
29 & 16 & 338\dg 07.3\pr & 29 & 17 & 353\dg 09.8\pr & 29 & 18 & 8\dg 12.3\pr & 29 & 19 & 23\dg 14.7\pr \\
29 & 20 & 38\dg 17.2\pr & 29 & 21 & 53\dg 19.7\pr & 29 & 22 & 68\dg 22.1\pr & 29 & 23 & 83\dg 24.6\pr \\
30 & 0 & 98\dg 27.0\pr & 30 & 1 & 113\dg 29.5\pr & 30 & 2 & 128\dg 32.0\pr & 30 & 3 & 143\dg 34.4\pr \\
30 & 4 & 158\dg 36.9\pr & 30 & 5 & 173\dg 39.4\pr & 30 & 6 & 188\dg 41.8\pr & 30 & 7 & 203\dg 44.3\pr \\
30 & 8 & 218\dg 46.8\pr & 30 & 9 & 233\dg 49.2\pr & 30 & 10 & 248\dg 51.7\pr & 30 & 11 & 263\dg 54.2\pr \\
30 & 12 & 278\dg 56.6\pr & 30 & 13 & 293\dg 59.1\pr & 30 & 14 & 309\dg 01.5\pr & 30 & 15 & 324\dg 04.0\pr \\
30 & 16 & 339\dg 06.5\pr & 30 & 17 & 354\dg 08.9\pr & 30 & 18 & 9\dg 11.4\pr & 30 & 19 & 24\dg 13.9\pr \\
30 & 20 & 39\dg 16.3\pr & 30 & 21 & 54\dg 18.8\pr & 30 & 22 & 69\dg 21.3\pr & 30 & 23 & 84\dg 23.7\pr \\
31 & 0 & 99\dg 26.2\pr & 31 & 1 & 114\dg 28.7\pr & 31 & 2 & 129\dg 31.1\pr & 31 & 3 & 144\dg 33.6\pr \\
31 & 4 & 159\dg 36.0\pr & 31 & 5 & 174\dg 38.5\pr & 31 & 6 & 189\dg 41.0\pr & 31 & 7 & 204\dg 43.4\pr \\
31 & 8 & 219\dg 45.9\pr & 31 & 9 & 234\dg 48.4\pr & 31 & 10 & 249\dg 50.8\pr & 31 & 11 & 264\dg 53.3\pr \\
31 & 12 & 279\dg 55.8\pr & 31 & 13 & 294\dg 58.2\pr & 31 & 14 & 310\dg 00.7\pr & 31 & 15 & 325\dg 03.1\pr \\
31 & 16 & 340\dg 05.6\pr & 31 & 17 & 355\dg 08.1\pr & 31 & 18 & 10\dg 10.5\pr & 31 & 19 & 25\dg 13.0\pr \\
31 & 20 & 40\dg 15.5\pr & 31 & 21 & 55\dg 17.9\pr & 31 & 22 & 70\dg 20.4\pr & 31 & 23 & 85\dg 22.9\pr \\
\end{longtable}}


% ============================================================
% PART IV --- CORRECTIONS & INCREMENTS
% ============================================================
\part{Correction \& Increment Tables}
\chapter{Altitude corrections}

\section{Dip of the horizon}
Subtract from the sextant altitude. $\text{Dip}=1.76\sqrt{h}\pr$, $h$ in metres.
\begin{center}\begin{tabular}{r r || r r}
\toprule
\textbf{Height (m)} & \textbf{Dip ($\pr$)} & \textbf{Height (m)} & \textbf{Dip ($\pr$)}\\
\midrule
1.0 & $-$1.8 & 7.0 & $-$4.7 \\
1.5 & $-$2.2 & 8.0 & $-$5.0 \\
2.0 & $-$2.5 & 9.0 & $-$5.3 \\
2.5 & $-$2.8 & 10.0 & $-$5.6 \\
3.0 & $-$3.0 & 11.0 & $-$5.8 \\
3.5 & $-$3.3 & 12.0 & $-$6.1 \\
4.0 & $-$3.5 & 15.0 & $-$6.8 \\
4.5 & $-$3.7 & 20.0 & $-$7.9 \\
5.0 & $-$3.9 & 25.0 & $-$8.8 \\
5.5 & $-$4.1 & 30.0 & $-$9.6 \\
6.0 & $-$4.3 &  &  \\
\bottomrule
\end{tabular}
\end{center}

\section{Atmospheric refraction}
Subtract from the apparent altitude $H_a$. Bennett's formula; standard atmosphere.
\begin{center}\begin{tabular}{r r || r r}
\toprule
\textbf{App.Alt ($\dg$)} & \textbf{Refr ($\pr$)} & \textbf{App.Alt ($\dg$)} & \textbf{Refr ($\pr$)}\\
\midrule
5 & $-$9.9 & 18 & $-$3.0 \\
6 & $-$8.5 & 19 & $-$2.9 \\
7 & $-$7.5 & 20 & $-$2.7 \\
8 & $-$6.6 & 22 & $-$2.4 \\
9 & $-$5.9 & 25 & $-$2.1 \\
10 & $-$5.4 & 30 & $-$1.7 \\
11 & $-$4.9 & 35 & $-$1.4 \\
12 & $-$4.5 & 40 & $-$1.2 \\
13 & $-$4.2 & 50 & $-$0.8 \\
14 & $-$3.9 & 60 & $-$0.6 \\
15 & $-$3.6 & 70 & $-$0.4 \\
16 & $-$3.4 & 80 & $-$0.2 \\
17 & $-$3.2 & 90 & 0.0 \\
\bottomrule
\end{tabular}
\end{center}

\section{Sun semidiameter}
Add for a lower-limb sight, subtract for an upper-limb sight.
\begin{center}\begin{tabular}{l r || l r}
\toprule
\textbf{Date} & \textbf{SD ($\pr$)} & \textbf{Date} & \textbf{SD ($\pr$)}\\
\midrule
Jan 1 & 16.3 & Jul 1 & 15.7 \\
Jan 15 & 16.3 & Jul 15 & 15.7 \\
Feb 1 & 16.2 & Aug 1 & 15.8 \\
Feb 15 & 16.2 & Aug 15 & 15.8 \\
Mar 1 & 16.1 & Sep 1 & 15.9 \\
Mar 15 & 16.1 & Sep 15 & 15.9 \\
Apr 1 & 16.0 & Oct 1 & 16.0 \\
Apr 15 & 15.9 & Oct 15 & 16.0 \\
May 1 & 15.9 & Nov 1 & 16.1 \\
May 15 & 15.8 & Nov 15 & 16.2 \\
Jun 1 & 15.8 & Dec 1 & 16.2 \\
Jun 15 & 15.8 & Dec 15 & 16.3 \\
\bottomrule
\end{tabular}
\end{center}

\chapter{GHA increments}

To find the GHA between whole-hour entries, add the increment for the elapsed
minutes. The Sun's GHA increases at $15\dg$ per hour; Aries at $15.041\dg$ per
hour. For the seconds, add $0.25\pr$ per second for the Sun and $0.2507\pr$ per
second for Aries.
\begin{center}{\footnotesize
\begin{tabular}{r r r || r r r}
\toprule
\textbf{m} & \textbf{Sun} & \textbf{Aries} & \textbf{m} & \textbf{Sun} & \textbf{Aries}\\
\midrule
0 & 0\dg 00.0\pr & 0\dg 00.0\pr & 30 & 7\dg 30.0\pr & 7\dg 31.2\pr \\
1 & 0\dg 15.0\pr & 0\dg 15.0\pr & 31 & 7\dg 45.0\pr & 7\dg 46.3\pr \\
2 & 0\dg 30.0\pr & 0\dg 30.1\pr & 32 & 8\dg 00.0\pr & 8\dg 01.3\pr \\
3 & 0\dg 45.0\pr & 0\dg 45.1\pr & 33 & 8\dg 15.0\pr & 8\dg 16.4\pr \\
4 & 1\dg 00.0\pr & 1\dg 00.2\pr & 34 & 8\dg 30.0\pr & 8\dg 31.4\pr \\
5 & 1\dg 15.0\pr & 1\dg 15.2\pr & 35 & 8\dg 45.0\pr & 8\dg 46.4\pr \\
6 & 1\dg 30.0\pr & 1\dg 30.2\pr & 36 & 9\dg 00.0\pr & 9\dg 01.5\pr \\
7 & 1\dg 45.0\pr & 1\dg 45.3\pr & 37 & 9\dg 15.0\pr & 9\dg 16.5\pr \\
8 & 2\dg 00.0\pr & 2\dg 00.3\pr & 38 & 9\dg 30.0\pr & 9\dg 31.6\pr \\
9 & 2\dg 15.0\pr & 2\dg 15.4\pr & 39 & 9\dg 45.0\pr & 9\dg 46.6\pr \\
10 & 2\dg 30.0\pr & 2\dg 30.4\pr & 40 & 10\dg 00.0\pr & 10\dg 01.6\pr \\
11 & 2\dg 45.0\pr & 2\dg 45.5\pr & 41 & 10\dg 15.0\pr & 10\dg 16.7\pr \\
12 & 3\dg 00.0\pr & 3\dg 00.5\pr & 42 & 10\dg 30.0\pr & 10\dg 31.7\pr \\
13 & 3\dg 15.0\pr & 3\dg 15.5\pr & 43 & 10\dg 45.0\pr & 10\dg 46.8\pr \\
14 & 3\dg 30.0\pr & 3\dg 30.6\pr & 44 & 11\dg 00.0\pr & 11\dg 01.8\pr \\
15 & 3\dg 45.0\pr & 3\dg 45.6\pr & 45 & 11\dg 15.0\pr & 11\dg 16.8\pr \\
16 & 4\dg 00.0\pr & 4\dg 00.7\pr & 46 & 11\dg 30.0\pr & 11\dg 31.9\pr \\
17 & 4\dg 15.0\pr & 4\dg 15.7\pr & 47 & 11\dg 45.0\pr & 11\dg 46.9\pr \\
18 & 4\dg 30.0\pr & 4\dg 30.7\pr & 48 & 12\dg 00.0\pr & 12\dg 02.0\pr \\
19 & 4\dg 45.0\pr & 4\dg 45.8\pr & 49 & 12\dg 15.0\pr & 12\dg 17.0\pr \\
20 & 5\dg 00.0\pr & 5\dg 00.8\pr & 50 & 12\dg 30.0\pr & 12\dg 32.1\pr \\
21 & 5\dg 15.0\pr & 5\dg 15.9\pr & 51 & 12\dg 45.0\pr & 12\dg 47.1\pr \\
22 & 5\dg 30.0\pr & 5\dg 30.9\pr & 52 & 13\dg 00.0\pr & 13\dg 02.1\pr \\
23 & 5\dg 45.0\pr & 5\dg 45.9\pr & 53 & 13\dg 15.0\pr & 13\dg 17.2\pr \\
24 & 6\dg 00.0\pr & 6\dg 01.0\pr & 54 & 13\dg 30.0\pr & 13\dg 32.2\pr \\
25 & 6\dg 15.0\pr & 6\dg 16.0\pr & 55 & 13\dg 45.0\pr & 13\dg 47.3\pr \\
26 & 6\dg 30.0\pr & 6\dg 31.1\pr & 56 & 14\dg 00.0\pr & 14\dg 02.3\pr \\
27 & 6\dg 45.0\pr & 6\dg 46.1\pr & 57 & 14\dg 15.0\pr & 14\dg 17.3\pr \\
28 & 7\dg 00.0\pr & 7\dg 01.1\pr & 58 & 14\dg 30.0\pr & 14\dg 32.4\pr \\
29 & 7\dg 15.0\pr & 7\dg 16.2\pr & 59 & 14\dg 45.0\pr & 14\dg 47.4\pr \\
\bottomrule
\end{tabular}}
\end{center}

\backmatter
\chapter{Provenance and verification}
\sloppy
Every figure in this book was generated by
\texttt{almanac/}\allowbreak\texttt{generate\_tables.wls}
from the package \texttt{wolfram/}\allowbreak\texttt{CelestialNavigation.wl} in the repository
\url{https://github.com/mthiel74/SextantNavigation}, whose 186-test suite checks
the ephemeris against JPL Horizons and Meeus. Independent sanity checks performed
at generation time confirm: the Sun's declination reaches $+23.44\dg$ at the June
solstice and $-23.44\dg$ at the December solstice; the Sun's GHA advances
$15.00\dg$ per hour and Aries $15.041\dg$ per hour; a bright star's SHA is stable
to under $1\pr$ across the year; and the three worked examples close on the
engine's own geometry to $0.0\pr$. Re-generate with
\texttt{wolframscript -file}\allowbreak\ \texttt{almanac/}\allowbreak\texttt{generate\_tables.wls}
and rebuild with
\texttt{latexmk -pdf}\allowbreak\ \texttt{almanac/}\allowbreak\texttt{SextantAlmanac2026.tex}.

\medskip
\noindent\textbf{Remember the disclaimer:} educational use only, approximately one
arcminute accuracy, not for sole-means marine navigation.

\end{document}
